\documentclass[11pt,a4paper]{article}
\usepackage[T1]{fontenc}
\usepackage[utf8]{inputenc}
\usepackage[margin=1in]{geometry}
\usepackage{amsmath,amssymb,booktabs,array}
\usepackage{xcolor}
\usepackage{tikz}
\usepackage{graphicx}
\usepackage{hyperref}
\hypersetup{colorlinks=true,linkcolor=blue!45!black,citecolor=blue!45!black,urlcolor=blue!45!black}
\setlength{\emergencystretch}{3em}
\renewcommand{\topfraction}{0.92}
\renewcommand{\textfraction}{0.05}
\renewcommand{\floatpagefraction}{0.75}
\newcommand{\E}{\mathbb{E}}
\newcommand{\R}{\mathbb{R}}
\newcommand{\Var}{\operatorname{Var}}
\newcommand{\Cov}{\operatorname{Cov}}
\newcommand{\code}[1]{\texttt{#1}}

% BEGIN GENERATED MACROS
\IfFileExists{generated-macros.tex}{% Generated from frozen supplement records; do not edit.
\newcommand{\RepairDeltaMean}{0.513}
\newcommand{\RepairDeltaLow}{0.212}
\newcommand{\RepairDeltaHigh}{0.814}
\newcommand{\RepairImproved}{21}
\newcommand{\CasesPerCondition}{64}
\newcommand{\CasesPerRow}{384}
\newcommand{\StudyPython}{3.13.2}
\newcommand{\StudyNumPy}{2.5.4}
\newcommand{\StudyPandas}{3.1.0}
\newcommand{\StudySciPy}{1.18.1}
\newcommand{\StudySklearn}{1.9.1}
\newcommand{\StudySklearnFits}{8,448}
\newcommand{\StudyNumPySolves}{384}
\newcommand{\StudyCovarianceChecks}{768}
\newcommand{\StudyCalibrationEnsembles}{152,832}
\newcommand{\CleanRawAnyLow}{14}
\newcommand{\CleanRawAnyHigh}{33}
\newcommand{\CleanHolmAnyLow}{2}
\newcommand{\CleanHolmAnyHigh}{13}
\newcommand{\SpreadRawAnyLow}{92}
\newcommand{\SpreadRawAnyHigh}{100}
\newcommand{\SpreadHolmAnyLow}{35}
\newcommand{\SpreadHolmAnyHigh}{59}
% Exploratory follow-up chosen after the Study 2 results were known (protocol 2.1,
% experiments/results-v2-followup); do not edit.
\newcommand{\AnswerKeySubtleAlarms}{148}
\newcommand{\LossSubtleAlarms}{36}
\newcommand{\SubtleCases}{640}
\newcommand{\AnswerKeyCleanAlarms}{12}
\newcommand{\LossCleanAlarms}{13}
\newcommand{\CleanNonSeasonalCases}{320}
\newcommand{\OwnLagSumFirst}{0.95}
\newcommand{\OwnLagSumSecond}{0.91}
\newcommand{\ReferenceDeltaMean}{0.354}
\newcommand{\ReferenceDeltaLow}{0.242}
\newcommand{\ReferenceDeltaHigh}{0.467}
\newcommand{\FaultyCloserCases}{53}
\newcommand{\ScorecardNoiseSD}{3}
\newcommand{\ScorecardMaxSD}{19}
\newcommand{\ScorecardNoiseScore}{84}
\newcommand{\ScorecardNoiseLoss}{16}
\newcommand{\TrendARPhiHigh}{0.77}
% Second exploratory record, chosen after the Study 2 results and the follow-up were known
% (protocol 2.2, experiments/results-v2-addendum); do not edit.
\newcommand{\FisherSpreadSuites}{63}
\newcommand{\HolmSpreadSuites}{30}
\newcommand{\FisherCleanSuites}{5}
\newcommand{\HolmCleanSuites}{3}
\newcommand{\HochbergSpreadSuites}{35}
\newcommand{\SimesSpreadSuites}{41}
\newcommand{\HochbergCleanSuites}{3}
\newcommand{\SimesCleanSuites}{3}
\newcommand{\SuitesPerVariant}{64}
\newcommand{\FisherCleanLow}{3}
\newcommand{\FisherCleanHigh}{17}
\newcommand{\SuiteAlphaPercent}{5}
\newcommand{\PairedFlagged}{21}
\newcommand{\PairedPairs}{24}
\newcommand{\PairedFlips}{9,999}
}{%
\PackageError{timegym}{Missing generated-macros.tex}{Run experiments.render_supplement before compiling.}}
% END GENERATED MACROS

\title{TimeGym: What Can Your Forecaster Learn,\newline and What Is Broken?}
\author{Diogo Seca}
\date{Working manuscript; versioned identity in the standalone build}

\begin{document}
\maketitle

\begin{abstract}
A forecast can score well and still be wrong: on the wrong dates, for the
wrong target, or produced by the wrong algorithm. Its error alone cannot
tell a software bug from a pattern the model cannot learn or from noise
nobody could predict. TimeGym is a gym for forecasting pipelines: a Python
package that tests them on synthetic series whose generating rules are
known. Because the rules are known, every test case comes with an answer
key, the best forecast anyone could make from the history, so a
pipeline's error splits into what it missed and what was unpredictable. A
scorecard turns that split into a score from 0 to 100 on ten kinds of
series at low and high noise. TimeGym also checks that forecasts carry the
requested dates, targets and shape, and that they move correctly when the
history is shifted. Two controlled studies test these checks. In the
first, every deliberately corrupted output was caught. In the second, four
bugs placed inside a fitted lag-Ridge pipeline all produced well-formed
forecasts: a common shift caught one, per-target shifts caught two, and
only a second, independent implementation caught all four in every
condition. Fixing one of these bugs made forecasts worse on a random walk
with drift, so a correct implementation and an accurate forecast are
different goals. Shuffling sample paths left every per-step distribution,
and so the CRPS, unchanged while breaking the joint paths; joint scores
and a covariance check detected it. The studies use selected bugs on
Gaussian scenarios; they do not measure how often such checks catch the
bugs found in practice.
\end{abstract}

\paragraph{Keywords.} Time series; forecasting; software testing; conditional
references; probabilistic verification; mutation testing; reproducibility.

\begin{figure}[p]
\centering
% BEGIN VECTOR FIGURE paper-scorecard
% Generated vector scene; SVG and PNG use the same plotted geometry.
\begingroup
\definecolor{tg000000}{HTML}{000000}
\definecolor{tg0000ff}{HTML}{0000ff}
\definecolor{tg0b0b0b}{HTML}{0b0b0b}
\definecolor{tg18314b}{HTML}{18314b}
\definecolor{tg626b73}{HTML}{626b73}
\definecolor{tg808080}{HTML}{808080}
\definecolor{tg8fd18f}{HTML}{8fd18f}
\definecolor{tg9da7b0}{HTML}{9da7b0}
\definecolor{tgc9cdd1}{HTML}{c9cdd1}
\definecolor{tgd03b3b}{HTML}{d03b3b}
\definecolor{tgeceef0}{HTML}{eceef0}
\definecolor{tgfab219}{HTML}{fab219}
\definecolor{tgffffff}{HTML}{ffffff}
\resizebox{\linewidth}{!}{%
\begin{tikzpicture}[x=1bp,y=-1bp]
\path[use as bounding box] (0,0) rectangle (448,590);
\node[anchor=west,inner sep=0pt,outer sep=0pt,text=tg18314b,font={\sffamily\fontsize{10.2bp}{12.24bp}\selectfont\bfseries}] at (0.000,7.000) {StandardScaler \ensuremath{\rightarrow} Ridge (recursive, window 24): 69/100};
\node[anchor=west,inner sep=0pt,outer sep=0pt,text=tg626b73,font={\sffamily\fontsize{6.8bp}{8.16bp}\selectfont}] at (0.000,20.000) {TimeGym scorecard \ensuremath{\cdot} 10 kinds of series \ensuremath{\times} 2 noise levels \ensuremath{\cdot} 10 instances each \ensuremath{\cdot} seed 0};
\node[anchor=west,inner sep=0pt,outer sep=0pt,text=tg18314b,font={\sffamily\fontsize{7bp}{8.4bp}\selectfont}] at (0.000,31.200) {Good at: Pure noise, Noise around a fixed level, Straight-line trend, Repeating seasonal cycle, Trend + seasonal cycle,};
\node[anchor=west,inner sep=0pt,outer sep=0pt,text=tg18314b,font={\sffamily\fontsize{7bp}{8.4bp}\selectfont}] at (0.000,39.600) {Mean reversion (AR(1)): shocks fade, Irregular cycles (AR(2)).};
\node[anchor=west,inner sep=0pt,outer sep=0pt,text=tg18314b,font={\sffamily\fontsize{7bp}{8.4bp}\selectfont}] at (0.000,48.000) {Fails at: Sudden level shift, Random walk with drift.};
\node[anchor=west,inner sep=0pt,outer sep=0pt,text=tg626b73,font={\sffamily\fontsize{6.2bp}{7.44bp}\selectfont}] at (0.000,57.420) {Output checks:  \ensuremath{\checkmark} Runs without errors   \ensuremath{\checkmark} Forecasts the requested dates   \ensuremath{\checkmark} Returns every series, in order   \ensuremath{\checkmark} Right number of steps and series};
\node[anchor=west,inner sep=0pt,outer sep=0pt,text=tg626b73,font={\sffamily\fontsize{6.2bp}{7.44bp}\selectfont}] at (0.000,64.860) {\ensuremath{\checkmark} No missing or infinite values};
\node[anchor=west,inner sep=0pt,outer sep=0pt,text=tg18314b,font={\sffamily\fontsize{7bp}{8.4bp}\selectfont\bfseries}] at (0.000,77.580) {What can and can't be predicted};
\node[anchor=west,inner sep=0pt,outer sep=0pt,text=tg18314b,font={\sffamily\fontsize{7bp}{8.4bp}\selectfont\bfseries}] at (168.000,77.580) {Low noise (sd 0.5)};
\node[anchor=west,inner sep=0pt,outer sep=0pt,text=tg18314b,font={\sffamily\fontsize{7bp}{8.4bp}\selectfont\bfseries}] at (255.000,77.580) {High noise (sd 2)};
\node[anchor=west,inner sep=0pt,outer sep=0pt,text=tg18314b,font={\sffamily\fontsize{7bp}{8.4bp}\selectfont\bfseries}] at (346.000,77.580) {Mistakes};
\draw[color=tgc9cdd1,line width=0.45bp,line cap=round] (0.000,83.580) -- (448.000,83.580);
\node[anchor=west,inner sep=0pt,outer sep=0pt,text=tg18314b,font={\sffamily\fontsize{7.2bp}{8.64bp}\selectfont\bfseries}] at (0.000,94.400) {Pure noise};
\node[anchor=west,inner sep=0pt,outer sep=0pt,text=tg18314b,font={\sffamily\fontsize{6.2bp}{7.44bp}\selectfont}] at (0.000,102.840) {\ensuremath{\checkmark} Only the average, zero: the best forecast is a flat};
\node[anchor=west,inner sep=0pt,outer sep=0pt,text=tg18314b,font={\sffamily\fontsize{6.2bp}{7.44bp}\selectfont}] at (0.000,110.280) {line at zero.};
\node[anchor=west,inner sep=0pt,outer sep=0pt,text=tg18314b,font={\sffamily\fontsize{6.2bp}{7.44bp}\selectfont}] at (0.000,118.120) {\ensuremath{\times} Every individual value. Any pattern a model finds};
\node[anchor=west,inner sep=0pt,outer sep=0pt,text=tg18314b,font={\sffamily\fontsize{6.2bp}{7.44bp}\selectfont}] at (0.000,125.560) {here is invented.};
\path[fill=tg8fd18f,draw=none,line width=0.5bp] (168.000,90.080) rectangle (205.969,98.680);
\node[anchor=center,inner sep=0pt,outer sep=0pt,text=tg0b0b0b,font={\sffamily\fontsize{6.6bp}{7.92bp}\selectfont\bfseries}] at (186.984,94.380) {84  Good};
\draw[color=tgc9cdd1,line width=0.45bp,line cap=round] (168.000,132.466) -- (248.000,132.466);
\draw[color=tg9da7b0,line width=0.5bp,line cap=round,dash pattern=on 2bp off 2bp] (232.000,101.080) -- (232.000,132.466);
\path[fill=tg808080,draw=none] (168.333,109.812) circle[radius=0.42bp] (169.000,115.006) circle[radius=0.42bp] (169.667,120.379) circle[radius=0.42bp] (170.333,121.523) circle[radius=0.42bp] (171.000,118.546) circle[radius=0.42bp] (171.667,114.196) circle[radius=0.42bp] (172.333,122.087) circle[radius=0.42bp] (173.000,116.951) circle[radius=0.42bp] (173.667,116.630) circle[radius=0.42bp] (174.333,112.139) circle[radius=0.42bp] (175.000,114.232) circle[radius=0.42bp] (175.667,113.329) circle[radius=0.42bp] (176.333,119.804) circle[radius=0.42bp] (177.000,116.440) circle[radius=0.42bp] (177.667,110.579) circle[radius=0.42bp] (178.333,106.027) circle[radius=0.42bp] (179.000,123.687) circle[radius=0.42bp] (179.667,105.895) circle[radius=0.42bp] (180.333,106.974) circle[radius=0.42bp] (181.000,110.596) circle[radius=0.42bp] (181.667,113.912) circle[radius=0.42bp] (182.333,117.623) circle[radius=0.42bp] (183.000,106.254) circle[radius=0.42bp] (183.667,103.032) circle[radius=0.42bp] (184.333,104.049) circle[radius=0.42bp] (185.000,107.171) circle[radius=0.42bp] (185.667,113.316) circle[radius=0.42bp] (186.333,123.361) circle[radius=0.42bp] (187.000,115.637) circle[radius=0.42bp] (187.667,111.397) circle[radius=0.42bp] (188.333,123.875) circle[radius=0.42bp] (189.000,113.074) circle[radius=0.42bp] (189.667,112.851) circle[radius=0.42bp] (190.333,111.143) circle[radius=0.42bp] (191.000,123.206) circle[radius=0.42bp] (191.667,119.854) circle[radius=0.42bp] (192.333,118.409) circle[radius=0.42bp] (193.000,123.114) circle[radius=0.42bp] (193.667,104.449) circle[radius=0.42bp] (194.333,118.791) circle[radius=0.42bp] (195.000,113.498) circle[radius=0.42bp] (195.667,117.268) circle[radius=0.42bp] (196.333,105.449) circle[radius=0.42bp] (197.000,107.137) circle[radius=0.42bp] (197.667,111.545) circle[radius=0.42bp] (198.333,129.747) circle[radius=0.42bp] (199.000,115.275) circle[radius=0.42bp] (199.667,111.222) circle[radius=0.42bp] (200.333,109.167) circle[radius=0.42bp] (201.000,119.573) circle[radius=0.42bp] (201.667,103.919) circle[radius=0.42bp] (202.333,124.081) circle[radius=0.42bp] (203.000,119.853) circle[radius=0.42bp] (203.667,109.609) circle[radius=0.42bp] (204.333,115.294) circle[radius=0.42bp] (205.000,102.761) circle[radius=0.42bp] (205.667,114.399) circle[radius=0.42bp] (206.333,119.671) circle[radius=0.42bp] (207.000,118.031) circle[radius=0.42bp] (207.667,122.610) circle[radius=0.42bp] (208.333,123.806) circle[radius=0.42bp] (209.000,111.564) circle[radius=0.42bp] (209.667,111.880) circle[radius=0.42bp] (210.333,107.303) circle[radius=0.42bp] (211.000,120.450) circle[radius=0.42bp] (211.667,104.771) circle[radius=0.42bp] (212.333,117.453) circle[radius=0.42bp] (213.000,105.507) circle[radius=0.42bp] (213.667,118.386) circle[radius=0.42bp] (214.333,120.328) circle[radius=0.42bp] (215.000,114.006) circle[radius=0.42bp] (215.667,108.991) circle[radius=0.42bp] (216.333,114.576) circle[radius=0.42bp] (217.000,119.366) circle[radius=0.42bp] (217.667,124.214) circle[radius=0.42bp] (218.333,124.601) circle[radius=0.42bp] (219.000,112.384) circle[radius=0.42bp] (219.667,109.259) circle[radius=0.42bp] (220.333,116.663) circle[radius=0.42bp] (221.000,122.502) circle[radius=0.42bp] (221.667,110.007) circle[radius=0.42bp] (222.333,123.824) circle[radius=0.42bp] (223.000,120.184) circle[radius=0.42bp] (223.667,111.624) circle[radius=0.42bp] (224.333,130.046) circle[radius=0.42bp] (225.000,113.130) circle[radius=0.42bp] (225.667,119.341) circle[radius=0.42bp] (226.333,114.908) circle[radius=0.42bp] (227.000,116.094) circle[radius=0.42bp] (227.667,114.312) circle[radius=0.42bp] (228.333,111.155) circle[radius=0.42bp] (229.000,120.474) circle[radius=0.42bp] (229.667,106.492) circle[radius=0.42bp] (230.333,110.950) circle[radius=0.42bp] (231.000,110.195) circle[radius=0.42bp] (231.667,108.135) circle[radius=0.42bp] (232.333,110.556) circle[radius=0.42bp] (233.000,110.193) circle[radius=0.42bp] (233.667,115.124) circle[radius=0.42bp] (234.333,124.763) circle[radius=0.42bp] (235.000,116.475) circle[radius=0.42bp] (235.667,120.546) circle[radius=0.42bp] (236.333,124.737) circle[radius=0.42bp] (237.000,113.951) circle[radius=0.42bp] (237.667,119.257) circle[radius=0.42bp] (238.333,122.216) circle[radius=0.42bp] (239.000,122.301) circle[radius=0.42bp] (239.667,113.887) circle[radius=0.42bp] (240.333,113.308) circle[radius=0.42bp] (241.000,107.124) circle[radius=0.42bp] (241.667,115.698) circle[radius=0.42bp] (242.333,108.924) circle[radius=0.42bp] (243.000,106.612) circle[radius=0.42bp] (243.667,108.229) circle[radius=0.42bp] (244.333,130.785) circle[radius=0.42bp] (245.000,107.726) circle[radius=0.42bp] (245.667,113.430) circle[radius=0.42bp] (246.333,112.890) circle[radius=0.42bp] (247.000,113.227) circle[radius=0.42bp] (247.667,113.153) circle[radius=0.42bp];
\draw[color=tg000000,line width=0.7bp,line cap=round] (232.333,115.609) -- (233.000,115.609) -- (233.667,115.609) -- (234.333,115.609) -- (235.000,115.609) -- (235.667,115.609) -- (236.333,115.609) -- (237.000,115.609) -- (237.667,115.609) -- (238.333,115.609) -- (239.000,115.609) -- (239.667,115.609) -- (240.333,115.609) -- (241.000,115.609) -- (241.667,115.609) -- (242.333,115.609) -- (243.000,115.609) -- (243.667,115.609) -- (244.333,115.609) -- (245.000,115.609) -- (245.667,115.609) -- (246.333,115.609) -- (247.000,115.609) -- (247.667,115.609);
\draw[color=tg0000ff,line width=0.3bp,line cap=round] (231.783,119.634) -- (232.883,120.734);
\draw[color=tg0000ff,line width=0.3bp,line cap=round] (231.783,120.734) -- (232.883,119.634);
\draw[color=tg0000ff,line width=0.3bp,line cap=round] (232.450,115.116) -- (233.550,116.216);
\draw[color=tg0000ff,line width=0.3bp,line cap=round] (232.450,116.216) -- (233.550,115.116);
\draw[color=tg0000ff,line width=0.3bp,line cap=round] (233.117,113.048) -- (234.217,114.148);
\draw[color=tg0000ff,line width=0.3bp,line cap=round] (233.117,114.148) -- (234.217,113.048);
\draw[color=tg0000ff,line width=0.3bp,line cap=round] (233.783,119.148) -- (234.883,120.248);
\draw[color=tg0000ff,line width=0.3bp,line cap=round] (233.783,120.248) -- (234.883,119.148);
\draw[color=tg0000ff,line width=0.3bp,line cap=round] (234.450,109.991) -- (235.550,111.091);
\draw[color=tg0000ff,line width=0.3bp,line cap=round] (234.450,111.091) -- (235.550,109.991);
\draw[color=tg0000ff,line width=0.3bp,line cap=round] (235.117,113.966) -- (236.217,115.066);
\draw[color=tg0000ff,line width=0.3bp,line cap=round] (235.117,115.066) -- (236.217,113.966);
\draw[color=tg0000ff,line width=0.3bp,line cap=round] (235.783,116.271) -- (236.883,117.371);
\draw[color=tg0000ff,line width=0.3bp,line cap=round] (235.783,117.371) -- (236.883,116.271);
\draw[color=tg0000ff,line width=0.3bp,line cap=round] (236.450,115.274) -- (237.550,116.374);
\draw[color=tg0000ff,line width=0.3bp,line cap=round] (236.450,116.374) -- (237.550,115.274);
\draw[color=tg0000ff,line width=0.3bp,line cap=round] (237.117,113.051) -- (238.217,114.151);
\draw[color=tg0000ff,line width=0.3bp,line cap=round] (237.117,114.151) -- (238.217,113.051);
\draw[color=tg0000ff,line width=0.3bp,line cap=round] (237.783,115.720) -- (238.883,116.820);
\draw[color=tg0000ff,line width=0.3bp,line cap=round] (237.783,116.820) -- (238.883,115.720);
\draw[color=tg0000ff,line width=0.3bp,line cap=round] (238.450,117.209) -- (239.550,118.309);
\draw[color=tg0000ff,line width=0.3bp,line cap=round] (238.450,118.309) -- (239.550,117.209);
\draw[color=tg0000ff,line width=0.3bp,line cap=round] (239.117,111.641) -- (240.217,112.741);
\draw[color=tg0000ff,line width=0.3bp,line cap=round] (239.117,112.741) -- (240.217,111.641);
\draw[color=tg0000ff,line width=0.3bp,line cap=round] (239.783,116.715) -- (240.883,117.815);
\draw[color=tg0000ff,line width=0.3bp,line cap=round] (239.783,117.815) -- (240.883,116.715);
\draw[color=tg0000ff,line width=0.3bp,line cap=round] (240.450,115.123) -- (241.550,116.223);
\draw[color=tg0000ff,line width=0.3bp,line cap=round] (240.450,116.223) -- (241.550,115.123);
\draw[color=tg0000ff,line width=0.3bp,line cap=round] (241.117,112.244) -- (242.217,113.344);
\draw[color=tg0000ff,line width=0.3bp,line cap=round] (241.117,113.344) -- (242.217,112.244);
\draw[color=tg0000ff,line width=0.3bp,line cap=round] (241.783,117.237) -- (242.883,118.337);
\draw[color=tg0000ff,line width=0.3bp,line cap=round] (241.783,118.337) -- (242.883,117.237);
\draw[color=tg0000ff,line width=0.3bp,line cap=round] (242.450,114.945) -- (243.550,116.045);
\draw[color=tg0000ff,line width=0.3bp,line cap=round] (242.450,116.045) -- (243.550,114.945);
\draw[color=tg0000ff,line width=0.3bp,line cap=round] (243.117,114.528) -- (244.217,115.628);
\draw[color=tg0000ff,line width=0.3bp,line cap=round] (243.117,115.628) -- (244.217,114.528);
\draw[color=tg0000ff,line width=0.3bp,line cap=round] (243.783,118.056) -- (244.883,119.156);
\draw[color=tg0000ff,line width=0.3bp,line cap=round] (243.783,119.156) -- (244.883,118.056);
\draw[color=tg0000ff,line width=0.3bp,line cap=round] (244.450,112.972) -- (245.550,114.072);
\draw[color=tg0000ff,line width=0.3bp,line cap=round] (244.450,114.072) -- (245.550,112.972);
\draw[color=tg0000ff,line width=0.3bp,line cap=round] (245.117,112.381) -- (246.217,113.481);
\draw[color=tg0000ff,line width=0.3bp,line cap=round] (245.117,113.481) -- (246.217,112.381);
\draw[color=tg0000ff,line width=0.3bp,line cap=round] (245.783,114.036) -- (246.883,115.136);
\draw[color=tg0000ff,line width=0.3bp,line cap=round] (245.783,115.136) -- (246.883,114.036);
\draw[color=tg0000ff,line width=0.3bp,line cap=round] (246.450,111.785) -- (247.550,112.885);
\draw[color=tg0000ff,line width=0.3bp,line cap=round] (246.450,112.885) -- (247.550,111.785);
\draw[color=tg0000ff,line width=0.3bp,line cap=round] (247.117,113.038) -- (248.217,114.138);
\draw[color=tg0000ff,line width=0.3bp,line cap=round] (247.117,114.138) -- (248.217,113.038);
\path[fill=tg8fd18f,draw=none,line width=0.5bp] (255.000,90.080) rectangle (292.969,98.680);
\node[anchor=center,inner sep=0pt,outer sep=0pt,text=tg0b0b0b,font={\sffamily\fontsize{6.6bp}{7.92bp}\selectfont\bfseries}] at (273.984,94.380) {84  Good};
\draw[color=tgc9cdd1,line width=0.45bp,line cap=round] (255.000,132.466) -- (335.000,132.466);
\draw[color=tg9da7b0,line width=0.5bp,line cap=round,dash pattern=on 2bp off 2bp] (319.000,101.080) -- (319.000,132.466);
\path[fill=tg808080,draw=none] (255.333,109.812) circle[radius=0.42bp] (256.000,115.006) circle[radius=0.42bp] (256.667,120.379) circle[radius=0.42bp] (257.333,121.523) circle[radius=0.42bp] (258.000,118.546) circle[radius=0.42bp] (258.667,114.196) circle[radius=0.42bp] (259.333,122.087) circle[radius=0.42bp] (260.000,116.951) circle[radius=0.42bp] (260.667,116.630) circle[radius=0.42bp] (261.333,112.139) circle[radius=0.42bp] (262.000,114.232) circle[radius=0.42bp] (262.667,113.329) circle[radius=0.42bp] (263.333,119.804) circle[radius=0.42bp] (264.000,116.440) circle[radius=0.42bp] (264.667,110.579) circle[radius=0.42bp] (265.333,106.027) circle[radius=0.42bp] (266.000,123.687) circle[radius=0.42bp] (266.667,105.895) circle[radius=0.42bp] (267.333,106.974) circle[radius=0.42bp] (268.000,110.596) circle[radius=0.42bp] (268.667,113.912) circle[radius=0.42bp] (269.333,117.623) circle[radius=0.42bp] (270.000,106.254) circle[radius=0.42bp] (270.667,103.032) circle[radius=0.42bp] (271.333,104.049) circle[radius=0.42bp] (272.000,107.171) circle[radius=0.42bp] (272.667,113.316) circle[radius=0.42bp] (273.333,123.361) circle[radius=0.42bp] (274.000,115.637) circle[radius=0.42bp] (274.667,111.397) circle[radius=0.42bp] (275.333,123.875) circle[radius=0.42bp] (276.000,113.074) circle[radius=0.42bp] (276.667,112.851) circle[radius=0.42bp] (277.333,111.143) circle[radius=0.42bp] (278.000,123.206) circle[radius=0.42bp] (278.667,119.854) circle[radius=0.42bp] (279.333,118.409) circle[radius=0.42bp] (280.000,123.114) circle[radius=0.42bp] (280.667,104.449) circle[radius=0.42bp] (281.333,118.791) circle[radius=0.42bp] (282.000,113.498) circle[radius=0.42bp] (282.667,117.268) circle[radius=0.42bp] (283.333,105.449) circle[radius=0.42bp] (284.000,107.137) circle[radius=0.42bp] (284.667,111.545) circle[radius=0.42bp] (285.333,129.747) circle[radius=0.42bp] (286.000,115.275) circle[radius=0.42bp] (286.667,111.222) circle[radius=0.42bp] (287.333,109.167) circle[radius=0.42bp] (288.000,119.573) circle[radius=0.42bp] (288.667,103.919) circle[radius=0.42bp] (289.333,124.081) circle[radius=0.42bp] (290.000,119.853) circle[radius=0.42bp] (290.667,109.609) circle[radius=0.42bp] (291.333,115.294) circle[radius=0.42bp] (292.000,102.761) circle[radius=0.42bp] (292.667,114.399) circle[radius=0.42bp] (293.333,119.671) circle[radius=0.42bp] (294.000,118.031) circle[radius=0.42bp] (294.667,122.610) circle[radius=0.42bp] (295.333,123.806) circle[radius=0.42bp] (296.000,111.564) circle[radius=0.42bp] (296.667,111.880) circle[radius=0.42bp] (297.333,107.303) circle[radius=0.42bp] (298.000,120.450) circle[radius=0.42bp] (298.667,104.771) circle[radius=0.42bp] (299.333,117.453) circle[radius=0.42bp] (300.000,105.507) circle[radius=0.42bp] (300.667,118.386) circle[radius=0.42bp] (301.333,120.328) circle[radius=0.42bp] (302.000,114.006) circle[radius=0.42bp] (302.667,108.991) circle[radius=0.42bp] (303.333,114.576) circle[radius=0.42bp] (304.000,119.366) circle[radius=0.42bp] (304.667,124.214) circle[radius=0.42bp] (305.333,124.601) circle[radius=0.42bp] (306.000,112.384) circle[radius=0.42bp] (306.667,109.259) circle[radius=0.42bp] (307.333,116.663) circle[radius=0.42bp] (308.000,122.502) circle[radius=0.42bp] (308.667,110.007) circle[radius=0.42bp] (309.333,123.824) circle[radius=0.42bp] (310.000,120.184) circle[radius=0.42bp] (310.667,111.624) circle[radius=0.42bp] (311.333,130.046) circle[radius=0.42bp] (312.000,113.130) circle[radius=0.42bp] (312.667,119.341) circle[radius=0.42bp] (313.333,114.908) circle[radius=0.42bp] (314.000,116.094) circle[radius=0.42bp] (314.667,114.312) circle[radius=0.42bp] (315.333,111.155) circle[radius=0.42bp] (316.000,120.474) circle[radius=0.42bp] (316.667,106.492) circle[radius=0.42bp] (317.333,110.950) circle[radius=0.42bp] (318.000,110.195) circle[radius=0.42bp] (318.667,108.135) circle[radius=0.42bp] (319.333,110.556) circle[radius=0.42bp] (320.000,110.193) circle[radius=0.42bp] (320.667,115.124) circle[radius=0.42bp] (321.333,124.763) circle[radius=0.42bp] (322.000,116.475) circle[radius=0.42bp] (322.667,120.546) circle[radius=0.42bp] (323.333,124.737) circle[radius=0.42bp] (324.000,113.951) circle[radius=0.42bp] (324.667,119.257) circle[radius=0.42bp] (325.333,122.216) circle[radius=0.42bp] (326.000,122.301) circle[radius=0.42bp] (326.667,113.887) circle[radius=0.42bp] (327.333,113.308) circle[radius=0.42bp] (328.000,107.124) circle[radius=0.42bp] (328.667,115.698) circle[radius=0.42bp] (329.333,108.924) circle[radius=0.42bp] (330.000,106.612) circle[radius=0.42bp] (330.667,108.229) circle[radius=0.42bp] (331.333,130.785) circle[radius=0.42bp] (332.000,107.726) circle[radius=0.42bp] (332.667,113.430) circle[radius=0.42bp] (333.333,112.890) circle[radius=0.42bp] (334.000,113.227) circle[radius=0.42bp] (334.667,113.153) circle[radius=0.42bp];
\draw[color=tg000000,line width=0.7bp,line cap=round] (319.333,115.609) -- (320.000,115.609) -- (320.667,115.609) -- (321.333,115.609) -- (322.000,115.609) -- (322.667,115.609) -- (323.333,115.609) -- (324.000,115.609) -- (324.667,115.609) -- (325.333,115.609) -- (326.000,115.609) -- (326.667,115.609) -- (327.333,115.609) -- (328.000,115.609) -- (328.667,115.609) -- (329.333,115.609) -- (330.000,115.609) -- (330.667,115.609) -- (331.333,115.609) -- (332.000,115.609) -- (332.667,115.609) -- (333.333,115.609) -- (334.000,115.609) -- (334.667,115.609);
\draw[color=tg0000ff,line width=0.3bp,line cap=round] (318.783,119.634) -- (319.883,120.734);
\draw[color=tg0000ff,line width=0.3bp,line cap=round] (318.783,120.734) -- (319.883,119.634);
\draw[color=tg0000ff,line width=0.3bp,line cap=round] (319.450,115.116) -- (320.550,116.216);
\draw[color=tg0000ff,line width=0.3bp,line cap=round] (319.450,116.216) -- (320.550,115.116);
\draw[color=tg0000ff,line width=0.3bp,line cap=round] (320.117,113.048) -- (321.217,114.148);
\draw[color=tg0000ff,line width=0.3bp,line cap=round] (320.117,114.148) -- (321.217,113.048);
\draw[color=tg0000ff,line width=0.3bp,line cap=round] (320.783,119.148) -- (321.883,120.248);
\draw[color=tg0000ff,line width=0.3bp,line cap=round] (320.783,120.248) -- (321.883,119.148);
\draw[color=tg0000ff,line width=0.3bp,line cap=round] (321.450,109.991) -- (322.550,111.091);
\draw[color=tg0000ff,line width=0.3bp,line cap=round] (321.450,111.091) -- (322.550,109.991);
\draw[color=tg0000ff,line width=0.3bp,line cap=round] (322.117,113.966) -- (323.217,115.066);
\draw[color=tg0000ff,line width=0.3bp,line cap=round] (322.117,115.066) -- (323.217,113.966);
\draw[color=tg0000ff,line width=0.3bp,line cap=round] (322.783,116.271) -- (323.883,117.371);
\draw[color=tg0000ff,line width=0.3bp,line cap=round] (322.783,117.371) -- (323.883,116.271);
\draw[color=tg0000ff,line width=0.3bp,line cap=round] (323.450,115.274) -- (324.550,116.374);
\draw[color=tg0000ff,line width=0.3bp,line cap=round] (323.450,116.374) -- (324.550,115.274);
\draw[color=tg0000ff,line width=0.3bp,line cap=round] (324.117,113.051) -- (325.217,114.151);
\draw[color=tg0000ff,line width=0.3bp,line cap=round] (324.117,114.151) -- (325.217,113.051);
\draw[color=tg0000ff,line width=0.3bp,line cap=round] (324.783,115.720) -- (325.883,116.820);
\draw[color=tg0000ff,line width=0.3bp,line cap=round] (324.783,116.820) -- (325.883,115.720);
\draw[color=tg0000ff,line width=0.3bp,line cap=round] (325.450,117.209) -- (326.550,118.309);
\draw[color=tg0000ff,line width=0.3bp,line cap=round] (325.450,118.309) -- (326.550,117.209);
\draw[color=tg0000ff,line width=0.3bp,line cap=round] (326.117,111.641) -- (327.217,112.741);
\draw[color=tg0000ff,line width=0.3bp,line cap=round] (326.117,112.741) -- (327.217,111.641);
\draw[color=tg0000ff,line width=0.3bp,line cap=round] (326.783,116.715) -- (327.883,117.815);
\draw[color=tg0000ff,line width=0.3bp,line cap=round] (326.783,117.815) -- (327.883,116.715);
\draw[color=tg0000ff,line width=0.3bp,line cap=round] (327.450,115.123) -- (328.550,116.223);
\draw[color=tg0000ff,line width=0.3bp,line cap=round] (327.450,116.223) -- (328.550,115.123);
\draw[color=tg0000ff,line width=0.3bp,line cap=round] (328.117,112.244) -- (329.217,113.344);
\draw[color=tg0000ff,line width=0.3bp,line cap=round] (328.117,113.344) -- (329.217,112.244);
\draw[color=tg0000ff,line width=0.3bp,line cap=round] (328.783,117.237) -- (329.883,118.337);
\draw[color=tg0000ff,line width=0.3bp,line cap=round] (328.783,118.337) -- (329.883,117.237);
\draw[color=tg0000ff,line width=0.3bp,line cap=round] (329.450,114.945) -- (330.550,116.045);
\draw[color=tg0000ff,line width=0.3bp,line cap=round] (329.450,116.045) -- (330.550,114.945);
\draw[color=tg0000ff,line width=0.3bp,line cap=round] (330.117,114.528) -- (331.217,115.628);
\draw[color=tg0000ff,line width=0.3bp,line cap=round] (330.117,115.628) -- (331.217,114.528);
\draw[color=tg0000ff,line width=0.3bp,line cap=round] (330.783,118.056) -- (331.883,119.156);
\draw[color=tg0000ff,line width=0.3bp,line cap=round] (330.783,119.156) -- (331.883,118.056);
\draw[color=tg0000ff,line width=0.3bp,line cap=round] (331.450,112.972) -- (332.550,114.072);
\draw[color=tg0000ff,line width=0.3bp,line cap=round] (331.450,114.072) -- (332.550,112.972);
\draw[color=tg0000ff,line width=0.3bp,line cap=round] (332.117,112.381) -- (333.217,113.481);
\draw[color=tg0000ff,line width=0.3bp,line cap=round] (332.117,113.481) -- (333.217,112.381);
\draw[color=tg0000ff,line width=0.3bp,line cap=round] (332.783,114.036) -- (333.883,115.136);
\draw[color=tg0000ff,line width=0.3bp,line cap=round] (332.783,115.136) -- (333.883,114.036);
\draw[color=tg0000ff,line width=0.3bp,line cap=round] (333.450,111.785) -- (334.550,112.885);
\draw[color=tg0000ff,line width=0.3bp,line cap=round] (333.450,112.885) -- (334.550,111.785);
\draw[color=tg0000ff,line width=0.3bp,line cap=round] (334.117,113.038) -- (335.217,114.138);
\draw[color=tg0000ff,line width=0.3bp,line cap=round] (334.117,114.138) -- (335.217,113.038);
\node[anchor=west,inner sep=0pt,outer sep=0pt,text=tg626b73,font={\sffamily\fontsize{6.2bp}{7.44bp}\selectfont}] at (346.000,93.800) {None found};
\draw[color=tgeceef0,line width=0.4bp,line cap=round] (0.000,135.466) -- (448.000,135.466);
\node[anchor=west,inner sep=0pt,outer sep=0pt,text=tg18314b,font={\sffamily\fontsize{7.2bp}{8.64bp}\selectfont\bfseries}] at (0.000,142.286) {Noise around a fixed level};
\node[anchor=west,inner sep=0pt,outer sep=0pt,text=tg18314b,font={\sffamily\fontsize{6.2bp}{7.44bp}\selectfont}] at (0.000,150.726) {\ensuremath{\checkmark} The level, estimated from the whole history.};
\node[anchor=west,inner sep=0pt,outer sep=0pt,text=tg18314b,font={\sffamily\fontsize{6.2bp}{7.44bp}\selectfont}] at (0.000,158.566) {\ensuremath{\times} Each wiggle around the level.};
\path[fill=tg8fd18f,draw=none,line width=0.5bp] (168.000,137.966) rectangle (205.969,146.566);
\node[anchor=center,inner sep=0pt,outer sep=0pt,text=tg0b0b0b,font={\sffamily\fontsize{6.6bp}{7.92bp}\selectfont\bfseries}] at (186.984,142.266) {84  Good};
\draw[color=tgc9cdd1,line width=0.45bp,line cap=round] (168.000,178.652) -- (248.000,178.652);
\draw[color=tg9da7b0,line width=0.5bp,line cap=round,dash pattern=on 2bp off 2bp] (232.000,148.966) -- (232.000,178.652);
\path[fill=tg808080,draw=none] (168.333,157.225) circle[radius=0.42bp] (169.000,162.137) circle[radius=0.42bp] (169.667,167.220) circle[radius=0.42bp] (170.333,168.301) circle[radius=0.42bp] (171.000,165.486) circle[radius=0.42bp] (171.667,161.372) circle[radius=0.42bp] (172.333,168.835) circle[radius=0.42bp] (173.000,163.977) circle[radius=0.42bp] (173.667,163.674) circle[radius=0.42bp] (174.333,159.426) circle[radius=0.42bp] (175.000,161.405) circle[radius=0.42bp] (175.667,160.551) circle[radius=0.42bp] (176.333,166.676) circle[radius=0.42bp] (177.000,163.494) circle[radius=0.42bp] (177.667,157.950) circle[radius=0.42bp] (178.333,153.645) circle[radius=0.42bp] (179.000,170.348) circle[radius=0.42bp] (179.667,153.521) circle[radius=0.42bp] (180.333,154.540) circle[radius=0.42bp] (181.000,157.966) circle[radius=0.42bp] (181.667,161.103) circle[radius=0.42bp] (182.333,164.613) circle[radius=0.42bp] (183.000,153.860) circle[radius=0.42bp] (183.667,150.812) circle[radius=0.42bp] (184.333,151.775) circle[radius=0.42bp] (185.000,154.727) circle[radius=0.42bp] (185.667,160.539) circle[radius=0.42bp] (186.333,170.041) circle[radius=0.42bp] (187.000,162.735) circle[radius=0.42bp] (187.667,158.724) circle[radius=0.42bp] (188.333,170.526) circle[radius=0.42bp] (189.000,160.310) circle[radius=0.42bp] (189.667,160.099) circle[radius=0.42bp] (190.333,158.484) circle[radius=0.42bp] (191.000,169.894) circle[radius=0.42bp] (191.667,166.723) circle[radius=0.42bp] (192.333,165.356) circle[radius=0.42bp] (193.000,169.807) circle[radius=0.42bp] (193.667,152.152) circle[radius=0.42bp] (194.333,165.717) circle[radius=0.42bp] (195.000,160.711) circle[radius=0.42bp] (195.667,164.277) circle[radius=0.42bp] (196.333,153.099) circle[radius=0.42bp] (197.000,154.695) circle[radius=0.42bp] (197.667,158.864) circle[radius=0.42bp] (198.333,176.080) circle[radius=0.42bp] (199.000,162.392) circle[radius=0.42bp] (199.667,158.559) circle[radius=0.42bp] (200.333,156.615) circle[radius=0.42bp] (201.000,166.458) circle[radius=0.42bp] (201.667,151.651) circle[radius=0.42bp] (202.333,170.721) circle[radius=0.42bp] (203.000,166.722) circle[radius=0.42bp] (203.667,157.033) circle[radius=0.42bp] (204.333,162.410) circle[radius=0.42bp] (205.000,150.556) circle[radius=0.42bp] (205.667,161.564) circle[radius=0.42bp] (206.333,166.550) circle[radius=0.42bp] (207.000,164.999) circle[radius=0.42bp] (207.667,169.329) circle[radius=0.42bp] (208.333,170.461) circle[radius=0.42bp] (209.000,158.882) circle[radius=0.42bp] (209.667,159.181) circle[radius=0.42bp] (210.333,154.852) circle[radius=0.42bp] (211.000,167.287) circle[radius=0.42bp] (211.667,152.457) circle[radius=0.42bp] (212.333,164.452) circle[radius=0.42bp] (213.000,153.154) circle[radius=0.42bp] (213.667,165.334) circle[radius=0.42bp] (214.333,167.171) circle[radius=0.42bp] (215.000,161.192) circle[radius=0.42bp] (215.667,156.448) circle[radius=0.42bp] (216.333,161.731) circle[radius=0.42bp] (217.000,166.261) circle[radius=0.42bp] (217.667,170.847) circle[radius=0.42bp] (218.333,171.213) circle[radius=0.42bp] (219.000,159.657) circle[radius=0.42bp] (219.667,156.702) circle[radius=0.42bp] (220.333,163.705) circle[radius=0.42bp] (221.000,169.228) circle[radius=0.42bp] (221.667,157.410) circle[radius=0.42bp] (222.333,170.478) circle[radius=0.42bp] (223.000,167.035) circle[radius=0.42bp] (223.667,158.939) circle[radius=0.42bp] (224.333,176.363) circle[radius=0.42bp] (225.000,160.363) circle[radius=0.42bp] (225.667,166.238) circle[radius=0.42bp] (226.333,162.045) circle[radius=0.42bp] (227.000,163.167) circle[radius=0.42bp] (227.667,161.481) circle[radius=0.42bp] (228.333,158.495) circle[radius=0.42bp] (229.000,167.310) circle[radius=0.42bp] (229.667,154.085) circle[radius=0.42bp] (230.333,158.302) circle[radius=0.42bp] (231.000,157.588) circle[radius=0.42bp] (231.667,155.639) circle[radius=0.42bp] (232.333,157.928) circle[radius=0.42bp] (233.000,157.586) circle[radius=0.42bp] (233.667,162.249) circle[radius=0.42bp] (234.333,171.366) circle[radius=0.42bp] (235.000,163.527) circle[radius=0.42bp] (235.667,167.378) circle[radius=0.42bp] (236.333,171.342) circle[radius=0.42bp] (237.000,161.139) circle[radius=0.42bp] (237.667,166.158) circle[radius=0.42bp] (238.333,168.957) circle[radius=0.42bp] (239.000,169.037) circle[radius=0.42bp] (239.667,161.079) circle[radius=0.42bp] (240.333,160.531) circle[radius=0.42bp] (241.000,154.683) circle[radius=0.42bp] (241.667,162.792) circle[radius=0.42bp] (242.333,156.385) circle[radius=0.42bp] (243.000,154.198) circle[radius=0.42bp] (243.667,155.728) circle[radius=0.42bp] (244.333,177.062) circle[radius=0.42bp] (245.000,155.252) circle[radius=0.42bp] (245.667,160.647) circle[radius=0.42bp] (246.333,160.136) circle[radius=0.42bp] (247.000,160.455) circle[radius=0.42bp] (247.667,160.385) circle[radius=0.42bp];
\draw[color=tg000000,line width=0.7bp,line cap=round] (232.333,162.708) -- (233.000,162.708) -- (233.667,162.708) -- (234.333,162.708) -- (235.000,162.708) -- (235.667,162.708) -- (236.333,162.708) -- (237.000,162.708) -- (237.667,162.708) -- (238.333,162.708) -- (239.000,162.708) -- (239.667,162.708) -- (240.333,162.708) -- (241.000,162.708) -- (241.667,162.708) -- (242.333,162.708) -- (243.000,162.708) -- (243.667,162.708) -- (244.333,162.708) -- (245.000,162.708) -- (245.667,162.708) -- (246.333,162.708) -- (247.000,162.708) -- (247.667,162.708);
\draw[color=tg0000ff,line width=0.3bp,line cap=round] (231.783,166.485) -- (232.883,167.585);
\draw[color=tg0000ff,line width=0.3bp,line cap=round] (231.783,167.585) -- (232.883,166.485);
\draw[color=tg0000ff,line width=0.3bp,line cap=round] (232.450,162.212) -- (233.550,163.312);
\draw[color=tg0000ff,line width=0.3bp,line cap=round] (232.450,163.312) -- (233.550,162.212);
\draw[color=tg0000ff,line width=0.3bp,line cap=round] (233.117,160.256) -- (234.217,161.356);
\draw[color=tg0000ff,line width=0.3bp,line cap=round] (233.117,161.356) -- (234.217,160.256);
\draw[color=tg0000ff,line width=0.3bp,line cap=round] (233.783,166.025) -- (234.883,167.125);
\draw[color=tg0000ff,line width=0.3bp,line cap=round] (233.783,167.125) -- (234.883,166.025);
\draw[color=tg0000ff,line width=0.3bp,line cap=round] (234.450,157.365) -- (235.550,158.465);
\draw[color=tg0000ff,line width=0.3bp,line cap=round] (234.450,158.465) -- (235.550,157.365);
\draw[color=tg0000ff,line width=0.3bp,line cap=round] (235.117,161.125) -- (236.217,162.225);
\draw[color=tg0000ff,line width=0.3bp,line cap=round] (235.117,162.225) -- (236.217,161.125);
\draw[color=tg0000ff,line width=0.3bp,line cap=round] (235.783,163.304) -- (236.883,164.404);
\draw[color=tg0000ff,line width=0.3bp,line cap=round] (235.783,164.404) -- (236.883,163.304);
\draw[color=tg0000ff,line width=0.3bp,line cap=round] (236.450,162.362) -- (237.550,163.462);
\draw[color=tg0000ff,line width=0.3bp,line cap=round] (236.450,163.462) -- (237.550,162.362);
\draw[color=tg0000ff,line width=0.3bp,line cap=round] (237.117,160.259) -- (238.217,161.359);
\draw[color=tg0000ff,line width=0.3bp,line cap=round] (237.117,161.359) -- (238.217,160.259);
\draw[color=tg0000ff,line width=0.3bp,line cap=round] (237.783,162.783) -- (238.883,163.883);
\draw[color=tg0000ff,line width=0.3bp,line cap=round] (237.783,163.883) -- (238.883,162.783);
\draw[color=tg0000ff,line width=0.3bp,line cap=round] (238.450,164.192) -- (239.550,165.292);
\draw[color=tg0000ff,line width=0.3bp,line cap=round] (238.450,165.292) -- (239.550,164.192);
\draw[color=tg0000ff,line width=0.3bp,line cap=round] (239.117,158.925) -- (240.217,160.025);
\draw[color=tg0000ff,line width=0.3bp,line cap=round] (239.117,160.025) -- (240.217,158.925);
\draw[color=tg0000ff,line width=0.3bp,line cap=round] (239.783,163.724) -- (240.883,164.824);
\draw[color=tg0000ff,line width=0.3bp,line cap=round] (239.783,164.824) -- (240.883,163.724);
\draw[color=tg0000ff,line width=0.3bp,line cap=round] (240.450,162.219) -- (241.550,163.319);
\draw[color=tg0000ff,line width=0.3bp,line cap=round] (240.450,163.319) -- (241.550,162.219);
\draw[color=tg0000ff,line width=0.3bp,line cap=round] (241.117,159.496) -- (242.217,160.596);
\draw[color=tg0000ff,line width=0.3bp,line cap=round] (241.117,160.596) -- (242.217,159.496);
\draw[color=tg0000ff,line width=0.3bp,line cap=round] (241.783,164.218) -- (242.883,165.318);
\draw[color=tg0000ff,line width=0.3bp,line cap=round] (241.783,165.318) -- (242.883,164.218);
\draw[color=tg0000ff,line width=0.3bp,line cap=round] (242.450,162.050) -- (243.550,163.150);
\draw[color=tg0000ff,line width=0.3bp,line cap=round] (242.450,163.150) -- (243.550,162.050);
\draw[color=tg0000ff,line width=0.3bp,line cap=round] (243.117,161.656) -- (244.217,162.756);
\draw[color=tg0000ff,line width=0.3bp,line cap=round] (243.117,162.756) -- (244.217,161.656);
\draw[color=tg0000ff,line width=0.3bp,line cap=round] (243.783,164.993) -- (244.883,166.093);
\draw[color=tg0000ff,line width=0.3bp,line cap=round] (243.783,166.093) -- (244.883,164.993);
\draw[color=tg0000ff,line width=0.3bp,line cap=round] (244.450,160.184) -- (245.550,161.284);
\draw[color=tg0000ff,line width=0.3bp,line cap=round] (244.450,161.284) -- (245.550,160.184);
\draw[color=tg0000ff,line width=0.3bp,line cap=round] (245.117,159.625) -- (246.217,160.725);
\draw[color=tg0000ff,line width=0.3bp,line cap=round] (245.117,160.725) -- (246.217,159.625);
\draw[color=tg0000ff,line width=0.3bp,line cap=round] (245.783,161.191) -- (246.883,162.291);
\draw[color=tg0000ff,line width=0.3bp,line cap=round] (245.783,162.291) -- (246.883,161.191);
\draw[color=tg0000ff,line width=0.3bp,line cap=round] (246.450,159.062) -- (247.550,160.162);
\draw[color=tg0000ff,line width=0.3bp,line cap=round] (246.450,160.162) -- (247.550,159.062);
\draw[color=tg0000ff,line width=0.3bp,line cap=round] (247.117,160.246) -- (248.217,161.346);
\draw[color=tg0000ff,line width=0.3bp,line cap=round] (247.117,161.346) -- (248.217,160.246);
\path[fill=tg8fd18f,draw=none,line width=0.5bp] (255.000,137.966) rectangle (292.969,146.566);
\node[anchor=center,inner sep=0pt,outer sep=0pt,text=tg0b0b0b,font={\sffamily\fontsize{6.6bp}{7.92bp}\selectfont\bfseries}] at (273.984,142.266) {84  Good};
\draw[color=tgc9cdd1,line width=0.45bp,line cap=round] (255.000,178.652) -- (335.000,178.652);
\draw[color=tg9da7b0,line width=0.5bp,line cap=round,dash pattern=on 2bp off 2bp] (319.000,148.966) -- (319.000,178.652);
\path[fill=tg808080,draw=none] (255.333,157.225) circle[radius=0.42bp] (256.000,162.137) circle[radius=0.42bp] (256.667,167.220) circle[radius=0.42bp] (257.333,168.301) circle[radius=0.42bp] (258.000,165.486) circle[radius=0.42bp] (258.667,161.372) circle[radius=0.42bp] (259.333,168.835) circle[radius=0.42bp] (260.000,163.977) circle[radius=0.42bp] (260.667,163.674) circle[radius=0.42bp] (261.333,159.426) circle[radius=0.42bp] (262.000,161.405) circle[radius=0.42bp] (262.667,160.551) circle[radius=0.42bp] (263.333,166.676) circle[radius=0.42bp] (264.000,163.494) circle[radius=0.42bp] (264.667,157.950) circle[radius=0.42bp] (265.333,153.645) circle[radius=0.42bp] (266.000,170.348) circle[radius=0.42bp] (266.667,153.521) circle[radius=0.42bp] (267.333,154.540) circle[radius=0.42bp] (268.000,157.966) circle[radius=0.42bp] (268.667,161.103) circle[radius=0.42bp] (269.333,164.613) circle[radius=0.42bp] (270.000,153.860) circle[radius=0.42bp] (270.667,150.812) circle[radius=0.42bp] (271.333,151.775) circle[radius=0.42bp] (272.000,154.727) circle[radius=0.42bp] (272.667,160.539) circle[radius=0.42bp] (273.333,170.041) circle[radius=0.42bp] (274.000,162.735) circle[radius=0.42bp] (274.667,158.724) circle[radius=0.42bp] (275.333,170.526) circle[radius=0.42bp] (276.000,160.310) circle[radius=0.42bp] (276.667,160.099) circle[radius=0.42bp] (277.333,158.484) circle[radius=0.42bp] (278.000,169.894) circle[radius=0.42bp] (278.667,166.723) circle[radius=0.42bp] (279.333,165.356) circle[radius=0.42bp] (280.000,169.807) circle[radius=0.42bp] (280.667,152.152) circle[radius=0.42bp] (281.333,165.717) circle[radius=0.42bp] (282.000,160.711) circle[radius=0.42bp] (282.667,164.277) circle[radius=0.42bp] (283.333,153.099) circle[radius=0.42bp] (284.000,154.695) circle[radius=0.42bp] (284.667,158.864) circle[radius=0.42bp] (285.333,176.080) circle[radius=0.42bp] (286.000,162.392) circle[radius=0.42bp] (286.667,158.559) circle[radius=0.42bp] (287.333,156.615) circle[radius=0.42bp] (288.000,166.458) circle[radius=0.42bp] (288.667,151.651) circle[radius=0.42bp] (289.333,170.721) circle[radius=0.42bp] (290.000,166.722) circle[radius=0.42bp] (290.667,157.033) circle[radius=0.42bp] (291.333,162.410) circle[radius=0.42bp] (292.000,150.556) circle[radius=0.42bp] (292.667,161.564) circle[radius=0.42bp] (293.333,166.550) circle[radius=0.42bp] (294.000,164.999) circle[radius=0.42bp] (294.667,169.329) circle[radius=0.42bp] (295.333,170.461) circle[radius=0.42bp] (296.000,158.882) circle[radius=0.42bp] (296.667,159.181) circle[radius=0.42bp] (297.333,154.852) circle[radius=0.42bp] (298.000,167.287) circle[radius=0.42bp] (298.667,152.457) circle[radius=0.42bp] (299.333,164.452) circle[radius=0.42bp] (300.000,153.154) circle[radius=0.42bp] (300.667,165.334) circle[radius=0.42bp] (301.333,167.171) circle[radius=0.42bp] (302.000,161.192) circle[radius=0.42bp] (302.667,156.448) circle[radius=0.42bp] (303.333,161.731) circle[radius=0.42bp] (304.000,166.261) circle[radius=0.42bp] (304.667,170.847) circle[radius=0.42bp] (305.333,171.213) circle[radius=0.42bp] (306.000,159.657) circle[radius=0.42bp] (306.667,156.702) circle[radius=0.42bp] (307.333,163.705) circle[radius=0.42bp] (308.000,169.228) circle[radius=0.42bp] (308.667,157.410) circle[radius=0.42bp] (309.333,170.478) circle[radius=0.42bp] (310.000,167.035) circle[radius=0.42bp] (310.667,158.939) circle[radius=0.42bp] (311.333,176.363) circle[radius=0.42bp] (312.000,160.363) circle[radius=0.42bp] (312.667,166.238) circle[radius=0.42bp] (313.333,162.045) circle[radius=0.42bp] (314.000,163.167) circle[radius=0.42bp] (314.667,161.481) circle[radius=0.42bp] (315.333,158.495) circle[radius=0.42bp] (316.000,167.310) circle[radius=0.42bp] (316.667,154.085) circle[radius=0.42bp] (317.333,158.302) circle[radius=0.42bp] (318.000,157.588) circle[radius=0.42bp] (318.667,155.639) circle[radius=0.42bp] (319.333,157.928) circle[radius=0.42bp] (320.000,157.586) circle[radius=0.42bp] (320.667,162.249) circle[radius=0.42bp] (321.333,171.366) circle[radius=0.42bp] (322.000,163.527) circle[radius=0.42bp] (322.667,167.378) circle[radius=0.42bp] (323.333,171.342) circle[radius=0.42bp] (324.000,161.139) circle[radius=0.42bp] (324.667,166.158) circle[radius=0.42bp] (325.333,168.957) circle[radius=0.42bp] (326.000,169.037) circle[radius=0.42bp] (326.667,161.079) circle[radius=0.42bp] (327.333,160.531) circle[radius=0.42bp] (328.000,154.683) circle[radius=0.42bp] (328.667,162.792) circle[radius=0.42bp] (329.333,156.385) circle[radius=0.42bp] (330.000,154.198) circle[radius=0.42bp] (330.667,155.728) circle[radius=0.42bp] (331.333,177.062) circle[radius=0.42bp] (332.000,155.252) circle[radius=0.42bp] (332.667,160.647) circle[radius=0.42bp] (333.333,160.136) circle[radius=0.42bp] (334.000,160.455) circle[radius=0.42bp] (334.667,160.385) circle[radius=0.42bp];
\draw[color=tg000000,line width=0.7bp,line cap=round] (319.333,162.708) -- (320.000,162.708) -- (320.667,162.708) -- (321.333,162.708) -- (322.000,162.708) -- (322.667,162.708) -- (323.333,162.708) -- (324.000,162.708) -- (324.667,162.708) -- (325.333,162.708) -- (326.000,162.708) -- (326.667,162.708) -- (327.333,162.708) -- (328.000,162.708) -- (328.667,162.708) -- (329.333,162.708) -- (330.000,162.708) -- (330.667,162.708) -- (331.333,162.708) -- (332.000,162.708) -- (332.667,162.708) -- (333.333,162.708) -- (334.000,162.708) -- (334.667,162.708);
\draw[color=tg0000ff,line width=0.3bp,line cap=round] (318.783,166.485) -- (319.883,167.585);
\draw[color=tg0000ff,line width=0.3bp,line cap=round] (318.783,167.585) -- (319.883,166.485);
\draw[color=tg0000ff,line width=0.3bp,line cap=round] (319.450,162.212) -- (320.550,163.312);
\draw[color=tg0000ff,line width=0.3bp,line cap=round] (319.450,163.312) -- (320.550,162.212);
\draw[color=tg0000ff,line width=0.3bp,line cap=round] (320.117,160.256) -- (321.217,161.356);
\draw[color=tg0000ff,line width=0.3bp,line cap=round] (320.117,161.356) -- (321.217,160.256);
\draw[color=tg0000ff,line width=0.3bp,line cap=round] (320.783,166.025) -- (321.883,167.125);
\draw[color=tg0000ff,line width=0.3bp,line cap=round] (320.783,167.125) -- (321.883,166.025);
\draw[color=tg0000ff,line width=0.3bp,line cap=round] (321.450,157.365) -- (322.550,158.465);
\draw[color=tg0000ff,line width=0.3bp,line cap=round] (321.450,158.465) -- (322.550,157.365);
\draw[color=tg0000ff,line width=0.3bp,line cap=round] (322.117,161.125) -- (323.217,162.225);
\draw[color=tg0000ff,line width=0.3bp,line cap=round] (322.117,162.225) -- (323.217,161.125);
\draw[color=tg0000ff,line width=0.3bp,line cap=round] (322.783,163.304) -- (323.883,164.404);
\draw[color=tg0000ff,line width=0.3bp,line cap=round] (322.783,164.404) -- (323.883,163.304);
\draw[color=tg0000ff,line width=0.3bp,line cap=round] (323.450,162.362) -- (324.550,163.462);
\draw[color=tg0000ff,line width=0.3bp,line cap=round] (323.450,163.462) -- (324.550,162.362);
\draw[color=tg0000ff,line width=0.3bp,line cap=round] (324.117,160.259) -- (325.217,161.359);
\draw[color=tg0000ff,line width=0.3bp,line cap=round] (324.117,161.359) -- (325.217,160.259);
\draw[color=tg0000ff,line width=0.3bp,line cap=round] (324.783,162.783) -- (325.883,163.883);
\draw[color=tg0000ff,line width=0.3bp,line cap=round] (324.783,163.883) -- (325.883,162.783);
\draw[color=tg0000ff,line width=0.3bp,line cap=round] (325.450,164.192) -- (326.550,165.292);
\draw[color=tg0000ff,line width=0.3bp,line cap=round] (325.450,165.292) -- (326.550,164.192);
\draw[color=tg0000ff,line width=0.3bp,line cap=round] (326.117,158.925) -- (327.217,160.025);
\draw[color=tg0000ff,line width=0.3bp,line cap=round] (326.117,160.025) -- (327.217,158.925);
\draw[color=tg0000ff,line width=0.3bp,line cap=round] (326.783,163.724) -- (327.883,164.824);
\draw[color=tg0000ff,line width=0.3bp,line cap=round] (326.783,164.824) -- (327.883,163.724);
\draw[color=tg0000ff,line width=0.3bp,line cap=round] (327.450,162.219) -- (328.550,163.319);
\draw[color=tg0000ff,line width=0.3bp,line cap=round] (327.450,163.319) -- (328.550,162.219);
\draw[color=tg0000ff,line width=0.3bp,line cap=round] (328.117,159.496) -- (329.217,160.596);
\draw[color=tg0000ff,line width=0.3bp,line cap=round] (328.117,160.596) -- (329.217,159.496);
\draw[color=tg0000ff,line width=0.3bp,line cap=round] (328.783,164.218) -- (329.883,165.318);
\draw[color=tg0000ff,line width=0.3bp,line cap=round] (328.783,165.318) -- (329.883,164.218);
\draw[color=tg0000ff,line width=0.3bp,line cap=round] (329.450,162.050) -- (330.550,163.150);
\draw[color=tg0000ff,line width=0.3bp,line cap=round] (329.450,163.150) -- (330.550,162.050);
\draw[color=tg0000ff,line width=0.3bp,line cap=round] (330.117,161.656) -- (331.217,162.756);
\draw[color=tg0000ff,line width=0.3bp,line cap=round] (330.117,162.756) -- (331.217,161.656);
\draw[color=tg0000ff,line width=0.3bp,line cap=round] (330.783,164.993) -- (331.883,166.093);
\draw[color=tg0000ff,line width=0.3bp,line cap=round] (330.783,166.093) -- (331.883,164.993);
\draw[color=tg0000ff,line width=0.3bp,line cap=round] (331.450,160.184) -- (332.550,161.284);
\draw[color=tg0000ff,line width=0.3bp,line cap=round] (331.450,161.284) -- (332.550,160.184);
\draw[color=tg0000ff,line width=0.3bp,line cap=round] (332.117,159.625) -- (333.217,160.725);
\draw[color=tg0000ff,line width=0.3bp,line cap=round] (332.117,160.725) -- (333.217,159.625);
\draw[color=tg0000ff,line width=0.3bp,line cap=round] (332.783,161.191) -- (333.883,162.291);
\draw[color=tg0000ff,line width=0.3bp,line cap=round] (332.783,162.291) -- (333.883,161.191);
\draw[color=tg0000ff,line width=0.3bp,line cap=round] (333.450,159.062) -- (334.550,160.162);
\draw[color=tg0000ff,line width=0.3bp,line cap=round] (333.450,160.162) -- (334.550,159.062);
\draw[color=tg0000ff,line width=0.3bp,line cap=round] (334.117,160.246) -- (335.217,161.346);
\draw[color=tg0000ff,line width=0.3bp,line cap=round] (334.117,161.346) -- (335.217,160.246);
\node[anchor=west,inner sep=0pt,outer sep=0pt,text=tg626b73,font={\sffamily\fontsize{6.2bp}{7.44bp}\selectfont}] at (346.000,141.686) {None found};
\draw[color=tgeceef0,line width=0.4bp,line cap=round] (0.000,181.652) -- (448.000,181.652);
\node[anchor=west,inner sep=0pt,outer sep=0pt,text=tg18314b,font={\sffamily\fontsize{7.2bp}{8.64bp}\selectfont\bfseries}] at (0.000,188.472) {Straight-line trend};
\node[anchor=west,inner sep=0pt,outer sep=0pt,text=tg18314b,font={\sffamily\fontsize{6.2bp}{7.44bp}\selectfont}] at (0.000,196.912) {\ensuremath{\checkmark} The line, extended any number of steps ahead.};
\node[anchor=west,inner sep=0pt,outer sep=0pt,text=tg18314b,font={\sffamily\fontsize{6.2bp}{7.44bp}\selectfont}] at (0.000,204.752) {\ensuremath{\times} The noise around the line.};
\path[fill=tg8fd18f,draw=none,line width=0.5bp] (168.000,184.152) rectangle (205.969,192.752);
\node[anchor=center,inner sep=0pt,outer sep=0pt,text=tg0b0b0b,font={\sffamily\fontsize{6.6bp}{7.92bp}\selectfont\bfseries}] at (186.984,188.452) {83  Good};
\draw[color=tgc9cdd1,line width=0.45bp,line cap=round] (168.000,224.838) -- (248.000,224.838);
\draw[color=tg9da7b0,line width=0.5bp,line cap=round,dash pattern=on 2bp off 2bp] (232.000,195.152) -- (232.000,224.838);
\path[fill=tg808080,draw=none] (168.333,221.903) circle[radius=0.42bp] (169.000,222.560) circle[radius=0.42bp] (169.667,223.248) circle[radius=0.42bp] (170.333,223.224) circle[radius=0.42bp] (171.000,222.508) circle[radius=0.42bp] (171.667,221.561) circle[radius=0.42bp] (172.333,222.672) circle[radius=0.42bp] (173.000,221.593) circle[radius=0.42bp] (173.667,221.323) circle[radius=0.42bp] (174.333,220.352) circle[radius=0.42bp] (175.000,220.488) circle[radius=0.42bp] (175.667,220.121) circle[radius=0.42bp] (176.333,220.994) circle[radius=0.42bp] (177.000,220.212) circle[radius=0.42bp] (177.667,219.011) circle[radius=0.42bp] (178.333,218.030) circle[radius=0.42bp] (179.000,220.784) circle[radius=0.42bp] (179.667,217.577) circle[radius=0.42bp] (180.333,217.542) circle[radius=0.42bp] (181.000,217.935) circle[radius=0.42bp] (181.667,218.277) circle[radius=0.42bp] (182.333,218.685) circle[radius=0.42bp] (183.000,216.558) circle[radius=0.42bp] (183.667,215.801) circle[radius=0.42bp] (184.333,215.756) circle[radius=0.42bp] (185.000,216.065) circle[radius=0.42bp] (185.667,216.883) circle[radius=0.42bp] (186.333,218.356) circle[radius=0.42bp] (187.000,216.841) circle[radius=0.42bp] (187.667,215.913) circle[radius=0.42bp] (188.333,217.795) circle[radius=0.42bp] (189.000,215.763) circle[radius=0.42bp] (189.667,215.510) circle[radius=0.42bp] (190.333,215.007) circle[radius=0.42bp] (191.000,216.819) circle[radius=0.42bp] (191.667,216.040) circle[radius=0.42bp] (192.333,215.581) circle[radius=0.42bp] (193.000,216.157) circle[radius=0.42bp] (193.667,212.803) circle[radius=0.42bp] (194.333,214.998) circle[radius=0.42bp] (195.000,213.893) circle[radius=0.42bp] (195.667,214.311) circle[radius=0.42bp] (196.333,212.108) circle[radius=0.42bp] (197.000,212.176) circle[radius=0.42bp] (197.667,212.702) circle[radius=0.42bp] (198.333,215.546) circle[radius=0.42bp] (199.000,212.897) circle[radius=0.42bp] (199.667,212.000) circle[radius=0.42bp] (200.333,211.439) circle[radius=0.42bp] (201.000,212.973) circle[radius=0.42bp] (201.667,210.125) circle[radius=0.42bp] (202.333,213.299) circle[radius=0.42bp] (203.000,212.373) circle[radius=0.42bp] (203.667,210.435) circle[radius=0.42bp] (204.333,211.175) circle[radius=0.42bp] (205.000,208.852) circle[radius=0.42bp] (205.667,210.593) circle[radius=0.42bp] (206.333,211.263) circle[radius=0.42bp] (207.000,210.772) circle[radius=0.42bp] (207.667,211.326) circle[radius=0.42bp] (208.333,211.311) circle[radius=0.42bp] (209.000,209.037) circle[radius=0.42bp] (209.667,208.875) circle[radius=0.42bp] (210.333,207.889) circle[radius=0.42bp] (211.000,209.884) circle[radius=0.42bp] (211.667,207.032) circle[radius=0.42bp] (212.333,208.949) circle[radius=0.42bp] (213.000,206.725) circle[radius=0.42bp] (213.667,208.674) circle[radius=0.42bp] (214.333,208.785) circle[radius=0.42bp] (215.000,207.506) circle[radius=0.42bp] (215.667,206.447) circle[radius=0.42bp] (216.333,207.171) circle[radius=0.42bp] (217.000,207.760) circle[radius=0.42bp] (217.667,208.360) circle[radius=0.42bp] (218.333,208.209) circle[radius=0.42bp] (219.000,205.939) circle[radius=0.42bp] (219.667,205.198) circle[radius=0.42bp] (220.333,206.227) circle[radius=0.42bp] (221.000,206.993) circle[radius=0.42bp] (221.667,204.677) circle[radius=0.42bp] (222.333,206.784) circle[radius=0.42bp] (223.000,205.956) circle[radius=0.42bp] (223.667,204.301) circle[radius=0.42bp] (224.333,207.183) circle[radius=0.42bp] (225.000,204.123) circle[radius=0.42bp] (225.667,204.951) circle[radius=0.42bp] (226.333,203.990) circle[radius=0.42bp] (227.000,203.974) circle[radius=0.42bp] (227.667,203.459) circle[radius=0.42bp] (228.333,202.712) circle[radius=0.42bp] (229.000,204.063) circle[radius=0.42bp] (229.667,201.497) circle[radius=0.42bp] (230.333,202.031) circle[radius=0.42bp] (231.000,201.688) circle[radius=0.42bp] (231.667,201.126) circle[radius=0.42bp] (232.333,201.317) circle[radius=0.42bp] (233.000,201.040) circle[radius=0.42bp] (233.667,201.654) circle[radius=0.42bp] (234.333,203.058) circle[radius=0.42bp] (235.000,201.449) circle[radius=0.42bp] (235.667,201.918) circle[radius=0.42bp] (236.333,202.407) circle[radius=0.42bp] (237.000,200.378) circle[radius=0.42bp] (237.667,201.054) circle[radius=0.42bp] (238.333,201.336) circle[radius=0.42bp] (239.000,201.134) circle[radius=0.42bp] (239.667,199.504) circle[radius=0.42bp] (240.333,199.191) circle[radius=0.42bp] (241.000,197.936) circle[radius=0.42bp] (241.667,199.161) circle[radius=0.42bp] (242.333,197.807) circle[radius=0.42bp] (243.000,197.202) circle[radius=0.42bp] (243.667,197.258) circle[radius=0.42bp] (244.333,200.835) circle[radius=0.42bp] (245.000,196.742) circle[radius=0.42bp] (245.667,197.486) circle[radius=0.42bp] (246.333,197.179) circle[radius=0.42bp] (247.000,197.020) circle[radius=0.42bp] (247.667,196.792) circle[radius=0.42bp];
\draw[color=tg000000,line width=0.7bp,line cap=round] (232.333,202.167) -- (233.000,201.951) -- (233.667,201.735) -- (234.333,201.519) -- (235.000,201.304) -- (235.667,201.088) -- (236.333,200.872) -- (237.000,200.656) -- (237.667,200.441) -- (238.333,200.225) -- (239.000,200.009) -- (239.667,199.794) -- (240.333,199.578) -- (241.000,199.362) -- (241.667,199.146) -- (242.333,198.931) -- (243.000,198.715) -- (243.667,198.499) -- (244.333,198.283) -- (245.000,198.068) -- (245.667,197.852) -- (246.333,197.636) -- (247.000,197.420) -- (247.667,197.205);
\draw[color=tg0000ff,line width=0.3bp,line cap=round] (231.783,202.401) -- (232.883,203.501);
\draw[color=tg0000ff,line width=0.3bp,line cap=round] (231.783,203.501) -- (232.883,202.401);
\draw[color=tg0000ff,line width=0.3bp,line cap=round] (232.450,201.740) -- (233.550,202.840);
\draw[color=tg0000ff,line width=0.3bp,line cap=round] (232.450,202.840) -- (233.550,201.740);
\draw[color=tg0000ff,line width=0.3bp,line cap=round] (233.117,201.348) -- (234.217,202.448);
\draw[color=tg0000ff,line width=0.3bp,line cap=round] (233.117,202.448) -- (234.217,201.348);
\draw[color=tg0000ff,line width=0.3bp,line cap=round] (233.783,201.720) -- (234.883,202.820);
\draw[color=tg0000ff,line width=0.3bp,line cap=round] (233.783,202.820) -- (234.883,201.720);
\draw[color=tg0000ff,line width=0.3bp,line cap=round] (234.450,200.341) -- (235.550,201.441);
\draw[color=tg0000ff,line width=0.3bp,line cap=round] (234.450,201.441) -- (235.550,200.341);
\draw[color=tg0000ff,line width=0.3bp,line cap=round] (235.117,200.665) -- (236.217,201.765);
\draw[color=tg0000ff,line width=0.3bp,line cap=round] (235.117,201.765) -- (236.217,200.665);
\draw[color=tg0000ff,line width=0.3bp,line cap=round] (235.783,200.789) -- (236.883,201.889);
\draw[color=tg0000ff,line width=0.3bp,line cap=round] (235.783,201.889) -- (236.883,200.789);
\draw[color=tg0000ff,line width=0.3bp,line cap=round] (236.450,200.513) -- (237.550,201.613);
\draw[color=tg0000ff,line width=0.3bp,line cap=round] (236.450,201.613) -- (237.550,200.513);
\draw[color=tg0000ff,line width=0.3bp,line cap=round] (237.117,199.985) -- (238.217,201.085);
\draw[color=tg0000ff,line width=0.3bp,line cap=round] (237.117,201.085) -- (238.217,199.985);
\draw[color=tg0000ff,line width=0.3bp,line cap=round] (237.783,200.126) -- (238.883,201.226);
\draw[color=tg0000ff,line width=0.3bp,line cap=round] (237.783,201.226) -- (238.883,200.126);
\draw[color=tg0000ff,line width=0.3bp,line cap=round] (238.450,200.047) -- (239.550,201.147);
\draw[color=tg0000ff,line width=0.3bp,line cap=round] (238.450,201.147) -- (239.550,200.047);
\draw[color=tg0000ff,line width=0.3bp,line cap=round] (239.117,199.164) -- (240.217,200.264);
\draw[color=tg0000ff,line width=0.3bp,line cap=round] (239.117,200.264) -- (240.217,199.164);
\draw[color=tg0000ff,line width=0.3bp,line cap=round] (239.783,199.582) -- (240.883,200.682);
\draw[color=tg0000ff,line width=0.3bp,line cap=round] (239.783,200.682) -- (240.883,199.582);
\draw[color=tg0000ff,line width=0.3bp,line cap=round] (240.450,199.020) -- (241.550,200.120);
\draw[color=tg0000ff,line width=0.3bp,line cap=round] (240.450,200.120) -- (241.550,199.020);
\draw[color=tg0000ff,line width=0.3bp,line cap=round] (241.117,198.518) -- (242.217,199.618);
\draw[color=tg0000ff,line width=0.3bp,line cap=round] (241.117,199.618) -- (242.217,198.518);
\draw[color=tg0000ff,line width=0.3bp,line cap=round] (241.783,198.851) -- (242.883,199.951);
\draw[color=tg0000ff,line width=0.3bp,line cap=round] (241.783,199.951) -- (242.883,198.851);
\draw[color=tg0000ff,line width=0.3bp,line cap=round] (242.450,198.429) -- (243.550,199.529);
\draw[color=tg0000ff,line width=0.3bp,line cap=round] (242.450,199.529) -- (243.550,198.429);
\draw[color=tg0000ff,line width=0.3bp,line cap=round] (243.117,198.147) -- (244.217,199.247);
\draw[color=tg0000ff,line width=0.3bp,line cap=round] (243.117,199.247) -- (244.217,198.147);
\draw[color=tg0000ff,line width=0.3bp,line cap=round] (243.783,198.278) -- (244.883,199.378);
\draw[color=tg0000ff,line width=0.3bp,line cap=round] (243.783,199.378) -- (244.883,198.278);
\draw[color=tg0000ff,line width=0.3bp,line cap=round] (244.450,197.627) -- (245.550,198.727);
\draw[color=tg0000ff,line width=0.3bp,line cap=round] (244.450,198.727) -- (245.550,197.627);
\draw[color=tg0000ff,line width=0.3bp,line cap=round] (245.117,197.194) -- (246.217,198.294);
\draw[color=tg0000ff,line width=0.3bp,line cap=round] (245.117,198.294) -- (246.217,197.194);
\draw[color=tg0000ff,line width=0.3bp,line cap=round] (245.783,197.172) -- (246.883,198.272);
\draw[color=tg0000ff,line width=0.3bp,line cap=round] (245.783,198.272) -- (246.883,197.172);
\draw[color=tg0000ff,line width=0.3bp,line cap=round] (246.450,196.812) -- (247.550,197.912);
\draw[color=tg0000ff,line width=0.3bp,line cap=round] (246.450,197.912) -- (247.550,196.812);
\draw[color=tg0000ff,line width=0.3bp,line cap=round] (247.117,196.808) -- (248.217,197.908);
\draw[color=tg0000ff,line width=0.3bp,line cap=round] (247.117,197.908) -- (248.217,196.808);
\path[fill=tg8fd18f,draw=none,line width=0.5bp] (255.000,184.152) rectangle (292.969,192.752);
\node[anchor=center,inner sep=0pt,outer sep=0pt,text=tg0b0b0b,font={\sffamily\fontsize{6.6bp}{7.92bp}\selectfont\bfseries}] at (273.984,188.452) {78  Good};
\draw[color=tgc9cdd1,line width=0.45bp,line cap=round] (255.000,224.838) -- (335.000,224.838);
\draw[color=tg9da7b0,line width=0.5bp,line cap=round,dash pattern=on 2bp off 2bp] (319.000,195.152) -- (319.000,224.838);
\path[fill=tg808080,draw=none] (255.333,217.536) circle[radius=0.42bp] (256.000,220.125) circle[radius=0.42bp] (256.667,222.810) circle[radius=0.42bp] (257.333,223.248) circle[radius=0.42bp] (258.000,221.495) circle[radius=0.42bp] (258.667,219.013) circle[radius=0.42bp] (259.333,223.036) circle[radius=0.42bp] (260.000,220.136) circle[radius=0.42bp] (260.667,219.796) circle[radius=0.42bp] (261.333,217.238) circle[radius=0.42bp] (262.000,218.180) circle[radius=0.42bp] (262.667,217.530) circle[radius=0.42bp] (263.333,220.800) circle[radius=0.42bp] (264.000,218.842) circle[radius=0.42bp] (264.667,215.557) circle[radius=0.42bp] (265.333,212.968) circle[radius=0.42bp] (266.000,222.182) circle[radius=0.42bp] (266.667,212.557) circle[radius=0.42bp] (267.333,212.959) circle[radius=0.42bp] (268.000,214.714) circle[radius=0.42bp] (268.667,216.305) circle[radius=0.42bp] (269.333,218.107) circle[radius=0.42bp] (270.000,211.895) circle[radius=0.42bp] (270.667,210.012) circle[radius=0.42bp] (271.333,210.382) circle[radius=0.42bp] (272.000,211.871) circle[radius=0.42bp] (272.667,214.966) circle[radius=0.42bp] (273.333,220.133) circle[radius=0.42bp] (274.000,215.858) circle[radius=0.42bp] (274.667,213.434) circle[radius=0.42bp] (275.333,219.895) circle[radius=0.42bp] (276.000,213.985) circle[radius=0.42bp] (276.667,213.696) circle[radius=0.42bp] (277.333,212.618) circle[radius=0.42bp] (278.000,218.858) circle[radius=0.42bp] (278.667,216.906) circle[radius=0.42bp] (279.333,215.967) circle[radius=0.42bp] (280.000,218.297) circle[radius=0.42bp] (280.667,208.208) circle[radius=0.42bp] (281.333,215.659) circle[radius=0.42bp] (282.000,212.676) circle[radius=0.42bp] (282.667,214.509) circle[radius=0.42bp] (283.333,208.058) circle[radius=0.42bp] (284.000,208.784) circle[radius=0.42bp] (284.667,210.956) circle[radius=0.42bp] (285.333,220.458) circle[radius=0.42bp] (286.000,212.597) circle[radius=0.42bp] (286.667,210.273) circle[radius=0.42bp] (287.333,209.011) circle[radius=0.42bp] (288.000,214.370) circle[radius=0.42bp] (288.667,205.881) circle[radius=0.42bp] (289.333,216.424) circle[radius=0.42bp] (290.000,214.007) circle[radius=0.42bp] (290.667,208.393) circle[radius=0.42bp] (291.333,211.244) circle[radius=0.42bp] (292.000,204.413) circle[radius=0.42bp] (292.667,210.427) circle[radius=0.42bp] (293.333,213.058) circle[radius=0.42bp] (294.000,212.016) circle[radius=0.42bp] (294.667,214.279) circle[radius=0.42bp] (295.333,214.744) circle[radius=0.42bp] (296.000,208.068) circle[radius=0.42bp] (296.667,208.066) circle[radius=0.42bp] (297.333,205.463) circle[radius=0.42bp] (298.000,212.279) circle[radius=0.42bp] (298.667,203.777) circle[radius=0.42bp] (299.333,210.345) circle[radius=0.42bp] (300.000,203.827) circle[radius=0.42bp] (300.667,210.500) circle[radius=0.42bp] (301.333,211.361) circle[radius=0.42bp] (302.000,207.832) circle[radius=0.42bp] (302.667,204.996) circle[radius=0.42bp] (303.333,207.793) circle[radius=0.42bp] (304.000,210.168) circle[radius=0.42bp] (304.667,212.574) circle[radius=0.42bp] (305.333,212.609) circle[radius=0.42bp] (306.000,205.947) circle[radius=0.42bp] (306.667,204.116) circle[radius=0.42bp] (307.333,207.880) circle[radius=0.42bp] (308.000,210.812) circle[radius=0.42bp] (308.667,204.002) circle[radius=0.42bp] (309.333,211.174) circle[radius=0.42bp] (310.000,209.069) circle[radius=0.42bp] (310.667,204.350) circle[radius=0.42bp] (311.333,213.968) circle[radius=0.42bp] (312.000,204.809) circle[radius=0.42bp] (312.667,207.939) circle[radius=0.42bp] (313.333,205.413) circle[radius=0.42bp] (314.000,205.873) circle[radius=0.42bp] (314.667,204.755) circle[radius=0.42bp] (315.333,202.907) circle[radius=0.42bp] (316.000,207.689) circle[radius=0.42bp] (316.667,200.088) circle[radius=0.42bp] (317.333,202.287) circle[radius=0.42bp] (318.000,201.715) circle[radius=0.42bp] (318.667,200.450) circle[radius=0.42bp] (319.333,201.566) circle[radius=0.42bp] (320.000,201.203) circle[radius=0.42bp] (320.667,203.652) circle[radius=0.42bp] (321.333,208.604) circle[radius=0.42bp] (322.000,204.029) circle[radius=0.42bp] (322.667,206.022) circle[radius=0.42bp] (323.333,208.079) circle[radius=0.42bp] (324.000,202.176) circle[radius=0.42bp] (324.667,204.826) circle[radius=0.42bp] (325.333,206.228) circle[radius=0.42bp] (326.000,206.102) circle[radius=0.42bp] (326.667,201.461) circle[radius=0.42bp] (327.333,200.982) circle[radius=0.42bp] (328.000,197.526) circle[radius=0.42bp] (328.667,201.912) circle[radius=0.42bp] (329.333,198.142) circle[radius=0.42bp] (330.000,196.742) circle[radius=0.42bp] (330.667,197.431) circle[radius=0.42bp] (331.333,209.247) circle[radius=0.42bp] (332.000,196.823) circle[radius=0.42bp] (332.667,199.684) circle[radius=0.42bp] (333.333,199.226) circle[radius=0.42bp] (334.000,199.235) circle[radius=0.42bp] (334.667,199.025) circle[radius=0.42bp];
\draw[color=tg000000,line width=0.7bp,line cap=round] (319.333,204.251) -- (320.000,204.080) -- (320.667,203.910) -- (321.333,203.740) -- (322.000,203.569) -- (322.667,203.399) -- (323.333,203.228) -- (324.000,203.058) -- (324.667,202.887) -- (325.333,202.717) -- (326.000,202.546) -- (326.667,202.376) -- (327.333,202.205) -- (328.000,202.035) -- (328.667,201.864) -- (329.333,201.694) -- (330.000,201.523) -- (330.667,201.353) -- (331.333,201.182) -- (332.000,201.012) -- (332.667,200.841) -- (333.333,200.671) -- (334.000,200.500) -- (334.667,200.330);
\draw[color=tg0000ff,line width=0.3bp,line cap=round] (318.783,207.472) -- (319.883,208.572);
\draw[color=tg0000ff,line width=0.3bp,line cap=round] (318.783,208.572) -- (319.883,207.472);
\draw[color=tg0000ff,line width=0.3bp,line cap=round] (319.450,204.995) -- (320.550,206.095);
\draw[color=tg0000ff,line width=0.3bp,line cap=round] (319.450,206.095) -- (320.550,204.995);
\draw[color=tg0000ff,line width=0.3bp,line cap=round] (320.117,204.065) -- (321.217,205.165);
\draw[color=tg0000ff,line width=0.3bp,line cap=round] (320.117,205.165) -- (321.217,204.065);
\draw[color=tg0000ff,line width=0.3bp,line cap=round] (320.783,206.781) -- (321.883,207.881);
\draw[color=tg0000ff,line width=0.3bp,line cap=round] (320.783,207.881) -- (321.883,206.781);
\draw[color=tg0000ff,line width=0.3bp,line cap=round] (321.450,201.602) -- (322.550,202.702);
\draw[color=tg0000ff,line width=0.3bp,line cap=round] (321.450,202.702) -- (322.550,201.602);
\draw[color=tg0000ff,line width=0.3bp,line cap=round] (322.117,203.802) -- (323.217,204.902);
\draw[color=tg0000ff,line width=0.3bp,line cap=round] (322.117,204.902) -- (323.217,203.802);
\draw[color=tg0000ff,line width=0.3bp,line cap=round] (322.783,205.047) -- (323.883,206.147);
\draw[color=tg0000ff,line width=0.3bp,line cap=round] (322.783,206.147) -- (323.883,205.047);
\draw[color=tg0000ff,line width=0.3bp,line cap=round] (323.450,204.535) -- (324.550,205.635);
\draw[color=tg0000ff,line width=0.3bp,line cap=round] (323.450,205.635) -- (324.550,204.535);
\draw[color=tg0000ff,line width=0.3bp,line cap=round] (324.117,202.861) -- (325.217,203.961);
\draw[color=tg0000ff,line width=0.3bp,line cap=round] (324.117,203.961) -- (325.217,202.861);
\draw[color=tg0000ff,line width=0.3bp,line cap=round] (324.783,204.163) -- (325.883,205.263);
\draw[color=tg0000ff,line width=0.3bp,line cap=round] (324.783,205.263) -- (325.883,204.163);
\draw[color=tg0000ff,line width=0.3bp,line cap=round] (325.450,204.629) -- (326.550,205.729);
\draw[color=tg0000ff,line width=0.3bp,line cap=round] (325.450,205.729) -- (326.550,204.629);
\draw[color=tg0000ff,line width=0.3bp,line cap=round] (326.117,201.438) -- (327.217,202.538);
\draw[color=tg0000ff,line width=0.3bp,line cap=round] (326.117,202.538) -- (327.217,201.438);
\draw[color=tg0000ff,line width=0.3bp,line cap=round] (326.783,204.108) -- (327.883,205.208);
\draw[color=tg0000ff,line width=0.3bp,line cap=round] (326.783,205.208) -- (327.883,204.108);
\draw[color=tg0000ff,line width=0.3bp,line cap=round] (327.450,202.612) -- (328.550,203.712);
\draw[color=tg0000ff,line width=0.3bp,line cap=round] (327.450,203.712) -- (328.550,202.612);
\draw[color=tg0000ff,line width=0.3bp,line cap=round] (328.117,201.102) -- (329.217,202.202);
\draw[color=tg0000ff,line width=0.3bp,line cap=round] (328.117,202.202) -- (329.217,201.102);
\draw[color=tg0000ff,line width=0.3bp,line cap=round] (328.783,203.732) -- (329.883,204.832);
\draw[color=tg0000ff,line width=0.3bp,line cap=round] (328.783,204.832) -- (329.883,203.732);
\draw[color=tg0000ff,line width=0.3bp,line cap=round] (329.450,202.268) -- (330.550,203.368);
\draw[color=tg0000ff,line width=0.3bp,line cap=round] (329.450,203.368) -- (330.550,202.268);
\draw[color=tg0000ff,line width=0.3bp,line cap=round] (330.117,201.976) -- (331.217,203.076);
\draw[color=tg0000ff,line width=0.3bp,line cap=round] (330.117,203.076) -- (331.217,201.976);
\draw[color=tg0000ff,line width=0.3bp,line cap=round] (330.783,203.647) -- (331.883,204.747);
\draw[color=tg0000ff,line width=0.3bp,line cap=round] (330.783,204.747) -- (331.883,203.647);
\draw[color=tg0000ff,line width=0.3bp,line cap=round] (331.450,200.796) -- (332.550,201.896);
\draw[color=tg0000ff,line width=0.3bp,line cap=round] (331.450,201.896) -- (332.550,200.796);
\draw[color=tg0000ff,line width=0.3bp,line cap=round] (332.117,200.110) -- (333.217,201.210);
\draw[color=tg0000ff,line width=0.3bp,line cap=round] (332.117,201.210) -- (333.217,200.110);
\draw[color=tg0000ff,line width=0.3bp,line cap=round] (332.783,200.987) -- (333.883,202.087);
\draw[color=tg0000ff,line width=0.3bp,line cap=round] (332.783,202.087) -- (333.883,200.987);
\draw[color=tg0000ff,line width=0.3bp,line cap=round] (333.450,199.821) -- (334.550,200.921);
\draw[color=tg0000ff,line width=0.3bp,line cap=round] (333.450,200.921) -- (334.550,199.821);
\draw[color=tg0000ff,line width=0.3bp,line cap=round] (334.117,200.651) -- (335.217,201.751);
\draw[color=tg0000ff,line width=0.3bp,line cap=round] (334.117,201.751) -- (335.217,200.651);
\node[anchor=west,inner sep=0pt,outer sep=0pt,text=tg626b73,font={\sffamily\fontsize{6.2bp}{7.44bp}\selectfont}] at (346.000,187.872) {None found};
\draw[color=tgeceef0,line width=0.4bp,line cap=round] (0.000,227.838) -- (448.000,227.838);
\node[anchor=west,inner sep=0pt,outer sep=0pt,text=tg18314b,font={\sffamily\fontsize{7.2bp}{8.64bp}\selectfont\bfseries}] at (0.000,234.658) {Repeating seasonal cycle};
\node[anchor=west,inner sep=0pt,outer sep=0pt,text=tg18314b,font={\sffamily\fontsize{6.2bp}{7.44bp}\selectfont}] at (0.000,243.098) {\ensuremath{\checkmark} The repeating shape, any number of steps};
\node[anchor=west,inner sep=0pt,outer sep=0pt,text=tg18314b,font={\sffamily\fontsize{6.2bp}{7.44bp}\selectfont}] at (0.000,250.538) {ahead.};
\node[anchor=west,inner sep=0pt,outer sep=0pt,text=tg18314b,font={\sffamily\fontsize{6.2bp}{7.44bp}\selectfont}] at (0.000,258.378) {\ensuremath{\times} The noise around the cycle.};
\path[fill=tg8fd18f,draw=none,line width=0.5bp] (168.000,230.338) rectangle (205.969,238.938);
\node[anchor=center,inner sep=0pt,outer sep=0pt,text=tg0b0b0b,font={\sffamily\fontsize{6.6bp}{7.92bp}\selectfont\bfseries}] at (186.984,234.638) {77  Good};
\draw[color=tgc9cdd1,line width=0.45bp,line cap=round] (168.000,271.024) -- (248.000,271.024);
\draw[color=tg9da7b0,line width=0.5bp,line cap=round,dash pattern=on 2bp off 2bp] (232.000,241.338) -- (232.000,271.024);
\path[fill=tg808080,draw=none] (168.333,254.056) circle[radius=0.42bp] (169.000,250.357) circle[radius=0.42bp] (169.667,248.063) circle[radius=0.42bp] (170.333,247.006) circle[radius=0.42bp] (171.000,247.580) circle[radius=0.42bp] (171.667,250.144) circle[radius=0.42bp] (172.333,257.286) circle[radius=0.42bp] (173.000,261.000) circle[radius=0.42bp] (173.667,264.624) circle[radius=0.42bp] (174.333,264.800) circle[radius=0.42bp] (175.000,263.993) circle[radius=0.42bp] (175.667,260.047) circle[radius=0.42bp] (176.333,256.686) circle[radius=0.42bp] (177.000,250.735) circle[radius=0.42bp] (177.667,245.484) circle[radius=0.42bp] (178.333,242.928) circle[radius=0.42bp] (179.000,248.933) circle[radius=0.42bp] (179.667,247.959) circle[radius=0.42bp] (180.333,253.309) circle[radius=0.42bp] (181.000,259.328) circle[radius=0.42bp] (181.667,263.909) circle[radius=0.42bp] (182.333,266.243) circle[radius=0.42bp] (183.000,261.894) circle[radius=0.42bp] (183.667,257.337) circle[radius=0.42bp] (184.333,252.539) circle[radius=0.42bp] (185.000,248.295) circle[radius=0.42bp] (185.667,246.204) circle[radius=0.42bp] (186.333,247.490) circle[radius=0.42bp] (187.000,246.815) circle[radius=0.42bp] (187.667,249.407) circle[radius=0.42bp] (188.333,257.757) circle[radius=0.42bp] (189.000,259.980) circle[radius=0.42bp] (189.667,263.630) circle[radius=0.42bp] (190.333,264.538) circle[radius=0.42bp] (191.000,266.355) circle[radius=0.42bp] (191.667,261.765) circle[radius=0.42bp] (192.333,256.318) circle[radius=0.42bp] (193.000,252.491) circle[radius=0.42bp] (193.667,243.870) circle[radius=0.42bp] (194.333,246.287) circle[radius=0.42bp] (195.000,246.252) circle[radius=0.42bp] (195.667,250.952) circle[radius=0.42bp] (196.333,252.908) circle[radius=0.42bp] (197.000,258.418) circle[radius=0.42bp] (197.667,263.286) circle[radius=0.42bp] (198.333,269.434) circle[radius=0.42bp] (199.000,264.268) circle[radius=0.42bp] (199.667,259.493) circle[radius=0.42bp] (200.333,253.886) circle[radius=0.42bp] (201.000,251.559) circle[radius=0.42bp] (201.667,243.731) circle[radius=0.42bp] (202.333,247.680) circle[radius=0.42bp] (203.000,247.924) circle[radius=0.42bp] (203.667,248.937) circle[radius=0.42bp] (204.333,255.499) circle[radius=0.42bp] (205.000,257.266) circle[radius=0.42bp] (205.667,264.037) circle[radius=0.42bp] (206.333,266.782) circle[radius=0.42bp] (207.000,264.993) circle[radius=0.42bp] (207.667,262.490) circle[radius=0.42bp] (208.333,257.739) circle[radius=0.42bp] (209.000,249.451) circle[radius=0.42bp] (209.667,245.826) circle[radius=0.42bp] (210.333,243.264) circle[radius=0.42bp] (211.000,248.082) circle[radius=0.42bp] (211.667,247.664) circle[radius=0.42bp] (212.333,256.067) circle[radius=0.42bp] (213.000,257.989) circle[radius=0.42bp] (213.667,265.086) circle[radius=0.42bp] (214.333,266.955) circle[radius=0.42bp] (215.000,263.934) circle[radius=0.42bp] (215.667,258.906) circle[radius=0.42bp] (216.333,255.310) circle[radius=0.42bp] (217.000,251.504) circle[radius=0.42bp] (217.667,249.072) circle[radius=0.42bp] (218.333,247.817) circle[radius=0.42bp] (219.000,245.959) circle[radius=0.42bp] (219.667,248.845) circle[radius=0.42bp] (220.333,255.859) circle[radius=0.42bp] (221.000,262.461) circle[radius=0.42bp] (221.667,262.881) circle[radius=0.42bp] (222.333,267.875) circle[radius=0.42bp] (223.000,265.560) circle[radius=0.42bp] (223.667,259.599) circle[radius=0.42bp] (224.333,259.381) circle[radius=0.42bp] (225.000,249.863) circle[radius=0.42bp] (225.667,247.790) circle[radius=0.42bp] (226.333,245.266) circle[radius=0.42bp] (227.000,246.935) circle[radius=0.42bp] (227.667,250.175) circle[radius=0.42bp] (228.333,254.409) circle[radius=0.42bp] (229.000,261.928) circle[radius=0.42bp] (229.667,261.956) circle[radius=0.42bp] (230.333,264.487) circle[radius=0.42bp] (231.000,262.931) circle[radius=0.42bp] (231.667,258.680) circle[radius=0.42bp] (232.333,254.252) circle[radius=0.42bp] (233.000,249.091) circle[radius=0.42bp] (233.667,246.680) circle[radius=0.42bp] (234.333,247.859) circle[radius=0.42bp] (235.000,247.035) circle[radius=0.42bp] (235.667,251.815) circle[radius=0.42bp] (236.333,257.984) circle[radius=0.42bp] (237.000,260.211) circle[radius=0.42bp] (237.667,265.316) circle[radius=0.42bp] (238.333,267.452) circle[radius=0.42bp] (239.000,266.117) circle[radius=0.42bp] (239.667,260.194) circle[radius=0.42bp] (240.333,254.976) circle[radius=0.42bp] (241.000,248.283) circle[radius=0.42bp] (241.667,246.831) circle[radius=0.42bp] (242.333,243.691) circle[radius=0.42bp] (243.000,244.440) circle[radius=0.42bp] (243.667,248.574) circle[radius=0.42bp] (244.333,259.575) circle[radius=0.42bp] (245.000,258.572) circle[radius=0.42bp] (245.667,263.782) circle[radius=0.42bp] (246.333,264.997) circle[radius=0.42bp] (247.000,263.729) circle[radius=0.42bp] (247.667,260.001) circle[radius=0.42bp];
\draw[color=tg000000,line width=0.7bp,line cap=round] (232.333,255.581) -- (233.000,250.516) -- (233.667,246.807) -- (234.333,245.450) -- (235.000,246.807) -- (235.667,250.516) -- (236.333,255.581) -- (237.000,260.647) -- (237.667,264.356) -- (238.333,265.713) -- (239.000,264.356) -- (239.667,260.647) -- (240.333,255.581) -- (241.000,250.516) -- (241.667,246.807) -- (242.333,245.450) -- (243.000,246.807) -- (243.667,250.516) -- (244.333,255.581) -- (245.000,260.647) -- (245.667,264.356) -- (246.333,265.713) -- (247.000,264.356) -- (247.667,260.647);
\draw[color=tg0000ff,line width=0.3bp,line cap=round] (231.783,256.485) -- (232.883,257.585);
\draw[color=tg0000ff,line width=0.3bp,line cap=round] (231.783,257.585) -- (232.883,256.485);
\draw[color=tg0000ff,line width=0.3bp,line cap=round] (232.450,250.600) -- (233.550,251.700);
\draw[color=tg0000ff,line width=0.3bp,line cap=round] (232.450,251.700) -- (233.550,250.600);
\draw[color=tg0000ff,line width=0.3bp,line cap=round] (233.117,246.370) -- (234.217,247.470);
\draw[color=tg0000ff,line width=0.3bp,line cap=round] (233.117,247.470) -- (234.217,246.370);
\draw[color=tg0000ff,line width=0.3bp,line cap=round] (233.783,245.927) -- (234.883,247.027);
\draw[color=tg0000ff,line width=0.3bp,line cap=round] (233.783,247.027) -- (234.883,245.927);
\draw[color=tg0000ff,line width=0.3bp,line cap=round] (234.450,245.051) -- (235.550,246.151);
\draw[color=tg0000ff,line width=0.3bp,line cap=round] (234.450,246.151) -- (235.550,245.051);
\draw[color=tg0000ff,line width=0.3bp,line cap=round] (235.117,249.488) -- (236.217,250.588);
\draw[color=tg0000ff,line width=0.3bp,line cap=round] (235.117,250.588) -- (236.217,249.488);
\draw[color=tg0000ff,line width=0.3bp,line cap=round] (235.783,254.957) -- (236.883,256.057);
\draw[color=tg0000ff,line width=0.3bp,line cap=round] (235.783,256.057) -- (236.883,254.957);
\draw[color=tg0000ff,line width=0.3bp,line cap=round] (236.450,259.772) -- (237.550,260.872);
\draw[color=tg0000ff,line width=0.3bp,line cap=round] (236.450,260.872) -- (237.550,259.772);
\draw[color=tg0000ff,line width=0.3bp,line cap=round] (237.117,263.006) -- (238.217,264.106);
\draw[color=tg0000ff,line width=0.3bp,line cap=round] (237.117,264.106) -- (238.217,263.006);
\draw[color=tg0000ff,line width=0.3bp,line cap=round] (237.783,265.075) -- (238.883,266.175);
\draw[color=tg0000ff,line width=0.3bp,line cap=round] (237.783,266.175) -- (238.883,265.075);
\draw[color=tg0000ff,line width=0.3bp,line cap=round] (238.450,264.284) -- (239.550,265.384);
\draw[color=tg0000ff,line width=0.3bp,line cap=round] (238.450,265.384) -- (239.550,264.284);
\draw[color=tg0000ff,line width=0.3bp,line cap=round] (239.117,259.650) -- (240.217,260.750);
\draw[color=tg0000ff,line width=0.3bp,line cap=round] (239.117,260.750) -- (240.217,259.650);
\draw[color=tg0000ff,line width=0.3bp,line cap=round] (239.783,255.874) -- (240.883,256.974);
\draw[color=tg0000ff,line width=0.3bp,line cap=round] (239.783,256.974) -- (240.883,255.874);
\draw[color=tg0000ff,line width=0.3bp,line cap=round] (240.450,250.457) -- (241.550,251.557);
\draw[color=tg0000ff,line width=0.3bp,line cap=round] (240.450,251.557) -- (241.550,250.457);
\draw[color=tg0000ff,line width=0.3bp,line cap=round] (241.117,246.179) -- (242.217,247.279);
\draw[color=tg0000ff,line width=0.3bp,line cap=round] (241.117,247.279) -- (242.217,246.179);
\draw[color=tg0000ff,line width=0.3bp,line cap=round] (241.783,245.765) -- (242.883,246.865);
\draw[color=tg0000ff,line width=0.3bp,line cap=round] (241.783,246.865) -- (242.883,245.765);
\draw[color=tg0000ff,line width=0.3bp,line cap=round] (242.450,246.420) -- (243.550,247.520);
\draw[color=tg0000ff,line width=0.3bp,line cap=round] (242.450,247.520) -- (243.550,246.420);
\draw[color=tg0000ff,line width=0.3bp,line cap=round] (243.117,249.698) -- (244.217,250.798);
\draw[color=tg0000ff,line width=0.3bp,line cap=round] (243.117,250.798) -- (244.217,249.698);
\draw[color=tg0000ff,line width=0.3bp,line cap=round] (243.783,255.163) -- (244.883,256.263);
\draw[color=tg0000ff,line width=0.3bp,line cap=round] (243.783,256.263) -- (244.883,255.163);
\draw[color=tg0000ff,line width=0.3bp,line cap=round] (244.450,259.114) -- (245.550,260.214);
\draw[color=tg0000ff,line width=0.3bp,line cap=round] (244.450,260.214) -- (245.550,259.114);
\draw[color=tg0000ff,line width=0.3bp,line cap=round] (245.117,262.489) -- (246.217,263.589);
\draw[color=tg0000ff,line width=0.3bp,line cap=round] (245.117,263.589) -- (246.217,262.489);
\draw[color=tg0000ff,line width=0.3bp,line cap=round] (245.783,264.248) -- (246.883,265.348);
\draw[color=tg0000ff,line width=0.3bp,line cap=round] (245.783,265.348) -- (246.883,264.248);
\draw[color=tg0000ff,line width=0.3bp,line cap=round] (246.450,262.775) -- (247.550,263.875);
\draw[color=tg0000ff,line width=0.3bp,line cap=round] (246.450,263.875) -- (247.550,262.775);
\draw[color=tg0000ff,line width=0.3bp,line cap=round] (247.117,259.683) -- (248.217,260.783);
\draw[color=tg0000ff,line width=0.3bp,line cap=round] (247.117,260.783) -- (248.217,259.683);
\path[fill=tg8fd18f,draw=none,line width=0.5bp] (255.000,230.338) rectangle (292.969,238.938);
\node[anchor=center,inner sep=0pt,outer sep=0pt,text=tg0b0b0b,font={\sffamily\fontsize{6.6bp}{7.92bp}\selectfont\bfseries}] at (273.984,234.638) {74  Good};
\draw[color=tgc9cdd1,line width=0.45bp,line cap=round] (255.000,271.024) -- (335.000,271.024);
\draw[color=tg9da7b0,line width=0.5bp,line cap=round,dash pattern=on 2bp off 2bp] (319.000,241.338) -- (319.000,271.024);
\path[fill=tg808080,draw=none] (255.333,251.541) circle[radius=0.42bp] (256.000,251.772) circle[radius=0.42bp] (256.667,252.892) circle[radius=0.42bp] (257.333,252.804) circle[radius=0.42bp] (258.000,251.782) circle[radius=0.42bp] (258.667,251.282) circle[radius=0.42bp] (259.333,258.971) circle[radius=0.42bp] (260.000,258.775) circle[radius=0.42bp] (260.667,260.714) circle[radius=0.42bp] (261.333,258.775) circle[radius=0.42bp] (262.000,259.261) circle[radius=0.42bp] (262.667,256.583) circle[radius=0.42bp] (263.333,257.589) circle[radius=0.42bp] (264.000,252.640) circle[radius=0.42bp] (264.667,246.960) circle[radius=0.42bp] (265.333,243.424) circle[radius=0.42bp] (266.000,254.895) circle[radius=0.42bp] (266.667,246.257) circle[radius=0.42bp] (267.333,249.823) circle[radius=0.42bp] (268.000,254.928) circle[radius=0.42bp] (268.667,259.068) circle[radius=0.42bp] (269.333,262.095) circle[radius=0.42bp] (270.000,254.432) circle[radius=0.42bp] (270.667,250.349) circle[radius=0.42bp] (271.333,248.053) circle[radius=0.42bp] (272.000,247.029) circle[radius=0.42bp] (272.667,248.617) circle[radius=0.42bp] (273.333,253.917) circle[radius=0.42bp] (274.000,250.022) circle[radius=0.42bp] (274.667,249.587) circle[radius=0.42bp] (275.333,260.054) circle[radius=0.42bp] (276.000,256.428) circle[radius=0.42bp] (276.667,258.426) circle[radius=0.42bp] (277.333,258.172) circle[radius=0.42bp] (278.000,264.694) circle[radius=0.42bp] (278.667,260.533) circle[radius=0.42bp] (279.333,256.745) circle[radius=0.42bp] (280.000,256.680) circle[radius=0.42bp] (280.667,243.249) circle[radius=0.42bp] (281.333,251.150) circle[radius=0.42bp] (282.000,248.727) circle[radius=0.42bp] (282.667,253.141) circle[radius=0.42bp] (283.333,248.900) circle[radius=0.42bp] (284.000,252.835) circle[radius=0.42bp] (284.667,257.635) circle[radius=0.42bp] (285.333,269.434) circle[radius=0.42bp] (286.000,259.893) circle[radius=0.42bp] (286.667,255.307) circle[radius=0.42bp] (287.333,251.151) circle[radius=0.42bp] (288.000,254.537) circle[radius=0.42bp] (288.667,242.928) circle[radius=0.42bp] (289.333,254.353) circle[radius=0.42bp] (290.000,252.574) circle[radius=0.42bp] (290.667,248.505) circle[radius=0.42bp] (291.333,254.859) circle[radius=0.42bp] (292.000,250.186) circle[radius=0.42bp] (292.667,259.363) circle[radius=0.42bp] (293.333,263.335) circle[radius=0.42bp] (294.000,261.562) circle[radius=0.42bp] (294.667,262.201) circle[radius=0.42bp] (295.333,260.012) circle[radius=0.42bp] (296.000,249.689) circle[radius=0.42bp] (296.667,247.748) circle[radius=0.42bp] (297.333,244.196) circle[radius=0.42bp] (298.000,252.935) circle[radius=0.42bp] (298.667,245.577) circle[radius=0.42bp] (299.333,256.166) circle[radius=0.42bp] (300.000,251.848) circle[radius=0.42bp] (300.667,261.776) circle[radius=0.42bp] (301.333,263.732) circle[radius=0.42bp] (302.000,259.125) circle[radius=0.42bp] (302.667,253.957) circle[radius=0.42bp] (303.333,254.425) circle[radius=0.42bp] (304.000,254.411) circle[radius=0.42bp] (304.667,255.214) circle[radius=0.42bp] (305.333,254.667) circle[radius=0.42bp] (306.000,248.052) circle[radius=0.42bp] (306.667,248.293) circle[radius=0.42bp] (307.333,255.688) circle[radius=0.42bp] (308.000,262.135) circle[radius=0.42bp] (308.667,256.704) circle[radius=0.42bp] (309.333,265.848) circle[radius=0.42bp] (310.000,262.865) circle[radius=0.42bp] (310.667,255.551) circle[radius=0.42bp] (311.333,263.789) circle[radius=0.42bp] (312.000,250.636) circle[radius=0.42bp] (312.667,252.264) circle[radius=0.42bp] (313.333,248.800) circle[radius=0.42bp] (314.000,250.299) circle[radius=0.42bp] (314.667,251.352) circle[radius=0.42bp] (315.333,252.354) circle[radius=0.42bp] (316.000,260.908) circle[radius=0.42bp] (316.667,254.576) circle[radius=0.42bp] (317.333,258.056) circle[radius=0.42bp] (318.000,256.818) circle[radius=0.42bp] (318.667,253.439) circle[radius=0.42bp] (319.333,251.991) circle[radius=0.42bp] (320.000,248.859) circle[radius=0.42bp] (320.667,249.711) circle[radius=0.42bp] (321.333,254.766) circle[radius=0.42bp] (322.000,250.529) circle[radius=0.42bp] (322.667,255.126) circle[radius=0.42bp] (323.333,260.576) circle[radius=0.42bp] (324.000,256.959) circle[radius=0.42bp] (324.667,262.303) circle[radius=0.42bp] (325.333,264.875) circle[radius=0.42bp] (326.000,264.146) circle[radius=0.42bp] (326.667,256.920) circle[radius=0.42bp] (327.333,253.657) circle[radius=0.42bp] (328.000,247.001) circle[radius=0.42bp] (328.667,250.059) circle[radius=0.42bp] (329.333,245.178) circle[radius=0.42bp] (330.000,244.559) circle[radius=0.42bp] (330.667,247.670) circle[radius=0.42bp] (331.333,264.236) circle[radius=0.42bp] (332.000,253.191) circle[radius=0.42bp] (332.667,258.776) circle[radius=0.42bp] (333.333,259.230) circle[radius=0.42bp] (334.000,258.653) circle[radius=0.42bp] (334.667,256.476) circle[radius=0.42bp];
\draw[color=tg000000,line width=0.7bp,line cap=round] (319.333,255.050) -- (320.000,252.137) -- (320.667,250.005) -- (321.333,249.224) -- (322.000,250.005) -- (322.667,252.137) -- (323.333,255.050) -- (324.000,257.963) -- (324.667,260.095) -- (325.333,260.876) -- (326.000,260.095) -- (326.667,257.963) -- (327.333,255.050) -- (328.000,252.137) -- (328.667,250.005) -- (329.333,249.224) -- (330.000,250.005) -- (330.667,252.137) -- (331.333,255.050) -- (332.000,257.963) -- (332.667,260.095) -- (333.333,260.876) -- (334.000,260.095) -- (334.667,257.963);
\draw[color=tg0000ff,line width=0.3bp,line cap=round] (318.783,258.874) -- (319.883,259.974);
\draw[color=tg0000ff,line width=0.3bp,line cap=round] (318.783,259.974) -- (319.883,258.874);
\draw[color=tg0000ff,line width=0.3bp,line cap=round] (319.450,253.347) -- (320.550,254.447);
\draw[color=tg0000ff,line width=0.3bp,line cap=round] (319.450,254.447) -- (320.550,253.347);
\draw[color=tg0000ff,line width=0.3bp,line cap=round] (320.117,249.640) -- (321.217,250.740);
\draw[color=tg0000ff,line width=0.3bp,line cap=round] (320.117,250.740) -- (321.217,249.640);
\draw[color=tg0000ff,line width=0.3bp,line cap=round] (320.783,251.675) -- (321.883,252.775);
\draw[color=tg0000ff,line width=0.3bp,line cap=round] (320.783,252.775) -- (321.883,251.675);
\draw[color=tg0000ff,line width=0.3bp,line cap=round] (321.450,246.327) -- (322.550,247.427);
\draw[color=tg0000ff,line width=0.3bp,line cap=round] (321.450,247.427) -- (322.550,246.327);
\draw[color=tg0000ff,line width=0.3bp,line cap=round] (322.117,250.455) -- (323.217,251.555);
\draw[color=tg0000ff,line width=0.3bp,line cap=round] (322.117,251.555) -- (323.217,250.455);
\draw[color=tg0000ff,line width=0.3bp,line cap=round] (322.783,254.420) -- (323.883,255.520);
\draw[color=tg0000ff,line width=0.3bp,line cap=round] (322.783,255.520) -- (323.883,254.420);
\draw[color=tg0000ff,line width=0.3bp,line cap=round] (323.450,256.567) -- (324.550,257.667);
\draw[color=tg0000ff,line width=0.3bp,line cap=round] (323.450,257.667) -- (324.550,256.567);
\draw[color=tg0000ff,line width=0.3bp,line cap=round] (324.117,257.338) -- (325.217,258.438);
\draw[color=tg0000ff,line width=0.3bp,line cap=round] (324.117,258.438) -- (325.217,257.338);
\draw[color=tg0000ff,line width=0.3bp,line cap=round] (324.783,260.056) -- (325.883,261.156);
\draw[color=tg0000ff,line width=0.3bp,line cap=round] (324.783,261.156) -- (325.883,260.056);
\draw[color=tg0000ff,line width=0.3bp,line cap=round] (325.450,261.009) -- (326.550,262.109);
\draw[color=tg0000ff,line width=0.3bp,line cap=round] (325.450,262.109) -- (326.550,261.009);
\draw[color=tg0000ff,line width=0.3bp,line cap=round] (326.117,256.146) -- (327.217,257.246);
\draw[color=tg0000ff,line width=0.3bp,line cap=round] (326.117,257.246) -- (327.217,256.146);
\draw[color=tg0000ff,line width=0.3bp,line cap=round] (326.783,256.813) -- (327.883,257.913);
\draw[color=tg0000ff,line width=0.3bp,line cap=round] (326.783,257.913) -- (327.883,256.813);
\draw[color=tg0000ff,line width=0.3bp,line cap=round] (327.450,253.074) -- (328.550,254.174);
\draw[color=tg0000ff,line width=0.3bp,line cap=round] (327.450,254.174) -- (328.550,253.074);
\draw[color=tg0000ff,line width=0.3bp,line cap=round] (328.117,249.155) -- (329.217,250.255);
\draw[color=tg0000ff,line width=0.3bp,line cap=round] (328.117,250.255) -- (329.217,249.155);
\draw[color=tg0000ff,line width=0.3bp,line cap=round] (328.783,251.323) -- (329.883,252.423);
\draw[color=tg0000ff,line width=0.3bp,line cap=round] (328.783,252.423) -- (329.883,251.323);
\draw[color=tg0000ff,line width=0.3bp,line cap=round] (329.450,249.717) -- (330.550,250.817);
\draw[color=tg0000ff,line width=0.3bp,line cap=round] (329.450,250.817) -- (330.550,249.717);
\draw[color=tg0000ff,line width=0.3bp,line cap=round] (330.117,250.742) -- (331.217,251.842);
\draw[color=tg0000ff,line width=0.3bp,line cap=round] (330.117,251.842) -- (331.217,250.742);
\draw[color=tg0000ff,line width=0.3bp,line cap=round] (330.783,255.123) -- (331.883,256.223);
\draw[color=tg0000ff,line width=0.3bp,line cap=round] (330.783,256.223) -- (331.883,255.123);
\draw[color=tg0000ff,line width=0.3bp,line cap=round] (331.450,254.386) -- (332.550,255.486);
\draw[color=tg0000ff,line width=0.3bp,line cap=round] (331.450,255.486) -- (332.550,254.386);
\draw[color=tg0000ff,line width=0.3bp,line cap=round] (332.117,256.000) -- (333.217,257.100);
\draw[color=tg0000ff,line width=0.3bp,line cap=round] (332.117,257.100) -- (333.217,256.000);
\draw[color=tg0000ff,line width=0.3bp,line cap=round] (332.783,258.109) -- (333.883,259.209);
\draw[color=tg0000ff,line width=0.3bp,line cap=round] (332.783,259.209) -- (333.883,258.109);
\draw[color=tg0000ff,line width=0.3bp,line cap=round] (333.450,256.842) -- (334.550,257.942);
\draw[color=tg0000ff,line width=0.3bp,line cap=round] (333.450,257.942) -- (334.550,256.842);
\draw[color=tg0000ff,line width=0.3bp,line cap=round] (334.117,256.506) -- (335.217,257.606);
\draw[color=tg0000ff,line width=0.3bp,line cap=round] (334.117,257.606) -- (335.217,256.506);
\node[anchor=west,inner sep=0pt,outer sep=0pt,text=tg626b73,font={\sffamily\fontsize{6.2bp}{7.44bp}\selectfont}] at (346.000,234.058) {None found};
\draw[color=tgeceef0,line width=0.4bp,line cap=round] (0.000,274.024) -- (448.000,274.024);
\node[anchor=west,inner sep=0pt,outer sep=0pt,text=tg18314b,font={\sffamily\fontsize{7.2bp}{8.64bp}\selectfont\bfseries}] at (0.000,280.844) {Trend + seasonal cycle};
\node[anchor=west,inner sep=0pt,outer sep=0pt,text=tg18314b,font={\sffamily\fontsize{6.2bp}{7.44bp}\selectfont}] at (0.000,289.284) {\ensuremath{\checkmark} Both the line and the repeating shape.};
\node[anchor=west,inner sep=0pt,outer sep=0pt,text=tg18314b,font={\sffamily\fontsize{6.2bp}{7.44bp}\selectfont}] at (0.000,297.124) {\ensuremath{\times} The noise around them.};
\path[fill=tg8fd18f,draw=none,line width=0.5bp] (168.000,276.524) rectangle (205.969,285.124);
\node[anchor=center,inner sep=0pt,outer sep=0pt,text=tg0b0b0b,font={\sffamily\fontsize{6.6bp}{7.92bp}\selectfont\bfseries}] at (186.984,280.824) {78  Good};
\draw[color=tgc9cdd1,line width=0.45bp,line cap=round] (168.000,317.210) -- (248.000,317.210);
\draw[color=tg9da7b0,line width=0.5bp,line cap=round,dash pattern=on 2bp off 2bp] (232.000,287.524) -- (232.000,317.210);
\path[fill=tg808080,draw=none] (168.333,312.544) circle[radius=0.42bp] (169.000,308.401) circle[radius=0.42bp] (169.667,307.505) circle[radius=0.42bp] (170.333,305.741) circle[radius=0.42bp] (171.000,306.917) circle[radius=0.42bp] (171.667,308.451) circle[radius=0.42bp] (172.333,310.376) circle[radius=0.42bp] (173.000,312.916) circle[radius=0.42bp] (173.667,314.792) circle[radius=0.42bp] (174.333,315.620) circle[radius=0.42bp] (175.000,313.534) circle[radius=0.42bp] (175.667,312.424) circle[radius=0.42bp] (176.333,310.226) circle[radius=0.42bp] (177.000,305.853) circle[radius=0.42bp] (177.667,305.162) circle[radius=0.42bp] (178.333,305.142) circle[radius=0.42bp] (179.000,305.377) circle[radius=0.42bp] (179.667,307.012) circle[radius=0.42bp] (180.333,309.545) circle[radius=0.42bp] (181.000,310.100) circle[radius=0.42bp] (181.667,312.852) circle[radius=0.42bp] (182.333,312.492) circle[radius=0.42bp] (183.000,312.429) circle[radius=0.42bp] (183.667,309.121) circle[radius=0.42bp] (184.333,308.118) circle[radius=0.42bp] (185.000,305.659) circle[radius=0.42bp] (185.667,305.415) circle[radius=0.42bp] (186.333,302.807) circle[radius=0.42bp] (187.000,301.845) circle[radius=0.42bp] (187.667,305.608) circle[radius=0.42bp] (188.333,306.824) circle[radius=0.42bp] (189.000,307.623) circle[radius=0.42bp] (189.667,310.439) circle[radius=0.42bp] (190.333,310.893) circle[radius=0.42bp] (191.000,311.700) circle[radius=0.42bp] (191.667,307.950) circle[radius=0.42bp] (192.333,306.555) circle[radius=0.42bp] (193.000,304.154) circle[radius=0.42bp] (193.667,302.090) circle[radius=0.42bp] (194.333,301.713) circle[radius=0.42bp] (195.000,301.351) circle[radius=0.42bp] (195.667,303.331) circle[radius=0.42bp] (196.333,305.303) circle[radius=0.42bp] (197.000,306.854) circle[radius=0.42bp] (197.667,307.567) circle[radius=0.42bp] (198.333,310.259) circle[radius=0.42bp] (199.000,309.242) circle[radius=0.42bp] (199.667,307.628) circle[radius=0.42bp] (200.333,303.711) circle[radius=0.42bp] (201.000,303.280) circle[radius=0.42bp] (201.667,300.620) circle[radius=0.42bp] (202.333,297.805) circle[radius=0.42bp] (203.000,299.966) circle[radius=0.42bp] (203.667,300.965) circle[radius=0.42bp] (204.333,303.558) circle[radius=0.42bp] (205.000,306.607) circle[radius=0.42bp] (205.667,307.311) circle[radius=0.42bp] (206.333,308.127) circle[radius=0.42bp] (207.000,306.820) circle[radius=0.42bp] (207.667,304.049) circle[radius=0.42bp] (208.333,302.423) circle[radius=0.42bp] (209.000,300.613) circle[radius=0.42bp] (209.667,299.214) circle[radius=0.42bp] (210.333,296.686) circle[radius=0.42bp] (211.000,298.232) circle[radius=0.42bp] (211.667,299.407) circle[radius=0.42bp] (212.333,299.766) circle[radius=0.42bp] (213.000,303.745) circle[radius=0.42bp] (213.667,305.385) circle[radius=0.42bp] (214.333,306.350) circle[radius=0.42bp] (215.000,305.123) circle[radius=0.42bp] (215.667,304.263) circle[radius=0.42bp] (216.333,301.497) circle[radius=0.42bp] (217.000,298.176) circle[radius=0.42bp] (217.667,295.038) circle[radius=0.42bp] (218.333,294.765) circle[radius=0.42bp] (219.000,296.316) circle[radius=0.42bp] (219.667,298.360) circle[radius=0.42bp] (220.333,298.587) circle[radius=0.42bp] (221.000,302.441) circle[radius=0.42bp] (221.667,303.206) circle[radius=0.42bp] (222.333,305.229) circle[radius=0.42bp] (223.000,302.576) circle[radius=0.42bp] (223.667,301.343) circle[radius=0.42bp] (224.333,300.084) circle[radius=0.42bp] (225.000,297.040) circle[radius=0.42bp] (225.667,294.246) circle[radius=0.42bp] (226.333,294.694) circle[radius=0.42bp] (227.000,294.559) circle[radius=0.42bp] (227.667,295.252) circle[radius=0.42bp] (228.333,298.293) circle[radius=0.42bp] (229.000,300.495) circle[radius=0.42bp] (229.667,302.474) circle[radius=0.42bp] (230.333,301.614) circle[radius=0.42bp] (231.000,301.996) circle[radius=0.42bp] (231.667,299.002) circle[radius=0.42bp] (232.333,296.927) circle[radius=0.42bp] (233.000,293.753) circle[radius=0.42bp] (233.667,293.984) circle[radius=0.42bp] (234.333,292.765) circle[radius=0.42bp] (235.000,294.338) circle[radius=0.42bp] (235.667,295.379) circle[radius=0.42bp] (236.333,295.163) circle[radius=0.42bp] (237.000,298.236) circle[radius=0.42bp] (237.667,297.970) circle[radius=0.42bp] (238.333,299.816) circle[radius=0.42bp] (239.000,299.963) circle[radius=0.42bp] (239.667,297.333) circle[radius=0.42bp] (240.333,293.757) circle[radius=0.42bp] (241.000,293.221) circle[radius=0.42bp] (241.667,291.215) circle[radius=0.42bp] (242.333,289.114) circle[radius=0.42bp] (243.000,290.410) circle[radius=0.42bp] (243.667,292.258) circle[radius=0.42bp] (244.333,294.893) circle[radius=0.42bp] (245.000,296.903) circle[radius=0.42bp] (245.667,298.682) circle[radius=0.42bp] (246.333,298.852) circle[radius=0.42bp] (247.000,297.291) circle[radius=0.42bp] (247.667,296.319) circle[radius=0.42bp];
\draw[color=tg000000,line width=0.7bp,line cap=round] (232.333,297.127) -- (233.000,294.707) -- (233.667,292.894) -- (234.333,292.135) -- (235.000,292.592) -- (235.667,294.102) -- (236.333,296.219) -- (237.000,298.337) -- (237.667,299.846) -- (238.333,300.303) -- (239.000,299.544) -- (239.667,297.732) -- (240.333,295.312) -- (241.000,292.891) -- (241.667,291.079) -- (242.333,290.320) -- (243.000,290.777) -- (243.667,292.286) -- (244.333,294.404) -- (245.000,296.522) -- (245.667,298.031) -- (246.333,298.488) -- (247.000,297.729) -- (247.667,295.917);
\draw[color=tg0000ff,line width=0.3bp,line cap=round] (231.783,296.977) -- (232.883,298.077);
\draw[color=tg0000ff,line width=0.3bp,line cap=round] (231.783,298.077) -- (232.883,296.977);
\draw[color=tg0000ff,line width=0.3bp,line cap=round] (232.450,294.129) -- (233.550,295.229);
\draw[color=tg0000ff,line width=0.3bp,line cap=round] (232.450,295.229) -- (233.550,294.129);
\draw[color=tg0000ff,line width=0.3bp,line cap=round] (233.117,292.223) -- (234.217,293.323);
\draw[color=tg0000ff,line width=0.3bp,line cap=round] (233.117,293.323) -- (234.217,292.223);
\draw[color=tg0000ff,line width=0.3bp,line cap=round] (233.783,292.117) -- (234.883,293.217);
\draw[color=tg0000ff,line width=0.3bp,line cap=round] (233.783,293.217) -- (234.883,292.117);
\draw[color=tg0000ff,line width=0.3bp,line cap=round] (234.450,292.676) -- (235.550,293.776);
\draw[color=tg0000ff,line width=0.3bp,line cap=round] (234.450,293.776) -- (235.550,292.676);
\draw[color=tg0000ff,line width=0.3bp,line cap=round] (235.117,293.468) -- (236.217,294.568);
\draw[color=tg0000ff,line width=0.3bp,line cap=round] (235.117,294.568) -- (236.217,293.468);
\draw[color=tg0000ff,line width=0.3bp,line cap=round] (235.783,296.092) -- (236.883,297.192);
\draw[color=tg0000ff,line width=0.3bp,line cap=round] (235.783,297.192) -- (236.883,296.092);
\draw[color=tg0000ff,line width=0.3bp,line cap=round] (236.450,297.995) -- (237.550,299.095);
\draw[color=tg0000ff,line width=0.3bp,line cap=round] (236.450,299.095) -- (237.550,297.995);
\draw[color=tg0000ff,line width=0.3bp,line cap=round] (237.117,299.712) -- (238.217,300.812);
\draw[color=tg0000ff,line width=0.3bp,line cap=round] (237.117,300.812) -- (238.217,299.712);
\draw[color=tg0000ff,line width=0.3bp,line cap=round] (237.783,299.555) -- (238.883,300.655);
\draw[color=tg0000ff,line width=0.3bp,line cap=round] (237.783,300.655) -- (238.883,299.555);
\draw[color=tg0000ff,line width=0.3bp,line cap=round] (238.450,299.225) -- (239.550,300.325);
\draw[color=tg0000ff,line width=0.3bp,line cap=round] (238.450,300.325) -- (239.550,299.225);
\draw[color=tg0000ff,line width=0.3bp,line cap=round] (239.117,297.665) -- (240.217,298.765);
\draw[color=tg0000ff,line width=0.3bp,line cap=round] (239.117,298.765) -- (240.217,297.665);
\draw[color=tg0000ff,line width=0.3bp,line cap=round] (239.783,294.923) -- (240.883,296.023);
\draw[color=tg0000ff,line width=0.3bp,line cap=round] (239.783,296.023) -- (240.883,294.923);
\draw[color=tg0000ff,line width=0.3bp,line cap=round] (240.450,292.432) -- (241.550,293.532);
\draw[color=tg0000ff,line width=0.3bp,line cap=round] (240.450,293.532) -- (241.550,292.432);
\draw[color=tg0000ff,line width=0.3bp,line cap=round] (241.117,290.755) -- (242.217,291.855);
\draw[color=tg0000ff,line width=0.3bp,line cap=round] (241.117,291.855) -- (242.217,290.755);
\draw[color=tg0000ff,line width=0.3bp,line cap=round] (241.783,290.098) -- (242.883,291.198);
\draw[color=tg0000ff,line width=0.3bp,line cap=round] (241.783,291.198) -- (242.883,290.098);
\draw[color=tg0000ff,line width=0.3bp,line cap=round] (242.450,290.521) -- (243.550,291.621);
\draw[color=tg0000ff,line width=0.3bp,line cap=round] (242.450,291.621) -- (243.550,290.521);
\draw[color=tg0000ff,line width=0.3bp,line cap=round] (243.117,292.126) -- (244.217,293.226);
\draw[color=tg0000ff,line width=0.3bp,line cap=round] (243.117,293.226) -- (244.217,292.126);
\draw[color=tg0000ff,line width=0.3bp,line cap=round] (243.783,294.350) -- (244.883,295.450);
\draw[color=tg0000ff,line width=0.3bp,line cap=round] (243.783,295.450) -- (244.883,294.350);
\draw[color=tg0000ff,line width=0.3bp,line cap=round] (244.450,296.611) -- (245.550,297.711);
\draw[color=tg0000ff,line width=0.3bp,line cap=round] (244.450,297.711) -- (245.550,296.611);
\draw[color=tg0000ff,line width=0.3bp,line cap=round] (245.117,297.698) -- (246.217,298.798);
\draw[color=tg0000ff,line width=0.3bp,line cap=round] (245.117,298.798) -- (246.217,297.698);
\draw[color=tg0000ff,line width=0.3bp,line cap=round] (245.783,298.290) -- (246.883,299.390);
\draw[color=tg0000ff,line width=0.3bp,line cap=round] (245.783,299.390) -- (246.883,298.290);
\draw[color=tg0000ff,line width=0.3bp,line cap=round] (246.450,297.181) -- (247.550,298.281);
\draw[color=tg0000ff,line width=0.3bp,line cap=round] (246.450,298.281) -- (247.550,297.181);
\draw[color=tg0000ff,line width=0.3bp,line cap=round] (247.117,295.696) -- (248.217,296.796);
\draw[color=tg0000ff,line width=0.3bp,line cap=round] (247.117,296.796) -- (248.217,295.696);
\path[fill=tg8fd18f,draw=none,line width=0.5bp] (255.000,276.524) rectangle (292.969,285.124);
\node[anchor=center,inner sep=0pt,outer sep=0pt,text=tg0b0b0b,font={\sffamily\fontsize{6.6bp}{7.92bp}\selectfont\bfseries}] at (273.984,280.824) {71  Good};
\draw[color=tgc9cdd1,line width=0.45bp,line cap=round] (255.000,317.210) -- (335.000,317.210);
\draw[color=tg9da7b0,line width=0.5bp,line cap=round,dash pattern=on 2bp off 2bp] (319.000,287.524) -- (319.000,317.210);
\path[fill=tg808080,draw=none] (255.333,313.359) circle[radius=0.42bp] (256.000,305.767) circle[radius=0.42bp] (256.667,307.278) circle[radius=0.42bp] (257.333,303.380) circle[radius=0.42bp] (258.000,306.098) circle[radius=0.42bp] (258.667,307.410) circle[radius=0.42bp] (259.333,308.508) circle[radius=0.42bp] (260.000,311.613) circle[radius=0.42bp] (260.667,314.036) circle[radius=0.42bp] (261.333,315.620) circle[radius=0.42bp] (262.000,310.674) circle[radius=0.42bp] (262.667,311.487) circle[radius=0.42bp] (263.333,310.238) circle[radius=0.42bp] (264.000,301.898) circle[radius=0.42bp] (264.667,304.077) circle[radius=0.42bp] (265.333,305.867) circle[radius=0.42bp] (266.000,305.515) circle[radius=0.42bp] (266.667,307.156) circle[radius=0.42bp] (267.333,310.239) circle[radius=0.42bp] (268.000,306.869) circle[radius=0.42bp] (268.667,312.151) circle[radius=0.42bp] (269.333,309.858) circle[radius=0.42bp] (270.000,311.510) circle[radius=0.42bp] (270.667,305.155) circle[radius=0.42bp] (271.333,307.803) circle[radius=0.42bp] (272.000,305.703) circle[radius=0.42bp] (272.667,309.341) circle[radius=0.42bp] (273.333,302.692) circle[radius=0.42bp] (274.000,298.437) circle[radius=0.42bp] (274.667,307.017) circle[radius=0.42bp] (275.333,305.802) circle[radius=0.42bp] (276.000,303.231) circle[radius=0.42bp] (276.667,308.722) circle[radius=0.42bp] (277.333,309.086) circle[radius=0.42bp] (278.000,313.574) circle[radius=0.42bp] (278.667,305.777) circle[radius=0.42bp] (279.333,307.147) circle[radius=0.42bp] (280.000,305.234) circle[radius=0.42bp] (280.667,302.936) circle[radius=0.42bp] (281.333,303.564) circle[radius=0.42bp] (282.000,301.266) circle[radius=0.42bp] (282.667,304.032) circle[radius=0.42bp] (283.333,305.282) circle[radius=0.42bp] (284.000,305.162) circle[radius=0.42bp] (284.667,303.793) circle[radius=0.42bp] (285.333,311.455) circle[radius=0.42bp] (286.000,309.995) circle[radius=0.42bp] (286.667,309.166) circle[radius=0.42bp] (287.333,302.311) circle[radius=0.42bp] (288.000,306.824) circle[radius=0.42bp] (288.667,302.582) circle[radius=0.42bp] (289.333,295.258) circle[radius=0.42bp] (290.000,301.187) circle[radius=0.42bp] (290.667,300.755) circle[radius=0.42bp] (291.333,304.030) circle[radius=0.42bp] (292.000,308.794) circle[radius=0.42bp] (292.667,307.400) circle[radius=0.42bp] (293.333,308.944) circle[radius=0.42bp] (294.000,306.536) circle[radius=0.42bp] (294.667,301.932) circle[radius=0.42bp] (295.333,302.549) circle[radius=0.42bp] (296.000,302.567) circle[radius=0.42bp] (296.667,302.435) circle[radius=0.42bp] (297.333,296.049) circle[radius=0.42bp] (298.000,299.972) circle[radius=0.42bp] (298.667,300.113) circle[radius=0.42bp] (299.333,296.106) circle[radius=0.42bp] (300.000,303.902) circle[radius=0.42bp] (300.667,305.556) circle[radius=0.42bp] (301.333,307.586) circle[radius=0.42bp] (302.000,305.442) circle[radius=0.42bp] (302.667,307.072) circle[radius=0.42bp] (303.333,303.969) circle[radius=0.42bp] (304.000,299.059) circle[radius=0.42bp] (304.667,293.257) circle[radius=0.42bp] (305.333,294.224) circle[radius=0.42bp] (306.000,298.165) circle[radius=0.42bp] (306.667,301.137) circle[radius=0.42bp] (307.333,296.698) circle[radius=0.42bp] (308.000,304.087) circle[radius=0.42bp] (308.667,302.890) circle[radius=0.42bp] (309.333,308.370) circle[radius=0.42bp] (310.000,301.576) circle[radius=0.42bp] (310.667,301.986) circle[radius=0.42bp] (311.333,303.800) circle[radius=0.42bp] (312.000,299.792) circle[radius=0.42bp] (312.667,295.114) circle[radius=0.42bp] (313.333,298.431) circle[radius=0.42bp] (314.000,296.873) circle[radius=0.42bp] (314.667,295.442) circle[radius=0.42bp] (315.333,300.179) circle[radius=0.42bp] (316.000,302.180) circle[radius=0.42bp] (316.667,304.944) circle[radius=0.42bp] (317.333,301.020) circle[radius=0.42bp] (318.000,304.123) circle[radius=0.42bp] (318.667,298.794) circle[radius=0.42bp] (319.333,297.945) circle[radius=0.42bp] (320.000,293.514) circle[radius=0.42bp] (320.667,298.699) circle[radius=0.42bp] (321.333,296.580) circle[radius=0.42bp] (322.000,300.592) circle[radius=0.42bp] (322.667,300.296) circle[radius=0.42bp] (323.333,294.411) circle[radius=0.42bp] (324.000,299.255) circle[radius=0.42bp] (324.667,294.694) circle[radius=0.42bp] (325.333,299.597) circle[radius=0.42bp] (326.000,301.933) circle[radius=0.42bp] (326.667,297.788) circle[radius=0.42bp] (327.333,292.046) circle[radius=0.42bp] (328.000,296.219) circle[radius=0.42bp] (328.667,294.108) circle[radius=0.42bp] (329.333,289.114) circle[radius=0.42bp] (330.000,292.223) circle[radius=0.42bp] (330.667,294.557) circle[radius=0.42bp] (331.333,297.971) circle[radius=0.42bp] (332.000,299.348) circle[radius=0.42bp] (332.667,301.454) circle[radius=0.42bp] (333.333,300.893) circle[radius=0.42bp] (334.000,297.660) circle[radius=0.42bp] (334.667,298.922) circle[radius=0.42bp];
\draw[color=tg000000,line width=0.7bp,line cap=round] (319.333,298.596) -- (320.000,296.623) -- (320.667,295.145) -- (321.333,294.526) -- (322.000,294.899) -- (322.667,296.129) -- (323.333,297.856) -- (324.000,299.582) -- (324.667,300.813) -- (325.333,301.185) -- (326.000,300.566) -- (326.667,299.089) -- (327.333,297.116) -- (328.000,295.143) -- (328.667,293.666) -- (329.333,293.047) -- (330.000,293.419) -- (330.667,294.650) -- (331.333,296.376) -- (332.000,298.102) -- (332.667,299.333) -- (333.333,299.706) -- (334.000,299.087) -- (334.667,297.609);
\draw[color=tg0000ff,line width=0.3bp,line cap=round] (318.783,300.521) -- (319.883,301.621);
\draw[color=tg0000ff,line width=0.3bp,line cap=round] (318.783,301.621) -- (319.883,300.521);
\draw[color=tg0000ff,line width=0.3bp,line cap=round] (319.450,297.214) -- (320.550,298.314);
\draw[color=tg0000ff,line width=0.3bp,line cap=round] (319.450,298.314) -- (320.550,297.214);
\draw[color=tg0000ff,line width=0.3bp,line cap=round] (320.117,295.618) -- (321.217,296.718);
\draw[color=tg0000ff,line width=0.3bp,line cap=round] (320.117,296.718) -- (321.217,295.618);
\draw[color=tg0000ff,line width=0.3bp,line cap=round] (320.783,297.073) -- (321.883,298.173);
\draw[color=tg0000ff,line width=0.3bp,line cap=round] (320.783,298.173) -- (321.883,297.073);
\draw[color=tg0000ff,line width=0.3bp,line cap=round] (321.450,297.438) -- (322.550,298.538);
\draw[color=tg0000ff,line width=0.3bp,line cap=round] (321.450,298.538) -- (322.550,297.438);
\draw[color=tg0000ff,line width=0.3bp,line cap=round] (322.117,295.006) -- (323.217,296.106);
\draw[color=tg0000ff,line width=0.3bp,line cap=round] (322.117,296.106) -- (323.217,295.006);
\draw[color=tg0000ff,line width=0.3bp,line cap=round] (322.783,299.070) -- (323.883,300.170);
\draw[color=tg0000ff,line width=0.3bp,line cap=round] (322.783,300.170) -- (323.883,299.070);
\draw[color=tg0000ff,line width=0.3bp,line cap=round] (323.450,298.461) -- (324.550,299.561);
\draw[color=tg0000ff,line width=0.3bp,line cap=round] (323.450,299.561) -- (324.550,298.461);
\draw[color=tg0000ff,line width=0.3bp,line cap=round] (324.117,301.006) -- (325.217,302.106);
\draw[color=tg0000ff,line width=0.3bp,line cap=round] (324.117,302.106) -- (325.217,301.006);
\draw[color=tg0000ff,line width=0.3bp,line cap=round] (324.783,299.039) -- (325.883,300.139);
\draw[color=tg0000ff,line width=0.3bp,line cap=round] (324.783,300.139) -- (325.883,299.039);
\draw[color=tg0000ff,line width=0.3bp,line cap=round] (325.450,301.922) -- (326.550,303.022);
\draw[color=tg0000ff,line width=0.3bp,line cap=round] (325.450,303.022) -- (326.550,301.922);
\draw[color=tg0000ff,line width=0.3bp,line cap=round] (326.117,301.236) -- (327.217,302.336);
\draw[color=tg0000ff,line width=0.3bp,line cap=round] (326.117,302.336) -- (327.217,301.236);
\draw[color=tg0000ff,line width=0.3bp,line cap=round] (326.783,298.115) -- (327.883,299.215);
\draw[color=tg0000ff,line width=0.3bp,line cap=round] (326.783,299.215) -- (327.883,298.115);
\draw[color=tg0000ff,line width=0.3bp,line cap=round] (327.450,296.528) -- (328.550,297.628);
\draw[color=tg0000ff,line width=0.3bp,line cap=round] (327.450,297.628) -- (328.550,296.528);
\draw[color=tg0000ff,line width=0.3bp,line cap=round] (328.117,295.490) -- (329.217,296.590);
\draw[color=tg0000ff,line width=0.3bp,line cap=round] (328.117,296.590) -- (329.217,295.490);
\draw[color=tg0000ff,line width=0.3bp,line cap=round] (328.783,294.844) -- (329.883,295.944);
\draw[color=tg0000ff,line width=0.3bp,line cap=round] (328.783,295.944) -- (329.883,294.844);
\draw[color=tg0000ff,line width=0.3bp,line cap=round] (329.450,294.812) -- (330.550,295.912);
\draw[color=tg0000ff,line width=0.3bp,line cap=round] (329.450,295.912) -- (330.550,294.812);
\draw[color=tg0000ff,line width=0.3bp,line cap=round] (330.117,295.918) -- (331.217,297.018);
\draw[color=tg0000ff,line width=0.3bp,line cap=round] (330.117,297.018) -- (331.217,295.918);
\draw[color=tg0000ff,line width=0.3bp,line cap=round] (330.783,297.302) -- (331.883,298.402);
\draw[color=tg0000ff,line width=0.3bp,line cap=round] (330.783,298.402) -- (331.883,297.302);
\draw[color=tg0000ff,line width=0.3bp,line cap=round] (331.450,299.355) -- (332.550,300.455);
\draw[color=tg0000ff,line width=0.3bp,line cap=round] (331.450,300.455) -- (332.550,299.355);
\draw[color=tg0000ff,line width=0.3bp,line cap=round] (332.117,298.590) -- (333.217,299.690);
\draw[color=tg0000ff,line width=0.3bp,line cap=round] (332.117,299.690) -- (333.217,298.590);
\draw[color=tg0000ff,line width=0.3bp,line cap=round] (332.783,300.357) -- (333.883,301.457);
\draw[color=tg0000ff,line width=0.3bp,line cap=round] (332.783,301.457) -- (333.883,300.357);
\draw[color=tg0000ff,line width=0.3bp,line cap=round] (333.450,298.342) -- (334.550,299.442);
\draw[color=tg0000ff,line width=0.3bp,line cap=round] (333.450,299.442) -- (334.550,298.342);
\draw[color=tg0000ff,line width=0.3bp,line cap=round] (334.117,299.803) -- (335.217,300.903);
\draw[color=tg0000ff,line width=0.3bp,line cap=round] (334.117,300.903) -- (335.217,299.803);
\node[anchor=west,inner sep=0pt,outer sep=0pt,text=tg626b73,font={\sffamily\fontsize{6.2bp}{7.44bp}\selectfont}] at (346.000,280.244) {None found};
\draw[color=tgeceef0,line width=0.4bp,line cap=round] (0.000,320.210) -- (448.000,320.210);
\node[anchor=west,inner sep=0pt,outer sep=0pt,text=tg18314b,font={\sffamily\fontsize{7.2bp}{8.64bp}\selectfont\bfseries}] at (0.000,327.030) {Sudden level shift};
\node[anchor=west,inner sep=0pt,outer sep=0pt,text=tg18314b,font={\sffamily\fontsize{6.2bp}{7.44bp}\selectfont}] at (0.000,335.470) {\ensuremath{\checkmark} The new level, from the history since the jump.};
\node[anchor=west,inner sep=0pt,outer sep=0pt,text=tg18314b,font={\sffamily\fontsize{6.2bp}{7.44bp}\selectfont}] at (0.000,343.310) {\ensuremath{\times} The noise. Older history sits at the old level and};
\node[anchor=west,inner sep=0pt,outer sep=0pt,text=tg18314b,font={\sffamily\fontsize{6.2bp}{7.44bp}\selectfont}] at (0.000,350.750) {misleads.};
\path[fill=tgd03b3b,draw=none,line width=0.5bp] (168.000,322.710) rectangle (203.797,331.310);
\node[anchor=center,inner sep=0pt,outer sep=0pt,text=tgffffff,font={\sffamily\fontsize{6.6bp}{7.92bp}\selectfont\bfseries}] at (185.898,327.010) {21  Fails};
\draw[color=tgc9cdd1,line width=0.45bp,line cap=round] (168.000,363.396) -- (248.000,363.396);
\draw[color=tg9da7b0,line width=0.5bp,line cap=round,dash pattern=on 2bp off 2bp] (232.000,333.710) -- (232.000,363.396);
\path[fill=tg808080,draw=none] (168.333,360.263) circle[radius=0.42bp] (169.000,356.640) circle[radius=0.42bp] (169.667,358.567) circle[radius=0.42bp] (170.333,356.453) circle[radius=0.42bp] (171.000,357.965) circle[radius=0.42bp] (171.667,358.018) circle[radius=0.42bp] (172.333,357.612) circle[radius=0.42bp] (173.000,358.501) circle[radius=0.42bp] (173.667,359.270) circle[radius=0.42bp] (174.333,360.051) circle[radius=0.42bp] (175.000,357.261) circle[radius=0.42bp] (175.667,358.738) circle[radius=0.42bp] (176.333,359.205) circle[radius=0.42bp] (177.000,355.099) circle[radius=0.42bp] (177.667,357.457) circle[radius=0.42bp] (178.333,359.011) circle[radius=0.42bp] (179.000,358.544) circle[radius=0.42bp] (179.667,358.808) circle[radius=0.42bp] (180.333,359.683) circle[radius=0.42bp] (181.000,356.396) circle[radius=0.42bp] (181.667,359.008) circle[radius=0.42bp] (182.333,357.290) circle[radius=0.42bp] (183.000,358.754) circle[radius=0.42bp] (183.667,355.609) circle[radius=0.42bp] (184.333,358.589) circle[radius=0.42bp] (185.000,358.507) circle[radius=0.42bp] (185.667,361.806) circle[radius=0.42bp] (186.333,357.917) circle[radius=0.42bp] (187.000,354.933) circle[radius=0.42bp] (187.667,359.673) circle[radius=0.42bp] (188.333,357.776) circle[radius=0.42bp] (189.000,355.004) circle[radius=0.42bp] (189.667,357.752) circle[radius=0.42bp] (190.333,357.746) circle[radius=0.42bp] (191.000,361.039) circle[radius=0.42bp] (191.667,356.965) circle[radius=0.42bp] (192.333,359.120) circle[radius=0.42bp] (193.000,359.159) circle[radius=0.42bp] (193.667,358.630) circle[radius=0.42bp] (194.333,359.434) circle[radius=0.42bp] (195.000,357.712) circle[radius=0.42bp] (195.667,358.702) circle[radius=0.42bp] (196.333,358.394) circle[radius=0.42bp] (197.000,357.204) circle[radius=0.42bp] (197.667,355.527) circle[radius=0.42bp] (198.333,360.228) circle[radius=0.42bp] (199.000,359.686) circle[radius=0.42bp] (199.667,360.104) circle[radius=0.42bp] (200.333,356.956) circle[radius=0.42bp] (201.000,361.138) circle[radius=0.42bp] (201.667,359.356) circle[radius=0.42bp] (202.333,355.032) circle[radius=0.42bp] (203.000,358.616) circle[radius=0.42bp] (203.667,357.543) circle[radius=0.42bp] (204.333,358.541) circle[radius=0.42bp] (205.000,360.500) circle[radius=0.42bp] (205.667,358.808) circle[radius=0.42bp] (206.333,359.563) circle[radius=0.42bp] (207.000,358.410) circle[radius=0.42bp] (207.667,356.394) circle[radius=0.42bp] (208.333,358.064) circle[radius=0.42bp] (209.000,359.348) circle[radius=0.42bp] (209.667,360.215) circle[radius=0.42bp] (210.333,356.496) circle[radius=0.42bp] (211.000,358.786) circle[radius=0.42bp] (211.667,358.083) circle[radius=0.42bp] (212.333,335.300) circle[radius=0.42bp] (213.000,339.214) circle[radius=0.42bp] (213.667,339.487) circle[radius=0.42bp] (214.333,340.557) circle[radius=0.42bp] (215.000,339.573) circle[radius=0.42bp] (215.667,341.577) circle[radius=0.42bp] (216.333,340.848) circle[radius=0.42bp] (217.000,338.955) circle[radius=0.42bp] (217.667,336.166) circle[radius=0.42bp] (218.333,337.188) circle[radius=0.42bp] (219.000,339.490) circle[radius=0.42bp] (219.667,340.612) circle[radius=0.42bp] (220.333,336.636) circle[radius=0.42bp] (221.000,340.288) circle[radius=0.42bp] (221.667,338.722) circle[radius=0.42bp] (222.333,342.016) circle[radius=0.42bp] (223.000,338.034) circle[radius=0.42bp] (223.667,339.252) circle[radius=0.42bp] (224.333,341.693) circle[radius=0.42bp] (225.000,340.381) circle[radius=0.42bp] (225.667,338.317) circle[radius=0.42bp] (226.333,340.856) circle[radius=0.42bp] (227.000,339.611) circle[radius=0.42bp] (227.667,337.894) circle[radius=0.42bp] (228.333,339.835) circle[radius=0.42bp] (229.000,340.013) circle[radius=0.42bp] (229.667,341.001) circle[radius=0.42bp] (230.333,338.231) circle[radius=0.42bp] (231.000,340.631) circle[radius=0.42bp] (231.667,338.147) circle[radius=0.42bp] (232.333,338.872) circle[radius=0.42bp] (233.000,337.287) circle[radius=0.42bp] (233.667,341.583) circle[radius=0.42bp] (234.333,340.616) circle[radius=0.42bp] (235.000,342.963) circle[radius=0.42bp] (235.667,341.979) circle[radius=0.42bp] (236.333,337.070) circle[radius=0.42bp] (237.000,339.081) circle[radius=0.42bp] (237.667,335.346) circle[radius=0.42bp] (238.333,338.267) circle[radius=0.42bp] (239.000,340.173) circle[radius=0.42bp] (239.667,338.453) circle[radius=0.42bp] (240.333,336.022) circle[radius=0.42bp] (241.000,339.985) circle[radius=0.42bp] (241.667,339.577) circle[radius=0.42bp] (242.333,336.756) circle[radius=0.42bp] (243.000,338.521) circle[radius=0.42bp] (243.667,339.232) circle[radius=0.42bp] (244.333,340.320) circle[radius=0.42bp] (245.000,340.095) circle[radius=0.42bp] (245.667,340.659) circle[radius=0.42bp] (246.333,340.058) circle[radius=0.42bp] (247.000,338.372) circle[radius=0.42bp] (247.667,340.138) circle[radius=0.42bp];
\draw[color=tg000000,line width=0.7bp,line cap=round] (232.333,339.292) -- (233.000,339.292) -- (233.667,339.292) -- (234.333,339.292) -- (235.000,339.292) -- (235.667,339.292) -- (236.333,339.292) -- (237.000,339.292) -- (237.667,339.292) -- (238.333,339.292) -- (239.000,339.292) -- (239.667,339.292) -- (240.333,339.292) -- (241.000,339.292) -- (241.667,339.292) -- (242.333,339.292) -- (243.000,339.292) -- (243.667,339.292) -- (244.333,339.292) -- (245.000,339.292) -- (245.667,339.292) -- (246.333,339.292) -- (247.000,339.292) -- (247.667,339.292);
\draw[color=tg0000ff,line width=0.3bp,line cap=round] (231.783,338.236) -- (232.883,339.336);
\draw[color=tg0000ff,line width=0.3bp,line cap=round] (231.783,339.336) -- (232.883,338.236);
\draw[color=tg0000ff,line width=0.3bp,line cap=round] (232.450,338.153) -- (233.550,339.253);
\draw[color=tg0000ff,line width=0.3bp,line cap=round] (232.450,339.253) -- (233.550,338.153);
\draw[color=tg0000ff,line width=0.3bp,line cap=round] (233.117,337.643) -- (234.217,338.743);
\draw[color=tg0000ff,line width=0.3bp,line cap=round] (233.117,338.743) -- (234.217,337.643);
\draw[color=tg0000ff,line width=0.3bp,line cap=round] (233.783,339.733) -- (234.883,340.833);
\draw[color=tg0000ff,line width=0.3bp,line cap=round] (233.783,340.833) -- (234.883,339.733);
\draw[color=tg0000ff,line width=0.3bp,line cap=round] (234.450,339.540) -- (235.550,340.640);
\draw[color=tg0000ff,line width=0.3bp,line cap=round] (234.450,340.640) -- (235.550,339.540);
\draw[color=tg0000ff,line width=0.3bp,line cap=round] (235.117,338.387) -- (236.217,339.487);
\draw[color=tg0000ff,line width=0.3bp,line cap=round] (235.117,339.487) -- (236.217,338.387);
\draw[color=tg0000ff,line width=0.3bp,line cap=round] (235.783,339.363) -- (236.883,340.463);
\draw[color=tg0000ff,line width=0.3bp,line cap=round] (235.783,340.463) -- (236.883,339.363);
\draw[color=tg0000ff,line width=0.3bp,line cap=round] (236.450,339.485) -- (237.550,340.585);
\draw[color=tg0000ff,line width=0.3bp,line cap=round] (236.450,340.585) -- (237.550,339.485);
\draw[color=tg0000ff,line width=0.3bp,line cap=round] (237.117,339.775) -- (238.217,340.875);
\draw[color=tg0000ff,line width=0.3bp,line cap=round] (237.117,340.875) -- (238.217,339.775);
\draw[color=tg0000ff,line width=0.3bp,line cap=round] (237.783,338.440) -- (238.883,339.540);
\draw[color=tg0000ff,line width=0.3bp,line cap=round] (237.783,339.540) -- (238.883,338.440);
\draw[color=tg0000ff,line width=0.3bp,line cap=round] (238.450,340.029) -- (239.550,341.129);
\draw[color=tg0000ff,line width=0.3bp,line cap=round] (238.450,341.129) -- (239.550,340.029);
\draw[color=tg0000ff,line width=0.3bp,line cap=round] (239.117,340.659) -- (240.217,341.759);
\draw[color=tg0000ff,line width=0.3bp,line cap=round] (239.117,341.759) -- (240.217,340.659);
\draw[color=tg0000ff,line width=0.3bp,line cap=round] (239.783,339.854) -- (240.883,340.954);
\draw[color=tg0000ff,line width=0.3bp,line cap=round] (239.783,340.954) -- (240.883,339.854);
\draw[color=tg0000ff,line width=0.3bp,line cap=round] (240.450,339.864) -- (241.550,340.964);
\draw[color=tg0000ff,line width=0.3bp,line cap=round] (240.450,340.964) -- (241.550,339.864);
\draw[color=tg0000ff,line width=0.3bp,line cap=round] (241.117,340.744) -- (242.217,341.844);
\draw[color=tg0000ff,line width=0.3bp,line cap=round] (241.117,341.844) -- (242.217,340.744);
\draw[color=tg0000ff,line width=0.3bp,line cap=round] (241.783,340.236) -- (242.883,341.336);
\draw[color=tg0000ff,line width=0.3bp,line cap=round] (241.783,341.336) -- (242.883,340.236);
\draw[color=tg0000ff,line width=0.3bp,line cap=round] (242.450,340.355) -- (243.550,341.455);
\draw[color=tg0000ff,line width=0.3bp,line cap=round] (242.450,341.455) -- (243.550,340.355);
\draw[color=tg0000ff,line width=0.3bp,line cap=round] (243.117,340.303) -- (244.217,341.403);
\draw[color=tg0000ff,line width=0.3bp,line cap=round] (243.117,341.403) -- (244.217,340.303);
\draw[color=tg0000ff,line width=0.3bp,line cap=round] (243.783,340.964) -- (244.883,342.064);
\draw[color=tg0000ff,line width=0.3bp,line cap=round] (243.783,342.064) -- (244.883,340.964);
\draw[color=tg0000ff,line width=0.3bp,line cap=round] (244.450,341.152) -- (245.550,342.252);
\draw[color=tg0000ff,line width=0.3bp,line cap=round] (244.450,342.252) -- (245.550,341.152);
\draw[color=tg0000ff,line width=0.3bp,line cap=round] (245.117,340.746) -- (246.217,341.846);
\draw[color=tg0000ff,line width=0.3bp,line cap=round] (245.117,341.846) -- (246.217,340.746);
\draw[color=tg0000ff,line width=0.3bp,line cap=round] (245.783,341.063) -- (246.883,342.163);
\draw[color=tg0000ff,line width=0.3bp,line cap=round] (245.783,342.163) -- (246.883,341.063);
\draw[color=tg0000ff,line width=0.3bp,line cap=round] (246.450,340.877) -- (247.550,341.977);
\draw[color=tg0000ff,line width=0.3bp,line cap=round] (246.450,341.977) -- (247.550,340.877);
\draw[color=tg0000ff,line width=0.3bp,line cap=round] (247.117,341.435) -- (248.217,342.535);
\draw[color=tg0000ff,line width=0.3bp,line cap=round] (247.117,342.535) -- (248.217,341.435);
\path[fill=tgfab219,draw=none,line width=0.5bp] (255.000,322.710) rectangle (309.719,331.310);
\node[anchor=center,inner sep=0pt,outer sep=0pt,text=tg0b0b0b,font={\sffamily\fontsize{6.6bp}{7.92bp}\selectfont\bfseries}] at (282.359,327.010) {45  Struggles};
\draw[color=tgc9cdd1,line width=0.45bp,line cap=round] (255.000,363.396) -- (335.000,363.396);
\draw[color=tg9da7b0,line width=0.5bp,line cap=round,dash pattern=on 2bp off 2bp] (319.000,333.710) -- (319.000,363.396);
\path[fill=tg808080,draw=none] (255.333,358.452) circle[radius=0.42bp] (256.000,350.574) circle[radius=0.42bp] (256.667,354.765) circle[radius=0.42bp] (257.333,350.167) circle[radius=0.42bp] (258.000,353.456) circle[radius=0.42bp] (258.667,353.570) circle[radius=0.42bp] (259.333,352.688) circle[radius=0.42bp] (260.000,354.621) circle[radius=0.42bp] (260.667,356.293) circle[radius=0.42bp] (261.333,357.991) circle[radius=0.42bp] (262.000,351.925) circle[radius=0.42bp] (262.667,355.137) circle[radius=0.42bp] (263.333,356.151) circle[radius=0.42bp] (264.000,347.225) circle[radius=0.42bp] (264.667,352.351) circle[radius=0.42bp] (265.333,355.729) circle[radius=0.42bp] (266.000,354.714) circle[radius=0.42bp] (266.667,355.288) circle[radius=0.42bp] (267.333,357.190) circle[radius=0.42bp] (268.000,350.045) circle[radius=0.42bp] (268.667,355.724) circle[radius=0.42bp] (269.333,351.988) circle[radius=0.42bp] (270.000,355.171) circle[radius=0.42bp] (270.667,348.334) circle[radius=0.42bp] (271.333,354.812) circle[radius=0.42bp] (272.000,354.634) circle[radius=0.42bp] (272.667,361.806) circle[radius=0.42bp] (273.333,353.352) circle[radius=0.42bp] (274.000,346.864) circle[radius=0.42bp] (274.667,357.168) circle[radius=0.42bp] (275.333,353.045) circle[radius=0.42bp] (276.000,347.019) circle[radius=0.42bp] (276.667,352.992) circle[radius=0.42bp] (277.333,352.980) circle[radius=0.42bp] (278.000,360.140) circle[radius=0.42bp] (278.667,351.281) circle[radius=0.42bp] (279.333,355.968) circle[radius=0.42bp] (280.000,356.051) circle[radius=0.42bp] (280.667,354.901) circle[radius=0.42bp] (281.333,356.649) circle[radius=0.42bp] (282.000,352.905) circle[radius=0.42bp] (282.667,355.057) circle[radius=0.42bp] (283.333,354.389) circle[radius=0.42bp] (284.000,351.801) circle[radius=0.42bp] (284.667,348.156) circle[radius=0.42bp] (285.333,358.376) circle[radius=0.42bp] (286.000,357.198) circle[radius=0.42bp] (286.667,358.107) circle[radius=0.42bp] (287.333,351.261) circle[radius=0.42bp] (288.000,360.355) circle[radius=0.42bp] (288.667,356.480) circle[radius=0.42bp] (289.333,347.078) circle[radius=0.42bp] (290.000,354.870) circle[radius=0.42bp] (290.667,352.538) circle[radius=0.42bp] (291.333,354.709) circle[radius=0.42bp] (292.000,358.968) circle[radius=0.42bp] (292.667,355.287) circle[radius=0.42bp] (293.333,356.931) circle[radius=0.42bp] (294.000,354.422) circle[radius=0.42bp] (294.667,350.039) circle[radius=0.42bp] (295.333,353.670) circle[radius=0.42bp] (296.000,356.462) circle[radius=0.42bp] (296.667,358.348) circle[radius=0.42bp] (297.333,350.262) circle[radius=0.42bp] (298.000,355.241) circle[radius=0.42bp] (298.667,353.713) circle[radius=0.42bp] (299.333,335.300) circle[radius=0.42bp] (300.000,343.810) circle[radius=0.42bp] (300.667,344.403) circle[radius=0.42bp] (301.333,346.728) circle[radius=0.42bp] (302.000,344.590) circle[radius=0.42bp] (302.667,348.946) circle[radius=0.42bp] (303.333,347.362) circle[radius=0.42bp] (304.000,343.245) circle[radius=0.42bp] (304.667,337.181) circle[radius=0.42bp] (305.333,339.404) circle[radius=0.42bp] (306.000,344.409) circle[radius=0.42bp] (306.667,346.849) circle[radius=0.42bp] (307.333,338.205) circle[radius=0.42bp] (308.000,346.144) circle[radius=0.42bp] (308.667,342.740) circle[radius=0.42bp] (309.333,349.901) circle[radius=0.42bp] (310.000,341.244) circle[radius=0.42bp] (310.667,343.891) circle[radius=0.42bp] (311.333,349.200) circle[radius=0.42bp] (312.000,346.346) circle[radius=0.42bp] (312.667,341.860) circle[radius=0.42bp] (313.333,347.378) circle[radius=0.42bp] (314.000,344.672) circle[radius=0.42bp] (314.667,340.939) circle[radius=0.42bp] (315.333,345.160) circle[radius=0.42bp] (316.000,345.545) circle[radius=0.42bp] (316.667,347.695) circle[radius=0.42bp] (317.333,341.672) circle[radius=0.42bp] (318.000,346.890) circle[radius=0.42bp] (318.667,341.489) circle[radius=0.42bp] (319.333,343.066) circle[radius=0.42bp] (320.000,339.620) circle[radius=0.42bp] (320.667,348.960) circle[radius=0.42bp] (321.333,346.857) circle[radius=0.42bp] (322.000,351.960) circle[radius=0.42bp] (322.667,349.820) circle[radius=0.42bp] (323.333,339.148) circle[radius=0.42bp] (324.000,343.519) circle[radius=0.42bp] (324.667,335.400) circle[radius=0.42bp] (325.333,341.751) circle[radius=0.42bp] (326.000,345.894) circle[radius=0.42bp] (326.667,342.154) circle[radius=0.42bp] (327.333,336.869) circle[radius=0.42bp] (328.000,345.486) circle[radius=0.42bp] (328.667,344.599) circle[radius=0.42bp] (329.333,338.465) circle[radius=0.42bp] (330.000,342.302) circle[radius=0.42bp] (330.667,343.848) circle[radius=0.42bp] (331.333,346.215) circle[radius=0.42bp] (332.000,345.724) circle[radius=0.42bp] (332.667,346.951) circle[radius=0.42bp] (333.333,345.643) circle[radius=0.42bp] (334.000,341.978) circle[radius=0.42bp] (334.667,345.819) circle[radius=0.42bp];
\draw[color=tg000000,line width=0.7bp,line cap=round] (319.333,343.978) -- (320.000,343.978) -- (320.667,343.978) -- (321.333,343.978) -- (322.000,343.978) -- (322.667,343.978) -- (323.333,343.978) -- (324.000,343.978) -- (324.667,343.978) -- (325.333,343.978) -- (326.000,343.978) -- (326.667,343.978) -- (327.333,343.978) -- (328.000,343.978) -- (328.667,343.978) -- (329.333,343.978) -- (330.000,343.978) -- (330.667,343.978) -- (331.333,343.978) -- (332.000,343.978) -- (332.667,343.978) -- (333.333,343.978) -- (334.000,343.978) -- (334.667,343.978);
\draw[color=tg0000ff,line width=0.3bp,line cap=round] (318.783,342.227) -- (319.883,343.327);
\draw[color=tg0000ff,line width=0.3bp,line cap=round] (318.783,343.327) -- (319.883,342.227);
\draw[color=tg0000ff,line width=0.3bp,line cap=round] (319.450,340.574) -- (320.550,341.674);
\draw[color=tg0000ff,line width=0.3bp,line cap=round] (319.450,341.674) -- (320.550,340.574);
\draw[color=tg0000ff,line width=0.3bp,line cap=round] (320.117,340.328) -- (321.217,341.428);
\draw[color=tg0000ff,line width=0.3bp,line cap=round] (320.117,341.428) -- (321.217,340.328);
\draw[color=tg0000ff,line width=0.3bp,line cap=round] (320.783,344.230) -- (321.883,345.330);
\draw[color=tg0000ff,line width=0.3bp,line cap=round] (320.783,345.330) -- (321.883,344.230);
\draw[color=tg0000ff,line width=0.3bp,line cap=round] (321.450,344.005) -- (322.550,345.105);
\draw[color=tg0000ff,line width=0.3bp,line cap=round] (321.450,345.105) -- (322.550,344.005);
\draw[color=tg0000ff,line width=0.3bp,line cap=round] (322.117,339.909) -- (323.217,341.009);
\draw[color=tg0000ff,line width=0.3bp,line cap=round] (322.117,341.009) -- (323.217,339.909);
\draw[color=tg0000ff,line width=0.3bp,line cap=round] (322.783,344.080) -- (323.883,345.180);
\draw[color=tg0000ff,line width=0.3bp,line cap=round] (322.783,345.180) -- (323.883,344.080);
\draw[color=tg0000ff,line width=0.3bp,line cap=round] (323.450,341.783) -- (324.550,342.883);
\draw[color=tg0000ff,line width=0.3bp,line cap=round] (323.450,342.883) -- (324.550,341.783);
\draw[color=tg0000ff,line width=0.3bp,line cap=round] (324.117,342.518) -- (325.217,343.618);
\draw[color=tg0000ff,line width=0.3bp,line cap=round] (324.117,343.618) -- (325.217,342.518);
\draw[color=tg0000ff,line width=0.3bp,line cap=round] (324.783,338.042) -- (325.883,339.142);
\draw[color=tg0000ff,line width=0.3bp,line cap=round] (324.783,339.142) -- (325.883,338.042);
\draw[color=tg0000ff,line width=0.3bp,line cap=round] (325.450,342.650) -- (326.550,343.750);
\draw[color=tg0000ff,line width=0.3bp,line cap=round] (325.450,343.750) -- (326.550,342.650);
\draw[color=tg0000ff,line width=0.3bp,line cap=round] (326.117,343.741) -- (327.217,344.841);
\draw[color=tg0000ff,line width=0.3bp,line cap=round] (326.117,344.841) -- (327.217,343.741);
\draw[color=tg0000ff,line width=0.3bp,line cap=round] (326.783,341.011) -- (327.883,342.111);
\draw[color=tg0000ff,line width=0.3bp,line cap=round] (326.783,342.111) -- (327.883,341.011);
\draw[color=tg0000ff,line width=0.3bp,line cap=round] (327.450,341.045) -- (328.550,342.145);
\draw[color=tg0000ff,line width=0.3bp,line cap=round] (327.450,342.145) -- (328.550,341.045);
\draw[color=tg0000ff,line width=0.3bp,line cap=round] (328.117,342.760) -- (329.217,343.860);
\draw[color=tg0000ff,line width=0.3bp,line cap=round] (328.117,343.860) -- (329.217,342.760);
\draw[color=tg0000ff,line width=0.3bp,line cap=round] (328.783,341.620) -- (329.883,342.720);
\draw[color=tg0000ff,line width=0.3bp,line cap=round] (328.783,342.720) -- (329.883,341.620);
\draw[color=tg0000ff,line width=0.3bp,line cap=round] (329.450,341.023) -- (330.550,342.123);
\draw[color=tg0000ff,line width=0.3bp,line cap=round] (329.450,342.123) -- (330.550,341.023);
\draw[color=tg0000ff,line width=0.3bp,line cap=round] (330.117,341.096) -- (331.217,342.196);
\draw[color=tg0000ff,line width=0.3bp,line cap=round] (330.117,342.196) -- (331.217,341.096);
\draw[color=tg0000ff,line width=0.3bp,line cap=round] (330.783,342.423) -- (331.883,343.523);
\draw[color=tg0000ff,line width=0.3bp,line cap=round] (330.783,343.523) -- (331.883,342.423);
\draw[color=tg0000ff,line width=0.3bp,line cap=round] (331.450,343.125) -- (332.550,344.225);
\draw[color=tg0000ff,line width=0.3bp,line cap=round] (331.450,344.225) -- (332.550,343.125);
\draw[color=tg0000ff,line width=0.3bp,line cap=round] (332.117,340.641) -- (333.217,341.741);
\draw[color=tg0000ff,line width=0.3bp,line cap=round] (332.117,341.741) -- (333.217,340.641);
\draw[color=tg0000ff,line width=0.3bp,line cap=round] (332.783,341.851) -- (333.883,342.951);
\draw[color=tg0000ff,line width=0.3bp,line cap=round] (332.783,342.951) -- (333.883,341.851);
\draw[color=tg0000ff,line width=0.3bp,line cap=round] (333.450,340.319) -- (334.550,341.419);
\draw[color=tg0000ff,line width=0.3bp,line cap=round] (333.450,341.419) -- (334.550,340.319);
\draw[color=tg0000ff,line width=0.3bp,line cap=round] (334.117,342.177) -- (335.217,343.277);
\draw[color=tg0000ff,line width=0.3bp,line cap=round] (334.117,343.277) -- (335.217,342.177);
\node[anchor=west,inner sep=0pt,outer sep=0pt,text=tg18314b,font={\sffamily\fontsize{6.2bp}{7.44bp}\selectfont}] at (346.000,326.430) {Low: Forecasts too low;};
\node[anchor=west,inner sep=0pt,outer sep=0pt,text=tg18314b,font={\sffamily\fontsize{6.2bp}{7.44bp}\selectfont}] at (352.000,333.870) {Invents a trend};
\node[anchor=west,inner sep=0pt,outer sep=0pt,text=tg18314b,font={\sffamily\fontsize{6.2bp}{7.44bp}\selectfont}] at (346.000,341.310) {High: Misreads the current level;};
\node[anchor=west,inner sep=0pt,outer sep=0pt,text=tg18314b,font={\sffamily\fontsize{6.2bp}{7.44bp}\selectfont}] at (352.000,348.750) {Invents a trend};
\draw[color=tgeceef0,line width=0.4bp,line cap=round] (0.000,366.396) -- (448.000,366.396);
\node[anchor=west,inner sep=0pt,outer sep=0pt,text=tg18314b,font={\sffamily\fontsize{7.2bp}{8.64bp}\selectfont\bfseries}] at (0.000,373.216) {Mean reversion (AR(1)): shocks fade};
\node[anchor=west,inner sep=0pt,outer sep=0pt,text=tg18314b,font={\sffamily\fontsize{6.2bp}{7.44bp}\selectfont}] at (0.000,381.656) {\ensuremath{\checkmark} The pull back to the average: a gap halves in};
\node[anchor=west,inner sep=0pt,outer sep=0pt,text=tg18314b,font={\sffamily\fontsize{6.2bp}{7.44bp}\selectfont}] at (0.000,389.096) {about three steps.};
\node[anchor=west,inner sep=0pt,outer sep=0pt,text=tg18314b,font={\sffamily\fontsize{6.2bp}{7.44bp}\selectfont}] at (0.000,396.936) {\ensuremath{\times} New shocks. After about ten steps only the};
\node[anchor=west,inner sep=0pt,outer sep=0pt,text=tg18314b,font={\sffamily\fontsize{6.2bp}{7.44bp}\selectfont}] at (0.000,404.376) {average is left; measurement noise hides where the};
\node[anchor=west,inner sep=0pt,outer sep=0pt,text=tg18314b,font={\sffamily\fontsize{6.2bp}{7.44bp}\selectfont}] at (0.000,411.816) {series is now.};
\path[fill=tg8fd18f,draw=none,line width=0.5bp] (168.000,368.896) rectangle (205.969,377.496);
\node[anchor=center,inner sep=0pt,outer sep=0pt,text=tg0b0b0b,font={\sffamily\fontsize{6.6bp}{7.92bp}\selectfont\bfseries}] at (186.984,373.196) {70  Good};
\draw[color=tgc9cdd1,line width=0.45bp,line cap=round] (168.000,418.722) -- (248.000,418.722);
\draw[color=tg9da7b0,line width=0.5bp,line cap=round,dash pattern=on 2bp off 2bp] (232.000,379.896) -- (232.000,418.722);
\path[fill=tg808080,draw=none] (168.333,396.047) circle[radius=0.42bp] (169.000,398.038) circle[radius=0.42bp] (169.667,401.432) circle[radius=0.42bp] (170.333,407.009) circle[radius=0.42bp] (171.000,406.472) circle[radius=0.42bp] (171.667,405.822) circle[radius=0.42bp] (172.333,408.710) circle[radius=0.42bp] (173.000,409.792) circle[radius=0.42bp] (173.667,411.824) circle[radius=0.42bp] (174.333,403.960) circle[radius=0.42bp] (175.000,408.417) circle[radius=0.42bp] (175.667,403.666) circle[radius=0.42bp] (176.333,406.253) circle[radius=0.42bp] (177.000,407.974) circle[radius=0.42bp] (177.667,401.405) circle[radius=0.42bp] (178.333,397.526) circle[radius=0.42bp] (179.000,404.161) circle[radius=0.42bp] (179.667,396.668) circle[radius=0.42bp] (180.333,394.686) circle[radius=0.42bp] (181.000,390.167) circle[radius=0.42bp] (181.667,393.370) circle[radius=0.42bp] (182.333,398.136) circle[radius=0.42bp] (183.000,396.718) circle[radius=0.42bp] (183.667,390.195) circle[radius=0.42bp] (184.333,381.976) circle[radius=0.42bp] (185.000,383.103) circle[radius=0.42bp] (185.667,386.374) circle[radius=0.42bp] (186.333,390.939) circle[radius=0.42bp] (187.000,398.532) circle[radius=0.42bp] (187.667,396.265) circle[radius=0.42bp] (188.333,402.744) circle[radius=0.42bp] (189.000,399.632) circle[radius=0.42bp] (189.667,401.022) circle[radius=0.42bp] (190.333,399.107) circle[radius=0.42bp] (191.000,406.742) circle[radius=0.42bp] (191.667,404.868) circle[radius=0.42bp] (192.333,405.933) circle[radius=0.42bp] (193.000,412.161) circle[radius=0.42bp] (193.667,402.827) circle[radius=0.42bp] (194.333,407.676) circle[radius=0.42bp] (195.000,404.601) circle[radius=0.42bp] (195.667,402.106) circle[radius=0.42bp] (196.333,398.269) circle[radius=0.42bp] (197.000,390.146) circle[radius=0.42bp] (197.667,395.529) circle[radius=0.42bp] (198.333,403.334) circle[radius=0.42bp] (199.000,404.524) circle[radius=0.42bp] (199.667,400.416) circle[radius=0.42bp] (200.333,398.166) circle[radius=0.42bp] (201.000,402.754) circle[radius=0.42bp] (201.667,396.827) circle[radius=0.42bp] (202.333,396.114) circle[radius=0.42bp] (203.000,405.057) circle[radius=0.42bp] (203.667,405.002) circle[radius=0.42bp] (204.333,402.763) circle[radius=0.42bp] (205.000,393.053) circle[radius=0.42bp] (205.667,396.510) circle[radius=0.42bp] (206.333,397.057) circle[radius=0.42bp] (207.000,400.049) circle[radius=0.42bp] (207.667,406.834) circle[radius=0.42bp] (208.333,411.820) circle[radius=0.42bp] (209.000,409.038) circle[radius=0.42bp] (209.667,405.602) circle[radius=0.42bp] (210.333,402.070) circle[radius=0.42bp] (211.000,400.779) circle[radius=0.42bp] (211.667,393.754) circle[radius=0.42bp] (212.333,401.735) circle[radius=0.42bp] (213.000,396.496) circle[radius=0.42bp] (213.667,401.222) circle[radius=0.42bp] (214.333,403.723) circle[radius=0.42bp] (215.000,407.285) circle[radius=0.42bp] (215.667,400.107) circle[radius=0.42bp] (216.333,396.034) circle[radius=0.42bp] (217.000,402.508) circle[radius=0.42bp] (217.667,405.769) circle[radius=0.42bp] (218.333,413.015) circle[radius=0.42bp] (219.000,405.166) circle[radius=0.42bp] (219.667,403.454) circle[radius=0.42bp] (220.333,405.222) circle[radius=0.42bp] (221.000,403.615) circle[radius=0.42bp] (221.667,404.873) circle[radius=0.42bp] (222.333,411.431) circle[radius=0.42bp] (223.000,410.412) circle[radius=0.42bp] (223.667,407.131) circle[radius=0.42bp] (224.333,413.009) circle[radius=0.42bp] (225.000,413.118) circle[radius=0.42bp] (225.667,410.564) circle[radius=0.42bp] (226.333,408.428) circle[radius=0.42bp] (227.000,410.409) circle[radius=0.42bp] (227.667,407.014) circle[radius=0.42bp] (228.333,405.588) circle[radius=0.42bp] (229.000,402.407) circle[radius=0.42bp] (229.667,402.220) circle[radius=0.42bp] (230.333,399.582) circle[radius=0.42bp] (231.000,398.577) circle[radius=0.42bp] (231.667,394.230) circle[radius=0.42bp] (232.333,392.646) circle[radius=0.42bp] (233.000,390.204) circle[radius=0.42bp] (233.667,395.011) circle[radius=0.42bp] (234.333,401.683) circle[radius=0.42bp] (235.000,405.083) circle[radius=0.42bp] (235.667,407.220) circle[radius=0.42bp] (236.333,405.475) circle[radius=0.42bp] (237.000,406.411) circle[radius=0.42bp] (237.667,411.187) circle[radius=0.42bp] (238.333,415.047) circle[radius=0.42bp] (239.000,416.642) circle[radius=0.42bp] (239.667,411.629) circle[radius=0.42bp] (240.333,408.626) circle[radius=0.42bp] (241.000,404.112) circle[radius=0.42bp] (241.667,405.459) circle[radius=0.42bp] (242.333,396.485) circle[radius=0.42bp] (243.000,393.930) circle[radius=0.42bp] (243.667,392.945) circle[radius=0.42bp] (244.333,404.049) circle[radius=0.42bp] (245.000,400.008) circle[radius=0.42bp] (245.667,404.308) circle[radius=0.42bp] (246.333,397.904) circle[radius=0.42bp] (247.000,399.922) circle[radius=0.42bp] (247.667,399.547) circle[radius=0.42bp];
\draw[color=tg000000,line width=0.7bp,line cap=round] (232.333,397.013) -- (233.000,398.404) -- (233.667,399.518) -- (234.333,400.408) -- (235.000,401.121) -- (235.667,401.691) -- (236.333,402.147) -- (237.000,402.512) -- (237.667,402.804) -- (238.333,403.037) -- (239.000,403.224) -- (239.667,403.374) -- (240.333,403.493) -- (241.000,403.589) -- (241.667,403.665) -- (242.333,403.726) -- (243.000,403.775) -- (243.667,403.815) -- (244.333,403.846) -- (245.000,403.871) -- (245.667,403.891) -- (246.333,403.907) -- (247.000,403.920) -- (247.667,403.930);
\draw[color=tg0000ff,line width=0.3bp,line cap=round] (231.783,396.443) -- (232.883,397.543);
\draw[color=tg0000ff,line width=0.3bp,line cap=round] (231.783,397.543) -- (232.883,396.443);
\draw[color=tg0000ff,line width=0.3bp,line cap=round] (232.450,396.022) -- (233.550,397.122);
\draw[color=tg0000ff,line width=0.3bp,line cap=round] (232.450,397.122) -- (233.550,396.022);
\draw[color=tg0000ff,line width=0.3bp,line cap=round] (233.117,397.177) -- (234.217,398.277);
\draw[color=tg0000ff,line width=0.3bp,line cap=round] (233.117,398.277) -- (234.217,397.177);
\draw[color=tg0000ff,line width=0.3bp,line cap=round] (233.783,400.187) -- (234.883,401.287);
\draw[color=tg0000ff,line width=0.3bp,line cap=round] (233.783,401.287) -- (234.883,400.187);
\draw[color=tg0000ff,line width=0.3bp,line cap=round] (234.450,399.485) -- (235.550,400.585);
\draw[color=tg0000ff,line width=0.3bp,line cap=round] (234.450,400.585) -- (235.550,399.485);
\draw[color=tg0000ff,line width=0.3bp,line cap=round] (235.117,398.640) -- (236.217,399.740);
\draw[color=tg0000ff,line width=0.3bp,line cap=round] (235.117,399.740) -- (236.217,398.640);
\draw[color=tg0000ff,line width=0.3bp,line cap=round] (235.783,399.976) -- (236.883,401.076);
\draw[color=tg0000ff,line width=0.3bp,line cap=round] (235.783,401.076) -- (236.883,399.976);
\draw[color=tg0000ff,line width=0.3bp,line cap=round] (236.450,401.875) -- (237.550,402.975);
\draw[color=tg0000ff,line width=0.3bp,line cap=round] (236.450,402.975) -- (237.550,401.875);
\draw[color=tg0000ff,line width=0.3bp,line cap=round] (237.117,401.332) -- (238.217,402.432);
\draw[color=tg0000ff,line width=0.3bp,line cap=round] (237.117,402.432) -- (238.217,401.332);
\draw[color=tg0000ff,line width=0.3bp,line cap=round] (237.783,402.230) -- (238.883,403.330);
\draw[color=tg0000ff,line width=0.3bp,line cap=round] (237.783,403.330) -- (238.883,402.230);
\draw[color=tg0000ff,line width=0.3bp,line cap=round] (238.450,403.470) -- (239.550,404.570);
\draw[color=tg0000ff,line width=0.3bp,line cap=round] (238.450,404.570) -- (239.550,403.470);
\draw[color=tg0000ff,line width=0.3bp,line cap=round] (239.117,403.676) -- (240.217,404.776);
\draw[color=tg0000ff,line width=0.3bp,line cap=round] (239.117,404.776) -- (240.217,403.676);
\draw[color=tg0000ff,line width=0.3bp,line cap=round] (239.783,403.526) -- (240.883,404.626);
\draw[color=tg0000ff,line width=0.3bp,line cap=round] (239.783,404.626) -- (240.883,403.526);
\draw[color=tg0000ff,line width=0.3bp,line cap=round] (240.450,402.742) -- (241.550,403.842);
\draw[color=tg0000ff,line width=0.3bp,line cap=round] (240.450,403.842) -- (241.550,402.742);
\draw[color=tg0000ff,line width=0.3bp,line cap=round] (241.117,402.807) -- (242.217,403.907);
\draw[color=tg0000ff,line width=0.3bp,line cap=round] (241.117,403.907) -- (242.217,402.807);
\draw[color=tg0000ff,line width=0.3bp,line cap=round] (241.783,403.757) -- (242.883,404.857);
\draw[color=tg0000ff,line width=0.3bp,line cap=round] (241.783,404.857) -- (242.883,403.757);
\draw[color=tg0000ff,line width=0.3bp,line cap=round] (242.450,402.941) -- (243.550,404.041);
\draw[color=tg0000ff,line width=0.3bp,line cap=round] (242.450,404.041) -- (243.550,402.941);
\draw[color=tg0000ff,line width=0.3bp,line cap=round] (243.117,403.065) -- (244.217,404.165);
\draw[color=tg0000ff,line width=0.3bp,line cap=round] (243.117,404.165) -- (244.217,403.065);
\draw[color=tg0000ff,line width=0.3bp,line cap=round] (243.783,403.503) -- (244.883,404.603);
\draw[color=tg0000ff,line width=0.3bp,line cap=round] (243.783,404.603) -- (244.883,403.503);
\draw[color=tg0000ff,line width=0.3bp,line cap=round] (244.450,402.963) -- (245.550,404.063);
\draw[color=tg0000ff,line width=0.3bp,line cap=round] (244.450,404.063) -- (245.550,402.963);
\draw[color=tg0000ff,line width=0.3bp,line cap=round] (245.117,402.279) -- (246.217,403.379);
\draw[color=tg0000ff,line width=0.3bp,line cap=round] (245.117,403.379) -- (246.217,402.279);
\draw[color=tg0000ff,line width=0.3bp,line cap=round] (245.783,402.127) -- (246.883,403.227);
\draw[color=tg0000ff,line width=0.3bp,line cap=round] (245.783,403.227) -- (246.883,402.127);
\draw[color=tg0000ff,line width=0.3bp,line cap=round] (246.450,401.560) -- (247.550,402.660);
\draw[color=tg0000ff,line width=0.3bp,line cap=round] (246.450,402.660) -- (247.550,401.560);
\draw[color=tg0000ff,line width=0.3bp,line cap=round] (247.117,400.251) -- (248.217,401.351);
\draw[color=tg0000ff,line width=0.3bp,line cap=round] (247.117,401.351) -- (248.217,400.251);
\path[fill=tg8fd18f,draw=none,line width=0.5bp] (255.000,368.896) rectangle (292.969,377.496);
\node[anchor=center,inner sep=0pt,outer sep=0pt,text=tg0b0b0b,font={\sffamily\fontsize{6.6bp}{7.92bp}\selectfont\bfseries}] at (273.984,373.196) {83  Good};
\draw[color=tgc9cdd1,line width=0.45bp,line cap=round] (255.000,418.722) -- (335.000,418.722);
\draw[color=tg9da7b0,line width=0.5bp,line cap=round,dash pattern=on 2bp off 2bp] (319.000,379.896) -- (319.000,418.722);
\path[fill=tg808080,draw=none] (255.333,389.596) circle[radius=0.42bp] (256.000,393.935) circle[radius=0.42bp] (256.667,395.717) circle[radius=0.42bp] (257.333,403.739) circle[radius=0.42bp] (258.000,399.860) circle[radius=0.42bp] (258.667,401.603) circle[radius=0.42bp] (259.333,402.654) circle[radius=0.42bp] (260.000,406.498) circle[radius=0.42bp] (260.667,413.251) circle[radius=0.42bp] (261.333,397.645) circle[radius=0.42bp] (262.000,412.717) circle[radius=0.42bp] (262.667,402.551) circle[radius=0.42bp] (263.333,404.224) circle[radius=0.42bp] (264.000,408.811) circle[radius=0.42bp] (264.667,397.526) circle[radius=0.42bp] (265.333,398.080) circle[radius=0.42bp] (266.000,403.692) circle[radius=0.42bp] (266.667,394.576) circle[radius=0.42bp] (267.333,397.200) circle[radius=0.42bp] (268.000,386.509) circle[radius=0.42bp] (268.667,392.939) circle[radius=0.42bp] (269.333,399.613) circle[radius=0.42bp] (270.000,404.541) circle[radius=0.42bp] (270.667,397.396) circle[radius=0.42bp] (271.333,381.976) circle[radius=0.42bp] (272.000,387.351) circle[radius=0.42bp] (272.667,390.584) circle[radius=0.42bp] (273.333,386.027) circle[radius=0.42bp] (274.000,403.148) circle[radius=0.42bp] (274.667,398.748) circle[radius=0.42bp] (275.333,402.841) circle[radius=0.42bp] (276.000,396.389) circle[radius=0.42bp] (276.667,402.402) circle[radius=0.42bp] (277.333,400.855) circle[radius=0.42bp] (278.000,410.212) circle[radius=0.42bp] (278.667,399.489) circle[radius=0.42bp] (279.333,399.927) circle[radius=0.42bp] (280.000,409.468) circle[radius=0.42bp] (280.667,400.337) circle[radius=0.42bp] (281.333,409.608) circle[radius=0.42bp] (282.000,404.141) circle[radius=0.42bp] (282.667,394.936) circle[radius=0.42bp] (283.333,397.248) circle[radius=0.42bp] (284.000,382.782) circle[radius=0.42bp] (284.667,399.043) circle[radius=0.42bp] (285.333,399.013) circle[radius=0.42bp] (286.000,402.975) circle[radius=0.42bp] (286.667,397.022) circle[radius=0.42bp] (287.333,397.751) circle[radius=0.42bp] (288.000,403.205) circle[radius=0.42bp] (288.667,400.278) circle[radius=0.42bp] (289.333,383.927) circle[radius=0.42bp] (290.000,402.873) circle[radius=0.42bp] (290.667,410.648) circle[radius=0.42bp] (291.333,403.563) circle[radius=0.42bp] (292.000,391.400) circle[radius=0.42bp] (292.667,398.574) circle[radius=0.42bp] (293.333,391.462) circle[radius=0.42bp] (294.000,395.005) circle[radius=0.42bp] (294.667,404.422) circle[radius=0.42bp] (295.333,409.212) circle[radius=0.42bp] (296.000,409.405) circle[radius=0.42bp] (296.667,405.764) circle[radius=0.42bp] (297.333,406.399) circle[radius=0.42bp] (298.000,394.704) circle[radius=0.42bp] (298.667,388.215) circle[radius=0.42bp] (299.333,405.290) circle[radius=0.42bp] (300.000,401.312) circle[radius=0.42bp] (300.667,406.995) circle[radius=0.42bp] (301.333,405.508) circle[radius=0.42bp] (302.000,416.642) circle[radius=0.42bp] (302.667,403.688) circle[radius=0.42bp] (303.333,391.000) circle[radius=0.42bp] (304.000,402.164) circle[radius=0.42bp] (304.667,399.617) circle[radius=0.42bp] (305.333,410.118) circle[radius=0.42bp] (306.000,395.543) circle[radius=0.42bp] (306.667,400.632) circle[radius=0.42bp] (307.333,404.216) circle[radius=0.42bp] (308.000,391.263) circle[radius=0.42bp] (308.667,403.739) circle[radius=0.42bp] (309.333,411.957) circle[radius=0.42bp] (310.000,405.507) circle[radius=0.42bp] (310.667,404.173) circle[radius=0.42bp] (311.333,403.847) circle[radius=0.42bp] (312.000,411.935) circle[radius=0.42bp] (312.667,403.185) circle[radius=0.42bp] (313.333,401.491) circle[radius=0.42bp] (314.000,409.009) circle[radius=0.42bp] (314.667,403.277) circle[radius=0.42bp] (315.333,406.281) circle[radius=0.42bp] (316.000,391.248) circle[radius=0.42bp] (316.667,403.407) circle[radius=0.42bp] (317.333,400.542) circle[radius=0.42bp] (318.000,402.306) circle[radius=0.42bp] (318.667,396.354) circle[radius=0.42bp] (319.333,393.871) circle[radius=0.42bp] (320.000,389.038) circle[radius=0.42bp] (320.667,397.967) circle[radius=0.42bp] (321.333,400.941) circle[radius=0.42bp] (322.000,408.287) circle[radius=0.42bp] (322.667,407.413) circle[radius=0.42bp] (323.333,391.750) circle[radius=0.42bp] (324.000,399.316) circle[radius=0.42bp] (324.667,410.014) circle[radius=0.42bp] (325.333,414.964) circle[radius=0.42bp] (326.000,414.626) circle[radius=0.42bp] (326.667,407.263) circle[radius=0.42bp] (327.333,404.926) circle[radius=0.42bp] (328.000,404.897) circle[radius=0.42bp] (328.667,408.027) circle[radius=0.42bp] (329.333,390.811) circle[radius=0.42bp] (330.000,392.928) circle[radius=0.42bp] (330.667,395.566) circle[radius=0.42bp] (331.333,402.742) circle[radius=0.42bp] (332.000,401.150) circle[radius=0.42bp] (332.667,413.866) circle[radius=0.42bp] (333.333,397.073) circle[radius=0.42bp] (334.000,403.269) circle[radius=0.42bp] (334.667,402.733) circle[radius=0.42bp];
\draw[color=tg000000,line width=0.7bp,line cap=round] (319.333,400.354) -- (320.000,400.643) -- (320.667,400.874) -- (321.333,401.059) -- (322.000,401.207) -- (322.667,401.325) -- (323.333,401.420) -- (324.000,401.496) -- (324.667,401.556) -- (325.333,401.605) -- (326.000,401.643) -- (326.667,401.674) -- (327.333,401.699) -- (328.000,401.719) -- (328.667,401.735) -- (329.333,401.748) -- (330.000,401.758) -- (330.667,401.766) -- (331.333,401.772) -- (332.000,401.778) -- (332.667,401.782) -- (333.333,401.785) -- (334.000,401.788) -- (334.667,401.790);
\draw[color=tg0000ff,line width=0.3bp,line cap=round] (318.783,396.143) -- (319.883,397.243);
\draw[color=tg0000ff,line width=0.3bp,line cap=round] (318.783,397.243) -- (319.883,396.143);
\draw[color=tg0000ff,line width=0.3bp,line cap=round] (319.450,400.017) -- (320.550,401.117);
\draw[color=tg0000ff,line width=0.3bp,line cap=round] (319.450,401.117) -- (320.550,400.017);
\draw[color=tg0000ff,line width=0.3bp,line cap=round] (320.117,396.730) -- (321.217,397.830);
\draw[color=tg0000ff,line width=0.3bp,line cap=round] (320.117,397.830) -- (321.217,396.730);
\draw[color=tg0000ff,line width=0.3bp,line cap=round] (320.783,400.551) -- (321.883,401.651);
\draw[color=tg0000ff,line width=0.3bp,line cap=round] (320.783,401.651) -- (321.883,400.551);
\draw[color=tg0000ff,line width=0.3bp,line cap=round] (321.450,400.781) -- (322.550,401.881);
\draw[color=tg0000ff,line width=0.3bp,line cap=round] (321.450,401.881) -- (322.550,400.781);
\draw[color=tg0000ff,line width=0.3bp,line cap=round] (322.117,398.361) -- (323.217,399.461);
\draw[color=tg0000ff,line width=0.3bp,line cap=round] (322.117,399.461) -- (323.217,398.361);
\draw[color=tg0000ff,line width=0.3bp,line cap=round] (322.783,402.969) -- (323.883,404.069);
\draw[color=tg0000ff,line width=0.3bp,line cap=round] (322.783,404.069) -- (323.883,402.969);
\draw[color=tg0000ff,line width=0.3bp,line cap=round] (323.450,398.134) -- (324.550,399.234);
\draw[color=tg0000ff,line width=0.3bp,line cap=round] (323.450,399.234) -- (324.550,398.134);
\draw[color=tg0000ff,line width=0.3bp,line cap=round] (324.117,402.221) -- (325.217,403.321);
\draw[color=tg0000ff,line width=0.3bp,line cap=round] (324.117,403.321) -- (325.217,402.221);
\draw[color=tg0000ff,line width=0.3bp,line cap=round] (324.783,400.672) -- (325.883,401.772);
\draw[color=tg0000ff,line width=0.3bp,line cap=round] (324.783,401.772) -- (325.883,400.672);
\draw[color=tg0000ff,line width=0.3bp,line cap=round] (325.450,402.085) -- (326.550,403.185);
\draw[color=tg0000ff,line width=0.3bp,line cap=round] (325.450,403.185) -- (326.550,402.085);
\draw[color=tg0000ff,line width=0.3bp,line cap=round] (326.117,404.069) -- (327.217,405.169);
\draw[color=tg0000ff,line width=0.3bp,line cap=round] (326.117,405.169) -- (327.217,404.069);
\draw[color=tg0000ff,line width=0.3bp,line cap=round] (326.783,399.847) -- (327.883,400.947);
\draw[color=tg0000ff,line width=0.3bp,line cap=round] (326.783,400.947) -- (327.883,399.847);
\draw[color=tg0000ff,line width=0.3bp,line cap=round] (327.450,402.554) -- (328.550,403.654);
\draw[color=tg0000ff,line width=0.3bp,line cap=round] (327.450,403.654) -- (328.550,402.554);
\draw[color=tg0000ff,line width=0.3bp,line cap=round] (328.117,399.781) -- (329.217,400.881);
\draw[color=tg0000ff,line width=0.3bp,line cap=round] (328.117,400.881) -- (329.217,399.781);
\draw[color=tg0000ff,line width=0.3bp,line cap=round] (328.783,401.538) -- (329.883,402.638);
\draw[color=tg0000ff,line width=0.3bp,line cap=round] (328.783,402.638) -- (329.883,401.538);
\draw[color=tg0000ff,line width=0.3bp,line cap=round] (329.450,401.208) -- (330.550,402.308);
\draw[color=tg0000ff,line width=0.3bp,line cap=round] (329.450,402.308) -- (330.550,401.208);
\draw[color=tg0000ff,line width=0.3bp,line cap=round] (330.117,399.324) -- (331.217,400.424);
\draw[color=tg0000ff,line width=0.3bp,line cap=round] (330.117,400.424) -- (331.217,399.324);
\draw[color=tg0000ff,line width=0.3bp,line cap=round] (330.783,401.622) -- (331.883,402.722);
\draw[color=tg0000ff,line width=0.3bp,line cap=round] (330.783,402.722) -- (331.883,401.622);
\draw[color=tg0000ff,line width=0.3bp,line cap=round] (331.450,397.888) -- (332.550,398.988);
\draw[color=tg0000ff,line width=0.3bp,line cap=round] (331.450,398.988) -- (332.550,397.888);
\draw[color=tg0000ff,line width=0.3bp,line cap=round] (332.117,400.355) -- (333.217,401.455);
\draw[color=tg0000ff,line width=0.3bp,line cap=round] (332.117,401.455) -- (333.217,400.355);
\draw[color=tg0000ff,line width=0.3bp,line cap=round] (332.783,398.678) -- (333.883,399.778);
\draw[color=tg0000ff,line width=0.3bp,line cap=round] (332.783,399.778) -- (333.883,398.678);
\draw[color=tg0000ff,line width=0.3bp,line cap=round] (333.450,399.362) -- (334.550,400.462);
\draw[color=tg0000ff,line width=0.3bp,line cap=round] (333.450,400.462) -- (334.550,399.362);
\draw[color=tg0000ff,line width=0.3bp,line cap=round] (334.117,399.656) -- (335.217,400.756);
\draw[color=tg0000ff,line width=0.3bp,line cap=round] (334.117,400.756) -- (335.217,399.656);
\node[anchor=west,inner sep=0pt,outer sep=0pt,text=tg626b73,font={\sffamily\fontsize{6.2bp}{7.44bp}\selectfont}] at (346.000,372.616) {None found};
\draw[color=tgeceef0,line width=0.4bp,line cap=round] (0.000,421.722) -- (448.000,421.722);
\node[anchor=west,inner sep=0pt,outer sep=0pt,text=tg18314b,font={\sffamily\fontsize{7.2bp}{8.64bp}\selectfont\bfseries}] at (0.000,428.542) {Irregular cycles (AR(2))};
\node[anchor=west,inner sep=0pt,outer sep=0pt,text=tg18314b,font={\sffamily\fontsize{6.2bp}{7.44bp}\selectfont}] at (0.000,436.982) {\ensuremath{\checkmark} The rest of the current swing, fading over the};
\node[anchor=west,inner sep=0pt,outer sep=0pt,text=tg18314b,font={\sffamily\fontsize{6.2bp}{7.44bp}\selectfont}] at (0.000,444.422) {next cycles.};
\node[anchor=west,inner sep=0pt,outer sep=0pt,text=tg18314b,font={\sffamily\fontsize{6.2bp}{7.44bp}\selectfont}] at (0.000,452.262) {\ensuremath{\times} When the next swings come and how big they};
\node[anchor=west,inner sep=0pt,outer sep=0pt,text=tg18314b,font={\sffamily\fontsize{6.2bp}{7.44bp}\selectfont}] at (0.000,459.702) {are; noise hides the current phase.};
\path[fill=tg8fd18f,draw=none,line width=0.5bp] (168.000,424.222) rectangle (205.969,432.822);
\node[anchor=center,inner sep=0pt,outer sep=0pt,text=tg0b0b0b,font={\sffamily\fontsize{6.6bp}{7.92bp}\selectfont\bfseries}] at (186.984,428.522) {72  Good};
\draw[color=tgc9cdd1,line width=0.45bp,line cap=round] (168.000,466.608) -- (248.000,466.608);
\draw[color=tg9da7b0,line width=0.5bp,line cap=round,dash pattern=on 2bp off 2bp] (232.000,435.222) -- (232.000,466.608);
\path[fill=tg808080,draw=none] (168.333,439.600) circle[radius=0.42bp] (169.000,441.625) circle[radius=0.42bp] (169.667,444.088) circle[radius=0.42bp] (170.333,449.168) circle[radius=0.42bp] (171.000,453.865) circle[radius=0.42bp] (171.667,458.347) circle[radius=0.42bp] (172.333,462.779) circle[radius=0.42bp] (173.000,462.497) circle[radius=0.42bp] (173.667,464.927) circle[radius=0.42bp] (174.333,461.415) circle[radius=0.42bp] (175.000,457.900) circle[radius=0.42bp] (175.667,454.270) circle[radius=0.42bp] (176.333,449.320) circle[radius=0.42bp] (177.000,447.126) circle[radius=0.42bp] (177.667,444.810) circle[radius=0.42bp] (178.333,441.627) circle[radius=0.42bp] (179.000,442.727) circle[radius=0.42bp] (179.667,441.622) circle[radius=0.42bp] (180.333,443.073) circle[radius=0.42bp] (181.000,444.894) circle[radius=0.42bp] (181.667,447.982) circle[radius=0.42bp] (182.333,450.217) circle[radius=0.42bp] (183.000,449.949) circle[radius=0.42bp] (183.667,450.936) circle[radius=0.42bp] (184.333,449.842) circle[radius=0.42bp] (185.000,446.362) circle[radius=0.42bp] (185.667,447.329) circle[radius=0.42bp] (186.333,447.276) circle[radius=0.42bp] (187.000,448.384) circle[radius=0.42bp] (187.667,448.626) circle[radius=0.42bp] (188.333,452.084) circle[radius=0.42bp] (189.000,453.957) circle[radius=0.42bp] (189.667,455.895) circle[radius=0.42bp] (190.333,454.033) circle[radius=0.42bp] (191.000,454.551) circle[radius=0.42bp] (191.667,456.155) circle[radius=0.42bp] (192.333,455.906) circle[radius=0.42bp] (193.000,457.665) circle[radius=0.42bp] (193.667,455.384) circle[radius=0.42bp] (194.333,452.352) circle[radius=0.42bp] (195.000,451.470) circle[radius=0.42bp] (195.667,448.228) circle[radius=0.42bp] (196.333,448.380) circle[radius=0.42bp] (197.000,444.644) circle[radius=0.42bp] (197.667,443.374) circle[radius=0.42bp] (198.333,444.092) circle[radius=0.42bp] (199.000,446.787) circle[radius=0.42bp] (199.667,449.427) circle[radius=0.42bp] (200.333,451.247) circle[radius=0.42bp] (201.000,450.145) circle[radius=0.42bp] (201.667,452.992) circle[radius=0.42bp] (202.333,455.736) circle[radius=0.42bp] (203.000,455.871) circle[radius=0.42bp] (203.667,454.631) circle[radius=0.42bp] (204.333,454.906) circle[radius=0.42bp] (205.000,450.740) circle[radius=0.42bp] (205.667,448.476) circle[radius=0.42bp] (206.333,447.992) circle[radius=0.42bp] (207.000,447.849) circle[radius=0.42bp] (207.667,449.443) circle[radius=0.42bp] (208.333,451.746) circle[radius=0.42bp] (209.000,454.253) circle[radius=0.42bp] (209.667,453.105) circle[radius=0.42bp] (210.333,452.328) circle[radius=0.42bp] (211.000,454.573) circle[radius=0.42bp] (211.667,452.809) circle[radius=0.42bp] (212.333,451.973) circle[radius=0.42bp] (213.000,449.033) circle[radius=0.42bp] (213.667,449.588) circle[radius=0.42bp] (214.333,448.207) circle[radius=0.42bp] (215.000,446.876) circle[radius=0.42bp] (215.667,448.532) circle[radius=0.42bp] (216.333,448.229) circle[radius=0.42bp] (217.000,450.829) circle[radius=0.42bp] (217.667,451.645) circle[radius=0.42bp] (218.333,456.259) circle[radius=0.42bp] (219.000,458.655) circle[radius=0.42bp] (219.667,456.227) circle[radius=0.42bp] (220.333,457.710) circle[radius=0.42bp] (221.000,457.706) circle[radius=0.42bp] (221.667,454.445) circle[radius=0.42bp] (222.333,453.372) circle[radius=0.42bp] (223.000,451.676) circle[radius=0.42bp] (223.667,451.960) circle[radius=0.42bp] (224.333,451.264) circle[radius=0.42bp] (225.000,451.530) circle[radius=0.42bp] (225.667,453.453) circle[radius=0.42bp] (226.333,453.138) circle[radius=0.42bp] (227.000,454.116) circle[radius=0.42bp] (227.667,451.270) circle[radius=0.42bp] (228.333,452.750) circle[radius=0.42bp] (229.000,452.198) circle[radius=0.42bp] (229.667,451.219) circle[radius=0.42bp] (230.333,448.916) circle[radius=0.42bp] (231.000,446.864) circle[radius=0.42bp] (231.667,444.407) circle[radius=0.42bp] (232.333,444.232) circle[radius=0.42bp] (233.000,443.210) circle[radius=0.42bp] (233.667,444.824) circle[radius=0.42bp] (234.333,447.251) circle[radius=0.42bp] (235.000,447.830) circle[radius=0.42bp] (235.667,453.261) circle[radius=0.42bp] (236.333,459.204) circle[radius=0.42bp] (237.000,462.470) circle[radius=0.42bp] (237.667,463.764) circle[radius=0.42bp] (238.333,463.613) circle[radius=0.42bp] (239.000,463.244) circle[radius=0.42bp] (239.667,461.676) circle[radius=0.42bp] (240.333,458.471) circle[radius=0.42bp] (241.000,451.065) circle[radius=0.42bp] (241.667,447.109) circle[radius=0.42bp] (242.333,443.281) circle[radius=0.42bp] (243.000,439.247) circle[radius=0.42bp] (243.667,436.903) circle[radius=0.42bp] (244.333,440.701) circle[radius=0.42bp] (245.000,440.853) circle[radius=0.42bp] (245.667,445.517) circle[radius=0.42bp] (246.333,449.503) circle[radius=0.42bp] (247.000,453.007) circle[radius=0.42bp] (247.667,454.902) circle[radius=0.42bp];
\draw[color=tg000000,line width=0.7bp,line cap=round] (232.333,444.342) -- (233.000,444.882) -- (233.667,446.383) -- (234.333,448.531) -- (235.000,450.948) -- (235.667,453.251) -- (236.333,455.112) -- (237.000,456.302) -- (237.667,456.709) -- (238.333,456.351) -- (239.000,455.355) -- (239.667,453.930) -- (240.333,452.327) -- (241.000,450.799) -- (241.667,449.564) -- (242.333,448.775) -- (243.000,448.505) -- (243.667,448.742) -- (244.333,449.403) -- (245.000,450.348) -- (245.667,451.412) -- (246.333,452.426) -- (247.000,453.245) -- (247.667,453.768);
\draw[color=tg0000ff,line width=0.3bp,line cap=round] (231.783,443.304) -- (232.883,444.404);
\draw[color=tg0000ff,line width=0.3bp,line cap=round] (231.783,444.404) -- (232.883,443.304);
\draw[color=tg0000ff,line width=0.3bp,line cap=round] (232.450,443.972) -- (233.550,445.072);
\draw[color=tg0000ff,line width=0.3bp,line cap=round] (232.450,445.072) -- (233.550,443.972);
\draw[color=tg0000ff,line width=0.3bp,line cap=round] (233.117,445.544) -- (234.217,446.644);
\draw[color=tg0000ff,line width=0.3bp,line cap=round] (233.117,446.644) -- (234.217,445.544);
\draw[color=tg0000ff,line width=0.3bp,line cap=round] (233.783,446.693) -- (234.883,447.793);
\draw[color=tg0000ff,line width=0.3bp,line cap=round] (233.783,447.793) -- (234.883,446.693);
\draw[color=tg0000ff,line width=0.3bp,line cap=round] (234.450,449.759) -- (235.550,450.859);
\draw[color=tg0000ff,line width=0.3bp,line cap=round] (234.450,450.859) -- (235.550,449.759);
\draw[color=tg0000ff,line width=0.3bp,line cap=round] (235.117,452.459) -- (236.217,453.559);
\draw[color=tg0000ff,line width=0.3bp,line cap=round] (235.117,453.559) -- (236.217,452.459);
\draw[color=tg0000ff,line width=0.3bp,line cap=round] (235.783,454.851) -- (236.883,455.951);
\draw[color=tg0000ff,line width=0.3bp,line cap=round] (235.783,455.951) -- (236.883,454.851);
\draw[color=tg0000ff,line width=0.3bp,line cap=round] (236.450,455.979) -- (237.550,457.079);
\draw[color=tg0000ff,line width=0.3bp,line cap=round] (236.450,457.079) -- (237.550,455.979);
\draw[color=tg0000ff,line width=0.3bp,line cap=round] (237.117,456.510) -- (238.217,457.610);
\draw[color=tg0000ff,line width=0.3bp,line cap=round] (237.117,457.610) -- (238.217,456.510);
\draw[color=tg0000ff,line width=0.3bp,line cap=round] (237.783,456.535) -- (238.883,457.635);
\draw[color=tg0000ff,line width=0.3bp,line cap=round] (237.783,457.635) -- (238.883,456.535);
\draw[color=tg0000ff,line width=0.3bp,line cap=round] (238.450,455.482) -- (239.550,456.582);
\draw[color=tg0000ff,line width=0.3bp,line cap=round] (238.450,456.582) -- (239.550,455.482);
\draw[color=tg0000ff,line width=0.3bp,line cap=round] (239.117,454.308) -- (240.217,455.408);
\draw[color=tg0000ff,line width=0.3bp,line cap=round] (239.117,455.408) -- (240.217,454.308);
\draw[color=tg0000ff,line width=0.3bp,line cap=round] (239.783,452.518) -- (240.883,453.618);
\draw[color=tg0000ff,line width=0.3bp,line cap=round] (239.783,453.618) -- (240.883,452.518);
\draw[color=tg0000ff,line width=0.3bp,line cap=round] (240.450,450.521) -- (241.550,451.621);
\draw[color=tg0000ff,line width=0.3bp,line cap=round] (240.450,451.621) -- (241.550,450.521);
\draw[color=tg0000ff,line width=0.3bp,line cap=round] (241.117,448.621) -- (242.217,449.721);
\draw[color=tg0000ff,line width=0.3bp,line cap=round] (241.117,449.721) -- (242.217,448.621);
\draw[color=tg0000ff,line width=0.3bp,line cap=round] (241.783,447.117) -- (242.883,448.217);
\draw[color=tg0000ff,line width=0.3bp,line cap=round] (241.783,448.217) -- (242.883,447.117);
\draw[color=tg0000ff,line width=0.3bp,line cap=round] (242.450,446.805) -- (243.550,447.905);
\draw[color=tg0000ff,line width=0.3bp,line cap=round] (242.450,447.905) -- (243.550,446.805);
\draw[color=tg0000ff,line width=0.3bp,line cap=round] (243.117,446.968) -- (244.217,448.068);
\draw[color=tg0000ff,line width=0.3bp,line cap=round] (243.117,448.068) -- (244.217,446.968);
\draw[color=tg0000ff,line width=0.3bp,line cap=round] (243.783,447.745) -- (244.883,448.845);
\draw[color=tg0000ff,line width=0.3bp,line cap=round] (243.783,448.845) -- (244.883,447.745);
\draw[color=tg0000ff,line width=0.3bp,line cap=round] (244.450,449.002) -- (245.550,450.102);
\draw[color=tg0000ff,line width=0.3bp,line cap=round] (244.450,450.102) -- (245.550,449.002);
\draw[color=tg0000ff,line width=0.3bp,line cap=round] (245.117,450.198) -- (246.217,451.298);
\draw[color=tg0000ff,line width=0.3bp,line cap=round] (245.117,451.298) -- (246.217,450.198);
\draw[color=tg0000ff,line width=0.3bp,line cap=round] (245.783,451.210) -- (246.883,452.310);
\draw[color=tg0000ff,line width=0.3bp,line cap=round] (245.783,452.310) -- (246.883,451.210);
\draw[color=tg0000ff,line width=0.3bp,line cap=round] (246.450,451.733) -- (247.550,452.833);
\draw[color=tg0000ff,line width=0.3bp,line cap=round] (246.450,452.833) -- (247.550,451.733);
\draw[color=tg0000ff,line width=0.3bp,line cap=round] (247.117,452.029) -- (248.217,453.129);
\draw[color=tg0000ff,line width=0.3bp,line cap=round] (247.117,453.129) -- (248.217,452.029);
\path[fill=tg8fd18f,draw=none,line width=0.5bp] (255.000,424.222) rectangle (292.969,432.822);
\node[anchor=center,inner sep=0pt,outer sep=0pt,text=tg0b0b0b,font={\sffamily\fontsize{6.6bp}{7.92bp}\selectfont\bfseries}] at (273.984,428.522) {82  Good};
\draw[color=tgc9cdd1,line width=0.45bp,line cap=round] (255.000,466.608) -- (335.000,466.608);
\draw[color=tg9da7b0,line width=0.5bp,line cap=round,dash pattern=on 2bp off 2bp] (319.000,435.222) -- (319.000,466.608);
\path[fill=tg808080,draw=none] (255.333,436.903) circle[radius=0.42bp] (256.000,440.778) circle[radius=0.42bp] (256.667,441.069) circle[radius=0.42bp] (257.333,446.571) circle[radius=0.42bp] (258.000,450.094) circle[radius=0.42bp] (258.667,455.543) circle[radius=0.42bp] (259.333,462.118) circle[radius=0.42bp] (260.000,456.627) circle[radius=0.42bp] (260.667,464.927) circle[radius=0.42bp] (261.333,458.360) circle[radius=0.42bp] (262.000,455.270) circle[radius=0.42bp] (262.667,453.914) circle[radius=0.42bp] (263.333,446.194) circle[radius=0.42bp] (264.000,446.018) circle[radius=0.42bp] (264.667,444.527) circle[radius=0.42bp] (265.333,439.790) circle[radius=0.42bp] (266.000,442.823) circle[radius=0.42bp] (266.667,437.963) circle[radius=0.42bp] (267.333,441.383) circle[radius=0.42bp] (268.000,444.677) circle[radius=0.42bp] (268.667,450.409) circle[radius=0.42bp] (269.333,450.991) circle[radius=0.42bp] (270.000,445.740) circle[radius=0.42bp] (270.667,448.973) circle[radius=0.42bp] (271.333,448.522) circle[radius=0.42bp] (272.000,441.496) circle[radius=0.42bp] (272.667,448.462) circle[radius=0.42bp] (273.333,446.954) circle[radius=0.42bp] (274.000,447.668) circle[radius=0.42bp] (274.667,445.654) circle[radius=0.42bp] (275.333,451.242) circle[radius=0.42bp] (276.000,452.744) circle[radius=0.42bp] (276.667,456.487) circle[radius=0.42bp] (277.333,449.993) circle[radius=0.42bp] (278.000,450.277) circle[radius=0.42bp] (278.667,454.342) circle[radius=0.42bp] (279.333,452.788) circle[radius=0.42bp] (280.000,457.349) circle[radius=0.42bp] (280.667,453.671) circle[radius=0.42bp] (281.333,447.191) circle[radius=0.42bp] (282.000,449.039) circle[radius=0.42bp] (282.667,441.698) circle[radius=0.42bp] (283.333,448.324) circle[radius=0.42bp] (284.000,442.252) circle[radius=0.42bp] (284.667,442.788) circle[radius=0.42bp] (285.333,441.808) circle[radius=0.42bp] (286.000,445.321) circle[radius=0.42bp] (286.667,448.748) circle[radius=0.42bp] (287.333,450.952) circle[radius=0.42bp] (288.000,441.758) circle[radius=0.42bp] (288.667,450.792) circle[radius=0.42bp] (289.333,457.213) circle[radius=0.42bp] (290.000,454.460) circle[radius=0.42bp] (290.667,450.614) circle[radius=0.42bp] (291.333,453.320) circle[radius=0.42bp] (292.000,445.845) circle[radius=0.42bp] (292.667,444.654) circle[radius=0.42bp] (293.333,446.660) circle[radius=0.42bp] (294.000,447.208) circle[radius=0.42bp] (294.667,449.692) circle[radius=0.42bp] (295.333,451.028) circle[radius=0.42bp] (296.000,454.801) circle[radius=0.42bp] (296.667,448.201) circle[radius=0.42bp] (297.333,446.725) circle[radius=0.42bp] (298.000,454.621) circle[radius=0.42bp] (298.667,452.951) circle[radius=0.42bp] (299.333,453.465) circle[radius=0.42bp] (300.000,449.254) circle[radius=0.42bp] (300.667,454.223) circle[radius=0.42bp] (301.333,448.920) circle[radius=0.42bp] (302.000,442.663) circle[radius=0.42bp] (302.667,447.529) circle[radius=0.42bp] (303.333,444.751) circle[radius=0.42bp] (304.000,449.799) circle[radius=0.42bp] (304.667,445.783) circle[radius=0.42bp] (305.333,453.015) circle[radius=0.42bp] (306.000,456.290) circle[radius=0.42bp] (306.667,448.053) circle[radius=0.42bp] (307.333,455.359) circle[radius=0.42bp] (308.000,457.235) circle[radius=0.42bp] (308.667,451.453) circle[radius=0.42bp] (309.333,451.037) circle[radius=0.42bp] (310.000,447.265) circle[radius=0.42bp] (310.667,451.677) circle[radius=0.42bp] (311.333,447.460) circle[radius=0.42bp] (312.000,447.588) circle[radius=0.42bp] (312.667,452.426) circle[radius=0.42bp] (313.333,450.413) circle[radius=0.42bp] (314.000,453.336) circle[radius=0.42bp] (314.667,444.234) circle[radius=0.42bp] (315.333,451.899) circle[radius=0.42bp] (316.000,450.899) circle[radius=0.42bp] (316.667,451.344) circle[radius=0.42bp] (317.333,447.859) circle[radius=0.42bp] (318.000,445.385) circle[radius=0.42bp] (318.667,441.655) circle[radius=0.42bp] (319.333,444.395) circle[radius=0.42bp] (320.000,442.627) circle[radius=0.42bp] (320.667,446.593) circle[radius=0.42bp] (321.333,447.439) circle[radius=0.42bp] (322.000,440.422) circle[radius=0.42bp] (322.667,448.678) circle[radius=0.42bp] (323.333,457.598) circle[radius=0.42bp] (324.000,461.566) circle[radius=0.42bp] (324.667,461.918) circle[radius=0.42bp] (325.333,459.788) circle[radius=0.42bp] (326.000,459.353) circle[radius=0.42bp] (326.667,459.642) circle[radius=0.42bp] (327.333,458.029) circle[radius=0.42bp] (328.000,445.994) circle[radius=0.42bp] (328.667,444.635) circle[radius=0.42bp] (329.333,443.085) circle[radius=0.42bp] (330.000,439.301) circle[radius=0.42bp] (330.667,437.938) circle[radius=0.42bp] (331.333,446.242) circle[radius=0.42bp] (332.000,439.973) circle[radius=0.42bp] (332.667,446.462) circle[radius=0.42bp] (333.333,449.739) circle[radius=0.42bp] (334.000,451.986) circle[radius=0.42bp] (334.667,450.733) circle[radius=0.42bp];
\draw[color=tg000000,line width=0.7bp,line cap=round] (319.333,444.924) -- (320.000,445.322) -- (320.667,446.309) -- (321.333,447.682) -- (322.000,449.202) -- (322.667,450.630) -- (323.333,451.766) -- (324.000,452.470) -- (324.667,452.682) -- (325.333,452.417) -- (326.000,451.763) -- (326.667,450.852) -- (327.333,449.843) -- (328.000,448.896) -- (328.667,448.143) -- (329.333,447.675) -- (330.000,447.535) -- (330.667,447.710) -- (331.333,448.145) -- (332.000,448.749) -- (332.667,449.418) -- (333.333,450.046) -- (334.000,450.546) -- (334.667,450.856);
\draw[color=tg0000ff,line width=0.3bp,line cap=round] (318.783,442.596) -- (319.883,443.696);
\draw[color=tg0000ff,line width=0.3bp,line cap=round] (318.783,443.696) -- (319.883,442.596);
\draw[color=tg0000ff,line width=0.3bp,line cap=round] (319.450,443.921) -- (320.550,445.021);
\draw[color=tg0000ff,line width=0.3bp,line cap=round] (319.450,445.021) -- (320.550,443.921);
\draw[color=tg0000ff,line width=0.3bp,line cap=round] (320.117,445.560) -- (321.217,446.660);
\draw[color=tg0000ff,line width=0.3bp,line cap=round] (320.117,446.660) -- (321.217,445.560);
\draw[color=tg0000ff,line width=0.3bp,line cap=round] (320.783,444.374) -- (321.883,445.474);
\draw[color=tg0000ff,line width=0.3bp,line cap=round] (320.783,445.474) -- (321.883,444.374);
\draw[color=tg0000ff,line width=0.3bp,line cap=round] (321.450,449.466) -- (322.550,450.566);
\draw[color=tg0000ff,line width=0.3bp,line cap=round] (321.450,450.566) -- (322.550,449.466);
\draw[color=tg0000ff,line width=0.3bp,line cap=round] (322.117,448.967) -- (323.217,450.067);
\draw[color=tg0000ff,line width=0.3bp,line cap=round] (322.117,450.067) -- (323.217,448.967);
\draw[color=tg0000ff,line width=0.3bp,line cap=round] (322.783,453.002) -- (323.883,454.102);
\draw[color=tg0000ff,line width=0.3bp,line cap=round] (322.783,454.102) -- (323.883,453.002);
\draw[color=tg0000ff,line width=0.3bp,line cap=round] (323.450,452.960) -- (324.550,454.060);
\draw[color=tg0000ff,line width=0.3bp,line cap=round] (323.450,454.060) -- (324.550,452.960);
\draw[color=tg0000ff,line width=0.3bp,line cap=round] (324.117,454.031) -- (325.217,455.131);
\draw[color=tg0000ff,line width=0.3bp,line cap=round] (324.117,455.131) -- (325.217,454.031);
\draw[color=tg0000ff,line width=0.3bp,line cap=round] (324.783,454.476) -- (325.883,455.576);
\draw[color=tg0000ff,line width=0.3bp,line cap=round] (324.783,455.576) -- (325.883,454.476);
\draw[color=tg0000ff,line width=0.3bp,line cap=round] (325.450,452.079) -- (326.550,453.179);
\draw[color=tg0000ff,line width=0.3bp,line cap=round] (325.450,453.179) -- (326.550,452.079);
\draw[color=tg0000ff,line width=0.3bp,line cap=round] (326.117,452.868) -- (327.217,453.968);
\draw[color=tg0000ff,line width=0.3bp,line cap=round] (326.117,453.968) -- (327.217,452.868);
\draw[color=tg0000ff,line width=0.3bp,line cap=round] (326.783,449.958) -- (327.883,451.058);
\draw[color=tg0000ff,line width=0.3bp,line cap=round] (326.783,451.058) -- (327.883,449.958);
\draw[color=tg0000ff,line width=0.3bp,line cap=round] (327.450,448.095) -- (328.550,449.195);
\draw[color=tg0000ff,line width=0.3bp,line cap=round] (327.450,449.195) -- (328.550,448.095);
\draw[color=tg0000ff,line width=0.3bp,line cap=round] (328.117,446.600) -- (329.217,447.700);
\draw[color=tg0000ff,line width=0.3bp,line cap=round] (328.117,447.700) -- (329.217,446.600);
\draw[color=tg0000ff,line width=0.3bp,line cap=round] (328.783,444.442) -- (329.883,445.542);
\draw[color=tg0000ff,line width=0.3bp,line cap=round] (328.783,445.542) -- (329.883,444.442);
\draw[color=tg0000ff,line width=0.3bp,line cap=round] (329.450,445.389) -- (330.550,446.489);
\draw[color=tg0000ff,line width=0.3bp,line cap=round] (329.450,446.489) -- (330.550,445.389);
\draw[color=tg0000ff,line width=0.3bp,line cap=round] (330.117,443.906) -- (331.217,445.006);
\draw[color=tg0000ff,line width=0.3bp,line cap=round] (330.117,445.006) -- (331.217,443.906);
\draw[color=tg0000ff,line width=0.3bp,line cap=round] (330.783,445.804) -- (331.883,446.904);
\draw[color=tg0000ff,line width=0.3bp,line cap=round] (330.783,446.904) -- (331.883,445.804);
\draw[color=tg0000ff,line width=0.3bp,line cap=round] (331.450,447.218) -- (332.550,448.318);
\draw[color=tg0000ff,line width=0.3bp,line cap=round] (331.450,448.318) -- (332.550,447.218);
\draw[color=tg0000ff,line width=0.3bp,line cap=round] (332.117,447.713) -- (333.217,448.813);
\draw[color=tg0000ff,line width=0.3bp,line cap=round] (332.117,448.813) -- (333.217,447.713);
\draw[color=tg0000ff,line width=0.3bp,line cap=round] (332.783,449.271) -- (333.883,450.371);
\draw[color=tg0000ff,line width=0.3bp,line cap=round] (332.783,450.371) -- (333.883,449.271);
\draw[color=tg0000ff,line width=0.3bp,line cap=round] (333.450,449.192) -- (334.550,450.292);
\draw[color=tg0000ff,line width=0.3bp,line cap=round] (333.450,450.292) -- (334.550,449.192);
\draw[color=tg0000ff,line width=0.3bp,line cap=round] (334.117,450.287) -- (335.217,451.387);
\draw[color=tg0000ff,line width=0.3bp,line cap=round] (334.117,451.387) -- (335.217,450.287);
\node[anchor=west,inner sep=0pt,outer sep=0pt,text=tg626b73,font={\sffamily\fontsize{6.2bp}{7.44bp}\selectfont}] at (346.000,427.942) {None found};
\draw[color=tgeceef0,line width=0.4bp,line cap=round] (0.000,469.608) -- (448.000,469.608);
\node[anchor=west,inner sep=0pt,outer sep=0pt,text=tg18314b,font={\sffamily\fontsize{7.2bp}{8.64bp}\selectfont\bfseries}] at (0.000,476.428) {Random walk: shocks never fade};
\node[anchor=west,inner sep=0pt,outer sep=0pt,text=tg18314b,font={\sffamily\fontsize{6.2bp}{7.44bp}\selectfont}] at (0.000,484.868) {\ensuremath{\checkmark} Where it is now: the best forecast is flat at the};
\node[anchor=west,inner sep=0pt,outer sep=0pt,text=tg18314b,font={\sffamily\fontsize{6.2bp}{7.44bp}\selectfont}] at (0.000,492.308) {current level, smoothed when noisy.};
\node[anchor=west,inner sep=0pt,outer sep=0pt,text=tg18314b,font={\sffamily\fontsize{6.2bp}{7.44bp}\selectfont}] at (0.000,500.148) {\ensuremath{\times} The direction of the next move. Uncertainty};
\node[anchor=west,inner sep=0pt,outer sep=0pt,text=tg18314b,font={\sffamily\fontsize{6.2bp}{7.44bp}\selectfont}] at (0.000,507.588) {grows without limit with the horizon.};
\path[fill=tgfab219,draw=none,line width=0.5bp] (168.000,472.108) rectangle (222.719,480.708);
\node[anchor=center,inner sep=0pt,outer sep=0pt,text=tg0b0b0b,font={\sffamily\fontsize{6.6bp}{7.92bp}\selectfont\bfseries}] at (195.359,476.408) {56  Struggles};
\draw[color=tgc9cdd1,line width=0.45bp,line cap=round] (168.000,514.494) -- (248.000,514.494);
\draw[color=tg9da7b0,line width=0.5bp,line cap=round,dash pattern=on 2bp off 2bp] (232.000,483.108) -- (232.000,514.494);
\path[fill=tg808080,draw=none] (168.333,503.621) circle[radius=0.42bp] (169.000,504.547) circle[radius=0.42bp] (169.667,505.413) circle[radius=0.42bp] (170.333,508.241) circle[radius=0.42bp] (171.000,507.642) circle[radius=0.42bp] (171.667,508.116) circle[radius=0.42bp] (172.333,509.235) circle[radius=0.42bp] (173.000,510.655) circle[radius=0.42bp] (173.667,512.813) circle[radius=0.42bp] (174.333,508.378) circle[radius=0.42bp] (175.000,512.384) circle[radius=0.42bp] (175.667,509.396) circle[radius=0.42bp] (176.333,510.297) circle[radius=0.42bp] (177.000,511.707) circle[radius=0.42bp] (177.667,508.152) circle[radius=0.42bp] (178.333,507.161) circle[radius=0.42bp] (179.000,509.340) circle[radius=0.42bp] (179.667,505.740) circle[radius=0.42bp] (180.333,505.231) circle[radius=0.42bp] (181.000,501.374) circle[radius=0.42bp] (181.667,502.501) circle[radius=0.42bp] (182.333,504.161) circle[radius=0.42bp] (183.000,504.179) circle[radius=0.42bp] (183.667,500.395) circle[radius=0.42bp] (184.333,494.280) circle[radius=0.42bp] (185.000,494.021) circle[radius=0.42bp] (185.667,493.840) circle[radius=0.42bp] (186.333,492.793) circle[radius=0.42bp] (187.000,497.063) circle[radius=0.42bp] (187.667,495.091) circle[radius=0.42bp] (188.333,496.804) circle[radius=0.42bp] (189.000,494.685) circle[radius=0.42bp] (189.667,495.880) circle[radius=0.42bp] (190.333,494.792) circle[radius=0.42bp] (191.000,497.963) circle[radius=0.42bp] (191.667,495.525) circle[radius=0.42bp] (192.333,496.049) circle[radius=0.42bp] (193.000,499.623) circle[radius=0.42bp] (193.667,496.156) circle[radius=0.42bp] (194.333,499.000) circle[radius=0.42bp] (195.000,497.325) circle[radius=0.42bp] (195.667,495.006) circle[radius=0.42bp] (196.333,494.509) circle[radius=0.42bp] (197.000,489.416) circle[radius=0.42bp] (197.667,493.007) circle[radius=0.42bp] (198.333,494.225) circle[radius=0.42bp] (199.000,495.280) circle[radius=0.42bp] (199.667,493.174) circle[radius=0.42bp] (200.333,492.546) circle[radius=0.42bp] (201.000,494.254) circle[radius=0.42bp] (201.667,492.046) circle[radius=0.42bp] (202.333,488.238) circle[radius=0.42bp] (203.000,493.782) circle[radius=0.42bp] (203.667,495.245) circle[radius=0.42bp] (204.333,493.180) circle[radius=0.42bp] (205.000,488.329) circle[radius=0.42bp] (205.667,489.771) circle[radius=0.42bp] (206.333,487.983) circle[radius=0.42bp] (207.000,489.055) circle[radius=0.42bp] (207.667,492.328) circle[radius=0.42bp] (208.333,494.656) circle[radius=0.42bp] (209.000,494.513) circle[radius=0.42bp] (209.667,493.222) circle[radius=0.42bp] (210.333,492.462) circle[radius=0.42bp] (211.000,489.660) circle[radius=0.42bp] (211.667,486.636) circle[radius=0.42bp] (212.333,491.203) circle[radius=0.42bp] (213.000,488.800) circle[radius=0.42bp] (213.667,490.286) circle[radius=0.42bp] (214.333,490.185) circle[radius=0.42bp] (215.000,492.943) circle[radius=0.42bp] (215.667,488.578) circle[radius=0.42bp] (216.333,484.789) circle[radius=0.42bp] (217.000,488.036) circle[radius=0.42bp] (217.667,488.242) circle[radius=0.42bp] (218.333,492.242) circle[radius=0.42bp] (219.000,488.231) circle[radius=0.42bp] (219.667,489.064) circle[radius=0.42bp] (220.333,490.150) circle[radius=0.42bp] (221.000,487.440) circle[radius=0.42bp] (221.667,490.360) circle[radius=0.42bp] (222.333,493.532) circle[radius=0.42bp] (223.000,492.498) circle[radius=0.42bp] (223.667,491.944) circle[radius=0.42bp] (224.333,493.612) circle[radius=0.42bp] (225.000,495.946) circle[radius=0.42bp] (225.667,494.278) circle[radius=0.42bp] (226.333,494.049) circle[radius=0.42bp] (227.000,496.374) circle[radius=0.42bp] (227.667,494.869) circle[radius=0.42bp] (228.333,495.306) circle[radius=0.42bp] (229.000,491.782) circle[radius=0.42bp] (229.667,494.156) circle[radius=0.42bp] (230.333,492.732) circle[radius=0.42bp] (231.000,492.395) circle[radius=0.42bp] (231.667,489.649) circle[radius=0.42bp] (232.333,488.006) circle[radius=0.42bp] (233.000,485.649) circle[radius=0.42bp] (233.667,487.547) circle[radius=0.42bp] (234.333,489.007) circle[radius=0.42bp] (235.000,490.997) circle[radius=0.42bp] (235.667,491.260) circle[radius=0.42bp] (236.333,488.112) circle[radius=0.42bp] (237.000,490.234) circle[radius=0.42bp] (237.667,493.718) circle[radius=0.42bp] (238.333,496.029) circle[radius=0.42bp] (239.000,497.021) circle[radius=0.42bp] (239.667,495.240) circle[radius=0.42bp] (240.333,494.635) circle[radius=0.42bp] (241.000,493.844) circle[radius=0.42bp] (241.667,494.654) circle[radius=0.42bp] (242.333,489.189) circle[radius=0.42bp] (243.000,488.560) circle[radius=0.42bp] (243.667,488.088) circle[radius=0.42bp] (244.333,491.310) circle[radius=0.42bp] (245.000,489.970) circle[radius=0.42bp] (245.667,492.981) circle[radius=0.42bp] (246.333,487.917) circle[radius=0.42bp] (247.000,489.105) circle[radius=0.42bp] (247.667,488.448) circle[radius=0.42bp];
\draw[color=tg000000,line width=0.7bp,line cap=round] (232.333,490.804) -- (233.000,490.804) -- (233.667,490.804) -- (234.333,490.804) -- (235.000,490.804) -- (235.667,490.804) -- (236.333,490.804) -- (237.000,490.804) -- (237.667,490.804) -- (238.333,490.804) -- (239.000,490.804) -- (239.667,490.804) -- (240.333,490.804) -- (241.000,490.804) -- (241.667,490.804) -- (242.333,490.804) -- (243.000,490.804) -- (243.667,490.804) -- (244.333,490.804) -- (245.000,490.804) -- (245.667,490.804) -- (246.333,490.804) -- (247.000,490.804) -- (247.667,490.804);
\draw[color=tg0000ff,line width=0.3bp,line cap=round] (231.783,489.689) -- (232.883,490.789);
\draw[color=tg0000ff,line width=0.3bp,line cap=round] (231.783,490.789) -- (232.883,489.689);
\draw[color=tg0000ff,line width=0.3bp,line cap=round] (232.450,489.920) -- (233.550,491.020);
\draw[color=tg0000ff,line width=0.3bp,line cap=round] (232.450,491.020) -- (233.550,489.920);
\draw[color=tg0000ff,line width=0.3bp,line cap=round] (233.117,489.262) -- (234.217,490.362);
\draw[color=tg0000ff,line width=0.3bp,line cap=round] (233.117,490.362) -- (234.217,489.262);
\draw[color=tg0000ff,line width=0.3bp,line cap=round] (233.783,490.084) -- (234.883,491.184);
\draw[color=tg0000ff,line width=0.3bp,line cap=round] (233.783,491.184) -- (234.883,490.084);
\draw[color=tg0000ff,line width=0.3bp,line cap=round] (234.450,489.935) -- (235.550,491.035);
\draw[color=tg0000ff,line width=0.3bp,line cap=round] (234.450,491.035) -- (235.550,489.935);
\draw[color=tg0000ff,line width=0.3bp,line cap=round] (235.117,489.301) -- (236.217,490.401);
\draw[color=tg0000ff,line width=0.3bp,line cap=round] (235.117,490.401) -- (236.217,489.301);
\draw[color=tg0000ff,line width=0.3bp,line cap=round] (235.783,489.978) -- (236.883,491.078);
\draw[color=tg0000ff,line width=0.3bp,line cap=round] (235.783,491.078) -- (236.883,489.978);
\draw[color=tg0000ff,line width=0.3bp,line cap=round] (236.450,489.721) -- (237.550,490.821);
\draw[color=tg0000ff,line width=0.3bp,line cap=round] (236.450,490.821) -- (237.550,489.721);
\draw[color=tg0000ff,line width=0.3bp,line cap=round] (237.117,490.207) -- (238.217,491.307);
\draw[color=tg0000ff,line width=0.3bp,line cap=round] (237.117,491.307) -- (238.217,490.207);
\draw[color=tg0000ff,line width=0.3bp,line cap=round] (237.783,489.946) -- (238.883,491.046);
\draw[color=tg0000ff,line width=0.3bp,line cap=round] (237.783,491.046) -- (238.883,489.946);
\draw[color=tg0000ff,line width=0.3bp,line cap=round] (238.450,490.561) -- (239.550,491.661);
\draw[color=tg0000ff,line width=0.3bp,line cap=round] (238.450,491.661) -- (239.550,490.561);
\draw[color=tg0000ff,line width=0.3bp,line cap=round] (239.117,490.894) -- (240.217,491.994);
\draw[color=tg0000ff,line width=0.3bp,line cap=round] (239.117,491.994) -- (240.217,490.894);
\draw[color=tg0000ff,line width=0.3bp,line cap=round] (239.783,490.148) -- (240.883,491.248);
\draw[color=tg0000ff,line width=0.3bp,line cap=round] (239.783,491.248) -- (240.883,490.148);
\draw[color=tg0000ff,line width=0.3bp,line cap=round] (240.450,490.162) -- (241.550,491.262);
\draw[color=tg0000ff,line width=0.3bp,line cap=round] (240.450,491.262) -- (241.550,490.162);
\draw[color=tg0000ff,line width=0.3bp,line cap=round] (241.117,489.921) -- (242.217,491.021);
\draw[color=tg0000ff,line width=0.3bp,line cap=round] (241.117,491.021) -- (242.217,489.921);
\draw[color=tg0000ff,line width=0.3bp,line cap=round] (241.783,490.372) -- (242.883,491.472);
\draw[color=tg0000ff,line width=0.3bp,line cap=round] (241.783,491.472) -- (242.883,490.372);
\draw[color=tg0000ff,line width=0.3bp,line cap=round] (242.450,489.949) -- (243.550,491.049);
\draw[color=tg0000ff,line width=0.3bp,line cap=round] (242.450,491.049) -- (243.550,489.949);
\draw[color=tg0000ff,line width=0.3bp,line cap=round] (243.117,489.745) -- (244.217,490.845);
\draw[color=tg0000ff,line width=0.3bp,line cap=round] (243.117,490.845) -- (244.217,489.745);
\draw[color=tg0000ff,line width=0.3bp,line cap=round] (243.783,490.146) -- (244.883,491.246);
\draw[color=tg0000ff,line width=0.3bp,line cap=round] (243.783,491.246) -- (244.883,490.146);
\draw[color=tg0000ff,line width=0.3bp,line cap=round] (244.450,489.904) -- (245.550,491.004);
\draw[color=tg0000ff,line width=0.3bp,line cap=round] (244.450,491.004) -- (245.550,489.904);
\draw[color=tg0000ff,line width=0.3bp,line cap=round] (245.117,490.027) -- (246.217,491.127);
\draw[color=tg0000ff,line width=0.3bp,line cap=round] (245.117,491.127) -- (246.217,490.027);
\draw[color=tg0000ff,line width=0.3bp,line cap=round] (245.783,489.715) -- (246.883,490.815);
\draw[color=tg0000ff,line width=0.3bp,line cap=round] (245.783,490.815) -- (246.883,489.715);
\draw[color=tg0000ff,line width=0.3bp,line cap=round] (246.450,489.958) -- (247.550,491.058);
\draw[color=tg0000ff,line width=0.3bp,line cap=round] (246.450,491.058) -- (247.550,489.958);
\draw[color=tg0000ff,line width=0.3bp,line cap=round] (247.117,489.858) -- (248.217,490.958);
\draw[color=tg0000ff,line width=0.3bp,line cap=round] (247.117,490.958) -- (248.217,489.858);
\path[fill=tg8fd18f,draw=none,line width=0.5bp] (255.000,472.108) rectangle (292.969,480.708);
\node[anchor=center,inner sep=0pt,outer sep=0pt,text=tg0b0b0b,font={\sffamily\fontsize{6.6bp}{7.92bp}\selectfont\bfseries}] at (273.984,476.408) {87  Good};
\draw[color=tgc9cdd1,line width=0.45bp,line cap=round] (255.000,514.494) -- (335.000,514.494);
\draw[color=tg9da7b0,line width=0.5bp,line cap=round,dash pattern=on 2bp off 2bp] (319.000,483.108) -- (319.000,514.494);
\path[fill=tg808080,draw=none] (255.333,497.937) circle[radius=0.42bp] (256.000,500.444) circle[radius=0.42bp] (256.667,500.574) circle[radius=0.42bp] (257.333,505.010) circle[radius=0.42bp] (258.000,502.332) circle[radius=0.42bp] (258.667,504.074) circle[radius=0.42bp] (259.333,504.110) circle[radius=0.42bp] (260.000,507.013) circle[radius=0.42bp] (260.667,511.846) circle[radius=0.42bp] (261.333,502.481) circle[radius=0.42bp] (262.000,512.813) circle[radius=0.42bp] (262.667,506.479) circle[radius=0.42bp] (263.333,506.934) circle[radius=0.42bp] (264.000,510.025) circle[radius=0.42bp] (264.667,503.450) circle[radius=0.42bp] (265.333,505.013) circle[radius=0.42bp] (266.000,506.970) circle[radius=0.42bp] (266.667,502.230) circle[radius=0.42bp] (267.333,504.573) circle[radius=0.42bp] (268.000,497.256) circle[radius=0.42bp] (268.667,500.705) circle[radius=0.42bp] (269.333,503.915) circle[radius=0.42bp] (270.000,507.840) circle[radius=0.42bp] (270.667,503.840) circle[radius=0.42bp] (271.333,493.740) circle[radius=0.42bp] (272.000,496.611) circle[radius=0.42bp] (272.667,497.119) circle[radius=0.42bp] (273.333,491.353) circle[radius=0.42bp] (274.000,501.744) circle[radius=0.42bp] (274.667,498.687) circle[radius=0.42bp] (275.333,499.429) circle[radius=0.42bp] (276.000,495.319) circle[radius=0.42bp] (276.667,499.375) circle[radius=0.42bp] (277.333,498.579) circle[radius=0.42bp] (278.000,503.151) circle[radius=0.42bp] (278.667,495.450) circle[radius=0.42bp] (279.333,495.565) circle[radius=0.42bp] (280.000,501.156) circle[radius=0.42bp] (280.667,497.341) circle[radius=0.42bp] (281.333,502.951) circle[radius=0.42bp] (282.000,499.745) circle[radius=0.42bp] (282.667,493.407) circle[radius=0.42bp] (283.333,496.412) circle[radius=0.42bp] (284.000,487.519) circle[radius=0.42bp] (284.667,497.954) circle[radius=0.42bp] (285.333,495.055) circle[radius=0.42bp] (286.000,497.763) circle[radius=0.42bp] (286.667,494.454) circle[radius=0.42bp] (287.333,495.597) circle[radius=0.42bp] (288.000,498.112) circle[radius=0.42bp] (288.667,497.583) circle[radius=0.42bp] (289.333,484.789) circle[radius=0.42bp] (290.000,496.536) circle[radius=0.42bp] (290.667,502.591) circle[radius=0.42bp] (291.333,497.691) circle[radius=0.42bp] (292.000,491.204) circle[radius=0.42bp] (292.667,495.219) circle[radius=0.42bp] (293.333,489.176) circle[radius=0.42bp] (294.000,490.782) circle[radius=0.42bp] (294.667,495.850) circle[radius=0.42bp] (295.333,498.120) circle[radius=0.42bp] (296.000,499.451) circle[radius=0.42bp] (296.667,497.853) circle[radius=0.42bp] (297.333,499.449) circle[radius=0.42bp] (298.000,490.613) circle[radius=0.42bp] (298.667,487.742) circle[radius=0.42bp] (299.333,498.104) circle[radius=0.42bp] (300.000,496.442) circle[radius=0.42bp] (300.667,498.910) circle[radius=0.42bp] (301.333,496.706) circle[radius=0.42bp] (302.000,504.091) circle[radius=0.42bp] (302.667,496.202) circle[radius=0.42bp] (303.333,487.400) circle[radius=0.42bp] (304.000,493.751) circle[radius=0.42bp] (304.667,490.680) circle[radius=0.42bp] (305.333,496.694) circle[radius=0.42bp] (306.000,488.261) circle[radius=0.42bp] (306.667,492.926) circle[radius=0.42bp] (307.333,495.135) circle[radius=0.42bp] (308.000,485.671) circle[radius=0.42bp] (308.667,495.110) circle[radius=0.42bp] (309.333,499.430) circle[radius=0.42bp] (310.000,495.010) circle[radius=0.42bp] (310.667,495.314) circle[radius=0.42bp] (311.333,493.387) circle[radius=0.42bp] (312.000,500.102) circle[radius=0.42bp] (312.667,494.490) circle[radius=0.42bp] (313.333,494.222) circle[radius=0.42bp] (314.000,499.687) circle[radius=0.42bp] (314.667,496.558) circle[radius=0.42bp] (315.333,499.471) circle[radius=0.42bp] (316.000,488.866) circle[radius=0.42bp] (316.667,498.433) circle[radius=0.42bp] (317.333,496.897) circle[radius=0.42bp] (318.000,498.324) circle[radius=0.42bp] (318.667,494.735) circle[radius=0.42bp] (319.333,492.828) circle[radius=0.42bp] (320.000,489.332) circle[radius=0.42bp] (320.667,494.158) circle[radius=0.42bp] (321.333,493.916) circle[radius=0.42bp] (322.000,498.406) circle[radius=0.42bp] (322.667,496.997) circle[radius=0.42bp] (323.333,485.539) circle[radius=0.42bp] (324.000,491.397) circle[radius=0.42bp] (324.667,498.372) circle[radius=0.42bp] (325.333,501.260) circle[radius=0.42bp] (326.000,500.901) circle[radius=0.42bp] (326.667,497.276) circle[radius=0.42bp] (327.333,496.758) circle[radius=0.42bp] (328.000,498.350) circle[radius=0.42bp] (328.667,500.287) circle[radius=0.42bp] (329.333,489.745) circle[radius=0.42bp] (330.000,491.947) circle[radius=0.42bp] (330.667,493.859) circle[radius=0.42bp] (331.333,495.412) circle[radius=0.42bp] (332.000,495.418) circle[radius=0.42bp] (332.667,503.660) circle[radius=0.42bp] (333.333,492.450) circle[radius=0.42bp] (334.000,496.332) circle[radius=0.42bp] (334.667,495.751) circle[radius=0.42bp];
\draw[color=tg000000,line width=0.7bp,line cap=round] (319.333,496.178) -- (320.000,496.178) -- (320.667,496.178) -- (321.333,496.178) -- (322.000,496.178) -- (322.667,496.178) -- (323.333,496.178) -- (324.000,496.178) -- (324.667,496.178) -- (325.333,496.178) -- (326.000,496.178) -- (326.667,496.178) -- (327.333,496.178) -- (328.000,496.178) -- (328.667,496.178) -- (329.333,496.178) -- (330.000,496.178) -- (330.667,496.178) -- (331.333,496.178) -- (332.000,496.178) -- (332.667,496.178) -- (333.333,496.178) -- (334.000,496.178) -- (334.667,496.178);
\draw[color=tg0000ff,line width=0.3bp,line cap=round] (318.783,493.138) -- (319.883,494.238);
\draw[color=tg0000ff,line width=0.3bp,line cap=round] (318.783,494.238) -- (319.883,493.138);
\draw[color=tg0000ff,line width=0.3bp,line cap=round] (319.450,495.824) -- (320.550,496.924);
\draw[color=tg0000ff,line width=0.3bp,line cap=round] (319.450,496.924) -- (320.550,495.824);
\draw[color=tg0000ff,line width=0.3bp,line cap=round] (320.117,492.534) -- (321.217,493.634);
\draw[color=tg0000ff,line width=0.3bp,line cap=round] (320.117,493.634) -- (321.217,492.534);
\draw[color=tg0000ff,line width=0.3bp,line cap=round] (320.783,494.431) -- (321.883,495.531);
\draw[color=tg0000ff,line width=0.3bp,line cap=round] (320.783,495.531) -- (321.883,494.431);
\draw[color=tg0000ff,line width=0.3bp,line cap=round] (321.450,495.220) -- (322.550,496.320);
\draw[color=tg0000ff,line width=0.3bp,line cap=round] (321.450,496.320) -- (322.550,495.220);
\draw[color=tg0000ff,line width=0.3bp,line cap=round] (322.117,492.894) -- (323.217,493.994);
\draw[color=tg0000ff,line width=0.3bp,line cap=round] (322.117,493.994) -- (323.217,492.894);
\draw[color=tg0000ff,line width=0.3bp,line cap=round] (322.783,496.487) -- (323.883,497.587);
\draw[color=tg0000ff,line width=0.3bp,line cap=round] (322.783,497.587) -- (323.883,496.487);
\draw[color=tg0000ff,line width=0.3bp,line cap=round] (323.450,491.954) -- (324.550,493.054);
\draw[color=tg0000ff,line width=0.3bp,line cap=round] (323.450,493.054) -- (324.550,491.954);
\draw[color=tg0000ff,line width=0.3bp,line cap=round] (324.117,495.062) -- (325.217,496.162);
\draw[color=tg0000ff,line width=0.3bp,line cap=round] (324.117,496.162) -- (325.217,495.062);
\draw[color=tg0000ff,line width=0.3bp,line cap=round] (324.783,493.557) -- (325.883,494.657);
\draw[color=tg0000ff,line width=0.3bp,line cap=round] (324.783,494.657) -- (325.883,493.557);
\draw[color=tg0000ff,line width=0.3bp,line cap=round] (325.450,493.932) -- (326.550,495.032);
\draw[color=tg0000ff,line width=0.3bp,line cap=round] (325.450,495.032) -- (326.550,493.932);
\draw[color=tg0000ff,line width=0.3bp,line cap=round] (326.117,495.883) -- (327.217,496.983);
\draw[color=tg0000ff,line width=0.3bp,line cap=round] (326.117,496.983) -- (327.217,495.883);
\draw[color=tg0000ff,line width=0.3bp,line cap=round] (326.783,492.888) -- (327.883,493.988);
\draw[color=tg0000ff,line width=0.3bp,line cap=round] (326.783,493.988) -- (327.883,492.888);
\draw[color=tg0000ff,line width=0.3bp,line cap=round] (327.450,495.807) -- (328.550,496.907);
\draw[color=tg0000ff,line width=0.3bp,line cap=round] (327.450,496.907) -- (328.550,495.807);
\draw[color=tg0000ff,line width=0.3bp,line cap=round] (328.117,493.646) -- (329.217,494.746);
\draw[color=tg0000ff,line width=0.3bp,line cap=round] (328.117,494.746) -- (329.217,493.646);
\draw[color=tg0000ff,line width=0.3bp,line cap=round] (328.783,494.607) -- (329.883,495.707);
\draw[color=tg0000ff,line width=0.3bp,line cap=round] (328.783,495.707) -- (329.883,494.607);
\draw[color=tg0000ff,line width=0.3bp,line cap=round] (329.450,494.732) -- (330.550,495.832);
\draw[color=tg0000ff,line width=0.3bp,line cap=round] (329.450,495.832) -- (330.550,494.732);
\draw[color=tg0000ff,line width=0.3bp,line cap=round] (330.117,493.335) -- (331.217,494.435);
\draw[color=tg0000ff,line width=0.3bp,line cap=round] (330.117,494.435) -- (331.217,493.335);
\draw[color=tg0000ff,line width=0.3bp,line cap=round] (330.783,495.664) -- (331.883,496.764);
\draw[color=tg0000ff,line width=0.3bp,line cap=round] (330.783,496.764) -- (331.883,495.664);
\draw[color=tg0000ff,line width=0.3bp,line cap=round] (331.450,492.913) -- (332.550,494.013);
\draw[color=tg0000ff,line width=0.3bp,line cap=round] (331.450,494.013) -- (332.550,492.913);
\draw[color=tg0000ff,line width=0.3bp,line cap=round] (332.117,494.814) -- (333.217,495.914);
\draw[color=tg0000ff,line width=0.3bp,line cap=round] (332.117,495.914) -- (333.217,494.814);
\draw[color=tg0000ff,line width=0.3bp,line cap=round] (332.783,493.925) -- (333.883,495.025);
\draw[color=tg0000ff,line width=0.3bp,line cap=round] (332.783,495.025) -- (333.883,493.925);
\draw[color=tg0000ff,line width=0.3bp,line cap=round] (333.450,493.985) -- (334.550,495.085);
\draw[color=tg0000ff,line width=0.3bp,line cap=round] (333.450,495.085) -- (334.550,493.985);
\draw[color=tg0000ff,line width=0.3bp,line cap=round] (334.117,494.692) -- (335.217,495.792);
\draw[color=tg0000ff,line width=0.3bp,line cap=round] (334.117,495.792) -- (335.217,494.692);
\node[anchor=west,inner sep=0pt,outer sep=0pt,text=tg18314b,font={\sffamily\fontsize{6.2bp}{7.44bp}\selectfont}] at (346.000,475.828) {Low: Misreads the current level;};
\node[anchor=west,inner sep=0pt,outer sep=0pt,text=tg18314b,font={\sffamily\fontsize{6.2bp}{7.44bp}\selectfont}] at (352.000,483.268) {Invents a trend};
\draw[color=tgeceef0,line width=0.4bp,line cap=round] (0.000,517.494) -- (448.000,517.494);
\node[anchor=west,inner sep=0pt,outer sep=0pt,text=tg18314b,font={\sffamily\fontsize{7.2bp}{8.64bp}\selectfont\bfseries}] at (0.000,524.314) {Random walk with drift};
\node[anchor=west,inner sep=0pt,outer sep=0pt,text=tg18314b,font={\sffamily\fontsize{6.2bp}{7.44bp}\selectfont}] at (0.000,532.754) {\ensuremath{\checkmark} The current level plus the steady drift.};
\node[anchor=west,inner sep=0pt,outer sep=0pt,text=tg18314b,font={\sffamily\fontsize{6.2bp}{7.44bp}\selectfont}] at (0.000,540.594) {\ensuremath{\times} The wandering around the drift line, which grows};
\node[anchor=west,inner sep=0pt,outer sep=0pt,text=tg18314b,font={\sffamily\fontsize{6.2bp}{7.44bp}\selectfont}] at (0.000,548.034) {with the horizon.};
\path[fill=tgfab219,draw=none,line width=0.5bp] (168.000,519.994) rectangle (222.719,528.594);
\node[anchor=center,inner sep=0pt,outer sep=0pt,text=tg0b0b0b,font={\sffamily\fontsize{6.6bp}{7.92bp}\selectfont\bfseries}] at (195.359,524.294) {41  Struggles};
\draw[color=tgc9cdd1,line width=0.45bp,line cap=round] (168.000,560.680) -- (248.000,560.680);
\draw[color=tg9da7b0,line width=0.5bp,line cap=round,dash pattern=on 2bp off 2bp] (232.000,530.994) -- (232.000,560.680);
\path[fill=tg808080,draw=none] (168.333,555.883) circle[radius=0.42bp] (169.000,556.180) circle[radius=0.42bp] (169.667,556.449) circle[radius=0.42bp] (170.333,557.634) circle[radius=0.42bp] (171.000,557.218) circle[radius=0.42bp] (171.667,557.303) circle[radius=0.42bp] (172.333,557.690) circle[radius=0.42bp] (173.000,558.217) circle[radius=0.42bp] (173.667,559.090) circle[radius=0.42bp] (174.333,556.881) circle[radius=0.42bp] (175.000,558.617) circle[radius=0.42bp] (175.667,557.084) circle[radius=0.42bp] (176.333,557.369) circle[radius=0.42bp] (177.000,557.892) circle[radius=0.42bp] (177.667,556.094) circle[radius=0.42bp] (178.333,555.495) circle[radius=0.42bp] (179.000,556.377) circle[radius=0.42bp] (179.667,554.559) circle[radius=0.42bp] (180.333,554.184) circle[radius=0.42bp] (181.000,552.246) circle[radius=0.42bp] (181.667,552.636) circle[radius=0.42bp] (182.333,553.276) circle[radius=0.42bp] (183.000,553.148) circle[radius=0.42bp] (183.667,551.243) circle[radius=0.42bp] (184.333,548.249) circle[radius=0.42bp] (185.000,547.992) circle[radius=0.42bp] (185.667,547.771) circle[radius=0.42bp] (186.333,547.146) circle[radius=0.42bp] (187.000,549.005) circle[radius=0.42bp] (187.667,547.947) circle[radius=0.42bp] (188.333,548.611) circle[radius=0.42bp] (189.000,547.485) circle[radius=0.42bp] (189.667,547.907) circle[radius=0.42bp] (190.333,547.263) circle[radius=0.42bp] (191.000,548.608) circle[radius=0.42bp] (191.667,547.333) circle[radius=0.42bp] (192.333,547.442) circle[radius=0.42bp] (193.000,548.976) circle[radius=0.42bp] (193.667,547.219) circle[radius=0.42bp] (194.333,548.412) circle[radius=0.42bp] (195.000,547.493) circle[radius=0.42bp] (195.667,546.273) circle[radius=0.42bp] (196.333,545.905) circle[radius=0.42bp] (197.000,543.388) circle[radius=0.42bp] (197.667,544.930) circle[radius=0.42bp] (198.333,545.363) circle[radius=0.42bp] (199.000,545.720) circle[radius=0.42bp] (199.667,544.600) circle[radius=0.42bp] (200.333,544.170) circle[radius=0.42bp] (201.000,544.832) circle[radius=0.42bp] (201.667,543.664) circle[radius=0.42bp] (202.333,541.748) circle[radius=0.42bp] (203.000,544.203) circle[radius=0.42bp] (203.667,544.750) circle[radius=0.42bp] (204.333,543.649) circle[radius=0.42bp] (205.000,541.246) circle[radius=0.42bp] (205.667,541.784) circle[radius=0.42bp] (206.333,540.812) circle[radius=0.42bp] (207.000,541.176) circle[radius=0.42bp] (207.667,542.570) circle[radius=0.42bp] (208.333,543.521) circle[radius=0.42bp] (209.000,543.319) circle[radius=0.42bp] (209.667,542.579) circle[radius=0.42bp] (210.333,542.088) circle[radius=0.42bp] (211.000,540.642) circle[radius=0.42bp] (211.667,539.092) circle[radius=0.42bp] (212.333,541.091) circle[radius=0.42bp] (213.000,539.831) circle[radius=0.42bp] (213.667,540.390) circle[radius=0.42bp] (214.333,540.206) circle[radius=0.42bp] (215.000,541.359) circle[radius=0.42bp] (215.667,539.183) circle[radius=0.42bp] (216.333,537.276) circle[radius=0.42bp] (217.000,538.657) circle[radius=0.42bp] (217.667,538.617) circle[radius=0.42bp] (218.333,540.350) circle[radius=0.42bp] (219.000,538.340) circle[radius=0.42bp] (219.667,538.593) circle[radius=0.42bp] (220.333,538.964) circle[radius=0.42bp] (221.000,537.561) circle[radius=0.42bp] (221.667,538.790) circle[radius=0.42bp] (222.333,540.136) circle[radius=0.42bp] (223.000,539.516) circle[radius=0.42bp] (223.667,539.121) circle[radius=0.42bp] (224.333,539.765) circle[radius=0.42bp] (225.000,540.719) circle[radius=0.42bp] (225.667,539.804) circle[radius=0.42bp] (226.333,539.560) circle[radius=0.42bp] (227.000,540.511) circle[radius=0.42bp] (227.667,539.671) circle[radius=0.42bp] (228.333,539.739) circle[radius=0.42bp] (229.000,537.956) circle[radius=0.42bp] (229.667,538.929) circle[radius=0.42bp] (230.333,538.128) circle[radius=0.42bp] (231.000,537.834) circle[radius=0.42bp] (231.667,536.415) circle[radius=0.42bp] (232.333,535.510) circle[radius=0.42bp] (233.000,534.272) circle[radius=0.42bp] (233.667,535.023) circle[radius=0.42bp] (234.333,535.570) circle[radius=0.42bp] (235.000,536.363) circle[radius=0.42bp] (235.667,536.350) circle[radius=0.42bp] (236.333,534.743) circle[radius=0.42bp] (237.000,535.598) circle[radius=0.42bp] (237.667,537.090) circle[radius=0.42bp] (238.333,538.034) circle[radius=0.42bp] (239.000,538.362) circle[radius=0.42bp] (239.667,537.393) circle[radius=0.42bp] (240.333,536.974) circle[radius=0.42bp] (241.000,536.468) circle[radius=0.42bp] (241.667,536.711) circle[radius=0.42bp] (242.333,534.020) circle[radius=0.42bp] (243.000,533.590) circle[radius=0.42bp] (243.667,533.233) circle[radius=0.42bp] (244.333,534.603) circle[radius=0.42bp] (245.000,533.840) circle[radius=0.42bp] (245.667,535.112) circle[radius=0.42bp] (246.333,532.609) circle[radius=0.42bp] (247.000,533.027) circle[radius=0.42bp] (247.667,532.584) circle[radius=0.42bp];
\draw[color=tg000000,line width=0.7bp,line cap=round] (232.333,536.818) -- (233.000,536.682) -- (233.667,536.546) -- (234.333,536.410) -- (235.000,536.273) -- (235.667,536.137) -- (236.333,536.001) -- (237.000,535.865) -- (237.667,535.729) -- (238.333,535.592) -- (239.000,535.456) -- (239.667,535.320) -- (240.333,535.184) -- (241.000,535.048) -- (241.667,534.911) -- (242.333,534.775) -- (243.000,534.639) -- (243.667,534.503) -- (244.333,534.367) -- (245.000,534.230) -- (245.667,534.094) -- (246.333,533.958) -- (247.000,533.822) -- (247.667,533.686);
\draw[color=tg0000ff,line width=0.3bp,line cap=round] (231.783,536.233) -- (232.883,537.333);
\draw[color=tg0000ff,line width=0.3bp,line cap=round] (231.783,537.333) -- (232.883,536.233);
\draw[color=tg0000ff,line width=0.3bp,line cap=round] (232.450,536.219) -- (233.550,537.319);
\draw[color=tg0000ff,line width=0.3bp,line cap=round] (232.450,537.319) -- (233.550,536.219);
\draw[color=tg0000ff,line width=0.3bp,line cap=round] (233.117,535.855) -- (234.217,536.955);
\draw[color=tg0000ff,line width=0.3bp,line cap=round] (233.117,536.955) -- (234.217,535.855);
\draw[color=tg0000ff,line width=0.3bp,line cap=round] (233.783,536.013) -- (234.883,537.113);
\draw[color=tg0000ff,line width=0.3bp,line cap=round] (233.783,537.113) -- (234.883,536.013);
\draw[color=tg0000ff,line width=0.3bp,line cap=round] (234.450,535.859) -- (235.550,536.959);
\draw[color=tg0000ff,line width=0.3bp,line cap=round] (234.450,536.959) -- (235.550,535.859);
\draw[color=tg0000ff,line width=0.3bp,line cap=round] (235.117,535.544) -- (236.217,536.644);
\draw[color=tg0000ff,line width=0.3bp,line cap=round] (235.117,536.644) -- (236.217,535.544);
\draw[color=tg0000ff,line width=0.3bp,line cap=round] (235.783,535.656) -- (236.883,536.756);
\draw[color=tg0000ff,line width=0.3bp,line cap=round] (235.783,536.756) -- (236.883,535.656);
\draw[color=tg0000ff,line width=0.3bp,line cap=round] (236.450,535.522) -- (237.550,536.622);
\draw[color=tg0000ff,line width=0.3bp,line cap=round] (236.450,536.622) -- (237.550,535.522);
\draw[color=tg0000ff,line width=0.3bp,line cap=round] (237.117,535.602) -- (238.217,536.702);
\draw[color=tg0000ff,line width=0.3bp,line cap=round] (237.117,536.702) -- (238.217,535.602);
\draw[color=tg0000ff,line width=0.3bp,line cap=round] (237.783,535.429) -- (238.883,536.529);
\draw[color=tg0000ff,line width=0.3bp,line cap=round] (237.783,536.529) -- (238.883,535.429);
\draw[color=tg0000ff,line width=0.3bp,line cap=round] (238.450,535.509) -- (239.550,536.609);
\draw[color=tg0000ff,line width=0.3bp,line cap=round] (238.450,536.609) -- (239.550,535.509);
\draw[color=tg0000ff,line width=0.3bp,line cap=round] (239.117,535.493) -- (240.217,536.593);
\draw[color=tg0000ff,line width=0.3bp,line cap=round] (239.117,536.593) -- (240.217,535.493);
\draw[color=tg0000ff,line width=0.3bp,line cap=round] (239.783,535.149) -- (240.883,536.249);
\draw[color=tg0000ff,line width=0.3bp,line cap=round] (239.783,536.249) -- (240.883,535.149);
\draw[color=tg0000ff,line width=0.3bp,line cap=round] (240.450,535.029) -- (241.550,536.129);
\draw[color=tg0000ff,line width=0.3bp,line cap=round] (240.450,536.129) -- (241.550,535.029);
\draw[color=tg0000ff,line width=0.3bp,line cap=round] (241.117,534.850) -- (242.217,535.950);
\draw[color=tg0000ff,line width=0.3bp,line cap=round] (241.117,535.950) -- (242.217,534.850);
\draw[color=tg0000ff,line width=0.3bp,line cap=round] (241.783,534.887) -- (242.883,535.987);
\draw[color=tg0000ff,line width=0.3bp,line cap=round] (241.783,535.987) -- (242.883,534.887);
\draw[color=tg0000ff,line width=0.3bp,line cap=round] (242.450,534.650) -- (243.550,535.750);
\draw[color=tg0000ff,line width=0.3bp,line cap=round] (242.450,535.750) -- (243.550,534.650);
\draw[color=tg0000ff,line width=0.3bp,line cap=round] (243.117,534.473) -- (244.217,535.573);
\draw[color=tg0000ff,line width=0.3bp,line cap=round] (243.117,535.573) -- (244.217,534.473);
\draw[color=tg0000ff,line width=0.3bp,line cap=round] (243.783,534.476) -- (244.883,535.576);
\draw[color=tg0000ff,line width=0.3bp,line cap=round] (243.783,535.576) -- (244.883,534.476);
\draw[color=tg0000ff,line width=0.3bp,line cap=round] (244.450,534.319) -- (245.550,535.419);
\draw[color=tg0000ff,line width=0.3bp,line cap=round] (244.450,535.419) -- (245.550,534.319);
\draw[color=tg0000ff,line width=0.3bp,line cap=round] (245.117,534.224) -- (246.217,535.324);
\draw[color=tg0000ff,line width=0.3bp,line cap=round] (245.117,535.324) -- (246.217,534.224);
\draw[color=tg0000ff,line width=0.3bp,line cap=round] (245.783,534.011) -- (246.883,535.111);
\draw[color=tg0000ff,line width=0.3bp,line cap=round] (245.783,535.111) -- (246.883,534.011);
\draw[color=tg0000ff,line width=0.3bp,line cap=round] (246.450,533.977) -- (247.550,535.077);
\draw[color=tg0000ff,line width=0.3bp,line cap=round] (246.450,535.077) -- (247.550,533.977);
\draw[color=tg0000ff,line width=0.3bp,line cap=round] (247.117,533.833) -- (248.217,534.933);
\draw[color=tg0000ff,line width=0.3bp,line cap=round] (247.117,534.933) -- (248.217,533.833);
\path[fill=tgd03b3b,draw=none,line width=0.5bp] (255.000,519.994) rectangle (290.797,528.594);
\node[anchor=center,inner sep=0pt,outer sep=0pt,text=tgffffff,font={\sffamily\fontsize{6.6bp}{7.92bp}\selectfont\bfseries}] at (272.898,524.294) {23  Fails};
\draw[color=tgc9cdd1,line width=0.45bp,line cap=round] (255.000,560.680) -- (335.000,560.680);
\draw[color=tg9da7b0,line width=0.5bp,line cap=round,dash pattern=on 2bp off 2bp] (319.000,530.994) -- (319.000,560.680);
\path[fill=tg808080,draw=none] (255.333,551.126) circle[radius=0.42bp] (256.000,552.543) circle[radius=0.42bp] (256.667,552.513) circle[radius=0.42bp] (257.333,555.102) circle[radius=0.42bp] (258.000,553.366) circle[radius=0.42bp] (258.667,554.317) circle[radius=0.42bp] (259.333,554.231) circle[radius=0.42bp] (260.000,555.888) circle[radius=0.42bp] (260.667,558.718) circle[radius=0.42bp] (261.333,552.917) circle[radius=0.42bp] (262.000,559.090) circle[radius=0.42bp] (262.667,555.132) circle[radius=0.42bp] (263.333,555.301) circle[radius=0.42bp] (264.000,557.072) circle[radius=0.42bp] (264.667,552.967) circle[radius=0.42bp] (265.333,553.809) circle[radius=0.42bp] (266.000,554.891) circle[radius=0.42bp] (266.667,551.902) circle[radius=0.42bp] (267.333,553.218) circle[radius=0.42bp] (268.000,548.663) circle[radius=0.42bp] (268.667,550.651) circle[radius=0.42bp] (269.333,552.495) circle[radius=0.42bp] (270.000,554.772) circle[radius=0.42bp] (270.667,552.233) circle[radius=0.42bp] (271.333,545.986) circle[radius=0.42bp] (272.000,547.623) circle[radius=0.42bp] (272.667,547.824) circle[radius=0.42bp] (273.333,544.211) circle[radius=0.42bp] (274.000,550.419) circle[radius=0.42bp] (274.667,548.453) circle[radius=0.42bp] (275.333,548.796) circle[radius=0.42bp] (276.000,546.190) circle[radius=0.42bp] (276.667,548.548) circle[radius=0.42bp] (277.333,547.956) circle[radius=0.42bp] (278.000,550.627) circle[radius=0.42bp] (278.667,545.838) circle[radius=0.42bp] (279.333,545.800) circle[radius=0.42bp] (280.000,549.091) circle[radius=0.42bp] (280.667,546.664) circle[radius=0.42bp] (281.333,549.966) circle[radius=0.42bp] (282.000,547.909) circle[radius=0.42bp] (282.667,543.949) circle[radius=0.42bp] (283.333,545.668) circle[radius=0.42bp] (284.000,540.154) circle[radius=0.42bp] (284.667,546.389) circle[radius=0.42bp] (285.333,544.519) circle[radius=0.42bp] (286.000,546.057) circle[radius=0.42bp] (286.667,543.938) circle[radius=0.42bp] (287.333,544.525) circle[radius=0.42bp] (288.000,545.946) circle[radius=0.42bp] (288.667,545.516) circle[radius=0.42bp] (289.333,537.632) circle[radius=0.42bp] (290.000,544.664) circle[radius=0.42bp] (290.667,548.237) circle[radius=0.42bp] (291.333,545.150) circle[radius=0.42bp] (292.000,541.099) circle[radius=0.42bp] (292.667,543.432) circle[radius=0.42bp] (293.333,539.651) circle[radius=0.42bp] (294.000,540.519) circle[radius=0.42bp] (294.667,543.492) circle[radius=0.42bp] (295.333,544.763) circle[radius=0.42bp] (296.000,545.465) circle[radius=0.42bp] (296.667,544.386) circle[radius=0.42bp] (297.333,545.248) circle[radius=0.42bp] (298.000,539.769) circle[radius=0.42bp] (298.667,537.916) circle[radius=0.42bp] (299.333,544.107) circle[radius=0.42bp] (300.000,542.989) circle[radius=0.42bp] (300.667,544.381) circle[radius=0.42bp] (301.333,542.933) circle[radius=0.42bp] (302.000,547.314) circle[radius=0.42bp] (302.667,542.411) circle[radius=0.42bp] (303.333,536.952) circle[radius=0.42bp] (304.000,540.705) circle[radius=0.42bp] (304.667,538.731) circle[radius=0.42bp] (305.333,542.278) circle[radius=0.42bp] (306.000,537.044) circle[radius=0.42bp] (306.667,539.772) circle[radius=0.42bp] (307.333,541.007) circle[radius=0.42bp] (308.000,535.146) circle[radius=0.42bp] (308.667,540.776) circle[radius=0.42bp] (309.333,543.294) circle[radius=0.42bp] (310.000,540.500) circle[radius=0.42bp] (310.667,540.576) circle[radius=0.42bp] (311.333,539.297) circle[radius=0.42bp] (312.000,543.271) circle[radius=0.42bp] (312.667,539.751) circle[radius=0.42bp] (313.333,539.481) circle[radius=0.42bp] (314.000,542.695) circle[radius=0.42bp] (314.667,540.685) circle[radius=0.42bp] (315.333,542.348) circle[radius=0.42bp] (316.000,535.794) circle[radius=0.42bp] (316.667,541.501) circle[radius=0.42bp] (317.333,540.460) circle[radius=0.42bp] (318.000,541.219) circle[radius=0.42bp] (318.667,538.929) circle[radius=0.42bp] (319.333,537.662) circle[radius=0.42bp] (320.000,535.429) circle[radius=0.42bp] (320.667,538.255) circle[radius=0.42bp] (321.333,538.000) circle[radius=0.42bp] (322.000,540.621) circle[radius=0.42bp] (322.667,539.657) circle[radius=0.42bp] (323.333,532.584) circle[radius=0.42bp] (324.000,536.037) circle[radius=0.42bp] (324.667,540.169) circle[radius=0.42bp] (325.333,541.817) circle[radius=0.42bp] (326.000,541.491) circle[radius=0.42bp] (326.667,539.179) circle[radius=0.42bp] (327.333,538.756) circle[radius=0.42bp] (328.000,539.616) circle[radius=0.42bp] (328.667,540.686) circle[radius=0.42bp] (329.333,534.170) circle[radius=0.42bp] (330.000,535.400) circle[radius=0.42bp] (330.667,536.455) circle[radius=0.42bp] (331.333,537.291) circle[radius=0.42bp] (332.000,537.187) circle[radius=0.42bp] (332.667,542.089) circle[radius=0.42bp] (333.333,535.167) circle[radius=0.42bp] (334.000,537.418) circle[radius=0.42bp] (334.667,536.958) circle[radius=0.42bp];
\draw[color=tg000000,line width=0.7bp,line cap=round] (319.333,539.699) -- (320.000,539.591) -- (320.667,539.483) -- (321.333,539.375) -- (322.000,539.267) -- (322.667,539.159) -- (323.333,539.051) -- (324.000,538.943) -- (324.667,538.836) -- (325.333,538.728) -- (326.000,538.620) -- (326.667,538.512) -- (327.333,538.404) -- (328.000,538.296) -- (328.667,538.188) -- (329.333,538.080) -- (330.000,537.972) -- (330.667,537.864) -- (331.333,537.757) -- (332.000,537.649) -- (332.667,537.541) -- (333.333,537.433) -- (334.000,537.325) -- (334.667,537.217);
\draw[color=tg0000ff,line width=0.3bp,line cap=round] (318.783,537.758) -- (319.883,538.858);
\draw[color=tg0000ff,line width=0.3bp,line cap=round] (318.783,538.858) -- (319.883,537.758);
\draw[color=tg0000ff,line width=0.3bp,line cap=round] (319.450,539.240) -- (320.550,540.340);
\draw[color=tg0000ff,line width=0.3bp,line cap=round] (319.450,540.340) -- (320.550,539.240);
\draw[color=tg0000ff,line width=0.3bp,line cap=round] (320.117,537.213) -- (321.217,538.313);
\draw[color=tg0000ff,line width=0.3bp,line cap=round] (320.117,538.313) -- (321.217,537.213);
\draw[color=tg0000ff,line width=0.3bp,line cap=round] (320.783,538.271) -- (321.883,539.371);
\draw[color=tg0000ff,line width=0.3bp,line cap=round] (320.783,539.371) -- (321.883,538.271);
\draw[color=tg0000ff,line width=0.3bp,line cap=round] (321.450,538.620) -- (322.550,539.720);
\draw[color=tg0000ff,line width=0.3bp,line cap=round] (321.450,539.720) -- (322.550,538.620);
\draw[color=tg0000ff,line width=0.3bp,line cap=round] (322.117,537.204) -- (323.217,538.304);
\draw[color=tg0000ff,line width=0.3bp,line cap=round] (322.117,538.304) -- (323.217,537.204);
\draw[color=tg0000ff,line width=0.3bp,line cap=round] (322.783,539.212) -- (323.883,540.312);
\draw[color=tg0000ff,line width=0.3bp,line cap=round] (322.783,540.312) -- (323.883,539.212);
\draw[color=tg0000ff,line width=0.3bp,line cap=round] (323.450,536.522) -- (324.550,537.622);
\draw[color=tg0000ff,line width=0.3bp,line cap=round] (323.450,537.622) -- (324.550,536.522);
\draw[color=tg0000ff,line width=0.3bp,line cap=round] (324.117,538.271) -- (325.217,539.371);
\draw[color=tg0000ff,line width=0.3bp,line cap=round] (324.117,539.371) -- (325.217,538.271);
\draw[color=tg0000ff,line width=0.3bp,line cap=round] (324.783,537.306) -- (325.883,538.406);
\draw[color=tg0000ff,line width=0.3bp,line cap=round] (324.783,538.406) -- (325.883,537.306);
\draw[color=tg0000ff,line width=0.3bp,line cap=round] (325.450,537.439) -- (326.550,538.539);
\draw[color=tg0000ff,line width=0.3bp,line cap=round] (325.450,538.539) -- (326.550,537.439);
\draw[color=tg0000ff,line width=0.3bp,line cap=round] (326.117,538.468) -- (327.217,539.568);
\draw[color=tg0000ff,line width=0.3bp,line cap=round] (326.117,539.568) -- (327.217,538.468);
\draw[color=tg0000ff,line width=0.3bp,line cap=round] (326.783,536.658) -- (327.883,537.758);
\draw[color=tg0000ff,line width=0.3bp,line cap=round] (326.783,537.758) -- (327.883,536.658);
\draw[color=tg0000ff,line width=0.3bp,line cap=round] (327.450,538.263) -- (328.550,539.363);
\draw[color=tg0000ff,line width=0.3bp,line cap=round] (327.450,539.363) -- (328.550,538.263);
\draw[color=tg0000ff,line width=0.3bp,line cap=round] (328.117,536.921) -- (329.217,538.021);
\draw[color=tg0000ff,line width=0.3bp,line cap=round] (328.117,538.021) -- (329.217,536.921);
\draw[color=tg0000ff,line width=0.3bp,line cap=round] (328.783,537.414) -- (329.883,538.514);
\draw[color=tg0000ff,line width=0.3bp,line cap=round] (328.783,538.514) -- (329.883,537.414);
\draw[color=tg0000ff,line width=0.3bp,line cap=round] (329.450,537.382) -- (330.550,538.482);
\draw[color=tg0000ff,line width=0.3bp,line cap=round] (329.450,538.482) -- (330.550,537.382);
\draw[color=tg0000ff,line width=0.3bp,line cap=round] (330.117,536.495) -- (331.217,537.595);
\draw[color=tg0000ff,line width=0.3bp,line cap=round] (330.117,537.595) -- (331.217,536.495);
\draw[color=tg0000ff,line width=0.3bp,line cap=round] (330.783,537.730) -- (331.883,538.830);
\draw[color=tg0000ff,line width=0.3bp,line cap=round] (330.783,538.830) -- (331.883,537.730);
\draw[color=tg0000ff,line width=0.3bp,line cap=round] (331.450,536.065) -- (332.550,537.165);
\draw[color=tg0000ff,line width=0.3bp,line cap=round] (331.450,537.165) -- (332.550,536.065);
\draw[color=tg0000ff,line width=0.3bp,line cap=round] (332.117,537.088) -- (333.217,538.188);
\draw[color=tg0000ff,line width=0.3bp,line cap=round] (332.117,538.188) -- (333.217,537.088);
\draw[color=tg0000ff,line width=0.3bp,line cap=round] (332.783,536.485) -- (333.883,537.585);
\draw[color=tg0000ff,line width=0.3bp,line cap=round] (332.783,537.585) -- (333.883,536.485);
\draw[color=tg0000ff,line width=0.3bp,line cap=round] (333.450,536.462) -- (334.550,537.562);
\draw[color=tg0000ff,line width=0.3bp,line cap=round] (333.450,537.562) -- (334.550,536.462);
\draw[color=tg0000ff,line width=0.3bp,line cap=round] (334.117,536.744) -- (335.217,537.844);
\draw[color=tg0000ff,line width=0.3bp,line cap=round] (334.117,537.844) -- (335.217,536.744);
\node[anchor=west,inner sep=0pt,outer sep=0pt,text=tg18314b,font={\sffamily\fontsize{6.2bp}{7.44bp}\selectfont}] at (346.000,523.714) {Low: Forecasts too low;};
\node[anchor=west,inner sep=0pt,outer sep=0pt,text=tg18314b,font={\sffamily\fontsize{6.2bp}{7.44bp}\selectfont}] at (352.000,531.154) {Misjudges the trend};
\node[anchor=west,inner sep=0pt,outer sep=0pt,text=tg18314b,font={\sffamily\fontsize{6.2bp}{7.44bp}\selectfont}] at (346.000,538.594) {High: Misreads the current level;};
\node[anchor=west,inner sep=0pt,outer sep=0pt,text=tg18314b,font={\sffamily\fontsize{6.2bp}{7.44bp}\selectfont}] at (352.000,546.034) {Misjudges the trend};
\draw[color=tgc9cdd1,line width=0.45bp,line cap=round] (0.000,565.180) -- (448.000,565.180);
\node[anchor=west,inner sep=0pt,outer sep=0pt,text=tg626b73,font={\sffamily\fontsize{6.2bp}{7.44bp}\selectfont}] at (0.000,570.400) {Score: 100 = as close as a forecaster that knows the true process; 50 = its avoidable error equals the unavoidable noise. Bands: 90+ nails it,};
\node[anchor=west,inner sep=0pt,outer sep=0pt,text=tg626b73,font={\sffamily\fontsize{6.2bp}{7.44bp}\selectfont}] at (0.000,577.840) {70+ good, 40+ struggles, below 40 fails. Each panel has its own vertical scale.};
\node[anchor=west,inner sep=0pt,outer sep=0pt,text=tg626b73,font={\sffamily\fontsize{6.2bp}{7.44bp}\selectfont}] at (0.000,585.280) {Panels show instance 1:};
\path[fill=tg808080,draw=none] (78.047,585.280) circle[radius=0.6bp] (80.547,585.280) circle[radius=0.6bp] (83.047,585.280) circle[radius=0.6bp];
\node[anchor=west,inner sep=0pt,outer sep=0pt,text=tg626b73,font={\sffamily\fontsize{6.2bp}{7.44bp}\selectfont}] at (87.047,585.280) {observations};
\draw[color=tg000000,line width=0.8bp,line cap=round] (131.844,585.280) -- (138.844,585.280);
\node[anchor=west,inner sep=0pt,outer sep=0pt,text=tg626b73,font={\sffamily\fontsize{6.2bp}{7.44bp}\selectfont}] at (142.844,585.280) {answer key};
\draw[color=tg0000ff,line width=0.45bp,line cap=round] (183.900,584.180) -- (186.100,586.380);
\draw[color=tg0000ff,line width=0.45bp,line cap=round] (183.900,586.380) -- (186.100,584.180);
\draw[color=tg0000ff,line width=0.45bp,line cap=round] (187.900,584.180) -- (190.100,586.380);
\draw[color=tg0000ff,line width=0.45bp,line cap=round] (187.900,586.380) -- (190.100,584.180);
\node[anchor=west,inner sep=0pt,outer sep=0pt,text=tg626b73,font={\sffamily\fontsize{6.2bp}{7.44bp}\selectfont}] at (194.000,585.280) {forecast};
\draw[color=tg9da7b0,line width=0.6bp,line cap=round,dash pattern=on 2bp off 2bp] (229.016,582.080) -- (229.016,588.480);
\node[anchor=west,inner sep=0pt,outer sep=0pt,text=tg626b73,font={\sffamily\fontsize{6.2bp}{7.44bp}\selectfont}] at (235.016,585.280) {forecast origin.};
\end{tikzpicture}%
}
\endgroup

% END VECTOR FIGURE paper-scorecard
\caption{\textbf{What a forecaster can and cannot learn.} TimeGym's
scorecard for the quick-start Ridge pipeline (24 lags, ridge penalty 1,
standardized lag inputs). Each row is a kind of series, with what can and
cannot be predicted by design; the columns are low and high noise. A
cell's score is $100\times$ unavoidable$/$(unavoidable $+$ avoidable)
error over ten series; 100 matches the answer key. Panels show one series:
gray dots are observations, the black line is the answer key, blue crosses
are the forecast.}
\label{fig:scorecard}
\end{figure}

\section{Introduction}

Reinforcement learning has gyms: standard environments with known rules,
in which agents are trained and tested~\cite{brockman2016gym}. Forecasting
benchmarks mostly use real series, where nobody knows the rules and so
nobody knows the best possible forecast. TimeGym is a gym for forecasting
pipelines. Its scenarios generate series from known rules, so every test
case comes with an answer key: the best forecast anyone could make from
the history alone. A pipeline's error then splits into two parts, what it
missed and noise nobody could predict, and a scorecard turns that split
into one picture (Figure~\ref{fig:scorecard}).

Consider three failures. A forecast has the wrong number of future steps;
it fails an output check. A last-value predictor returns
the required shape but misses a deterministic slope; this can be an honest
capability limitation. A sample-based predictor retains the correct
distribution at every future coordinate but corrupts the pairing between
coordinates; point and marginal uncertainty scores can remain unchanged.
These failures require different evidence and different remedies.

\begin{table}[t]
\centering\small
\begin{tabular}{@{}>{\raggedright\arraybackslash}p{0.21\linewidth}>{\raggedright\arraybackslash}p{0.25\linewidth}>{\raggedright\arraybackslash}p{0.46\linewidth}@{}}
\toprule
Term used here & Formal name & What it means \\
\midrule
Output checks & Forecast contract, output contract & Right number of steps; the requested dates and series, in order; no missing or infinite values; quantiles in order \\[0.3em]
Answer key & Conditional reference & The best forecast anyone could make knowing the rules but not the future noise; for squared error, the conditional mean \\[0.3em]
Avoidable error & Mean-reference MSE & How far the forecast is from the answer key \\[0.3em]
Unavoidable error & $\operatorname{tr}(\Sigma_F)/K$ & The error the answer key itself makes, because the future noise has not happened yet \\[0.3em]
Shift check & Declared, behavioral or metamorphic relation & A what-if rule: add 10 to the history and the forecast should rise by 10 \\[0.3em]
Check & Comparator & Any of the five checks compared in Study 2 \\[0.3em]
Second implementation & Pseudo-oracle & The same algorithm written again, independently \\[0.3em]
Known only to the evaluator & Privileged information & True parameters, future change dates \\[0.3em]
What the forecast stands for & Estimand, functional & A mean, a median or a quantile \\[0.3em]
Can't learn it & Capability gap & The model cannot learn the pattern; not a bug \\[0.3em]
The numbers behind the verdict & Diagnostic evidence & The measured values a report stores with each verdict \\
\bottomrule
\end{tabular}
\caption{Plain terms used in this paper and the formal names they stand
for. The formal names appear where the mathematics or the cited
literature needs them.}
\label{tab:terms}
\end{table}

TimeGym keeps them apart. Its intended use is to test a forecasting
pipeline before deployment and to reproduce a failure after one. A
scenario defines a controlled data-generating process (DGP), what the
forecaster may see, and the answer key held by the evaluator. Output
checks define a well-formed forecast. A report records which check failed
or how much avoidable error was observed, together with the case
configuration and seed, so that the case can be replayed.

The paper makes three contributions.
\begin{itemize}
\item A specification of forecasting tests that separates three kinds of
evidence: output checks, agreement with a known answer key, and shift
checks that compare a task with a transformed copy of itself
(Sections~\ref{sec:references} and~\ref{sec:contracts}).
\item An open-source implementation whose scorecard and reports record
verdicts, the numbers behind them, and replay information. The
scorecard's ten kinds of series are built from five Gaussian families, an
AR(2) for irregular cycles, measurement noise on the autoregressive and
random-walk rows, and two composite rows (a trend with a cycle, and a
sudden level shift). Both studies use the five families only. Composed,
financial and joint scenarios extend the catalog
(Appendix~\ref{app:extended}).
\item Two controlled studies test the checks themselves
(Sections~\ref{sec:protocol} and~\ref{sec:results}). What they found:
\begin{itemize}
\item Bugs that keep the output well formed passed every output check.
Reversed lags and stale recursive state also passed both shift checks.
Only the independent implementation caught them in all six conditions;
the loss alarm caught them only on the seasonal one.
\item Swapping two target columns passed a common shift but failed
per-target shifts in every case.
\item Fixing the stale-state bug made the average error worse on a random
walk with drift ($+$\RepairDeltaMean{} MSE; only \RepairImproved{} of
\CasesPerCondition{} cases improved).
\item Shuffling sample paths left every per-step distribution, and so the
CRPS, exactly unchanged, while energy and variogram scores got worse.
\item Four covariance checks per suite gave at least one false alarm in
14 of 64 clean suites (22\%). Holm's correction cut that to 3 of 64 but
also cut detection of shrunken spread from 63 to 30 of 64.
\end{itemize}
\end{itemize}

TimeGym sits between three lines of work: synthetic benchmarks that map
what forecasting models can
learn~\cite{tan2025syntsbench,ghoroghchian2026simpletimebench,sun2026finstressts},
metamorphic testing of forecasting
software~\cite{dwarakanath2019metamorphic}, and interface conformance
tests such as sktime's~\cite{sktimeConformance}. It adds scores stated
against an answer key, separate checks for point, marginal and joint
forecasts, and repair rechecks. Section~\ref{sec:related} compares them
feature by feature.

Sections~\ref{sec:references} and~\ref{sec:contracts} define the tasks,
answer keys, output checks, and reports. Sections~\ref{sec:protocol}
and~\ref{sec:results} describe the two studies and their results.
Section~\ref{sec:related} places the work among prior testing and
forecasting research, and Section~\ref{sec:discussion} discusses what the
studies establish and what they do not. Table~\ref{tab:terms} lists the
plain terms used throughout and the formal names they stand for.

\paragraph{What changed since 2021.}
The 2021 version introduced synthetic challenges for debugging forecasting
pipelines~\cite{seca2021timegym}. It compared forecasts with the
noise-free part of the generator using sMAPE, itself an asymmetric
measure~\cite{goodwin1999asymmetry}, and called that part the optimal
forecast. Here the answer key depends on what the forecast stands for (a
mean, a median or a quantile), its loss, and what the forecaster may see,
so future noise is averaged out rather than treated as known. The 2021
version also argued that fresh synthetic series make multiple-testing
corrections unnecessary, and that a pattern learned on synthetic series
will be learned on real data. Fresh draws do not remove the need for
those corrections, and synthetic tests speak only about the mechanisms
they contain. Everything described here is new work.

\section{Forecasting tasks and answer keys}
\label{sec:references}

\subsection{What the forecaster sees and what it is asked for}

Let $Y_t\in\R^D$ denote the target vector and let the observed context be
$y_{0:L-1}$. The task requests $H$ future steps,
$Y_{L:L+H-1}$. The forecaster sees this context and
the forecast indices. In the implemented generators, the evaluator also
knows the fixed DGP parameters $\theta$. Its reference distribution is
\begin{equation}
 F_\theta = \mathcal{L}\!\left(Y_{L:L+H-1}
       \mid Y_{0:L-1}=y_{0:L-1},\theta\right).
\end{equation}
Future innovations are integrated out, not exposed to the forecaster.
The ordinary adapter receives neither the realized future nor the
answer key. Declared change schedules are also known only to the
evaluator. Controls that are deliberately given this information, to test
the evaluator itself, are identified separately.

For a conditional-mean forecast, the answer key at lead $h$ and target $d$ is
$m_{hd}=\E_\theta[Y_{L+h-1,d}\mid y_{0:L-1}]$, with $h=1,\ldots,H$.
Let $K=HD$, let $\Sigma_F$ be the covariance of the flattened future vector,
and let $a$ be a forecast based only on the allowed information and
independent model randomness. Conditional squared-error risk decomposes as
\begin{equation}
 \frac{1}{K}\E_\theta[\|Y-a\|_2^2\mid y_{0:L-1},a]
 = \frac{1}{K}\|m-a\|_2^2 + \frac{1}{K}\operatorname{tr}(\Sigma_F).
 \label{eq:risk}
\end{equation}
The first term is the avoidable error: how far the forecast is from the
answer key. The second is the unavoidable error, from noise no forecaster
can predict. The scorecard reports their ratio as a score,
\begin{equation}
 \mathrm{score}=100\,\frac{U}{U+A},
 \label{eq:score}
\end{equation}
where $A$ and $U$ are the avoidable and unavoidable errors averaged over
the cases of one cell: 100 matches the answer key, and 50 means the
avoidable error equals the unavoidable one. The avoidable error differs
from realized test-set MSE, which also contains the future noise. The
answer key is known only to the evaluator: a learned forecaster must infer
structure and parameters from a finite history, so a nonzero avoidable
error alone does not establish an implementation fault.

For level $q\in(0,1)$, the answer key is instead the marginal conditional
quantile $F_{\theta,hd}^{-1}(q)$, evaluated using pinball loss. A median
forecast corresponds to absolute error. Conditional means and medians
coincide in the five core Gaussian families. They generally differ in the
asymmetric price and jump distributions; what the forecast stands for must
remain explicit.
Marginal quantiles do not determine the joint distribution of trajectories.

\subsection{Kinds of series}

Figure~\ref{fig:scorecard} shows the kinds of series at work. Five
Gaussian families supply the answer keys used in both studies
(Table~\ref{tab:families}).
Let $\epsilon_t\stackrel{\mathrm{iid}}{\sim}\mathcal{N}(0,\Omega)$, where
$\Omega=\operatorname{diag}(\sigma)R\operatorname{diag}(\sigma)$ and $R$
is a valid target-correlation matrix. Innovations are independent across
time but can be correlated across targets. Vector operations in
Table~\ref{tab:families} are componentwise, and $\phi$ is scalar.

\begin{table}[!htb]
\centering\footnotesize
\setlength{\tabcolsep}{4pt}
\begin{tabular}{@{}>{\raggedright\arraybackslash}p{0.15\linewidth}>{\raggedright\arraybackslash}p{0.27\linewidth}>{\raggedright\arraybackslash}p{0.18\linewidth}>{\raggedright\arraybackslash}p{0.16\linewidth}>{\raggedright\arraybackslash}p{0.15\linewidth}@{}}
\toprule
Kind of series & Equation; answer key & Predictable & Unpredictable by design & Unavoidable error as the horizon grows \\
\midrule
Noise around a fixed level & $Y_t=\mu+\epsilon_t$; $m_h=\mu$ & The level & Every wiggle around it & Flat at $\sigma^2$ \\[0.5em]
Straight-line trend & $Y_t=b+ct+\epsilon_t$; $m_h=b+c(L+h-1)$ & The line, any number of steps ahead & The noise around the line & Flat at $\sigma^2$ \\[0.5em]
Repeating seasonal cycle & $Y_t=\mu+A\sin(2\pi t/p)+\epsilon_t$; the same curve at $t=L+h-1$ & The repeating shape, any number of steps ahead & The noise around the cycle & Flat at $\sigma^2$ \\[0.5em]
Mean reversion (AR(1)): shocks fade & $Y_t=\mu+\phi(Y_{t-1}-\mu)+\epsilon_t$, $|\phi|<1$; $m_h=\mu+\phi^h(y_{L-1}-\mu)$ & The pull back to the average: at $\phi=0.8$ a gap halves in about three steps & New shocks; after about ten steps only the average is left & Rises from $\sigma^2$ towards $\sigma^2/(1-\phi^2)$, $2.8\sigma^2$ at $\phi=0.8$ \\[0.5em]
Random walk: shocks never fade (optional drift) & $Y_t=Y_{t-1}+c+\epsilon_t$; $m_h=y_{L-1}+hc$ & Where it is now, plus the drift per step & The direction of the next move & Grows without limit: $h\sigma^2$ \\
\bottomrule
\end{tabular}
\caption{The five core Gaussian families under the scorecard's names:
their answer keys (the conditional mean $m_h$ at lead $h$), what can and
cannot be predicted by design, and the unavoidable error at lead $h$ for
a scalar innovation variance $\sigma^2$.}
\label{tab:families}
\end{table}

\paragraph{Shocks that vanish, fade, or stay.}
In the first three rows a shock at time $t$ does not carry into later
observations: the series fluctuates around a level, a line, or a
repeating curve. The last two rows are both autoregressive: each value is
a weighted sum of earlier values plus a new random shock. In mean
reversion a shock fades by the factor $\phi$ each step, so the series
keeps returning to its mean; for positive $\phi$ this is the
discrete-time Ornstein--Uhlenbeck process. A random walk is the same
equation with $\phi=1$: a shock never fades and there is no mean to
return to, so the best forecast is the current level plus the drift, and
its uncertainty grows without limit. Its centered increments
$Y_t-Y_{t-1}-c$, rather than its levels, are white noise. The parameter
$c$ is the optional \emph{drift}, an expected change per step; $c=0$ does
not prevent an individual path from wandering upward or downward. A
trend and a random walk with drift therefore differ in whether shocks
persist and how uncertainty grows, not in a visible slope.

The AR(1) context starts from the stationary law
$Y_0\sim\mathcal{N}(\mu,\Omega/(1-\phi^2))$; the random walk starts at
a specified deterministic intercept. For white noise, trend, and
seasonality, future covariance blocks are
$\Cov(Y_{L+h-1},Y_{L+j-1}\mid y)=\mathbf{1}_{h=j}\Omega$.
For AR(1), define the lower-triangular innovation map
$B_{hk}=\phi^{h-k}\mathbf{1}_{k\leq h}$, for $h,k=1,\ldots,H$.
The full covariance is $(BB^\top)\otimes\Omega$, using horizon-major
flattening. For a random walk, covariance block $(h,j)$ is
$\min(h,j)\Omega$. The conditional marginal quantiles are
\begin{equation}
 Q_{hd}(q)=m_{hd}+\sqrt{(\Sigma_F)_{(h,d),(h,d)}}\,\Phi^{-1}(q).
\end{equation}
Zero innovation variance gives a deterministic answer key.

\paragraph{Cycles and measurement noise.}
The scorecard adds two constructions to these families. The first is an
AR(2),
$Y_t=\mu+\phi_1(Y_{t-1}-\mu)+\phi_2(Y_{t-2}-\mu)+\epsilon_t$. With
$\phi_1^2+4\phi_2<0$ its characteristic roots are complex, so a shock
rings as a damped oscillation instead of fading
monotonically~\cite{yule1927method}: a damping factor $r$ per step and a
period of $P$ steps give $\phi_1=2r\cos(2\pi/P)$ and $\phi_2=-r^2$, and
the scorecard uses $r=0.95$ and $P=16$. The cycles are irregular because
every new shock moves their phase and size. The second construction is
measurement noise. On the AR(1), AR(2), and random-walk rows the forecaster
sees $X_t=Y_t+\eta_t$ with independent
$\eta_t\sim\mathcal{N}(0,\tau^2)$, so the current state of the process is
hidden; the innovation scale is fixed and the noise level sets $\tau$.
These are linear Gaussian state-space models, so the answer key stays
exact: a Kalman filter over the context gives the conditional mean and
covariance of the future observations~\cite{kalman1960new}, and $\tau=0$
recovers the formulas above. The two remaining scorecard rows, a trend
with a seasonal cycle and a sudden level shift, use the composite scenario
of Appendix~\ref{app:extended} with fresh Gaussian noise, whose answer key
is its deterministic part.

Reference validation compares analytic quantities with conditional
simulations and tests boundary cases such as zero noise and correlation
validity. Simulation agreement is an implementation check; it is not an
independent proof when simulator and reference share code. Algebraic
checks and explicit small cases are therefore also needed. Every
synthetic condition in both studies comes from the five families of
Table~\ref{tab:families}.
Appendix~\ref{app:extended} describes the composed, financial, and joint
scenarios in the software catalog, which neither study uses. Missingness,
exogenous covariates, and irregular sampling are outside the implemented
scenario scope.

\section{Output checks and reports}
\label{sec:contracts}

\subsection{What a run checks}

Table~\ref{tab:checks} lists the questions a run asks about a forecast.
The first six are output checks and are always asked: a forecast that
fails one is not compared with the answer key. The avoidable error is
measured whenever the scenario has an answer key. The shift check and the
joint-spread check run on request. A model can break the shift check by
design, for example through regularization, clipping, or features of the
absolute level, and its per-series form exists only in the Study 2 code;
the joint-spread check needs sample paths and an exact sampler
(Section~\ref{sec:joint-spread}). Each check
ends in one of five verdicts: pass, fail, can't tell
(\code{INCONCLUSIVE}), not applicable (\code{UNSUPPORTED}), or crashed
(\code{ERROR}). A model exception is not silently converted into a
numerical loss, and a kind of output the model does not return is
recorded as such instead of being counted as a pass.

\begin{table}[t]
\centering\small
\begin{tabular}{@{}>{\raggedright\arraybackslash}p{0.29\linewidth}>{\raggedright\arraybackslash}p{0.35\linewidth}>{\raggedright\arraybackslash}p{0.28\linewidth}@{}}
\toprule
Check & Question it answers & Example failure it catches \\
\midrule
Runs without errors & Do fit and predict finish? & An exception inside the model \\[0.3em]
Forecasts the requested dates & Are the dates the requested ones, in order? & Dates shifted one step early \\[0.3em]
Returns every series, in order & Are the series the requested ones, in order? & Two series labels swapped \\[0.3em]
Right number of steps and series & Is the shape right, for points, quantiles and samples? & One step short; sample axes transposed \\[0.3em]
No missing or infinite values & Are all values finite? & A NaN \\[0.3em]
Quantiles in order & Do the quantiles never cross? & The 90\% quantile below the 10\% one \\[0.3em]
Shift check & Add $c$ to the history: does the forecast rise by $c$? & Inverse scaling left out \\[0.3em]
Per-series shift check (experimental) & Shift each series by a different amount: does each forecast follow its own? & Two series' values swapped under correct labels \\[0.3em]
Avoidable error & How far is the forecast from the answer key? & Misses the trend (scored 0--100) \\[0.3em]
Joint spread matches the answer key & Do sample paths spread and move together like the true future? & Spread shrunk; paths shuffled \\
\bottomrule
\end{tabular}
\caption{What a run checks. The first six rows are the output checks,
under the titles the package reports.}
\label{tab:checks}
\end{table}

Point values have shape $(H,D)$, marginal quantiles $(H,D,Q)$ at strictly
increasing levels in $(0,1)$, and sample trajectories $(S,H,D)$; every
sample index denotes one complete joint future across both horizon and
target axes. Exact shapes prevent accidental NumPy broadcasting from
silently changing what the forecast stands for. The output checks catch
an explicitly displaced date or a reversed list of series names. Values
shifted internally under apparently correct labels pass them; catching
those takes an answer key or a shift check. The implementation supplies a
bounded set of checks rather than an exhaustive information-flow
verifier.

\subsection{Grades, seeds, and reports}

The scorecard grades capability without calling it a bug. It turns a
cell's avoidable error into a score from 0 to 100
(Equation~\eqref{eq:score}) and a band: 90 or more nails it, 70 good, 40
struggles, below 40 fails. A band describes what the model can learn. The
underlying diagnostic still returns \code{INCONCLUSIVE}, because no
single tolerance separates a weak model from a faulty one. A forecast
that crashes or fails an output check scores zero. A low score can
reflect statistical estimation, a deliberately limited model class, or an
implementation fault.
Coverage computed across one path's dependent horizon-target coordinates
is also descriptive; it does not establish calibration. The report
identifies the affected property and a reproducible case, not a unique
source-code cause. Suspected explanations require an intervention.

\code{Pipeline(model).scorecard()} runs the scenario grid and returns the
picture together with a machine-readable report (\code{card.report}) and
every series and forecast (\code{card.run}), so a forecast plot and an
error plot describe the same execution without refitting.
\code{Pipeline.run} does the same for any list of scenarios, and
\code{Pipeline.test} returns the report alone. Each run takes one seed:
instance $i$ of every kind of series uses \code{seed + i}, so the noise
levels are a paired comparison, and \code{seed=None} draws a fresh seed
and records it. Reproducibility requires more than a seed: the
case parameters, model identification, forecast dimensions, implementation
and environment versions, diagnostic settings, and all outcomes must be
retained. Experiment artifacts record these together with the generated
results. These metadata do not make external training code
deterministic.

\subsection{A reproducible diagnosis and repair workflow}

Table~\ref{tab:workflow} summarizes the operational sequence. A failed
case can be regenerated from its recorded scenario
configuration. Reproduction establishes that a property fails under those
conditions; investigating and modifying the implementation remain explicit
developer actions.

The numbers behind a verdict determine the next action. A horizon or date
mismatch is a failed output check. A covariance rejection is evidence
against the specified sampling law under the stated Monte Carlo null. A
failed shift check contradicts that relation, provided the
transformation preserves the model's intended task. A large avoidable
error without such a failure is a capability observation. A useful
repair recheck repeats the original case, its transformed counterpart
when applicable, and clean controls, retaining failures as well as passes.

\begin{table}[t]
\centering\small
\begin{tabular}{@{}p{0.21\linewidth}p{0.21\linewidth}p{0.21\linewidth}p{0.25\linewidth}@{}}
\toprule
1. Specify & 2. Execute & 3. Inspect & 4. Intervene and recheck \\
\midrule
Context and forecast axes; permitted information; required outputs;
answer key or shift relation &
Fit and forecast; run the output checks; evaluate the other applicable checks &
Verdicts, the numbers behind them, saved predictions, and a reproducible case &
Change a candidate cause; rerun the paired case and clean controls \\
\bottomrule
\end{tabular}
\caption{TimeGym's diagnostic workflow. The software records failed
properties and replay information; a developer supplies the candidate
explanation and repair.}
\label{tab:workflow}
\end{table}

\subsection{The joint-spread check}
\label{sec:joint-spread}

The check reported as ``Joint spread matches the answer key'' is an
optional Monte Carlo test of the ensemble covariance. It runs when a
reference covariance, an exact conditional sampler, and at least 20 joint
sample paths with $K\geq2$ are available. Study 1 uses
analytical covariance; Appendix~\ref{sec:joint-stresses} describes an
independently estimated, fixed target for general CCC-GARCH cross-moments. Let $\widehat\Sigma_X$ use the
usual $S-1$ denominator. The discrepancy is
\begin{equation}
 T(X)=\frac{\|\widehat\Sigma_X-\Sigma_F\|_F}{\|\Sigma_F\|_F}.
\end{equation}
A zero reference norm is reported as \code{UNSUPPORTED}. Draw $B$
independent calibration ensembles, each containing $S$ independent paths
from the same conditional law, and calculate
\begin{equation}
 p_{\mathrm{MC}}=\frac{1+\sum_{b=1}^{B}
                \mathbf{1}\{T(X_b)\geq T(X)\}}{B+1}.
 \label{eq:mc}
\end{equation}
Under the null that the submitted ensemble itself consists of independent
draws from that reference law, the $B+1$ statistics are exchangeable.
The Monte Carlo rank rule therefore controls the per-check rejection
probability at the attainable discrete levels, with conservative treatment
of ties~\cite{hope1968montecarlo}.
Study 1 uses $B=99$ and rejects when $p_{\mathrm{MC}}\leq0.05$.
Insufficient rank resolution for the requested level is
\code{INCONCLUSIVE}. No adjustment across multiple cases is automatic.

This is a reference-sampling consistency check, not a test of all forms of
forecast calibration. Variance errors as well as dependence errors can
trigger it. Failure to reject does not establish equality of covariance
or distributions. In particular, an ensemble generated by deterministic
quadrature, dependent Monte Carlo, or a learned parameter estimate does
not satisfy the stated null merely because it is a useful forecast.
The Gaussian reference is fully specified in Study 1, which makes
null simulation feasible but limits direct application to real data.

\subsection{Scores and a dependence counterexample}

\begin{figure}[!t]
\centering
% BEGIN VECTOR FIGURE paper-dependence
% Generated vector scene; SVG and PNG use the same plotted geometry.
\begingroup
\definecolor{tg0000ff}{HTML}{0000ff}
\definecolor{tg18314b}{HTML}{18314b}
\definecolor{tg526577}{HTML}{526577}
\definecolor{tg808080}{HTML}{808080}
\definecolor{tgc65c16}{HTML}{c65c16}
\definecolor{tge1e7ec}{HTML}{e1e7ec}
\resizebox{\linewidth}{!}{%
\begin{tikzpicture}[x=1bp,y=-1bp]
\path[use as bounding box] (0,0) rectangle (448,376);
\node[anchor=west,inner sep=0pt,outer sep=0pt,text=tg18314b,font={\sffamily\fontsize{10bp}{12bp}\selectfont\bfseries}] at (12.000,12.000) {Dependence changes while every marginal stays identical};
\node[anchor=west,inner sep=0pt,outer sep=0pt,text=tg526577,font={\sffamily\fontsize{7.5bp}{9bp}\selectfont}] at (12.000,28.000) {128 joint samples from one predefined AR(1) case; independent coordinate-wise permutations.};
\node[anchor=center,inner sep=0pt,outer sep=0pt,text=tg18314b,font={\sffamily\fontsize{8bp}{9.6bp}\selectfont\bfseries}] at (84.500,46.000) {Reference samples};
\node[anchor=center,inner sep=0pt,outer sep=0pt,text=tg526577,font={\sffamily\fontsize{7bp}{8.4bp}\selectfont}] at (84.500,60.000) {First-horizon correlation: +0.67};
\draw[color=tge1e7ec,line width=0.4bp,line cap=round] (32.000,157.000) -- (137.000,157.000);
\node[anchor=east,inner sep=0pt,outer sep=0pt,text=tg526577,font={\sffamily\fontsize{7bp}{8.4bp}\selectfont}] at (27.000,157.000) {-4};
\draw[color=tge1e7ec,line width=0.4bp,line cap=round] (32.000,115.000) -- (137.000,115.000);
\node[anchor=east,inner sep=0pt,outer sep=0pt,text=tg526577,font={\sffamily\fontsize{7bp}{8.4bp}\selectfont}] at (27.000,115.000) {0};
\draw[color=tge1e7ec,line width=0.4bp,line cap=round] (32.000,73.000) -- (137.000,73.000);
\node[anchor=east,inner sep=0pt,outer sep=0pt,text=tg526577,font={\sffamily\fontsize{7bp}{8.4bp}\selectfont}] at (27.000,73.000) {4};
\draw[color=tg808080,line width=0.5bp,line cap=round] (32.000,73.000) -- (32.000,157.000);
\draw[color=tg808080,line width=0.5bp,line cap=round] (32.000,157.000) -- (137.000,157.000);
\draw[color=tg808080,line width=0.5bp,line cap=round] (32.000,157.000) -- (32.000,159.000);
\node[anchor=center,inner sep=0pt,outer sep=0pt,text=tg526577,font={\sffamily\fontsize{7bp}{8.4bp}\selectfont}] at (32.000,167.000) {-4};
\draw[color=tg808080,line width=0.5bp,line cap=round] (84.500,157.000) -- (84.500,159.000);
\node[anchor=center,inner sep=0pt,outer sep=0pt,text=tg526577,font={\sffamily\fontsize{7bp}{8.4bp}\selectfont}] at (84.500,167.000) {0};
\draw[color=tg808080,line width=0.5bp,line cap=round] (137.000,157.000) -- (137.000,159.000);
\node[anchor=center,inner sep=0pt,outer sep=0pt,text=tg526577,font={\sffamily\fontsize{7bp}{8.4bp}\selectfont}] at (137.000,167.000) {4};
\fill[tg0000ff] (96.578,114.576) circle[radius=0.9bp];
\fill[tg0000ff] (86.190,119.291) circle[radius=0.9bp];
\fill[tg0000ff] (76.975,139.657) circle[radius=0.9bp];
\fill[tg0000ff] (83.064,123.449) circle[radius=0.9bp];
\fill[tg0000ff] (83.741,133.982) circle[radius=0.9bp];
\fill[tg0000ff] (94.187,114.718) circle[radius=0.9bp];
\fill[tg0000ff] (79.184,127.437) circle[radius=0.9bp];
\fill[tg0000ff] (72.580,125.321) circle[radius=0.9bp];
\fill[tg0000ff] (76.736,126.785) circle[radius=0.9bp];
\fill[tg0000ff] (112.017,104.610) circle[radius=0.9bp];
\fill[tg0000ff] (55.007,124.975) circle[radius=0.9bp];
\fill[tg0000ff] (102.412,117.811) circle[radius=0.9bp];
\fill[tg0000ff] (66.855,141.650) circle[radius=0.9bp];
\fill[tg0000ff] (70.727,122.601) circle[radius=0.9bp];
\fill[tg0000ff] (79.091,124.428) circle[radius=0.9bp];
\fill[tg0000ff] (80.728,131.387) circle[radius=0.9bp];
\fill[tg0000ff] (80.001,134.680) circle[radius=0.9bp];
\fill[tg0000ff] (84.155,118.101) circle[radius=0.9bp];
\fill[tg0000ff] (78.655,127.928) circle[radius=0.9bp];
\fill[tg0000ff] (72.134,133.167) circle[radius=0.9bp];
\fill[tg0000ff] (74.491,127.945) circle[radius=0.9bp];
\fill[tg0000ff] (78.212,125.087) circle[radius=0.9bp];
\fill[tg0000ff] (72.722,125.185) circle[radius=0.9bp];
\fill[tg0000ff] (80.976,122.675) circle[radius=0.9bp];
\fill[tg0000ff] (90.795,102.245) circle[radius=0.9bp];
\fill[tg0000ff] (92.155,113.554) circle[radius=0.9bp];
\fill[tg0000ff] (96.097,115.163) circle[radius=0.9bp];
\fill[tg0000ff] (68.854,132.913) circle[radius=0.9bp];
\fill[tg0000ff] (66.992,121.149) circle[radius=0.9bp];
\fill[tg0000ff] (69.461,120.124) circle[radius=0.9bp];
\fill[tg0000ff] (104.253,111.576) circle[radius=0.9bp];
\fill[tg0000ff] (75.102,126.742) circle[radius=0.9bp];
\fill[tg0000ff] (86.347,114.734) circle[radius=0.9bp];
\fill[tg0000ff] (60.346,134.581) circle[radius=0.9bp];
\fill[tg0000ff] (67.352,135.315) circle[radius=0.9bp];
\fill[tg0000ff] (73.537,144.499) circle[radius=0.9bp];
\fill[tg0000ff] (75.559,145.407) circle[radius=0.9bp];
\fill[tg0000ff] (113.129,108.787) circle[radius=0.9bp];
\fill[tg0000ff] (81.005,120.224) circle[radius=0.9bp];
\fill[tg0000ff] (100.967,120.008) circle[radius=0.9bp];
\fill[tg0000ff] (93.672,113.308) circle[radius=0.9bp];
\fill[tg0000ff] (90.637,131.708) circle[radius=0.9bp];
\fill[tg0000ff] (80.733,116.520) circle[radius=0.9bp];
\fill[tg0000ff] (98.034,110.522) circle[radius=0.9bp];
\fill[tg0000ff] (73.750,116.837) circle[radius=0.9bp];
\fill[tg0000ff] (81.779,119.025) circle[radius=0.9bp];
\fill[tg0000ff] (72.916,119.089) circle[radius=0.9bp];
\fill[tg0000ff] (82.933,130.645) circle[radius=0.9bp];
\fill[tg0000ff] (90.128,117.725) circle[radius=0.9bp];
\fill[tg0000ff] (70.207,126.339) circle[radius=0.9bp];
\fill[tg0000ff] (89.973,105.872) circle[radius=0.9bp];
\fill[tg0000ff] (94.482,117.198) circle[radius=0.9bp];
\fill[tg0000ff] (87.159,122.325) circle[radius=0.9bp];
\fill[tg0000ff] (60.642,149.184) circle[radius=0.9bp];
\fill[tg0000ff] (83.501,127.309) circle[radius=0.9bp];
\fill[tg0000ff] (93.578,121.703) circle[radius=0.9bp];
\fill[tg0000ff] (88.665,108.343) circle[radius=0.9bp];
\fill[tg0000ff] (78.903,127.227) circle[radius=0.9bp];
\fill[tg0000ff] (77.888,128.270) circle[radius=0.9bp];
\fill[tg0000ff] (78.283,137.683) circle[radius=0.9bp];
\fill[tg0000ff] (112.076,120.319) circle[radius=0.9bp];
\fill[tg0000ff] (97.830,124.088) circle[radius=0.9bp];
\fill[tg0000ff] (46.026,146.218) circle[radius=0.9bp];
\fill[tg0000ff] (78.377,119.504) circle[radius=0.9bp];
\fill[tg0000ff] (87.546,123.732) circle[radius=0.9bp];
\fill[tg0000ff] (66.045,127.721) circle[radius=0.9bp];
\fill[tg0000ff] (106.384,109.484) circle[radius=0.9bp];
\fill[tg0000ff] (96.865,124.390) circle[radius=0.9bp];
\fill[tg0000ff] (81.989,131.947) circle[radius=0.9bp];
\fill[tg0000ff] (82.344,119.394) circle[radius=0.9bp];
\fill[tg0000ff] (80.087,128.444) circle[radius=0.9bp];
\fill[tg0000ff] (87.583,134.975) circle[radius=0.9bp];
\fill[tg0000ff] (73.915,124.472) circle[radius=0.9bp];
\fill[tg0000ff] (100.697,122.716) circle[radius=0.9bp];
\fill[tg0000ff] (83.153,129.810) circle[radius=0.9bp];
\fill[tg0000ff] (134.227,89.861) circle[radius=0.9bp];
\fill[tg0000ff] (73.746,129.939) circle[radius=0.9bp];
\fill[tg0000ff] (96.011,121.588) circle[radius=0.9bp];
\fill[tg0000ff] (75.501,121.511) circle[radius=0.9bp];
\fill[tg0000ff] (84.416,112.937) circle[radius=0.9bp];
\fill[tg0000ff] (80.965,122.100) circle[radius=0.9bp];
\fill[tg0000ff] (113.210,114.124) circle[radius=0.9bp];
\fill[tg0000ff] (85.572,128.411) circle[radius=0.9bp];
\fill[tg0000ff] (66.381,130.992) circle[radius=0.9bp];
\fill[tg0000ff] (110.500,101.802) circle[radius=0.9bp];
\fill[tg0000ff] (77.397,127.481) circle[radius=0.9bp];
\fill[tg0000ff] (74.640,114.396) circle[radius=0.9bp];
\fill[tg0000ff] (79.777,116.629) circle[radius=0.9bp];
\fill[tg0000ff] (82.223,117.095) circle[radius=0.9bp];
\fill[tg0000ff] (67.855,128.101) circle[radius=0.9bp];
\fill[tg0000ff] (105.255,103.951) circle[radius=0.9bp];
\fill[tg0000ff] (78.887,141.904) circle[radius=0.9bp];
\fill[tg0000ff] (94.087,108.083) circle[radius=0.9bp];
\fill[tg0000ff] (49.793,134.346) circle[radius=0.9bp];
\fill[tg0000ff] (93.585,112.980) circle[radius=0.9bp];
\fill[tg0000ff] (95.102,99.312) circle[radius=0.9bp];
\fill[tg0000ff] (80.972,122.081) circle[radius=0.9bp];
\fill[tg0000ff] (94.987,123.405) circle[radius=0.9bp];
\fill[tg0000ff] (100.662,133.059) circle[radius=0.9bp];
\fill[tg0000ff] (96.531,124.419) circle[radius=0.9bp];
\fill[tg0000ff] (72.954,140.171) circle[radius=0.9bp];
\fill[tg0000ff] (74.963,128.171) circle[radius=0.9bp];
\fill[tg0000ff] (86.666,118.942) circle[radius=0.9bp];
\fill[tg0000ff] (90.416,108.032) circle[radius=0.9bp];
\fill[tg0000ff] (74.544,132.626) circle[radius=0.9bp];
\fill[tg0000ff] (68.519,128.843) circle[radius=0.9bp];
\fill[tg0000ff] (84.526,120.634) circle[radius=0.9bp];
\fill[tg0000ff] (90.743,114.588) circle[radius=0.9bp];
\fill[tg0000ff] (104.335,107.571) circle[radius=0.9bp];
\fill[tg0000ff] (110.833,127.110) circle[radius=0.9bp];
\fill[tg0000ff] (99.266,120.366) circle[radius=0.9bp];
\fill[tg0000ff] (101.941,117.683) circle[radius=0.9bp];
\fill[tg0000ff] (82.357,123.499) circle[radius=0.9bp];
\fill[tg0000ff] (92.948,116.235) circle[radius=0.9bp];
\fill[tg0000ff] (75.833,125.398) circle[radius=0.9bp];
\fill[tg0000ff] (91.963,124.522) circle[radius=0.9bp];
\fill[tg0000ff] (74.003,115.870) circle[radius=0.9bp];
\fill[tg0000ff] (116.407,105.730) circle[radius=0.9bp];
\fill[tg0000ff] (93.133,111.529) circle[radius=0.9bp];
\fill[tg0000ff] (101.041,110.543) circle[radius=0.9bp];
\fill[tg0000ff] (74.891,127.260) circle[radius=0.9bp];
\fill[tg0000ff] (66.143,136.320) circle[radius=0.9bp];
\fill[tg0000ff] (83.065,110.550) circle[radius=0.9bp];
\fill[tg0000ff] (83.971,114.056) circle[radius=0.9bp];
\fill[tg0000ff] (70.485,123.377) circle[radius=0.9bp];
\fill[tg0000ff] (92.107,113.048) circle[radius=0.9bp];
\fill[tg0000ff] (98.193,101.815) circle[radius=0.9bp];
\fill[tg0000ff] (92.177,124.494) circle[radius=0.9bp];
\node[anchor=center,inner sep=0pt,outer sep=0pt,text=tg526577,font={\sffamily\fontsize{7bp}{8.4bp}\selectfont}] at (84.500,183.000) {Target 1 (x), target 2 (y)};
\node[anchor=center,inner sep=0pt,outer sep=0pt,text=tg18314b,font={\sffamily\fontsize{8bp}{9.6bp}\selectfont\bfseries}] at (231.500,46.000) {Shuffled samples};
\node[anchor=center,inner sep=0pt,outer sep=0pt,text=tg526577,font={\sffamily\fontsize{7bp}{8.4bp}\selectfont}] at (231.500,60.000) {First-horizon correlation: -0.05};
\draw[color=tge1e7ec,line width=0.4bp,line cap=round] (179.000,157.000) -- (284.000,157.000);
\node[anchor=east,inner sep=0pt,outer sep=0pt,text=tg526577,font={\sffamily\fontsize{7bp}{8.4bp}\selectfont}] at (174.000,157.000) {-4};
\draw[color=tge1e7ec,line width=0.4bp,line cap=round] (179.000,115.000) -- (284.000,115.000);
\node[anchor=east,inner sep=0pt,outer sep=0pt,text=tg526577,font={\sffamily\fontsize{7bp}{8.4bp}\selectfont}] at (174.000,115.000) {0};
\draw[color=tge1e7ec,line width=0.4bp,line cap=round] (179.000,73.000) -- (284.000,73.000);
\node[anchor=east,inner sep=0pt,outer sep=0pt,text=tg526577,font={\sffamily\fontsize{7bp}{8.4bp}\selectfont}] at (174.000,73.000) {4};
\draw[color=tg808080,line width=0.5bp,line cap=round] (179.000,73.000) -- (179.000,157.000);
\draw[color=tg808080,line width=0.5bp,line cap=round] (179.000,157.000) -- (284.000,157.000);
\draw[color=tg808080,line width=0.5bp,line cap=round] (179.000,157.000) -- (179.000,159.000);
\node[anchor=center,inner sep=0pt,outer sep=0pt,text=tg526577,font={\sffamily\fontsize{7bp}{8.4bp}\selectfont}] at (179.000,167.000) {-4};
\draw[color=tg808080,line width=0.5bp,line cap=round] (231.500,157.000) -- (231.500,159.000);
\node[anchor=center,inner sep=0pt,outer sep=0pt,text=tg526577,font={\sffamily\fontsize{7bp}{8.4bp}\selectfont}] at (231.500,167.000) {0};
\draw[color=tg808080,line width=0.5bp,line cap=round] (284.000,157.000) -- (284.000,159.000);
\node[anchor=center,inner sep=0pt,outer sep=0pt,text=tg526577,font={\sffamily\fontsize{7bp}{8.4bp}\selectfont}] at (284.000,167.000) {4};
\fill[tgc65c16] (221.003,116.629) circle[radius=0.9bp];
\fill[tgc65c16] (243.531,120.634) circle[radius=0.9bp];
\fill[tgc65c16] (226.184,132.626) circle[radius=0.9bp];
\fill[tgc65c16] (215.854,133.059) circle[radius=0.9bp];
\fill[tgc65c16] (234.546,122.100) circle[radius=0.9bp];
\fill[tgc65c16] (221.891,146.218) circle[radius=0.9bp];
\fill[tgc65c16] (215.519,111.529) circle[radius=0.9bp];
\fill[tgc65c16] (213.045,103.951) circle[radius=0.9bp];
\fill[tgc65c16] (213.855,124.975) circle[radius=0.9bp];
\fill[tgc65c16] (221.544,121.511) circle[radius=0.9bp];
\fill[tgc65c16] (237.743,124.494) circle[radius=0.9bp];
\fill[tgc65c16] (213.143,117.683) circle[radius=0.9bp];
\fill[tgc65c16] (235.665,117.811) circle[radius=0.9bp];
\fill[tgc65c16] (222.833,139.657) circle[radius=0.9bp];
\fill[tgc65c16] (213.381,127.309) circle[radius=0.9bp];
\fill[tgc65c16] (223.975,117.198) circle[radius=0.9bp];
\fill[tgc65c16] (228.989,120.124) circle[radius=0.9bp];
\fill[tgc65c16] (227.972,121.149) circle[radius=0.9bp];
\fill[tgc65c16] (214.855,133.982) circle[radius=0.9bp];
\fill[tgc65c16] (243.865,104.610) circle[radius=0.9bp];
\fill[tgc65c16] (217.727,114.124) circle[radius=0.9bp];
\fill[tgc65c16] (225.903,114.056) circle[radius=0.9bp];
\fill[tgc65c16] (241.087,123.377) circle[radius=0.9bp];
\fill[tgc65c16] (230.153,115.870) circle[radius=0.9bp];
\fill[tgc65c16] (251.335,112.937) circle[radius=0.9bp];
\fill[tgc65c16] (236.973,123.449) circle[radius=0.9bp];
\fill[tgc65c16] (220.537,108.787) circle[radius=0.9bp];
\fill[tgc65c16] (263.407,128.270) circle[radius=0.9bp];
\fill[tgc65c16] (207.346,125.321) circle[radius=0.9bp];
\fill[tgc65c16] (227.087,116.520) circle[radius=0.9bp];
\fill[tgc65c16] (225.655,125.087) circle[radius=0.9bp];
\fill[tgc65c16] (227.001,114.718) circle[radius=0.9bp];
\fill[tgc65c16] (245.193,124.428) circle[radius=0.9bp];
\fill[tgc65c16] (207.642,124.390) circle[radius=0.9bp];
\fill[tgc65c16] (239.177,129.939) circle[radius=0.9bp];
\fill[tgc65c16] (237.128,105.872) circle[radius=0.9bp];
\fill[tgc65c16] (247.697,149.184) circle[radius=0.9bp];
\fill[tgc65c16] (223.736,108.343) circle[radius=0.9bp];
\fill[tgc65c16] (241.482,133.167) circle[radius=0.9bp];
\fill[tgc65c16] (241.987,99.312) circle[radius=0.9bp];
\fill[tgc65c16] (243.097,141.650) circle[radius=0.9bp];
\fill[tgc65c16] (247.662,108.083) circle[radius=0.9bp];
\fill[tgc65c16] (238.963,134.581) circle[radius=0.9bp];
\fill[tgc65c16] (217.207,102.245) circle[radius=0.9bp];
\fill[tgc65c16] (220.750,126.785) circle[radius=0.9bp];
\fill[tgc65c16] (229.344,108.032) circle[radius=0.9bp];
\fill[tgc65c16] (237.795,126.742) circle[radius=0.9bp];
\fill[tgc65c16] (222.501,115.163) circle[radius=0.9bp];
\fill[tgc65c16] (219.580,128.411) circle[radius=0.9bp];
\fill[tgc65c16] (244.830,101.815) circle[radius=0.9bp];
\fill[tgc65c16] (237.416,131.947) circle[radius=0.9bp];
\fill[tgc65c16] (230.065,121.588) circle[radius=0.9bp];
\fill[tgc65c16] (228.005,113.308) circle[radius=0.9bp];
\fill[tgc65c16] (230.741,127.721) circle[radius=0.9bp];
\fill[tgc65c16] (214.352,118.101) circle[radius=0.9bp];
\fill[tgc65c16] (220.746,128.101) circle[radius=0.9bp];
\fill[tgc65c16] (196.793,124.522) circle[radius=0.9bp];
\fill[tgc65c16] (246.266,114.576) circle[radius=0.9bp];
\fill[tgc65c16] (231.416,116.837) circle[radius=0.9bp];
\fill[tgc65c16] (220.915,118.942) circle[radius=0.9bp];
\fill[tgc65c16] (239.107,116.235) circle[radius=0.9bp];
\fill[tgc65c16] (248.041,105.730) circle[radius=0.9bp];
\fill[tgc65c16] (221.640,120.366) circle[radius=0.9bp];
\fill[tgc65c16] (229.357,120.319) circle[radius=0.9bp];
\fill[tgc65c16] (219.134,127.437) circle[radius=0.9bp];
\fill[tgc65c16] (217.485,140.171) circle[radius=0.9bp];
\fill[tgc65c16] (230.501,119.291) circle[radius=0.9bp];
\fill[tgc65c16] (219.916,113.554) circle[radius=0.9bp];
\fill[tgc65c16] (219.722,128.171) circle[radius=0.9bp];
\fill[tgc65c16] (243.578,114.588) circle[radius=0.9bp];
\fill[tgc65c16] (229.933,122.325) circle[radius=0.9bp];
\fill[tgc65c16] (241.187,131.387) circle[radius=0.9bp];
\fill[tgc65c16] (245.034,119.504) circle[radius=0.9bp];
\fill[tgc65c16] (227.728,112.980) circle[radius=0.9bp];
\fill[tgc65c16] (229.223,119.394) circle[radius=0.9bp];
\fill[tgc65c16] (248.941,136.320) circle[radius=0.9bp];
\fill[tgc65c16] (227.976,114.734) circle[radius=0.9bp];
\fill[tgc65c16] (243.011,107.571) circle[radius=0.9bp];
\fill[tgc65c16] (219.954,124.472) circle[radius=0.9bp];
\fill[tgc65c16] (237.637,132.913) circle[radius=0.9bp];
\fill[tgc65c16] (234.583,89.861) circle[radius=0.9bp];
\fill[tgc65c16] (239.155,127.110) circle[radius=0.9bp];
\fill[tgc65c16] (247.967,125.398) circle[radius=0.9bp];
\fill[tgc65c16] (224.397,101.802) circle[radius=0.9bp];
\fill[tgc65c16] (260.129,124.419) circle[radius=0.9bp];
\fill[tgc65c16] (225.377,120.008) circle[radius=0.9bp];
\fill[tgc65c16] (225.283,131.708) circle[radius=0.9bp];
\fill[tgc65c16] (193.026,110.543) circle[radius=0.9bp];
\fill[tgc65c16] (228.779,119.025) circle[radius=0.9bp];
\fill[tgc65c16] (240.672,117.095) circle[radius=0.9bp];
\fill[tgc65c16] (257.500,122.601) circle[radius=0.9bp];
\fill[tgc65c16] (227.733,117.725) circle[radius=0.9bp];
\fill[tgc65c16] (226.091,125.185) circle[radius=0.9bp];
\fill[tgc65c16] (259.076,127.481) circle[radius=0.9bp];
\fill[tgc65c16] (231.155,137.683) circle[radius=0.9bp];
\fill[tgc65c16] (239.948,134.975) circle[radius=0.9bp];
\fill[tgc65c16] (227.965,123.499) circle[radius=0.9bp];
\fill[tgc65c16] (213.992,127.928) circle[radius=0.9bp];
\fill[tgc65c16] (233.666,145.407) circle[radius=0.9bp];
\fill[tgc65c16] (232.572,110.522) circle[radius=0.9bp];
\fill[tgc65c16] (221.491,128.444) circle[radius=0.9bp];
\fill[tgc65c16] (221.963,127.260) circle[radius=0.9bp];
\fill[tgc65c16] (216.461,135.315) circle[radius=0.9bp];
\fill[tgc65c16] (240.585,110.550) circle[radius=0.9bp];
\fill[tgc65c16] (230.064,130.645) circle[radius=0.9bp];
\fill[tgc65c16] (231.526,141.904) circle[radius=0.9bp];
\fill[tgc65c16] (202.007,122.716) circle[radius=0.9bp];
\fill[tgc65c16] (225.212,127.945) circle[radius=0.9bp];
\fill[tgc65c16] (259.017,123.405) circle[radius=0.9bp];
\fill[tgc65c16] (222.559,129.810) circle[radius=0.9bp];
\fill[tgc65c16] (242.102,123.732) circle[radius=0.9bp];
\fill[tgc65c16] (226.777,128.843) circle[radius=0.9bp];
\fill[tgc65c16] (251.253,111.576) circle[radius=0.9bp];
\fill[tgc65c16] (222.102,114.396) circle[radius=0.9bp];
\fill[tgc65c16] (240.578,144.499) circle[radius=0.9bp];
\fill[tgc65c16] (260.210,124.088) circle[radius=0.9bp];
\fill[tgc65c16] (253.384,134.346) circle[radius=0.9bp];
\fill[tgc65c16] (281.227,130.992) circle[radius=0.9bp];
\fill[tgc65c16] (230.971,127.227) circle[radius=0.9bp];
\fill[tgc65c16] (224.888,121.703) circle[radius=0.9bp];
\fill[tgc65c16] (240.133,113.048) circle[radius=0.9bp];
\fill[tgc65c16] (234.159,109.484) circle[radius=0.9bp];
\fill[tgc65c16] (249.412,134.680) circle[radius=0.9bp];
\fill[tgc65c16] (252.255,120.224) circle[radius=0.9bp];
\fill[tgc65c16] (257.833,122.081) circle[radius=0.9bp];
\fill[tgc65c16] (225.887,122.675) circle[radius=0.9bp];
\fill[tgc65c16] (233.347,126.339) circle[radius=0.9bp];
\fill[tgc65c16] (233.190,119.089) circle[radius=0.9bp];
\node[anchor=center,inner sep=0pt,outer sep=0pt,text=tg526577,font={\sffamily\fontsize{7bp}{8.4bp}\selectfont}] at (231.500,183.000) {Target 1 (x), target 2 (y)};
\node[anchor=center,inner sep=0pt,outer sep=0pt,text=tg18314b,font={\sffamily\fontsize{8bp}{9.6bp}\selectfont\bfseries}] at (379.000,46.000) {Marginals overlap};
\node[anchor=center,inner sep=0pt,outer sep=0pt,text=tg526577,font={\sffamily\fontsize{7bp}{8.4bp}\selectfont}] at (379.000,60.000) {Target 1 empirical CDF};
\draw[color=tge1e7ec,line width=0.4bp,line cap=round] (326.000,157.000) -- (432.000,157.000);
\node[anchor=east,inner sep=0pt,outer sep=0pt,text=tg526577,font={\sffamily\fontsize{7bp}{8.4bp}\selectfont}] at (321.000,157.000) {0};
\draw[color=tge1e7ec,line width=0.4bp,line cap=round] (326.000,115.000) -- (432.000,115.000);
\node[anchor=east,inner sep=0pt,outer sep=0pt,text=tg526577,font={\sffamily\fontsize{7bp}{8.4bp}\selectfont}] at (321.000,115.000) {0.5};
\draw[color=tge1e7ec,line width=0.4bp,line cap=round] (326.000,73.000) -- (432.000,73.000);
\node[anchor=east,inner sep=0pt,outer sep=0pt,text=tg526577,font={\sffamily\fontsize{7bp}{8.4bp}\selectfont}] at (321.000,73.000) {1};
\draw[color=tg808080,line width=0.5bp,line cap=round] (326.000,73.000) -- (326.000,157.000);
\draw[color=tg808080,line width=0.5bp,line cap=round] (326.000,157.000) -- (432.000,157.000);
\draw[color=tg808080,line width=0.5bp,line cap=round] (326.000,157.000) -- (326.000,159.000);
\node[anchor=center,inner sep=0pt,outer sep=0pt,text=tg526577,font={\sffamily\fontsize{7bp}{8.4bp}\selectfont}] at (326.000,167.000) {-4};
\draw[color=tg808080,line width=0.5bp,line cap=round] (379.000,157.000) -- (379.000,159.000);
\node[anchor=center,inner sep=0pt,outer sep=0pt,text=tg526577,font={\sffamily\fontsize{7bp}{8.4bp}\selectfont}] at (379.000,167.000) {0};
\draw[color=tg808080,line width=0.5bp,line cap=round] (432.000,157.000) -- (432.000,159.000);
\node[anchor=center,inner sep=0pt,outer sep=0pt,text=tg526577,font={\sffamily\fontsize{7bp}{8.4bp}\selectfont}] at (432.000,167.000) {4};
\draw[color=tg0000ff,line width=1.15bp,line cap=round] (326.000,157.000) -- (340.160,157.000) -- (340.160,156.344) -- (343.963,156.344) -- (343.963,155.688) -- (349.226,155.688) -- (349.226,155.031) -- (354.616,155.031) -- (354.616,154.375) -- (354.915,154.375) -- (354.915,153.719) -- (360.369,153.719) -- (360.369,153.062) -- (360.468,153.062) -- (360.468,152.406) -- (360.709,152.406) -- (360.709,151.750) -- (361.187,151.750) -- (361.187,151.094) -- (361.325,151.094) -- (361.325,150.438) -- (361.689,150.438) -- (361.689,149.781) -- (362.197,149.781) -- (362.197,149.125) -- (362.867,149.125) -- (362.867,148.469) -- (363.205,148.469) -- (363.205,147.812) -- (363.818,147.812) -- (363.818,147.156) -- (364.571,147.156) -- (364.571,146.500) -- (364.852,146.500) -- (364.852,145.844) -- (365.096,145.844) -- (365.096,145.188) -- (366.516,145.188) -- (366.516,144.531) -- (366.966,144.531) -- (366.966,143.875) -- (367.110,143.875) -- (367.110,143.219) -- (367.306,143.219) -- (367.306,142.562) -- (367.344,142.562) -- (367.344,141.906) -- (367.933,141.906) -- (367.933,141.250) -- (368.144,141.250) -- (368.144,140.594) -- (368.147,140.594) -- (368.147,139.938) -- (368.314,139.938) -- (368.314,139.281) -- (368.403,139.281) -- (368.403,138.625) -- (368.896,138.625) -- (368.896,137.969) -- (368.949,137.969) -- (368.949,137.312) -- (369.046,137.312) -- (369.046,136.656) -- (369.300,136.656) -- (369.300,136.000) -- (369.372,136.000) -- (369.372,135.344) -- (369.513,135.344) -- (369.513,134.688) -- (369.916,134.688) -- (369.916,134.031) -- (369.974,134.031) -- (369.974,133.375) -- (370.251,133.375) -- (370.251,132.719) -- (371.162,132.719) -- (371.162,132.062) -- (371.403,132.062) -- (371.403,131.406) -- (371.830,131.406) -- (371.830,130.750) -- (372.325,130.750) -- (372.325,130.094) -- (372.652,130.094) -- (372.652,129.438) -- (372.724,129.438) -- (372.724,128.781) -- (372.818,128.781) -- (372.818,128.125) -- (373.100,128.125) -- (373.100,127.469) -- (373.333,127.469) -- (373.333,126.812) -- (373.349,126.812) -- (373.349,126.156) -- (373.539,126.156) -- (373.539,125.500) -- (373.633,125.500) -- (373.633,124.844) -- (374.232,124.844) -- (374.232,124.188) -- (374.458,124.188) -- (374.458,123.531) -- (374.545,123.531) -- (374.545,122.875) -- (375.192,122.875) -- (375.192,122.219) -- (375.197,122.219) -- (375.197,121.562) -- (375.431,121.562) -- (375.431,120.906) -- (375.439,120.906) -- (375.439,120.250) -- (375.442,120.250) -- (375.442,119.594) -- (375.472,119.594) -- (375.472,118.938) -- (376.253,118.938) -- (376.253,118.281) -- (376.465,118.281) -- (376.465,117.625) -- (376.701,117.625) -- (376.701,116.969) -- (376.823,116.969) -- (376.823,116.312) -- (376.836,116.312) -- (376.836,115.656) -- (377.418,115.656) -- (377.418,115.000) -- (377.550,115.000) -- (377.550,114.344) -- (377.552,114.344) -- (377.552,113.688) -- (377.641,113.688) -- (377.641,113.031) -- (377.991,113.031) -- (377.991,112.375) -- (378.234,112.375) -- (378.234,111.719) -- (378.466,111.719) -- (378.466,111.062) -- (378.652,111.062) -- (378.652,110.406) -- (378.915,110.406) -- (378.915,109.750) -- (379.026,109.750) -- (379.026,109.094) -- (380.082,109.094) -- (380.082,108.438) -- (380.706,108.438) -- (380.706,107.781) -- (380.865,107.781) -- (380.865,107.125) -- (381.187,107.125) -- (381.187,106.469) -- (381.685,106.469) -- (381.685,105.812) -- (382.075,105.812) -- (382.075,105.156) -- (382.112,105.156) -- (382.112,104.500) -- (383.204,104.500) -- (383.204,103.844) -- (384.526,103.844) -- (384.526,103.188) -- (384.681,103.188) -- (384.681,102.531) -- (384.973,102.531) -- (384.973,101.875) -- (385.195,101.875) -- (385.195,101.219) -- (385.303,101.219) -- (385.303,100.562) -- (385.355,100.562) -- (385.355,99.906) -- (386.534,99.906) -- (386.534,99.250) -- (386.679,99.250) -- (386.679,98.594) -- (386.728,98.594) -- (386.728,97.938) -- (386.750,97.938) -- (386.750,97.281) -- (387.528,97.281) -- (387.528,96.625) -- (387.715,96.625) -- (387.715,95.969) -- (388.165,95.969) -- (388.165,95.312) -- (388.171,95.312) -- (388.171,94.656) -- (388.259,94.656) -- (388.259,94.000) -- (388.678,94.000) -- (388.678,93.344) -- (388.779,93.344) -- (388.779,92.688) -- (389.077,92.688) -- (389.077,92.031) -- (389.586,92.031) -- (389.586,91.375) -- (389.703,91.375) -- (389.703,90.719) -- (390.621,90.719) -- (390.621,90.062) -- (390.707,90.062) -- (390.707,89.406) -- (391.145,89.406) -- (391.145,88.750) -- (391.193,88.750) -- (391.193,88.094) -- (391.483,88.094) -- (391.483,87.438) -- (392.457,87.438) -- (392.457,86.781) -- (392.663,86.781) -- (392.663,86.125) -- (392.823,86.125) -- (392.823,85.469) -- (393.906,85.469) -- (393.906,84.812) -- (395.315,84.812) -- (395.315,84.156) -- (395.351,84.156) -- (395.351,83.500) -- (395.624,83.500) -- (395.624,82.844) -- (395.699,82.844) -- (395.699,82.188) -- (396.607,82.188) -- (396.607,81.531) -- (397.082,81.531) -- (397.082,80.875) -- (398.941,80.875) -- (398.941,80.219) -- (399.024,80.219) -- (399.024,79.562) -- (399.953,79.562) -- (399.953,78.906) -- (401.093,78.906) -- (401.093,78.250) -- (405.247,78.250) -- (405.247,77.594) -- (405.584,77.594) -- (405.584,76.938) -- (406.779,76.938) -- (406.779,76.281) -- (406.839,76.281) -- (406.839,75.625) -- (407.902,75.625) -- (407.902,74.969) -- (407.984,74.969) -- (407.984,74.312) -- (411.211,74.312) -- (411.211,73.656) -- (429.201,73.656) -- (429.201,73.000) -- (432.000,73.000);
\draw[color=tgc65c16,line width=0.75bp,line cap=round,dash pattern=on 2bp off 2bp] (326.000,157.000) -- (340.160,157.000) -- (340.160,156.344) -- (343.963,156.344) -- (343.963,155.688) -- (349.226,155.688) -- (349.226,155.031) -- (354.616,155.031) -- (354.616,154.375) -- (354.915,154.375) -- (354.915,153.719) -- (360.369,153.719) -- (360.369,153.062) -- (360.468,153.062) -- (360.468,152.406) -- (360.709,152.406) -- (360.709,151.750) -- (361.187,151.750) -- (361.187,151.094) -- (361.325,151.094) -- (361.325,150.438) -- (361.689,150.438) -- (361.689,149.781) -- (362.197,149.781) -- (362.197,149.125) -- (362.867,149.125) -- (362.867,148.469) -- (363.205,148.469) -- (363.205,147.812) -- (363.818,147.812) -- (363.818,147.156) -- (364.571,147.156) -- (364.571,146.500) -- (364.852,146.500) -- (364.852,145.844) -- (365.096,145.844) -- (365.096,145.188) -- (366.516,145.188) -- (366.516,144.531) -- (366.966,144.531) -- (366.966,143.875) -- (367.110,143.875) -- (367.110,143.219) -- (367.306,143.219) -- (367.306,142.562) -- (367.344,142.562) -- (367.344,141.906) -- (367.933,141.906) -- (367.933,141.250) -- (368.144,141.250) -- (368.144,140.594) -- (368.147,140.594) -- (368.147,139.938) -- (368.314,139.938) -- (368.314,139.281) -- (368.403,139.281) -- (368.403,138.625) -- (368.896,138.625) -- (368.896,137.969) -- (368.949,137.969) -- (368.949,137.312) -- (369.046,137.312) -- (369.046,136.656) -- (369.300,136.656) -- (369.300,136.000) -- (369.372,136.000) -- (369.372,135.344) -- (369.513,135.344) -- (369.513,134.688) -- (369.916,134.688) -- (369.916,134.031) -- (369.974,134.031) -- (369.974,133.375) -- (370.251,133.375) -- (370.251,132.719) -- (371.162,132.719) -- (371.162,132.062) -- (371.403,132.062) -- (371.403,131.406) -- (371.830,131.406) -- (371.830,130.750) -- (372.325,130.750) -- (372.325,130.094) -- (372.652,130.094) -- (372.652,129.438) -- (372.724,129.438) -- (372.724,128.781) -- (372.818,128.781) -- (372.818,128.125) -- (373.100,128.125) -- (373.100,127.469) -- (373.333,127.469) -- (373.333,126.812) -- (373.349,126.812) -- (373.349,126.156) -- (373.539,126.156) -- (373.539,125.500) -- (373.633,125.500) -- (373.633,124.844) -- (374.232,124.844) -- (374.232,124.188) -- (374.458,124.188) -- (374.458,123.531) -- (374.545,123.531) -- (374.545,122.875) -- (375.192,122.875) -- (375.192,122.219) -- (375.197,122.219) -- (375.197,121.562) -- (375.431,121.562) -- (375.431,120.906) -- (375.439,120.906) -- (375.439,120.250) -- (375.442,120.250) -- (375.442,119.594) -- (375.472,119.594) -- (375.472,118.938) -- (376.253,118.938) -- (376.253,118.281) -- (376.465,118.281) -- (376.465,117.625) -- (376.701,117.625) -- (376.701,116.969) -- (376.823,116.969) -- (376.823,116.312) -- (376.836,116.312) -- (376.836,115.656) -- (377.418,115.656) -- (377.418,115.000) -- (377.550,115.000) -- (377.550,114.344) -- (377.552,114.344) -- (377.552,113.688) -- (377.641,113.688) -- (377.641,113.031) -- (377.991,113.031) -- (377.991,112.375) -- (378.234,112.375) -- (378.234,111.719) -- (378.466,111.719) -- (378.466,111.062) -- (378.652,111.062) -- (378.652,110.406) -- (378.915,110.406) -- (378.915,109.750) -- (379.026,109.750) -- (379.026,109.094) -- (380.082,109.094) -- (380.082,108.438) -- (380.706,108.438) -- (380.706,107.781) -- (380.865,107.781) -- (380.865,107.125) -- (381.187,107.125) -- (381.187,106.469) -- (381.685,106.469) -- (381.685,105.812) -- (382.075,105.812) -- (382.075,105.156) -- (382.112,105.156) -- (382.112,104.500) -- (383.204,104.500) -- (383.204,103.844) -- (384.526,103.844) -- (384.526,103.188) -- (384.681,103.188) -- (384.681,102.531) -- (384.973,102.531) -- (384.973,101.875) -- (385.195,101.875) -- (385.195,101.219) -- (385.303,101.219) -- (385.303,100.562) -- (385.355,100.562) -- (385.355,99.906) -- (386.534,99.906) -- (386.534,99.250) -- (386.679,99.250) -- (386.679,98.594) -- (386.728,98.594) -- (386.728,97.938) -- (386.750,97.938) -- (386.750,97.281) -- (387.528,97.281) -- (387.528,96.625) -- (387.715,96.625) -- (387.715,95.969) -- (388.165,95.969) -- (388.165,95.312) -- (388.171,95.312) -- (388.171,94.656) -- (388.259,94.656) -- (388.259,94.000) -- (388.678,94.000) -- (388.678,93.344) -- (388.779,93.344) -- (388.779,92.688) -- (389.077,92.688) -- (389.077,92.031) -- (389.586,92.031) -- (389.586,91.375) -- (389.703,91.375) -- (389.703,90.719) -- (390.621,90.719) -- (390.621,90.062) -- (390.707,90.062) -- (390.707,89.406) -- (391.145,89.406) -- (391.145,88.750) -- (391.193,88.750) -- (391.193,88.094) -- (391.483,88.094) -- (391.483,87.438) -- (392.457,87.438) -- (392.457,86.781) -- (392.663,86.781) -- (392.663,86.125) -- (392.823,86.125) -- (392.823,85.469) -- (393.906,85.469) -- (393.906,84.812) -- (395.315,84.812) -- (395.315,84.156) -- (395.351,84.156) -- (395.351,83.500) -- (395.624,83.500) -- (395.624,82.844) -- (395.699,82.844) -- (395.699,82.188) -- (396.607,82.188) -- (396.607,81.531) -- (397.082,81.531) -- (397.082,80.875) -- (398.941,80.875) -- (398.941,80.219) -- (399.024,80.219) -- (399.024,79.562) -- (399.953,79.562) -- (399.953,78.906) -- (401.093,78.906) -- (401.093,78.250) -- (405.247,78.250) -- (405.247,77.594) -- (405.584,77.594) -- (405.584,76.938) -- (406.779,76.938) -- (406.779,76.281) -- (406.839,76.281) -- (406.839,75.625) -- (407.902,75.625) -- (407.902,74.969) -- (407.984,74.969) -- (407.984,74.312) -- (411.211,74.312) -- (411.211,73.656) -- (429.201,73.656) -- (429.201,73.000) -- (432.000,73.000);
\node[anchor=center,inner sep=0pt,outer sep=0pt,text=tg526577,font={\sffamily\fontsize{7bp}{8.4bp}\selectfont}] at (379.000,183.000) {Identical sorted values};
\draw[color=tg0000ff,line width=1.2bp,line cap=round] (93.000,198.000) -- (108.000,198.000);
\node[anchor=west,inner sep=0pt,outer sep=0pt,text=tg526577,font={\sffamily\fontsize{7bp}{8.4bp}\selectfont}] at (113.000,198.000) {Reference};
\draw[color=tgc65c16,line width=0.8bp,line cap=round,dash pattern=on 2bp off 2bp] (183.000,198.000) -- (198.000,198.000);
\node[anchor=west,inner sep=0pt,outer sep=0pt,text=tg526577,font={\sffamily\fontsize{7bp}{8.4bp}\selectfont}] at (203.000,198.000) {Shuffled};
\node[anchor=west,inner sep=0pt,outer sep=0pt,text=tg18314b,font={\sffamily\fontsize{7.5bp}{9bp}\selectfont\bfseries}] at (12.000,214.000) {All 16 horizon--target marginal sample sets are exactly unchanged.};
\draw[color=tge1e7ec,line width=0.6bp,line cap=round] (12.000,225.000) -- (436.000,225.000);
\node[anchor=center,inner sep=0pt,outer sep=0pt,text=tg18314b,font={\sffamily\fontsize{7.7bp}{9.24bp}\selectfont\bfseries}] at (84.500,240.000) {Marginal CRPS (\ensuremath{\times}10\textsuperscript{-16})};
\draw[color=tge1e7ec,line width=0.4bp,line cap=round] (32.000,326.000) -- (137.000,326.000);
\node[anchor=east,inner sep=0pt,outer sep=0pt,text=tg526577,font={\sffamily\fontsize{7bp}{8.4bp}\selectfont}] at (27.000,326.000) {-1};
\draw[color=tge1e7ec,line width=0.4bp,line cap=round] (32.000,291.000) -- (137.000,291.000);
\node[anchor=east,inner sep=0pt,outer sep=0pt,text=tg526577,font={\sffamily\fontsize{7bp}{8.4bp}\selectfont}] at (27.000,291.000) {0};
\draw[color=tge1e7ec,line width=0.4bp,line cap=round] (32.000,256.000) -- (137.000,256.000);
\node[anchor=east,inner sep=0pt,outer sep=0pt,text=tg526577,font={\sffamily\fontsize{7bp}{8.4bp}\selectfont}] at (27.000,256.000) {1};
\draw[color=tg808080,line width=0.5bp,line cap=round] (32.000,256.000) -- (32.000,326.000);
\draw[color=tg808080,line width=0.5bp,line cap=round] (32.000,326.000) -- (137.000,326.000);
\draw[color=tg808080,line width=0.5bp,line cap=round] (32.000,326.000) -- (32.000,328.000);
\node[anchor=center,inner sep=0pt,outer sep=0pt,text=tg526577,font={\sffamily\fontsize{7bp}{8.4bp}\selectfont}] at (32.000,336.000) {0};
\draw[color=tg808080,line width=0.5bp,line cap=round] (84.500,326.000) -- (84.500,328.000);
\node[anchor=center,inner sep=0pt,outer sep=0pt,text=tg526577,font={\sffamily\fontsize{7bp}{8.4bp}\selectfont}] at (84.500,336.000) {0.5};
\draw[color=tg808080,line width=0.5bp,line cap=round] (137.000,326.000) -- (137.000,328.000);
\node[anchor=center,inner sep=0pt,outer sep=0pt,text=tg526577,font={\sffamily\fontsize{7bp}{8.4bp}\selectfont}] at (137.000,336.000) {1};
\draw[color=tg808080,line width=0.6bp,line cap=round,dash pattern=on 2bp off 2bp] (32.000,291.000) -- (137.000,291.000);
\draw[color=tg0000ff,line width=0.85bp,line cap=round] (32.000,291.000) -- (32.000,291.000);
\draw[color=tg0000ff,line width=0.85bp,line cap=round] (30.000,291.000) -- (34.000,291.000);
\draw[color=tg0000ff,line width=0.85bp,line cap=round] (30.000,291.000) -- (34.000,291.000);
\fill[tg0000ff] (32.000,291.000) circle[radius=1.8bp];
\draw[color=tg0000ff,line width=0.85bp,line cap=round] (58.250,308.198) -- (58.250,279.266);
\draw[color=tg0000ff,line width=0.85bp,line cap=round] (56.250,308.198) -- (60.250,308.198);
\draw[color=tg0000ff,line width=0.85bp,line cap=round] (56.250,279.266) -- (60.250,279.266);
\fill[tg0000ff] (58.250,293.732) circle[radius=1.8bp];
\draw[color=tg0000ff,line width=0.85bp,line cap=round] (84.500,306.301) -- (84.500,271.449);
\draw[color=tg0000ff,line width=0.85bp,line cap=round] (82.500,306.301) -- (86.500,306.301);
\draw[color=tg0000ff,line width=0.85bp,line cap=round] (82.500,271.449) -- (86.500,271.449);
\fill[tg0000ff] (84.500,288.875) circle[radius=1.8bp];
\draw[color=tg0000ff,line width=0.85bp,line cap=round] (110.750,308.951) -- (110.750,273.656);
\draw[color=tg0000ff,line width=0.85bp,line cap=round] (108.750,308.951) -- (112.750,308.951);
\draw[color=tg0000ff,line width=0.85bp,line cap=round] (108.750,273.656) -- (112.750,273.656);
\fill[tg0000ff] (110.750,291.304) circle[radius=1.8bp];
\draw[color=tg0000ff,line width=0.85bp,line cap=round] (137.000,312.431) -- (137.000,272.605);
\draw[color=tg0000ff,line width=0.85bp,line cap=round] (135.000,312.431) -- (139.000,312.431);
\draw[color=tg0000ff,line width=0.85bp,line cap=round] (135.000,272.605) -- (139.000,272.605);
\fill[tg0000ff] (137.000,292.518) circle[radius=1.8bp];
\node[anchor=center,inner sep=0pt,outer sep=0pt,text=tg18314b,font={\sffamily\fontsize{7.7bp}{9.24bp}\selectfont\bfseries}] at (231.500,240.000) {Joint energy score};
\draw[color=tge1e7ec,line width=0.4bp,line cap=round] (179.000,321.333) -- (284.000,321.333);
\node[anchor=east,inner sep=0pt,outer sep=0pt,text=tg526577,font={\sffamily\fontsize{7bp}{8.4bp}\selectfont}] at (174.000,321.333) {0};
\draw[color=tge1e7ec,line width=0.4bp,line cap=round] (179.000,292.167) -- (284.000,292.167);
\node[anchor=east,inner sep=0pt,outer sep=0pt,text=tg526577,font={\sffamily\fontsize{7bp}{8.4bp}\selectfont}] at (174.000,292.167) {0.05};
\draw[color=tge1e7ec,line width=0.4bp,line cap=round] (179.000,263.000) -- (284.000,263.000);
\node[anchor=east,inner sep=0pt,outer sep=0pt,text=tg526577,font={\sffamily\fontsize{7bp}{8.4bp}\selectfont}] at (174.000,263.000) {0.1};
\draw[color=tg808080,line width=0.5bp,line cap=round] (179.000,256.000) -- (179.000,326.000);
\draw[color=tg808080,line width=0.5bp,line cap=round] (179.000,326.000) -- (284.000,326.000);
\draw[color=tg808080,line width=0.5bp,line cap=round] (179.000,326.000) -- (179.000,328.000);
\node[anchor=center,inner sep=0pt,outer sep=0pt,text=tg526577,font={\sffamily\fontsize{7bp}{8.4bp}\selectfont}] at (179.000,336.000) {0};
\draw[color=tg808080,line width=0.5bp,line cap=round] (231.500,326.000) -- (231.500,328.000);
\node[anchor=center,inner sep=0pt,outer sep=0pt,text=tg526577,font={\sffamily\fontsize{7bp}{8.4bp}\selectfont}] at (231.500,336.000) {0.5};
\draw[color=tg808080,line width=0.5bp,line cap=round] (284.000,326.000) -- (284.000,328.000);
\node[anchor=center,inner sep=0pt,outer sep=0pt,text=tg526577,font={\sffamily\fontsize{7bp}{8.4bp}\selectfont}] at (284.000,336.000) {1};
\draw[color=tg808080,line width=0.6bp,line cap=round,dash pattern=on 2bp off 2bp] (179.000,321.333) -- (284.000,321.333);
\draw[color=tg0000ff,line width=0.85bp,line cap=round] (179.000,321.333) -- (179.000,321.333);
\draw[color=tg0000ff,line width=0.85bp,line cap=round] (177.000,321.333) -- (181.000,321.333);
\draw[color=tg0000ff,line width=0.85bp,line cap=round] (177.000,321.333) -- (181.000,321.333);
\fill[tg0000ff] (179.000,321.333) circle[radius=1.8bp];
\draw[color=tg0000ff,line width=0.85bp,line cap=round] (205.250,319.180) -- (205.250,305.993);
\draw[color=tg0000ff,line width=0.85bp,line cap=round] (203.250,319.180) -- (207.250,319.180);
\draw[color=tg0000ff,line width=0.85bp,line cap=round] (203.250,305.993) -- (207.250,305.993);
\fill[tg0000ff] (205.250,312.586) circle[radius=1.8bp];
\draw[color=tg0000ff,line width=0.85bp,line cap=round] (231.500,309.164) -- (231.500,288.112);
\draw[color=tg0000ff,line width=0.85bp,line cap=round] (229.500,309.164) -- (233.500,309.164);
\draw[color=tg0000ff,line width=0.85bp,line cap=round] (229.500,288.112) -- (233.500,288.112);
\fill[tg0000ff] (231.500,298.638) circle[radius=1.8bp];
\draw[color=tg0000ff,line width=0.85bp,line cap=round] (257.750,294.733) -- (257.750,268.717);
\draw[color=tg0000ff,line width=0.85bp,line cap=round] (255.750,294.733) -- (259.750,294.733);
\draw[color=tg0000ff,line width=0.85bp,line cap=round] (255.750,268.717) -- (259.750,268.717);
\fill[tg0000ff] (257.750,281.725) circle[radius=1.8bp];
\draw[color=tg0000ff,line width=0.85bp,line cap=round] (284.000,290.204) -- (284.000,262.962);
\draw[color=tg0000ff,line width=0.85bp,line cap=round] (282.000,290.204) -- (286.000,290.204);
\draw[color=tg0000ff,line width=0.85bp,line cap=round] (282.000,262.962) -- (286.000,262.962);
\fill[tg0000ff] (284.000,276.583) circle[radius=1.8bp];
\node[anchor=center,inner sep=0pt,outer sep=0pt,text=tg18314b,font={\sffamily\fontsize{7.7bp}{9.24bp}\selectfont\bfseries}] at (378.500,240.000) {Joint variogram score};
\draw[color=tge1e7ec,line width=0.4bp,line cap=round] (326.000,321.254) -- (431.000,321.254);
\node[anchor=east,inner sep=0pt,outer sep=0pt,text=tg526577,font={\sffamily\fontsize{7bp}{8.4bp}\selectfont}] at (321.000,321.254) {0};
\draw[color=tge1e7ec,line width=0.4bp,line cap=round] (326.000,291.593) -- (431.000,291.593);
\node[anchor=east,inner sep=0pt,outer sep=0pt,text=tg526577,font={\sffamily\fontsize{7bp}{8.4bp}\selectfont}] at (321.000,291.593) {0.025};
\draw[color=tge1e7ec,line width=0.4bp,line cap=round] (326.000,261.932) -- (431.000,261.932);
\node[anchor=east,inner sep=0pt,outer sep=0pt,text=tg526577,font={\sffamily\fontsize{7bp}{8.4bp}\selectfont}] at (321.000,261.932) {0.05};
\draw[color=tg808080,line width=0.5bp,line cap=round] (326.000,256.000) -- (326.000,326.000);
\draw[color=tg808080,line width=0.5bp,line cap=round] (326.000,326.000) -- (431.000,326.000);
\draw[color=tg808080,line width=0.5bp,line cap=round] (326.000,326.000) -- (326.000,328.000);
\node[anchor=center,inner sep=0pt,outer sep=0pt,text=tg526577,font={\sffamily\fontsize{7bp}{8.4bp}\selectfont}] at (326.000,336.000) {0};
\draw[color=tg808080,line width=0.5bp,line cap=round] (378.500,326.000) -- (378.500,328.000);
\node[anchor=center,inner sep=0pt,outer sep=0pt,text=tg526577,font={\sffamily\fontsize{7bp}{8.4bp}\selectfont}] at (378.500,336.000) {0.5};
\draw[color=tg808080,line width=0.5bp,line cap=round] (431.000,326.000) -- (431.000,328.000);
\node[anchor=center,inner sep=0pt,outer sep=0pt,text=tg526577,font={\sffamily\fontsize{7bp}{8.4bp}\selectfont}] at (431.000,336.000) {1};
\draw[color=tg808080,line width=0.6bp,line cap=round,dash pattern=on 2bp off 2bp] (326.000,321.254) -- (431.000,321.254);
\draw[color=tg0000ff,line width=0.85bp,line cap=round] (326.000,321.254) -- (326.000,321.254);
\draw[color=tg0000ff,line width=0.85bp,line cap=round] (324.000,321.254) -- (328.000,321.254);
\draw[color=tg0000ff,line width=0.85bp,line cap=round] (324.000,321.254) -- (328.000,321.254);
\fill[tg0000ff] (326.000,321.254) circle[radius=1.8bp];
\draw[color=tg0000ff,line width=0.85bp,line cap=round] (352.250,319.058) -- (352.250,307.018);
\draw[color=tg0000ff,line width=0.85bp,line cap=round] (350.250,319.058) -- (354.250,319.058);
\draw[color=tg0000ff,line width=0.85bp,line cap=round] (350.250,307.018) -- (354.250,307.018);
\fill[tg0000ff] (352.250,313.038) circle[radius=1.8bp];
\draw[color=tg0000ff,line width=0.85bp,line cap=round] (378.500,309.153) -- (378.500,287.294);
\draw[color=tg0000ff,line width=0.85bp,line cap=round] (376.500,309.153) -- (380.500,309.153);
\draw[color=tg0000ff,line width=0.85bp,line cap=round] (376.500,287.294) -- (380.500,287.294);
\fill[tg0000ff] (378.500,298.224) circle[radius=1.8bp];
\draw[color=tg0000ff,line width=0.85bp,line cap=round] (404.750,299.742) -- (404.750,272.854);
\draw[color=tg0000ff,line width=0.85bp,line cap=round] (402.750,299.742) -- (406.750,299.742);
\draw[color=tg0000ff,line width=0.85bp,line cap=round] (402.750,272.854) -- (406.750,272.854);
\fill[tg0000ff] (404.750,286.298) circle[radius=1.8bp];
\draw[color=tg0000ff,line width=0.85bp,line cap=round] (431.000,293.948) -- (431.000,265.289);
\draw[color=tg0000ff,line width=0.85bp,line cap=round] (429.000,293.948) -- (433.000,293.948);
\draw[color=tg0000ff,line width=0.85bp,line cap=round] (429.000,265.289) -- (433.000,265.289);
\fill[tg0000ff] (431.000,279.618) circle[radius=1.8bp];
\node[anchor=center,inner sep=0pt,outer sep=0pt,text=tg18314b,font={\sffamily\fontsize{7.5bp}{9bp}\selectfont}] at (224.000,352.000) {Shuffle strength (0 = unchanged; 1 = fully shuffled)};
\node[anchor=center,inner sep=0pt,outer sep=0pt,text=tg526577,font={\sffamily\fontsize{7bp}{8.4bp}\selectfont}] at (224.000,368.000) {Score change from clean: paired mean and 95\% t interval over 64 cases; CRPS shows roundoff only.};
\end{tikzpicture}%
}
\endgroup

% END VECTOR FIGURE paper-dependence
\caption{\textbf{A marginally identical forecast can have different joint dependence.}
Top: the first predefined Study 1 case (seed 10000), with 128 exact-reference
samples and their fully coordinate-permuted counterparts. The first-horizon
target-pair clouds change, while the displayed marginal empirical CDFs
coincide; all 16 horizon--target marginal sample sets are exactly unchanged.
Bottom: mutated-minus-clean score means and paired 95\% t intervals over
64 independent cases, at the protocol's fixed permutation strengths. Points
are unconnected and score panels use different scales. CRPS changes are
floating-point roundoff; joint-score changes expose the altered
dependence.}
\label{fig:dependence}
\end{figure}


Figure~\ref{fig:dependence} makes the marginal--joint distinction visible
before its algebraic formulation. All reported scores are losses, so lower is better. Point results include
realized MSE and the avoidable error. Marginal quantiles
use pinball loss. Ensemble CRPS is averaged over individual horizon-target
coordinates. For flattened ensemble paths $x^{(s)}\in\R^K$ and observation
$y$, the empirical energy score is
\begin{equation}
 \widehat{\mathrm{ES}} = \frac1S\sum_s\|x^{(s)}-y\|_2
 -\frac{1}{2S^2}\sum_{s,r}\|x^{(s)}-x^{(r)}\|_2.
\end{equation}
The implementation includes diagonal sample pairs and scores the finite
empirical distribution; it is not the fair-ensemble estimator. Euclidean
norms are not divided by $K$. The variogram score uses $p=1/2$ and equal
weights over distinct unordered coordinates:
\begin{equation}
 \widehat{\mathrm{VS}}_{p} = \binom K2^{-1}
 \sum_{i<j}\left(|y_i-y_j|^p-
       \frac1S\sum_s|x_i^{(s)}-x_j^{(s)}|^p\right)^2,
 \qquad K\geq2.
\end{equation}
The score is undefined at $K=1$. Scores are not compared across dimensions
or normalizations as if they shared one scale. Variogram scores are not
strictly proper over unrestricted joint distributions and cannot detect
all defects; they complement rather than replace other diagnostics.

To construct a controlled dependence defect, choose a separate permutation
$\pi_k$ of $\{1,\ldots,S\}$ for each flattened coordinate $k$ and set
\begin{equation}
 \widetilde{x}^{(s)}_k=x^{(\pi_k(s))}_k.
 \label{eq:permutation}
\end{equation}
Shuffling each coordinate's samples separately reorders, but does not
change, the values at that coordinate, so every per-coordinate statistic,
including the average CRPS, is unchanged; only the pairing between
coordinates changes. Independent random
permutations tend to disrupt the dependence in a correlated ensemble,
although a finite permutation need not eliminate it entirely. This
identity gives marginal diagnostics a negative control that can be run.
If the reference itself factorizes across coordinates, independent
permutations do not create the same population dependence defect.

\section{Controlled evaluation protocol}
\label{sec:protocol}

\subsection{Two studies and their questions}

Two controlled studies evaluate the checks themselves. Study 1 corrupts
exact answer-key forecasts, permutes ensemble
coordinates, runs unmodified baselines, and repairs one preprocessing
fault. Study 2 inserts internal regressions into a fitted pipeline and
evaluates covariance checks under multiplicity. Every fault in both
studies was chosen by hand to exercise an implemented check. Each study's parameters,
seeds, transformations, and tolerances are fixed in its own protocol
before its evaluation run. The studies ask:
\begin{description}
\item[RQ1 (Study 1).] Do the output checks and the covariance check
detect the corruptions, and how often do they alarm on clean
forecasts?
\item[RQ2 (Study 1).] Does removing a known corruption restore the
affected checks without worsening the tested clean behavior?
\item[RQ3 (Study 1).] Do joint scores respond when every marginal sample
distribution is preserved?
\item[RQ4 (Study 2).] Which checks detect internal regressions that
pass the output checks, and does restoring the intended algorithm
lower forecast loss?
\item[RQ5 (Study 2).] How do raw and Holm-adjusted covariance decisions
trade suite-level false alarms against power?
\end{description}

The unit of replication is an independently generated scenario-seed
case. Multiple horizons and targets within one case are not treated as
independent replications. Each case fixes one forecast origin. Where a
model is fitted (the capability sweep and Study 2), it is refitted per
case, a fixed-origin design with recalibration in Tashman's terms; for a
single real series, rolling origins over multiple test periods are more
efficient~\cite{tashman2000out}, and ordinary $K$-fold cross-validation is
valid for purely autoregressive models with uncorrelated
errors~\cite{bergmeir2018note}. Independent synthetic cases replace neither
procedure. Calibration cases and evaluation seeds are disjoint; held-out
seeds provide new stochastic draws, not held-out mechanisms. Reports retain
misses, unsupported cases, exceptions, and clean false alarms. Statistical
uncertainty is computed over independent cases, as recommended for
randomized procedures~\cite{arcuri2011practical}, and each protocol records
its exact convention with the results.

\subsection{Study 1: output corruptions and dependence permutation}

Study 1 uses 64 evaluation seeds (10000--10063), disjoint from development
seeds 0--15. Its fault and dependence experiments use one bivariate AR(1)
task with a context of 120 observations, $H=8$, $D=2$, and $S=128$
samples; the AR(1) coefficient is 0.8 and the innovation correlation is
0.7. Each corruption is applied to the exact reference forecast, one at a
time, and the clean, faulty, and repaired variants use the same generated
case. Answer-key forecasts use information known only to the evaluator;
they test the mechanics of evaluation and
are not counted as learned forecasting achievements. Repair means removing
the specified corruption and rerunning the same checks; there is no
automated repair discovery.

Six output corruptions displace the forecast index, swap target labels,
truncate the horizon, insert a nonfinite point, cross quantiles, or
transpose the sample axes. These are intentionally direct violations of
the output checks; detecting them tests that the checks are implemented
correctly, not that they find subtle bugs. Two further
mutations shrink ensemble variation or permute trajectory coordinates.
They retain well-formed outputs and exercise the reference-covariance
check.

The output checks provide a baseline for malformed outputs. A
conventional realized-loss threshold, calibrated on separate clean cases,
supplies a second comparison. A reported advantage over that threshold
applies only to the specified calibration and test distribution, not to
every possible holdout test. We do not report a head-to-head comparison
with sktime conformance tests or the recent synthetic benchmarks.

The sample-dependence experiment begins with correct conditional paths,
applies Equation~\eqref{eq:permutation}, and compares the resulting
paired marginal and joint scores. Clean and corrupted paths share the
same empirical marginals by construction. The question is the response
of the evaluation measures to the known mutation, rather than the
relative performance of learned probabilistic models.

The unmodified last-value, seasonal-naive (period 12), and fitted Gaussian
AR(1) baselines illustrate capability limitations without intentionally
inserted software defects. A separate preprocessing example omits inverse
scaling in a standardized historical-mean predictor and rechecks a declared
shift-equivariance relation after restoring that step, first on synthetic
white noise and then on one real series.

\subsection{Study 2: implementation regressions and calibrated loss alarms}
\label{sec:implementation-study}

Study 2 moves the injected faults inside a fitted forecasting pipeline. A
six-lag recursive lag-Ridge predictor uses an unpenalized intercept, penalty 1, and standardized targets.
Four branches reverse the prediction-time lag order, retain stale
recursive state, omit inverse target scaling, or silently swap the two
numeric target columns while preserving their labels. Each changes an
implementation operation and returns a well-shaped forecast. The repair
is the clean implementation fitted to the identical context.

Six fixed conditions cover white noise with unequal target locations,
positive and negative AR(1), opposing trends, seasonality, and random
walks with drift. Every condition has 64 evaluation cases and 128 separate
clean calibration cases.

Five checks ask different questions: the output checks; the shift check
with a common shift of 10; the per-series shift check with
shifts $(7,-11)$; agreement with an independently implemented NumPy
lag-Ridge, the second implementation; and a calibrated realized-MSE alarm. The common
shift uses the public diagnostic. Per-series shifts and implementation
agreement are experimental checks. The NumPy implementation independently
constructs lag columns and solves augmented normal equations with an
unpenalized intercept in the original target units. It receives neither
the future nor generator parameters and checks algorithm agreement, not
forecast optimality. Numerical comparisons use fixed absolute tolerance
$10^{-8}$ and relative tolerance $10^{-7}$.

For each condition let $\ell_1,\ldots,\ell_n$, $n=128$, be realized MSEs
from clean calibration cases and let $\ell$ be the evaluation-case MSE.
The split-calibration rank rule is
\begin{equation}
 p_{\mathrm{loss}}=\frac{1+\sum_{i=1}^n\mathbf{1}\{\ell_i\geq\ell\}}{n+1},
 \qquad \text{alarm if }p_{\mathrm{loss}}\leq0.05.
 \label{eq:loss-rank}
\end{equation}
Equivalently, reject above calibration order statistic
$\lceil(n+1)(1-0.05)\rceil$, with conservative ties. For a fixed clean
algorithm under the same case-generating condition, calibration and test
losses are exchangeable, yielding a marginal finite-sample rejection
bound by the split-conformal rank argument~\cite{lei2018conformal}.
This probability averages over calibration and test samples. It does not
guarantee a 5\% rejection probability conditional on the one calibration
sample actually drawn. The reported Wilson intervals condition on that
frozen calibration sample; no common threshold is pooled across conditions.

Each faulty case is paired with the clean run on the same case. We retain
new alarms, lost alarms, misses, and clean-minus-faulty losses, rather
than interpreting an unpaired fault alarm rate as diagnostic improvement.
Loss intervals use case-level paired differences within a condition.
Restoring agreement with the clean algorithm need not improve realized
loss on every noisy future.

\subsection{Study 2: covariance conditions and multiplicity}

Study 2 also evaluates 64 independent four-task suites. Each
suite contains factorized white noise, correlated white noise, negative
AR(1), and a correlated random walk with drift, with $H=6$, $D=2$, $S=64$, and
$B=199$ null ensembles per check. Clean reference samples are paired with
a partial coordinate shuffle and with a 0.75 multiplier on deviations
from the conditional mean. The shuffle permutes half the sample indices
within each coordinate while preserving its complete empirical marginal.

For each variant, the four condition-specific tests form one declared
family. We report individual rejections, any raw rejection in a suite,
and any rejection after Holm's step-down adjustment~\cite{holm1979multiple}.
With sorted valid null $p$-values, Holm compares the $j$th value against
$0.05/(4-j+1)$ and stops at the first non-rejection. Its family-wise
guarantee concerns this fixed family; it does not extend to comparisons
across variants, study sections, or adaptive revisions of the protocol.
The factorized white-noise shuffle is a negative control: independent
coordinates remain independent under these sample permutations, so a
rejection is a false alarm rather than successful defect detection.

\section{Results}
\label{sec:results}

The first four subsections report Study 1: fault detection and false
alarms (RQ1 and RQ2), marginal invariance (RQ3), the capability baselines,
and the inverse-scaling repair. The last two report Study 2: implementation
regressions (RQ4) and covariance checks under multiplicity (RQ5).

% BEGIN GENERATED RESULTS
\IfFileExists{generated-results.tex}{% Generated by experiments/render_results.py; do not edit numerical tables by hand.

\subsection{Study 1: fault detection and false alarms}

Table~\ref{tab:faults} reports alarms for each corruption. The six explicit output violations were detected in every evaluation case. Both well-formed ensemble mutations were rejected by the covariance check in every case.

\begin{table}[htbp]\centering\scriptsize

\caption{Alarm counts and Wilson 95\% intervals (percent). Covariance is not run on malformed outputs; --- denotes an undefined result, not a negative result.}

\label{tab:faults}

\begin{tabular}{@{}lccc@{}}\toprule

Case & Point-MSE alarm & Output contract & Covariance \\ \midrule

Clean control & 5/64 [3.4, 17.0] & 0/64 [0.0, 5.7] & 6/64 [4.4, 19.0] \\

Index shift & 5/64 [3.4, 17.0] & 64/64 [94.3, 100.0] & --- \\

Target-label swap & 5/64 [3.4, 17.0] & 64/64 [94.3, 100.0] & --- \\

Horizon truncation & 64/64 [94.3, 100.0] & 64/64 [94.3, 100.0] & --- \\

Nonfinite point & 64/64 [94.3, 100.0] & 64/64 [94.3, 100.0] & --- \\

Quantile crossing & 5/64 [3.4, 17.0] & 64/64 [94.3, 100.0] & --- \\

Sample-axis transpose & 5/64 [3.4, 17.0] & 64/64 [94.3, 100.0] & --- \\

Variance shrinkage & 5/64 [3.4, 17.0] & 0/64 [0.0, 5.7] & 64/64 [94.3, 100.0] \\

Dependence shuffle & 5/64 [3.4, 17.0] & 0/64 [0.0, 5.7] & 64/64 [94.3, 100.0] \\

\bottomrule\end{tabular}\end{table}

The clean covariance check rejected 6/64 cases (9.4\%; Wilson interval [4.4, 19.0]\%). The observed rate is above the nominal 5\% per-check level, with an interval that includes 5\%. No threshold was changed to remove these false alarms. This run does not establish precise calibration or a family-wise error rate.

The point-MSE cutoff was 4.1853, the empirical 95th percentile from 16 development cases. It alarmed on 5/64 clean cases. It detects invalid point arrays, but ignores index labels, target names, quantiles, and trajectory samples. Where a mutation leaves the point values unchanged, its alarms are identical to the clean alarms and supply no mutation-specific evidence.

Disabling an output mutation restored the exact clean output, including its sample array. Repaired outputs had zero contract violations and the same 6/64 covariance alarms as the clean control. These paired rechecks confirm removal of the corruption; they do not imply that every repaired case receives a statistical pass.

\subsection{Study 1: marginal invariance and joint-score response}

Table~\ref{tab:dependence} varies the proportion of ensemble rows selected for coordinate-wise permutation. Every empirical marginal is preserved at every strength, as checked directly in the runner. The maximum absolute CRPS difference across all cases and strengths was $8.88\times10^{-16}$. Energy and variogram scores both worsened on average at full permutation; their paired intervals exclude zero for this setting. Their raw effect magnitudes have different units and must not be used to rank their sensitivity.

\begin{table}[htbp]\centering\small

\caption{Mutated minus clean score, mean and paired 95\% t interval over 64 independent cases. CRPS differences are floating-point roundoff. Intervals are exploratory and not multiplicity-adjusted.}

\label{tab:dependence}

\begin{tabular}{@{}rccc@{}}\toprule

Strength & CRPS mean & Energy & Variogram \\ \midrule

0.00 & $0$ & +0.0000 [+0.0000, +0.0000] & +0.0000 [+0.0000, +0.0000] \\

0.25 & $-7.81\times10^{-18}$ & +0.0150 [+0.0037, +0.0263] & +0.0069 [+0.0019, +0.0120] \\

0.50 & $6.07\times10^{-18}$ & +0.0389 [+0.0209, +0.0570] & +0.0194 [+0.0102, +0.0286] \\

0.75 & $-8.67\times10^{-19}$ & +0.0679 [+0.0456, +0.0902] & +0.0295 [+0.0181, +0.0408] \\

1.00 & $-4.34\times10^{-18}$ & +0.0767 [+0.0534, +0.1001] & +0.0351 [+0.0230, +0.0472] \\

\bottomrule\end{tabular}\end{table}

\subsection{Study 1: capability profiles}

The univariate capability sweep uses 16 evaluation seeds, three innovation scales, five kinds of series, and three unmodified models. Table~\ref{tab:capability} scores each model against the answer key at the lowest and highest scale: 100 times the unavoidable error over the sum of the unavoidable and avoidable errors (Section~\ref{sec:references}). The full three-scale sweep is in \code{figures/capability-curves.pdf} and \code{capability-summary.csv} of the reproducibility artifact. The fitted AR(1) nails white noise and mean reversion and fails on the trend at both scales. For the zero-drift random walk, the last value is the answer key. Noise matters where a fixed-size pattern competes with it: on the seasonal cycle the fitted AR(1) scores 9 at low noise and 91 at high noise, where the cycle is mostly buried and ignoring it costs little.

\begin{table}[htbp]\centering\small

\caption{Score out of 100 (100 matches the answer key) at low ($\sigma=0.25$) and high ($\sigma=2$) innovation noise, over 16 cases each. N nails it, G good, S struggles, F fails. White noise, mean reversion and the random walk score the same at both levels: with a zero mean and no drift, a larger $\sigma$ only rescales the series. The package's scorecard adds measurement noise to tell those levels apart.}

\label{tab:capability}

\begin{tabular}{@{}lrrrrrr@{}}\toprule
 & \multicolumn{2}{c}{Last value} & \multicolumn{2}{c}{Seasonal naive (12)} & \multicolumn{2}{c}{Fitted AR(1)} \\
\cmidrule(lr){2-3}\cmidrule(lr){4-5}\cmidrule(lr){6-7}
Kind of series & Low & High & Low & High & Low & High \\ \midrule

Noise around a fixed level & 51 S & 51 S & 55 S & 55 S & 99 N & 99 N \\

Straight-line trend & 17 F & 50 S & 4 F & 47 S & 39 F & 18 F \\

Repeating seasonal cycle & 5 F & 46 S & 55 S & 55 S & 9 F & 91 N \\

Mean reversion (AR(1)) & 68 S & 68 S & 40 S & 40 S & 95 N & 95 N \\

Random walk & 100 N & 100 N & 37 F & 37 F & 87 G & 87 G \\

\bottomrule\end{tabular}\end{table}

\subsection{Study 1: an inverse-scaling repair}

A separate predictor standardizes its training context and forecasts the mean in standardized units. The fault omits the inverse transform; the repair applies the training mean and scale. The historical-mean predictor is shift equivariant, so this particular metamorphic relation is a valid declared contract. On white-noise cases with mean 5 and innovation scale 2, the relation failed for 64/64 faulty fits (Wilson 95\% interval [94.3, 100.0]\%) and 0/64 repaired fits ([0.0, 5.7]\%). The paired repaired-minus-faulty MAE difference was -3.4711 [-3.6333, -3.3090] (95\% t interval).

We then apply the same known omission and repair to AirPassengers, the monthly 1949--1960 international-airline passenger series in thousands~\cite{rAirPassengers}. All scaling statistics are fitted on each training prefix. The four ordered origins each forecast the next 12 months; their results are reported separately because they are not independent series replicates.

\begin{table}[htbp]\centering\small

\caption{AirPassengers inverse-scaling illustration. MAE is in thousands of passengers. The declared shift relation fails before and passes after repair at all four origins.}

\label{tab:real}

\begin{tabular}{@{}lrrr@{}}\toprule

Last training month & Faulty MAE & Repaired MAE & Difference \\ \midrule

1956-12 & 368.42 & 154.71 & -213.71 \\

1957-12 & 381.00 & 150.10 & -230.90 \\

1958-12 & 428.33 & 182.43 & -245.91 \\

1959-12 & 476.17 & 213.67 & -262.49 \\

\bottomrule\end{tabular}\end{table}

Repair improves MAE at every displayed origin, but the repaired historical mean remains a poor model of this growing seasonal series.
}{%
\PackageError{timegym}{Missing generated-results.tex}{Run experiments.render_results before compiling.}}
% END GENERATED RESULTS

\begin{figure}[!t]
\centering
% BEGIN VECTOR FIGURE paper-diagnostic-evidence
% Generated vector scene; SVG and PNG use the same plotted geometry.
\begingroup
\definecolor{tg18314b}{HTML}{18314b}
\definecolor{tg19324b}{HTML}{19324b}
\definecolor{tg19324c}{HTML}{19324c}
\definecolor{tg626b73}{HTML}{626b73}
\definecolor{tg808080}{HTML}{808080}
\definecolor{tga2aeb9}{HTML}{a2aeb9}
\definecolor{tga4b0bb}{HTML}{a4b0bb}
\definecolor{tgc2cad2}{HTML}{c2cad2}
\definecolor{tgc4ccd4}{HTML}{c4ccd4}
\definecolor{tgc65c16}{HTML}{c65c16}
\definecolor{tge1e4e7}{HTML}{e1e4e7}
\definecolor{tge6ebef}{HTML}{e6ebef}
\definecolor{tgeaeff3}{HTML}{eaeff3}
\definecolor{tgf1f5f8}{HTML}{f1f5f8}
\definecolor{tgffffff}{HTML}{ffffff}
\resizebox{\linewidth}{!}{%
\begin{tikzpicture}[x=1bp,y=-1bp]
\path[use as bounding box] (0,0) rectangle (448,328);
\node[anchor=west,inner sep=0pt,outer sep=0pt,text=tg18314b,font={\sffamily\fontsize{10bp}{12bp}\selectfont\bfseries}] at (12.000,12.000) {Different checks expose different implementation faults};
\node[anchor=west,inner sep=0pt,outer sep=0pt,text=tg626b73,font={\sffamily\fontsize{7.5bp}{9bp}\selectfont}] at (12.000,28.000) {Alarms / 384 cases per row: 6 fixed conditions, 64 independent cases per condition.};
\node[anchor=center,inner sep=0pt,outer sep=0pt,text=tg18314b,font={\sffamily\fontsize{7.5bp}{9bp}\selectfont\bfseries}] at (172.000,43.500) {Output};
\node[anchor=center,inner sep=0pt,outer sep=0pt,text=tg18314b,font={\sffamily\fontsize{7.5bp}{9bp}\selectfont\bfseries}] at (172.000,52.500) {checks};
\node[anchor=center,inner sep=0pt,outer sep=0pt,text=tg18314b,font={\sffamily\fontsize{7.5bp}{9bp}\selectfont\bfseries}] at (220.000,43.500) {Scalar};
\node[anchor=center,inner sep=0pt,outer sep=0pt,text=tg18314b,font={\sffamily\fontsize{7.5bp}{9bp}\selectfont\bfseries}] at (220.000,52.500) {shift};
\node[anchor=center,inner sep=0pt,outer sep=0pt,text=tg18314b,font={\sffamily\fontsize{7.5bp}{9bp}\selectfont\bfseries}] at (268.000,43.500) {Target};
\node[anchor=center,inner sep=0pt,outer sep=0pt,text=tg18314b,font={\sffamily\fontsize{7.5bp}{9bp}\selectfont\bfseries}] at (268.000,52.500) {shift*};
\node[anchor=center,inner sep=0pt,outer sep=0pt,text=tg18314b,font={\sffamily\fontsize{7.5bp}{9bp}\selectfont\bfseries}] at (316.000,43.500) {NumPy};
\node[anchor=center,inner sep=0pt,outer sep=0pt,text=tg18314b,font={\sffamily\fontsize{7.5bp}{9bp}\selectfont\bfseries}] at (316.000,52.500) {agreement*};
\node[anchor=center,inner sep=0pt,outer sep=0pt,text=tg18314b,font={\sffamily\fontsize{7.5bp}{9bp}\selectfont\bfseries}] at (364.000,43.500) {Calibrated};
\node[anchor=center,inner sep=0pt,outer sep=0pt,text=tg18314b,font={\sffamily\fontsize{7.5bp}{9bp}\selectfont\bfseries}] at (364.000,52.500) {MSE*};
\node[anchor=center,inner sep=0pt,outer sep=0pt,text=tg18314b,font={\sffamily\fontsize{7.5bp}{9bp}\selectfont\bfseries}] at (412.000,43.500) {Answer};
\node[anchor=center,inner sep=0pt,outer sep=0pt,text=tg18314b,font={\sffamily\fontsize{7.5bp}{9bp}\selectfont\bfseries}] at (412.000,52.500) {key*†};
\node[anchor=east,inner sep=0pt,outer sep=0pt,text=tg18314b,font={\sffamily\fontsize{7.4bp}{8.88bp}\selectfont\bfseries}] at (139.000,77.000) {Correct implementation};
\path[fill=tgf1f5f8,draw=none,line width=0.5bp] (148.700,65.700) rectangle (195.300,88.300);
\node[anchor=center,inner sep=0pt,outer sep=0pt,text=tg18314b,font={\sffamily\fontsize{8bp}{9.6bp}\selectfont}] at (172.000,77.000) {0/384};
\path[fill=tgf1f5f8,draw=none,line width=0.5bp] (196.700,65.700) rectangle (243.300,88.300);
\node[anchor=center,inner sep=0pt,outer sep=0pt,text=tg18314b,font={\sffamily\fontsize{8bp}{9.6bp}\selectfont}] at (220.000,77.000) {0/384};
\path[fill=tgf1f5f8,draw=none,line width=0.5bp] (244.700,65.700) rectangle (291.300,88.300);
\node[anchor=center,inner sep=0pt,outer sep=0pt,text=tg18314b,font={\sffamily\fontsize{8bp}{9.6bp}\selectfont}] at (268.000,77.000) {0/384};
\path[fill=tgf1f5f8,draw=none,line width=0.5bp] (292.700,65.700) rectangle (339.300,88.300);
\node[anchor=center,inner sep=0pt,outer sep=0pt,text=tg18314b,font={\sffamily\fontsize{8bp}{9.6bp}\selectfont}] at (316.000,77.000) {0/384};
\path[fill=tge6ebef,draw=none,line width=0.5bp] (340.700,65.700) rectangle (387.300,88.300);
\node[anchor=center,inner sep=0pt,outer sep=0pt,text=tg18314b,font={\sffamily\fontsize{8bp}{9.6bp}\selectfont\bfseries}] at (364.000,77.000) {20/384};
\path[fill=tgeaeff3,draw=none,line width=0.5bp] (388.700,65.700) rectangle (435.300,88.300);
\node[anchor=center,inner sep=0pt,outer sep=0pt,text=tg18314b,font={\sffamily\fontsize{8bp}{9.6bp}\selectfont\bfseries}] at (412.000,77.000) {12/384};
\node[anchor=east,inner sep=0pt,outer sep=0pt,text=tg18314b,font={\sffamily\fontsize{7.4bp}{8.88bp}\selectfont}] at (139.000,101.000) {Reversed prediction lags};
\path[fill=tgf1f5f8,draw=none,line width=0.5bp] (148.700,89.700) rectangle (195.300,112.300);
\node[anchor=center,inner sep=0pt,outer sep=0pt,text=tg18314b,font={\sffamily\fontsize{8bp}{9.6bp}\selectfont}] at (172.000,101.000) {0/384};
\path[fill=tgf1f5f8,draw=none,line width=0.5bp] (196.700,89.700) rectangle (243.300,112.300);
\node[anchor=center,inner sep=0pt,outer sep=0pt,text=tg18314b,font={\sffamily\fontsize{8bp}{9.6bp}\selectfont}] at (220.000,101.000) {0/384};
\path[fill=tgf1f5f8,draw=none,line width=0.5bp] (244.700,89.700) rectangle (291.300,112.300);
\node[anchor=center,inner sep=0pt,outer sep=0pt,text=tg18314b,font={\sffamily\fontsize{8bp}{9.6bp}\selectfont}] at (268.000,101.000) {0/384};
\path[fill=tg18314b,draw=none,line width=0.5bp] (292.700,89.700) rectangle (339.300,112.300);
\node[anchor=center,inner sep=0pt,outer sep=0pt,text=tgffffff,font={\sffamily\fontsize{8bp}{9.6bp}\selectfont\bfseries}] at (316.000,101.000) {384/384};
\path[fill=tgc2cad2,draw=none,line width=0.5bp] (340.700,89.700) rectangle (387.300,112.300);
\node[anchor=center,inner sep=0pt,outer sep=0pt,text=tg18314b,font={\sffamily\fontsize{8bp}{9.6bp}\selectfont\bfseries}] at (364.000,101.000) {84/384};
\path[fill=tga4b0bb,draw=none,line width=0.5bp] (388.700,89.700) rectangle (435.300,112.300);
\node[anchor=center,inner sep=0pt,outer sep=0pt,text=tg18314b,font={\sffamily\fontsize{8bp}{9.6bp}\selectfont\bfseries}] at (412.000,101.000) {136/384};
\node[anchor=east,inner sep=0pt,outer sep=0pt,text=tg18314b,font={\sffamily\fontsize{7.4bp}{8.88bp}\selectfont}] at (139.000,125.000) {Stale recursive state};
\path[fill=tgf1f5f8,draw=none,line width=0.5bp] (148.700,113.700) rectangle (195.300,136.300);
\node[anchor=center,inner sep=0pt,outer sep=0pt,text=tg18314b,font={\sffamily\fontsize{8bp}{9.6bp}\selectfont}] at (172.000,125.000) {0/384};
\path[fill=tgf1f5f8,draw=none,line width=0.5bp] (196.700,113.700) rectangle (243.300,136.300);
\node[anchor=center,inner sep=0pt,outer sep=0pt,text=tg18314b,font={\sffamily\fontsize{8bp}{9.6bp}\selectfont}] at (220.000,125.000) {0/384};
\path[fill=tgf1f5f8,draw=none,line width=0.5bp] (244.700,113.700) rectangle (291.300,136.300);
\node[anchor=center,inner sep=0pt,outer sep=0pt,text=tg18314b,font={\sffamily\fontsize{8bp}{9.6bp}\selectfont}] at (268.000,125.000) {0/384};
\path[fill=tg18314b,draw=none,line width=0.5bp] (292.700,113.700) rectangle (339.300,136.300);
\node[anchor=center,inner sep=0pt,outer sep=0pt,text=tgffffff,font={\sffamily\fontsize{8bp}{9.6bp}\selectfont\bfseries}] at (316.000,125.000) {384/384};
\path[fill=tgc4ccd4,draw=none,line width=0.5bp] (340.700,113.700) rectangle (387.300,136.300);
\node[anchor=center,inner sep=0pt,outer sep=0pt,text=tg18314b,font={\sffamily\fontsize{8bp}{9.6bp}\selectfont\bfseries}] at (364.000,125.000) {80/384};
\path[fill=tga2aeb9,draw=none,line width=0.5bp] (388.700,113.700) rectangle (435.300,136.300);
\node[anchor=center,inner sep=0pt,outer sep=0pt,text=tg18314b,font={\sffamily\fontsize{8bp}{9.6bp}\selectfont\bfseries}] at (412.000,125.000) {140/384};
\node[anchor=east,inner sep=0pt,outer sep=0pt,text=tg18314b,font={\sffamily\fontsize{7.4bp}{8.88bp}\selectfont}] at (139.000,149.000) {Missing inverse scaling};
\path[fill=tgf1f5f8,draw=none,line width=0.5bp] (148.700,137.700) rectangle (195.300,160.300);
\node[anchor=center,inner sep=0pt,outer sep=0pt,text=tg18314b,font={\sffamily\fontsize{8bp}{9.6bp}\selectfont}] at (172.000,149.000) {0/384};
\path[fill=tg18314b,draw=none,line width=0.5bp] (196.700,137.700) rectangle (243.300,160.300);
\node[anchor=center,inner sep=0pt,outer sep=0pt,text=tgffffff,font={\sffamily\fontsize{8bp}{9.6bp}\selectfont\bfseries}] at (220.000,149.000) {384/384};
\path[fill=tg18314b,draw=none,line width=0.5bp] (244.700,137.700) rectangle (291.300,160.300);
\node[anchor=center,inner sep=0pt,outer sep=0pt,text=tgffffff,font={\sffamily\fontsize{8bp}{9.6bp}\selectfont\bfseries}] at (268.000,149.000) {384/384};
\path[fill=tg18314b,draw=none,line width=0.5bp] (292.700,137.700) rectangle (339.300,160.300);
\node[anchor=center,inner sep=0pt,outer sep=0pt,text=tgffffff,font={\sffamily\fontsize{8bp}{9.6bp}\selectfont\bfseries}] at (316.000,149.000) {384/384};
\path[fill=tg18314b,draw=none,line width=0.5bp] (340.700,137.700) rectangle (387.300,160.300);
\node[anchor=center,inner sep=0pt,outer sep=0pt,text=tgffffff,font={\sffamily\fontsize{8bp}{9.6bp}\selectfont\bfseries}] at (364.000,149.000) {384/384};
\path[fill=tg18314b,draw=none,line width=0.5bp] (388.700,137.700) rectangle (435.300,160.300);
\node[anchor=center,inner sep=0pt,outer sep=0pt,text=tgffffff,font={\sffamily\fontsize{8bp}{9.6bp}\selectfont\bfseries}] at (412.000,149.000) {384/384};
\node[anchor=east,inner sep=0pt,outer sep=0pt,text=tg18314b,font={\sffamily\fontsize{7.4bp}{8.88bp}\selectfont}] at (139.000,173.000) {Swapped recursive targets};
\path[fill=tgf1f5f8,draw=none,line width=0.5bp] (148.700,161.700) rectangle (195.300,184.300);
\node[anchor=center,inner sep=0pt,outer sep=0pt,text=tg18314b,font={\sffamily\fontsize{8bp}{9.6bp}\selectfont}] at (172.000,173.000) {0/384};
\path[fill=tgf1f5f8,draw=none,line width=0.5bp] (196.700,161.700) rectangle (243.300,184.300);
\node[anchor=center,inner sep=0pt,outer sep=0pt,text=tg18314b,font={\sffamily\fontsize{8bp}{9.6bp}\selectfont}] at (220.000,173.000) {0/384};
\path[fill=tg18314b,draw=none,line width=0.5bp] (244.700,161.700) rectangle (291.300,184.300);
\node[anchor=center,inner sep=0pt,outer sep=0pt,text=tgffffff,font={\sffamily\fontsize{8bp}{9.6bp}\selectfont\bfseries}] at (268.000,173.000) {384/384};
\path[fill=tg18314b,draw=none,line width=0.5bp] (292.700,161.700) rectangle (339.300,184.300);
\node[anchor=center,inner sep=0pt,outer sep=0pt,text=tgffffff,font={\sffamily\fontsize{8bp}{9.6bp}\selectfont\bfseries}] at (316.000,173.000) {384/384};
\path[fill=tg19324c,draw=none,line width=0.5bp] (340.700,161.700) rectangle (387.300,184.300);
\node[anchor=center,inner sep=0pt,outer sep=0pt,text=tgffffff,font={\sffamily\fontsize{8bp}{9.6bp}\selectfont\bfseries}] at (364.000,173.000) {382/384};
\path[fill=tg19324b,draw=none,line width=0.5bp] (388.700,161.700) rectangle (435.300,184.300);
\node[anchor=center,inner sep=0pt,outer sep=0pt,text=tgffffff,font={\sffamily\fontsize{8bp}{9.6bp}\selectfont\bfseries}] at (412.000,173.000) {383/384};
\draw[color=tg808080,line width=0.65bp,line cap=round] (12.000,89.000) -- (436.000,89.000);
\node[anchor=west,inner sep=0pt,outer sep=0pt,text=tg626b73,font={\sffamily\fontsize{7bp}{8.4bp}\selectfont}] at (12.000,199.000) {* Checks used only in this experiment. † Exploratory addition, chosen after the results were known.};
\node[anchor=west,inner sep=0pt,outer sep=0pt,text=tg626b73,font={\sffamily\fontsize{7bp}{8.4bp}\selectfont}] at (12.000,213.000) {Dark cells mean more alarms. The clean row retains false alarms; the output checks accept all four faults.};
\draw[color=tge1e4e7,line width=0.6bp,line cap=round] (12.000,225.000) -- (436.000,225.000);
\node[anchor=west,inner sep=0pt,outer sep=0pt,text=tg18314b,font={\sffamily\fontsize{8.5bp}{10.2bp}\selectfont\bfseries}] at (12.000,240.000) {Restoring the intended algorithm can still increase forecast error};
\node[anchor=west,inner sep=0pt,outer sep=0pt,text=tg626b73,font={\sffamily\fontsize{7.5bp}{9bp}\selectfont}] at (12.000,255.000) {Random walk with drift / stale recursive state};
\node[anchor=east,inner sep=0pt,outer sep=0pt,text=tgc65c16,font={\sffamily\fontsize{7.5bp}{9bp}\selectfont}] at (436.000,255.000) {Only 21/64 forecasts improve};
\draw[color=tg808080,line width=0.5bp,line cap=round] (40.000,283.000) -- (408.000,283.000);
\draw[color=tg808080,line width=0.65bp,line cap=round,dash pattern=on 2bp off 2bp] (224.000,268.000) -- (224.000,292.000);
\draw[color=tg808080,line width=0.5bp,line cap=round] (40.000,283.000) -- (40.000,286.000);
\node[anchor=center,inner sep=0pt,outer sep=0pt,text=tg626b73,font={\sffamily\fontsize{7bp}{8.4bp}\selectfont}] at (40.000,297.000) {-1};
\draw[color=tg808080,line width=0.5bp,line cap=round] (132.000,283.000) -- (132.000,286.000);
\node[anchor=center,inner sep=0pt,outer sep=0pt,text=tg626b73,font={\sffamily\fontsize{7bp}{8.4bp}\selectfont}] at (132.000,297.000) {-0.5};
\draw[color=tg808080,line width=0.5bp,line cap=round] (224.000,283.000) -- (224.000,286.000);
\node[anchor=center,inner sep=0pt,outer sep=0pt,text=tg626b73,font={\sffamily\fontsize{7bp}{8.4bp}\selectfont}] at (224.000,297.000) {0};
\draw[color=tg808080,line width=0.5bp,line cap=round] (316.000,283.000) -- (316.000,286.000);
\node[anchor=center,inner sep=0pt,outer sep=0pt,text=tg626b73,font={\sffamily\fontsize{7bp}{8.4bp}\selectfont}] at (316.000,297.000) {+0.5};
\draw[color=tg808080,line width=0.5bp,line cap=round] (408.000,283.000) -- (408.000,286.000);
\node[anchor=center,inner sep=0pt,outer sep=0pt,text=tg626b73,font={\sffamily\fontsize{7bp}{8.4bp}\selectfont}] at (408.000,297.000) {+1};
\draw[color=tgc65c16,line width=1.4bp,line cap=round] (263.048,283.000) -- (373.793,283.000);
\draw[color=tgc65c16,line width=1.1bp,line cap=round] (263.048,280.000) -- (263.048,286.000);
\draw[color=tgc65c16,line width=1.1bp,line cap=round] (373.793,280.000) -- (373.793,286.000);
\path[fill=tgc65c16,draw=none] (318.421,283.000) circle[radius=2.5bp];
\node[anchor=center,inner sep=0pt,outer sep=0pt,text=tgc65c16,font={\sffamily\fontsize{7.5bp}{9bp}\selectfont}] at (318.421,270.000) {+0.513 [0.212, 0.814]};
\node[anchor=center,inner sep=0pt,outer sep=0pt,text=tg18314b,font={\sffamily\fontsize{7.5bp}{9bp}\selectfont}] at (224.000,317.000) {Realized MSE: clean - faulty (paired mean and unadjusted 95\% t interval)};
\end{tikzpicture}%
}
\endgroup

% END VECTOR FIGURE paper-diagnostic-evidence
\caption{\textbf{Diagnostic blind spots and the limits of repair.}
Top: alarm counts from four internal lag-Ridge regressions, using the same
six conditions and \CasesPerCondition{} cases per condition
(\CasesPerRow{} per row). Darker cells mean
more alarms, including false alarms in the correct-implementation row.
All four regressions pass the output checks. Asterisks mark experiment-only
checks.
Bottom: correcting stale recursion on the random-walk condition with drift restores
agreement with the declared algorithm, yet increases average realized MSE
by \RepairDeltaMean{} (paired 95\% t interval
\RepairDeltaLow{}--\RepairDeltaHigh{}); only \RepairImproved{} of
\CasesPerCondition{} cases improve.
Agreement with the intended implementation and forecasting accuracy are
separate properties. Table~\ref{tab:supplement-complete} gives the full
condition-specific alarm counts; accompanying records retain the paired losses.}
\label{fig:diagnostic-evidence}
\end{figure}

% BEGIN GENERATED SUPPLEMENT
\IfFileExists{generated-supplement.tex}{% Generated from experiments/results-v2; do not edit numerical statements manually.

\subsection{Study 2: implementation regressions}

Study 2 evaluates 4 implementation regressions under 6 fixed conditions, with 128 independent clean calibration cases and 64 disjoint evaluation cases each (768 calibration and 384 evaluation cases in total). The ridge penalty is 1 and the lag window is 6. Figure~\ref{fig:diagnostic-evidence} gives the pooled alarm counts, and Table~\ref{tab:supplement-complete} gives them for each condition. The paired comparison is clean-minus-faulty, with no separate repaired arm. Any output-contract violation terminates the study.

Across the fixed condition--fault pairs, the calibrated loss alarm changed from no alarm to alarm in 861 paired comparisons and from alarm to no alarm in 11. The same clean case appears once per fault in these descriptive paired counts; they are not independent replicates of a pooled statistical test.

The clean implementation reduced mean realized MSE in 22/24 condition--fault pairs; 18/24 paired 95\% t intervals were wholly below zero. The full paired loss differences, intervals, and numbers of individually improved cases are retained in \code{model-repair-losses.csv}. These exploratory intervals are not adjusted across comparisons.

For random walk with drift with stale state, clean-minus-faulty mean MSE was +0.5132 [+0.2122, +0.8141] (paired 95\% t interval); only 21/64 realized forecasts improved. Restoring the declared algorithm can therefore worsen forecast accuracy on a particular condition.

\begin{table}[p]\centering\scriptsize

\caption{Complete condition-specific implementation alarm counts. Each entry is out of 64 paired cases. C, S, T, N, and L denote contract, scalar shift, target shift, NumPy agreement, and calibrated loss respectively. Clean rows are negative controls; there is no separate corrected-refit row.}

\label{tab:supplement-complete}

\begin{tabular}{@{}llrrrrr@{}}\toprule

Condition & Branch & C & S & T & N & L \\ \midrule

White noise & Clean & 0 & 0 & 0 & 0 & 1 \\

 & Reversed lags & 0 & 0 & 0 & 64 & 2 \\

 & Stale state & 0 & 0 & 0 & 64 & 2 \\

 & No inverse scale & 0 & 64 & 64 & 64 & 64 \\

 & Target swap & 0 & 0 & 64 & 64 & 64 \\

\addlinespace

Positive AR(1) & Clean & 0 & 0 & 0 & 0 & 3 \\

 & Reversed lags & 0 & 0 & 0 & 64 & 5 \\

 & Stale state & 0 & 0 & 0 & 64 & 5 \\

 & No inverse scale & 0 & 64 & 64 & 64 & 64 \\

 & Target swap & 0 & 0 & 64 & 64 & 64 \\

\addlinespace

Negative AR(1) & Clean & 0 & 0 & 0 & 0 & 1 \\

 & Reversed lags & 0 & 0 & 0 & 64 & 6 \\

 & Stale state & 0 & 0 & 0 & 64 & 2 \\

 & No inverse scale & 0 & 64 & 64 & 64 & 64 \\

 & Target swap & 0 & 0 & 64 & 64 & 64 \\

\addlinespace

Opposing trends & Clean & 0 & 0 & 0 & 0 & 3 \\

 & Reversed lags & 0 & 0 & 0 & 64 & 2 \\

 & Stale state & 0 & 0 & 0 & 64 & 6 \\

 & No inverse scale & 0 & 64 & 64 & 64 & 64 \\

 & Target swap & 0 & 0 & 64 & 64 & 64 \\

\addlinespace

Seasonal scales & Clean & 0 & 0 & 0 & 0 & 7 \\

 & Reversed lags & 0 & 0 & 0 & 64 & 64 \\

 & Stale state & 0 & 0 & 0 & 64 & 64 \\

 & No inverse scale & 0 & 64 & 64 & 64 & 64 \\

 & Target swap & 0 & 0 & 64 & 64 & 64 \\

\addlinespace

Random walk with drift & Clean & 0 & 0 & 0 & 0 & 5 \\

 & Reversed lags & 0 & 0 & 0 & 64 & 5 \\

 & Stale state & 0 & 0 & 0 & 64 & 1 \\

 & No inverse scale & 0 & 64 & 64 & 64 & 64 \\

 & Target swap & 0 & 0 & 64 & 64 & 62 \\

\addlinespace

\bottomrule\end{tabular}\end{table}

\subsection{Study 2: covariance checks under multiplicity}

For clean suites, any false alarm occurred in 14/64 raw suites versus 3/64 after Holm correction. For reduced-spread suites, detection of at least one altered condition fell from 63/64 to 30/64. The correction thus reduces family-wise false alarms at a substantial power cost in this setting.

\begin{table}[htbp]\centering\small

\caption{Any false alarm per four-check suite, counts and Wilson 95\% intervals (percent). Clean suites have four true nulls, partial-shuffle suites one, and reduced-spread suites none. --- denotes an undefined result, not a negative result.}

\label{tab:supplement-familywise}

\begin{tabular}{@{}lrr@{}}\toprule

Variant & Raw & Holm \\ \midrule

Exact-reference control & 14/64 [13.5, 33.4] & 3/64 [1.6, 12.9] \\

Partial dependence shuffle & 4/64 [2.5, 15.0] & 4/64 [2.5, 15.0] \\

Reduced spread & --- & --- \\

\bottomrule\end{tabular}\end{table}

\begin{table}[htbp]\centering\small

\caption{All condition-specific covariance rejections, raw and Holm, out of 64 suites. Clean columns and the first shuffled row are negative controls; other entries measure power on these alternatives.}

\label{tab:supplement-covariance}

\begin{tabular}{@{}p{0.34\linewidth}rrrrrr@{}}\toprule

 & \multicolumn{2}{c}{Clean} & \multicolumn{2}{c}{Partial shuffle} & \multicolumn{2}{c}{Reduced spread} \\

Condition & Raw & Holm & Raw & Holm & Raw & Holm \\ \midrule

White noise: independent targets & 3 & 0 & 4 & 4 & 24 & 6 \\

White noise: correlated targets & 5 & 2 & 59 & 51 & 46 & 13 \\

Negative AR(1) & 6 & 1 & 64 & 63 & 42 & 17 \\

Correlated random walk with drift & 2 & 0 & 64 & 64 & 53 & 21 \\

\bottomrule\end{tabular}\end{table}

Coordinate-wise sorted sample arrays were identical before and after every partial permutation. The largest absolute marginal CRPS change was $8.88\times10^{-16}$. Paired energy and variogram score differences are retained for every condition.
}{%
\PackageError{timegym}{Missing generated-supplement.tex}{Run experiments.render_supplement before compiling.}}
% END GENERATED SUPPLEMENT

\section{Related work and positioning}
\label{sec:related}

\subsection{The oracle problem}

Testing a forecasting pipeline is an instance of the test oracle problem.
Weyuker called a program non-testable when no oracle exists or the oracle
is too costly to use, and listed the remedies that remain available: a
partial oracle, tests on simplified data with known results, and known
properties of the output~\cite{weyuker1982testing}. Davis and Weyuker
had already proposed the pseudo-oracle, an independently produced implementation of the
same specification whose outputs are compared with those of the program
under test, within a tolerance for numerical programs~\cite{davis1981pseudo}.
Guderlei and Mayer extended metamorphic relations to programs with random
output by comparing the output distributions of related executions with
hypothesis tests~\cite{guderlei2007statistical}. Oracle issues specific to
machine learning were reviewed in~\cite{seca2021review}. Table~\ref{tab:oracles}
maps the evidence used in this paper onto these oracle types.

Pham et al.'s CRADLE detects bugs in deep-learning libraries by comparing
the outputs of different backends~\cite{pham2019cradle}. Tolerance-based
assertions are common in those libraries' test suites and a recurring
source of flaky tests~\cite{nejadgholi2019study,dutta2020detecting}; the
explicit \code{INCONCLUSIVE} outcome keeps a tolerance choice from being
reported as a verdict. Shift checks are properties in the sense of
property-based testing, stated beside the code and checked on generated
inputs~\cite{claessen2000quickcheck}; here the scenario generators supply
those inputs.

\begin{table}[!htb]
\centering\small
\begin{tabular}{@{}p{0.26\linewidth}p{0.40\linewidth}p{0.26\linewidth}@{}}
\toprule
Oracle type & TimeGym mechanism & Source \\
\midrule
Partial oracle & Output checks on shape, dates, finiteness, and quantile order & \cite{weyuker1982testing} \\[0.4em]
Simplified data with known results & Synthetic scenarios with answer keys & \cite{weyuker1982testing} \\[0.4em]
Known properties & Shift checks & \cite{weyuker1982testing,claessen2000quickcheck} \\[0.4em]
Pseudo-oracle & Independent NumPy lag-Ridge implementation (experimental) & \cite{davis1981pseudo,pham2019cradle} \\[0.4em]
Statistical test oracle & Reference-covariance rank check; calibrated loss alarm & \cite{guderlei2007statistical} \\
\bottomrule
\end{tabular}
\caption{The evidence in this paper by oracle type.}
\label{tab:oracles}
\end{table}

\subsection{Controlled tasks, conformance, and metamorphic testing}

Controlled forecasting tasks, software conformance, and metamorphic
testing are the three neighbouring lines of work
(Table~\ref{tab:positioning}). SynTSBench connects configurable temporal
features to model capability, irregularity robustness, and theoretical
performance boundaries~\cite{tan2025syntsbench}: it derives a theoretical
optimum for each synthetic task and reports it beside the models' raw
errors. SimpleTimeBench uses elementary patterns and leading covariates
to diagnose time-series foundation
models~\cite{ghoroghchian2026simpletimebench}; it describes the ideal
forecast for each process and scores models against naive baselines.

For the financial scenarios in the catalog
(Appendix~\ref{app:extended}), the direct prior work is
FinStressTS~\cite{sun2026finstressts}, a synthetic benchmark published in
2026 that compares 15 forecasting models on 30 environments built from
six mechanism families, including conditional volatility, heavy tails,
regime switching, and jumps. It is prior work for those scenarios only:
TimeGym dates from 2021~\cite{seca2021timegym} and has a different
purpose. FinStressTS is a benchmark for comparing models on financial
mechanisms. TimeGym is a general tool for testing one forecasting
pipeline, with output checks, injected faults, and repair rechecks, and
financial scenarios are one part of its catalog. Real financial evaluation
also needs target-specific protocols; FinTSB includes market constraints
and transaction costs~\cite{hu2026fintsb}.

Dwarakanath et al. specified 19 metamorphic relations for a forecasting
application, reported eight previously unknown issues, and evaluated
mutation detection in a reference LSTM~\cite{dwarakanath2019metamorphic}.
This is the direct methodological predecessor of the shift checks, and it
includes evidence on actual application defects that our studies of
hand-picked faults do not provide. TimeGym adds answer keys and separate
checks for marginal and joint forecasts to that testing workflow; it does
not introduce forecasting metamorphic testing, and it implements far
fewer relations.

Framework conformance and evaluation practice also matter. The sktime
\code{check\_estimator} utility runs tests appropriate to an estimator's
unified interface~\cite{sktimeConformance}. It exercises that interface
on unstructured noise series and does not judge what the forecast says.
TFB standardizes evaluation pipelines to make forecasting comparisons
more consistent and reproducible~\cite{qiu2024tfb}. TimeGym can sit beside
these checks: an API can be well formed while an application-specific
transformation fails a shift check.
Production checklists for machine-learning systems are written as
assertions and scored~\cite{breck2017test}; TimeGym supplies
forecasting-specific assertions for the model evaluation stage of that
workflow~\cite{amershi2019software}.

\begin{table}[t]
\centering\small
\newcommand{\yes}{\checkmark}
\newcommand{\partly}{(\checkmark)}
\newcommand{\no}{--}
\newcommand{\work}[1]{\makebox[1.8em][l]{\rotatebox{55}{#1}}}
\begin{tabular}{@{}p{0.46\linewidth}cccccc@{}}
\toprule
 & \work{SynTSBench} & \work{SimpleTimeBench} & \work{FinStressTS} & \work{Dwarakanath et al.} & \work{sktime conformance} & \work{TimeGym} \\
\midrule
\multicolumn{7}{@{}l}{\emph{Use}} \\
Compares a set of trained models & \yes & \yes & \yes & \no & \no & \no \\
Tests the software of one pipeline or estimator & \no & \no & \no & \yes & \yes & \yes \\
\addlinespace
\multicolumn{7}{@{}l}{\emph{Data}} \\
Synthetic series with known rules & \yes & \yes & \yes & \no & \partly & \yes \\
Real series & \partly & \yes & \no & \yes & \no & \partly \\
Financial mechanisms & \partly & \partly & \yes & \no & \no & \yes \\
\addlinespace
\multicolumn{7}{@{}l}{\emph{Scoring}} \\
Best possible forecast for each case & \yes & \partly & \partly & \no & \no & \yes \\
Scores stated against that forecast & \no & \no & \no & \no & \no & \yes \\
Probabilistic forecasts scored & \no & \partly & \yes & \no & \no & \yes \\
Joint dependence checked & \no & \no & \partly & \no & \no & \yes \\
\addlinespace
\multicolumn{7}{@{}l}{\emph{Software checks}} \\
Output or interface checks & \no & \no & \no & \partly & \yes & \yes \\
What-if (metamorphic) relations & \no & \no & \no & \yes & \no & \yes \\
Injected faults & \no & \no & \no & \yes & \no & \yes \\
Previously unknown bugs reported & \no & \no & \no & \yes & \no & \no \\
\bottomrule
\end{tabular}
\caption{What each work is for and what it can do, as read from each
paper and, for sktime, from the source of its conformance suite.
\checkmark: yes; (\checkmark): in part; --: not reported. In part means:
SynTSBench validates on real datasets in an appendix and has one
stock-price simulator; SimpleTimeBench describes each process's ideal
forecast without scoring against it, uses the CRPS only on real datasets,
and includes geometric Brownian motion and heavy-tailed noise; FinStressTS
has the true conditional distribution but scores against realized values,
and judges dependence through the CRPS of the cross-sectional sum;
Dwarakanath et al. test how degenerate inputs are handled; sktime
generates unstructured noise to exercise the interface; TimeGym uses one
real series.}
\label{tab:positioning}
\end{table}

Table~\ref{tab:positioning} shows where the works differ in use as well
as in capability. The three benchmarks compare many trained models and
are the place to look for a ranking; TimeGym ranks none. Dwarakanath et
al., sktime, and TimeGym instead test one piece of software. Only
SynTSBench also derives the best possible forecast for each task, and it
reports that optimum beside raw errors, where TimeGym states every score
against it. Only FinStressTS also scores probabilistic forecasts on
synthetic series, through the CRPS of the sum across series, where
TimeGym scores each series and checks the joint covariance of the sample
paths. The comparison
runs in both directions. FinStressTS covers more financial mechanisms
than TimeGym's catalog and benchmarks models on them, which neither study
here does. Dwarakanath et al. define nineteen relations to TimeGym's two
shift relations, generate their mutants with a tool where ours are
hand-picked, and found bugs in a deployed application. The sktime suite
covers far more of an estimator's interface, including refitting,
updating, and persistence, than TimeGym's six output checks.

\subsection{Statistical foundations}

Point forecasts must be scored with a loss that matches what they stand
for: squared error for the mean, absolute error for the median, and the
pinball loss for a
quantile~\cite{gneiting2011point,gneiting2011quantiles}. Probabilistic
forecasts should be as sharp as possible while staying
calibrated~\cite{gneiting2007probabilistic,gneiting2014probabilistic},
and proper scoring rules reward both at once~\cite{gneiting2007proper}.
TimeGym computes the averaged CRPS~\cite{matheson1976scoring}, the energy
score~\cite{gneiting2008assessing} and the variogram score, which responds
to errors in dependence~\cite{scheuerer2015variogram}. The CRPS
decomposition~\cite{hersbach2000decomposition}, weighted
CRPS~\cite{gneiting2011comparing}, multivariate rank
histograms~\cite{gneiting2008assessing} and conditional coverage
tests~\cite{christoffersen1998evaluating} are outside its scope.

None of these ingredients is new: the oracle types, metamorphic relations,
financial stress mechanisms, and proper scores all have the precedents
cited above. The contribution is their combination into executable
forecasting checks with recorded assumptions.

\section{Discussion}
\label{sec:discussion}

\subsection{What the checks establish}

Study 1 confirms the mechanics (RQ1 and RQ2): every corruption
was detected, and repaired outputs reproduced the clean control exactly,
including its false alarms. Because the six output corruptions are direct
violations, this shows that the output checks are implemented correctly,
not that they find subtle faults.

For probabilistic forecasts (RQ3), marginal-preserving permutations isolate a
blind spot without relying on a noisy difference in marginal quality.
The factorized negative control is equally important: shuffling samples
is not itself evidence of a fault when the intended law has independent
coordinates.

For RQ4, the implementation results (Figure~\ref{fig:diagnostic-evidence}
and Table~\ref{tab:supplement-complete}) make the diagnostic distinction
concrete. Reversing lags and retaining stale
state pass both shift checks even though they disagree with
the intended algorithm. A common shift also misses a numeric target swap;
different target shifts expose it.
Agreement with a second implementation is more discriminating in this
corpus, but requires a second specification of the intended algorithm
and does not establish forecasting optimality.

The repair counterexample further separates verification from model
selection. Correcting the stale-state regression restores the declared
recursive algorithm, while the random-walk condition with drift favors the faulty
variant's realized loss. A score can therefore reward a specification
violation, just as a conforming model can forecast poorly. A useful report
retains both observations rather than relabeling one as the other.

For RQ5, the covariance results (Tables~\ref{tab:supplement-familywise}
and~\ref{tab:supplement-covariance}) also show why per-check rejection and
suite-level false alarms must be reported separately, together with the
power lost after multiplicity adjustment.

\subsection{Limitations}

The evidence concerns specified properties under controlled conditions.
The fault-study corpus is Gaussian and hand-picked.
The software catalog also includes composed, financial, and multivariate
mechanisms; no fault study covers them, and Appendix~\ref{app:extended}
states their own limits. Known DGP parameters and declared change
schedules make answer keys more informative than anything available to an
ordinary fitted forecaster, so a large avoidable error can reflect
estimation or model limitations. Marginal invariance under coordinate
permutations is an exact algebraic property; the power of a particular
joint check remains an empirical question.

\paragraph{Fault selection.}
All faults were chosen by the developer to illustrate the implemented
checks, which favors detectability, and all repairs were specified in
advance. Detection rates therefore describe this corpus, not the
frequency or detectability of naturally occurring faults.
Mutation analysis rests on the coupling of simple seeded faults to real
ones~\cite{jia2011analysis}. Generated mutants have been found similar
in detectability to real faults, whereas hand-seeded faults were harder
to detect~\cite{andrews2005mutation}; in a corpus of 357 real faults, 73\%
were coupled to common mutants and 17\% to none~\cite{just2014mutants}.
The present corpus is hand-selected and therefore closer to the
hand-seeded case. In deep learning, mutation operators have instead been
derived from a taxonomy of real
faults~\cite{humbatova2020taxonomy,humbatova2021deepcrime}, and mutant
killing is defined statistically because retraining alone moves
accuracy~\cite{jahangirova2020empirical}; no comparable taxonomy of
forecasting-pipeline faults informs the present selection, and the
calibrated loss alarm of Section~\ref{sec:implementation-study} plays the
corresponding statistical role here.
Study 2 moves beyond output corruption but does not supply an independent
corpus of naturally occurring defects. The second implementation verifies
one intended algorithm; it cannot establish that the algorithm is
appropriate for every forecasting task. These studies assess detection and
repair rechecks, not automatic source-code localization or repair
discovery.

\paragraph{Statistical scope.}
Statistical guarantees apply to their stated nulls and sampling units.
The covariance rank check (Equation~\eqref{eq:mc}) assumes independent
reference-law samples; the loss rank check (Equation~\eqref{eq:loss-rank})
assumes exchangeable clean cases within a condition.
Numerical conformance tolerances are engineering choices. Holm adjustment
protects only the declared four-check family. Other reported comparisons
are descriptive, and fresh evaluation seeds do not undo adaptive choice
of scenarios, mutations, or diagnostics.

\subsection{Toward a broader evaluation}

AirPassengers supplies one known preprocessing repair at four origins
from one real series (Table~\ref{tab:real}). It does not establish
synthetic-to-real transfer. The scenario families were chosen for
transparent references, not for representativeness; even forecasting
benchmark collections differ from the distribution of real
series~\cite{spiliotis2020forecasting}. A useful development sequence isolates mechanisms, compares matched pairs,
then adds supported interactions. Seeds, parameter regions, and compositions
must be reserved from model development; a synthetic suite can itself be
overfit. TimeGym can also serve as a training ground: a model can be tuned
against the scorecard with 20 or 30 instances per cell, so that close
candidates are told apart, or with \code{seed=None}, so that every
evaluation draws fresh series. Fresh seeds give new draws of the same
mechanisms, not new mechanisms, so a model tuned this way still needs
held-out kinds of series and untouched real data. After pipeline selection, untouched real data and rolling forecast
origins provide a separate transfer assessment.
A broader study should prospectively collect faults from independent
repositories, compare directly with relevant framework tests, and
evaluate documented repairs on additional real tasks and unseen
mechanisms. Those results are needed to assess practical coverage and
external validity.

\section{Conclusion}

TimeGym connects forecast specifications to reproducible evidence. Its
answer keys separate avoidable error from future noise, and a scorecard
shows that split for ten kinds of series; output checks and shift checks
tell different kinds of failure apart. Study 2 shows that
well-formed outputs and passed shift checks can coexist with algorithmic
regressions, and that restoring an algorithm need not improve its forecast
loss. The probabilistic experiments make the corresponding distinction between
marginal quality and joint dependence, including a null-preserving
negative control and explicit multiplicity handling. Together these
results support a practical discipline: state the property, reproduce the
case, inspect the evidence, and verify a proposed repair against both the
original failure and clean controls. The accompanying software and paired
records provide a basis for testing that discipline on independently
collected forecasting defects.

\section*{Code and data availability}

The implementation and executable studies accompany this manuscript in the
\href{https://github.com/diogoseca/timegym}{TimeGym repository}; its README
gives installation, test, and experiment commands. The package is developed
in executable source notebooks and exported to modules with nbdev. An
explicit build checks complete module ownership and detects
notebook--module drift without executing examples. The tutorials start
from a scorecard quick start and cover custom models, univariate and
multivariate forecasts, marginal and joint uncertainty, fault diagnosis
and repair, replay, real data, capability sweeps, target transforms,
reductions, and the composed and financial scenarios; a separate checker
executes every active notebook in a fresh kernel. These teaching examples are distinct from the two studies.

Both studies ran with Python \StudyPython, NumPy \StudyNumPy, pandas
\StudyPandas, and SciPy \StudySciPy; Study 2 also used scikit-learn
\StudySklearn. The retained Study 2 design requires \StudySklearnFits{}
scikit-learn fits,
\StudyNumPySolves{} independent NumPy solves, and \StudyCovarianceChecks{}
covariance checks containing \StudyCalibrationEnsembles{} calibration
ensembles. Each study's protocol, raw case records, paired score records,
source and data fingerprints, and summary tables are saved under
\code{experiments/results/} and \code{experiments/results-v2/}, and every
table in Section~\ref{sec:results} is generated from them. Following the
reporting recommendations for forecasting research, the records give a
well-defined data-generating process with its seeds, described methods and
parameterization, software versions, and defined
measures~\cite{makridakis2018objectivity}. Results correspond to the result
digests \texttt{cb8b53b50f30c6b6} (Study 1) and \texttt{13b1228e7124b95f}
(Study 2, which was \texttt{ba9b2e90145ca50d} before two offline record
migrations that removed duplicate records and renamed fields without
rerunning the study or changing any measurement); the complete SHA-256 values are recorded
in the experiment manifests.

% The standalone builder renders references.bib into a BibTeX-free inline list.
\bibliographystyle{plain}
\bibliography{references}

\appendix

\section{Composed, financial, and joint scenarios}
\label{app:extended}

Beyond the five core families and the AR(2), the software catalog has
four further scenario kinds: \code{composite},
\code{merton\_jump\_diffusion}, \code{garch}, and
\code{markov\_switching}. The scorecard draws two of its rows from
\code{composite}; nothing else here is part of the scorecard. All four use the same
\code{Scenario} and \code{Pipeline} interfaces as the core families and
support multivariate targets under explicit dependence laws. Neither study
uses them: their executable examples and mathematical validation do not
expand the Gaussian fault-study corpus of Sections~\ref{sec:protocol}
and~\ref{sec:results}.

\subsection{Seasonality, process drift, and changing mechanisms}

Seasonality is recurring structure; a fixed sinusoid is not the drift term
of a random walk. Nor is a stable seasonal law automatically concept drift
when calendar phase is included in the conditioning information. We distinguish
\emph{process drift}, an expected increment, from changes in the generating
parameters or predictive relationship~\cite{gama2014drift}. The latter can
change a trend slope, seasonal amplitude, period or phase, innovation scale,
or dependence. Likewise, a GARCH variance state or a hidden Markov regime
can change under a fixed law: a state transition and a parameter change are
separate experimental interventions.

The \code{composite} scenario specifies a deterministic component
$g_t=b+ct+A\sin(2\pi t/p+\phi)$ and one residual process: fresh Gaussian
noise or an accumulated-shock process. The latter may include independent
compound-Poisson increments. Declared changes can replace mean parameters
or noise scale at an explicit index. The before/after mean uses absolute
$t$, so changing the slope need not preserve level continuity. These are
additive level constructions, not multiplicative financial price models.
The evaluator knows the change schedule; an ordinary fitted model does not.

\subsection{Financial stress mechanisms and non-Gaussian references}

Figure~\ref{fig:financial-stresses} illustrates the financial families,
which isolate established stochastic mechanisms and selected interactions.
Their targets are price, log price, or log return, which have different
conditional means.

For \code{merton\_jump\_diffusion}, let $J_i\sim\mathcal N(a,b_J^2)$,
$N\sim\operatorname{Poisson}(\lambda\Delta t)$, and
$\kappa=\exp(a+b_J^2/2)-1$. Independent one-step log returns follow
\begin{equation}
 R=(\mu-\tfrac12\sigma^2-\lambda\kappa)\Delta t
   +\sigma\sqrt{\Delta t}\,Z+\sum_{i=1}^{N}J_i,
 \qquad Z\sim\mathcal N(0,1).
\end{equation}
This is the Merton construction~\cite{merton1976jumps}, using a physical
expected-price-growth convention rather than a risk-neutral pricing
calibration. Zero jump intensity reduces it to geometric Brownian motion.
The cumulant $K(u)=\log\E[e^{uR}]$ gives the conditional price mean
$S_T\exp(hK(1))$ and price covariance; exponentiating a conditional
log-price mean generally gives the wrong mean on the price scale.

For \code{garch}, Gaussian innovations follow
$r_t=\mu+\sqrt{v_t}z_t$ and
$v_{t+1}=\omega+\alpha(r_t-\mu)^2+\beta v_t$, with
$\omega>0$, $\alpha,\beta\geq0$, and $\alpha+\beta<1$; this is the GARCH
generalization~\cite{bollerslev1986garch} of Engle's ARCH
process~\cite{engle1982autoregressive}.
The next variance is reconstructed from observed returns. Each simulated
future recursively updates its own variance; replacing that recursion by
independent Gaussian draws with expected variances loses part of the joint
law. The implementation supplies log-return or log-price tasks, not price
moments obtained by an unjustified exponentiation.

For \code{markov\_switching}, a finite hidden Markov chain controls return
means, diffusion scales, and optionally jump parameters. The reference
filters $\Pr(S_T\mid y_{0:T},\theta)$ from observed returns and propagates
that distribution, following the hidden-regime principle~\cite{hamilton1989regimes}.
It never conditions on the simulator's actual terminal regime. Gaussian
emission filtering supports an exact conditional sampler under the declared
law. With jumps, the emission likelihood is a truncated Poisson mixture:
its omitted probability mass is recorded, but is not a posterior-error
bound. Those references are marked approximate.

\code{SimulationReference} retains analytical means and either analytical or
explicitly estimated covariance under its declared conditioning distribution,
and draws joint futures from an explicit sampler.
Its marginal quantiles are reproducible Monte Carlo estimates, with sample
count and seed recorded, rather than Gaussian quantiles inferred from two
moments. The covariance rank check returns \texttt{UNSUPPORTED} for
approximate conditional samplers. Shape checks and descriptive forecast
scores remain available; no exact test guarantee is attached to the
approximate hidden-regime reference.

\begin{figure}[p]
\centering
% BEGIN VECTOR FIGURE paper-financial-stresses
% Generated vector scene; SVG and PNG use the same plotted geometry.
\begingroup
\definecolor{tg0000ff}{HTML}{0000ff}
\definecolor{tg18314b}{HTML}{18314b}
\definecolor{tg575757}{HTML}{575757}
\definecolor{tg626b73}{HTML}{626b73}
\definecolor{tg808080}{HTML}{808080}
\definecolor{tg9da7b0}{HTML}{9da7b0}
\definecolor{tgc9cdd1}{HTML}{c9cdd1}
\resizebox{\linewidth}{!}{%
\begin{tikzpicture}[x=1bp,y=-1bp]
\path[use as bounding box] (0,0) rectangle (448,484);
\node[anchor=west,inner sep=0pt,outer sep=0pt,text=tg18314b,font={\sffamily\fontsize{10.2bp}{12.24bp}\selectfont\bfseries}] at (0.000,10.000) {Six financial configurations. Three seeds each.};
\fill[tg808080] (10.000,29.000) circle[radius=1.1bp];
\node[anchor=west,inner sep=0pt,outer sep=0pt,text=tg18314b,font={\sffamily\fontsize{7.4bp}{8.88bp}\selectfont}] at (22.000,29.000) {Observed values};
\draw[color=tg0000ff,line width=0.65bp,line cap=round] (136.000,27.000) -- (140.000,31.000);
\draw[color=tg0000ff,line width=0.65bp,line cap=round] (136.000,31.000) -- (140.000,27.000);
\node[anchor=west,inner sep=0pt,outer sep=0pt,text=tg18314b,font={\sffamily\fontsize{7.4bp}{8.88bp}\selectfont}] at (150.000,29.000) {Last-value forecast};
\draw[color=tg575757,line width=0.8bp,line cap=round] (284.000,29.000) -- (296.000,29.000);
\node[anchor=west,inner sep=0pt,outer sep=0pt,text=tg18314b,font={\sffamily\fontsize{7.4bp}{8.88bp}\selectfont}] at (301.000,29.000) {Conditional mean};
\node[anchor=west,inner sep=0pt,outer sep=0pt,text=tg18314b,font={\sffamily\fontsize{8bp}{9.6bp}\selectfont\bfseries}] at (0.000,52.000) {Geometric Brownian motion (log price)};
\node[anchor=east,inner sep=0pt,outer sep=0pt,text=tg626b73,font={\sffamily\fontsize{7bp}{8.4bp}\selectfont}] at (448.000,52.000) {GBM limit; zero jumps};
\draw[color=tgc9cdd1,line width=0.45bp,line cap=round] (28.000,61.000) -- (28.000,104.000);
\draw[color=tgc9cdd1,line width=0.45bp,line cap=round] (28.000,104.000) -- (147.000,104.000);
\node[anchor=east,inner sep=0pt,outer sep=0pt,text=tg626b73,font={\sffamily\fontsize{6.7bp}{8.04bp}\selectfont}] at (25.000,97.980) {4.56};
\node[anchor=east,inner sep=0pt,outer sep=0pt,text=tg626b73,font={\sffamily\fontsize{6.7bp}{8.04bp}\selectfont}] at (25.000,67.020) {4.65};
\draw[color=tg9da7b0,line width=0.6bp,line cap=round,dash pattern=on 2bp off 2bp] (117.250,61.000) -- (117.250,104.000);
\fill[tg808080] (28.000,68.776) circle[radius=0.52bp];
\fill[tg808080] (29.240,72.176) circle[radius=0.52bp];
\fill[tg808080] (30.479,77.755) circle[radius=0.52bp];
\fill[tg808080] (31.719,74.489) circle[radius=0.52bp];
\fill[tg808080] (32.958,76.026) circle[radius=0.52bp];
\fill[tg808080] (34.198,78.251) circle[radius=0.52bp];
\fill[tg808080] (35.438,76.651) circle[radius=0.52bp];
\fill[tg808080] (36.677,78.308) circle[radius=0.52bp];
\fill[tg808080] (37.917,82.076) circle[radius=0.52bp];
\fill[tg808080] (39.156,78.113) circle[radius=0.52bp];
\fill[tg808080] (40.396,77.107) circle[radius=0.52bp];
\fill[tg808080] (41.635,77.785) circle[radius=0.52bp];
\fill[tg808080] (42.875,74.919) circle[radius=0.52bp];
\fill[tg808080] (44.115,76.744) circle[radius=0.52bp];
\fill[tg808080] (45.354,77.328) circle[radius=0.52bp];
\fill[tg808080] (46.594,76.960) circle[radius=0.52bp];
\fill[tg808080] (47.833,69.085) circle[radius=0.52bp];
\fill[tg808080] (49.073,69.969) circle[radius=0.52bp];
\fill[tg808080] (50.312,65.481) circle[radius=0.52bp];
\fill[tg808080] (51.552,65.161) circle[radius=0.52bp];
\fill[tg808080] (52.792,67.488) circle[radius=0.52bp];
\fill[tg808080] (54.031,70.852) circle[radius=0.52bp];
\fill[tg808080] (55.271,71.882) circle[radius=0.52bp];
\fill[tg808080] (56.510,73.588) circle[radius=0.52bp];
\fill[tg808080] (57.750,78.648) circle[radius=0.52bp];
\fill[tg808080] (58.990,78.194) circle[radius=0.52bp];
\fill[tg808080] (60.229,72.839) circle[radius=0.52bp];
\fill[tg808080] (61.469,70.998) circle[radius=0.52bp];
\fill[tg808080] (62.708,74.731) circle[radius=0.52bp];
\fill[tg808080] (63.948,71.409) circle[radius=0.52bp];
\fill[tg808080] (65.188,74.059) circle[radius=0.52bp];
\fill[tg808080] (66.427,79.311) circle[radius=0.52bp];
\fill[tg808080] (67.667,79.534) circle[radius=0.52bp];
\fill[tg808080] (68.906,82.022) circle[radius=0.52bp];
\fill[tg808080] (70.146,84.878) circle[radius=0.52bp];
\fill[tg808080] (71.385,83.607) circle[radius=0.52bp];
\fill[tg808080] (72.625,81.692) circle[radius=0.52bp];
\fill[tg808080] (73.865,81.456) circle[radius=0.52bp];
\fill[tg808080] (75.104,84.572) circle[radius=0.52bp];
\fill[tg808080] (76.344,85.599) circle[radius=0.52bp];
\fill[tg808080] (77.583,84.161) circle[radius=0.52bp];
\fill[tg808080] (78.823,80.496) circle[radius=0.52bp];
\fill[tg808080] (80.062,78.901) circle[radius=0.52bp];
\fill[tg808080] (81.302,75.326) circle[radius=0.52bp];
\fill[tg808080] (82.542,73.846) circle[radius=0.52bp];
\fill[tg808080] (83.781,71.115) circle[radius=0.52bp];
\fill[tg808080] (85.021,73.263) circle[radius=0.52bp];
\fill[tg808080] (86.260,73.002) circle[radius=0.52bp];
\fill[tg808080] (87.500,77.936) circle[radius=0.52bp];
\fill[tg808080] (88.740,77.516) circle[radius=0.52bp];
\fill[tg808080] (89.979,77.408) circle[radius=0.52bp];
\fill[tg808080] (91.219,80.332) circle[radius=0.52bp];
\fill[tg808080] (92.458,81.627) circle[radius=0.52bp];
\fill[tg808080] (93.698,85.953) circle[radius=0.52bp];
\fill[tg808080] (94.938,85.631) circle[radius=0.52bp];
\fill[tg808080] (96.177,82.999) circle[radius=0.52bp];
\fill[tg808080] (97.417,81.016) circle[radius=0.52bp];
\fill[tg808080] (98.656,76.662) circle[radius=0.52bp];
\fill[tg808080] (99.896,76.235) circle[radius=0.52bp];
\fill[tg808080] (101.135,78.508) circle[radius=0.52bp];
\fill[tg808080] (102.375,78.447) circle[radius=0.52bp];
\fill[tg808080] (103.615,76.421) circle[radius=0.52bp];
\fill[tg808080] (104.854,83.302) circle[radius=0.52bp];
\fill[tg808080] (106.094,84.687) circle[radius=0.52bp];
\fill[tg808080] (107.333,84.446) circle[radius=0.52bp];
\fill[tg808080] (108.573,84.032) circle[radius=0.52bp];
\fill[tg808080] (109.812,83.285) circle[radius=0.52bp];
\fill[tg808080] (111.052,79.812) circle[radius=0.52bp];
\fill[tg808080] (112.292,79.188) circle[radius=0.52bp];
\fill[tg808080] (113.531,76.937) circle[radius=0.52bp];
\fill[tg808080] (114.771,78.938) circle[radius=0.52bp];
\fill[tg808080] (116.010,83.951) circle[radius=0.52bp];
\fill[tg808080] (117.250,84.767) circle[radius=0.52bp];
\fill[tg808080] (118.490,83.346) circle[radius=0.52bp];
\fill[tg808080] (119.729,82.896) circle[radius=0.52bp];
\fill[tg808080] (120.969,85.471) circle[radius=0.52bp];
\fill[tg808080] (122.208,83.468) circle[radius=0.52bp];
\fill[tg808080] (123.448,85.809) circle[radius=0.52bp];
\fill[tg808080] (124.688,93.657) circle[radius=0.52bp];
\fill[tg808080] (125.927,95.543) circle[radius=0.52bp];
\fill[tg808080] (127.167,96.978) circle[radius=0.52bp];
\fill[tg808080] (128.406,95.040) circle[radius=0.52bp];
\fill[tg808080] (129.646,99.839) circle[radius=0.52bp];
\fill[tg808080] (130.885,99.208) circle[radius=0.52bp];
\fill[tg808080] (132.125,99.832) circle[radius=0.52bp];
\fill[tg808080] (133.365,97.881) circle[radius=0.52bp];
\fill[tg808080] (134.604,98.198) circle[radius=0.52bp];
\fill[tg808080] (135.844,98.234) circle[radius=0.52bp];
\fill[tg808080] (137.083,98.847) circle[radius=0.52bp];
\fill[tg808080] (138.323,95.159) circle[radius=0.52bp];
\fill[tg808080] (139.562,94.639) circle[radius=0.52bp];
\fill[tg808080] (140.802,91.976) circle[radius=0.52bp];
\fill[tg808080] (142.042,90.933) circle[radius=0.52bp];
\fill[tg808080] (143.281,89.970) circle[radius=0.52bp];
\fill[tg808080] (144.521,89.665) circle[radius=0.52bp];
\fill[tg808080] (145.760,91.873) circle[radius=0.52bp];
\draw[color=tg575757,line width=0.65bp,line cap=round] (117.250,83.913) -- (118.490,83.876) -- (119.729,83.838) -- (120.969,83.801) -- (122.208,83.763) -- (123.448,83.725) -- (124.688,83.688) -- (125.927,83.650) -- (127.167,83.613) -- (128.406,83.575) -- (129.646,83.537) -- (130.885,83.500) -- (132.125,83.462) -- (133.365,83.425) -- (134.604,83.387) -- (135.844,83.350) -- (137.083,83.312) -- (138.323,83.274) -- (139.562,83.237) -- (140.802,83.199) -- (142.042,83.162) -- (143.281,83.124) -- (144.521,83.086) -- (145.760,83.049);
\draw[color=tg0000ff,line width=0.38bp,line cap=round] (116.750,83.451) -- (117.750,84.451);
\draw[color=tg0000ff,line width=0.38bp,line cap=round] (116.750,84.451) -- (117.750,83.451);
\draw[color=tg0000ff,line width=0.38bp,line cap=round] (117.990,83.451) -- (118.990,84.451);
\draw[color=tg0000ff,line width=0.38bp,line cap=round] (117.990,84.451) -- (118.990,83.451);
\draw[color=tg0000ff,line width=0.38bp,line cap=round] (119.229,83.451) -- (120.229,84.451);
\draw[color=tg0000ff,line width=0.38bp,line cap=round] (119.229,84.451) -- (120.229,83.451);
\draw[color=tg0000ff,line width=0.38bp,line cap=round] (120.469,83.451) -- (121.469,84.451);
\draw[color=tg0000ff,line width=0.38bp,line cap=round] (120.469,84.451) -- (121.469,83.451);
\draw[color=tg0000ff,line width=0.38bp,line cap=round] (121.708,83.451) -- (122.708,84.451);
\draw[color=tg0000ff,line width=0.38bp,line cap=round] (121.708,84.451) -- (122.708,83.451);
\draw[color=tg0000ff,line width=0.38bp,line cap=round] (122.948,83.451) -- (123.948,84.451);
\draw[color=tg0000ff,line width=0.38bp,line cap=round] (122.948,84.451) -- (123.948,83.451);
\draw[color=tg0000ff,line width=0.38bp,line cap=round] (124.188,83.451) -- (125.188,84.451);
\draw[color=tg0000ff,line width=0.38bp,line cap=round] (124.188,84.451) -- (125.188,83.451);
\draw[color=tg0000ff,line width=0.38bp,line cap=round] (125.427,83.451) -- (126.427,84.451);
\draw[color=tg0000ff,line width=0.38bp,line cap=round] (125.427,84.451) -- (126.427,83.451);
\draw[color=tg0000ff,line width=0.38bp,line cap=round] (126.667,83.451) -- (127.667,84.451);
\draw[color=tg0000ff,line width=0.38bp,line cap=round] (126.667,84.451) -- (127.667,83.451);
\draw[color=tg0000ff,line width=0.38bp,line cap=round] (127.906,83.451) -- (128.906,84.451);
\draw[color=tg0000ff,line width=0.38bp,line cap=round] (127.906,84.451) -- (128.906,83.451);
\draw[color=tg0000ff,line width=0.38bp,line cap=round] (129.146,83.451) -- (130.146,84.451);
\draw[color=tg0000ff,line width=0.38bp,line cap=round] (129.146,84.451) -- (130.146,83.451);
\draw[color=tg0000ff,line width=0.38bp,line cap=round] (130.385,83.451) -- (131.385,84.451);
\draw[color=tg0000ff,line width=0.38bp,line cap=round] (130.385,84.451) -- (131.385,83.451);
\draw[color=tg0000ff,line width=0.38bp,line cap=round] (131.625,83.451) -- (132.625,84.451);
\draw[color=tg0000ff,line width=0.38bp,line cap=round] (131.625,84.451) -- (132.625,83.451);
\draw[color=tg0000ff,line width=0.38bp,line cap=round] (132.865,83.451) -- (133.865,84.451);
\draw[color=tg0000ff,line width=0.38bp,line cap=round] (132.865,84.451) -- (133.865,83.451);
\draw[color=tg0000ff,line width=0.38bp,line cap=round] (134.104,83.451) -- (135.104,84.451);
\draw[color=tg0000ff,line width=0.38bp,line cap=round] (134.104,84.451) -- (135.104,83.451);
\draw[color=tg0000ff,line width=0.38bp,line cap=round] (135.344,83.451) -- (136.344,84.451);
\draw[color=tg0000ff,line width=0.38bp,line cap=round] (135.344,84.451) -- (136.344,83.451);
\draw[color=tg0000ff,line width=0.38bp,line cap=round] (136.583,83.451) -- (137.583,84.451);
\draw[color=tg0000ff,line width=0.38bp,line cap=round] (136.583,84.451) -- (137.583,83.451);
\draw[color=tg0000ff,line width=0.38bp,line cap=round] (137.823,83.451) -- (138.823,84.451);
\draw[color=tg0000ff,line width=0.38bp,line cap=round] (137.823,84.451) -- (138.823,83.451);
\draw[color=tg0000ff,line width=0.38bp,line cap=round] (139.062,83.451) -- (140.062,84.451);
\draw[color=tg0000ff,line width=0.38bp,line cap=round] (139.062,84.451) -- (140.062,83.451);
\draw[color=tg0000ff,line width=0.38bp,line cap=round] (140.302,83.451) -- (141.302,84.451);
\draw[color=tg0000ff,line width=0.38bp,line cap=round] (140.302,84.451) -- (141.302,83.451);
\draw[color=tg0000ff,line width=0.38bp,line cap=round] (141.542,83.451) -- (142.542,84.451);
\draw[color=tg0000ff,line width=0.38bp,line cap=round] (141.542,84.451) -- (142.542,83.451);
\draw[color=tg0000ff,line width=0.38bp,line cap=round] (142.781,83.451) -- (143.781,84.451);
\draw[color=tg0000ff,line width=0.38bp,line cap=round] (142.781,84.451) -- (143.781,83.451);
\draw[color=tg0000ff,line width=0.38bp,line cap=round] (144.021,83.451) -- (145.021,84.451);
\draw[color=tg0000ff,line width=0.38bp,line cap=round] (144.021,84.451) -- (145.021,83.451);
\draw[color=tg0000ff,line width=0.38bp,line cap=round] (145.260,83.451) -- (146.260,84.451);
\draw[color=tg0000ff,line width=0.38bp,line cap=round] (145.260,84.451) -- (146.260,83.451);
\node[anchor=center,inner sep=0pt,outer sep=0pt,text=tg18314b,font={\sffamily\fontsize{7.4bp}{8.88bp}\selectfont\bfseries}] at (87.500,40.000) {Seed 11};
\draw[color=tgc9cdd1,line width=0.45bp,line cap=round] (174.000,61.000) -- (174.000,104.000);
\draw[color=tgc9cdd1,line width=0.45bp,line cap=round] (174.000,104.000) -- (293.000,104.000);
\node[anchor=east,inner sep=0pt,outer sep=0pt,text=tg626b73,font={\sffamily\fontsize{6.7bp}{8.04bp}\selectfont}] at (171.000,97.980) {4.53};
\node[anchor=east,inner sep=0pt,outer sep=0pt,text=tg626b73,font={\sffamily\fontsize{6.7bp}{8.04bp}\selectfont}] at (171.000,67.020) {4.66};
\draw[color=tg9da7b0,line width=0.6bp,line cap=round,dash pattern=on 2bp off 2bp] (263.250,61.000) -- (263.250,104.000);
\fill[tg808080] (174.000,79.459) circle[radius=0.52bp];
\fill[tg808080] (175.240,83.444) circle[radius=0.52bp];
\fill[tg808080] (176.479,76.304) circle[radius=0.52bp];
\fill[tg808080] (177.719,78.238) circle[radius=0.52bp];
\fill[tg808080] (178.958,79.773) circle[radius=0.52bp];
\fill[tg808080] (180.198,79.328) circle[radius=0.52bp];
\fill[tg808080] (181.438,76.520) circle[radius=0.52bp];
\fill[tg808080] (182.677,71.685) circle[radius=0.52bp];
\fill[tg808080] (183.917,71.796) circle[radius=0.52bp];
\fill[tg808080] (185.156,68.269) circle[radius=0.52bp];
\fill[tg808080] (186.396,66.229) circle[radius=0.52bp];
\fill[tg808080] (187.635,65.677) circle[radius=0.52bp];
\fill[tg808080] (188.875,65.161) circle[radius=0.52bp];
\fill[tg808080] (190.115,67.090) circle[radius=0.52bp];
\fill[tg808080] (191.354,68.300) circle[radius=0.52bp];
\fill[tg808080] (192.594,66.723) circle[radius=0.52bp];
\fill[tg808080] (193.833,65.991) circle[radius=0.52bp];
\fill[tg808080] (195.073,67.605) circle[radius=0.52bp];
\fill[tg808080] (196.312,72.097) circle[radius=0.52bp];
\fill[tg808080] (197.552,68.951) circle[radius=0.52bp];
\fill[tg808080] (198.792,73.534) circle[radius=0.52bp];
\fill[tg808080] (200.031,71.553) circle[radius=0.52bp];
\fill[tg808080] (201.271,72.514) circle[radius=0.52bp];
\fill[tg808080] (202.510,72.419) circle[radius=0.52bp];
\fill[tg808080] (203.750,70.278) circle[radius=0.52bp];
\fill[tg808080] (204.990,68.438) circle[radius=0.52bp];
\fill[tg808080] (206.229,68.521) circle[radius=0.52bp];
\fill[tg808080] (207.469,65.597) circle[radius=0.52bp];
\fill[tg808080] (208.708,66.657) circle[radius=0.52bp];
\fill[tg808080] (209.948,66.527) circle[radius=0.52bp];
\fill[tg808080] (211.188,69.750) circle[radius=0.52bp];
\fill[tg808080] (212.427,68.578) circle[radius=0.52bp];
\fill[tg808080] (213.667,68.214) circle[radius=0.52bp];
\fill[tg808080] (214.906,68.953) circle[radius=0.52bp];
\fill[tg808080] (216.146,69.438) circle[radius=0.52bp];
\fill[tg808080] (217.385,70.402) circle[radius=0.52bp];
\fill[tg808080] (218.625,73.812) circle[radius=0.52bp];
\fill[tg808080] (219.865,72.054) circle[radius=0.52bp];
\fill[tg808080] (221.104,73.798) circle[radius=0.52bp];
\fill[tg808080] (222.344,74.469) circle[radius=0.52bp];
\fill[tg808080] (223.583,73.471) circle[radius=0.52bp];
\fill[tg808080] (224.823,72.688) circle[radius=0.52bp];
\fill[tg808080] (226.062,74.337) circle[radius=0.52bp];
\fill[tg808080] (227.302,76.869) circle[radius=0.52bp];
\fill[tg808080] (228.542,74.074) circle[radius=0.52bp];
\fill[tg808080] (229.781,76.609) circle[radius=0.52bp];
\fill[tg808080] (231.021,79.110) circle[radius=0.52bp];
\fill[tg808080] (232.260,79.665) circle[radius=0.52bp];
\fill[tg808080] (233.500,78.167) circle[radius=0.52bp];
\fill[tg808080] (234.740,77.292) circle[radius=0.52bp];
\fill[tg808080] (235.979,73.670) circle[radius=0.52bp];
\fill[tg808080] (237.219,76.658) circle[radius=0.52bp];
\fill[tg808080] (238.458,74.651) circle[radius=0.52bp];
\fill[tg808080] (239.698,75.975) circle[radius=0.52bp];
\fill[tg808080] (240.938,78.855) circle[radius=0.52bp];
\fill[tg808080] (242.177,78.975) circle[radius=0.52bp];
\fill[tg808080] (243.417,84.984) circle[radius=0.52bp];
\fill[tg808080] (244.656,83.316) circle[radius=0.52bp];
\fill[tg808080] (245.896,83.278) circle[radius=0.52bp];
\fill[tg808080] (247.135,84.784) circle[radius=0.52bp];
\fill[tg808080] (248.375,82.617) circle[radius=0.52bp];
\fill[tg808080] (249.615,87.307) circle[radius=0.52bp];
\fill[tg808080] (250.854,87.541) circle[radius=0.52bp];
\fill[tg808080] (252.094,87.453) circle[radius=0.52bp];
\fill[tg808080] (253.333,86.653) circle[radius=0.52bp];
\fill[tg808080] (254.573,87.611) circle[radius=0.52bp];
\fill[tg808080] (255.812,90.625) circle[radius=0.52bp];
\fill[tg808080] (257.052,93.351) circle[radius=0.52bp];
\fill[tg808080] (258.292,96.834) circle[radius=0.52bp];
\fill[tg808080] (259.531,96.989) circle[radius=0.52bp];
\fill[tg808080] (260.771,93.532) circle[radius=0.52bp];
\fill[tg808080] (262.010,94.232) circle[radius=0.52bp];
\fill[tg808080] (263.250,92.216) circle[radius=0.52bp];
\fill[tg808080] (264.490,94.253) circle[radius=0.52bp];
\fill[tg808080] (265.729,94.331) circle[radius=0.52bp];
\fill[tg808080] (266.969,94.858) circle[radius=0.52bp];
\fill[tg808080] (268.208,94.824) circle[radius=0.52bp];
\fill[tg808080] (269.448,96.517) circle[radius=0.52bp];
\fill[tg808080] (270.688,99.839) circle[radius=0.52bp];
\fill[tg808080] (271.927,98.384) circle[radius=0.52bp];
\fill[tg808080] (273.167,94.325) circle[radius=0.52bp];
\fill[tg808080] (274.406,94.836) circle[radius=0.52bp];
\fill[tg808080] (275.646,95.478) circle[radius=0.52bp];
\fill[tg808080] (276.885,92.690) circle[radius=0.52bp];
\fill[tg808080] (278.125,92.200) circle[radius=0.52bp];
\fill[tg808080] (279.365,95.288) circle[radius=0.52bp];
\fill[tg808080] (280.604,96.356) circle[radius=0.52bp];
\fill[tg808080] (281.844,95.102) circle[radius=0.52bp];
\fill[tg808080] (283.083,95.179) circle[radius=0.52bp];
\fill[tg808080] (284.323,94.573) circle[radius=0.52bp];
\fill[tg808080] (285.562,89.799) circle[radius=0.52bp];
\fill[tg808080] (286.802,88.594) circle[radius=0.52bp];
\fill[tg808080] (288.042,88.595) circle[radius=0.52bp];
\fill[tg808080] (289.281,87.393) circle[radius=0.52bp];
\fill[tg808080] (290.521,87.120) circle[radius=0.52bp];
\fill[tg808080] (291.760,89.045) circle[radius=0.52bp];
\draw[color=tg575757,line width=0.65bp,line cap=round] (263.250,94.207) -- (264.490,94.181) -- (265.729,94.155) -- (266.969,94.130) -- (268.208,94.104) -- (269.448,94.078) -- (270.688,94.052) -- (271.927,94.027) -- (273.167,94.001) -- (274.406,93.975) -- (275.646,93.949) -- (276.885,93.924) -- (278.125,93.898) -- (279.365,93.872) -- (280.604,93.846) -- (281.844,93.821) -- (283.083,93.795) -- (284.323,93.769) -- (285.562,93.744) -- (286.802,93.718) -- (288.042,93.692) -- (289.281,93.666) -- (290.521,93.641) -- (291.760,93.615);
\draw[color=tg0000ff,line width=0.38bp,line cap=round] (262.750,93.732) -- (263.750,94.732);
\draw[color=tg0000ff,line width=0.38bp,line cap=round] (262.750,94.732) -- (263.750,93.732);
\draw[color=tg0000ff,line width=0.38bp,line cap=round] (263.990,93.732) -- (264.990,94.732);
\draw[color=tg0000ff,line width=0.38bp,line cap=round] (263.990,94.732) -- (264.990,93.732);
\draw[color=tg0000ff,line width=0.38bp,line cap=round] (265.229,93.732) -- (266.229,94.732);
\draw[color=tg0000ff,line width=0.38bp,line cap=round] (265.229,94.732) -- (266.229,93.732);
\draw[color=tg0000ff,line width=0.38bp,line cap=round] (266.469,93.732) -- (267.469,94.732);
\draw[color=tg0000ff,line width=0.38bp,line cap=round] (266.469,94.732) -- (267.469,93.732);
\draw[color=tg0000ff,line width=0.38bp,line cap=round] (267.708,93.732) -- (268.708,94.732);
\draw[color=tg0000ff,line width=0.38bp,line cap=round] (267.708,94.732) -- (268.708,93.732);
\draw[color=tg0000ff,line width=0.38bp,line cap=round] (268.948,93.732) -- (269.948,94.732);
\draw[color=tg0000ff,line width=0.38bp,line cap=round] (268.948,94.732) -- (269.948,93.732);
\draw[color=tg0000ff,line width=0.38bp,line cap=round] (270.188,93.732) -- (271.188,94.732);
\draw[color=tg0000ff,line width=0.38bp,line cap=round] (270.188,94.732) -- (271.188,93.732);
\draw[color=tg0000ff,line width=0.38bp,line cap=round] (271.427,93.732) -- (272.427,94.732);
\draw[color=tg0000ff,line width=0.38bp,line cap=round] (271.427,94.732) -- (272.427,93.732);
\draw[color=tg0000ff,line width=0.38bp,line cap=round] (272.667,93.732) -- (273.667,94.732);
\draw[color=tg0000ff,line width=0.38bp,line cap=round] (272.667,94.732) -- (273.667,93.732);
\draw[color=tg0000ff,line width=0.38bp,line cap=round] (273.906,93.732) -- (274.906,94.732);
\draw[color=tg0000ff,line width=0.38bp,line cap=round] (273.906,94.732) -- (274.906,93.732);
\draw[color=tg0000ff,line width=0.38bp,line cap=round] (275.146,93.732) -- (276.146,94.732);
\draw[color=tg0000ff,line width=0.38bp,line cap=round] (275.146,94.732) -- (276.146,93.732);
\draw[color=tg0000ff,line width=0.38bp,line cap=round] (276.385,93.732) -- (277.385,94.732);
\draw[color=tg0000ff,line width=0.38bp,line cap=round] (276.385,94.732) -- (277.385,93.732);
\draw[color=tg0000ff,line width=0.38bp,line cap=round] (277.625,93.732) -- (278.625,94.732);
\draw[color=tg0000ff,line width=0.38bp,line cap=round] (277.625,94.732) -- (278.625,93.732);
\draw[color=tg0000ff,line width=0.38bp,line cap=round] (278.865,93.732) -- (279.865,94.732);
\draw[color=tg0000ff,line width=0.38bp,line cap=round] (278.865,94.732) -- (279.865,93.732);
\draw[color=tg0000ff,line width=0.38bp,line cap=round] (280.104,93.732) -- (281.104,94.732);
\draw[color=tg0000ff,line width=0.38bp,line cap=round] (280.104,94.732) -- (281.104,93.732);
\draw[color=tg0000ff,line width=0.38bp,line cap=round] (281.344,93.732) -- (282.344,94.732);
\draw[color=tg0000ff,line width=0.38bp,line cap=round] (281.344,94.732) -- (282.344,93.732);
\draw[color=tg0000ff,line width=0.38bp,line cap=round] (282.583,93.732) -- (283.583,94.732);
\draw[color=tg0000ff,line width=0.38bp,line cap=round] (282.583,94.732) -- (283.583,93.732);
\draw[color=tg0000ff,line width=0.38bp,line cap=round] (283.823,93.732) -- (284.823,94.732);
\draw[color=tg0000ff,line width=0.38bp,line cap=round] (283.823,94.732) -- (284.823,93.732);
\draw[color=tg0000ff,line width=0.38bp,line cap=round] (285.062,93.732) -- (286.062,94.732);
\draw[color=tg0000ff,line width=0.38bp,line cap=round] (285.062,94.732) -- (286.062,93.732);
\draw[color=tg0000ff,line width=0.38bp,line cap=round] (286.302,93.732) -- (287.302,94.732);
\draw[color=tg0000ff,line width=0.38bp,line cap=round] (286.302,94.732) -- (287.302,93.732);
\draw[color=tg0000ff,line width=0.38bp,line cap=round] (287.542,93.732) -- (288.542,94.732);
\draw[color=tg0000ff,line width=0.38bp,line cap=round] (287.542,94.732) -- (288.542,93.732);
\draw[color=tg0000ff,line width=0.38bp,line cap=round] (288.781,93.732) -- (289.781,94.732);
\draw[color=tg0000ff,line width=0.38bp,line cap=round] (288.781,94.732) -- (289.781,93.732);
\draw[color=tg0000ff,line width=0.38bp,line cap=round] (290.021,93.732) -- (291.021,94.732);
\draw[color=tg0000ff,line width=0.38bp,line cap=round] (290.021,94.732) -- (291.021,93.732);
\draw[color=tg0000ff,line width=0.38bp,line cap=round] (291.260,93.732) -- (292.260,94.732);
\draw[color=tg0000ff,line width=0.38bp,line cap=round] (291.260,94.732) -- (292.260,93.732);
\node[anchor=center,inner sep=0pt,outer sep=0pt,text=tg18314b,font={\sffamily\fontsize{7.4bp}{8.88bp}\selectfont\bfseries}] at (233.500,40.000) {Seed 22};
\draw[color=tgc9cdd1,line width=0.45bp,line cap=round] (320.000,61.000) -- (320.000,104.000);
\draw[color=tgc9cdd1,line width=0.45bp,line cap=round] (320.000,104.000) -- (439.000,104.000);
\node[anchor=east,inner sep=0pt,outer sep=0pt,text=tg626b73,font={\sffamily\fontsize{6.7bp}{8.04bp}\selectfont}] at (317.000,97.980) {4.62};
\node[anchor=east,inner sep=0pt,outer sep=0pt,text=tg626b73,font={\sffamily\fontsize{6.7bp}{8.04bp}\selectfont}] at (317.000,67.020) {4.73};
\draw[color=tg9da7b0,line width=0.6bp,line cap=round,dash pattern=on 2bp off 2bp] (409.250,61.000) -- (409.250,104.000);
\fill[tg808080] (320.000,85.461) circle[radius=0.52bp];
\fill[tg808080] (321.240,84.635) circle[radius=0.52bp];
\fill[tg808080] (322.479,82.823) circle[radius=0.52bp];
\fill[tg808080] (323.719,80.970) circle[radius=0.52bp];
\fill[tg808080] (324.958,79.092) circle[radius=0.52bp];
\fill[tg808080] (326.198,80.798) circle[radius=0.52bp];
\fill[tg808080] (327.438,83.337) circle[radius=0.52bp];
\fill[tg808080] (328.677,86.110) circle[radius=0.52bp];
\fill[tg808080] (329.917,81.776) circle[radius=0.52bp];
\fill[tg808080] (331.156,84.574) circle[radius=0.52bp];
\fill[tg808080] (332.396,88.308) circle[radius=0.52bp];
\fill[tg808080] (333.635,89.414) circle[radius=0.52bp];
\fill[tg808080] (334.875,93.212) circle[radius=0.52bp];
\fill[tg808080] (336.115,95.892) circle[radius=0.52bp];
\fill[tg808080] (337.354,97.403) circle[radius=0.52bp];
\fill[tg808080] (338.594,95.484) circle[radius=0.52bp];
\fill[tg808080] (339.833,93.420) circle[radius=0.52bp];
\fill[tg808080] (341.073,91.694) circle[radius=0.52bp];
\fill[tg808080] (342.312,93.244) circle[radius=0.52bp];
\fill[tg808080] (343.552,91.488) circle[radius=0.52bp];
\fill[tg808080] (344.792,91.212) circle[radius=0.52bp];
\fill[tg808080] (346.031,88.601) circle[radius=0.52bp];
\fill[tg808080] (347.271,90.560) circle[radius=0.52bp];
\fill[tg808080] (348.510,87.571) circle[radius=0.52bp];
\fill[tg808080] (349.750,84.874) circle[radius=0.52bp];
\fill[tg808080] (350.990,83.227) circle[radius=0.52bp];
\fill[tg808080] (352.229,80.321) circle[radius=0.52bp];
\fill[tg808080] (353.469,77.987) circle[radius=0.52bp];
\fill[tg808080] (354.708,76.082) circle[radius=0.52bp];
\fill[tg808080] (355.948,73.143) circle[radius=0.52bp];
\fill[tg808080] (357.188,75.773) circle[radius=0.52bp];
\fill[tg808080] (358.427,74.885) circle[radius=0.52bp];
\fill[tg808080] (359.667,73.709) circle[radius=0.52bp];
\fill[tg808080] (360.906,69.518) circle[radius=0.52bp];
\fill[tg808080] (362.146,68.704) circle[radius=0.52bp];
\fill[tg808080] (363.385,76.099) circle[radius=0.52bp];
\fill[tg808080] (364.625,77.256) circle[radius=0.52bp];
\fill[tg808080] (365.865,77.393) circle[radius=0.52bp];
\fill[tg808080] (367.104,81.984) circle[radius=0.52bp];
\fill[tg808080] (368.344,81.656) circle[radius=0.52bp];
\fill[tg808080] (369.583,79.538) circle[radius=0.52bp];
\fill[tg808080] (370.823,74.160) circle[radius=0.52bp];
\fill[tg808080] (372.062,81.818) circle[radius=0.52bp];
\fill[tg808080] (373.302,85.155) circle[radius=0.52bp];
\fill[tg808080] (374.542,85.860) circle[radius=0.52bp];
\fill[tg808080] (375.781,82.845) circle[radius=0.52bp];
\fill[tg808080] (377.021,83.730) circle[radius=0.52bp];
\fill[tg808080] (378.260,81.025) circle[radius=0.52bp];
\fill[tg808080] (379.500,82.652) circle[radius=0.52bp];
\fill[tg808080] (380.740,84.038) circle[radius=0.52bp];
\fill[tg808080] (381.979,82.443) circle[radius=0.52bp];
\fill[tg808080] (383.219,82.923) circle[radius=0.52bp];
\fill[tg808080] (384.458,85.436) circle[radius=0.52bp];
\fill[tg808080] (385.698,81.730) circle[radius=0.52bp];
\fill[tg808080] (386.938,78.304) circle[radius=0.52bp];
\fill[tg808080] (388.177,76.706) circle[radius=0.52bp];
\fill[tg808080] (389.417,75.836) circle[radius=0.52bp];
\fill[tg808080] (390.656,69.737) circle[radius=0.52bp];
\fill[tg808080] (391.896,65.161) circle[radius=0.52bp];
\fill[tg808080] (393.135,65.716) circle[radius=0.52bp];
\fill[tg808080] (394.375,68.045) circle[radius=0.52bp];
\fill[tg808080] (395.615,69.382) circle[radius=0.52bp];
\fill[tg808080] (396.854,69.677) circle[radius=0.52bp];
\fill[tg808080] (398.094,74.287) circle[radius=0.52bp];
\fill[tg808080] (399.333,69.871) circle[radius=0.52bp];
\fill[tg808080] (400.573,70.474) circle[radius=0.52bp];
\fill[tg808080] (401.812,71.465) circle[radius=0.52bp];
\fill[tg808080] (403.052,71.875) circle[radius=0.52bp];
\fill[tg808080] (404.292,71.358) circle[radius=0.52bp];
\fill[tg808080] (405.531,69.293) circle[radius=0.52bp];
\fill[tg808080] (406.771,70.560) circle[radius=0.52bp];
\fill[tg808080] (408.010,70.466) circle[radius=0.52bp];
\fill[tg808080] (409.250,77.191) circle[radius=0.52bp];
\fill[tg808080] (410.490,73.493) circle[radius=0.52bp];
\fill[tg808080] (411.729,70.832) circle[radius=0.52bp];
\fill[tg808080] (412.969,75.673) circle[radius=0.52bp];
\fill[tg808080] (414.208,74.517) circle[radius=0.52bp];
\fill[tg808080] (415.448,74.193) circle[radius=0.52bp];
\fill[tg808080] (416.688,76.218) circle[radius=0.52bp];
\fill[tg808080] (417.927,79.442) circle[radius=0.52bp];
\fill[tg808080] (419.167,78.601) circle[radius=0.52bp];
\fill[tg808080] (420.406,77.775) circle[radius=0.52bp];
\fill[tg808080] (421.646,77.160) circle[radius=0.52bp];
\fill[tg808080] (422.885,81.257) circle[radius=0.52bp];
\fill[tg808080] (424.125,79.865) circle[radius=0.52bp];
\fill[tg808080] (425.365,82.496) circle[radius=0.52bp];
\fill[tg808080] (426.604,87.043) circle[radius=0.52bp];
\fill[tg808080] (427.844,86.598) circle[radius=0.52bp];
\fill[tg808080] (429.083,88.758) circle[radius=0.52bp];
\fill[tg808080] (430.323,88.559) circle[radius=0.52bp];
\fill[tg808080] (431.562,94.444) circle[radius=0.52bp];
\fill[tg808080] (432.802,93.696) circle[radius=0.52bp];
\fill[tg808080] (434.042,92.169) circle[radius=0.52bp];
\fill[tg808080] (435.281,93.035) circle[radius=0.52bp];
\fill[tg808080] (436.521,95.013) circle[radius=0.52bp];
\fill[tg808080] (437.760,99.839) circle[radius=0.52bp];
\draw[color=tg575757,line width=0.65bp,line cap=round] (409.250,70.437) -- (410.490,70.407) -- (411.729,70.377) -- (412.969,70.347) -- (414.208,70.318) -- (415.448,70.288) -- (416.688,70.258) -- (417.927,70.228) -- (419.167,70.199) -- (420.406,70.169) -- (421.646,70.139) -- (422.885,70.109) -- (424.125,70.080) -- (425.365,70.050) -- (426.604,70.020) -- (427.844,69.990) -- (429.083,69.960) -- (430.323,69.931) -- (431.562,69.901) -- (432.802,69.871) -- (434.042,69.841) -- (435.281,69.812) -- (436.521,69.782) -- (437.760,69.752);
\draw[color=tg0000ff,line width=0.38bp,line cap=round] (408.750,69.966) -- (409.750,70.966);
\draw[color=tg0000ff,line width=0.38bp,line cap=round] (408.750,70.966) -- (409.750,69.966);
\draw[color=tg0000ff,line width=0.38bp,line cap=round] (409.990,69.966) -- (410.990,70.966);
\draw[color=tg0000ff,line width=0.38bp,line cap=round] (409.990,70.966) -- (410.990,69.966);
\draw[color=tg0000ff,line width=0.38bp,line cap=round] (411.229,69.966) -- (412.229,70.966);
\draw[color=tg0000ff,line width=0.38bp,line cap=round] (411.229,70.966) -- (412.229,69.966);
\draw[color=tg0000ff,line width=0.38bp,line cap=round] (412.469,69.966) -- (413.469,70.966);
\draw[color=tg0000ff,line width=0.38bp,line cap=round] (412.469,70.966) -- (413.469,69.966);
\draw[color=tg0000ff,line width=0.38bp,line cap=round] (413.708,69.966) -- (414.708,70.966);
\draw[color=tg0000ff,line width=0.38bp,line cap=round] (413.708,70.966) -- (414.708,69.966);
\draw[color=tg0000ff,line width=0.38bp,line cap=round] (414.948,69.966) -- (415.948,70.966);
\draw[color=tg0000ff,line width=0.38bp,line cap=round] (414.948,70.966) -- (415.948,69.966);
\draw[color=tg0000ff,line width=0.38bp,line cap=round] (416.188,69.966) -- (417.188,70.966);
\draw[color=tg0000ff,line width=0.38bp,line cap=round] (416.188,70.966) -- (417.188,69.966);
\draw[color=tg0000ff,line width=0.38bp,line cap=round] (417.427,69.966) -- (418.427,70.966);
\draw[color=tg0000ff,line width=0.38bp,line cap=round] (417.427,70.966) -- (418.427,69.966);
\draw[color=tg0000ff,line width=0.38bp,line cap=round] (418.667,69.966) -- (419.667,70.966);
\draw[color=tg0000ff,line width=0.38bp,line cap=round] (418.667,70.966) -- (419.667,69.966);
\draw[color=tg0000ff,line width=0.38bp,line cap=round] (419.906,69.966) -- (420.906,70.966);
\draw[color=tg0000ff,line width=0.38bp,line cap=round] (419.906,70.966) -- (420.906,69.966);
\draw[color=tg0000ff,line width=0.38bp,line cap=round] (421.146,69.966) -- (422.146,70.966);
\draw[color=tg0000ff,line width=0.38bp,line cap=round] (421.146,70.966) -- (422.146,69.966);
\draw[color=tg0000ff,line width=0.38bp,line cap=round] (422.385,69.966) -- (423.385,70.966);
\draw[color=tg0000ff,line width=0.38bp,line cap=round] (422.385,70.966) -- (423.385,69.966);
\draw[color=tg0000ff,line width=0.38bp,line cap=round] (423.625,69.966) -- (424.625,70.966);
\draw[color=tg0000ff,line width=0.38bp,line cap=round] (423.625,70.966) -- (424.625,69.966);
\draw[color=tg0000ff,line width=0.38bp,line cap=round] (424.865,69.966) -- (425.865,70.966);
\draw[color=tg0000ff,line width=0.38bp,line cap=round] (424.865,70.966) -- (425.865,69.966);
\draw[color=tg0000ff,line width=0.38bp,line cap=round] (426.104,69.966) -- (427.104,70.966);
\draw[color=tg0000ff,line width=0.38bp,line cap=round] (426.104,70.966) -- (427.104,69.966);
\draw[color=tg0000ff,line width=0.38bp,line cap=round] (427.344,69.966) -- (428.344,70.966);
\draw[color=tg0000ff,line width=0.38bp,line cap=round] (427.344,70.966) -- (428.344,69.966);
\draw[color=tg0000ff,line width=0.38bp,line cap=round] (428.583,69.966) -- (429.583,70.966);
\draw[color=tg0000ff,line width=0.38bp,line cap=round] (428.583,70.966) -- (429.583,69.966);
\draw[color=tg0000ff,line width=0.38bp,line cap=round] (429.823,69.966) -- (430.823,70.966);
\draw[color=tg0000ff,line width=0.38bp,line cap=round] (429.823,70.966) -- (430.823,69.966);
\draw[color=tg0000ff,line width=0.38bp,line cap=round] (431.062,69.966) -- (432.062,70.966);
\draw[color=tg0000ff,line width=0.38bp,line cap=round] (431.062,70.966) -- (432.062,69.966);
\draw[color=tg0000ff,line width=0.38bp,line cap=round] (432.302,69.966) -- (433.302,70.966);
\draw[color=tg0000ff,line width=0.38bp,line cap=round] (432.302,70.966) -- (433.302,69.966);
\draw[color=tg0000ff,line width=0.38bp,line cap=round] (433.542,69.966) -- (434.542,70.966);
\draw[color=tg0000ff,line width=0.38bp,line cap=round] (433.542,70.966) -- (434.542,69.966);
\draw[color=tg0000ff,line width=0.38bp,line cap=round] (434.781,69.966) -- (435.781,70.966);
\draw[color=tg0000ff,line width=0.38bp,line cap=round] (434.781,70.966) -- (435.781,69.966);
\draw[color=tg0000ff,line width=0.38bp,line cap=round] (436.021,69.966) -- (437.021,70.966);
\draw[color=tg0000ff,line width=0.38bp,line cap=round] (436.021,70.966) -- (437.021,69.966);
\draw[color=tg0000ff,line width=0.38bp,line cap=round] (437.260,69.966) -- (438.260,70.966);
\draw[color=tg0000ff,line width=0.38bp,line cap=round] (437.260,70.966) -- (438.260,69.966);
\node[anchor=center,inner sep=0pt,outer sep=0pt,text=tg18314b,font={\sffamily\fontsize{7.4bp}{8.88bp}\selectfont\bfseries}] at (379.500,40.000) {Seed 33};
\node[anchor=west,inner sep=0pt,outer sep=0pt,text=tg18314b,font={\sffamily\fontsize{8bp}{9.6bp}\selectfont\bfseries}] at (0.000,119.000) {Merton jump diffusion (log price)};
\node[anchor=east,inner sep=0pt,outer sep=0pt,text=tg626b73,font={\sffamily\fontsize{7bp}{8.4bp}\selectfont}] at (448.000,119.000) {Compound Poisson jumps};
\draw[color=tgc9cdd1,line width=0.45bp,line cap=round] (28.000,128.000) -- (28.000,171.000);
\draw[color=tgc9cdd1,line width=0.45bp,line cap=round] (28.000,171.000) -- (147.000,171.000);
\node[anchor=east,inner sep=0pt,outer sep=0pt,text=tg626b73,font={\sffamily\fontsize{6.7bp}{8.04bp}\selectfont}] at (25.000,164.980) {4.23};
\node[anchor=east,inner sep=0pt,outer sep=0pt,text=tg626b73,font={\sffamily\fontsize{6.7bp}{8.04bp}\selectfont}] at (25.000,134.020) {4.74};
\draw[color=tg9da7b0,line width=0.6bp,line cap=round,dash pattern=on 2bp off 2bp] (117.250,128.000) -- (117.250,171.000);
\fill[tg808080] (28.000,134.365) circle[radius=0.52bp];
\fill[tg808080] (29.240,134.904) circle[radius=0.52bp];
\fill[tg808080] (30.479,135.828) circle[radius=0.52bp];
\fill[tg808080] (31.719,135.190) circle[radius=0.52bp];
\fill[tg808080] (32.958,135.400) circle[radius=0.52bp];
\fill[tg808080] (34.198,135.732) circle[radius=0.52bp];
\fill[tg808080] (35.438,135.388) circle[radius=0.52bp];
\fill[tg808080] (36.677,135.619) circle[radius=0.52bp];
\fill[tg808080] (37.917,136.223) circle[radius=0.52bp];
\fill[tg808080] (39.156,135.462) circle[radius=0.52bp];
\fill[tg808080] (40.396,135.223) circle[radius=0.52bp];
\fill[tg808080] (41.635,135.281) circle[radius=0.52bp];
\fill[tg808080] (42.875,134.858) circle[radius=0.52bp];
\fill[tg808080] (44.115,135.119) circle[radius=0.52bp];
\fill[tg808080] (45.354,135.160) circle[radius=0.52bp];
\fill[tg808080] (46.594,135.034) circle[radius=0.52bp];
\fill[tg808080] (47.833,133.582) circle[radius=0.52bp];
\fill[tg808080] (49.073,133.677) circle[radius=0.52bp];
\fill[tg808080] (50.312,132.823) circle[radius=0.52bp];
\fill[tg808080] (51.552,132.705) circle[radius=0.52bp];
\fill[tg808080] (52.792,133.054) circle[radius=0.52bp];
\fill[tg808080] (54.031,133.587) circle[radius=0.52bp];
\fill[tg808080] (55.271,133.708) circle[radius=0.52bp];
\fill[tg808080] (56.510,133.948) circle[radius=0.52bp];
\fill[tg808080] (57.750,134.780) circle[radius=0.52bp];
\fill[tg808080] (58.990,134.639) circle[radius=0.52bp];
\fill[tg808080] (60.229,133.631) circle[radius=0.52bp];
\fill[tg808080] (61.469,133.245) circle[radius=0.52bp];
\fill[tg808080] (62.708,133.843) circle[radius=0.52bp];
\fill[tg808080] (63.948,133.195) circle[radius=0.52bp];
\fill[tg808080] (65.188,133.601) circle[radius=0.52bp];
\fill[tg808080] (66.427,134.468) circle[radius=0.52bp];
\fill[tg808080] (67.667,134.446) circle[radius=0.52bp];
\fill[tg808080] (68.906,134.824) circle[radius=0.52bp];
\fill[tg808080] (70.146,135.267) circle[radius=0.52bp];
\fill[tg808080] (71.385,134.981) circle[radius=0.52bp];
\fill[tg808080] (72.625,134.581) circle[radius=0.52bp];
\fill[tg808080] (73.865,134.478) circle[radius=0.52bp];
\fill[tg808080] (75.104,134.967) circle[radius=0.52bp];
\fill[tg808080] (76.344,135.087) circle[radius=0.52bp];
\fill[tg808080] (77.583,134.772) circle[radius=0.52bp];
\fill[tg808080] (78.823,134.063) circle[radius=0.52bp];
\fill[tg808080] (80.062,133.720) circle[radius=0.52bp];
\fill[tg808080] (81.302,133.028) circle[radius=0.52bp];
\fill[tg808080] (82.542,132.705) circle[radius=0.52bp];
\fill[tg808080] (83.781,132.161) circle[radius=0.52bp];
\fill[tg808080] (85.021,132.479) circle[radius=0.52bp];
\fill[tg808080] (86.260,132.372) circle[radius=0.52bp];
\fill[tg808080] (87.500,133.182) circle[radius=0.52bp];
\fill[tg808080] (88.740,133.046) circle[radius=0.52bp];
\fill[tg808080] (89.979,132.966) circle[radius=0.52bp];
\fill[tg808080] (91.219,133.421) circle[radius=0.52bp];
\fill[tg808080] (92.458,133.588) circle[radius=0.52bp];
\fill[tg808080] (93.698,134.291) circle[radius=0.52bp];
\fill[tg808080] (94.938,134.173) circle[radius=0.52bp];
\fill[tg808080] (96.177,133.646) circle[radius=0.52bp];
\fill[tg808080] (97.417,133.235) circle[radius=0.52bp];
\fill[tg808080] (98.656,132.405) circle[radius=0.52bp];
\fill[tg808080] (99.896,132.268) circle[radius=0.52bp];
\fill[tg808080] (101.135,149.192) circle[radius=0.52bp];
\fill[tg808080] (102.375,149.120) circle[radius=0.52bp];
\fill[tg808080] (103.615,148.701) circle[radius=0.52bp];
\fill[tg808080] (104.854,149.855) circle[radius=0.52bp];
\fill[tg808080] (106.094,150.038) circle[radius=0.52bp];
\fill[tg808080] (107.333,149.934) circle[radius=0.52bp];
\fill[tg808080] (108.573,149.800) circle[radius=0.52bp];
\fill[tg808080] (109.812,149.606) circle[radius=0.52bp];
\fill[tg808080] (111.052,148.932) circle[radius=0.52bp];
\fill[tg808080] (112.292,148.760) circle[radius=0.52bp];
\fill[tg808080] (113.531,148.301) circle[radius=0.52bp];
\fill[tg808080] (114.771,148.593) circle[radius=0.52bp];
\fill[tg808080] (116.010,149.417) circle[radius=0.52bp];
\fill[tg808080] (117.250,149.500) circle[radius=0.52bp];
\fill[tg808080] (118.490,149.188) circle[radius=0.52bp];
\fill[tg808080] (119.729,149.047) circle[radius=0.52bp];
\fill[tg808080] (120.969,149.440) circle[radius=0.52bp];
\fill[tg808080] (122.208,149.025) circle[radius=0.52bp];
\fill[tg808080] (123.448,149.377) circle[radius=0.52bp];
\fill[tg808080] (124.688,150.702) circle[radius=0.52bp];
\fill[tg808080] (125.927,166.264) circle[radius=0.52bp];
\fill[tg808080] (127.167,166.456) circle[radius=0.52bp];
\fill[tg808080] (128.406,166.053) circle[radius=0.52bp];
\fill[tg808080] (129.646,166.839) circle[radius=0.52bp];
\fill[tg808080] (130.885,166.666) circle[radius=0.52bp];
\fill[tg808080] (132.125,166.715) circle[radius=0.52bp];
\fill[tg808080] (133.365,166.309) circle[radius=0.52bp];
\fill[tg808080] (134.604,166.304) circle[radius=0.52bp];
\fill[tg808080] (135.844,166.248) circle[radius=0.52bp];
\fill[tg808080] (137.083,166.296) circle[radius=0.52bp];
\fill[tg808080] (138.323,165.583) circle[radius=0.52bp];
\fill[tg808080] (139.562,165.430) circle[radius=0.52bp];
\fill[tg808080] (140.802,164.898) circle[radius=0.52bp];
\fill[tg808080] (142.042,164.652) circle[radius=0.52bp];
\fill[tg808080] (143.281,164.421) circle[radius=0.52bp];
\fill[tg808080] (144.521,164.306) circle[radius=0.52bp];
\fill[tg808080] (145.760,164.634) circle[radius=0.52bp];
\draw[color=tg575757,line width=0.65bp,line cap=round] (117.250,149.422) -- (118.490,149.428) -- (119.729,149.433) -- (120.969,149.438) -- (122.208,149.443) -- (123.448,149.449) -- (124.688,149.454) -- (125.927,149.459) -- (127.167,149.464) -- (128.406,149.469) -- (129.646,149.475) -- (130.885,149.480) -- (132.125,149.485) -- (133.365,149.490) -- (134.604,149.496) -- (135.844,149.501) -- (137.083,149.506) -- (138.323,149.511) -- (139.562,149.517) -- (140.802,149.522) -- (142.042,149.527) -- (143.281,149.532) -- (144.521,149.537) -- (145.760,149.543);
\draw[color=tg0000ff,line width=0.38bp,line cap=round] (116.750,148.917) -- (117.750,149.917);
\draw[color=tg0000ff,line width=0.38bp,line cap=round] (116.750,149.917) -- (117.750,148.917);
\draw[color=tg0000ff,line width=0.38bp,line cap=round] (117.990,148.917) -- (118.990,149.917);
\draw[color=tg0000ff,line width=0.38bp,line cap=round] (117.990,149.917) -- (118.990,148.917);
\draw[color=tg0000ff,line width=0.38bp,line cap=round] (119.229,148.917) -- (120.229,149.917);
\draw[color=tg0000ff,line width=0.38bp,line cap=round] (119.229,149.917) -- (120.229,148.917);
\draw[color=tg0000ff,line width=0.38bp,line cap=round] (120.469,148.917) -- (121.469,149.917);
\draw[color=tg0000ff,line width=0.38bp,line cap=round] (120.469,149.917) -- (121.469,148.917);
\draw[color=tg0000ff,line width=0.38bp,line cap=round] (121.708,148.917) -- (122.708,149.917);
\draw[color=tg0000ff,line width=0.38bp,line cap=round] (121.708,149.917) -- (122.708,148.917);
\draw[color=tg0000ff,line width=0.38bp,line cap=round] (122.948,148.917) -- (123.948,149.917);
\draw[color=tg0000ff,line width=0.38bp,line cap=round] (122.948,149.917) -- (123.948,148.917);
\draw[color=tg0000ff,line width=0.38bp,line cap=round] (124.188,148.917) -- (125.188,149.917);
\draw[color=tg0000ff,line width=0.38bp,line cap=round] (124.188,149.917) -- (125.188,148.917);
\draw[color=tg0000ff,line width=0.38bp,line cap=round] (125.427,148.917) -- (126.427,149.917);
\draw[color=tg0000ff,line width=0.38bp,line cap=round] (125.427,149.917) -- (126.427,148.917);
\draw[color=tg0000ff,line width=0.38bp,line cap=round] (126.667,148.917) -- (127.667,149.917);
\draw[color=tg0000ff,line width=0.38bp,line cap=round] (126.667,149.917) -- (127.667,148.917);
\draw[color=tg0000ff,line width=0.38bp,line cap=round] (127.906,148.917) -- (128.906,149.917);
\draw[color=tg0000ff,line width=0.38bp,line cap=round] (127.906,149.917) -- (128.906,148.917);
\draw[color=tg0000ff,line width=0.38bp,line cap=round] (129.146,148.917) -- (130.146,149.917);
\draw[color=tg0000ff,line width=0.38bp,line cap=round] (129.146,149.917) -- (130.146,148.917);
\draw[color=tg0000ff,line width=0.38bp,line cap=round] (130.385,148.917) -- (131.385,149.917);
\draw[color=tg0000ff,line width=0.38bp,line cap=round] (130.385,149.917) -- (131.385,148.917);
\draw[color=tg0000ff,line width=0.38bp,line cap=round] (131.625,148.917) -- (132.625,149.917);
\draw[color=tg0000ff,line width=0.38bp,line cap=round] (131.625,149.917) -- (132.625,148.917);
\draw[color=tg0000ff,line width=0.38bp,line cap=round] (132.865,148.917) -- (133.865,149.917);
\draw[color=tg0000ff,line width=0.38bp,line cap=round] (132.865,149.917) -- (133.865,148.917);
\draw[color=tg0000ff,line width=0.38bp,line cap=round] (134.104,148.917) -- (135.104,149.917);
\draw[color=tg0000ff,line width=0.38bp,line cap=round] (134.104,149.917) -- (135.104,148.917);
\draw[color=tg0000ff,line width=0.38bp,line cap=round] (135.344,148.917) -- (136.344,149.917);
\draw[color=tg0000ff,line width=0.38bp,line cap=round] (135.344,149.917) -- (136.344,148.917);
\draw[color=tg0000ff,line width=0.38bp,line cap=round] (136.583,148.917) -- (137.583,149.917);
\draw[color=tg0000ff,line width=0.38bp,line cap=round] (136.583,149.917) -- (137.583,148.917);
\draw[color=tg0000ff,line width=0.38bp,line cap=round] (137.823,148.917) -- (138.823,149.917);
\draw[color=tg0000ff,line width=0.38bp,line cap=round] (137.823,149.917) -- (138.823,148.917);
\draw[color=tg0000ff,line width=0.38bp,line cap=round] (139.062,148.917) -- (140.062,149.917);
\draw[color=tg0000ff,line width=0.38bp,line cap=round] (139.062,149.917) -- (140.062,148.917);
\draw[color=tg0000ff,line width=0.38bp,line cap=round] (140.302,148.917) -- (141.302,149.917);
\draw[color=tg0000ff,line width=0.38bp,line cap=round] (140.302,149.917) -- (141.302,148.917);
\draw[color=tg0000ff,line width=0.38bp,line cap=round] (141.542,148.917) -- (142.542,149.917);
\draw[color=tg0000ff,line width=0.38bp,line cap=round] (141.542,149.917) -- (142.542,148.917);
\draw[color=tg0000ff,line width=0.38bp,line cap=round] (142.781,148.917) -- (143.781,149.917);
\draw[color=tg0000ff,line width=0.38bp,line cap=round] (142.781,149.917) -- (143.781,148.917);
\draw[color=tg0000ff,line width=0.38bp,line cap=round] (144.021,148.917) -- (145.021,149.917);
\draw[color=tg0000ff,line width=0.38bp,line cap=round] (144.021,149.917) -- (145.021,148.917);
\draw[color=tg0000ff,line width=0.38bp,line cap=round] (145.260,148.917) -- (146.260,149.917);
\draw[color=tg0000ff,line width=0.38bp,line cap=round] (145.260,149.917) -- (146.260,148.917);
\draw[color=tgc9cdd1,line width=0.45bp,line cap=round] (174.000,128.000) -- (174.000,171.000);
\draw[color=tgc9cdd1,line width=0.45bp,line cap=round] (174.000,171.000) -- (293.000,171.000);
\node[anchor=east,inner sep=0pt,outer sep=0pt,text=tg626b73,font={\sffamily\fontsize{6.7bp}{8.04bp}\selectfont}] at (171.000,164.980) {4.49};
\node[anchor=east,inner sep=0pt,outer sep=0pt,text=tg626b73,font={\sffamily\fontsize{6.7bp}{8.04bp}\selectfont}] at (171.000,134.020) {4.58};
\draw[color=tg9da7b0,line width=0.6bp,line cap=round,dash pattern=on 2bp off 2bp] (263.250,128.000) -- (263.250,171.000);
\fill[tg808080] (174.000,161.430) circle[radius=0.52bp];
\fill[tg808080] (175.240,166.839) circle[radius=0.52bp];
\fill[tg808080] (176.479,156.192) circle[radius=0.52bp];
\fill[tg808080] (177.719,158.639) circle[radius=0.52bp];
\fill[tg808080] (178.958,160.512) circle[radius=0.52bp];
\fill[tg808080] (180.198,159.526) circle[radius=0.52bp];
\fill[tg808080] (181.438,155.131) circle[radius=0.52bp];
\fill[tg808080] (182.677,147.809) circle[radius=0.52bp];
\fill[tg808080] (183.917,147.627) circle[radius=0.52bp];
\fill[tg808080] (185.156,142.194) circle[radius=0.52bp];
\fill[tg808080] (186.396,138.907) circle[radius=0.52bp];
\fill[tg808080] (187.635,137.767) circle[radius=0.52bp];
\fill[tg808080] (188.875,136.680) circle[radius=0.52bp];
\fill[tg808080] (190.115,139.120) circle[radius=0.52bp];
\fill[tg808080] (191.354,140.523) circle[radius=0.52bp];
\fill[tg808080] (192.594,137.904) circle[radius=0.52bp];
\fill[tg808080] (193.833,136.504) circle[radius=0.52bp];
\fill[tg808080] (195.073,138.490) circle[radius=0.52bp];
\fill[tg808080] (196.312,144.629) circle[radius=0.52bp];
\fill[tg808080] (197.552,139.747) circle[radius=0.52bp];
\fill[tg808080] (198.792,146.017) circle[radius=0.52bp];
\fill[tg808080] (200.031,142.815) circle[radius=0.52bp];
\fill[tg808080] (201.271,143.858) circle[radius=0.52bp];
\fill[tg808080] (202.510,143.378) circle[radius=0.52bp];
\fill[tg808080] (203.750,139.945) circle[radius=0.52bp];
\fill[tg808080] (204.990,136.947) circle[radius=0.52bp];
\fill[tg808080] (206.229,136.723) circle[radius=0.52bp];
\fill[tg808080] (207.469,132.161) circle[radius=0.52bp];
\fill[tg808080] (208.708,133.348) circle[radius=0.52bp];
\fill[tg808080] (209.948,132.816) circle[radius=0.52bp];
\fill[tg808080] (211.188,137.125) circle[radius=0.52bp];
\fill[tg808080] (212.427,135.089) circle[radius=0.52bp];
\fill[tg808080] (213.667,134.221) circle[radius=0.52bp];
\fill[tg808080] (214.906,134.944) circle[radius=0.52bp];
\fill[tg808080] (216.146,135.302) circle[radius=0.52bp];
\fill[tg808080] (217.385,136.349) circle[radius=0.52bp];
\fill[tg808080] (218.625,140.927) circle[radius=0.52bp];
\fill[tg808080] (219.865,138.047) circle[radius=0.52bp];
\fill[tg808080] (221.104,140.220) circle[radius=0.52bp];
\fill[tg808080] (222.344,140.846) circle[radius=0.52bp];
\fill[tg808080] (223.583,139.062) circle[radius=0.52bp];
\fill[tg808080] (224.823,137.589) circle[radius=0.52bp];
\fill[tg808080] (226.062,139.626) circle[radius=0.52bp];
\fill[tg808080] (227.302,142.935) circle[radius=0.52bp];
\fill[tg808080] (228.542,138.560) circle[radius=0.52bp];
\fill[tg808080] (229.781,141.874) circle[radius=0.52bp];
\fill[tg808080] (231.021,145.141) circle[radius=0.52bp];
\fill[tg808080] (232.260,145.598) circle[radius=0.52bp];
\fill[tg808080] (233.500,143.093) circle[radius=0.52bp];
\fill[tg808080] (234.740,141.488) circle[radius=0.52bp];
\fill[tg808080] (235.979,135.917) circle[radius=0.52bp];
\fill[tg808080] (237.219,139.886) circle[radius=0.52bp];
\fill[tg808080] (238.458,136.647) circle[radius=0.52bp];
\fill[tg808080] (239.698,138.214) circle[radius=0.52bp];
\fill[tg808080] (240.938,142.027) circle[radius=0.52bp];
\fill[tg808080] (242.177,141.856) circle[radius=0.52bp];
\fill[tg808080] (243.417,150.185) circle[radius=0.52bp];
\fill[tg808080] (244.656,147.434) circle[radius=0.52bp];
\fill[tg808080] (245.896,147.038) circle[radius=0.52bp];
\fill[tg808080] (247.135,148.866) circle[radius=0.52bp];
\fill[tg808080] (248.375,145.397) circle[radius=0.52bp];
\fill[tg808080] (249.615,151.822) circle[radius=0.52bp];
\fill[tg808080] (250.854,151.816) circle[radius=0.52bp];
\fill[tg808080] (252.094,151.346) circle[radius=0.52bp];
\fill[tg808080] (253.333,149.848) circle[radius=0.52bp];
\fill[tg808080] (254.573,150.887) circle[radius=0.52bp];
\fill[tg808080] (255.812,154.894) circle[radius=0.52bp];
\fill[tg808080] (257.052,158.485) circle[radius=0.52bp];
\fill[tg808080] (258.292,163.168) circle[radius=0.52bp];
\fill[tg808080] (259.531,163.048) circle[radius=0.52bp];
\fill[tg808080] (260.771,157.716) circle[radius=0.52bp];
\fill[tg808080] (262.010,158.384) circle[radius=0.52bp];
\fill[tg808080] (263.250,155.131) circle[radius=0.52bp];
\fill[tg808080] (264.490,157.726) circle[radius=0.52bp];
\fill[tg808080] (265.729,157.496) circle[radius=0.52bp];
\fill[tg808080] (266.969,157.913) circle[radius=0.52bp];
\fill[tg808080] (268.208,157.522) circle[radius=0.52bp];
\fill[tg808080] (269.448,159.621) circle[radius=0.52bp];
\fill[tg808080] (270.688,164.072) circle[radius=0.52bp];
\fill[tg808080] (271.927,161.628) circle[radius=0.52bp];
\fill[tg808080] (273.167,155.428) circle[radius=0.52bp];
\fill[tg808080] (274.406,155.823) circle[radius=0.52bp];
\fill[tg808080] (275.646,156.406) circle[radius=0.52bp];
\fill[tg808080] (276.885,152.039) circle[radius=0.52bp];
\fill[tg808080] (278.125,150.989) circle[radius=0.52bp];
\fill[tg808080] (279.365,155.102) circle[radius=0.52bp];
\fill[tg808080] (280.604,156.300) circle[radius=0.52bp];
\fill[tg808080] (281.844,154.147) circle[radius=0.52bp];
\fill[tg808080] (283.083,153.915) circle[radius=0.52bp];
\fill[tg808080] (284.323,152.697) circle[radius=0.52bp];
\fill[tg808080] (285.562,145.465) circle[radius=0.52bp];
\fill[tg808080] (286.802,143.382) circle[radius=0.52bp];
\fill[tg808080] (288.042,143.040) circle[radius=0.52bp];
\fill[tg808080] (289.281,140.964) circle[radius=0.52bp];
\fill[tg808080] (290.521,140.226) circle[radius=0.52bp];
\fill[tg808080] (291.760,142.662) circle[radius=0.52bp];
\draw[color=tg575757,line width=0.65bp,line cap=round] (263.250,158.413) -- (264.490,158.442) -- (265.729,158.471) -- (266.969,158.500) -- (268.208,158.530) -- (269.448,158.559) -- (270.688,158.588) -- (271.927,158.617) -- (273.167,158.647) -- (274.406,158.676) -- (275.646,158.705) -- (276.885,158.734) -- (278.125,158.764) -- (279.365,158.793) -- (280.604,158.822) -- (281.844,158.851) -- (283.083,158.881) -- (284.323,158.910) -- (285.562,158.939) -- (286.802,158.968) -- (288.042,158.998) -- (289.281,159.027) -- (290.521,159.056) -- (291.760,159.085);
\draw[color=tg0000ff,line width=0.38bp,line cap=round] (262.750,157.884) -- (263.750,158.884);
\draw[color=tg0000ff,line width=0.38bp,line cap=round] (262.750,158.884) -- (263.750,157.884);
\draw[color=tg0000ff,line width=0.38bp,line cap=round] (263.990,157.884) -- (264.990,158.884);
\draw[color=tg0000ff,line width=0.38bp,line cap=round] (263.990,158.884) -- (264.990,157.884);
\draw[color=tg0000ff,line width=0.38bp,line cap=round] (265.229,157.884) -- (266.229,158.884);
\draw[color=tg0000ff,line width=0.38bp,line cap=round] (265.229,158.884) -- (266.229,157.884);
\draw[color=tg0000ff,line width=0.38bp,line cap=round] (266.469,157.884) -- (267.469,158.884);
\draw[color=tg0000ff,line width=0.38bp,line cap=round] (266.469,158.884) -- (267.469,157.884);
\draw[color=tg0000ff,line width=0.38bp,line cap=round] (267.708,157.884) -- (268.708,158.884);
\draw[color=tg0000ff,line width=0.38bp,line cap=round] (267.708,158.884) -- (268.708,157.884);
\draw[color=tg0000ff,line width=0.38bp,line cap=round] (268.948,157.884) -- (269.948,158.884);
\draw[color=tg0000ff,line width=0.38bp,line cap=round] (268.948,158.884) -- (269.948,157.884);
\draw[color=tg0000ff,line width=0.38bp,line cap=round] (270.188,157.884) -- (271.188,158.884);
\draw[color=tg0000ff,line width=0.38bp,line cap=round] (270.188,158.884) -- (271.188,157.884);
\draw[color=tg0000ff,line width=0.38bp,line cap=round] (271.427,157.884) -- (272.427,158.884);
\draw[color=tg0000ff,line width=0.38bp,line cap=round] (271.427,158.884) -- (272.427,157.884);
\draw[color=tg0000ff,line width=0.38bp,line cap=round] (272.667,157.884) -- (273.667,158.884);
\draw[color=tg0000ff,line width=0.38bp,line cap=round] (272.667,158.884) -- (273.667,157.884);
\draw[color=tg0000ff,line width=0.38bp,line cap=round] (273.906,157.884) -- (274.906,158.884);
\draw[color=tg0000ff,line width=0.38bp,line cap=round] (273.906,158.884) -- (274.906,157.884);
\draw[color=tg0000ff,line width=0.38bp,line cap=round] (275.146,157.884) -- (276.146,158.884);
\draw[color=tg0000ff,line width=0.38bp,line cap=round] (275.146,158.884) -- (276.146,157.884);
\draw[color=tg0000ff,line width=0.38bp,line cap=round] (276.385,157.884) -- (277.385,158.884);
\draw[color=tg0000ff,line width=0.38bp,line cap=round] (276.385,158.884) -- (277.385,157.884);
\draw[color=tg0000ff,line width=0.38bp,line cap=round] (277.625,157.884) -- (278.625,158.884);
\draw[color=tg0000ff,line width=0.38bp,line cap=round] (277.625,158.884) -- (278.625,157.884);
\draw[color=tg0000ff,line width=0.38bp,line cap=round] (278.865,157.884) -- (279.865,158.884);
\draw[color=tg0000ff,line width=0.38bp,line cap=round] (278.865,158.884) -- (279.865,157.884);
\draw[color=tg0000ff,line width=0.38bp,line cap=round] (280.104,157.884) -- (281.104,158.884);
\draw[color=tg0000ff,line width=0.38bp,line cap=round] (280.104,158.884) -- (281.104,157.884);
\draw[color=tg0000ff,line width=0.38bp,line cap=round] (281.344,157.884) -- (282.344,158.884);
\draw[color=tg0000ff,line width=0.38bp,line cap=round] (281.344,158.884) -- (282.344,157.884);
\draw[color=tg0000ff,line width=0.38bp,line cap=round] (282.583,157.884) -- (283.583,158.884);
\draw[color=tg0000ff,line width=0.38bp,line cap=round] (282.583,158.884) -- (283.583,157.884);
\draw[color=tg0000ff,line width=0.38bp,line cap=round] (283.823,157.884) -- (284.823,158.884);
\draw[color=tg0000ff,line width=0.38bp,line cap=round] (283.823,158.884) -- (284.823,157.884);
\draw[color=tg0000ff,line width=0.38bp,line cap=round] (285.062,157.884) -- (286.062,158.884);
\draw[color=tg0000ff,line width=0.38bp,line cap=round] (285.062,158.884) -- (286.062,157.884);
\draw[color=tg0000ff,line width=0.38bp,line cap=round] (286.302,157.884) -- (287.302,158.884);
\draw[color=tg0000ff,line width=0.38bp,line cap=round] (286.302,158.884) -- (287.302,157.884);
\draw[color=tg0000ff,line width=0.38bp,line cap=round] (287.542,157.884) -- (288.542,158.884);
\draw[color=tg0000ff,line width=0.38bp,line cap=round] (287.542,158.884) -- (288.542,157.884);
\draw[color=tg0000ff,line width=0.38bp,line cap=round] (288.781,157.884) -- (289.781,158.884);
\draw[color=tg0000ff,line width=0.38bp,line cap=round] (288.781,158.884) -- (289.781,157.884);
\draw[color=tg0000ff,line width=0.38bp,line cap=round] (290.021,157.884) -- (291.021,158.884);
\draw[color=tg0000ff,line width=0.38bp,line cap=round] (290.021,158.884) -- (291.021,157.884);
\draw[color=tg0000ff,line width=0.38bp,line cap=round] (291.260,157.884) -- (292.260,158.884);
\draw[color=tg0000ff,line width=0.38bp,line cap=round] (291.260,158.884) -- (292.260,157.884);
\draw[color=tgc9cdd1,line width=0.45bp,line cap=round] (320.000,128.000) -- (320.000,171.000);
\draw[color=tgc9cdd1,line width=0.45bp,line cap=round] (320.000,171.000) -- (439.000,171.000);
\node[anchor=east,inner sep=0pt,outer sep=0pt,text=tg626b73,font={\sffamily\fontsize{6.7bp}{8.04bp}\selectfont}] at (317.000,164.980) {4.36};
\node[anchor=east,inner sep=0pt,outer sep=0pt,text=tg626b73,font={\sffamily\fontsize{6.7bp}{8.04bp}\selectfont}] at (317.000,134.020) {4.61};
\draw[color=tg9da7b0,line width=0.6bp,line cap=round,dash pattern=on 2bp off 2bp] (409.250,128.000) -- (409.250,171.000);
\fill[tg808080] (320.000,133.615) circle[radius=0.52bp];
\fill[tg808080] (321.240,133.113) circle[radius=0.52bp];
\fill[tg808080] (322.479,132.161) circle[radius=0.52bp];
\fill[tg808080] (323.719,151.450) circle[radius=0.52bp];
\fill[tg808080] (324.958,150.468) circle[radius=0.52bp];
\fill[tg808080] (326.198,151.121) circle[radius=0.52bp];
\fill[tg808080] (327.438,152.153) circle[radius=0.52bp];
\fill[tg808080] (328.677,153.293) circle[radius=0.52bp];
\fill[tg808080] (329.917,151.191) circle[radius=0.52bp];
\fill[tg808080] (331.156,152.341) circle[radius=0.52bp];
\fill[tg808080] (332.396,153.919) circle[radius=0.52bp];
\fill[tg808080] (333.635,154.298) circle[radius=0.52bp];
\fill[tg808080] (334.875,155.905) circle[radius=0.52bp];
\fill[tg808080] (336.115,152.804) circle[radius=0.52bp];
\fill[tg808080] (337.354,153.368) circle[radius=0.52bp];
\fill[tg808080] (338.594,152.367) circle[radius=0.52bp];
\fill[tg808080] (339.833,151.301) circle[radius=0.52bp];
\fill[tg808080] (341.073,150.388) circle[radius=0.52bp];
\fill[tg808080] (342.312,150.969) circle[radius=0.52bp];
\fill[tg808080] (343.552,166.839) circle[radius=0.52bp];
\fill[tg808080] (344.792,166.587) circle[radius=0.52bp];
\fill[tg808080] (346.031,165.271) circle[radius=0.52bp];
\fill[tg808080] (347.271,166.040) circle[radius=0.52bp];
\fill[tg808080] (348.510,164.551) circle[radius=0.52bp];
\fill[tg808080] (349.750,163.195) circle[radius=0.52bp];
\fill[tg808080] (350.990,162.319) circle[radius=0.52bp];
\fill[tg808080] (352.229,160.868) circle[radius=0.52bp];
\fill[tg808080] (353.469,159.678) circle[radius=0.52bp];
\fill[tg808080] (354.708,158.684) circle[radius=0.52bp];
\fill[tg808080] (355.948,157.218) circle[radius=0.52bp];
\fill[tg808080] (357.188,158.292) circle[radius=0.52bp];
\fill[tg808080] (358.427,157.762) circle[radius=0.52bp];
\fill[tg808080] (359.667,157.100) circle[radius=0.52bp];
\fill[tg808080] (360.906,155.063) circle[radius=0.52bp];
\fill[tg808080] (362.146,154.566) circle[radius=0.52bp];
\fill[tg808080] (363.385,157.813) circle[radius=0.52bp];
\fill[tg808080] (364.625,158.215) circle[radius=0.52bp];
\fill[tg808080] (365.865,158.152) circle[radius=0.52bp];
\fill[tg808080] (367.104,160.121) circle[radius=0.52bp];
\fill[tg808080] (368.344,159.846) circle[radius=0.52bp];
\fill[tg808080] (369.583,158.755) circle[radius=0.52bp];
\fill[tg808080] (370.823,156.176) circle[radius=0.52bp];
\fill[tg808080] (372.062,159.543) circle[radius=0.52bp];
\fill[tg808080] (373.302,160.940) circle[radius=0.52bp];
\fill[tg808080] (374.542,161.136) circle[radius=0.52bp];
\fill[tg808080] (375.781,159.635) circle[radius=0.52bp];
\fill[tg808080] (377.021,159.914) circle[radius=0.52bp];
\fill[tg808080] (378.260,158.555) circle[radius=0.52bp];
\fill[tg808080] (379.500,159.799) circle[radius=0.52bp];
\fill[tg808080] (380.740,160.306) circle[radius=0.52bp];
\fill[tg808080] (381.979,159.453) circle[radius=0.52bp];
\fill[tg808080] (383.219,159.547) circle[radius=0.52bp];
\fill[tg808080] (384.458,160.567) circle[radius=0.52bp];
\fill[tg808080] (385.698,158.752) circle[radius=0.52bp];
\fill[tg808080] (386.938,157.063) circle[radius=0.52bp];
\fill[tg808080] (388.177,156.209) circle[radius=0.52bp];
\fill[tg808080] (389.417,155.687) circle[radius=0.52bp];
\fill[tg808080] (390.656,152.780) circle[radius=0.52bp];
\fill[tg808080] (391.896,150.568) circle[radius=0.52bp];
\fill[tg808080] (393.135,150.695) circle[radius=0.52bp];
\fill[tg808080] (394.375,151.632) circle[radius=0.52bp];
\fill[tg808080] (395.615,152.116) circle[radius=0.52bp];
\fill[tg808080] (396.854,152.125) circle[radius=0.52bp];
\fill[tg808080] (398.094,154.103) circle[radius=0.52bp];
\fill[tg808080] (399.333,151.963) circle[radius=0.52bp];
\fill[tg808080] (400.573,152.113) circle[radius=0.52bp];
\fill[tg808080] (401.812,152.439) circle[radius=0.52bp];
\fill[tg808080] (403.052,152.501) circle[radius=0.52bp];
\fill[tg808080] (404.292,152.140) circle[radius=0.52bp];
\fill[tg808080] (405.531,151.072) circle[radius=0.52bp];
\fill[tg808080] (406.771,151.524) circle[radius=0.52bp];
\fill[tg808080] (408.010,151.356) circle[radius=0.52bp];
\fill[tg808080] (409.250,154.298) circle[radius=0.52bp];
\fill[tg808080] (410.490,152.486) circle[radius=0.52bp];
\fill[tg808080] (411.729,151.147) circle[radius=0.52bp];
\fill[tg808080] (412.969,153.229) circle[radius=0.52bp];
\fill[tg808080] (414.208,152.577) circle[radius=0.52bp];
\fill[tg808080] (415.448,152.304) circle[radius=0.52bp];
\fill[tg808080] (416.688,153.102) circle[radius=0.52bp];
\fill[tg808080] (417.927,154.447) circle[radius=0.52bp];
\fill[tg808080] (419.167,153.938) circle[radius=0.52bp];
\fill[tg808080] (420.406,153.435) circle[radius=0.52bp];
\fill[tg808080] (421.646,153.030) circle[radius=0.52bp];
\fill[tg808080] (422.885,154.772) circle[radius=0.52bp];
\fill[tg808080] (424.125,154.012) circle[radius=0.52bp];
\fill[tg808080] (425.365,155.087) circle[radius=0.52bp];
\fill[tg808080] (426.604,157.035) circle[radius=0.52bp];
\fill[tg808080] (427.844,156.707) circle[radius=0.52bp];
\fill[tg808080] (429.083,157.566) circle[radius=0.52bp];
\fill[tg808080] (430.323,157.350) circle[radius=0.52bp];
\fill[tg808080] (431.562,159.909) circle[radius=0.52bp];
\fill[tg808080] (432.802,159.442) circle[radius=0.52bp];
\fill[tg808080] (434.042,158.620) circle[radius=0.52bp];
\fill[tg808080] (435.281,158.890) circle[radius=0.52bp];
\fill[tg808080] (436.521,159.666) circle[radius=0.52bp];
\fill[tg808080] (437.760,161.742) circle[radius=0.52bp];
\draw[color=tg575757,line width=0.65bp,line cap=round] (409.250,151.367) -- (410.490,151.378) -- (411.729,151.389) -- (412.969,151.399) -- (414.208,151.410) -- (415.448,151.421) -- (416.688,151.431) -- (417.927,151.442) -- (419.167,151.453) -- (420.406,151.463) -- (421.646,151.474) -- (422.885,151.485) -- (424.125,151.495) -- (425.365,151.506) -- (426.604,151.517) -- (427.844,151.527) -- (429.083,151.538) -- (430.323,151.549) -- (431.562,151.560) -- (432.802,151.570) -- (434.042,151.581) -- (435.281,151.592) -- (436.521,151.602) -- (437.760,151.613);
\draw[color=tg0000ff,line width=0.38bp,line cap=round] (408.750,150.856) -- (409.750,151.856);
\draw[color=tg0000ff,line width=0.38bp,line cap=round] (408.750,151.856) -- (409.750,150.856);
\draw[color=tg0000ff,line width=0.38bp,line cap=round] (409.990,150.856) -- (410.990,151.856);
\draw[color=tg0000ff,line width=0.38bp,line cap=round] (409.990,151.856) -- (410.990,150.856);
\draw[color=tg0000ff,line width=0.38bp,line cap=round] (411.229,150.856) -- (412.229,151.856);
\draw[color=tg0000ff,line width=0.38bp,line cap=round] (411.229,151.856) -- (412.229,150.856);
\draw[color=tg0000ff,line width=0.38bp,line cap=round] (412.469,150.856) -- (413.469,151.856);
\draw[color=tg0000ff,line width=0.38bp,line cap=round] (412.469,151.856) -- (413.469,150.856);
\draw[color=tg0000ff,line width=0.38bp,line cap=round] (413.708,150.856) -- (414.708,151.856);
\draw[color=tg0000ff,line width=0.38bp,line cap=round] (413.708,151.856) -- (414.708,150.856);
\draw[color=tg0000ff,line width=0.38bp,line cap=round] (414.948,150.856) -- (415.948,151.856);
\draw[color=tg0000ff,line width=0.38bp,line cap=round] (414.948,151.856) -- (415.948,150.856);
\draw[color=tg0000ff,line width=0.38bp,line cap=round] (416.188,150.856) -- (417.188,151.856);
\draw[color=tg0000ff,line width=0.38bp,line cap=round] (416.188,151.856) -- (417.188,150.856);
\draw[color=tg0000ff,line width=0.38bp,line cap=round] (417.427,150.856) -- (418.427,151.856);
\draw[color=tg0000ff,line width=0.38bp,line cap=round] (417.427,151.856) -- (418.427,150.856);
\draw[color=tg0000ff,line width=0.38bp,line cap=round] (418.667,150.856) -- (419.667,151.856);
\draw[color=tg0000ff,line width=0.38bp,line cap=round] (418.667,151.856) -- (419.667,150.856);
\draw[color=tg0000ff,line width=0.38bp,line cap=round] (419.906,150.856) -- (420.906,151.856);
\draw[color=tg0000ff,line width=0.38bp,line cap=round] (419.906,151.856) -- (420.906,150.856);
\draw[color=tg0000ff,line width=0.38bp,line cap=round] (421.146,150.856) -- (422.146,151.856);
\draw[color=tg0000ff,line width=0.38bp,line cap=round] (421.146,151.856) -- (422.146,150.856);
\draw[color=tg0000ff,line width=0.38bp,line cap=round] (422.385,150.856) -- (423.385,151.856);
\draw[color=tg0000ff,line width=0.38bp,line cap=round] (422.385,151.856) -- (423.385,150.856);
\draw[color=tg0000ff,line width=0.38bp,line cap=round] (423.625,150.856) -- (424.625,151.856);
\draw[color=tg0000ff,line width=0.38bp,line cap=round] (423.625,151.856) -- (424.625,150.856);
\draw[color=tg0000ff,line width=0.38bp,line cap=round] (424.865,150.856) -- (425.865,151.856);
\draw[color=tg0000ff,line width=0.38bp,line cap=round] (424.865,151.856) -- (425.865,150.856);
\draw[color=tg0000ff,line width=0.38bp,line cap=round] (426.104,150.856) -- (427.104,151.856);
\draw[color=tg0000ff,line width=0.38bp,line cap=round] (426.104,151.856) -- (427.104,150.856);
\draw[color=tg0000ff,line width=0.38bp,line cap=round] (427.344,150.856) -- (428.344,151.856);
\draw[color=tg0000ff,line width=0.38bp,line cap=round] (427.344,151.856) -- (428.344,150.856);
\draw[color=tg0000ff,line width=0.38bp,line cap=round] (428.583,150.856) -- (429.583,151.856);
\draw[color=tg0000ff,line width=0.38bp,line cap=round] (428.583,151.856) -- (429.583,150.856);
\draw[color=tg0000ff,line width=0.38bp,line cap=round] (429.823,150.856) -- (430.823,151.856);
\draw[color=tg0000ff,line width=0.38bp,line cap=round] (429.823,151.856) -- (430.823,150.856);
\draw[color=tg0000ff,line width=0.38bp,line cap=round] (431.062,150.856) -- (432.062,151.856);
\draw[color=tg0000ff,line width=0.38bp,line cap=round] (431.062,151.856) -- (432.062,150.856);
\draw[color=tg0000ff,line width=0.38bp,line cap=round] (432.302,150.856) -- (433.302,151.856);
\draw[color=tg0000ff,line width=0.38bp,line cap=round] (432.302,151.856) -- (433.302,150.856);
\draw[color=tg0000ff,line width=0.38bp,line cap=round] (433.542,150.856) -- (434.542,151.856);
\draw[color=tg0000ff,line width=0.38bp,line cap=round] (433.542,151.856) -- (434.542,150.856);
\draw[color=tg0000ff,line width=0.38bp,line cap=round] (434.781,150.856) -- (435.781,151.856);
\draw[color=tg0000ff,line width=0.38bp,line cap=round] (434.781,151.856) -- (435.781,150.856);
\draw[color=tg0000ff,line width=0.38bp,line cap=round] (436.021,150.856) -- (437.021,151.856);
\draw[color=tg0000ff,line width=0.38bp,line cap=round] (436.021,151.856) -- (437.021,150.856);
\draw[color=tg0000ff,line width=0.38bp,line cap=round] (437.260,150.856) -- (438.260,151.856);
\draw[color=tg0000ff,line width=0.38bp,line cap=round] (437.260,151.856) -- (438.260,150.856);
\node[anchor=west,inner sep=0pt,outer sep=0pt,text=tg18314b,font={\sffamily\fontsize{8bp}{9.6bp}\selectfont\bfseries}] at (0.000,186.000) {GARCH(1,1) (log returns)};
\node[anchor=east,inner sep=0pt,outer sep=0pt,text=tg626b73,font={\sffamily\fontsize{7bp}{8.4bp}\selectfont}] at (448.000,186.000) {Volatility depends on past shocks};
\draw[color=tgc9cdd1,line width=0.45bp,line cap=round] (28.000,195.000) -- (28.000,238.000);
\draw[color=tgc9cdd1,line width=0.45bp,line cap=round] (28.000,238.000) -- (147.000,238.000);
\node[anchor=east,inner sep=0pt,outer sep=0pt,text=tg626b73,font={\sffamily\fontsize{6.7bp}{8.04bp}\selectfont}] at (25.000,231.980) {-0.012};
\node[anchor=east,inner sep=0pt,outer sep=0pt,text=tg626b73,font={\sffamily\fontsize{6.7bp}{8.04bp}\selectfont}] at (25.000,201.020) {0.011};
\draw[color=tg9da7b0,line width=0.6bp,line cap=round,dash pattern=on 2bp off 2bp] (117.250,195.000) -- (117.250,238.000);
\fill[tg808080] (28.000,221.928) circle[radius=0.52bp];
\fill[tg808080] (29.240,214.408) circle[radius=0.52bp];
\fill[tg808080] (30.479,219.185) circle[radius=0.52bp];
\fill[tg808080] (31.719,210.760) circle[radius=0.52bp];
\fill[tg808080] (32.958,209.375) circle[radius=0.52bp];
\fill[tg808080] (34.198,213.173) circle[radius=0.52bp];
\fill[tg808080] (35.438,213.219) circle[radius=0.52bp];
\fill[tg808080] (36.677,214.836) circle[radius=0.52bp];
\fill[tg808080] (37.917,218.423) circle[radius=0.52bp];
\fill[tg808080] (39.156,216.536) circle[radius=0.52bp];
\fill[tg808080] (40.396,218.574) circle[radius=0.52bp];
\fill[tg808080] (41.635,213.071) circle[radius=0.52bp];
\fill[tg808080] (42.875,215.388) circle[radius=0.52bp];
\fill[tg808080] (44.115,212.016) circle[radius=0.52bp];
\fill[tg808080] (45.354,223.417) circle[radius=0.52bp];
\fill[tg808080] (46.594,209.606) circle[radius=0.52bp];
\fill[tg808080] (47.833,220.633) circle[radius=0.52bp];
\fill[tg808080] (49.073,220.466) circle[radius=0.52bp];
\fill[tg808080] (50.312,216.818) circle[radius=0.52bp];
\fill[tg808080] (51.552,219.168) circle[radius=0.52bp];
\fill[tg808080] (52.792,213.586) circle[radius=0.52bp];
\fill[tg808080] (54.031,219.097) circle[radius=0.52bp];
\fill[tg808080] (55.271,222.192) circle[radius=0.52bp];
\fill[tg808080] (56.510,221.083) circle[radius=0.52bp];
\fill[tg808080] (57.750,206.622) circle[radius=0.52bp];
\fill[tg808080] (58.990,215.149) circle[radius=0.52bp];
\fill[tg808080] (60.229,223.622) circle[radius=0.52bp];
\fill[tg808080] (61.469,215.412) circle[radius=0.52bp];
\fill[tg808080] (62.708,226.396) circle[radius=0.52bp];
\fill[tg808080] (63.948,201.276) circle[radius=0.52bp];
\fill[tg808080] (65.188,229.630) circle[radius=0.52bp];
\fill[tg808080] (66.427,210.576) circle[radius=0.52bp];
\fill[tg808080] (67.667,210.232) circle[radius=0.52bp];
\fill[tg808080] (68.906,228.774) circle[radius=0.52bp];
\fill[tg808080] (70.146,212.479) circle[radius=0.52bp];
\fill[tg808080] (71.385,206.197) circle[radius=0.52bp];
\fill[tg808080] (72.625,211.093) circle[radius=0.52bp];
\fill[tg808080] (73.865,213.158) circle[radius=0.52bp];
\fill[tg808080] (75.104,207.566) circle[radius=0.52bp];
\fill[tg808080] (76.344,214.123) circle[radius=0.52bp];
\fill[tg808080] (77.583,212.818) circle[radius=0.52bp];
\fill[tg808080] (78.823,216.427) circle[radius=0.52bp];
\fill[tg808080] (80.062,210.961) circle[radius=0.52bp];
\fill[tg808080] (81.302,222.098) circle[radius=0.52bp];
\fill[tg808080] (82.542,209.774) circle[radius=0.52bp];
\fill[tg808080] (83.781,219.117) circle[radius=0.52bp];
\fill[tg808080] (85.021,223.124) circle[radius=0.52bp];
\fill[tg808080] (86.260,212.796) circle[radius=0.52bp];
\fill[tg808080] (87.500,203.536) circle[radius=0.52bp];
\fill[tg808080] (88.740,207.615) circle[radius=0.52bp];
\fill[tg808080] (89.979,219.737) circle[radius=0.52bp];
\fill[tg808080] (91.219,209.932) circle[radius=0.52bp];
\fill[tg808080] (92.458,219.012) circle[radius=0.52bp];
\fill[tg808080] (93.698,216.403) circle[radius=0.52bp];
\fill[tg808080] (94.938,214.837) circle[radius=0.52bp];
\fill[tg808080] (96.177,231.219) circle[radius=0.52bp];
\fill[tg808080] (97.417,213.763) circle[radius=0.52bp];
\fill[tg808080] (98.656,224.081) circle[radius=0.52bp];
\fill[tg808080] (99.896,221.418) circle[radius=0.52bp];
\fill[tg808080] (101.135,227.726) circle[radius=0.52bp];
\fill[tg808080] (102.375,229.323) circle[radius=0.52bp];
\fill[tg808080] (103.615,224.959) circle[radius=0.52bp];
\fill[tg808080] (104.854,207.218) circle[radius=0.52bp];
\fill[tg808080] (106.094,199.161) circle[radius=0.52bp];
\fill[tg808080] (107.333,215.753) circle[radius=0.52bp];
\fill[tg808080] (108.573,204.370) circle[radius=0.52bp];
\fill[tg808080] (109.812,218.203) circle[radius=0.52bp];
\fill[tg808080] (111.052,233.839) circle[radius=0.52bp];
\fill[tg808080] (112.292,213.765) circle[radius=0.52bp];
\fill[tg808080] (113.531,210.798) circle[radius=0.52bp];
\fill[tg808080] (114.771,219.695) circle[radius=0.52bp];
\fill[tg808080] (116.010,212.673) circle[radius=0.52bp];
\fill[tg808080] (117.250,210.498) circle[radius=0.52bp];
\fill[tg808080] (118.490,222.686) circle[radius=0.52bp];
\fill[tg808080] (119.729,227.764) circle[radius=0.52bp];
\fill[tg808080] (120.969,217.498) circle[radius=0.52bp];
\fill[tg808080] (122.208,217.280) circle[radius=0.52bp];
\fill[tg808080] (123.448,218.310) circle[radius=0.52bp];
\fill[tg808080] (124.688,207.918) circle[radius=0.52bp];
\fill[tg808080] (125.927,201.990) circle[radius=0.52bp];
\fill[tg808080] (127.167,219.242) circle[radius=0.52bp];
\fill[tg808080] (128.406,209.496) circle[radius=0.52bp];
\fill[tg808080] (129.646,207.628) circle[radius=0.52bp];
\fill[tg808080] (130.885,209.487) circle[radius=0.52bp];
\fill[tg808080] (132.125,206.874) circle[radius=0.52bp];
\fill[tg808080] (133.365,218.746) circle[radius=0.52bp];
\fill[tg808080] (134.604,209.567) circle[radius=0.52bp];
\fill[tg808080] (135.844,199.549) circle[radius=0.52bp];
\fill[tg808080] (137.083,207.717) circle[radius=0.52bp];
\fill[tg808080] (138.323,221.329) circle[radius=0.52bp];
\fill[tg808080] (139.562,209.967) circle[radius=0.52bp];
\fill[tg808080] (140.802,204.477) circle[radius=0.52bp];
\fill[tg808080] (142.042,212.083) circle[radius=0.52bp];
\fill[tg808080] (143.281,220.352) circle[radius=0.52bp];
\fill[tg808080] (144.521,202.642) circle[radius=0.52bp];
\fill[tg808080] (145.760,227.041) circle[radius=0.52bp];
\draw[color=tg575757,line width=0.65bp,line cap=round] (117.250,215.474) -- (118.490,215.474) -- (119.729,215.474) -- (120.969,215.474) -- (122.208,215.474) -- (123.448,215.474) -- (124.688,215.474) -- (125.927,215.474) -- (127.167,215.474) -- (128.406,215.474) -- (129.646,215.474) -- (130.885,215.474) -- (132.125,215.474) -- (133.365,215.474) -- (134.604,215.474) -- (135.844,215.474) -- (137.083,215.474) -- (138.323,215.474) -- (139.562,215.474) -- (140.802,215.474) -- (142.042,215.474) -- (143.281,215.474) -- (144.521,215.474) -- (145.760,215.474);
\draw[color=tg0000ff,line width=0.38bp,line cap=round] (116.750,212.173) -- (117.750,213.173);
\draw[color=tg0000ff,line width=0.38bp,line cap=round] (116.750,213.173) -- (117.750,212.173);
\draw[color=tg0000ff,line width=0.38bp,line cap=round] (117.990,212.173) -- (118.990,213.173);
\draw[color=tg0000ff,line width=0.38bp,line cap=round] (117.990,213.173) -- (118.990,212.173);
\draw[color=tg0000ff,line width=0.38bp,line cap=round] (119.229,212.173) -- (120.229,213.173);
\draw[color=tg0000ff,line width=0.38bp,line cap=round] (119.229,213.173) -- (120.229,212.173);
\draw[color=tg0000ff,line width=0.38bp,line cap=round] (120.469,212.173) -- (121.469,213.173);
\draw[color=tg0000ff,line width=0.38bp,line cap=round] (120.469,213.173) -- (121.469,212.173);
\draw[color=tg0000ff,line width=0.38bp,line cap=round] (121.708,212.173) -- (122.708,213.173);
\draw[color=tg0000ff,line width=0.38bp,line cap=round] (121.708,213.173) -- (122.708,212.173);
\draw[color=tg0000ff,line width=0.38bp,line cap=round] (122.948,212.173) -- (123.948,213.173);
\draw[color=tg0000ff,line width=0.38bp,line cap=round] (122.948,213.173) -- (123.948,212.173);
\draw[color=tg0000ff,line width=0.38bp,line cap=round] (124.188,212.173) -- (125.188,213.173);
\draw[color=tg0000ff,line width=0.38bp,line cap=round] (124.188,213.173) -- (125.188,212.173);
\draw[color=tg0000ff,line width=0.38bp,line cap=round] (125.427,212.173) -- (126.427,213.173);
\draw[color=tg0000ff,line width=0.38bp,line cap=round] (125.427,213.173) -- (126.427,212.173);
\draw[color=tg0000ff,line width=0.38bp,line cap=round] (126.667,212.173) -- (127.667,213.173);
\draw[color=tg0000ff,line width=0.38bp,line cap=round] (126.667,213.173) -- (127.667,212.173);
\draw[color=tg0000ff,line width=0.38bp,line cap=round] (127.906,212.173) -- (128.906,213.173);
\draw[color=tg0000ff,line width=0.38bp,line cap=round] (127.906,213.173) -- (128.906,212.173);
\draw[color=tg0000ff,line width=0.38bp,line cap=round] (129.146,212.173) -- (130.146,213.173);
\draw[color=tg0000ff,line width=0.38bp,line cap=round] (129.146,213.173) -- (130.146,212.173);
\draw[color=tg0000ff,line width=0.38bp,line cap=round] (130.385,212.173) -- (131.385,213.173);
\draw[color=tg0000ff,line width=0.38bp,line cap=round] (130.385,213.173) -- (131.385,212.173);
\draw[color=tg0000ff,line width=0.38bp,line cap=round] (131.625,212.173) -- (132.625,213.173);
\draw[color=tg0000ff,line width=0.38bp,line cap=round] (131.625,213.173) -- (132.625,212.173);
\draw[color=tg0000ff,line width=0.38bp,line cap=round] (132.865,212.173) -- (133.865,213.173);
\draw[color=tg0000ff,line width=0.38bp,line cap=round] (132.865,213.173) -- (133.865,212.173);
\draw[color=tg0000ff,line width=0.38bp,line cap=round] (134.104,212.173) -- (135.104,213.173);
\draw[color=tg0000ff,line width=0.38bp,line cap=round] (134.104,213.173) -- (135.104,212.173);
\draw[color=tg0000ff,line width=0.38bp,line cap=round] (135.344,212.173) -- (136.344,213.173);
\draw[color=tg0000ff,line width=0.38bp,line cap=round] (135.344,213.173) -- (136.344,212.173);
\draw[color=tg0000ff,line width=0.38bp,line cap=round] (136.583,212.173) -- (137.583,213.173);
\draw[color=tg0000ff,line width=0.38bp,line cap=round] (136.583,213.173) -- (137.583,212.173);
\draw[color=tg0000ff,line width=0.38bp,line cap=round] (137.823,212.173) -- (138.823,213.173);
\draw[color=tg0000ff,line width=0.38bp,line cap=round] (137.823,213.173) -- (138.823,212.173);
\draw[color=tg0000ff,line width=0.38bp,line cap=round] (139.062,212.173) -- (140.062,213.173);
\draw[color=tg0000ff,line width=0.38bp,line cap=round] (139.062,213.173) -- (140.062,212.173);
\draw[color=tg0000ff,line width=0.38bp,line cap=round] (140.302,212.173) -- (141.302,213.173);
\draw[color=tg0000ff,line width=0.38bp,line cap=round] (140.302,213.173) -- (141.302,212.173);
\draw[color=tg0000ff,line width=0.38bp,line cap=round] (141.542,212.173) -- (142.542,213.173);
\draw[color=tg0000ff,line width=0.38bp,line cap=round] (141.542,213.173) -- (142.542,212.173);
\draw[color=tg0000ff,line width=0.38bp,line cap=round] (142.781,212.173) -- (143.781,213.173);
\draw[color=tg0000ff,line width=0.38bp,line cap=round] (142.781,213.173) -- (143.781,212.173);
\draw[color=tg0000ff,line width=0.38bp,line cap=round] (144.021,212.173) -- (145.021,213.173);
\draw[color=tg0000ff,line width=0.38bp,line cap=round] (144.021,213.173) -- (145.021,212.173);
\draw[color=tg0000ff,line width=0.38bp,line cap=round] (145.260,212.173) -- (146.260,213.173);
\draw[color=tg0000ff,line width=0.38bp,line cap=round] (145.260,213.173) -- (146.260,212.173);
\draw[color=tgc9cdd1,line width=0.45bp,line cap=round] (174.000,195.000) -- (174.000,238.000);
\draw[color=tgc9cdd1,line width=0.45bp,line cap=round] (174.000,238.000) -- (293.000,238.000);
\node[anchor=east,inner sep=0pt,outer sep=0pt,text=tg626b73,font={\sffamily\fontsize{6.7bp}{8.04bp}\selectfont}] at (171.000,231.980) {-0.008};
\node[anchor=east,inner sep=0pt,outer sep=0pt,text=tg626b73,font={\sffamily\fontsize{6.7bp}{8.04bp}\selectfont}] at (171.000,201.020) {0.011};
\draw[color=tg9da7b0,line width=0.6bp,line cap=round,dash pattern=on 2bp off 2bp] (263.250,195.000) -- (263.250,238.000);
\fill[tg808080] (174.000,207.947) circle[radius=0.52bp];
\fill[tg808080] (175.240,201.217) circle[radius=0.52bp];
\fill[tg808080] (176.479,217.027) circle[radius=0.52bp];
\fill[tg808080] (177.719,219.852) circle[radius=0.52bp];
\fill[tg808080] (178.958,216.979) circle[radius=0.52bp];
\fill[tg808080] (180.198,222.818) circle[radius=0.52bp];
\fill[tg808080] (181.438,226.720) circle[radius=0.52bp];
\fill[tg808080] (182.677,213.193) circle[radius=0.52bp];
\fill[tg808080] (183.917,212.665) circle[radius=0.52bp];
\fill[tg808080] (185.156,223.374) circle[radius=0.52bp];
\fill[tg808080] (186.396,222.422) circle[radius=0.52bp];
\fill[tg808080] (187.635,224.077) circle[radius=0.52bp];
\fill[tg808080] (188.875,224.634) circle[radius=0.52bp];
\fill[tg808080] (190.115,226.742) circle[radius=0.52bp];
\fill[tg808080] (191.354,231.388) circle[radius=0.52bp];
\fill[tg808080] (192.594,210.702) circle[radius=0.52bp];
\fill[tg808080] (193.833,222.247) circle[radius=0.52bp];
\fill[tg808080] (195.073,218.408) circle[radius=0.52bp];
\fill[tg808080] (196.312,204.057) circle[radius=0.52bp];
\fill[tg808080] (197.552,223.786) circle[radius=0.52bp];
\fill[tg808080] (198.792,213.591) circle[radius=0.52bp];
\fill[tg808080] (200.031,231.287) circle[radius=0.52bp];
\fill[tg808080] (201.271,225.157) circle[radius=0.52bp];
\fill[tg808080] (202.510,212.448) circle[radius=0.52bp];
\fill[tg808080] (203.750,221.460) circle[radius=0.52bp];
\fill[tg808080] (204.990,231.619) circle[radius=0.52bp];
\fill[tg808080] (206.229,233.839) circle[radius=0.52bp];
\fill[tg808080] (207.469,215.480) circle[radius=0.52bp];
\fill[tg808080] (208.708,228.435) circle[radius=0.52bp];
\fill[tg808080] (209.948,212.255) circle[radius=0.52bp];
\fill[tg808080] (211.188,216.425) circle[radius=0.52bp];
\fill[tg808080] (212.427,222.045) circle[radius=0.52bp];
\fill[tg808080] (213.667,219.683) circle[radius=0.52bp];
\fill[tg808080] (214.906,212.874) circle[radius=0.52bp];
\fill[tg808080] (216.146,218.833) circle[radius=0.52bp];
\fill[tg808080] (217.385,226.565) circle[radius=0.52bp];
\fill[tg808080] (218.625,216.412) circle[radius=0.52bp];
\fill[tg808080] (219.865,221.858) circle[radius=0.52bp];
\fill[tg808080] (221.104,220.851) circle[radius=0.52bp];
\fill[tg808080] (222.344,214.923) circle[radius=0.52bp];
\fill[tg808080] (223.583,216.717) circle[radius=0.52bp];
\fill[tg808080] (224.823,211.872) circle[radius=0.52bp];
\fill[tg808080] (226.062,220.481) circle[radius=0.52bp];
\fill[tg808080] (227.302,222.381) circle[radius=0.52bp];
\fill[tg808080] (228.542,221.510) circle[radius=0.52bp];
\fill[tg808080] (229.781,216.459) circle[radius=0.52bp];
\fill[tg808080] (231.021,224.625) circle[radius=0.52bp];
\fill[tg808080] (232.260,221.631) circle[radius=0.52bp];
\fill[tg808080] (233.500,224.945) circle[radius=0.52bp];
\fill[tg808080] (234.740,224.629) circle[radius=0.52bp];
\fill[tg808080] (235.979,220.378) circle[radius=0.52bp];
\fill[tg808080] (237.219,202.058) circle[radius=0.52bp];
\fill[tg808080] (238.458,222.694) circle[radius=0.52bp];
\fill[tg808080] (239.698,216.770) circle[radius=0.52bp];
\fill[tg808080] (240.938,218.271) circle[radius=0.52bp];
\fill[tg808080] (242.177,220.480) circle[radius=0.52bp];
\fill[tg808080] (243.417,218.827) circle[radius=0.52bp];
\fill[tg808080] (244.656,204.534) circle[radius=0.52bp];
\fill[tg808080] (245.896,221.420) circle[radius=0.52bp];
\fill[tg808080] (247.135,205.334) circle[radius=0.52bp];
\fill[tg808080] (248.375,212.812) circle[radius=0.52bp];
\fill[tg808080] (249.615,227.168) circle[radius=0.52bp];
\fill[tg808080] (250.854,225.087) circle[radius=0.52bp];
\fill[tg808080] (252.094,223.217) circle[radius=0.52bp];
\fill[tg808080] (253.333,228.465) circle[radius=0.52bp];
\fill[tg808080] (254.573,208.780) circle[radius=0.52bp];
\fill[tg808080] (255.812,215.217) circle[radius=0.52bp];
\fill[tg808080] (257.052,214.366) circle[radius=0.52bp];
\fill[tg808080] (258.292,227.825) circle[radius=0.52bp];
\fill[tg808080] (259.531,214.963) circle[radius=0.52bp];
\fill[tg808080] (260.771,223.566) circle[radius=0.52bp];
\fill[tg808080] (262.010,199.161) circle[radius=0.52bp];
\fill[tg808080] (263.250,220.030) circle[radius=0.52bp];
\fill[tg808080] (264.490,216.196) circle[radius=0.52bp];
\fill[tg808080] (265.729,215.123) circle[radius=0.52bp];
\fill[tg808080] (266.969,223.307) circle[radius=0.52bp];
\fill[tg808080] (268.208,222.414) circle[radius=0.52bp];
\fill[tg808080] (269.448,222.297) circle[radius=0.52bp];
\fill[tg808080] (270.688,212.253) circle[radius=0.52bp];
\fill[tg808080] (271.927,224.895) circle[radius=0.52bp];
\fill[tg808080] (273.167,220.318) circle[radius=0.52bp];
\fill[tg808080] (274.406,216.525) circle[radius=0.52bp];
\fill[tg808080] (275.646,222.072) circle[radius=0.52bp];
\fill[tg808080] (276.885,228.933) circle[radius=0.52bp];
\fill[tg808080] (278.125,228.519) circle[radius=0.52bp];
\fill[tg808080] (279.365,222.542) circle[radius=0.52bp];
\fill[tg808080] (280.604,226.872) circle[radius=0.52bp];
\fill[tg808080] (281.844,211.099) circle[radius=0.52bp];
\fill[tg808080] (283.083,218.936) circle[radius=0.52bp];
\fill[tg808080] (284.323,219.278) circle[radius=0.52bp];
\fill[tg808080] (285.562,226.772) circle[radius=0.52bp];
\fill[tg808080] (286.802,214.937) circle[radius=0.52bp];
\fill[tg808080] (288.042,221.835) circle[radius=0.52bp];
\fill[tg808080] (289.281,226.195) circle[radius=0.52bp];
\fill[tg808080] (290.521,224.003) circle[radius=0.52bp];
\fill[tg808080] (291.760,200.953) circle[radius=0.52bp];
\draw[color=tg575757,line width=0.65bp,line cap=round] (263.250,219.033) -- (264.490,219.033) -- (265.729,219.033) -- (266.969,219.033) -- (268.208,219.033) -- (269.448,219.033) -- (270.688,219.033) -- (271.927,219.033) -- (273.167,219.033) -- (274.406,219.033) -- (275.646,219.033) -- (276.885,219.033) -- (278.125,219.033) -- (279.365,219.033) -- (280.604,219.033) -- (281.844,219.033) -- (283.083,219.033) -- (284.323,219.033) -- (285.562,219.033) -- (286.802,219.033) -- (288.042,219.033) -- (289.281,219.033) -- (290.521,219.033) -- (291.760,219.033);
\draw[color=tg0000ff,line width=0.38bp,line cap=round] (262.750,198.661) -- (263.750,199.661);
\draw[color=tg0000ff,line width=0.38bp,line cap=round] (262.750,199.661) -- (263.750,198.661);
\draw[color=tg0000ff,line width=0.38bp,line cap=round] (263.990,198.661) -- (264.990,199.661);
\draw[color=tg0000ff,line width=0.38bp,line cap=round] (263.990,199.661) -- (264.990,198.661);
\draw[color=tg0000ff,line width=0.38bp,line cap=round] (265.229,198.661) -- (266.229,199.661);
\draw[color=tg0000ff,line width=0.38bp,line cap=round] (265.229,199.661) -- (266.229,198.661);
\draw[color=tg0000ff,line width=0.38bp,line cap=round] (266.469,198.661) -- (267.469,199.661);
\draw[color=tg0000ff,line width=0.38bp,line cap=round] (266.469,199.661) -- (267.469,198.661);
\draw[color=tg0000ff,line width=0.38bp,line cap=round] (267.708,198.661) -- (268.708,199.661);
\draw[color=tg0000ff,line width=0.38bp,line cap=round] (267.708,199.661) -- (268.708,198.661);
\draw[color=tg0000ff,line width=0.38bp,line cap=round] (268.948,198.661) -- (269.948,199.661);
\draw[color=tg0000ff,line width=0.38bp,line cap=round] (268.948,199.661) -- (269.948,198.661);
\draw[color=tg0000ff,line width=0.38bp,line cap=round] (270.188,198.661) -- (271.188,199.661);
\draw[color=tg0000ff,line width=0.38bp,line cap=round] (270.188,199.661) -- (271.188,198.661);
\draw[color=tg0000ff,line width=0.38bp,line cap=round] (271.427,198.661) -- (272.427,199.661);
\draw[color=tg0000ff,line width=0.38bp,line cap=round] (271.427,199.661) -- (272.427,198.661);
\draw[color=tg0000ff,line width=0.38bp,line cap=round] (272.667,198.661) -- (273.667,199.661);
\draw[color=tg0000ff,line width=0.38bp,line cap=round] (272.667,199.661) -- (273.667,198.661);
\draw[color=tg0000ff,line width=0.38bp,line cap=round] (273.906,198.661) -- (274.906,199.661);
\draw[color=tg0000ff,line width=0.38bp,line cap=round] (273.906,199.661) -- (274.906,198.661);
\draw[color=tg0000ff,line width=0.38bp,line cap=round] (275.146,198.661) -- (276.146,199.661);
\draw[color=tg0000ff,line width=0.38bp,line cap=round] (275.146,199.661) -- (276.146,198.661);
\draw[color=tg0000ff,line width=0.38bp,line cap=round] (276.385,198.661) -- (277.385,199.661);
\draw[color=tg0000ff,line width=0.38bp,line cap=round] (276.385,199.661) -- (277.385,198.661);
\draw[color=tg0000ff,line width=0.38bp,line cap=round] (277.625,198.661) -- (278.625,199.661);
\draw[color=tg0000ff,line width=0.38bp,line cap=round] (277.625,199.661) -- (278.625,198.661);
\draw[color=tg0000ff,line width=0.38bp,line cap=round] (278.865,198.661) -- (279.865,199.661);
\draw[color=tg0000ff,line width=0.38bp,line cap=round] (278.865,199.661) -- (279.865,198.661);
\draw[color=tg0000ff,line width=0.38bp,line cap=round] (280.104,198.661) -- (281.104,199.661);
\draw[color=tg0000ff,line width=0.38bp,line cap=round] (280.104,199.661) -- (281.104,198.661);
\draw[color=tg0000ff,line width=0.38bp,line cap=round] (281.344,198.661) -- (282.344,199.661);
\draw[color=tg0000ff,line width=0.38bp,line cap=round] (281.344,199.661) -- (282.344,198.661);
\draw[color=tg0000ff,line width=0.38bp,line cap=round] (282.583,198.661) -- (283.583,199.661);
\draw[color=tg0000ff,line width=0.38bp,line cap=round] (282.583,199.661) -- (283.583,198.661);
\draw[color=tg0000ff,line width=0.38bp,line cap=round] (283.823,198.661) -- (284.823,199.661);
\draw[color=tg0000ff,line width=0.38bp,line cap=round] (283.823,199.661) -- (284.823,198.661);
\draw[color=tg0000ff,line width=0.38bp,line cap=round] (285.062,198.661) -- (286.062,199.661);
\draw[color=tg0000ff,line width=0.38bp,line cap=round] (285.062,199.661) -- (286.062,198.661);
\draw[color=tg0000ff,line width=0.38bp,line cap=round] (286.302,198.661) -- (287.302,199.661);
\draw[color=tg0000ff,line width=0.38bp,line cap=round] (286.302,199.661) -- (287.302,198.661);
\draw[color=tg0000ff,line width=0.38bp,line cap=round] (287.542,198.661) -- (288.542,199.661);
\draw[color=tg0000ff,line width=0.38bp,line cap=round] (287.542,199.661) -- (288.542,198.661);
\draw[color=tg0000ff,line width=0.38bp,line cap=round] (288.781,198.661) -- (289.781,199.661);
\draw[color=tg0000ff,line width=0.38bp,line cap=round] (288.781,199.661) -- (289.781,198.661);
\draw[color=tg0000ff,line width=0.38bp,line cap=round] (290.021,198.661) -- (291.021,199.661);
\draw[color=tg0000ff,line width=0.38bp,line cap=round] (290.021,199.661) -- (291.021,198.661);
\draw[color=tg0000ff,line width=0.38bp,line cap=round] (291.260,198.661) -- (292.260,199.661);
\draw[color=tg0000ff,line width=0.38bp,line cap=round] (291.260,199.661) -- (292.260,198.661);
\draw[color=tgc9cdd1,line width=0.45bp,line cap=round] (320.000,195.000) -- (320.000,238.000);
\draw[color=tgc9cdd1,line width=0.45bp,line cap=round] (320.000,238.000) -- (439.000,238.000);
\node[anchor=east,inner sep=0pt,outer sep=0pt,text=tg626b73,font={\sffamily\fontsize{6.7bp}{8.04bp}\selectfont}] at (317.000,231.980) {-0.018};
\node[anchor=east,inner sep=0pt,outer sep=0pt,text=tg626b73,font={\sffamily\fontsize{6.7bp}{8.04bp}\selectfont}] at (317.000,201.020) {0.017};
\draw[color=tg9da7b0,line width=0.6bp,line cap=round,dash pattern=on 2bp off 2bp] (409.250,195.000) -- (409.250,238.000);
\fill[tg808080] (320.000,204.947) circle[radius=0.52bp];
\fill[tg808080] (321.240,210.399) circle[radius=0.52bp];
\fill[tg808080] (322.479,223.580) circle[radius=0.52bp];
\fill[tg808080] (323.719,213.730) circle[radius=0.52bp];
\fill[tg808080] (324.958,208.392) circle[radius=0.52bp];
\fill[tg808080] (326.198,217.477) circle[radius=0.52bp];
\fill[tg808080] (327.438,203.755) circle[radius=0.52bp];
\fill[tg808080] (328.677,221.731) circle[radius=0.52bp];
\fill[tg808080] (329.917,200.045) circle[radius=0.52bp];
\fill[tg808080] (331.156,201.783) circle[radius=0.52bp];
\fill[tg808080] (332.396,210.838) circle[radius=0.52bp];
\fill[tg808080] (333.635,213.966) circle[radius=0.52bp];
\fill[tg808080] (334.875,215.792) circle[radius=0.52bp];
\fill[tg808080] (336.115,223.370) circle[radius=0.52bp];
\fill[tg808080] (337.354,208.511) circle[radius=0.52bp];
\fill[tg808080] (338.594,206.348) circle[radius=0.52bp];
\fill[tg808080] (339.833,209.838) circle[radius=0.52bp];
\fill[tg808080] (341.073,213.905) circle[radius=0.52bp];
\fill[tg808080] (342.312,215.302) circle[radius=0.52bp];
\fill[tg808080] (343.552,222.096) circle[radius=0.52bp];
\fill[tg808080] (344.792,220.275) circle[radius=0.52bp];
\fill[tg808080] (346.031,224.993) circle[radius=0.52bp];
\fill[tg808080] (347.271,212.803) circle[radius=0.52bp];
\fill[tg808080] (348.510,221.748) circle[radius=0.52bp];
\fill[tg808080] (349.750,210.798) circle[radius=0.52bp];
\fill[tg808080] (350.990,220.115) circle[radius=0.52bp];
\fill[tg808080] (352.229,207.630) circle[radius=0.52bp];
\fill[tg808080] (353.469,206.014) circle[radius=0.52bp];
\fill[tg808080] (354.708,205.206) circle[radius=0.52bp];
\fill[tg808080] (355.948,207.204) circle[radius=0.52bp];
\fill[tg808080] (357.188,218.449) circle[radius=0.52bp];
\fill[tg808080] (358.427,218.329) circle[radius=0.52bp];
\fill[tg808080] (359.667,221.007) circle[radius=0.52bp];
\fill[tg808080] (360.906,213.907) circle[radius=0.52bp];
\fill[tg808080] (362.146,205.872) circle[radius=0.52bp];
\fill[tg808080] (363.385,223.538) circle[radius=0.52bp];
\fill[tg808080] (364.625,199.161) circle[radius=0.52bp];
\fill[tg808080] (365.865,221.264) circle[radius=0.52bp];
\fill[tg808080] (367.104,226.077) circle[radius=0.52bp];
\fill[tg808080] (368.344,213.098) circle[radius=0.52bp];
\fill[tg808080] (369.583,207.404) circle[radius=0.52bp];
\fill[tg808080] (370.823,218.951) circle[radius=0.52bp];
\fill[tg808080] (372.062,223.679) circle[radius=0.52bp];
\fill[tg808080] (373.302,202.967) circle[radius=0.52bp];
\fill[tg808080] (374.542,225.428) circle[radius=0.52bp];
\fill[tg808080] (375.781,207.526) circle[radius=0.52bp];
\fill[tg808080] (377.021,214.775) circle[radius=0.52bp];
\fill[tg808080] (378.260,218.952) circle[radius=0.52bp];
\fill[tg808080] (379.500,215.711) circle[radius=0.52bp];
\fill[tg808080] (380.740,223.243) circle[radius=0.52bp];
\fill[tg808080] (381.979,223.523) circle[radius=0.52bp];
\fill[tg808080] (383.219,213.950) circle[radius=0.52bp];
\fill[tg808080] (384.458,222.653) circle[radius=0.52bp];
\fill[tg808080] (385.698,220.998) circle[radius=0.52bp];
\fill[tg808080] (386.938,206.983) circle[radius=0.52bp];
\fill[tg808080] (388.177,203.871) circle[radius=0.52bp];
\fill[tg808080] (389.417,207.393) circle[radius=0.52bp];
\fill[tg808080] (390.656,208.981) circle[radius=0.52bp];
\fill[tg808080] (391.896,217.995) circle[radius=0.52bp];
\fill[tg808080] (393.135,215.399) circle[radius=0.52bp];
\fill[tg808080] (394.375,220.877) circle[radius=0.52bp];
\fill[tg808080] (395.615,233.839) circle[radius=0.52bp];
\fill[tg808080] (396.854,214.280) circle[radius=0.52bp];
\fill[tg808080] (398.094,214.219) circle[radius=0.52bp];
\fill[tg808080] (399.333,208.563) circle[radius=0.52bp];
\fill[tg808080] (400.573,225.078) circle[radius=0.52bp];
\fill[tg808080] (401.812,207.165) circle[radius=0.52bp];
\fill[tg808080] (403.052,220.314) circle[radius=0.52bp];
\fill[tg808080] (404.292,225.234) circle[radius=0.52bp];
\fill[tg808080] (405.531,218.063) circle[radius=0.52bp];
\fill[tg808080] (406.771,209.511) circle[radius=0.52bp];
\fill[tg808080] (408.010,216.167) circle[radius=0.52bp];
\fill[tg808080] (409.250,213.250) circle[radius=0.52bp];
\fill[tg808080] (410.490,221.409) circle[radius=0.52bp];
\fill[tg808080] (411.729,213.229) circle[radius=0.52bp];
\fill[tg808080] (412.969,210.504) circle[radius=0.52bp];
\fill[tg808080] (414.208,221.140) circle[radius=0.52bp];
\fill[tg808080] (415.448,213.699) circle[radius=0.52bp];
\fill[tg808080] (416.688,223.326) circle[radius=0.52bp];
\fill[tg808080] (417.927,214.769) circle[radius=0.52bp];
\fill[tg808080] (419.167,224.196) circle[radius=0.52bp];
\fill[tg808080] (420.406,213.670) circle[radius=0.52bp];
\fill[tg808080] (421.646,216.127) circle[radius=0.52bp];
\fill[tg808080] (422.885,222.000) circle[radius=0.52bp];
\fill[tg808080] (424.125,214.115) circle[radius=0.52bp];
\fill[tg808080] (425.365,221.184) circle[radius=0.52bp];
\fill[tg808080] (426.604,215.634) circle[radius=0.52bp];
\fill[tg808080] (427.844,220.600) circle[radius=0.52bp];
\fill[tg808080] (429.083,213.386) circle[radius=0.52bp];
\fill[tg808080] (430.323,218.779) circle[radius=0.52bp];
\fill[tg808080] (431.562,209.313) circle[radius=0.52bp];
\fill[tg808080] (432.802,213.710) circle[radius=0.52bp];
\fill[tg808080] (434.042,213.395) circle[radius=0.52bp];
\fill[tg808080] (435.281,222.434) circle[radius=0.52bp];
\fill[tg808080] (436.521,226.863) circle[radius=0.52bp];
\fill[tg808080] (437.760,221.185) circle[radius=0.52bp];
\draw[color=tg575757,line width=0.65bp,line cap=round] (409.250,216.013) -- (410.490,216.013) -- (411.729,216.013) -- (412.969,216.013) -- (414.208,216.013) -- (415.448,216.013) -- (416.688,216.013) -- (417.927,216.013) -- (419.167,216.013) -- (420.406,216.013) -- (421.646,216.013) -- (422.885,216.013) -- (424.125,216.013) -- (425.365,216.013) -- (426.604,216.013) -- (427.844,216.013) -- (429.083,216.013) -- (430.323,216.013) -- (431.562,216.013) -- (432.802,216.013) -- (434.042,216.013) -- (435.281,216.013) -- (436.521,216.013) -- (437.760,216.013);
\draw[color=tg0000ff,line width=0.38bp,line cap=round] (408.750,215.667) -- (409.750,216.667);
\draw[color=tg0000ff,line width=0.38bp,line cap=round] (408.750,216.667) -- (409.750,215.667);
\draw[color=tg0000ff,line width=0.38bp,line cap=round] (409.990,215.667) -- (410.990,216.667);
\draw[color=tg0000ff,line width=0.38bp,line cap=round] (409.990,216.667) -- (410.990,215.667);
\draw[color=tg0000ff,line width=0.38bp,line cap=round] (411.229,215.667) -- (412.229,216.667);
\draw[color=tg0000ff,line width=0.38bp,line cap=round] (411.229,216.667) -- (412.229,215.667);
\draw[color=tg0000ff,line width=0.38bp,line cap=round] (412.469,215.667) -- (413.469,216.667);
\draw[color=tg0000ff,line width=0.38bp,line cap=round] (412.469,216.667) -- (413.469,215.667);
\draw[color=tg0000ff,line width=0.38bp,line cap=round] (413.708,215.667) -- (414.708,216.667);
\draw[color=tg0000ff,line width=0.38bp,line cap=round] (413.708,216.667) -- (414.708,215.667);
\draw[color=tg0000ff,line width=0.38bp,line cap=round] (414.948,215.667) -- (415.948,216.667);
\draw[color=tg0000ff,line width=0.38bp,line cap=round] (414.948,216.667) -- (415.948,215.667);
\draw[color=tg0000ff,line width=0.38bp,line cap=round] (416.188,215.667) -- (417.188,216.667);
\draw[color=tg0000ff,line width=0.38bp,line cap=round] (416.188,216.667) -- (417.188,215.667);
\draw[color=tg0000ff,line width=0.38bp,line cap=round] (417.427,215.667) -- (418.427,216.667);
\draw[color=tg0000ff,line width=0.38bp,line cap=round] (417.427,216.667) -- (418.427,215.667);
\draw[color=tg0000ff,line width=0.38bp,line cap=round] (418.667,215.667) -- (419.667,216.667);
\draw[color=tg0000ff,line width=0.38bp,line cap=round] (418.667,216.667) -- (419.667,215.667);
\draw[color=tg0000ff,line width=0.38bp,line cap=round] (419.906,215.667) -- (420.906,216.667);
\draw[color=tg0000ff,line width=0.38bp,line cap=round] (419.906,216.667) -- (420.906,215.667);
\draw[color=tg0000ff,line width=0.38bp,line cap=round] (421.146,215.667) -- (422.146,216.667);
\draw[color=tg0000ff,line width=0.38bp,line cap=round] (421.146,216.667) -- (422.146,215.667);
\draw[color=tg0000ff,line width=0.38bp,line cap=round] (422.385,215.667) -- (423.385,216.667);
\draw[color=tg0000ff,line width=0.38bp,line cap=round] (422.385,216.667) -- (423.385,215.667);
\draw[color=tg0000ff,line width=0.38bp,line cap=round] (423.625,215.667) -- (424.625,216.667);
\draw[color=tg0000ff,line width=0.38bp,line cap=round] (423.625,216.667) -- (424.625,215.667);
\draw[color=tg0000ff,line width=0.38bp,line cap=round] (424.865,215.667) -- (425.865,216.667);
\draw[color=tg0000ff,line width=0.38bp,line cap=round] (424.865,216.667) -- (425.865,215.667);
\draw[color=tg0000ff,line width=0.38bp,line cap=round] (426.104,215.667) -- (427.104,216.667);
\draw[color=tg0000ff,line width=0.38bp,line cap=round] (426.104,216.667) -- (427.104,215.667);
\draw[color=tg0000ff,line width=0.38bp,line cap=round] (427.344,215.667) -- (428.344,216.667);
\draw[color=tg0000ff,line width=0.38bp,line cap=round] (427.344,216.667) -- (428.344,215.667);
\draw[color=tg0000ff,line width=0.38bp,line cap=round] (428.583,215.667) -- (429.583,216.667);
\draw[color=tg0000ff,line width=0.38bp,line cap=round] (428.583,216.667) -- (429.583,215.667);
\draw[color=tg0000ff,line width=0.38bp,line cap=round] (429.823,215.667) -- (430.823,216.667);
\draw[color=tg0000ff,line width=0.38bp,line cap=round] (429.823,216.667) -- (430.823,215.667);
\draw[color=tg0000ff,line width=0.38bp,line cap=round] (431.062,215.667) -- (432.062,216.667);
\draw[color=tg0000ff,line width=0.38bp,line cap=round] (431.062,216.667) -- (432.062,215.667);
\draw[color=tg0000ff,line width=0.38bp,line cap=round] (432.302,215.667) -- (433.302,216.667);
\draw[color=tg0000ff,line width=0.38bp,line cap=round] (432.302,216.667) -- (433.302,215.667);
\draw[color=tg0000ff,line width=0.38bp,line cap=round] (433.542,215.667) -- (434.542,216.667);
\draw[color=tg0000ff,line width=0.38bp,line cap=round] (433.542,216.667) -- (434.542,215.667);
\draw[color=tg0000ff,line width=0.38bp,line cap=round] (434.781,215.667) -- (435.781,216.667);
\draw[color=tg0000ff,line width=0.38bp,line cap=round] (434.781,216.667) -- (435.781,215.667);
\draw[color=tg0000ff,line width=0.38bp,line cap=round] (436.021,215.667) -- (437.021,216.667);
\draw[color=tg0000ff,line width=0.38bp,line cap=round] (436.021,216.667) -- (437.021,215.667);
\draw[color=tg0000ff,line width=0.38bp,line cap=round] (437.260,215.667) -- (438.260,216.667);
\draw[color=tg0000ff,line width=0.38bp,line cap=round] (437.260,216.667) -- (438.260,215.667);
\node[anchor=west,inner sep=0pt,outer sep=0pt,text=tg18314b,font={\sffamily\fontsize{8bp}{9.6bp}\selectfont\bfseries}] at (0.000,253.000) {Hidden Markov regimes (log returns)};
\node[anchor=east,inner sep=0pt,outer sep=0pt,text=tg626b73,font={\sffamily\fontsize{7bp}{8.4bp}\selectfont}] at (448.000,253.000) {Hidden states inferred from history};
\draw[color=tgc9cdd1,line width=0.45bp,line cap=round] (28.000,262.000) -- (28.000,305.000);
\draw[color=tgc9cdd1,line width=0.45bp,line cap=round] (28.000,305.000) -- (147.000,305.000);
\node[anchor=east,inner sep=0pt,outer sep=0pt,text=tg626b73,font={\sffamily\fontsize{6.7bp}{8.04bp}\selectfont}] at (25.000,298.980) {-0.08};
\node[anchor=east,inner sep=0pt,outer sep=0pt,text=tg626b73,font={\sffamily\fontsize{6.7bp}{8.04bp}\selectfont}] at (25.000,268.020) {0.03};
\draw[color=tg9da7b0,line width=0.6bp,line cap=round,dash pattern=on 2bp off 2bp] (117.250,262.000) -- (117.250,305.000);
\fill[tg808080] (28.000,278.507) circle[radius=0.52bp];
\fill[tg808080] (29.240,279.599) circle[radius=0.52bp];
\fill[tg808080] (30.479,275.168) circle[radius=0.52bp];
\fill[tg808080] (31.719,277.574) circle[radius=0.52bp];
\fill[tg808080] (32.958,277.919) circle[radius=0.52bp];
\fill[tg808080] (34.198,276.002) circle[radius=0.52bp];
\fill[tg808080] (35.438,277.634) circle[radius=0.52bp];
\fill[tg808080] (36.677,278.692) circle[radius=0.52bp];
\fill[tg808080] (37.917,274.818) circle[radius=0.52bp];
\fill[tg808080] (39.156,276.300) circle[radius=0.52bp];
\fill[tg808080] (40.396,277.144) circle[radius=0.52bp];
\fill[tg808080] (41.635,275.368) circle[radius=0.52bp];
\fill[tg808080] (42.875,277.718) circle[radius=0.52bp];
\fill[tg808080] (44.115,277.096) circle[radius=0.52bp];
\fill[tg808080] (45.354,276.620) circle[radius=0.52bp];
\fill[tg808080] (46.594,272.858) circle[radius=0.52bp];
\fill[tg808080] (47.833,277.247) circle[radius=0.52bp];
\fill[tg808080] (49.073,274.555) circle[radius=0.52bp];
\fill[tg808080] (50.312,276.644) circle[radius=0.52bp];
\fill[tg808080] (51.552,277.969) circle[radius=0.52bp];
\fill[tg808080] (52.792,278.489) circle[radius=0.52bp];
\fill[tg808080] (54.031,277.320) circle[radius=0.52bp];
\fill[tg808080] (55.271,277.658) circle[radius=0.52bp];
\fill[tg808080] (56.510,279.339) circle[radius=0.52bp];
\fill[tg808080] (57.750,276.577) circle[radius=0.52bp];
\fill[tg808080] (58.990,274.121) circle[radius=0.52bp];
\fill[tg808080] (60.229,275.881) circle[radius=0.52bp];
\fill[tg808080] (61.469,278.674) circle[radius=0.52bp];
\fill[tg808080] (62.708,275.139) circle[radius=0.52bp];
\fill[tg808080] (63.948,278.132) circle[radius=0.52bp];
\fill[tg808080] (65.188,279.435) circle[radius=0.52bp];
\fill[tg808080] (66.427,276.916) circle[radius=0.52bp];
\fill[tg808080] (67.667,278.051) circle[radius=0.52bp];
\fill[tg808080] (68.906,278.234) circle[radius=0.52bp];
\fill[tg808080] (70.146,276.167) circle[radius=0.52bp];
\fill[tg808080] (71.385,275.844) circle[radius=0.52bp];
\fill[tg808080] (72.625,276.686) circle[radius=0.52bp];
\fill[tg808080] (73.865,278.365) circle[radius=0.52bp];
\fill[tg808080] (75.104,277.319) circle[radius=0.52bp];
\fill[tg808080] (76.344,276.083) circle[radius=0.52bp];
\fill[tg808080] (77.583,274.968) circle[radius=0.52bp];
\fill[tg808080] (78.823,276.004) circle[radius=0.52bp];
\fill[tg808080] (80.062,275.013) circle[radius=0.52bp];
\fill[tg808080] (81.302,276.062) circle[radius=0.52bp];
\fill[tg808080] (82.542,275.436) circle[radius=0.52bp];
\fill[tg808080] (83.781,277.880) circle[radius=0.52bp];
\fill[tg808080] (85.021,276.673) circle[radius=0.52bp];
\fill[tg808080] (86.260,279.276) circle[radius=0.52bp];
\fill[tg808080] (87.500,276.593) circle[radius=0.52bp];
\fill[tg808080] (88.740,276.750) circle[radius=0.52bp];
\fill[tg808080] (89.979,278.269) circle[radius=0.52bp];
\fill[tg808080] (91.219,277.453) circle[radius=0.52bp];
\fill[tg808080] (92.458,278.971) circle[radius=0.52bp];
\fill[tg808080] (93.698,276.643) circle[radius=0.52bp];
\fill[tg808080] (94.938,275.485) circle[radius=0.52bp];
\fill[tg808080] (96.177,275.810) circle[radius=0.52bp];
\fill[tg808080] (97.417,274.623) circle[radius=0.52bp];
\fill[tg808080] (98.656,276.590) circle[radius=0.52bp];
\fill[tg808080] (99.896,277.943) circle[radius=0.52bp];
\fill[tg808080] (101.135,276.773) circle[radius=0.52bp];
\fill[tg808080] (102.375,275.789) circle[radius=0.52bp];
\fill[tg808080] (103.615,280.251) circle[radius=0.52bp];
\fill[tg808080] (104.854,277.498) circle[radius=0.52bp];
\fill[tg808080] (106.094,276.683) circle[radius=0.52bp];
\fill[tg808080] (107.333,276.596) circle[radius=0.52bp];
\fill[tg808080] (108.573,276.429) circle[radius=0.52bp];
\fill[tg808080] (109.812,275.064) circle[radius=0.52bp];
\fill[tg808080] (111.052,276.492) circle[radius=0.52bp];
\fill[tg808080] (112.292,275.676) circle[radius=0.52bp];
\fill[tg808080] (113.531,277.806) circle[radius=0.52bp];
\fill[tg808080] (114.771,292.315) circle[radius=0.52bp];
\fill[tg808080] (116.010,279.700) circle[radius=0.52bp];
\fill[tg808080] (117.250,272.977) circle[radius=0.52bp];
\fill[tg808080] (118.490,275.894) circle[radius=0.52bp];
\fill[tg808080] (119.729,284.986) circle[radius=0.52bp];
\fill[tg808080] (120.969,271.226) circle[radius=0.52bp];
\fill[tg808080] (122.208,284.281) circle[radius=0.52bp];
\fill[tg808080] (123.448,300.839) circle[radius=0.52bp];
\fill[tg808080] (124.688,282.916) circle[radius=0.52bp];
\fill[tg808080] (125.927,281.561) circle[radius=0.52bp];
\fill[tg808080] (127.167,271.420) circle[radius=0.52bp];
\fill[tg808080] (128.406,291.671) circle[radius=0.52bp];
\fill[tg808080] (129.646,275.351) circle[radius=0.52bp];
\fill[tg808080] (130.885,279.122) circle[radius=0.52bp];
\fill[tg808080] (132.125,271.384) circle[radius=0.52bp];
\fill[tg808080] (133.365,278.199) circle[radius=0.52bp];
\fill[tg808080] (134.604,277.353) circle[radius=0.52bp];
\fill[tg808080] (135.844,279.092) circle[radius=0.52bp];
\fill[tg808080] (137.083,266.161) circle[radius=0.52bp];
\fill[tg808080] (138.323,275.684) circle[radius=0.52bp];
\fill[tg808080] (139.562,269.241) circle[radius=0.52bp];
\fill[tg808080] (140.802,274.112) circle[radius=0.52bp];
\fill[tg808080] (142.042,274.354) circle[radius=0.52bp];
\fill[tg808080] (143.281,276.329) circle[radius=0.52bp];
\fill[tg808080] (144.521,283.884) circle[radius=0.52bp];
\fill[tg808080] (145.760,275.126) circle[radius=0.52bp];
\draw[color=tg575757,line width=0.65bp,line cap=round] (117.250,277.069) -- (118.490,277.048) -- (119.729,277.030) -- (120.969,277.013) -- (122.208,276.999) -- (123.448,276.986) -- (124.688,276.974) -- (125.927,276.964) -- (127.167,276.955) -- (128.406,276.946) -- (129.646,276.939) -- (130.885,276.933) -- (132.125,276.927) -- (133.365,276.922) -- (134.604,276.917) -- (135.844,276.913) -- (137.083,276.909) -- (138.323,276.906) -- (139.562,276.903) -- (140.802,276.901) -- (142.042,276.899) -- (143.281,276.897) -- (144.521,276.895) -- (145.760,276.893);
\draw[color=tg0000ff,line width=0.38bp,line cap=round] (116.750,279.200) -- (117.750,280.200);
\draw[color=tg0000ff,line width=0.38bp,line cap=round] (116.750,280.200) -- (117.750,279.200);
\draw[color=tg0000ff,line width=0.38bp,line cap=round] (117.990,279.200) -- (118.990,280.200);
\draw[color=tg0000ff,line width=0.38bp,line cap=round] (117.990,280.200) -- (118.990,279.200);
\draw[color=tg0000ff,line width=0.38bp,line cap=round] (119.229,279.200) -- (120.229,280.200);
\draw[color=tg0000ff,line width=0.38bp,line cap=round] (119.229,280.200) -- (120.229,279.200);
\draw[color=tg0000ff,line width=0.38bp,line cap=round] (120.469,279.200) -- (121.469,280.200);
\draw[color=tg0000ff,line width=0.38bp,line cap=round] (120.469,280.200) -- (121.469,279.200);
\draw[color=tg0000ff,line width=0.38bp,line cap=round] (121.708,279.200) -- (122.708,280.200);
\draw[color=tg0000ff,line width=0.38bp,line cap=round] (121.708,280.200) -- (122.708,279.200);
\draw[color=tg0000ff,line width=0.38bp,line cap=round] (122.948,279.200) -- (123.948,280.200);
\draw[color=tg0000ff,line width=0.38bp,line cap=round] (122.948,280.200) -- (123.948,279.200);
\draw[color=tg0000ff,line width=0.38bp,line cap=round] (124.188,279.200) -- (125.188,280.200);
\draw[color=tg0000ff,line width=0.38bp,line cap=round] (124.188,280.200) -- (125.188,279.200);
\draw[color=tg0000ff,line width=0.38bp,line cap=round] (125.427,279.200) -- (126.427,280.200);
\draw[color=tg0000ff,line width=0.38bp,line cap=round] (125.427,280.200) -- (126.427,279.200);
\draw[color=tg0000ff,line width=0.38bp,line cap=round] (126.667,279.200) -- (127.667,280.200);
\draw[color=tg0000ff,line width=0.38bp,line cap=round] (126.667,280.200) -- (127.667,279.200);
\draw[color=tg0000ff,line width=0.38bp,line cap=round] (127.906,279.200) -- (128.906,280.200);
\draw[color=tg0000ff,line width=0.38bp,line cap=round] (127.906,280.200) -- (128.906,279.200);
\draw[color=tg0000ff,line width=0.38bp,line cap=round] (129.146,279.200) -- (130.146,280.200);
\draw[color=tg0000ff,line width=0.38bp,line cap=round] (129.146,280.200) -- (130.146,279.200);
\draw[color=tg0000ff,line width=0.38bp,line cap=round] (130.385,279.200) -- (131.385,280.200);
\draw[color=tg0000ff,line width=0.38bp,line cap=round] (130.385,280.200) -- (131.385,279.200);
\draw[color=tg0000ff,line width=0.38bp,line cap=round] (131.625,279.200) -- (132.625,280.200);
\draw[color=tg0000ff,line width=0.38bp,line cap=round] (131.625,280.200) -- (132.625,279.200);
\draw[color=tg0000ff,line width=0.38bp,line cap=round] (132.865,279.200) -- (133.865,280.200);
\draw[color=tg0000ff,line width=0.38bp,line cap=round] (132.865,280.200) -- (133.865,279.200);
\draw[color=tg0000ff,line width=0.38bp,line cap=round] (134.104,279.200) -- (135.104,280.200);
\draw[color=tg0000ff,line width=0.38bp,line cap=round] (134.104,280.200) -- (135.104,279.200);
\draw[color=tg0000ff,line width=0.38bp,line cap=round] (135.344,279.200) -- (136.344,280.200);
\draw[color=tg0000ff,line width=0.38bp,line cap=round] (135.344,280.200) -- (136.344,279.200);
\draw[color=tg0000ff,line width=0.38bp,line cap=round] (136.583,279.200) -- (137.583,280.200);
\draw[color=tg0000ff,line width=0.38bp,line cap=round] (136.583,280.200) -- (137.583,279.200);
\draw[color=tg0000ff,line width=0.38bp,line cap=round] (137.823,279.200) -- (138.823,280.200);
\draw[color=tg0000ff,line width=0.38bp,line cap=round] (137.823,280.200) -- (138.823,279.200);
\draw[color=tg0000ff,line width=0.38bp,line cap=round] (139.062,279.200) -- (140.062,280.200);
\draw[color=tg0000ff,line width=0.38bp,line cap=round] (139.062,280.200) -- (140.062,279.200);
\draw[color=tg0000ff,line width=0.38bp,line cap=round] (140.302,279.200) -- (141.302,280.200);
\draw[color=tg0000ff,line width=0.38bp,line cap=round] (140.302,280.200) -- (141.302,279.200);
\draw[color=tg0000ff,line width=0.38bp,line cap=round] (141.542,279.200) -- (142.542,280.200);
\draw[color=tg0000ff,line width=0.38bp,line cap=round] (141.542,280.200) -- (142.542,279.200);
\draw[color=tg0000ff,line width=0.38bp,line cap=round] (142.781,279.200) -- (143.781,280.200);
\draw[color=tg0000ff,line width=0.38bp,line cap=round] (142.781,280.200) -- (143.781,279.200);
\draw[color=tg0000ff,line width=0.38bp,line cap=round] (144.021,279.200) -- (145.021,280.200);
\draw[color=tg0000ff,line width=0.38bp,line cap=round] (144.021,280.200) -- (145.021,279.200);
\draw[color=tg0000ff,line width=0.38bp,line cap=round] (145.260,279.200) -- (146.260,280.200);
\draw[color=tg0000ff,line width=0.38bp,line cap=round] (145.260,280.200) -- (146.260,279.200);
\draw[color=tgc9cdd1,line width=0.45bp,line cap=round] (174.000,262.000) -- (174.000,305.000);
\draw[color=tgc9cdd1,line width=0.45bp,line cap=round] (174.000,305.000) -- (293.000,305.000);
\node[anchor=east,inner sep=0pt,outer sep=0pt,text=tg626b73,font={\sffamily\fontsize{6.7bp}{8.04bp}\selectfont}] at (171.000,298.980) {-0.05};
\node[anchor=east,inner sep=0pt,outer sep=0pt,text=tg626b73,font={\sffamily\fontsize{6.7bp}{8.04bp}\selectfont}] at (171.000,268.020) {0.07};
\draw[color=tg9da7b0,line width=0.6bp,line cap=round,dash pattern=on 2bp off 2bp] (263.250,262.000) -- (263.250,305.000);
\fill[tg808080] (174.000,288.776) circle[radius=0.52bp];
\fill[tg808080] (175.240,280.919) circle[radius=0.52bp];
\fill[tg808080] (176.479,287.327) circle[radius=0.52bp];
\fill[tg808080] (177.719,287.046) circle[radius=0.52bp];
\fill[tg808080] (178.958,285.647) circle[radius=0.52bp];
\fill[tg808080] (180.198,283.979) circle[radius=0.52bp];
\fill[tg808080] (181.438,282.547) circle[radius=0.52bp];
\fill[tg808080] (182.677,286.040) circle[radius=0.52bp];
\fill[tg808080] (183.917,283.471) circle[radius=0.52bp];
\fill[tg808080] (185.156,284.521) circle[radius=0.52bp];
\fill[tg808080] (186.396,285.572) circle[radius=0.52bp];
\fill[tg808080] (187.635,285.598) circle[radius=0.52bp];
\fill[tg808080] (188.875,287.324) circle[radius=0.52bp];
\fill[tg808080] (190.115,286.816) circle[radius=0.52bp];
\fill[tg808080] (191.354,284.848) circle[radius=0.52bp];
\fill[tg808080] (192.594,285.444) circle[radius=0.52bp];
\fill[tg808080] (193.833,287.102) circle[radius=0.52bp];
\fill[tg808080] (195.073,289.134) circle[radius=0.52bp];
\fill[tg808080] (196.312,283.740) circle[radius=0.52bp];
\fill[tg808080] (197.552,289.198) circle[radius=0.52bp];
\fill[tg808080] (198.792,284.563) circle[radius=0.52bp];
\fill[tg808080] (200.031,286.640) circle[radius=0.52bp];
\fill[tg808080] (201.271,285.895) circle[radius=0.52bp];
\fill[tg808080] (202.510,277.315) circle[radius=0.52bp];
\fill[tg808080] (203.750,278.594) circle[radius=0.52bp];
\fill[tg808080] (204.990,286.741) circle[radius=0.52bp];
\fill[tg808080] (206.229,274.001) circle[radius=0.52bp];
\fill[tg808080] (207.469,290.880) circle[radius=0.52bp];
\fill[tg808080] (208.708,285.836) circle[radius=0.52bp];
\fill[tg808080] (209.948,300.048) circle[radius=0.52bp];
\fill[tg808080] (211.188,281.421) circle[radius=0.52bp];
\fill[tg808080] (212.427,284.849) circle[radius=0.52bp];
\fill[tg808080] (213.667,289.519) circle[radius=0.52bp];
\fill[tg808080] (214.906,288.447) circle[radius=0.52bp];
\fill[tg808080] (216.146,290.473) circle[radius=0.52bp];
\fill[tg808080] (217.385,300.839) circle[radius=0.52bp];
\fill[tg808080] (218.625,278.939) circle[radius=0.52bp];
\fill[tg808080] (219.865,293.778) circle[radius=0.52bp];
\fill[tg808080] (221.104,289.234) circle[radius=0.52bp];
\fill[tg808080] (222.344,282.161) circle[radius=0.52bp];
\fill[tg808080] (223.583,283.072) circle[radius=0.52bp];
\fill[tg808080] (224.823,293.377) circle[radius=0.52bp];
\fill[tg808080] (226.062,297.115) circle[radius=0.52bp];
\fill[tg808080] (227.302,274.550) circle[radius=0.52bp];
\fill[tg808080] (228.542,297.129) circle[radius=0.52bp];
\fill[tg808080] (229.781,296.988) circle[radius=0.52bp];
\fill[tg808080] (231.021,288.741) circle[radius=0.52bp];
\fill[tg808080] (232.260,280.041) circle[radius=0.52bp];
\fill[tg808080] (233.500,282.683) circle[radius=0.52bp];
\fill[tg808080] (234.740,271.041) circle[radius=0.52bp];
\fill[tg808080] (235.979,299.051) circle[radius=0.52bp];
\fill[tg808080] (237.219,277.885) circle[radius=0.52bp];
\fill[tg808080] (238.458,286.897) circle[radius=0.52bp];
\fill[tg808080] (239.698,287.995) circle[radius=0.52bp];
\fill[tg808080] (240.938,286.046) circle[radius=0.52bp];
\fill[tg808080] (242.177,290.205) circle[radius=0.52bp];
\fill[tg808080] (243.417,284.784) circle[radius=0.52bp];
\fill[tg808080] (244.656,285.936) circle[radius=0.52bp];
\fill[tg808080] (245.896,287.025) circle[radius=0.52bp];
\fill[tg808080] (247.135,284.432) circle[radius=0.52bp];
\fill[tg808080] (248.375,289.274) circle[radius=0.52bp];
\fill[tg808080] (249.615,286.127) circle[radius=0.52bp];
\fill[tg808080] (250.854,285.900) circle[radius=0.52bp];
\fill[tg808080] (252.094,285.397) circle[radius=0.52bp];
\fill[tg808080] (253.333,286.638) circle[radius=0.52bp];
\fill[tg808080] (254.573,288.091) circle[radius=0.52bp];
\fill[tg808080] (255.812,287.887) circle[radius=0.52bp];
\fill[tg808080] (257.052,288.421) circle[radius=0.52bp];
\fill[tg808080] (258.292,286.071) circle[radius=0.52bp];
\fill[tg808080] (259.531,283.520) circle[radius=0.52bp];
\fill[tg808080] (260.771,286.456) circle[radius=0.52bp];
\fill[tg808080] (262.010,284.538) circle[radius=0.52bp];
\fill[tg808080] (263.250,287.400) circle[radius=0.52bp];
\fill[tg808080] (264.490,286.017) circle[radius=0.52bp];
\fill[tg808080] (265.729,286.334) circle[radius=0.52bp];
\fill[tg808080] (266.969,285.938) circle[radius=0.52bp];
\fill[tg808080] (268.208,287.157) circle[radius=0.52bp];
\fill[tg808080] (269.448,288.307) circle[radius=0.52bp];
\fill[tg808080] (270.688,284.934) circle[radius=0.52bp];
\fill[tg808080] (271.927,283.095) circle[radius=0.52bp];
\fill[tg808080] (273.167,286.323) circle[radius=0.52bp];
\fill[tg808080] (274.406,286.415) circle[radius=0.52bp];
\fill[tg808080] (275.646,283.992) circle[radius=0.52bp];
\fill[tg808080] (276.885,285.616) circle[radius=0.52bp];
\fill[tg808080] (278.125,288.142) circle[radius=0.52bp];
\fill[tg808080] (279.365,290.913) circle[radius=0.52bp];
\fill[tg808080] (280.604,281.076) circle[radius=0.52bp];
\fill[tg808080] (281.844,286.715) circle[radius=0.52bp];
\fill[tg808080] (283.083,283.823) circle[radius=0.52bp];
\fill[tg808080] (284.323,266.161) circle[radius=0.52bp];
\fill[tg808080] (285.562,281.281) circle[radius=0.52bp];
\fill[tg808080] (286.802,286.393) circle[radius=0.52bp];
\fill[tg808080] (288.042,281.299) circle[radius=0.52bp];
\fill[tg808080] (289.281,285.231) circle[radius=0.52bp];
\fill[tg808080] (290.521,294.549) circle[radius=0.52bp];
\fill[tg808080] (291.760,288.286) circle[radius=0.52bp];
\draw[color=tg575757,line width=0.65bp,line cap=round] (263.250,285.956) -- (264.490,285.965) -- (265.729,285.973) -- (266.969,285.980) -- (268.208,285.986) -- (269.448,285.991) -- (270.688,285.996) -- (271.927,286.000) -- (273.167,286.004) -- (274.406,286.008) -- (275.646,286.011) -- (276.885,286.013) -- (278.125,286.016) -- (279.365,286.018) -- (280.604,286.020) -- (281.844,286.022) -- (283.083,286.023) -- (284.323,286.024) -- (285.562,286.026) -- (286.802,286.027) -- (288.042,286.028) -- (289.281,286.029) -- (290.521,286.029) -- (291.760,286.030);
\draw[color=tg0000ff,line width=0.38bp,line cap=round] (262.750,284.038) -- (263.750,285.038);
\draw[color=tg0000ff,line width=0.38bp,line cap=round] (262.750,285.038) -- (263.750,284.038);
\draw[color=tg0000ff,line width=0.38bp,line cap=round] (263.990,284.038) -- (264.990,285.038);
\draw[color=tg0000ff,line width=0.38bp,line cap=round] (263.990,285.038) -- (264.990,284.038);
\draw[color=tg0000ff,line width=0.38bp,line cap=round] (265.229,284.038) -- (266.229,285.038);
\draw[color=tg0000ff,line width=0.38bp,line cap=round] (265.229,285.038) -- (266.229,284.038);
\draw[color=tg0000ff,line width=0.38bp,line cap=round] (266.469,284.038) -- (267.469,285.038);
\draw[color=tg0000ff,line width=0.38bp,line cap=round] (266.469,285.038) -- (267.469,284.038);
\draw[color=tg0000ff,line width=0.38bp,line cap=round] (267.708,284.038) -- (268.708,285.038);
\draw[color=tg0000ff,line width=0.38bp,line cap=round] (267.708,285.038) -- (268.708,284.038);
\draw[color=tg0000ff,line width=0.38bp,line cap=round] (268.948,284.038) -- (269.948,285.038);
\draw[color=tg0000ff,line width=0.38bp,line cap=round] (268.948,285.038) -- (269.948,284.038);
\draw[color=tg0000ff,line width=0.38bp,line cap=round] (270.188,284.038) -- (271.188,285.038);
\draw[color=tg0000ff,line width=0.38bp,line cap=round] (270.188,285.038) -- (271.188,284.038);
\draw[color=tg0000ff,line width=0.38bp,line cap=round] (271.427,284.038) -- (272.427,285.038);
\draw[color=tg0000ff,line width=0.38bp,line cap=round] (271.427,285.038) -- (272.427,284.038);
\draw[color=tg0000ff,line width=0.38bp,line cap=round] (272.667,284.038) -- (273.667,285.038);
\draw[color=tg0000ff,line width=0.38bp,line cap=round] (272.667,285.038) -- (273.667,284.038);
\draw[color=tg0000ff,line width=0.38bp,line cap=round] (273.906,284.038) -- (274.906,285.038);
\draw[color=tg0000ff,line width=0.38bp,line cap=round] (273.906,285.038) -- (274.906,284.038);
\draw[color=tg0000ff,line width=0.38bp,line cap=round] (275.146,284.038) -- (276.146,285.038);
\draw[color=tg0000ff,line width=0.38bp,line cap=round] (275.146,285.038) -- (276.146,284.038);
\draw[color=tg0000ff,line width=0.38bp,line cap=round] (276.385,284.038) -- (277.385,285.038);
\draw[color=tg0000ff,line width=0.38bp,line cap=round] (276.385,285.038) -- (277.385,284.038);
\draw[color=tg0000ff,line width=0.38bp,line cap=round] (277.625,284.038) -- (278.625,285.038);
\draw[color=tg0000ff,line width=0.38bp,line cap=round] (277.625,285.038) -- (278.625,284.038);
\draw[color=tg0000ff,line width=0.38bp,line cap=round] (278.865,284.038) -- (279.865,285.038);
\draw[color=tg0000ff,line width=0.38bp,line cap=round] (278.865,285.038) -- (279.865,284.038);
\draw[color=tg0000ff,line width=0.38bp,line cap=round] (280.104,284.038) -- (281.104,285.038);
\draw[color=tg0000ff,line width=0.38bp,line cap=round] (280.104,285.038) -- (281.104,284.038);
\draw[color=tg0000ff,line width=0.38bp,line cap=round] (281.344,284.038) -- (282.344,285.038);
\draw[color=tg0000ff,line width=0.38bp,line cap=round] (281.344,285.038) -- (282.344,284.038);
\draw[color=tg0000ff,line width=0.38bp,line cap=round] (282.583,284.038) -- (283.583,285.038);
\draw[color=tg0000ff,line width=0.38bp,line cap=round] (282.583,285.038) -- (283.583,284.038);
\draw[color=tg0000ff,line width=0.38bp,line cap=round] (283.823,284.038) -- (284.823,285.038);
\draw[color=tg0000ff,line width=0.38bp,line cap=round] (283.823,285.038) -- (284.823,284.038);
\draw[color=tg0000ff,line width=0.38bp,line cap=round] (285.062,284.038) -- (286.062,285.038);
\draw[color=tg0000ff,line width=0.38bp,line cap=round] (285.062,285.038) -- (286.062,284.038);
\draw[color=tg0000ff,line width=0.38bp,line cap=round] (286.302,284.038) -- (287.302,285.038);
\draw[color=tg0000ff,line width=0.38bp,line cap=round] (286.302,285.038) -- (287.302,284.038);
\draw[color=tg0000ff,line width=0.38bp,line cap=round] (287.542,284.038) -- (288.542,285.038);
\draw[color=tg0000ff,line width=0.38bp,line cap=round] (287.542,285.038) -- (288.542,284.038);
\draw[color=tg0000ff,line width=0.38bp,line cap=round] (288.781,284.038) -- (289.781,285.038);
\draw[color=tg0000ff,line width=0.38bp,line cap=round] (288.781,285.038) -- (289.781,284.038);
\draw[color=tg0000ff,line width=0.38bp,line cap=round] (290.021,284.038) -- (291.021,285.038);
\draw[color=tg0000ff,line width=0.38bp,line cap=round] (290.021,285.038) -- (291.021,284.038);
\draw[color=tg0000ff,line width=0.38bp,line cap=round] (291.260,284.038) -- (292.260,285.038);
\draw[color=tg0000ff,line width=0.38bp,line cap=round] (291.260,285.038) -- (292.260,284.038);
\draw[color=tgc9cdd1,line width=0.45bp,line cap=round] (320.000,262.000) -- (320.000,305.000);
\draw[color=tgc9cdd1,line width=0.45bp,line cap=round] (320.000,305.000) -- (439.000,305.000);
\node[anchor=east,inner sep=0pt,outer sep=0pt,text=tg626b73,font={\sffamily\fontsize{6.7bp}{8.04bp}\selectfont}] at (317.000,298.980) {-0.022};
\node[anchor=east,inner sep=0pt,outer sep=0pt,text=tg626b73,font={\sffamily\fontsize{6.7bp}{8.04bp}\selectfont}] at (317.000,268.020) {0.012};
\draw[color=tg9da7b0,line width=0.6bp,line cap=round,dash pattern=on 2bp off 2bp] (409.250,262.000) -- (409.250,305.000);
\fill[tg808080] (320.000,277.202) circle[radius=0.52bp];
\fill[tg808080] (321.240,275.139) circle[radius=0.52bp];
\fill[tg808080] (322.479,275.052) circle[radius=0.52bp];
\fill[tg808080] (323.719,275.000) circle[radius=0.52bp];
\fill[tg808080] (324.958,282.505) circle[radius=0.52bp];
\fill[tg808080] (326.198,284.248) circle[radius=0.52bp];
\fill[tg808080] (327.438,284.739) circle[radius=0.52bp];
\fill[tg808080] (328.677,269.858) circle[radius=0.52bp];
\fill[tg808080] (329.917,284.790) circle[radius=0.52bp];
\fill[tg808080] (331.156,286.750) circle[radius=0.52bp];
\fill[tg808080] (332.396,281.249) circle[radius=0.52bp];
\fill[tg808080] (333.635,286.884) circle[radius=0.52bp];
\fill[tg808080] (334.875,284.545) circle[radius=0.52bp];
\fill[tg808080] (336.115,282.095) circle[radius=0.52bp];
\fill[tg808080] (337.354,274.915) circle[radius=0.52bp];
\fill[tg808080] (338.594,274.612) circle[radius=0.52bp];
\fill[tg808080] (339.833,275.318) circle[radius=0.52bp];
\fill[tg808080] (341.073,282.177) circle[radius=0.52bp];
\fill[tg808080] (342.312,275.254) circle[radius=0.52bp];
\fill[tg808080] (343.552,278.355) circle[radius=0.52bp];
\fill[tg808080] (344.792,273.465) circle[radius=0.52bp];
\fill[tg808080] (346.031,283.036) circle[radius=0.52bp];
\fill[tg808080] (347.271,272.672) circle[radius=0.52bp];
\fill[tg808080] (348.510,273.286) circle[radius=0.52bp];
\fill[tg808080] (349.750,275.483) circle[radius=0.52bp];
\fill[tg808080] (350.990,272.847) circle[radius=0.52bp];
\fill[tg808080] (352.229,274.045) circle[radius=0.52bp];
\fill[tg808080] (353.469,274.945) circle[radius=0.52bp];
\fill[tg808080] (354.708,272.779) circle[radius=0.52bp];
\fill[tg808080] (355.948,284.438) circle[radius=0.52bp];
\fill[tg808080] (357.188,277.074) circle[radius=0.52bp];
\fill[tg808080] (358.427,276.469) circle[radius=0.52bp];
\fill[tg808080] (359.667,270.157) circle[radius=0.52bp];
\fill[tg808080] (360.906,277.227) circle[radius=0.52bp];
\fill[tg808080] (362.146,294.416) circle[radius=0.52bp];
\fill[tg808080] (363.385,281.354) circle[radius=0.52bp];
\fill[tg808080] (364.625,279.219) circle[radius=0.52bp];
\fill[tg808080] (365.865,288.546) circle[radius=0.52bp];
\fill[tg808080] (367.104,278.247) circle[radius=0.52bp];
\fill[tg808080] (368.344,274.497) circle[radius=0.52bp];
\fill[tg808080] (369.583,267.671) circle[radius=0.52bp];
\fill[tg808080] (370.823,294.967) circle[radius=0.52bp];
\fill[tg808080] (372.062,285.919) circle[radius=0.52bp];
\fill[tg808080] (373.302,280.408) circle[radius=0.52bp];
\fill[tg808080] (374.542,272.619) circle[radius=0.52bp];
\fill[tg808080] (375.781,280.786) circle[radius=0.52bp];
\fill[tg808080] (377.021,273.270) circle[radius=0.52bp];
\fill[tg808080] (378.260,300.839) circle[radius=0.52bp];
\fill[tg808080] (379.500,297.811) circle[radius=0.52bp];
\fill[tg808080] (380.740,275.592) circle[radius=0.52bp];
\fill[tg808080] (381.979,279.937) circle[radius=0.52bp];
\fill[tg808080] (383.219,284.193) circle[radius=0.52bp];
\fill[tg808080] (384.458,271.174) circle[radius=0.52bp];
\fill[tg808080] (385.698,271.757) circle[radius=0.52bp];
\fill[tg808080] (386.938,275.586) circle[radius=0.52bp];
\fill[tg808080] (388.177,277.111) circle[radius=0.52bp];
\fill[tg808080] (389.417,266.161) circle[radius=0.52bp];
\fill[tg808080] (390.656,269.352) circle[radius=0.52bp];
\fill[tg808080] (391.896,280.094) circle[radius=0.52bp];
\fill[tg808080] (393.135,283.809) circle[radius=0.52bp];
\fill[tg808080] (394.375,281.733) circle[radius=0.52bp];
\fill[tg808080] (395.615,279.549) circle[radius=0.52bp];
\fill[tg808080] (396.854,288.585) circle[radius=0.52bp];
\fill[tg808080] (398.094,269.684) circle[radius=0.52bp];
\fill[tg808080] (399.333,280.196) circle[radius=0.52bp];
\fill[tg808080] (400.573,281.008) circle[radius=0.52bp];
\fill[tg808080] (401.812,279.790) circle[radius=0.52bp];
\fill[tg808080] (403.052,277.850) circle[radius=0.52bp];
\fill[tg808080] (404.292,274.609) circle[radius=0.52bp];
\fill[tg808080] (405.531,296.306) circle[radius=0.52bp];
\fill[tg808080] (406.771,278.737) circle[radius=0.52bp];
\fill[tg808080] (408.010,293.013) circle[radius=0.52bp];
\fill[tg808080] (409.250,271.190) circle[radius=0.52bp];
\fill[tg808080] (410.490,273.361) circle[radius=0.52bp];
\fill[tg808080] (411.729,289.068) circle[radius=0.52bp];
\fill[tg808080] (412.969,276.511) circle[radius=0.52bp];
\fill[tg808080] (414.208,278.255) circle[radius=0.52bp];
\fill[tg808080] (415.448,283.171) circle[radius=0.52bp];
\fill[tg808080] (416.688,285.683) circle[radius=0.52bp];
\fill[tg808080] (417.927,277.172) circle[radius=0.52bp];
\fill[tg808080] (419.167,277.203) circle[radius=0.52bp];
\fill[tg808080] (420.406,277.645) circle[radius=0.52bp];
\fill[tg808080] (421.646,287.509) circle[radius=0.52bp];
\fill[tg808080] (422.885,276.018) circle[radius=0.52bp];
\fill[tg808080] (424.125,284.442) circle[radius=0.52bp];
\fill[tg808080] (425.365,288.453) circle[radius=0.52bp];
\fill[tg808080] (426.604,278.000) circle[radius=0.52bp];
\fill[tg808080] (427.844,283.454) circle[radius=0.52bp];
\fill[tg808080] (429.083,278.517) circle[radius=0.52bp];
\fill[tg808080] (430.323,291.254) circle[radius=0.52bp];
\fill[tg808080] (431.562,277.365) circle[radius=0.52bp];
\fill[tg808080] (432.802,275.735) circle[radius=0.52bp];
\fill[tg808080] (434.042,280.745) circle[radius=0.52bp];
\fill[tg808080] (435.281,283.074) circle[radius=0.52bp];
\fill[tg808080] (436.521,289.037) circle[radius=0.52bp];
\fill[tg808080] (437.760,284.613) circle[radius=0.52bp];
\draw[color=tg575757,line width=0.65bp,line cap=round] (409.250,279.156) -- (410.490,279.159) -- (411.729,279.162) -- (412.969,279.164) -- (414.208,279.167) -- (415.448,279.169) -- (416.688,279.170) -- (417.927,279.172) -- (419.167,279.173) -- (420.406,279.175) -- (421.646,279.176) -- (422.885,279.177) -- (424.125,279.178) -- (425.365,279.178) -- (426.604,279.179) -- (427.844,279.180) -- (429.083,279.180) -- (430.323,279.181) -- (431.562,279.181) -- (432.802,279.182) -- (434.042,279.182) -- (435.281,279.182) -- (436.521,279.183) -- (437.760,279.183);
\draw[color=tg0000ff,line width=0.38bp,line cap=round] (408.750,292.513) -- (409.750,293.513);
\draw[color=tg0000ff,line width=0.38bp,line cap=round] (408.750,293.513) -- (409.750,292.513);
\draw[color=tg0000ff,line width=0.38bp,line cap=round] (409.990,292.513) -- (410.990,293.513);
\draw[color=tg0000ff,line width=0.38bp,line cap=round] (409.990,293.513) -- (410.990,292.513);
\draw[color=tg0000ff,line width=0.38bp,line cap=round] (411.229,292.513) -- (412.229,293.513);
\draw[color=tg0000ff,line width=0.38bp,line cap=round] (411.229,293.513) -- (412.229,292.513);
\draw[color=tg0000ff,line width=0.38bp,line cap=round] (412.469,292.513) -- (413.469,293.513);
\draw[color=tg0000ff,line width=0.38bp,line cap=round] (412.469,293.513) -- (413.469,292.513);
\draw[color=tg0000ff,line width=0.38bp,line cap=round] (413.708,292.513) -- (414.708,293.513);
\draw[color=tg0000ff,line width=0.38bp,line cap=round] (413.708,293.513) -- (414.708,292.513);
\draw[color=tg0000ff,line width=0.38bp,line cap=round] (414.948,292.513) -- (415.948,293.513);
\draw[color=tg0000ff,line width=0.38bp,line cap=round] (414.948,293.513) -- (415.948,292.513);
\draw[color=tg0000ff,line width=0.38bp,line cap=round] (416.188,292.513) -- (417.188,293.513);
\draw[color=tg0000ff,line width=0.38bp,line cap=round] (416.188,293.513) -- (417.188,292.513);
\draw[color=tg0000ff,line width=0.38bp,line cap=round] (417.427,292.513) -- (418.427,293.513);
\draw[color=tg0000ff,line width=0.38bp,line cap=round] (417.427,293.513) -- (418.427,292.513);
\draw[color=tg0000ff,line width=0.38bp,line cap=round] (418.667,292.513) -- (419.667,293.513);
\draw[color=tg0000ff,line width=0.38bp,line cap=round] (418.667,293.513) -- (419.667,292.513);
\draw[color=tg0000ff,line width=0.38bp,line cap=round] (419.906,292.513) -- (420.906,293.513);
\draw[color=tg0000ff,line width=0.38bp,line cap=round] (419.906,293.513) -- (420.906,292.513);
\draw[color=tg0000ff,line width=0.38bp,line cap=round] (421.146,292.513) -- (422.146,293.513);
\draw[color=tg0000ff,line width=0.38bp,line cap=round] (421.146,293.513) -- (422.146,292.513);
\draw[color=tg0000ff,line width=0.38bp,line cap=round] (422.385,292.513) -- (423.385,293.513);
\draw[color=tg0000ff,line width=0.38bp,line cap=round] (422.385,293.513) -- (423.385,292.513);
\draw[color=tg0000ff,line width=0.38bp,line cap=round] (423.625,292.513) -- (424.625,293.513);
\draw[color=tg0000ff,line width=0.38bp,line cap=round] (423.625,293.513) -- (424.625,292.513);
\draw[color=tg0000ff,line width=0.38bp,line cap=round] (424.865,292.513) -- (425.865,293.513);
\draw[color=tg0000ff,line width=0.38bp,line cap=round] (424.865,293.513) -- (425.865,292.513);
\draw[color=tg0000ff,line width=0.38bp,line cap=round] (426.104,292.513) -- (427.104,293.513);
\draw[color=tg0000ff,line width=0.38bp,line cap=round] (426.104,293.513) -- (427.104,292.513);
\draw[color=tg0000ff,line width=0.38bp,line cap=round] (427.344,292.513) -- (428.344,293.513);
\draw[color=tg0000ff,line width=0.38bp,line cap=round] (427.344,293.513) -- (428.344,292.513);
\draw[color=tg0000ff,line width=0.38bp,line cap=round] (428.583,292.513) -- (429.583,293.513);
\draw[color=tg0000ff,line width=0.38bp,line cap=round] (428.583,293.513) -- (429.583,292.513);
\draw[color=tg0000ff,line width=0.38bp,line cap=round] (429.823,292.513) -- (430.823,293.513);
\draw[color=tg0000ff,line width=0.38bp,line cap=round] (429.823,293.513) -- (430.823,292.513);
\draw[color=tg0000ff,line width=0.38bp,line cap=round] (431.062,292.513) -- (432.062,293.513);
\draw[color=tg0000ff,line width=0.38bp,line cap=round] (431.062,293.513) -- (432.062,292.513);
\draw[color=tg0000ff,line width=0.38bp,line cap=round] (432.302,292.513) -- (433.302,293.513);
\draw[color=tg0000ff,line width=0.38bp,line cap=round] (432.302,293.513) -- (433.302,292.513);
\draw[color=tg0000ff,line width=0.38bp,line cap=round] (433.542,292.513) -- (434.542,293.513);
\draw[color=tg0000ff,line width=0.38bp,line cap=round] (433.542,293.513) -- (434.542,292.513);
\draw[color=tg0000ff,line width=0.38bp,line cap=round] (434.781,292.513) -- (435.781,293.513);
\draw[color=tg0000ff,line width=0.38bp,line cap=round] (434.781,293.513) -- (435.781,292.513);
\draw[color=tg0000ff,line width=0.38bp,line cap=round] (436.021,292.513) -- (437.021,293.513);
\draw[color=tg0000ff,line width=0.38bp,line cap=round] (436.021,293.513) -- (437.021,292.513);
\draw[color=tg0000ff,line width=0.38bp,line cap=round] (437.260,292.513) -- (438.260,293.513);
\draw[color=tg0000ff,line width=0.38bp,line cap=round] (437.260,293.513) -- (438.260,292.513);
\node[anchor=west,inner sep=0pt,outer sep=0pt,text=tg18314b,font={\sffamily\fontsize{8bp}{9.6bp}\selectfont\bfseries}] at (0.000,320.000) {Hidden regimes + jumps (log returns)};
\node[anchor=east,inner sep=0pt,outer sep=0pt,text=tg626b73,font={\sffamily\fontsize{7bp}{8.4bp}\selectfont}] at (448.000,320.000) {Approximate filtered reference};
\draw[color=tgc9cdd1,line width=0.45bp,line cap=round] (28.000,329.000) -- (28.000,372.000);
\draw[color=tgc9cdd1,line width=0.45bp,line cap=round] (28.000,372.000) -- (147.000,372.000);
\node[anchor=east,inner sep=0pt,outer sep=0pt,text=tg626b73,font={\sffamily\fontsize{6.7bp}{8.04bp}\selectfont}] at (25.000,365.980) {-0.23};
\node[anchor=east,inner sep=0pt,outer sep=0pt,text=tg626b73,font={\sffamily\fontsize{6.7bp}{8.04bp}\selectfont}] at (25.000,335.020) {0.03};
\draw[color=tg9da7b0,line width=0.6bp,line cap=round,dash pattern=on 2bp off 2bp] (117.250,329.000) -- (117.250,372.000);
\fill[tg808080] (28.000,338.843) circle[radius=0.52bp];
\fill[tg808080] (29.240,339.318) circle[radius=0.52bp];
\fill[tg808080] (30.479,337.390) circle[radius=0.52bp];
\fill[tg808080] (31.719,338.437) circle[radius=0.52bp];
\fill[tg808080] (32.958,338.587) circle[radius=0.52bp];
\fill[tg808080] (34.198,337.753) circle[radius=0.52bp];
\fill[tg808080] (35.438,338.463) circle[radius=0.52bp];
\fill[tg808080] (36.677,338.923) circle[radius=0.52bp];
\fill[tg808080] (37.917,337.238) circle[radius=0.52bp];
\fill[tg808080] (39.156,337.882) circle[radius=0.52bp];
\fill[tg808080] (40.396,338.250) circle[radius=0.52bp];
\fill[tg808080] (41.635,337.477) circle[radius=0.52bp];
\fill[tg808080] (42.875,338.500) circle[radius=0.52bp];
\fill[tg808080] (44.115,338.229) circle[radius=0.52bp];
\fill[tg808080] (45.354,338.022) circle[radius=0.52bp];
\fill[tg808080] (46.594,336.385) circle[radius=0.52bp];
\fill[tg808080] (47.833,338.295) circle[radius=0.52bp];
\fill[tg808080] (49.073,337.123) circle[radius=0.52bp];
\fill[tg808080] (50.312,338.032) circle[radius=0.52bp];
\fill[tg808080] (51.552,338.609) circle[radius=0.52bp];
\fill[tg808080] (52.792,338.835) circle[radius=0.52bp];
\fill[tg808080] (54.031,338.326) circle[radius=0.52bp];
\fill[tg808080] (55.271,338.473) circle[radius=0.52bp];
\fill[tg808080] (56.510,339.205) circle[radius=0.52bp];
\fill[tg808080] (57.750,338.003) circle[radius=0.52bp];
\fill[tg808080] (58.990,336.935) circle[radius=0.52bp];
\fill[tg808080] (60.229,337.700) circle[radius=0.52bp];
\fill[tg808080] (61.469,338.916) circle[radius=0.52bp];
\fill[tg808080] (62.708,337.378) circle[radius=0.52bp];
\fill[tg808080] (63.948,338.679) circle[radius=0.52bp];
\fill[tg808080] (65.188,339.247) circle[radius=0.52bp];
\fill[tg808080] (66.427,338.150) circle[radius=0.52bp];
\fill[tg808080] (67.667,338.644) circle[radius=0.52bp];
\fill[tg808080] (68.906,338.724) circle[radius=0.52bp];
\fill[tg808080] (70.146,337.825) circle[radius=0.52bp];
\fill[tg808080] (71.385,337.684) circle[radius=0.52bp];
\fill[tg808080] (72.625,338.050) circle[radius=0.52bp];
\fill[tg808080] (73.865,338.781) circle[radius=0.52bp];
\fill[tg808080] (75.104,338.326) circle[radius=0.52bp];
\fill[tg808080] (76.344,337.788) circle[radius=0.52bp];
\fill[tg808080] (77.583,337.303) circle[radius=0.52bp];
\fill[tg808080] (78.823,337.754) circle[radius=0.52bp];
\fill[tg808080] (80.062,337.323) circle[radius=0.52bp];
\fill[tg808080] (81.302,337.779) circle[radius=0.52bp];
\fill[tg808080] (82.542,337.507) circle[radius=0.52bp];
\fill[tg808080] (83.781,338.570) circle[radius=0.52bp];
\fill[tg808080] (85.021,338.045) circle[radius=0.52bp];
\fill[tg808080] (86.260,339.177) circle[radius=0.52bp];
\fill[tg808080] (87.500,338.010) circle[radius=0.52bp];
\fill[tg808080] (88.740,338.078) circle[radius=0.52bp];
\fill[tg808080] (89.979,338.739) circle[radius=0.52bp];
\fill[tg808080] (91.219,338.384) circle[radius=0.52bp];
\fill[tg808080] (92.458,339.045) circle[radius=0.52bp];
\fill[tg808080] (93.698,338.032) circle[radius=0.52bp];
\fill[tg808080] (94.938,337.528) circle[radius=0.52bp];
\fill[tg808080] (96.177,337.670) circle[radius=0.52bp];
\fill[tg808080] (97.417,337.153) circle[radius=0.52bp];
\fill[tg808080] (98.656,338.009) circle[radius=0.52bp];
\fill[tg808080] (99.896,345.266) circle[radius=0.52bp];
\fill[tg808080] (101.135,338.088) circle[radius=0.52bp];
\fill[tg808080] (102.375,337.660) circle[radius=0.52bp];
\fill[tg808080] (103.615,339.602) circle[radius=0.52bp];
\fill[tg808080] (104.854,338.404) circle[radius=0.52bp];
\fill[tg808080] (106.094,338.049) circle[radius=0.52bp];
\fill[tg808080] (107.333,338.011) circle[radius=0.52bp];
\fill[tg808080] (108.573,337.939) circle[radius=0.52bp];
\fill[tg808080] (109.812,337.345) circle[radius=0.52bp];
\fill[tg808080] (111.052,337.966) circle[radius=0.52bp];
\fill[tg808080] (112.292,337.611) circle[radius=0.52bp];
\fill[tg808080] (113.531,338.538) circle[radius=0.52bp];
\fill[tg808080] (114.771,344.540) circle[radius=0.52bp];
\fill[tg808080] (116.010,339.052) circle[radius=0.52bp];
\fill[tg808080] (117.250,336.126) circle[radius=0.52bp];
\fill[tg808080] (118.490,337.396) circle[radius=0.52bp];
\fill[tg808080] (119.729,341.352) circle[radius=0.52bp];
\fill[tg808080] (120.969,335.365) circle[radius=0.52bp];
\fill[tg808080] (122.208,341.045) circle[radius=0.52bp];
\fill[tg808080] (123.448,348.249) circle[radius=0.52bp];
\fill[tg808080] (124.688,367.839) circle[radius=0.52bp];
\fill[tg808080] (125.927,339.861) circle[radius=0.52bp];
\fill[tg808080] (127.167,335.449) circle[radius=0.52bp];
\fill[tg808080] (128.406,344.260) circle[radius=0.52bp];
\fill[tg808080] (129.646,337.159) circle[radius=0.52bp];
\fill[tg808080] (130.885,338.800) circle[radius=0.52bp];
\fill[tg808080] (132.125,335.434) circle[radius=0.52bp];
\fill[tg808080] (133.365,338.399) circle[radius=0.52bp];
\fill[tg808080] (134.604,338.031) circle[radius=0.52bp];
\fill[tg808080] (135.844,338.787) circle[radius=0.52bp];
\fill[tg808080] (137.083,333.161) circle[radius=0.52bp];
\fill[tg808080] (138.323,337.304) circle[radius=0.52bp];
\fill[tg808080] (139.562,334.501) circle[radius=0.52bp];
\fill[tg808080] (140.802,336.620) circle[radius=0.52bp];
\fill[tg808080] (142.042,336.726) circle[radius=0.52bp];
\fill[tg808080] (143.281,337.585) circle[radius=0.52bp];
\fill[tg808080] (144.521,340.872) circle[radius=0.52bp];
\fill[tg808080] (145.760,337.372) circle[radius=0.52bp];
\draw[color=tg575757,line width=0.65bp,line cap=round] (117.250,338.184) -- (118.490,338.179) -- (119.729,338.176) -- (120.969,338.172) -- (122.208,338.169) -- (123.448,338.167) -- (124.688,338.164) -- (125.927,338.162) -- (127.167,338.160) -- (128.406,338.159) -- (129.646,338.157) -- (130.885,338.156) -- (132.125,338.155) -- (133.365,338.154) -- (134.604,338.153) -- (135.844,338.152) -- (137.083,338.151) -- (138.323,338.150) -- (139.562,338.150) -- (140.802,338.149) -- (142.042,338.149) -- (143.281,338.148) -- (144.521,338.148) -- (145.760,338.148);
\draw[color=tg0000ff,line width=0.38bp,line cap=round] (116.750,338.552) -- (117.750,339.552);
\draw[color=tg0000ff,line width=0.38bp,line cap=round] (116.750,339.552) -- (117.750,338.552);
\draw[color=tg0000ff,line width=0.38bp,line cap=round] (117.990,338.552) -- (118.990,339.552);
\draw[color=tg0000ff,line width=0.38bp,line cap=round] (117.990,339.552) -- (118.990,338.552);
\draw[color=tg0000ff,line width=0.38bp,line cap=round] (119.229,338.552) -- (120.229,339.552);
\draw[color=tg0000ff,line width=0.38bp,line cap=round] (119.229,339.552) -- (120.229,338.552);
\draw[color=tg0000ff,line width=0.38bp,line cap=round] (120.469,338.552) -- (121.469,339.552);
\draw[color=tg0000ff,line width=0.38bp,line cap=round] (120.469,339.552) -- (121.469,338.552);
\draw[color=tg0000ff,line width=0.38bp,line cap=round] (121.708,338.552) -- (122.708,339.552);
\draw[color=tg0000ff,line width=0.38bp,line cap=round] (121.708,339.552) -- (122.708,338.552);
\draw[color=tg0000ff,line width=0.38bp,line cap=round] (122.948,338.552) -- (123.948,339.552);
\draw[color=tg0000ff,line width=0.38bp,line cap=round] (122.948,339.552) -- (123.948,338.552);
\draw[color=tg0000ff,line width=0.38bp,line cap=round] (124.188,338.552) -- (125.188,339.552);
\draw[color=tg0000ff,line width=0.38bp,line cap=round] (124.188,339.552) -- (125.188,338.552);
\draw[color=tg0000ff,line width=0.38bp,line cap=round] (125.427,338.552) -- (126.427,339.552);
\draw[color=tg0000ff,line width=0.38bp,line cap=round] (125.427,339.552) -- (126.427,338.552);
\draw[color=tg0000ff,line width=0.38bp,line cap=round] (126.667,338.552) -- (127.667,339.552);
\draw[color=tg0000ff,line width=0.38bp,line cap=round] (126.667,339.552) -- (127.667,338.552);
\draw[color=tg0000ff,line width=0.38bp,line cap=round] (127.906,338.552) -- (128.906,339.552);
\draw[color=tg0000ff,line width=0.38bp,line cap=round] (127.906,339.552) -- (128.906,338.552);
\draw[color=tg0000ff,line width=0.38bp,line cap=round] (129.146,338.552) -- (130.146,339.552);
\draw[color=tg0000ff,line width=0.38bp,line cap=round] (129.146,339.552) -- (130.146,338.552);
\draw[color=tg0000ff,line width=0.38bp,line cap=round] (130.385,338.552) -- (131.385,339.552);
\draw[color=tg0000ff,line width=0.38bp,line cap=round] (130.385,339.552) -- (131.385,338.552);
\draw[color=tg0000ff,line width=0.38bp,line cap=round] (131.625,338.552) -- (132.625,339.552);
\draw[color=tg0000ff,line width=0.38bp,line cap=round] (131.625,339.552) -- (132.625,338.552);
\draw[color=tg0000ff,line width=0.38bp,line cap=round] (132.865,338.552) -- (133.865,339.552);
\draw[color=tg0000ff,line width=0.38bp,line cap=round] (132.865,339.552) -- (133.865,338.552);
\draw[color=tg0000ff,line width=0.38bp,line cap=round] (134.104,338.552) -- (135.104,339.552);
\draw[color=tg0000ff,line width=0.38bp,line cap=round] (134.104,339.552) -- (135.104,338.552);
\draw[color=tg0000ff,line width=0.38bp,line cap=round] (135.344,338.552) -- (136.344,339.552);
\draw[color=tg0000ff,line width=0.38bp,line cap=round] (135.344,339.552) -- (136.344,338.552);
\draw[color=tg0000ff,line width=0.38bp,line cap=round] (136.583,338.552) -- (137.583,339.552);
\draw[color=tg0000ff,line width=0.38bp,line cap=round] (136.583,339.552) -- (137.583,338.552);
\draw[color=tg0000ff,line width=0.38bp,line cap=round] (137.823,338.552) -- (138.823,339.552);
\draw[color=tg0000ff,line width=0.38bp,line cap=round] (137.823,339.552) -- (138.823,338.552);
\draw[color=tg0000ff,line width=0.38bp,line cap=round] (139.062,338.552) -- (140.062,339.552);
\draw[color=tg0000ff,line width=0.38bp,line cap=round] (139.062,339.552) -- (140.062,338.552);
\draw[color=tg0000ff,line width=0.38bp,line cap=round] (140.302,338.552) -- (141.302,339.552);
\draw[color=tg0000ff,line width=0.38bp,line cap=round] (140.302,339.552) -- (141.302,338.552);
\draw[color=tg0000ff,line width=0.38bp,line cap=round] (141.542,338.552) -- (142.542,339.552);
\draw[color=tg0000ff,line width=0.38bp,line cap=round] (141.542,339.552) -- (142.542,338.552);
\draw[color=tg0000ff,line width=0.38bp,line cap=round] (142.781,338.552) -- (143.781,339.552);
\draw[color=tg0000ff,line width=0.38bp,line cap=round] (142.781,339.552) -- (143.781,338.552);
\draw[color=tg0000ff,line width=0.38bp,line cap=round] (144.021,338.552) -- (145.021,339.552);
\draw[color=tg0000ff,line width=0.38bp,line cap=round] (144.021,339.552) -- (145.021,338.552);
\draw[color=tg0000ff,line width=0.38bp,line cap=round] (145.260,338.552) -- (146.260,339.552);
\draw[color=tg0000ff,line width=0.38bp,line cap=round] (145.260,339.552) -- (146.260,338.552);
\draw[color=tgc9cdd1,line width=0.45bp,line cap=round] (174.000,329.000) -- (174.000,372.000);
\draw[color=tgc9cdd1,line width=0.45bp,line cap=round] (174.000,372.000) -- (293.000,372.000);
\node[anchor=east,inner sep=0pt,outer sep=0pt,text=tg626b73,font={\sffamily\fontsize{6.7bp}{8.04bp}\selectfont}] at (171.000,365.980) {-0.05};
\node[anchor=east,inner sep=0pt,outer sep=0pt,text=tg626b73,font={\sffamily\fontsize{6.7bp}{8.04bp}\selectfont}] at (171.000,335.020) {0.07};
\draw[color=tg9da7b0,line width=0.6bp,line cap=round,dash pattern=on 2bp off 2bp] (263.250,329.000) -- (263.250,372.000);
\fill[tg808080] (174.000,356.464) circle[radius=0.52bp];
\fill[tg808080] (175.240,348.607) circle[radius=0.52bp];
\fill[tg808080] (176.479,355.015) circle[radius=0.52bp];
\fill[tg808080] (177.719,354.734) circle[radius=0.52bp];
\fill[tg808080] (178.958,353.335) circle[radius=0.52bp];
\fill[tg808080] (180.198,351.667) circle[radius=0.52bp];
\fill[tg808080] (181.438,350.235) circle[radius=0.52bp];
\fill[tg808080] (182.677,353.728) circle[radius=0.52bp];
\fill[tg808080] (183.917,351.159) circle[radius=0.52bp];
\fill[tg808080] (185.156,352.209) circle[radius=0.52bp];
\fill[tg808080] (186.396,353.260) circle[radius=0.52bp];
\fill[tg808080] (187.635,353.285) circle[radius=0.52bp];
\fill[tg808080] (188.875,355.012) circle[radius=0.52bp];
\fill[tg808080] (190.115,354.504) circle[radius=0.52bp];
\fill[tg808080] (191.354,352.536) circle[radius=0.52bp];
\fill[tg808080] (192.594,353.132) circle[radius=0.52bp];
\fill[tg808080] (193.833,354.790) circle[radius=0.52bp];
\fill[tg808080] (195.073,356.822) circle[radius=0.52bp];
\fill[tg808080] (196.312,351.428) circle[radius=0.52bp];
\fill[tg808080] (197.552,356.886) circle[radius=0.52bp];
\fill[tg808080] (198.792,352.251) circle[radius=0.52bp];
\fill[tg808080] (200.031,354.328) circle[radius=0.52bp];
\fill[tg808080] (201.271,353.583) circle[radius=0.52bp];
\fill[tg808080] (202.510,344.315) circle[radius=0.52bp];
\fill[tg808080] (203.750,345.594) circle[radius=0.52bp];
\fill[tg808080] (204.990,353.741) circle[radius=0.52bp];
\fill[tg808080] (206.229,341.001) circle[radius=0.52bp];
\fill[tg808080] (207.469,357.880) circle[radius=0.52bp];
\fill[tg808080] (208.708,352.836) circle[radius=0.52bp];
\fill[tg808080] (209.948,367.048) circle[radius=0.52bp];
\fill[tg808080] (211.188,348.421) circle[radius=0.52bp];
\fill[tg808080] (212.427,351.849) circle[radius=0.52bp];
\fill[tg808080] (213.667,356.519) circle[radius=0.52bp];
\fill[tg808080] (214.906,355.447) circle[radius=0.52bp];
\fill[tg808080] (216.146,357.473) circle[radius=0.52bp];
\fill[tg808080] (217.385,367.839) circle[radius=0.52bp];
\fill[tg808080] (218.625,345.939) circle[radius=0.52bp];
\fill[tg808080] (219.865,360.778) circle[radius=0.52bp];
\fill[tg808080] (221.104,356.234) circle[radius=0.52bp];
\fill[tg808080] (222.344,349.161) circle[radius=0.52bp];
\fill[tg808080] (223.583,350.072) circle[radius=0.52bp];
\fill[tg808080] (224.823,360.377) circle[radius=0.52bp];
\fill[tg808080] (226.062,364.115) circle[radius=0.52bp];
\fill[tg808080] (227.302,341.550) circle[radius=0.52bp];
\fill[tg808080] (228.542,364.129) circle[radius=0.52bp];
\fill[tg808080] (229.781,363.988) circle[radius=0.52bp];
\fill[tg808080] (231.021,355.741) circle[radius=0.52bp];
\fill[tg808080] (232.260,347.041) circle[radius=0.52bp];
\fill[tg808080] (233.500,349.683) circle[radius=0.52bp];
\fill[tg808080] (234.740,338.041) circle[radius=0.52bp];
\fill[tg808080] (235.979,366.051) circle[radius=0.52bp];
\fill[tg808080] (237.219,344.885) circle[radius=0.52bp];
\fill[tg808080] (238.458,354.585) circle[radius=0.52bp];
\fill[tg808080] (239.698,355.683) circle[radius=0.52bp];
\fill[tg808080] (240.938,353.734) circle[radius=0.52bp];
\fill[tg808080] (242.177,357.893) circle[radius=0.52bp];
\fill[tg808080] (243.417,352.472) circle[radius=0.52bp];
\fill[tg808080] (244.656,353.623) circle[radius=0.52bp];
\fill[tg808080] (245.896,354.713) circle[radius=0.52bp];
\fill[tg808080] (247.135,352.120) circle[radius=0.52bp];
\fill[tg808080] (248.375,356.962) circle[radius=0.52bp];
\fill[tg808080] (249.615,353.815) circle[radius=0.52bp];
\fill[tg808080] (250.854,353.587) circle[radius=0.52bp];
\fill[tg808080] (252.094,353.084) circle[radius=0.52bp];
\fill[tg808080] (253.333,354.326) circle[radius=0.52bp];
\fill[tg808080] (254.573,355.779) circle[radius=0.52bp];
\fill[tg808080] (255.812,355.575) circle[radius=0.52bp];
\fill[tg808080] (257.052,356.109) circle[radius=0.52bp];
\fill[tg808080] (258.292,353.759) circle[radius=0.52bp];
\fill[tg808080] (259.531,351.208) circle[radius=0.52bp];
\fill[tg808080] (260.771,354.144) circle[radius=0.52bp];
\fill[tg808080] (262.010,352.226) circle[radius=0.52bp];
\fill[tg808080] (263.250,355.088) circle[radius=0.52bp];
\fill[tg808080] (264.490,353.705) circle[radius=0.52bp];
\fill[tg808080] (265.729,354.022) circle[radius=0.52bp];
\fill[tg808080] (266.969,353.626) circle[radius=0.52bp];
\fill[tg808080] (268.208,354.845) circle[radius=0.52bp];
\fill[tg808080] (269.448,355.995) circle[radius=0.52bp];
\fill[tg808080] (270.688,352.622) circle[radius=0.52bp];
\fill[tg808080] (271.927,350.783) circle[radius=0.52bp];
\fill[tg808080] (273.167,354.011) circle[radius=0.52bp];
\fill[tg808080] (274.406,354.103) circle[radius=0.52bp];
\fill[tg808080] (275.646,351.680) circle[radius=0.52bp];
\fill[tg808080] (276.885,353.304) circle[radius=0.52bp];
\fill[tg808080] (278.125,355.830) circle[radius=0.52bp];
\fill[tg808080] (279.365,357.913) circle[radius=0.52bp];
\fill[tg808080] (280.604,348.076) circle[radius=0.52bp];
\fill[tg808080] (281.844,353.715) circle[radius=0.52bp];
\fill[tg808080] (283.083,350.823) circle[radius=0.52bp];
\fill[tg808080] (284.323,333.161) circle[radius=0.52bp];
\fill[tg808080] (285.562,348.281) circle[radius=0.52bp];
\fill[tg808080] (286.802,353.393) circle[radius=0.52bp];
\fill[tg808080] (288.042,348.299) circle[radius=0.52bp];
\fill[tg808080] (289.281,352.231) circle[radius=0.52bp];
\fill[tg808080] (290.521,361.549) circle[radius=0.52bp];
\fill[tg808080] (291.760,355.286) circle[radius=0.52bp];
\draw[color=tg575757,line width=0.65bp,line cap=round] (263.250,353.647) -- (264.490,353.658) -- (265.729,353.668) -- (266.969,353.676) -- (268.208,353.684) -- (269.448,353.691) -- (270.688,353.697) -- (271.927,353.702) -- (273.167,353.707) -- (274.406,353.711) -- (275.646,353.715) -- (276.885,353.718) -- (278.125,353.721) -- (279.365,353.724) -- (280.604,353.726) -- (281.844,353.728) -- (283.083,353.730) -- (284.323,353.732) -- (285.562,353.734) -- (286.802,353.735) -- (288.042,353.736) -- (289.281,353.737) -- (290.521,353.738) -- (291.760,353.739);
\draw[color=tg0000ff,line width=0.38bp,line cap=round] (262.750,351.726) -- (263.750,352.726);
\draw[color=tg0000ff,line width=0.38bp,line cap=round] (262.750,352.726) -- (263.750,351.726);
\draw[color=tg0000ff,line width=0.38bp,line cap=round] (263.990,351.726) -- (264.990,352.726);
\draw[color=tg0000ff,line width=0.38bp,line cap=round] (263.990,352.726) -- (264.990,351.726);
\draw[color=tg0000ff,line width=0.38bp,line cap=round] (265.229,351.726) -- (266.229,352.726);
\draw[color=tg0000ff,line width=0.38bp,line cap=round] (265.229,352.726) -- (266.229,351.726);
\draw[color=tg0000ff,line width=0.38bp,line cap=round] (266.469,351.726) -- (267.469,352.726);
\draw[color=tg0000ff,line width=0.38bp,line cap=round] (266.469,352.726) -- (267.469,351.726);
\draw[color=tg0000ff,line width=0.38bp,line cap=round] (267.708,351.726) -- (268.708,352.726);
\draw[color=tg0000ff,line width=0.38bp,line cap=round] (267.708,352.726) -- (268.708,351.726);
\draw[color=tg0000ff,line width=0.38bp,line cap=round] (268.948,351.726) -- (269.948,352.726);
\draw[color=tg0000ff,line width=0.38bp,line cap=round] (268.948,352.726) -- (269.948,351.726);
\draw[color=tg0000ff,line width=0.38bp,line cap=round] (270.188,351.726) -- (271.188,352.726);
\draw[color=tg0000ff,line width=0.38bp,line cap=round] (270.188,352.726) -- (271.188,351.726);
\draw[color=tg0000ff,line width=0.38bp,line cap=round] (271.427,351.726) -- (272.427,352.726);
\draw[color=tg0000ff,line width=0.38bp,line cap=round] (271.427,352.726) -- (272.427,351.726);
\draw[color=tg0000ff,line width=0.38bp,line cap=round] (272.667,351.726) -- (273.667,352.726);
\draw[color=tg0000ff,line width=0.38bp,line cap=round] (272.667,352.726) -- (273.667,351.726);
\draw[color=tg0000ff,line width=0.38bp,line cap=round] (273.906,351.726) -- (274.906,352.726);
\draw[color=tg0000ff,line width=0.38bp,line cap=round] (273.906,352.726) -- (274.906,351.726);
\draw[color=tg0000ff,line width=0.38bp,line cap=round] (275.146,351.726) -- (276.146,352.726);
\draw[color=tg0000ff,line width=0.38bp,line cap=round] (275.146,352.726) -- (276.146,351.726);
\draw[color=tg0000ff,line width=0.38bp,line cap=round] (276.385,351.726) -- (277.385,352.726);
\draw[color=tg0000ff,line width=0.38bp,line cap=round] (276.385,352.726) -- (277.385,351.726);
\draw[color=tg0000ff,line width=0.38bp,line cap=round] (277.625,351.726) -- (278.625,352.726);
\draw[color=tg0000ff,line width=0.38bp,line cap=round] (277.625,352.726) -- (278.625,351.726);
\draw[color=tg0000ff,line width=0.38bp,line cap=round] (278.865,351.726) -- (279.865,352.726);
\draw[color=tg0000ff,line width=0.38bp,line cap=round] (278.865,352.726) -- (279.865,351.726);
\draw[color=tg0000ff,line width=0.38bp,line cap=round] (280.104,351.726) -- (281.104,352.726);
\draw[color=tg0000ff,line width=0.38bp,line cap=round] (280.104,352.726) -- (281.104,351.726);
\draw[color=tg0000ff,line width=0.38bp,line cap=round] (281.344,351.726) -- (282.344,352.726);
\draw[color=tg0000ff,line width=0.38bp,line cap=round] (281.344,352.726) -- (282.344,351.726);
\draw[color=tg0000ff,line width=0.38bp,line cap=round] (282.583,351.726) -- (283.583,352.726);
\draw[color=tg0000ff,line width=0.38bp,line cap=round] (282.583,352.726) -- (283.583,351.726);
\draw[color=tg0000ff,line width=0.38bp,line cap=round] (283.823,351.726) -- (284.823,352.726);
\draw[color=tg0000ff,line width=0.38bp,line cap=round] (283.823,352.726) -- (284.823,351.726);
\draw[color=tg0000ff,line width=0.38bp,line cap=round] (285.062,351.726) -- (286.062,352.726);
\draw[color=tg0000ff,line width=0.38bp,line cap=round] (285.062,352.726) -- (286.062,351.726);
\draw[color=tg0000ff,line width=0.38bp,line cap=round] (286.302,351.726) -- (287.302,352.726);
\draw[color=tg0000ff,line width=0.38bp,line cap=round] (286.302,352.726) -- (287.302,351.726);
\draw[color=tg0000ff,line width=0.38bp,line cap=round] (287.542,351.726) -- (288.542,352.726);
\draw[color=tg0000ff,line width=0.38bp,line cap=round] (287.542,352.726) -- (288.542,351.726);
\draw[color=tg0000ff,line width=0.38bp,line cap=round] (288.781,351.726) -- (289.781,352.726);
\draw[color=tg0000ff,line width=0.38bp,line cap=round] (288.781,352.726) -- (289.781,351.726);
\draw[color=tg0000ff,line width=0.38bp,line cap=round] (290.021,351.726) -- (291.021,352.726);
\draw[color=tg0000ff,line width=0.38bp,line cap=round] (290.021,352.726) -- (291.021,351.726);
\draw[color=tg0000ff,line width=0.38bp,line cap=round] (291.260,351.726) -- (292.260,352.726);
\draw[color=tg0000ff,line width=0.38bp,line cap=round] (291.260,352.726) -- (292.260,351.726);
\draw[color=tgc9cdd1,line width=0.45bp,line cap=round] (320.000,329.000) -- (320.000,372.000);
\draw[color=tgc9cdd1,line width=0.45bp,line cap=round] (320.000,372.000) -- (439.000,372.000);
\node[anchor=east,inner sep=0pt,outer sep=0pt,text=tg626b73,font={\sffamily\fontsize{6.7bp}{8.04bp}\selectfont}] at (317.000,365.980) {-0.019};
\node[anchor=east,inner sep=0pt,outer sep=0pt,text=tg626b73,font={\sffamily\fontsize{6.7bp}{8.04bp}\selectfont}] at (317.000,335.020) {0.012};
\draw[color=tg9da7b0,line width=0.6bp,line cap=round,dash pattern=on 2bp off 2bp] (409.250,329.000) -- (409.250,372.000);
\fill[tg808080] (320.000,345.008) circle[radius=0.52bp];
\fill[tg808080] (321.240,342.794) circle[radius=0.52bp];
\fill[tg808080] (322.479,342.701) circle[radius=0.52bp];
\fill[tg808080] (323.719,342.645) circle[radius=0.52bp];
\fill[tg808080] (324.958,350.698) circle[radius=0.52bp];
\fill[tg808080] (326.198,352.568) circle[radius=0.52bp];
\fill[tg808080] (327.438,353.095) circle[radius=0.52bp];
\fill[tg808080] (328.677,337.128) circle[radius=0.52bp];
\fill[tg808080] (329.917,353.150) circle[radius=0.52bp];
\fill[tg808080] (331.156,355.253) circle[radius=0.52bp];
\fill[tg808080] (332.396,349.351) circle[radius=0.52bp];
\fill[tg808080] (333.635,355.396) circle[radius=0.52bp];
\fill[tg808080] (334.875,352.887) circle[radius=0.52bp];
\fill[tg808080] (336.115,350.258) circle[radius=0.52bp];
\fill[tg808080] (337.354,342.554) circle[radius=0.52bp];
\fill[tg808080] (338.594,342.228) circle[radius=0.52bp];
\fill[tg808080] (339.833,342.986) circle[radius=0.52bp];
\fill[tg808080] (341.073,350.346) circle[radius=0.52bp];
\fill[tg808080] (342.312,342.918) circle[radius=0.52bp];
\fill[tg808080] (343.552,346.244) circle[radius=0.52bp];
\fill[tg808080] (344.792,340.998) circle[radius=0.52bp];
\fill[tg808080] (346.031,351.267) circle[radius=0.52bp];
\fill[tg808080] (347.271,340.147) circle[radius=0.52bp];
\fill[tg808080] (348.510,340.806) circle[radius=0.52bp];
\fill[tg808080] (349.750,343.164) circle[radius=0.52bp];
\fill[tg808080] (350.990,340.335) circle[radius=0.52bp];
\fill[tg808080] (352.229,341.620) circle[radius=0.52bp];
\fill[tg808080] (353.469,342.586) circle[radius=0.52bp];
\fill[tg808080] (354.708,340.262) circle[radius=0.52bp];
\fill[tg808080] (355.948,352.772) circle[radius=0.52bp];
\fill[tg808080] (357.188,344.871) circle[radius=0.52bp];
\fill[tg808080] (358.427,344.221) circle[radius=0.52bp];
\fill[tg808080] (359.667,337.449) circle[radius=0.52bp];
\fill[tg808080] (360.906,345.035) circle[radius=0.52bp];
\fill[tg808080] (362.146,363.479) circle[radius=0.52bp];
\fill[tg808080] (363.385,349.463) circle[radius=0.52bp];
\fill[tg808080] (364.625,347.172) circle[radius=0.52bp];
\fill[tg808080] (365.865,357.180) circle[radius=0.52bp];
\fill[tg808080] (367.104,346.129) circle[radius=0.52bp];
\fill[tg808080] (368.344,342.105) circle[radius=0.52bp];
\fill[tg808080] (369.583,334.782) circle[radius=0.52bp];
\fill[tg808080] (370.823,364.069) circle[radius=0.52bp];
\fill[tg808080] (372.062,354.362) circle[radius=0.52bp];
\fill[tg808080] (373.302,348.448) circle[radius=0.52bp];
\fill[tg808080] (374.542,340.091) circle[radius=0.52bp];
\fill[tg808080] (375.781,348.853) circle[radius=0.52bp];
\fill[tg808080] (377.021,340.788) circle[radius=0.52bp];
\fill[tg808080] (378.260,367.839) circle[radius=0.52bp];
\fill[tg808080] (379.500,364.590) circle[radius=0.52bp];
\fill[tg808080] (380.740,343.280) circle[radius=0.52bp];
\fill[tg808080] (381.979,347.942) circle[radius=0.52bp];
\fill[tg808080] (383.219,352.510) circle[radius=0.52bp];
\fill[tg808080] (384.458,338.540) circle[radius=0.52bp];
\fill[tg808080] (385.698,339.166) circle[radius=0.52bp];
\fill[tg808080] (386.938,343.274) circle[radius=0.52bp];
\fill[tg808080] (388.177,344.911) circle[radius=0.52bp];
\fill[tg808080] (389.417,333.161) circle[radius=0.52bp];
\fill[tg808080] (390.656,336.585) circle[radius=0.52bp];
\fill[tg808080] (391.896,348.110) circle[radius=0.52bp];
\fill[tg808080] (393.135,352.097) circle[radius=0.52bp];
\fill[tg808080] (394.375,349.869) circle[radius=0.52bp];
\fill[tg808080] (395.615,347.527) circle[radius=0.52bp];
\fill[tg808080] (396.854,357.222) circle[radius=0.52bp];
\fill[tg808080] (398.094,336.941) circle[radius=0.52bp];
\fill[tg808080] (399.333,348.220) circle[radius=0.52bp];
\fill[tg808080] (400.573,349.091) circle[radius=0.52bp];
\fill[tg808080] (401.812,347.785) circle[radius=0.52bp];
\fill[tg808080] (403.052,345.703) circle[radius=0.52bp];
\fill[tg808080] (404.292,342.225) circle[radius=0.52bp];
\fill[tg808080] (405.531,362.975) circle[radius=0.52bp];
\fill[tg808080] (406.771,346.655) circle[radius=0.52bp];
\fill[tg808080] (408.010,361.972) circle[radius=0.52bp];
\fill[tg808080] (409.250,338.557) circle[radius=0.52bp];
\fill[tg808080] (410.490,340.886) circle[radius=0.52bp];
\fill[tg808080] (411.729,357.740) circle[radius=0.52bp];
\fill[tg808080] (412.969,344.266) circle[radius=0.52bp];
\fill[tg808080] (414.208,346.138) circle[radius=0.52bp];
\fill[tg808080] (415.448,351.413) circle[radius=0.52bp];
\fill[tg808080] (416.688,354.108) circle[radius=0.52bp];
\fill[tg808080] (417.927,344.976) circle[radius=0.52bp];
\fill[tg808080] (419.167,345.009) circle[radius=0.52bp];
\fill[tg808080] (420.406,345.483) circle[radius=0.52bp];
\fill[tg808080] (421.646,356.068) circle[radius=0.52bp];
\fill[tg808080] (422.885,343.738) circle[radius=0.52bp];
\fill[tg808080] (424.125,352.776) circle[radius=0.52bp];
\fill[tg808080] (425.365,357.080) circle[radius=0.52bp];
\fill[tg808080] (426.604,345.864) circle[radius=0.52bp];
\fill[tg808080] (427.844,351.716) circle[radius=0.52bp];
\fill[tg808080] (429.083,346.419) circle[radius=0.52bp];
\fill[tg808080] (430.323,360.086) circle[radius=0.52bp];
\fill[tg808080] (431.562,345.182) circle[radius=0.52bp];
\fill[tg808080] (432.802,343.434) circle[radius=0.52bp];
\fill[tg808080] (434.042,348.809) circle[radius=0.52bp];
\fill[tg808080] (435.281,351.309) circle[radius=0.52bp];
\fill[tg808080] (436.521,357.706) circle[radius=0.52bp];
\fill[tg808080] (437.760,352.959) circle[radius=0.52bp];
\draw[color=tg575757,line width=0.65bp,line cap=round] (409.250,347.030) -- (410.490,347.050) -- (411.729,347.069) -- (412.969,347.085) -- (414.208,347.100) -- (415.448,347.113) -- (416.688,347.124) -- (417.927,347.134) -- (419.167,347.144) -- (420.406,347.152) -- (421.646,347.159) -- (422.885,347.165) -- (424.125,347.171) -- (425.365,347.176) -- (426.604,347.181) -- (427.844,347.185) -- (429.083,347.188) -- (430.323,347.192) -- (431.562,347.194) -- (432.802,347.197) -- (434.042,347.199) -- (435.281,347.201) -- (436.521,347.203) -- (437.760,347.205);
\draw[color=tg0000ff,line width=0.38bp,line cap=round] (408.750,361.472) -- (409.750,362.472);
\draw[color=tg0000ff,line width=0.38bp,line cap=round] (408.750,362.472) -- (409.750,361.472);
\draw[color=tg0000ff,line width=0.38bp,line cap=round] (409.990,361.472) -- (410.990,362.472);
\draw[color=tg0000ff,line width=0.38bp,line cap=round] (409.990,362.472) -- (410.990,361.472);
\draw[color=tg0000ff,line width=0.38bp,line cap=round] (411.229,361.472) -- (412.229,362.472);
\draw[color=tg0000ff,line width=0.38bp,line cap=round] (411.229,362.472) -- (412.229,361.472);
\draw[color=tg0000ff,line width=0.38bp,line cap=round] (412.469,361.472) -- (413.469,362.472);
\draw[color=tg0000ff,line width=0.38bp,line cap=round] (412.469,362.472) -- (413.469,361.472);
\draw[color=tg0000ff,line width=0.38bp,line cap=round] (413.708,361.472) -- (414.708,362.472);
\draw[color=tg0000ff,line width=0.38bp,line cap=round] (413.708,362.472) -- (414.708,361.472);
\draw[color=tg0000ff,line width=0.38bp,line cap=round] (414.948,361.472) -- (415.948,362.472);
\draw[color=tg0000ff,line width=0.38bp,line cap=round] (414.948,362.472) -- (415.948,361.472);
\draw[color=tg0000ff,line width=0.38bp,line cap=round] (416.188,361.472) -- (417.188,362.472);
\draw[color=tg0000ff,line width=0.38bp,line cap=round] (416.188,362.472) -- (417.188,361.472);
\draw[color=tg0000ff,line width=0.38bp,line cap=round] (417.427,361.472) -- (418.427,362.472);
\draw[color=tg0000ff,line width=0.38bp,line cap=round] (417.427,362.472) -- (418.427,361.472);
\draw[color=tg0000ff,line width=0.38bp,line cap=round] (418.667,361.472) -- (419.667,362.472);
\draw[color=tg0000ff,line width=0.38bp,line cap=round] (418.667,362.472) -- (419.667,361.472);
\draw[color=tg0000ff,line width=0.38bp,line cap=round] (419.906,361.472) -- (420.906,362.472);
\draw[color=tg0000ff,line width=0.38bp,line cap=round] (419.906,362.472) -- (420.906,361.472);
\draw[color=tg0000ff,line width=0.38bp,line cap=round] (421.146,361.472) -- (422.146,362.472);
\draw[color=tg0000ff,line width=0.38bp,line cap=round] (421.146,362.472) -- (422.146,361.472);
\draw[color=tg0000ff,line width=0.38bp,line cap=round] (422.385,361.472) -- (423.385,362.472);
\draw[color=tg0000ff,line width=0.38bp,line cap=round] (422.385,362.472) -- (423.385,361.472);
\draw[color=tg0000ff,line width=0.38bp,line cap=round] (423.625,361.472) -- (424.625,362.472);
\draw[color=tg0000ff,line width=0.38bp,line cap=round] (423.625,362.472) -- (424.625,361.472);
\draw[color=tg0000ff,line width=0.38bp,line cap=round] (424.865,361.472) -- (425.865,362.472);
\draw[color=tg0000ff,line width=0.38bp,line cap=round] (424.865,362.472) -- (425.865,361.472);
\draw[color=tg0000ff,line width=0.38bp,line cap=round] (426.104,361.472) -- (427.104,362.472);
\draw[color=tg0000ff,line width=0.38bp,line cap=round] (426.104,362.472) -- (427.104,361.472);
\draw[color=tg0000ff,line width=0.38bp,line cap=round] (427.344,361.472) -- (428.344,362.472);
\draw[color=tg0000ff,line width=0.38bp,line cap=round] (427.344,362.472) -- (428.344,361.472);
\draw[color=tg0000ff,line width=0.38bp,line cap=round] (428.583,361.472) -- (429.583,362.472);
\draw[color=tg0000ff,line width=0.38bp,line cap=round] (428.583,362.472) -- (429.583,361.472);
\draw[color=tg0000ff,line width=0.38bp,line cap=round] (429.823,361.472) -- (430.823,362.472);
\draw[color=tg0000ff,line width=0.38bp,line cap=round] (429.823,362.472) -- (430.823,361.472);
\draw[color=tg0000ff,line width=0.38bp,line cap=round] (431.062,361.472) -- (432.062,362.472);
\draw[color=tg0000ff,line width=0.38bp,line cap=round] (431.062,362.472) -- (432.062,361.472);
\draw[color=tg0000ff,line width=0.38bp,line cap=round] (432.302,361.472) -- (433.302,362.472);
\draw[color=tg0000ff,line width=0.38bp,line cap=round] (432.302,362.472) -- (433.302,361.472);
\draw[color=tg0000ff,line width=0.38bp,line cap=round] (433.542,361.472) -- (434.542,362.472);
\draw[color=tg0000ff,line width=0.38bp,line cap=round] (433.542,362.472) -- (434.542,361.472);
\draw[color=tg0000ff,line width=0.38bp,line cap=round] (434.781,361.472) -- (435.781,362.472);
\draw[color=tg0000ff,line width=0.38bp,line cap=round] (434.781,362.472) -- (435.781,361.472);
\draw[color=tg0000ff,line width=0.38bp,line cap=round] (436.021,361.472) -- (437.021,362.472);
\draw[color=tg0000ff,line width=0.38bp,line cap=round] (436.021,362.472) -- (437.021,361.472);
\draw[color=tg0000ff,line width=0.38bp,line cap=round] (437.260,361.472) -- (438.260,362.472);
\draw[color=tg0000ff,line width=0.38bp,line cap=round] (437.260,362.472) -- (438.260,361.472);
\node[anchor=west,inner sep=0pt,outer sep=0pt,text=tg18314b,font={\sffamily\fontsize{8bp}{9.6bp}\selectfont\bfseries}] at (0.000,387.000) {Hidden regimes + jumps (price)};
\node[anchor=east,inner sep=0pt,outer sep=0pt,text=tg626b73,font={\sffamily\fontsize{7bp}{8.4bp}\selectfont}] at (448.000,387.000) {Approximate filtered reference};
\draw[color=tgc9cdd1,line width=0.45bp,line cap=round] (28.000,396.000) -- (28.000,439.000);
\draw[color=tgc9cdd1,line width=0.45bp,line cap=round] (28.000,439.000) -- (147.000,439.000);
\node[anchor=east,inner sep=0pt,outer sep=0pt,text=tg626b73,font={\sffamily\fontsize{6.7bp}{8.04bp}\selectfont}] at (25.000,432.980) {74};
\node[anchor=east,inner sep=0pt,outer sep=0pt,text=tg626b73,font={\sffamily\fontsize{6.7bp}{8.04bp}\selectfont}] at (25.000,402.020) {117};
\draw[color=tg9da7b0,line width=0.6bp,line cap=round,dash pattern=on 2bp off 2bp] (117.250,396.000) -- (117.250,439.000);
\fill[tg808080] (28.000,400.826) circle[radius=0.52bp];
\fill[tg808080] (29.240,401.341) circle[radius=0.52bp];
\fill[tg808080] (30.479,402.183) circle[radius=0.52bp];
\fill[tg808080] (31.719,401.682) circle[radius=0.52bp];
\fill[tg808080] (32.958,401.910) circle[radius=0.52bp];
\fill[tg808080] (34.198,402.241) circle[radius=0.52bp];
\fill[tg808080] (35.438,401.994) circle[radius=0.52bp];
\fill[tg808080] (36.677,402.239) circle[radius=0.52bp];
\fill[tg808080] (37.917,402.800) circle[radius=0.52bp];
\fill[tg808080] (39.156,402.198) circle[radius=0.52bp];
\fill[tg808080] (40.396,402.040) circle[radius=0.52bp];
\fill[tg808080] (41.635,402.137) circle[radius=0.52bp];
\fill[tg808080] (42.875,401.697) circle[radius=0.52bp];
\fill[tg808080] (44.115,401.969) circle[radius=0.52bp];
\fill[tg808080] (45.354,402.051) circle[radius=0.52bp];
\fill[tg808080] (46.594,401.990) circle[radius=0.52bp];
\fill[tg808080] (47.833,400.783) circle[radius=0.52bp];
\fill[tg808080] (49.073,400.913) circle[radius=0.52bp];
\fill[tg808080] (50.312,400.216) circle[radius=0.52bp];
\fill[tg808080] (51.552,400.161) circle[radius=0.52bp];
\fill[tg808080] (52.792,400.515) circle[radius=0.52bp];
\fill[tg808080] (54.031,401.026) circle[radius=0.52bp];
\fill[tg808080] (55.271,401.178) circle[radius=0.52bp];
\fill[tg808080] (56.510,401.433) circle[radius=0.52bp];
\fill[tg808080] (57.750,402.195) circle[radius=0.52bp];
\fill[tg808080] (58.990,402.121) circle[radius=0.52bp];
\fill[tg808080] (60.229,401.302) circle[radius=0.52bp];
\fill[tg808080] (61.469,401.015) circle[radius=0.52bp];
\fill[tg808080] (62.708,401.579) circle[radius=0.52bp];
\fill[tg808080] (63.948,401.066) circle[radius=0.52bp];
\fill[tg808080] (65.188,401.466) circle[radius=0.52bp];
\fill[tg808080] (66.427,402.257) circle[radius=0.52bp];
\fill[tg808080] (67.667,402.285) circle[radius=0.52bp];
\fill[tg808080] (68.906,402.654) circle[radius=0.52bp];
\fill[tg808080] (70.146,403.076) circle[radius=0.52bp];
\fill[tg808080] (71.385,402.881) circle[radius=0.52bp];
\fill[tg808080] (72.625,402.588) circle[radius=0.52bp];
\fill[tg808080] (73.865,402.547) circle[radius=0.52bp];
\fill[tg808080] (75.104,403.009) circle[radius=0.52bp];
\fill[tg808080] (76.344,403.157) circle[radius=0.52bp];
\fill[tg808080] (77.583,402.936) circle[radius=0.52bp];
\fill[tg808080] (78.823,402.380) circle[radius=0.52bp];
\fill[tg808080] (80.062,402.133) circle[radius=0.52bp];
\fill[tg808080] (81.302,401.586) circle[radius=0.52bp];
\fill[tg808080] (82.542,401.354) circle[radius=0.52bp];
\fill[tg808080] (83.781,400.931) circle[radius=0.52bp];
\fill[tg808080] (85.021,401.254) circle[radius=0.52bp];
\fill[tg808080] (86.260,401.209) circle[radius=0.52bp];
\fill[tg808080] (87.500,401.954) circle[radius=0.52bp];
\fill[tg808080] (88.740,401.885) circle[radius=0.52bp];
\fill[tg808080] (89.979,401.863) circle[radius=0.52bp];
\fill[tg808080] (91.219,402.300) circle[radius=0.52bp];
\fill[tg808080] (92.458,402.489) circle[radius=0.52bp];
\fill[tg808080] (93.698,403.133) circle[radius=0.52bp];
\fill[tg808080] (94.938,403.079) circle[radius=0.52bp];
\fill[tg808080] (96.177,402.679) circle[radius=0.52bp];
\fill[tg808080] (97.417,402.375) circle[radius=0.52bp];
\fill[tg808080] (98.656,401.710) circle[radius=0.52bp];
\fill[tg808080] (99.896,401.640) circle[radius=0.52bp];
\fill[tg808080] (101.135,406.489) circle[radius=0.52bp];
\fill[tg808080] (102.375,406.475) circle[radius=0.52bp];
\fill[tg808080] (103.615,406.179) circle[radius=0.52bp];
\fill[tg808080] (104.854,407.158) circle[radius=0.52bp];
\fill[tg808080] (106.094,407.350) circle[radius=0.52bp];
\fill[tg808080] (107.333,407.310) circle[radius=0.52bp];
\fill[tg808080] (108.573,407.246) circle[radius=0.52bp];
\fill[tg808080] (109.812,407.135) circle[radius=0.52bp];
\fill[tg808080] (111.052,406.634) circle[radius=0.52bp];
\fill[tg808080] (112.292,406.539) circle[radius=0.52bp];
\fill[tg808080] (113.531,406.211) circle[radius=0.52bp];
\fill[tg808080] (114.771,406.493) circle[radius=0.52bp];
\fill[tg808080] (116.010,410.613) circle[radius=0.52bp];
\fill[tg808080] (117.250,411.198) circle[radius=0.52bp];
\fill[tg808080] (118.490,409.961) circle[radius=0.52bp];
\fill[tg808080] (119.729,409.510) circle[radius=0.52bp];
\fill[tg808080] (120.969,411.535) circle[radius=0.52bp];
\fill[tg808080] (122.208,409.824) circle[radius=0.52bp];
\fill[tg808080] (123.448,411.652) circle[radius=0.52bp];
\fill[tg808080] (124.688,417.638) circle[radius=0.52bp];
\fill[tg808080] (125.927,432.565) circle[radius=0.52bp];
\fill[tg808080] (127.167,433.336) circle[radius=0.52bp];
\fill[tg808080] (128.406,432.161) circle[radius=0.52bp];
\fill[tg808080] (129.646,434.839) circle[radius=0.52bp];
\fill[tg808080] (130.885,434.434) circle[radius=0.52bp];
\fill[tg808080] (132.125,434.728) circle[radius=0.52bp];
\fill[tg808080] (133.365,433.577) circle[radius=0.52bp];
\fill[tg808080] (134.604,433.703) circle[radius=0.52bp];
\fill[tg808080] (135.844,433.669) circle[radius=0.52bp];
\fill[tg808080] (137.083,433.962) circle[radius=0.52bp];
\fill[tg808080] (138.323,431.782) circle[radius=0.52bp];
\fill[tg808080] (139.562,431.418) circle[radius=0.52bp];
\fill[tg808080] (140.802,429.760) circle[radius=0.52bp];
\fill[tg808080] (142.042,429.061) circle[radius=0.52bp];
\fill[tg808080] (143.281,428.404) circle[radius=0.52bp];
\fill[tg808080] (144.521,428.153) circle[radius=0.52bp];
\fill[tg808080] (145.760,429.463) circle[radius=0.52bp];
\draw[color=tg575757,line width=0.65bp,line cap=round] (117.250,410.632) -- (118.490,410.649) -- (119.729,410.665) -- (120.969,410.678) -- (122.208,410.690) -- (123.448,410.701) -- (124.688,410.710) -- (125.927,410.718) -- (127.167,410.725) -- (128.406,410.732) -- (129.646,410.737) -- (130.885,410.742) -- (132.125,410.747) -- (133.365,410.750) -- (134.604,410.754) -- (135.844,410.757) -- (137.083,410.759) -- (138.323,410.761) -- (139.562,410.763) -- (140.802,410.765) -- (142.042,410.766) -- (143.281,410.767) -- (144.521,410.768) -- (145.760,410.769);
\draw[color=tg0000ff,line width=0.38bp,line cap=round] (116.750,410.113) -- (117.750,411.113);
\draw[color=tg0000ff,line width=0.38bp,line cap=round] (116.750,411.113) -- (117.750,410.113);
\draw[color=tg0000ff,line width=0.38bp,line cap=round] (117.990,410.113) -- (118.990,411.113);
\draw[color=tg0000ff,line width=0.38bp,line cap=round] (117.990,411.113) -- (118.990,410.113);
\draw[color=tg0000ff,line width=0.38bp,line cap=round] (119.229,410.113) -- (120.229,411.113);
\draw[color=tg0000ff,line width=0.38bp,line cap=round] (119.229,411.113) -- (120.229,410.113);
\draw[color=tg0000ff,line width=0.38bp,line cap=round] (120.469,410.113) -- (121.469,411.113);
\draw[color=tg0000ff,line width=0.38bp,line cap=round] (120.469,411.113) -- (121.469,410.113);
\draw[color=tg0000ff,line width=0.38bp,line cap=round] (121.708,410.113) -- (122.708,411.113);
\draw[color=tg0000ff,line width=0.38bp,line cap=round] (121.708,411.113) -- (122.708,410.113);
\draw[color=tg0000ff,line width=0.38bp,line cap=round] (122.948,410.113) -- (123.948,411.113);
\draw[color=tg0000ff,line width=0.38bp,line cap=round] (122.948,411.113) -- (123.948,410.113);
\draw[color=tg0000ff,line width=0.38bp,line cap=round] (124.188,410.113) -- (125.188,411.113);
\draw[color=tg0000ff,line width=0.38bp,line cap=round] (124.188,411.113) -- (125.188,410.113);
\draw[color=tg0000ff,line width=0.38bp,line cap=round] (125.427,410.113) -- (126.427,411.113);
\draw[color=tg0000ff,line width=0.38bp,line cap=round] (125.427,411.113) -- (126.427,410.113);
\draw[color=tg0000ff,line width=0.38bp,line cap=round] (126.667,410.113) -- (127.667,411.113);
\draw[color=tg0000ff,line width=0.38bp,line cap=round] (126.667,411.113) -- (127.667,410.113);
\draw[color=tg0000ff,line width=0.38bp,line cap=round] (127.906,410.113) -- (128.906,411.113);
\draw[color=tg0000ff,line width=0.38bp,line cap=round] (127.906,411.113) -- (128.906,410.113);
\draw[color=tg0000ff,line width=0.38bp,line cap=round] (129.146,410.113) -- (130.146,411.113);
\draw[color=tg0000ff,line width=0.38bp,line cap=round] (129.146,411.113) -- (130.146,410.113);
\draw[color=tg0000ff,line width=0.38bp,line cap=round] (130.385,410.113) -- (131.385,411.113);
\draw[color=tg0000ff,line width=0.38bp,line cap=round] (130.385,411.113) -- (131.385,410.113);
\draw[color=tg0000ff,line width=0.38bp,line cap=round] (131.625,410.113) -- (132.625,411.113);
\draw[color=tg0000ff,line width=0.38bp,line cap=round] (131.625,411.113) -- (132.625,410.113);
\draw[color=tg0000ff,line width=0.38bp,line cap=round] (132.865,410.113) -- (133.865,411.113);
\draw[color=tg0000ff,line width=0.38bp,line cap=round] (132.865,411.113) -- (133.865,410.113);
\draw[color=tg0000ff,line width=0.38bp,line cap=round] (134.104,410.113) -- (135.104,411.113);
\draw[color=tg0000ff,line width=0.38bp,line cap=round] (134.104,411.113) -- (135.104,410.113);
\draw[color=tg0000ff,line width=0.38bp,line cap=round] (135.344,410.113) -- (136.344,411.113);
\draw[color=tg0000ff,line width=0.38bp,line cap=round] (135.344,411.113) -- (136.344,410.113);
\draw[color=tg0000ff,line width=0.38bp,line cap=round] (136.583,410.113) -- (137.583,411.113);
\draw[color=tg0000ff,line width=0.38bp,line cap=round] (136.583,411.113) -- (137.583,410.113);
\draw[color=tg0000ff,line width=0.38bp,line cap=round] (137.823,410.113) -- (138.823,411.113);
\draw[color=tg0000ff,line width=0.38bp,line cap=round] (137.823,411.113) -- (138.823,410.113);
\draw[color=tg0000ff,line width=0.38bp,line cap=round] (139.062,410.113) -- (140.062,411.113);
\draw[color=tg0000ff,line width=0.38bp,line cap=round] (139.062,411.113) -- (140.062,410.113);
\draw[color=tg0000ff,line width=0.38bp,line cap=round] (140.302,410.113) -- (141.302,411.113);
\draw[color=tg0000ff,line width=0.38bp,line cap=round] (140.302,411.113) -- (141.302,410.113);
\draw[color=tg0000ff,line width=0.38bp,line cap=round] (141.542,410.113) -- (142.542,411.113);
\draw[color=tg0000ff,line width=0.38bp,line cap=round] (141.542,411.113) -- (142.542,410.113);
\draw[color=tg0000ff,line width=0.38bp,line cap=round] (142.781,410.113) -- (143.781,411.113);
\draw[color=tg0000ff,line width=0.38bp,line cap=round] (142.781,411.113) -- (143.781,410.113);
\draw[color=tg0000ff,line width=0.38bp,line cap=round] (144.021,410.113) -- (145.021,411.113);
\draw[color=tg0000ff,line width=0.38bp,line cap=round] (144.021,411.113) -- (145.021,410.113);
\draw[color=tg0000ff,line width=0.38bp,line cap=round] (145.260,410.113) -- (146.260,411.113);
\draw[color=tg0000ff,line width=0.38bp,line cap=round] (145.260,411.113) -- (146.260,410.113);
\node[anchor=center,inner sep=0pt,outer sep=0pt,text=tg626b73,font={\sffamily\fontsize{6.8bp}{8.16bp}\selectfont}] at (28.000,448.000) {-72};
\node[anchor=center,inner sep=0pt,outer sep=0pt,text=tg626b73,font={\sffamily\fontsize{6.8bp}{8.16bp}\selectfont}] at (117.250,448.000) {0};
\node[anchor=center,inner sep=0pt,outer sep=0pt,text=tg626b73,font={\sffamily\fontsize{6.8bp}{8.16bp}\selectfont}] at (147.000,448.000) {+24};
\draw[color=tgc9cdd1,line width=0.45bp,line cap=round] (174.000,396.000) -- (174.000,439.000);
\draw[color=tgc9cdd1,line width=0.45bp,line cap=round] (174.000,439.000) -- (293.000,439.000);
\node[anchor=east,inner sep=0pt,outer sep=0pt,text=tg626b73,font={\sffamily\fontsize{6.7bp}{8.04bp}\selectfont}] at (171.000,432.980) {81};
\node[anchor=east,inner sep=0pt,outer sep=0pt,text=tg626b73,font={\sffamily\fontsize{6.7bp}{8.04bp}\selectfont}] at (171.000,402.020) {97};
\draw[color=tg9da7b0,line width=0.6bp,line cap=round,dash pattern=on 2bp off 2bp] (263.250,396.000) -- (263.250,439.000);
\fill[tg808080] (174.000,423.651) circle[radius=0.52bp];
\fill[tg808080] (175.240,425.379) circle[radius=0.52bp];
\fill[tg808080] (176.479,422.240) circle[radius=0.52bp];
\fill[tg808080] (177.719,423.082) circle[radius=0.52bp];
\fill[tg808080] (178.958,423.745) circle[radius=0.52bp];
\fill[tg808080] (180.198,423.539) circle[radius=0.52bp];
\fill[tg808080] (181.438,422.291) circle[radius=0.52bp];
\fill[tg808080] (182.677,420.128) circle[radius=0.52bp];
\fill[tg808080] (183.917,420.167) circle[radius=0.52bp];
\fill[tg808080] (185.156,418.569) circle[radius=0.52bp];
\fill[tg808080] (186.396,417.634) circle[radius=0.52bp];
\fill[tg808080] (187.635,417.371) circle[radius=0.52bp];
\fill[tg808080] (188.875,417.125) circle[radius=0.52bp];
\fill[tg808080] (190.115,417.990) circle[radius=0.52bp];
\fill[tg808080] (191.354,418.526) circle[radius=0.52bp];
\fill[tg808080] (192.594,417.801) circle[radius=0.52bp];
\fill[tg808080] (193.833,417.456) circle[radius=0.52bp];
\fill[tg808080] (195.073,418.178) circle[radius=0.52bp];
\fill[tg808080] (196.312,420.189) circle[radius=0.52bp];
\fill[tg808080] (197.552,418.764) circle[radius=0.52bp];
\fill[tg808080] (198.792,420.808) circle[radius=0.52bp];
\fill[tg808080] (200.031,419.911) circle[radius=0.52bp];
\fill[tg808080] (201.271,420.330) circle[radius=0.52bp];
\fill[tg808080] (202.510,420.276) circle[radius=0.52bp];
\fill[tg808080] (203.750,414.248) circle[radius=0.52bp];
\fill[tg808080] (204.990,408.873) circle[radius=0.52bp];
\fill[tg808080] (206.229,408.923) circle[radius=0.52bp];
\fill[tg808080] (207.469,400.161) circle[radius=0.52bp];
\fill[tg808080] (208.708,403.121) circle[radius=0.52bp];
\fill[tg808080] (209.948,402.541) circle[radius=0.52bp];
\fill[tg808080] (211.188,411.669) circle[radius=0.52bp];
\fill[tg808080] (212.427,408.144) circle[radius=0.52bp];
\fill[tg808080] (213.667,406.906) circle[radius=0.52bp];
\fill[tg808080] (214.906,408.844) circle[radius=0.52bp];
\fill[tg808080] (216.146,410.043) circle[radius=0.52bp];
\fill[tg808080] (217.385,412.580) circle[radius=0.52bp];
\fill[tg808080] (218.625,421.720) circle[radius=0.52bp];
\fill[tg808080] (219.865,416.795) circle[radius=0.52bp];
\fill[tg808080] (221.104,421.331) circle[radius=0.52bp];
\fill[tg808080] (222.344,422.939) circle[radius=0.52bp];
\fill[tg808080] (223.583,420.105) circle[radius=0.52bp];
\fill[tg808080] (224.823,417.810) circle[radius=0.52bp];
\fill[tg808080] (226.062,422.068) circle[radius=0.52bp];
\fill[tg808080] (227.302,428.491) circle[radius=0.52bp];
\fill[tg808080] (228.542,421.018) circle[radius=0.52bp];
\fill[tg808080] (229.781,427.490) circle[radius=0.52bp];
\fill[tg808080] (231.021,433.633) circle[radius=0.52bp];
\fill[tg808080] (232.260,434.839) circle[radius=0.52bp];
\fill[tg808080] (233.500,430.951) circle[radius=0.52bp];
\fill[tg808080] (234.740,428.568) circle[radius=0.52bp];
\fill[tg808080] (235.979,418.872) circle[radius=0.52bp];
\fill[tg808080] (237.219,426.602) circle[radius=0.52bp];
\fill[tg808080] (238.458,421.156) circle[radius=0.52bp];
\fill[tg808080] (239.698,421.733) circle[radius=0.52bp];
\fill[tg808080] (240.938,422.995) circle[radius=0.52bp];
\fill[tg808080] (242.177,423.036) circle[radius=0.52bp];
\fill[tg808080] (243.417,425.649) circle[radius=0.52bp];
\fill[tg808080] (244.656,424.914) circle[radius=0.52bp];
\fill[tg808080] (245.896,424.887) circle[radius=0.52bp];
\fill[tg808080] (247.135,425.530) circle[radius=0.52bp];
\fill[tg808080] (248.375,424.577) circle[radius=0.52bp];
\fill[tg808080] (249.615,426.599) circle[radius=0.52bp];
\fill[tg808080] (250.854,426.689) circle[radius=0.52bp];
\fill[tg808080] (252.094,426.640) circle[radius=0.52bp];
\fill[tg808080] (253.333,426.284) circle[radius=0.52bp];
\fill[tg808080] (254.573,426.686) circle[radius=0.52bp];
\fill[tg808080] (255.812,427.968) circle[radius=0.52bp];
\fill[tg808080] (257.052,429.118) circle[radius=0.52bp];
\fill[tg808080] (258.292,430.578) circle[radius=0.52bp];
\fill[tg808080] (259.531,430.632) circle[radius=0.52bp];
\fill[tg808080] (260.771,429.162) circle[radius=0.52bp];
\fill[tg808080] (262.010,429.448) circle[radius=0.52bp];
\fill[tg808080] (263.250,428.582) circle[radius=0.52bp];
\fill[tg808080] (264.490,429.436) circle[radius=0.52bp];
\fill[tg808080] (265.729,429.458) circle[radius=0.52bp];
\fill[tg808080] (266.969,429.670) circle[radius=0.52bp];
\fill[tg808080] (268.208,429.645) circle[radius=0.52bp];
\fill[tg808080] (269.448,430.349) circle[radius=0.52bp];
\fill[tg808080] (270.688,431.731) circle[radius=0.52bp];
\fill[tg808080] (271.927,431.112) circle[radius=0.52bp];
\fill[tg808080] (273.167,429.391) circle[radius=0.52bp];
\fill[tg808080] (274.406,429.597) circle[radius=0.52bp];
\fill[tg808080] (275.646,429.858) circle[radius=0.52bp];
\fill[tg808080] (276.885,428.665) circle[radius=0.52bp];
\fill[tg808080] (278.125,428.446) circle[radius=0.52bp];
\fill[tg808080] (279.365,429.746) circle[radius=0.52bp];
\fill[tg808080] (280.604,432.266) circle[radius=0.52bp];
\fill[tg808080] (281.844,428.938) circle[radius=0.52bp];
\fill[tg808080] (283.083,428.967) circle[radius=0.52bp];
\fill[tg808080] (284.323,427.248) circle[radius=0.52bp];
\fill[tg808080] (285.562,414.302) circle[radius=0.52bp];
\fill[tg808080] (286.802,410.735) circle[radius=0.52bp];
\fill[tg808080] (288.042,410.552) circle[radius=0.52bp];
\fill[tg808080] (289.281,406.921) circle[radius=0.52bp];
\fill[tg808080] (290.521,405.937) circle[radius=0.52bp];
\fill[tg808080] (291.760,411.271) circle[radius=0.52bp];
\draw[color=tg575757,line width=0.65bp,line cap=round] (263.250,429.426) -- (264.490,429.406) -- (265.729,429.388) -- (266.969,429.372) -- (268.208,429.358) -- (269.448,429.345) -- (270.688,429.333) -- (271.927,429.322) -- (273.167,429.313) -- (274.406,429.304) -- (275.646,429.296) -- (276.885,429.289) -- (278.125,429.283) -- (279.365,429.277) -- (280.604,429.272) -- (281.844,429.267) -- (283.083,429.262) -- (284.323,429.258) -- (285.562,429.254) -- (286.802,429.251) -- (288.042,429.247) -- (289.281,429.244) -- (290.521,429.242) -- (291.760,429.239);
\draw[color=tg0000ff,line width=0.38bp,line cap=round] (262.750,428.948) -- (263.750,429.948);
\draw[color=tg0000ff,line width=0.38bp,line cap=round] (262.750,429.948) -- (263.750,428.948);
\draw[color=tg0000ff,line width=0.38bp,line cap=round] (263.990,428.948) -- (264.990,429.948);
\draw[color=tg0000ff,line width=0.38bp,line cap=round] (263.990,429.948) -- (264.990,428.948);
\draw[color=tg0000ff,line width=0.38bp,line cap=round] (265.229,428.948) -- (266.229,429.948);
\draw[color=tg0000ff,line width=0.38bp,line cap=round] (265.229,429.948) -- (266.229,428.948);
\draw[color=tg0000ff,line width=0.38bp,line cap=round] (266.469,428.948) -- (267.469,429.948);
\draw[color=tg0000ff,line width=0.38bp,line cap=round] (266.469,429.948) -- (267.469,428.948);
\draw[color=tg0000ff,line width=0.38bp,line cap=round] (267.708,428.948) -- (268.708,429.948);
\draw[color=tg0000ff,line width=0.38bp,line cap=round] (267.708,429.948) -- (268.708,428.948);
\draw[color=tg0000ff,line width=0.38bp,line cap=round] (268.948,428.948) -- (269.948,429.948);
\draw[color=tg0000ff,line width=0.38bp,line cap=round] (268.948,429.948) -- (269.948,428.948);
\draw[color=tg0000ff,line width=0.38bp,line cap=round] (270.188,428.948) -- (271.188,429.948);
\draw[color=tg0000ff,line width=0.38bp,line cap=round] (270.188,429.948) -- (271.188,428.948);
\draw[color=tg0000ff,line width=0.38bp,line cap=round] (271.427,428.948) -- (272.427,429.948);
\draw[color=tg0000ff,line width=0.38bp,line cap=round] (271.427,429.948) -- (272.427,428.948);
\draw[color=tg0000ff,line width=0.38bp,line cap=round] (272.667,428.948) -- (273.667,429.948);
\draw[color=tg0000ff,line width=0.38bp,line cap=round] (272.667,429.948) -- (273.667,428.948);
\draw[color=tg0000ff,line width=0.38bp,line cap=round] (273.906,428.948) -- (274.906,429.948);
\draw[color=tg0000ff,line width=0.38bp,line cap=round] (273.906,429.948) -- (274.906,428.948);
\draw[color=tg0000ff,line width=0.38bp,line cap=round] (275.146,428.948) -- (276.146,429.948);
\draw[color=tg0000ff,line width=0.38bp,line cap=round] (275.146,429.948) -- (276.146,428.948);
\draw[color=tg0000ff,line width=0.38bp,line cap=round] (276.385,428.948) -- (277.385,429.948);
\draw[color=tg0000ff,line width=0.38bp,line cap=round] (276.385,429.948) -- (277.385,428.948);
\draw[color=tg0000ff,line width=0.38bp,line cap=round] (277.625,428.948) -- (278.625,429.948);
\draw[color=tg0000ff,line width=0.38bp,line cap=round] (277.625,429.948) -- (278.625,428.948);
\draw[color=tg0000ff,line width=0.38bp,line cap=round] (278.865,428.948) -- (279.865,429.948);
\draw[color=tg0000ff,line width=0.38bp,line cap=round] (278.865,429.948) -- (279.865,428.948);
\draw[color=tg0000ff,line width=0.38bp,line cap=round] (280.104,428.948) -- (281.104,429.948);
\draw[color=tg0000ff,line width=0.38bp,line cap=round] (280.104,429.948) -- (281.104,428.948);
\draw[color=tg0000ff,line width=0.38bp,line cap=round] (281.344,428.948) -- (282.344,429.948);
\draw[color=tg0000ff,line width=0.38bp,line cap=round] (281.344,429.948) -- (282.344,428.948);
\draw[color=tg0000ff,line width=0.38bp,line cap=round] (282.583,428.948) -- (283.583,429.948);
\draw[color=tg0000ff,line width=0.38bp,line cap=round] (282.583,429.948) -- (283.583,428.948);
\draw[color=tg0000ff,line width=0.38bp,line cap=round] (283.823,428.948) -- (284.823,429.948);
\draw[color=tg0000ff,line width=0.38bp,line cap=round] (283.823,429.948) -- (284.823,428.948);
\draw[color=tg0000ff,line width=0.38bp,line cap=round] (285.062,428.948) -- (286.062,429.948);
\draw[color=tg0000ff,line width=0.38bp,line cap=round] (285.062,429.948) -- (286.062,428.948);
\draw[color=tg0000ff,line width=0.38bp,line cap=round] (286.302,428.948) -- (287.302,429.948);
\draw[color=tg0000ff,line width=0.38bp,line cap=round] (286.302,429.948) -- (287.302,428.948);
\draw[color=tg0000ff,line width=0.38bp,line cap=round] (287.542,428.948) -- (288.542,429.948);
\draw[color=tg0000ff,line width=0.38bp,line cap=round] (287.542,429.948) -- (288.542,428.948);
\draw[color=tg0000ff,line width=0.38bp,line cap=round] (288.781,428.948) -- (289.781,429.948);
\draw[color=tg0000ff,line width=0.38bp,line cap=round] (288.781,429.948) -- (289.781,428.948);
\draw[color=tg0000ff,line width=0.38bp,line cap=round] (290.021,428.948) -- (291.021,429.948);
\draw[color=tg0000ff,line width=0.38bp,line cap=round] (290.021,429.948) -- (291.021,428.948);
\draw[color=tg0000ff,line width=0.38bp,line cap=round] (291.260,428.948) -- (292.260,429.948);
\draw[color=tg0000ff,line width=0.38bp,line cap=round] (291.260,429.948) -- (292.260,428.948);
\node[anchor=center,inner sep=0pt,outer sep=0pt,text=tg626b73,font={\sffamily\fontsize{6.8bp}{8.16bp}\selectfont}] at (174.000,448.000) {-72};
\node[anchor=center,inner sep=0pt,outer sep=0pt,text=tg626b73,font={\sffamily\fontsize{6.8bp}{8.16bp}\selectfont}] at (263.250,448.000) {0};
\node[anchor=center,inner sep=0pt,outer sep=0pt,text=tg626b73,font={\sffamily\fontsize{6.8bp}{8.16bp}\selectfont}] at (293.000,448.000) {+24};
\draw[color=tgc9cdd1,line width=0.45bp,line cap=round] (320.000,396.000) -- (320.000,439.000);
\draw[color=tgc9cdd1,line width=0.45bp,line cap=round] (320.000,439.000) -- (439.000,439.000);
\node[anchor=east,inner sep=0pt,outer sep=0pt,text=tg626b73,font={\sffamily\fontsize{6.7bp}{8.04bp}\selectfont}] at (317.000,432.980) {34.2};
\node[anchor=east,inner sep=0pt,outer sep=0pt,text=tg626b73,font={\sffamily\fontsize{6.7bp}{8.04bp}\selectfont}] at (317.000,402.020) {37.8};
\draw[color=tg9da7b0,line width=0.6bp,line cap=round,dash pattern=on 2bp off 2bp] (409.250,396.000) -- (409.250,439.000);
\fill[tg808080] (320.000,413.051) circle[radius=0.52bp];
\fill[tg808080] (321.240,412.441) circle[radius=0.52bp];
\fill[tg808080] (322.479,411.126) circle[radius=0.52bp];
\fill[tg808080] (323.719,409.775) circle[radius=0.52bp];
\fill[tg808080] (324.958,408.401) circle[radius=0.52bp];
\fill[tg808080] (326.198,409.610) circle[radius=0.52bp];
\fill[tg808080] (327.438,411.410) circle[radius=0.52bp];
\fill[tg808080] (328.677,413.367) circle[radius=0.52bp];
\fill[tg808080] (329.917,410.250) circle[radius=0.52bp];
\fill[tg808080] (331.156,412.232) circle[radius=0.52bp];
\fill[tg808080] (332.396,414.864) circle[radius=0.52bp];
\fill[tg808080] (333.635,415.626) circle[radius=0.52bp];
\fill[tg808080] (334.875,418.274) circle[radius=0.52bp];
\fill[tg808080] (336.115,420.125) circle[radius=0.52bp];
\fill[tg808080] (337.354,421.153) circle[radius=0.52bp];
\fill[tg808080] (338.594,419.800) circle[radius=0.52bp];
\fill[tg808080] (339.833,418.339) circle[radius=0.52bp];
\fill[tg808080] (341.073,417.108) circle[radius=0.52bp];
\fill[tg808080] (342.312,418.175) circle[radius=0.52bp];
\fill[tg808080] (343.552,416.922) circle[radius=0.52bp];
\fill[tg808080] (344.792,416.708) circle[radius=0.52bp];
\fill[tg808080] (346.031,414.845) circle[radius=0.52bp];
\fill[tg808080] (347.271,416.208) circle[radius=0.52bp];
\fill[tg808080] (348.510,414.074) circle[radius=0.52bp];
\fill[tg808080] (349.750,412.134) circle[radius=0.52bp];
\fill[tg808080] (350.990,410.935) circle[radius=0.52bp];
\fill[tg808080] (352.229,408.825) circle[radius=0.52bp];
\fill[tg808080] (353.469,407.116) circle[radius=0.52bp];
\fill[tg808080] (354.708,405.710) circle[radius=0.52bp];
\fill[tg808080] (355.948,403.541) circle[radius=0.52bp];
\fill[tg808080] (357.188,405.442) circle[radius=0.52bp];
\fill[tg808080] (358.427,404.774) circle[radius=0.52bp];
\fill[tg808080] (359.667,403.892) circle[radius=0.52bp];
\fill[tg808080] (360.906,400.786) circle[radius=0.52bp];
\fill[tg808080] (362.146,400.161) circle[radius=0.52bp];
\fill[tg808080] (363.385,405.574) circle[radius=0.52bp];
\fill[tg808080] (364.625,406.394) circle[radius=0.52bp];
\fill[tg808080] (365.865,406.473) circle[radius=0.52bp];
\fill[tg808080] (367.104,409.771) circle[radius=0.52bp];
\fill[tg808080] (368.344,409.514) circle[radius=0.52bp];
\fill[tg808080] (369.583,407.965) circle[radius=0.52bp];
\fill[tg808080] (370.823,404.030) circle[radius=0.52bp];
\fill[tg808080] (372.062,409.567) circle[radius=0.52bp];
\fill[tg808080] (373.302,411.939) circle[radius=0.52bp];
\fill[tg808080] (374.542,412.421) circle[radius=0.52bp];
\fill[tg808080] (375.781,410.243) circle[radius=0.52bp];
\fill[tg808080] (377.021,410.857) circle[radius=0.52bp];
\fill[tg808080] (378.260,408.892) circle[radius=0.52bp];
\fill[tg808080] (379.500,415.533) circle[radius=0.52bp];
\fill[tg808080] (380.740,421.034) circle[radius=0.52bp];
\fill[tg808080] (381.979,419.905) circle[radius=0.52bp];
\fill[tg808080] (383.219,420.219) circle[radius=0.52bp];
\fill[tg808080] (384.458,421.941) circle[radius=0.52bp];
\fill[tg808080] (385.698,419.348) circle[radius=0.52bp];
\fill[tg808080] (386.938,416.928) circle[radius=0.52bp];
\fill[tg808080] (388.177,415.782) circle[radius=0.52bp];
\fill[tg808080] (389.417,415.147) circle[radius=0.52bp];
\fill[tg808080] (390.656,410.784) circle[radius=0.52bp];
\fill[tg808080] (391.896,407.466) circle[radius=0.52bp];
\fill[tg808080] (393.135,407.846) circle[radius=0.52bp];
\fill[tg808080] (394.375,409.506) circle[radius=0.52bp];
\fill[tg808080] (395.615,410.446) circle[radius=0.52bp];
\fill[tg808080] (396.854,410.637) circle[radius=0.52bp];
\fill[tg808080] (398.094,413.906) circle[radius=0.52bp];
\fill[tg808080] (399.333,410.734) circle[radius=0.52bp];
\fill[tg808080] (400.573,411.145) circle[radius=0.52bp];
\fill[tg808080] (401.812,411.834) circle[radius=0.52bp];
\fill[tg808080] (403.052,412.106) circle[radius=0.52bp];
\fill[tg808080] (404.292,411.715) circle[radius=0.52bp];
\fill[tg808080] (405.531,410.215) circle[radius=0.52bp];
\fill[tg808080] (406.771,415.303) circle[radius=0.52bp];
\fill[tg808080] (408.010,415.216) circle[radius=0.52bp];
\fill[tg808080] (409.250,419.913) circle[radius=0.52bp];
\fill[tg808080] (410.490,417.308) circle[radius=0.52bp];
\fill[tg808080] (411.729,415.413) circle[radius=0.52bp];
\fill[tg808080] (412.969,418.794) circle[radius=0.52bp];
\fill[tg808080] (414.208,417.964) circle[radius=0.52bp];
\fill[tg808080] (415.448,417.717) circle[radius=0.52bp];
\fill[tg808080] (416.688,419.113) circle[radius=0.52bp];
\fill[tg808080] (417.927,421.335) circle[radius=0.52bp];
\fill[tg808080] (419.167,420.732) circle[radius=0.52bp];
\fill[tg808080] (420.406,420.137) circle[radius=0.52bp];
\fill[tg808080] (421.646,419.689) circle[radius=0.52bp];
\fill[tg808080] (422.885,422.510) circle[radius=0.52bp];
\fill[tg808080] (424.125,421.527) circle[radius=0.52bp];
\fill[tg808080] (425.365,423.324) circle[radius=0.52bp];
\fill[tg808080] (426.604,426.418) circle[radius=0.52bp];
\fill[tg808080] (427.844,426.095) circle[radius=0.52bp];
\fill[tg808080] (429.083,427.545) circle[radius=0.52bp];
\fill[tg808080] (430.323,427.390) circle[radius=0.52bp];
\fill[tg808080] (431.562,431.341) circle[radius=0.52bp];
\fill[tg808080] (432.802,430.820) circle[radius=0.52bp];
\fill[tg808080] (434.042,429.773) circle[radius=0.52bp];
\fill[tg808080] (435.281,430.336) circle[radius=0.52bp];
\fill[tg808080] (436.521,431.644) circle[radius=0.52bp];
\fill[tg808080] (437.760,434.839) circle[radius=0.52bp];
\draw[color=tg575757,line width=0.65bp,line cap=round] (409.250,415.176) -- (410.490,415.140) -- (411.729,415.108) -- (412.969,415.079) -- (414.208,415.053) -- (415.448,415.030) -- (416.688,415.009) -- (417.927,414.990) -- (419.167,414.972) -- (420.406,414.957) -- (421.646,414.943) -- (422.885,414.930) -- (424.125,414.918) -- (425.365,414.907) -- (426.604,414.898) -- (427.844,414.889) -- (429.083,414.881) -- (430.323,414.873) -- (431.562,414.866) -- (432.802,414.859) -- (434.042,414.853) -- (435.281,414.848) -- (436.521,414.842) -- (437.760,414.837);
\draw[color=tg0000ff,line width=0.38bp,line cap=round] (408.750,414.716) -- (409.750,415.716);
\draw[color=tg0000ff,line width=0.38bp,line cap=round] (408.750,415.716) -- (409.750,414.716);
\draw[color=tg0000ff,line width=0.38bp,line cap=round] (409.990,414.716) -- (410.990,415.716);
\draw[color=tg0000ff,line width=0.38bp,line cap=round] (409.990,415.716) -- (410.990,414.716);
\draw[color=tg0000ff,line width=0.38bp,line cap=round] (411.229,414.716) -- (412.229,415.716);
\draw[color=tg0000ff,line width=0.38bp,line cap=round] (411.229,415.716) -- (412.229,414.716);
\draw[color=tg0000ff,line width=0.38bp,line cap=round] (412.469,414.716) -- (413.469,415.716);
\draw[color=tg0000ff,line width=0.38bp,line cap=round] (412.469,415.716) -- (413.469,414.716);
\draw[color=tg0000ff,line width=0.38bp,line cap=round] (413.708,414.716) -- (414.708,415.716);
\draw[color=tg0000ff,line width=0.38bp,line cap=round] (413.708,415.716) -- (414.708,414.716);
\draw[color=tg0000ff,line width=0.38bp,line cap=round] (414.948,414.716) -- (415.948,415.716);
\draw[color=tg0000ff,line width=0.38bp,line cap=round] (414.948,415.716) -- (415.948,414.716);
\draw[color=tg0000ff,line width=0.38bp,line cap=round] (416.188,414.716) -- (417.188,415.716);
\draw[color=tg0000ff,line width=0.38bp,line cap=round] (416.188,415.716) -- (417.188,414.716);
\draw[color=tg0000ff,line width=0.38bp,line cap=round] (417.427,414.716) -- (418.427,415.716);
\draw[color=tg0000ff,line width=0.38bp,line cap=round] (417.427,415.716) -- (418.427,414.716);
\draw[color=tg0000ff,line width=0.38bp,line cap=round] (418.667,414.716) -- (419.667,415.716);
\draw[color=tg0000ff,line width=0.38bp,line cap=round] (418.667,415.716) -- (419.667,414.716);
\draw[color=tg0000ff,line width=0.38bp,line cap=round] (419.906,414.716) -- (420.906,415.716);
\draw[color=tg0000ff,line width=0.38bp,line cap=round] (419.906,415.716) -- (420.906,414.716);
\draw[color=tg0000ff,line width=0.38bp,line cap=round] (421.146,414.716) -- (422.146,415.716);
\draw[color=tg0000ff,line width=0.38bp,line cap=round] (421.146,415.716) -- (422.146,414.716);
\draw[color=tg0000ff,line width=0.38bp,line cap=round] (422.385,414.716) -- (423.385,415.716);
\draw[color=tg0000ff,line width=0.38bp,line cap=round] (422.385,415.716) -- (423.385,414.716);
\draw[color=tg0000ff,line width=0.38bp,line cap=round] (423.625,414.716) -- (424.625,415.716);
\draw[color=tg0000ff,line width=0.38bp,line cap=round] (423.625,415.716) -- (424.625,414.716);
\draw[color=tg0000ff,line width=0.38bp,line cap=round] (424.865,414.716) -- (425.865,415.716);
\draw[color=tg0000ff,line width=0.38bp,line cap=round] (424.865,415.716) -- (425.865,414.716);
\draw[color=tg0000ff,line width=0.38bp,line cap=round] (426.104,414.716) -- (427.104,415.716);
\draw[color=tg0000ff,line width=0.38bp,line cap=round] (426.104,415.716) -- (427.104,414.716);
\draw[color=tg0000ff,line width=0.38bp,line cap=round] (427.344,414.716) -- (428.344,415.716);
\draw[color=tg0000ff,line width=0.38bp,line cap=round] (427.344,415.716) -- (428.344,414.716);
\draw[color=tg0000ff,line width=0.38bp,line cap=round] (428.583,414.716) -- (429.583,415.716);
\draw[color=tg0000ff,line width=0.38bp,line cap=round] (428.583,415.716) -- (429.583,414.716);
\draw[color=tg0000ff,line width=0.38bp,line cap=round] (429.823,414.716) -- (430.823,415.716);
\draw[color=tg0000ff,line width=0.38bp,line cap=round] (429.823,415.716) -- (430.823,414.716);
\draw[color=tg0000ff,line width=0.38bp,line cap=round] (431.062,414.716) -- (432.062,415.716);
\draw[color=tg0000ff,line width=0.38bp,line cap=round] (431.062,415.716) -- (432.062,414.716);
\draw[color=tg0000ff,line width=0.38bp,line cap=round] (432.302,414.716) -- (433.302,415.716);
\draw[color=tg0000ff,line width=0.38bp,line cap=round] (432.302,415.716) -- (433.302,414.716);
\draw[color=tg0000ff,line width=0.38bp,line cap=round] (433.542,414.716) -- (434.542,415.716);
\draw[color=tg0000ff,line width=0.38bp,line cap=round] (433.542,415.716) -- (434.542,414.716);
\draw[color=tg0000ff,line width=0.38bp,line cap=round] (434.781,414.716) -- (435.781,415.716);
\draw[color=tg0000ff,line width=0.38bp,line cap=round] (434.781,415.716) -- (435.781,414.716);
\draw[color=tg0000ff,line width=0.38bp,line cap=round] (436.021,414.716) -- (437.021,415.716);
\draw[color=tg0000ff,line width=0.38bp,line cap=round] (436.021,415.716) -- (437.021,414.716);
\draw[color=tg0000ff,line width=0.38bp,line cap=round] (437.260,414.716) -- (438.260,415.716);
\draw[color=tg0000ff,line width=0.38bp,line cap=round] (437.260,415.716) -- (438.260,414.716);
\node[anchor=center,inner sep=0pt,outer sep=0pt,text=tg626b73,font={\sffamily\fontsize{6.8bp}{8.16bp}\selectfont}] at (320.000,448.000) {-72};
\node[anchor=center,inner sep=0pt,outer sep=0pt,text=tg626b73,font={\sffamily\fontsize{6.8bp}{8.16bp}\selectfont}] at (409.250,448.000) {0};
\node[anchor=center,inner sep=0pt,outer sep=0pt,text=tg626b73,font={\sffamily\fontsize{6.8bp}{8.16bp}\selectfont}] at (439.000,448.000) {+24};
\node[anchor=center,inner sep=0pt,outer sep=0pt,text=tg18314b,font={\sffamily\fontsize{7.3bp}{8.76bp}\selectfont}] at (224.000,461.000) {Time relative to forecast origin  |  Each panel has its own value scale};
\node[anchor=center,inner sep=0pt,outer sep=0pt,text=tg18314b,font={\sffamily\fontsize{7.4bp}{8.88bp}\selectfont\bfseries}] at (224.000,477.000) {Known parameters; regime+jump reference means use approximate filtering.};
\end{tikzpicture}%
}
\endgroup

% END VECTOR FIGURE paper-financial-stresses
\caption{\textbf{Financial mechanisms, isolated and composed.} The same
last-value forecaster is fitted to six configurations, three independent
series each.
The first pair contrasts geometric Brownian motion with Merton jumps in log
prices. Subsequent rows show GARCH log returns, hidden return regimes,
regimes combined with jumps, and the combined process on the price scale.
Gray dots are observations and blue crosses are separate point forecasts;
thin curves show the answer key (the conditional mean). The last two rows use an
approximate numerically filtered regime posterior, visibly labeled. Each
model observes 240 values and predicts 24; only the final 72 context values
are displayed. Units and vertical scales vary by row and panel.}
\label{fig:financial-stresses}
\end{figure}


\subsection{Joint mechanisms and dependence-aware uncertainty}
\label{sec:joint-stresses}

The same four identifiers accept $D>1$ targets. Each sampled future is an
$H\times D$ array, and covariance uses horizon-major vectorization. A sample
index identifies one complete future; fitting each target separately and
independently drawing its uncertainty generally does not reproduce that law.
Scalar location and scale parameters broadcast over targets, while vectors
specify target-specific values. The implementation separates ordinary-shock
correlation, jump-size correlation, shared jump arrivals, and shared regimes.

For additive compositions, innovation covariance is
$\operatorname{diag}(\boldsymbol\sigma_t)R
\operatorname{diag}(\boldsymbol\sigma_t)$. An optional shared Poisson clock
adds vector marks $\boldsymbol J\sim\mathcal N(\boldsymbol a,B)$, giving
increment covariance
\begin{equation}
 Q_t=\operatorname{diag}(\boldsymbol\sigma_t)R
       \operatorname{diag}(\boldsymbol\sigma_t)
       +\lambda(B+\boldsymbol a\boldsymbol a^\top).
\end{equation}
Fresh residuals have block-diagonal time covariance. Accumulated residuals
sum $Q_t$ through the earlier of the two forecast horizons. Scheduled changes
may affect different targets differently; their future schedule remains
an evaluator information advantage.

Joint Merton returns use correlated diffusion and one shared Poisson arrival
process, with a separate correlation matrix for Gaussian log-jump marks.
The cumulant construction extends to vectors. For
$K(\boldsymbol u)=\log\E\exp(\boldsymbol u^\top\boldsymbol R)$ and coordinate
vectors $\boldsymbol e_i$, conditional price covariance is
\begin{equation}
 \operatorname{Cov}(S_{T+h,i},S_{T+k,j}\mid\mathcal I_T)
 =m_{h,i}m_{k,j}\left[
 \exp\!\left(\min(h,k)
 \{K(\boldsymbol e_i+\boldsymbol e_j)-K(\boldsymbol e_i)
 -K(\boldsymbol e_j)\}\right)-1\right].
\end{equation}
Here $K$ includes one observation step and $m_{h,i}$ is the conditional price
mean. Zero ordinary-shock correlation does not remove dependence from a
shared jump clock. The declared construction has common jump arrivals;
it does not add an unspecified independent jump process for each target.

The GARCH extension is constant-conditional-correlation GARCH
(CCC-GARCH)~\cite{bollerslev1990ccc}:
$\boldsymbol\epsilon_t=\operatorname{diag}(\sqrt{\boldsymbol v_t})
\boldsymbol z_t$, with $\boldsymbol z_t\sim\mathcal N(0,R)$ and one
GARCH(1,1) recursion per target. Conditional simulation updates all variance
states along each future trajectory. Means and marginal variances are
analytical; cross-time return covariances vanish. However, future same-time
cross-series covariance contains
\begin{equation}
 R_{ij}\,\E[\sqrt{v_{T+h,i}v_{T+h,j}}\mid\mathcal I_T],
\end{equation}
which is generally not $R_{ij}\sqrt{\E[v_{T+h,i}]\E[v_{T+h,j}]}$.
An independent Monte Carlo ensemble estimates the expected conditional
covariance, rescaled to the analytical diagonal to preserve both known
marginal variances and positive semidefiniteness. The sample count and seed
are recorded; the default count is 16,384. First-step, zero-correlation,
and fully deterministic-variance cases use analytical covariance.
Log-level covariance accumulates the return blocks.

These numerical moments are labeled
\code{analytic\_mean\_monte\_carlo\_covariance}; the conditional sampler
remains exact under the declared law. Covariance rank calibration holds the
estimated target matrix fixed, using independent forecast and null ensembles.
Conditional on that fixed target, the statistics remain exchangeable under
the null. This separates the validity of rank calibration from the precision
of the numerical covariance, which may affect diagnostic power. The check
still cannot detect every form of higher-order dependence.

Joint Markov scenarios share one hidden chain. Regime parameters may be
$K\times D$ arrays, and within-regime correlation may differ by state. The
forward filter uses the complete observed return vector, never the simulated
state. Analytical recursions propagate the filtered distribution and all
cross-target/cross-horizon moments. With shared jumps, a truncated multivariate
Gaussian mixture approximates each emission likelihood, so the covariance
rank check again returns \code{UNSUPPORTED}. Multivariate filtering
requires positive diffusion scales and positive-definite state correlation
matrices, while the other joint families also support singular correlations.

Figure~\ref{fig:joint-stresses} shows a targeted pairing fault: permuting one
target's entire sample trajectories preserves its marginal distributions and
within-target temporal paths, while changing dependence between targets.

\begin{figure}[p]
\centering
% BEGIN VECTOR FIGURE paper-joint-stresses
% Generated vector scene; SVG and PNG use the same plotted geometry.
\begingroup
\definecolor{tg0000ff}{HTML}{0000ff}
\definecolor{tg053061}{HTML}{053061}
\definecolor{tg08366a}{HTML}{08366a}
\definecolor{tg0c3d73}{HTML}{0c3d73}
\definecolor{tg0f437b}{HTML}{0f437b}
\definecolor{tg124984}{HTML}{124984}
\definecolor{tg175290}{HTML}{175290}
\definecolor{tg18314b}{HTML}{18314b}
\definecolor{tg1a5899}{HTML}{1a5899}
\definecolor{tg1d5fa2}{HTML}{1d5fa2}
\definecolor{tg2065ab}{HTML}{2065ab}
\definecolor{tg266caf}{HTML}{266caf}
\definecolor{tg2a71b2}{HTML}{2a71b2}
\definecolor{tg2e77b5}{HTML}{2e77b5}
\definecolor{tg327cb7}{HTML}{327cb7}
\definecolor{tg3783bb}{HTML}{3783bb}
\definecolor{tg3b88be}{HTML}{3b88be}
\definecolor{tg3f8ec0}{HTML}{3f8ec0}
\definecolor{tg4393c3}{HTML}{4393c3}
\definecolor{tg4f9bc7}{HTML}{4f9bc7}
\definecolor{tg59a1ca}{HTML}{59a1ca}
\definecolor{tg626b73}{HTML}{626b73}
\definecolor{tg62a7ce}{HTML}{62a7ce}
\definecolor{tg67001f}{HTML}{67001f}
\definecolor{tg6bacd1}{HTML}{6bacd1}
\definecolor{tg700320}{HTML}{700320}
\definecolor{tg78b4d5}{HTML}{78b4d5}
\definecolor{tg790622}{HTML}{790622}
\definecolor{tg7c0722}{HTML}{7c0722}
\definecolor{tg808080}{HTML}{808080}
\definecolor{tg810823}{HTML}{810823}
\definecolor{tg81bad8}{HTML}{81bad8}
\definecolor{tg8a0b25}{HTML}{8a0b25}
\definecolor{tg8ac0db}{HTML}{8ac0db}
\definecolor{tg930e26}{HTML}{930e26}
\definecolor{tg93c6de}{HTML}{93c6de}
\definecolor{tg960f27}{HTML}{960f27}
\definecolor{tg991027}{HTML}{991027}
\definecolor{tg9c1127}{HTML}{9c1127}
\definecolor{tg9dcbe1}{HTML}{9dcbe1}
\definecolor{tg9f1228}{HTML}{9f1228}
\definecolor{tga21328}{HTML}{a21328}
\definecolor{tga51429}{HTML}{a51429}
\definecolor{tga5cee3}{HTML}{a5cee3}
\definecolor{tga81529}{HTML}{a81529}
\definecolor{tgab162a}{HTML}{ab162a}
\definecolor{tgacd2e5}{HTML}{acd2e5}
\definecolor{tgb1182b}{HTML}{b1182b}
\definecolor{tgb3d6e8}{HTML}{b3d6e8}
\definecolor{tgb72230}{HTML}{b72230}
\definecolor{tgb82531}{HTML}{b82531}
\definecolor{tgba2832}{HTML}{ba2832}
\definecolor{tgbb2a34}{HTML}{bb2a34}
\definecolor{tgbbdaea}{HTML}{bbdaea}
\definecolor{tgbd2d35}{HTML}{bd2d35}
\definecolor{tgbf3338}{HTML}{bf3338}
\definecolor{tgc13639}{HTML}{c13639}
\definecolor{tgc2383a}{HTML}{c2383a}
\definecolor{tgc43b3c}{HTML}{c43b3c}
\definecolor{tgc53e3d}{HTML}{c53e3d}
\definecolor{tgc5dfec}{HTML}{c5dfec}
\definecolor{tgc6413e}{HTML}{c6413e}
\definecolor{tgc94741}{HTML}{c94741}
\definecolor{tgc9cdd1}{HTML}{c9cdd1}
\definecolor{tgcce2ef}{HTML}{cce2ef}
\definecolor{tgce4f45}{HTML}{ce4f45}
\definecolor{tgd05548}{HTML}{d05548}
\definecolor{tgd25849}{HTML}{d25849}
\definecolor{tgd2e6f0}{HTML}{d2e6f0}
\definecolor{tgd35a4a}{HTML}{d35a4a}
\definecolor{tgd55d4c}{HTML}{d55d4c}
\definecolor{tgd6604d}{HTML}{d6604d}
\definecolor{tgd7634f}{HTML}{d7634f}
\definecolor{tgd7e8f1}{HTML}{d7e8f1}
\definecolor{tgd86551}{HTML}{d86551}
\definecolor{tgda6853}{HTML}{da6853}
\definecolor{tgdb6b55}{HTML}{db6b55}
\definecolor{tgddebf2}{HTML}{ddebf2}
\definecolor{tgde735c}{HTML}{de735c}
\definecolor{tgdf765e}{HTML}{df765e}
\definecolor{tge17860}{HTML}{e17860}
\definecolor{tge1e4e7}{HTML}{e1e4e7}
\definecolor{tge1edf3}{HTML}{e1edf3}
\definecolor{tge27b62}{HTML}{e27b62}
\definecolor{tge48066}{HTML}{e48066}
\definecolor{tge58368}{HTML}{e58368}
\definecolor{tge6eff4}{HTML}{e6eff4}
\definecolor{tge98b6e}{HTML}{e98b6e}
\definecolor{tgea8e70}{HTML}{ea8e70}
\definecolor{tgeaf1f5}{HTML}{eaf1f5}
\definecolor{tgedf2f5}{HTML}{edf2f5}
\definecolor{tgee9677}{HTML}{ee9677}
\definecolor{tgef9979}{HTML}{ef9979}
\definecolor{tgf0f4f6}{HTML}{f0f4f6}
\definecolor{tgf19e7d}{HTML}{f19e7d}
\definecolor{tgf2a17f}{HTML}{f2a17f}
\definecolor{tgf2f5f6}{HTML}{f2f5f6}
\definecolor{tgf3a481}{HTML}{f3a481}
\definecolor{tgf3f5f6}{HTML}{f3f5f6}
\definecolor{tgf4a683}{HTML}{f4a683}
\definecolor{tgf5a886}{HTML}{f5a886}
\definecolor{tgf5ac8b}{HTML}{f5ac8b}
\definecolor{tgf5f6f7}{HTML}{f5f6f7}
\definecolor{tgf6af8e}{HTML}{f6af8e}
\definecolor{tgf6b394}{HTML}{f6b394}
\definecolor{tgf6f7f7}{HTML}{f6f7f7}
\definecolor{tgf7b596}{HTML}{f7b596}
\definecolor{tgf7b799}{HTML}{f7b799}
\definecolor{tgf7f5f4}{HTML}{f7f5f4}
\definecolor{tgf7f6f6}{HTML}{f7f6f6}
\definecolor{tgf8bb9e}{HTML}{f8bb9e}
\definecolor{tgf8f1ed}{HTML}{f8f1ed}
\definecolor{tgf8f2ef}{HTML}{f8f2ef}
\definecolor{tgf8f3f0}{HTML}{f8f3f0}
\definecolor{tgf8f4f2}{HTML}{f8f4f2}
\definecolor{tgf9c2a7}{HTML}{f9c2a7}
\definecolor{tgf9ebe3}{HTML}{f9ebe3}
\definecolor{tgf9eee7}{HTML}{f9eee7}
\definecolor{tgf9efe9}{HTML}{f9efe9}
\definecolor{tgf9f0eb}{HTML}{f9f0eb}
\definecolor{tgfac8af}{HTML}{fac8af}
\definecolor{tgfae7dc}{HTML}{fae7dc}
\definecolor{tgfaeae1}{HTML}{faeae1}
\definecolor{tgfbd0b9}{HTML}{fbd0b9}
\definecolor{tgfbe4d6}{HTML}{fbe4d6}
\definecolor{tgfcd7c2}{HTML}{fcd7c2}
\definecolor{tgfcdfcf}{HTML}{fcdfcf}
\definecolor{tgfddcc9}{HTML}{fddcc9}
\resizebox{\linewidth}{!}{%
\begin{tikzpicture}[x=1bp,y=-1bp]
\path[use as bounding box] (0,0) rectangle (448,412);
\node[anchor=west,inner sep=0pt,outer sep=0pt,text=tg18314b,font={\sffamily\fontsize{10.2bp}{12.24bp}\selectfont\bfseries}] at (0.000,11.000) {Same marginal uncertainty. Different joint futures.};
\node[anchor=west,inner sep=0pt,outer sep=0pt,text=tg18314b,font={\sffamily\fontsize{7.5bp}{9bp}\selectfont}] at (0.000,29.000) {Merton log prices | Whole target trajectories are re-paired across samples};
\node[anchor=center,inner sep=0pt,outer sep=0pt,text=tg18314b,font={\sffamily\fontsize{9bp}{10.8bp}\selectfont\bfseries}] at (117.000,48.000) {Paired futures};
\draw[color=tgc9cdd1,line width=0.5bp,line cap=round] (61.000,63.000) -- (61.000,175.000);
\draw[color=tgc9cdd1,line width=0.5bp,line cap=round] (61.000,175.000) -- (173.000,175.000);
\draw[color=tge1e4e7,line width=0.45bp,line cap=round] (117.000,63.000) -- (117.000,175.000);
\draw[color=tge1e4e7,line width=0.45bp,line cap=round] (61.000,119.000) -- (173.000,119.000);
\fill[tg808080] (107.864,120.688) circle[radius=0.6bp];
\fill[tg808080] (87.994,150.559) circle[radius=0.6bp];
\fill[tg808080] (116.241,118.093) circle[radius=0.6bp];
\fill[tg808080] (111.487,124.020) circle[radius=0.6bp];
\fill[tg808080] (97.312,125.851) circle[radius=0.6bp];
\fill[tg808080] (122.594,113.012) circle[radius=0.6bp];
\fill[tg808080] (109.453,122.760) circle[radius=0.6bp];
\fill[tg808080] (123.685,110.791) circle[radius=0.6bp];
\fill[tg808080] (117.025,120.677) circle[radius=0.6bp];
\fill[tg808080] (124.456,110.707) circle[radius=0.6bp];
\fill[tg808080] (101.965,129.814) circle[radius=0.6bp];
\fill[tg808080] (118.365,123.323) circle[radius=0.6bp];
\fill[tg808080] (105.526,131.059) circle[radius=0.6bp];
\fill[tg808080] (109.584,134.545) circle[radius=0.6bp];
\fill[tg808080] (105.196,123.283) circle[radius=0.6bp];
\fill[tg808080] (114.130,117.194) circle[radius=0.6bp];
\fill[tg808080] (116.941,114.324) circle[radius=0.6bp];
\fill[tg808080] (114.535,119.017) circle[radius=0.6bp];
\fill[tg808080] (124.735,105.728) circle[radius=0.6bp];
\fill[tg808080] (129.222,112.631) circle[radius=0.6bp];
\fill[tg808080] (117.202,118.259) circle[radius=0.6bp];
\fill[tg808080] (108.715,132.358) circle[radius=0.6bp];
\fill[tg808080] (130.933,106.608) circle[radius=0.6bp];
\fill[tg808080] (130.765,110.929) circle[radius=0.6bp];
\fill[tg808080] (110.051,123.954) circle[radius=0.6bp];
\fill[tg808080] (114.734,120.304) circle[radius=0.6bp];
\fill[tg808080] (108.417,112.125) circle[radius=0.6bp];
\fill[tg808080] (115.271,122.034) circle[radius=0.6bp];
\fill[tg808080] (119.220,115.999) circle[radius=0.6bp];
\fill[tg808080] (110.688,121.093) circle[radius=0.6bp];
\fill[tg808080] (107.345,124.131) circle[radius=0.6bp];
\fill[tg808080] (109.815,115.624) circle[radius=0.6bp];
\fill[tg808080] (122.422,115.010) circle[radius=0.6bp];
\fill[tg808080] (122.530,119.695) circle[radius=0.6bp];
\fill[tg808080] (136.031,108.671) circle[radius=0.6bp];
\fill[tg808080] (117.437,119.328) circle[radius=0.6bp];
\fill[tg808080] (123.898,113.981) circle[radius=0.6bp];
\fill[tg808080] (119.257,114.582) circle[radius=0.6bp];
\fill[tg808080] (129.526,112.904) circle[radius=0.6bp];
\fill[tg808080] (122.434,119.076) circle[radius=0.6bp];
\fill[tg808080] (126.384,115.268) circle[radius=0.6bp];
\fill[tg808080] (105.337,124.339) circle[radius=0.6bp];
\fill[tg808080] (115.470,125.408) circle[radius=0.6bp];
\fill[tg808080] (122.921,110.880) circle[radius=0.6bp];
\fill[tg808080] (116.482,123.899) circle[radius=0.6bp];
\fill[tg808080] (125.241,112.472) circle[radius=0.6bp];
\fill[tg808080] (139.160,97.738) circle[radius=0.6bp];
\fill[tg808080] (112.949,122.403) circle[radius=0.6bp];
\fill[tg808080] (115.054,110.890) circle[radius=0.6bp];
\fill[tg808080] (106.756,133.094) circle[radius=0.6bp];
\fill[tg808080] (115.380,122.075) circle[radius=0.6bp];
\fill[tg808080] (115.301,120.586) circle[radius=0.6bp];
\fill[tg808080] (112.546,126.072) circle[radius=0.6bp];
\fill[tg808080] (119.673,120.175) circle[radius=0.6bp];
\fill[tg808080] (138.932,96.560) circle[radius=0.6bp];
\fill[tg808080] (119.169,119.794) circle[radius=0.6bp];
\fill[tg808080] (114.005,111.874) circle[radius=0.6bp];
\fill[tg808080] (130.449,106.317) circle[radius=0.6bp];
\fill[tg808080] (129.113,100.262) circle[radius=0.6bp];
\fill[tg808080] (113.619,120.897) circle[radius=0.6bp];
\fill[tg808080] (120.638,118.733) circle[radius=0.6bp];
\fill[tg808080] (114.812,118.187) circle[radius=0.6bp];
\fill[tg808080] (118.440,113.057) circle[radius=0.6bp];
\fill[tg808080] (112.088,121.813) circle[radius=0.6bp];
\fill[tg808080] (105.714,127.893) circle[radius=0.6bp];
\fill[tg808080] (113.407,126.431) circle[radius=0.6bp];
\fill[tg808080] (118.845,114.335) circle[radius=0.6bp];
\fill[tg808080] (110.023,125.603) circle[radius=0.6bp];
\fill[tg808080] (121.112,116.404) circle[radius=0.6bp];
\fill[tg808080] (125.635,110.899) circle[radius=0.6bp];
\fill[tg808080] (123.216,117.980) circle[radius=0.6bp];
\fill[tg808080] (127.311,115.435) circle[radius=0.6bp];
\fill[tg808080] (122.266,112.373) circle[radius=0.6bp];
\fill[tg808080] (124.407,104.183) circle[radius=0.6bp];
\fill[tg808080] (109.912,131.229) circle[radius=0.6bp];
\fill[tg808080] (118.139,123.375) circle[radius=0.6bp];
\fill[tg808080] (116.627,116.395) circle[radius=0.6bp];
\fill[tg808080] (111.430,128.183) circle[radius=0.6bp];
\fill[tg808080] (122.382,114.551) circle[radius=0.6bp];
\fill[tg808080] (118.730,118.431) circle[radius=0.6bp];
\fill[tg808080] (109.151,126.625) circle[radius=0.6bp];
\fill[tg808080] (119.579,128.032) circle[radius=0.6bp];
\fill[tg808080] (108.847,123.749) circle[radius=0.6bp];
\fill[tg808080] (131.018,105.870) circle[radius=0.6bp];
\fill[tg808080] (118.195,118.531) circle[radius=0.6bp];
\fill[tg808080] (136.689,106.256) circle[radius=0.6bp];
\fill[tg808080] (104.622,123.077) circle[radius=0.6bp];
\fill[tg808080] (124.581,113.731) circle[radius=0.6bp];
\fill[tg808080] (117.622,122.740) circle[radius=0.6bp];
\fill[tg808080] (113.899,120.489) circle[radius=0.6bp];
\fill[tg808080] (136.029,103.775) circle[radius=0.6bp];
\fill[tg808080] (123.204,116.271) circle[radius=0.6bp];
\fill[tg808080] (122.627,118.225) circle[radius=0.6bp];
\fill[tg808080] (122.901,113.709) circle[radius=0.6bp];
\fill[tg808080] (117.441,120.645) circle[radius=0.6bp];
\fill[tg808080] (106.507,134.066) circle[radius=0.6bp];
\fill[tg808080] (118.651,116.143) circle[radius=0.6bp];
\fill[tg808080] (124.591,110.432) circle[radius=0.6bp];
\fill[tg808080] (115.964,113.555) circle[radius=0.6bp];
\fill[tg808080] (119.571,115.739) circle[radius=0.6bp];
\fill[tg808080] (114.099,118.562) circle[radius=0.6bp];
\fill[tg808080] (106.353,126.005) circle[radius=0.6bp];
\fill[tg808080] (126.593,113.765) circle[radius=0.6bp];
\fill[tg808080] (115.800,116.210) circle[radius=0.6bp];
\fill[tg808080] (111.958,115.308) circle[radius=0.6bp];
\fill[tg808080] (116.198,117.836) circle[radius=0.6bp];
\fill[tg808080] (110.405,117.326) circle[radius=0.6bp];
\fill[tg808080] (117.931,118.193) circle[radius=0.6bp];
\fill[tg808080] (127.765,106.880) circle[radius=0.6bp];
\fill[tg808080] (123.905,111.484) circle[radius=0.6bp];
\fill[tg808080] (125.445,111.300) circle[radius=0.6bp];
\fill[tg808080] (116.022,117.506) circle[radius=0.6bp];
\fill[tg808080] (122.456,120.114) circle[radius=0.6bp];
\fill[tg808080] (110.835,130.952) circle[radius=0.6bp];
\fill[tg808080] (116.183,121.755) circle[radius=0.6bp];
\fill[tg808080] (104.909,129.123) circle[radius=0.6bp];
\fill[tg808080] (116.238,119.313) circle[radius=0.6bp];
\fill[tg808080] (118.711,126.814) circle[radius=0.6bp];
\fill[tg808080] (135.346,108.928) circle[radius=0.6bp];
\fill[tg808080] (119.476,114.993) circle[radius=0.6bp];
\fill[tg808080] (119.578,119.299) circle[radius=0.6bp];
\fill[tg808080] (131.338,108.971) circle[radius=0.6bp];
\fill[tg808080] (114.949,128.947) circle[radius=0.6bp];
\fill[tg808080] (122.380,110.696) circle[radius=0.6bp];
\fill[tg808080] (106.730,129.181) circle[radius=0.6bp];
\fill[tg808080] (122.176,115.255) circle[radius=0.6bp];
\fill[tg808080] (129.928,103.457) circle[radius=0.6bp];
\fill[tg808080] (123.730,112.502) circle[radius=0.6bp];
\fill[tg808080] (107.808,130.471) circle[radius=0.6bp];
\fill[tg808080] (113.221,119.159) circle[radius=0.6bp];
\fill[tg808080] (118.787,126.596) circle[radius=0.6bp];
\fill[tg808080] (96.757,137.522) circle[radius=0.6bp];
\fill[tg808080] (117.607,125.950) circle[radius=0.6bp];
\fill[tg808080] (117.321,117.001) circle[radius=0.6bp];
\fill[tg808080] (131.755,107.817) circle[radius=0.6bp];
\fill[tg808080] (128.555,113.506) circle[radius=0.6bp];
\fill[tg808080] (118.492,114.891) circle[radius=0.6bp];
\fill[tg808080] (114.931,117.949) circle[radius=0.6bp];
\fill[tg808080] (109.282,122.324) circle[radius=0.6bp];
\fill[tg808080] (127.583,109.565) circle[radius=0.6bp];
\fill[tg808080] (119.443,121.584) circle[radius=0.6bp];
\fill[tg808080] (110.870,125.814) circle[radius=0.6bp];
\fill[tg808080] (133.159,112.145) circle[radius=0.6bp];
\fill[tg808080] (109.349,123.660) circle[radius=0.6bp];
\fill[tg808080] (131.019,113.266) circle[radius=0.6bp];
\fill[tg808080] (118.254,120.180) circle[radius=0.6bp];
\fill[tg808080] (120.752,114.570) circle[radius=0.6bp];
\fill[tg808080] (116.337,115.134) circle[radius=0.6bp];
\fill[tg808080] (123.976,114.189) circle[radius=0.6bp];
\fill[tg808080] (127.756,110.824) circle[radius=0.6bp];
\fill[tg808080] (126.429,112.103) circle[radius=0.6bp];
\fill[tg808080] (131.508,107.349) circle[radius=0.6bp];
\fill[tg808080] (114.898,116.164) circle[radius=0.6bp];
\fill[tg808080] (130.970,102.172) circle[radius=0.6bp];
\fill[tg808080] (108.354,131.672) circle[radius=0.6bp];
\fill[tg808080] (117.689,116.187) circle[radius=0.6bp];
\fill[tg808080] (115.240,118.349) circle[radius=0.6bp];
\fill[tg808080] (127.114,112.353) circle[radius=0.6bp];
\fill[tg808080] (107.789,119.805) circle[radius=0.6bp];
\fill[tg808080] (121.637,116.209) circle[radius=0.6bp];
\fill[tg808080] (66.824,172.333) circle[radius=0.6bp];
\fill[tg808080] (107.678,123.042) circle[radius=0.6bp];
\fill[tg808080] (126.880,111.911) circle[radius=0.6bp];
\fill[tg808080] (123.099,119.123) circle[radius=0.6bp];
\fill[tg808080] (117.077,121.204) circle[radius=0.6bp];
\fill[tg808080] (116.360,128.900) circle[radius=0.6bp];
\fill[tg808080] (130.621,110.478) circle[radius=0.6bp];
\fill[tg808080] (117.660,118.736) circle[radius=0.6bp];
\fill[tg808080] (125.257,109.007) circle[radius=0.6bp];
\fill[tg808080] (118.698,118.410) circle[radius=0.6bp];
\fill[tg808080] (120.530,124.369) circle[radius=0.6bp];
\fill[tg808080] (113.317,124.046) circle[radius=0.6bp];
\fill[tg808080] (129.960,112.208) circle[radius=0.6bp];
\fill[tg808080] (114.088,130.859) circle[radius=0.6bp];
\fill[tg808080] (103.053,133.027) circle[radius=0.6bp];
\fill[tg808080] (117.391,119.393) circle[radius=0.6bp];
\fill[tg808080] (133.981,106.361) circle[radius=0.6bp];
\fill[tg808080] (124.532,113.157) circle[radius=0.6bp];
\fill[tg808080] (108.797,126.129) circle[radius=0.6bp];
\fill[tg808080] (128.627,144.796) circle[radius=0.6bp];
\fill[tg808080] (112.350,125.726) circle[radius=0.6bp];
\fill[tg808080] (115.627,111.121) circle[radius=0.6bp];
\fill[tg808080] (116.683,118.382) circle[radius=0.6bp];
\fill[tg808080] (116.789,116.864) circle[radius=0.6bp];
\fill[tg808080] (122.084,120.759) circle[radius=0.6bp];
\fill[tg808080] (115.330,123.356) circle[radius=0.6bp];
\fill[tg808080] (119.729,116.631) circle[radius=0.6bp];
\fill[tg808080] (126.418,114.538) circle[radius=0.6bp];
\fill[tg808080] (114.480,121.738) circle[radius=0.6bp];
\fill[tg808080] (118.969,115.871) circle[radius=0.6bp];
\fill[tg808080] (116.149,114.032) circle[radius=0.6bp];
\fill[tg808080] (103.541,134.284) circle[radius=0.6bp];
\fill[tg808080] (111.426,122.803) circle[radius=0.6bp];
\fill[tg808080] (132.985,105.327) circle[radius=0.6bp];
\fill[tg808080] (121.175,116.277) circle[radius=0.6bp];
\fill[tg808080] (124.369,128.362) circle[radius=0.6bp];
\fill[tg808080] (100.751,122.566) circle[radius=0.6bp];
\fill[tg808080] (119.309,119.362) circle[radius=0.6bp];
\fill[tg808080] (126.834,113.040) circle[radius=0.6bp];
\fill[tg808080] (120.888,113.011) circle[radius=0.6bp];
\fill[tg808080] (119.432,118.544) circle[radius=0.6bp];
\fill[tg808080] (117.474,115.595) circle[radius=0.6bp];
\fill[tg808080] (125.501,113.810) circle[radius=0.6bp];
\fill[tg808080] (127.297,116.682) circle[radius=0.6bp];
\fill[tg808080] (116.784,121.389) circle[radius=0.6bp];
\fill[tg808080] (122.323,114.869) circle[radius=0.6bp];
\fill[tg808080] (116.599,117.585) circle[radius=0.6bp];
\fill[tg808080] (125.224,111.458) circle[radius=0.6bp];
\fill[tg808080] (112.717,113.426) circle[radius=0.6bp];
\fill[tg808080] (114.435,112.670) circle[radius=0.6bp];
\fill[tg808080] (128.059,114.388) circle[radius=0.6bp];
\fill[tg808080] (113.530,123.042) circle[radius=0.6bp];
\fill[tg808080] (124.615,112.761) circle[radius=0.6bp];
\fill[tg808080] (110.314,125.184) circle[radius=0.6bp];
\fill[tg808080] (120.890,119.088) circle[radius=0.6bp];
\fill[tg808080] (116.359,123.934) circle[radius=0.6bp];
\fill[tg808080] (112.933,132.028) circle[radius=0.6bp];
\fill[tg808080] (110.495,119.482) circle[radius=0.6bp];
\fill[tg808080] (124.346,110.019) circle[radius=0.6bp];
\fill[tg808080] (117.226,114.303) circle[radius=0.6bp];
\fill[tg808080] (114.115,119.901) circle[radius=0.6bp];
\fill[tg808080] (123.893,113.318) circle[radius=0.6bp];
\fill[tg808080] (130.798,107.593) circle[radius=0.6bp];
\fill[tg808080] (110.228,128.355) circle[radius=0.6bp];
\fill[tg808080] (117.574,116.779) circle[radius=0.6bp];
\fill[tg808080] (122.380,114.407) circle[radius=0.6bp];
\fill[tg808080] (126.306,106.611) circle[radius=0.6bp];
\fill[tg808080] (135.038,105.793) circle[radius=0.6bp];
\fill[tg808080] (122.750,109.272) circle[radius=0.6bp];
\fill[tg808080] (118.404,120.905) circle[radius=0.6bp];
\fill[tg808080] (109.057,110.513) circle[radius=0.6bp];
\fill[tg808080] (117.390,119.684) circle[radius=0.6bp];
\fill[tg808080] (122.759,117.912) circle[radius=0.6bp];
\fill[tg808080] (115.846,117.352) circle[radius=0.6bp];
\fill[tg808080] (123.323,113.031) circle[radius=0.6bp];
\fill[tg808080] (112.589,125.638) circle[radius=0.6bp];
\fill[tg808080] (105.313,134.114) circle[radius=0.6bp];
\fill[tg808080] (121.429,106.914) circle[radius=0.6bp];
\fill[tg808080] (109.797,125.420) circle[radius=0.6bp];
\fill[tg808080] (130.990,110.422) circle[radius=0.6bp];
\fill[tg808080] (120.873,117.614) circle[radius=0.6bp];
\fill[tg808080] (120.409,124.849) circle[radius=0.6bp];
\fill[tg808080] (124.227,110.988) circle[radius=0.6bp];
\fill[tg808080] (118.452,116.085) circle[radius=0.6bp];
\fill[tg808080] (101.190,131.491) circle[radius=0.6bp];
\fill[tg808080] (129.844,105.283) circle[radius=0.6bp];
\fill[tg808080] (134.162,111.602) circle[radius=0.6bp];
\fill[tg808080] (116.583,119.030) circle[radius=0.6bp];
\fill[tg808080] (126.807,123.990) circle[radius=0.6bp];
\fill[tg808080] (120.157,114.917) circle[radius=0.6bp];
\fill[tg808080] (104.722,128.656) circle[radius=0.6bp];
\fill[tg808080] (114.995,121.430) circle[radius=0.6bp];
\fill[tg808080] (128.505,111.351) circle[radius=0.6bp];
\fill[tg808080] (118.149,119.484) circle[radius=0.6bp];
\fill[tg808080] (114.605,125.708) circle[radius=0.6bp];
\fill[tg808080] (113.245,117.740) circle[radius=0.6bp];
\fill[tg808080] (110.116,126.827) circle[radius=0.6bp];
\fill[tg808080] (109.494,127.315) circle[radius=0.6bp];
\fill[tg808080] (132.019,96.529) circle[radius=0.6bp];
\fill[tg808080] (114.136,121.365) circle[radius=0.6bp];
\fill[tg808080] (115.544,117.402) circle[radius=0.6bp];
\fill[tg808080] (119.452,120.389) circle[radius=0.6bp];
\fill[tg808080] (103.056,125.843) circle[radius=0.6bp];
\fill[tg808080] (131.020,103.636) circle[radius=0.6bp];
\fill[tg808080] (121.538,110.398) circle[radius=0.6bp];
\fill[tg808080] (124.999,103.130) circle[radius=0.6bp];
\fill[tg808080] (119.230,120.917) circle[radius=0.6bp];
\fill[tg808080] (77.962,163.641) circle[radius=0.6bp];
\fill[tg808080] (126.316,113.392) circle[radius=0.6bp];
\fill[tg808080] (123.398,110.051) circle[radius=0.6bp];
\fill[tg808080] (83.851,155.233) circle[radius=0.6bp];
\fill[tg808080] (118.432,121.589) circle[radius=0.6bp];
\fill[tg808080] (114.097,122.693) circle[radius=0.6bp];
\fill[tg808080] (112.373,123.985) circle[radius=0.6bp];
\fill[tg808080] (101.575,126.486) circle[radius=0.6bp];
\fill[tg808080] (122.213,111.690) circle[radius=0.6bp];
\fill[tg808080] (128.404,114.561) circle[radius=0.6bp];
\fill[tg808080] (111.898,121.543) circle[radius=0.6bp];
\fill[tg808080] (125.675,110.945) circle[radius=0.6bp];
\fill[tg808080] (120.310,114.723) circle[radius=0.6bp];
\fill[tg808080] (118.908,110.669) circle[radius=0.6bp];
\fill[tg808080] (109.136,122.378) circle[radius=0.6bp];
\fill[tg808080] (104.247,126.937) circle[radius=0.6bp];
\fill[tg808080] (113.967,113.931) circle[radius=0.6bp];
\fill[tg808080] (120.798,110.407) circle[radius=0.6bp];
\fill[tg808080] (118.513,113.260) circle[radius=0.6bp];
\fill[tg808080] (117.291,123.075) circle[radius=0.6bp];
\fill[tg808080] (119.269,131.214) circle[radius=0.6bp];
\fill[tg808080] (112.460,124.104) circle[radius=0.6bp];
\fill[tg808080] (120.586,109.281) circle[radius=0.6bp];
\fill[tg808080] (114.253,116.602) circle[radius=0.6bp];
\fill[tg808080] (114.795,114.767) circle[radius=0.6bp];
\fill[tg808080] (114.151,117.921) circle[radius=0.6bp];
\fill[tg808080] (100.989,83.284) circle[radius=0.6bp];
\fill[tg808080] (111.760,126.411) circle[radius=0.6bp];
\fill[tg808080] (127.383,116.284) circle[radius=0.6bp];
\fill[tg808080] (123.842,116.572) circle[radius=0.6bp];
\fill[tg808080] (117.617,113.245) circle[radius=0.6bp];
\fill[tg808080] (129.337,107.368) circle[radius=0.6bp];
\fill[tg808080] (138.360,106.317) circle[radius=0.6bp];
\fill[tg808080] (133.439,119.479) circle[radius=0.6bp];
\fill[tg808080] (120.401,111.540) circle[radius=0.6bp];
\fill[tg808080] (110.982,122.112) circle[radius=0.6bp];
\fill[tg808080] (113.613,128.304) circle[radius=0.6bp];
\fill[tg808080] (103.713,121.612) circle[radius=0.6bp];
\fill[tg808080] (133.239,107.722) circle[radius=0.6bp];
\fill[tg808080] (125.592,107.275) circle[radius=0.6bp];
\fill[tg808080] (105.036,124.406) circle[radius=0.6bp];
\fill[tg808080] (118.826,124.541) circle[radius=0.6bp];
\fill[tg808080] (118.744,123.188) circle[radius=0.6bp];
\fill[tg808080] (104.858,124.841) circle[radius=0.6bp];
\fill[tg808080] (113.473,115.016) circle[radius=0.6bp];
\fill[tg808080] (108.967,133.222) circle[radius=0.6bp];
\fill[tg808080] (125.844,105.190) circle[radius=0.6bp];
\fill[tg808080] (119.823,109.497) circle[radius=0.6bp];
\fill[tg808080] (112.214,124.401) circle[radius=0.6bp];
\fill[tg808080] (112.485,124.313) circle[radius=0.6bp];
\fill[tg808080] (115.019,110.412) circle[radius=0.6bp];
\fill[tg808080] (116.151,121.196) circle[radius=0.6bp];
\fill[tg808080] (119.507,113.077) circle[radius=0.6bp];
\fill[tg808080] (125.840,110.003) circle[radius=0.6bp];
\fill[tg808080] (124.091,117.111) circle[radius=0.6bp];
\fill[tg808080] (128.950,108.650) circle[radius=0.6bp];
\fill[tg808080] (125.741,117.952) circle[radius=0.6bp];
\fill[tg808080] (123.698,115.596) circle[radius=0.6bp];
\fill[tg808080] (120.431,112.461) circle[radius=0.6bp];
\fill[tg808080] (119.678,116.025) circle[radius=0.6bp];
\fill[tg808080] (125.421,113.442) circle[radius=0.6bp];
\fill[tg808080] (122.233,110.339) circle[radius=0.6bp];
\fill[tg808080] (125.391,111.204) circle[radius=0.6bp];
\fill[tg808080] (99.766,126.249) circle[radius=0.6bp];
\fill[tg808080] (89.325,112.632) circle[radius=0.6bp];
\fill[tg808080] (101.543,124.373) circle[radius=0.6bp];
\fill[tg808080] (117.916,115.425) circle[radius=0.6bp];
\fill[tg808080] (120.934,120.564) circle[radius=0.6bp];
\fill[tg808080] (114.913,119.891) circle[radius=0.6bp];
\fill[tg808080] (114.557,118.592) circle[radius=0.6bp];
\fill[tg808080] (113.072,111.152) circle[radius=0.6bp];
\fill[tg808080] (126.811,109.192) circle[radius=0.6bp];
\fill[tg808080] (141.575,98.894) circle[radius=0.6bp];
\fill[tg808080] (127.125,109.752) circle[radius=0.6bp];
\fill[tg808080] (127.847,107.773) circle[radius=0.6bp];
\fill[tg808080] (111.952,129.791) circle[radius=0.6bp];
\fill[tg808080] (109.172,126.508) circle[radius=0.6bp];
\fill[tg808080] (106.434,124.783) circle[radius=0.6bp];
\fill[tg808080] (123.042,111.717) circle[radius=0.6bp];
\fill[tg808080] (120.352,118.930) circle[radius=0.6bp];
\fill[tg808080] (114.299,122.205) circle[radius=0.6bp];
\fill[tg808080] (120.535,117.805) circle[radius=0.6bp];
\fill[tg808080] (132.622,102.837) circle[radius=0.6bp];
\fill[tg808080] (123.229,118.824) circle[radius=0.6bp];
\fill[tg808080] (122.219,107.609) circle[radius=0.6bp];
\fill[tg808080] (121.701,118.736) circle[radius=0.6bp];
\fill[tg808080] (114.540,126.378) circle[radius=0.6bp];
\fill[tg808080] (126.963,112.219) circle[radius=0.6bp];
\fill[tg808080] (123.062,114.537) circle[radius=0.6bp];
\fill[tg808080] (123.783,108.517) circle[radius=0.6bp];
\fill[tg808080] (114.714,118.778) circle[radius=0.6bp];
\fill[tg808080] (120.228,118.473) circle[radius=0.6bp];
\fill[tg808080] (127.504,122.096) circle[radius=0.6bp];
\fill[tg808080] (108.971,118.964) circle[radius=0.6bp];
\fill[tg808080] (126.947,115.276) circle[radius=0.6bp];
\fill[tg808080] (111.939,117.733) circle[radius=0.6bp];
\fill[tg808080] (119.272,119.408) circle[radius=0.6bp];
\fill[tg808080] (109.187,116.526) circle[radius=0.6bp];
\fill[tg808080] (121.341,117.267) circle[radius=0.6bp];
\fill[tg808080] (116.024,121.476) circle[radius=0.6bp];
\fill[tg808080] (122.262,108.601) circle[radius=0.6bp];
\fill[tg808080] (114.872,120.691) circle[radius=0.6bp];
\fill[tg808080] (146.750,96.194) circle[radius=0.6bp];
\fill[tg808080] (115.396,118.313) circle[radius=0.6bp];
\fill[tg808080] (118.368,112.616) circle[radius=0.6bp];
\fill[tg808080] (108.883,131.424) circle[radius=0.6bp];
\fill[tg808080] (122.583,110.240) circle[radius=0.6bp];
\fill[tg808080] (119.926,118.987) circle[radius=0.6bp];
\fill[tg808080] (108.795,127.840) circle[radius=0.6bp];
\fill[tg808080] (128.200,114.120) circle[radius=0.6bp];
\fill[tg808080] (130.513,110.403) circle[radius=0.6bp];
\fill[tg808080] (111.063,122.871) circle[radius=0.6bp];
\fill[tg808080] (110.857,118.722) circle[radius=0.6bp];
\draw[color=tg0000ff,line width=0.7bp,line cap=round] (114.800,116.800) -- (119.200,121.200);
\draw[color=tg0000ff,line width=0.7bp,line cap=round] (114.800,121.200) -- (119.200,116.800);
\node[anchor=center,inner sep=0pt,outer sep=0pt,text=tg18314b,font={\sffamily\fontsize{6.5bp}{7.8bp}\selectfont}] at (67.222,183.000) {-4};
\node[anchor=east,inner sep=0pt,outer sep=0pt,text=tg18314b,font={\sffamily\fontsize{6.5bp}{7.8bp}\selectfont}] at (56.000,168.778) {-4};
\node[anchor=center,inner sep=0pt,outer sep=0pt,text=tg18314b,font={\sffamily\fontsize{6.5bp}{7.8bp}\selectfont}] at (117.000,183.000) {0};
\node[anchor=east,inner sep=0pt,outer sep=0pt,text=tg18314b,font={\sffamily\fontsize{6.5bp}{7.8bp}\selectfont}] at (56.000,119.000) {0};
\node[anchor=center,inner sep=0pt,outer sep=0pt,text=tg18314b,font={\sffamily\fontsize{6.5bp}{7.8bp}\selectfont}] at (166.778,183.000) {4};
\node[anchor=east,inner sep=0pt,outer sep=0pt,text=tg18314b,font={\sffamily\fontsize{6.5bp}{7.8bp}\selectfont}] at (56.000,69.222) {4};
\node[anchor=center,inner sep=0pt,outer sep=0pt,text=tg18314b,font={\sffamily\fontsize{7bp}{8.4bp}\selectfont}] at (117.000,195.000) {Standardized target 0};
\node[anchor=west,inner sep=0pt,outer sep=0pt,text=tg626b73,font={\sffamily\fontsize{6.5bp}{7.8bp}\selectfont}] at (61.000,57.000) {Target 1};
\node[anchor=center,inner sep=0pt,outer sep=0pt,text=tg18314b,font={\sffamily\fontsize{8bp}{9.6bp}\selectfont}] at (117.000,209.000) {Across horizons and targets};
\path[fill=tg67001f,draw=none,line width=0.5bp] (61.000,232.000) rectangle (70.383,241.383);
\path[fill=tg991027,draw=none,line width=0.5bp] (70.333,232.000) rectangle (79.717,241.383);
\path[fill=tgba2832,draw=none,line width=0.5bp] (79.667,232.000) rectangle (89.050,241.383);
\path[fill=tgd05548,draw=none,line width=0.5bp] (89.000,232.000) rectangle (98.383,241.383);
\path[fill=tgde735c,draw=none,line width=0.5bp] (98.333,232.000) rectangle (107.717,241.383);
\path[fill=tgea8e70,draw=none,line width=0.5bp] (107.667,232.000) rectangle (117.050,241.383);
\path[fill=tgea8e70,draw=none,line width=0.5bp] (117.000,232.000) rectangle (126.383,241.383);
\path[fill=tgf5a886,draw=none,line width=0.5bp] (126.333,232.000) rectangle (135.717,241.383);
\path[fill=tgf4a683,draw=none,line width=0.5bp] (135.667,232.000) rectangle (145.050,241.383);
\path[fill=tgf7b799,draw=none,line width=0.5bp] (145.000,232.000) rectangle (154.383,241.383);
\path[fill=tgf5ac8b,draw=none,line width=0.5bp] (154.333,232.000) rectangle (163.717,241.383);
\path[fill=tgf8bb9e,draw=none,line width=0.5bp] (163.667,232.000) rectangle (173.050,241.383);
\path[fill=tg991027,draw=none,line width=0.5bp] (61.000,241.333) rectangle (70.383,250.717);
\path[fill=tg67001f,draw=none,line width=0.5bp] (70.333,241.333) rectangle (79.717,250.717);
\path[fill=tgce4f45,draw=none,line width=0.5bp] (79.667,241.333) rectangle (89.050,250.717);
\path[fill=tgc2383a,draw=none,line width=0.5bp] (89.000,241.333) rectangle (98.383,250.717);
\path[fill=tge98b6e,draw=none,line width=0.5bp] (98.333,241.333) rectangle (107.717,250.717);
\path[fill=tgdf765e,draw=none,line width=0.5bp] (107.667,241.333) rectangle (117.050,250.717);
\path[fill=tgf2a17f,draw=none,line width=0.5bp] (117.000,241.333) rectangle (126.383,250.717);
\path[fill=tgea8e70,draw=none,line width=0.5bp] (126.333,241.333) rectangle (135.717,250.717);
\path[fill=tgf6b394,draw=none,line width=0.5bp] (135.667,241.333) rectangle (145.050,250.717);
\path[fill=tgf3a481,draw=none,line width=0.5bp] (145.000,241.333) rectangle (154.383,250.717);
\path[fill=tgf8bb9e,draw=none,line width=0.5bp] (154.333,241.333) rectangle (163.717,250.717);
\path[fill=tgf6af8e,draw=none,line width=0.5bp] (163.667,241.333) rectangle (173.050,250.717);
\path[fill=tgba2832,draw=none,line width=0.5bp] (61.000,250.667) rectangle (70.383,260.050);
\path[fill=tgce4f45,draw=none,line width=0.5bp] (70.333,250.667) rectangle (79.717,260.050);
\path[fill=tg67001f,draw=none,line width=0.5bp] (79.667,250.667) rectangle (89.050,260.050);
\path[fill=tgab162a,draw=none,line width=0.5bp] (89.000,250.667) rectangle (98.383,260.050);
\path[fill=tgc13639,draw=none,line width=0.5bp] (98.333,250.667) rectangle (107.717,260.050);
\path[fill=tgd6604d,draw=none,line width=0.5bp] (107.667,250.667) rectangle (117.050,260.050);
\path[fill=tgd55d4c,draw=none,line width=0.5bp] (117.000,250.667) rectangle (126.383,260.050);
\path[fill=tge58368,draw=none,line width=0.5bp] (126.333,250.667) rectangle (135.717,260.050);
\path[fill=tge17860,draw=none,line width=0.5bp] (135.667,250.667) rectangle (145.050,260.050);
\path[fill=tgef9979,draw=none,line width=0.5bp] (145.000,250.667) rectangle (154.383,260.050);
\path[fill=tge58368,draw=none,line width=0.5bp] (154.333,250.667) rectangle (163.717,260.050);
\path[fill=tgf2a17f,draw=none,line width=0.5bp] (163.667,250.667) rectangle (173.050,260.050);
\path[fill=tgd05548,draw=none,line width=0.5bp] (61.000,260.000) rectangle (70.383,269.383);
\path[fill=tgc2383a,draw=none,line width=0.5bp] (70.333,260.000) rectangle (79.717,269.383);
\path[fill=tgab162a,draw=none,line width=0.5bp] (79.667,260.000) rectangle (89.050,269.383);
\path[fill=tg67001f,draw=none,line width=0.5bp] (89.000,260.000) rectangle (98.383,269.383);
\path[fill=tgd7634f,draw=none,line width=0.5bp] (98.333,260.000) rectangle (107.717,269.383);
\path[fill=tgba2832,draw=none,line width=0.5bp] (107.667,260.000) rectangle (117.050,269.383);
\path[fill=tge48066,draw=none,line width=0.5bp] (117.000,260.000) rectangle (126.383,269.383);
\path[fill=tgce4f45,draw=none,line width=0.5bp] (126.333,260.000) rectangle (135.717,269.383);
\path[fill=tgee9677,draw=none,line width=0.5bp] (135.667,260.000) rectangle (145.050,269.383);
\path[fill=tgdb6b55,draw=none,line width=0.5bp] (145.000,260.000) rectangle (154.383,269.383);
\path[fill=tgf2a17f,draw=none,line width=0.5bp] (154.333,260.000) rectangle (163.717,269.383);
\path[fill=tge17860,draw=none,line width=0.5bp] (163.667,260.000) rectangle (173.050,269.383);
\path[fill=tgde735c,draw=none,line width=0.5bp] (61.000,269.333) rectangle (70.383,278.717);
\path[fill=tge98b6e,draw=none,line width=0.5bp] (70.333,269.333) rectangle (79.717,278.717);
\path[fill=tgc13639,draw=none,line width=0.5bp] (79.667,269.333) rectangle (89.050,278.717);
\path[fill=tgd7634f,draw=none,line width=0.5bp] (89.000,269.333) rectangle (98.383,278.717);
\path[fill=tg67001f,draw=none,line width=0.5bp] (98.333,269.333) rectangle (107.717,278.717);
\path[fill=tga51429,draw=none,line width=0.5bp] (107.667,269.333) rectangle (117.050,278.717);
\path[fill=tg991027,draw=none,line width=0.5bp] (117.000,269.333) rectangle (126.383,278.717);
\path[fill=tgc2383a,draw=none,line width=0.5bp] (126.333,269.333) rectangle (135.717,278.717);
\path[fill=tgbd2d35,draw=none,line width=0.5bp] (135.667,269.333) rectangle (145.050,278.717);
\path[fill=tgd35a4a,draw=none,line width=0.5bp] (145.000,269.333) rectangle (154.383,278.717);
\path[fill=tgc53e3d,draw=none,line width=0.5bp] (154.333,269.333) rectangle (163.717,278.717);
\path[fill=tgd86551,draw=none,line width=0.5bp] (163.667,269.333) rectangle (173.050,278.717);
\path[fill=tgea8e70,draw=none,line width=0.5bp] (61.000,278.667) rectangle (70.383,288.050);
\path[fill=tgdf765e,draw=none,line width=0.5bp] (70.333,278.667) rectangle (79.717,288.050);
\path[fill=tgd6604d,draw=none,line width=0.5bp] (79.667,278.667) rectangle (89.050,288.050);
\path[fill=tgba2832,draw=none,line width=0.5bp] (89.000,278.667) rectangle (98.383,288.050);
\path[fill=tga51429,draw=none,line width=0.5bp] (98.333,278.667) rectangle (107.717,288.050);
\path[fill=tg67001f,draw=none,line width=0.5bp] (107.667,278.667) rectangle (117.050,288.050);
\path[fill=tgbf3338,draw=none,line width=0.5bp] (117.000,278.667) rectangle (126.383,288.050);
\path[fill=tg991027,draw=none,line width=0.5bp] (126.333,278.667) rectangle (135.717,288.050);
\path[fill=tgd35a4a,draw=none,line width=0.5bp] (135.667,278.667) rectangle (145.050,288.050);
\path[fill=tgb82531,draw=none,line width=0.5bp] (145.000,278.667) rectangle (154.383,288.050);
\path[fill=tgda6853,draw=none,line width=0.5bp] (154.333,278.667) rectangle (163.717,288.050);
\path[fill=tgc13639,draw=none,line width=0.5bp] (163.667,278.667) rectangle (173.050,288.050);
\path[fill=tgea8e70,draw=none,line width=0.5bp] (61.000,288.000) rectangle (70.383,297.383);
\path[fill=tgf2a17f,draw=none,line width=0.5bp] (70.333,288.000) rectangle (79.717,297.383);
\path[fill=tgd55d4c,draw=none,line width=0.5bp] (79.667,288.000) rectangle (89.050,297.383);
\path[fill=tge48066,draw=none,line width=0.5bp] (89.000,288.000) rectangle (98.383,297.383);
\path[fill=tg991027,draw=none,line width=0.5bp] (98.333,288.000) rectangle (107.717,297.383);
\path[fill=tgbf3338,draw=none,line width=0.5bp] (107.667,288.000) rectangle (117.050,297.383);
\path[fill=tg67001f,draw=none,line width=0.5bp] (117.000,288.000) rectangle (126.383,297.383);
\path[fill=tga21328,draw=none,line width=0.5bp] (126.333,288.000) rectangle (135.717,297.383);
\path[fill=tg960f27,draw=none,line width=0.5bp] (135.667,288.000) rectangle (145.050,297.383);
\path[fill=tgbd2d35,draw=none,line width=0.5bp] (145.000,288.000) rectangle (154.383,297.383);
\path[fill=tgab162a,draw=none,line width=0.5bp] (154.333,288.000) rectangle (163.717,297.383);
\path[fill=tgc2383a,draw=none,line width=0.5bp] (163.667,288.000) rectangle (173.050,297.383);
\path[fill=tgf5a886,draw=none,line width=0.5bp] (61.000,297.333) rectangle (70.383,306.717);
\path[fill=tgea8e70,draw=none,line width=0.5bp] (70.333,297.333) rectangle (79.717,306.717);
\path[fill=tge58368,draw=none,line width=0.5bp] (79.667,297.333) rectangle (89.050,306.717);
\path[fill=tgce4f45,draw=none,line width=0.5bp] (89.000,297.333) rectangle (98.383,306.717);
\path[fill=tgc2383a,draw=none,line width=0.5bp] (98.333,297.333) rectangle (107.717,306.717);
\path[fill=tg991027,draw=none,line width=0.5bp] (107.667,297.333) rectangle (117.050,306.717);
\path[fill=tga21328,draw=none,line width=0.5bp] (117.000,297.333) rectangle (126.383,306.717);
\path[fill=tg67001f,draw=none,line width=0.5bp] (126.333,297.333) rectangle (135.717,306.717);
\path[fill=tgbf3338,draw=none,line width=0.5bp] (135.667,297.333) rectangle (145.050,306.717);
\path[fill=tg930e26,draw=none,line width=0.5bp] (145.000,297.333) rectangle (154.383,306.717);
\path[fill=tgc6413e,draw=none,line width=0.5bp] (154.333,297.333) rectangle (163.717,306.717);
\path[fill=tga51429,draw=none,line width=0.5bp] (163.667,297.333) rectangle (173.050,306.717);
\path[fill=tgf4a683,draw=none,line width=0.5bp] (61.000,306.667) rectangle (70.383,316.050);
\path[fill=tgf6b394,draw=none,line width=0.5bp] (70.333,306.667) rectangle (79.717,316.050);
\path[fill=tge17860,draw=none,line width=0.5bp] (79.667,306.667) rectangle (89.050,316.050);
\path[fill=tgee9677,draw=none,line width=0.5bp] (89.000,306.667) rectangle (98.383,316.050);
\path[fill=tgbd2d35,draw=none,line width=0.5bp] (98.333,306.667) rectangle (107.717,316.050);
\path[fill=tgd35a4a,draw=none,line width=0.5bp] (107.667,306.667) rectangle (117.050,316.050);
\path[fill=tg960f27,draw=none,line width=0.5bp] (117.000,306.667) rectangle (126.383,316.050);
\path[fill=tgbf3338,draw=none,line width=0.5bp] (126.333,306.667) rectangle (135.717,316.050);
\path[fill=tg67001f,draw=none,line width=0.5bp] (135.667,306.667) rectangle (145.050,316.050);
\path[fill=tg9c1127,draw=none,line width=0.5bp] (145.000,306.667) rectangle (154.383,316.050);
\path[fill=tg7c0722,draw=none,line width=0.5bp] (154.333,306.667) rectangle (163.717,316.050);
\path[fill=tgab162a,draw=none,line width=0.5bp] (163.667,306.667) rectangle (173.050,316.050);
\path[fill=tgf7b799,draw=none,line width=0.5bp] (61.000,316.000) rectangle (70.383,325.383);
\path[fill=tgf3a481,draw=none,line width=0.5bp] (70.333,316.000) rectangle (79.717,325.383);
\path[fill=tgef9979,draw=none,line width=0.5bp] (79.667,316.000) rectangle (89.050,325.383);
\path[fill=tgdb6b55,draw=none,line width=0.5bp] (89.000,316.000) rectangle (98.383,325.383);
\path[fill=tgd35a4a,draw=none,line width=0.5bp] (98.333,316.000) rectangle (107.717,325.383);
\path[fill=tgb82531,draw=none,line width=0.5bp] (107.667,316.000) rectangle (117.050,325.383);
\path[fill=tgbd2d35,draw=none,line width=0.5bp] (117.000,316.000) rectangle (126.383,325.383);
\path[fill=tg930e26,draw=none,line width=0.5bp] (126.333,316.000) rectangle (135.717,325.383);
\path[fill=tg9c1127,draw=none,line width=0.5bp] (135.667,316.000) rectangle (145.050,325.383);
\path[fill=tg67001f,draw=none,line width=0.5bp] (145.000,316.000) rectangle (154.383,325.383);
\path[fill=tgb1182b,draw=none,line width=0.5bp] (154.333,316.000) rectangle (163.717,325.383);
\path[fill=tg790622,draw=none,line width=0.5bp] (163.667,316.000) rectangle (173.050,325.383);
\path[fill=tgf5ac8b,draw=none,line width=0.5bp] (61.000,325.333) rectangle (70.383,334.717);
\path[fill=tgf8bb9e,draw=none,line width=0.5bp] (70.333,325.333) rectangle (79.717,334.717);
\path[fill=tge58368,draw=none,line width=0.5bp] (79.667,325.333) rectangle (89.050,334.717);
\path[fill=tgf2a17f,draw=none,line width=0.5bp] (89.000,325.333) rectangle (98.383,334.717);
\path[fill=tgc53e3d,draw=none,line width=0.5bp] (98.333,325.333) rectangle (107.717,334.717);
\path[fill=tgda6853,draw=none,line width=0.5bp] (107.667,325.333) rectangle (117.050,334.717);
\path[fill=tgab162a,draw=none,line width=0.5bp] (117.000,325.333) rectangle (126.383,334.717);
\path[fill=tgc6413e,draw=none,line width=0.5bp] (126.333,325.333) rectangle (135.717,334.717);
\path[fill=tg7c0722,draw=none,line width=0.5bp] (135.667,325.333) rectangle (145.050,334.717);
\path[fill=tgb1182b,draw=none,line width=0.5bp] (145.000,325.333) rectangle (154.383,334.717);
\path[fill=tg67001f,draw=none,line width=0.5bp] (154.333,325.333) rectangle (163.717,334.717);
\path[fill=tg9c1127,draw=none,line width=0.5bp] (163.667,325.333) rectangle (173.050,334.717);
\path[fill=tgf8bb9e,draw=none,line width=0.5bp] (61.000,334.667) rectangle (70.383,344.050);
\path[fill=tgf6af8e,draw=none,line width=0.5bp] (70.333,334.667) rectangle (79.717,344.050);
\path[fill=tgf2a17f,draw=none,line width=0.5bp] (79.667,334.667) rectangle (89.050,344.050);
\path[fill=tge17860,draw=none,line width=0.5bp] (89.000,334.667) rectangle (98.383,344.050);
\path[fill=tgd86551,draw=none,line width=0.5bp] (98.333,334.667) rectangle (107.717,344.050);
\path[fill=tgc13639,draw=none,line width=0.5bp] (107.667,334.667) rectangle (117.050,344.050);
\path[fill=tgc2383a,draw=none,line width=0.5bp] (117.000,334.667) rectangle (126.383,344.050);
\path[fill=tga51429,draw=none,line width=0.5bp] (126.333,334.667) rectangle (135.717,344.050);
\path[fill=tgab162a,draw=none,line width=0.5bp] (135.667,334.667) rectangle (145.050,344.050);
\path[fill=tg790622,draw=none,line width=0.5bp] (145.000,334.667) rectangle (154.383,344.050);
\path[fill=tg9c1127,draw=none,line width=0.5bp] (154.333,334.667) rectangle (163.717,344.050);
\path[fill=tg67001f,draw=none,line width=0.5bp] (163.667,334.667) rectangle (173.050,344.050);
\node[anchor=east,inner sep=0pt,outer sep=0pt,text=tg18314b,font={\sffamily\fontsize{6bp}{7.2bp}\selectfont}] at (57.000,236.667) {h1:y0};
\node[anchor=center,inner sep=0pt,outer sep=0pt,text=tg18314b,font={\sffamily\fontsize{6bp}{7.2bp}\selectfont}] at (65.667,352.000) {h1:y0};
\node[anchor=east,inner sep=0pt,outer sep=0pt,text=tg18314b,font={\sffamily\fontsize{6bp}{7.2bp}\selectfont}] at (57.000,274.000) {h3:y0};
\node[anchor=center,inner sep=0pt,outer sep=0pt,text=tg18314b,font={\sffamily\fontsize{6bp}{7.2bp}\selectfont}] at (103.000,352.000) {h3:y0};
\node[anchor=east,inner sep=0pt,outer sep=0pt,text=tg18314b,font={\sffamily\fontsize{6bp}{7.2bp}\selectfont}] at (57.000,330.000) {h6:y0};
\node[anchor=center,inner sep=0pt,outer sep=0pt,text=tg18314b,font={\sffamily\fontsize{6bp}{7.2bp}\selectfont}] at (159.000,352.000) {h6:y0};
\node[anchor=center,inner sep=0pt,outer sep=0pt,text=tg18314b,font={\sffamily\fontsize{8bp}{9.6bp}\selectfont\bfseries}] at (117.000,369.000) {Aggregate SD: 0.91};
\node[anchor=center,inner sep=0pt,outer sep=0pt,text=tg18314b,font={\sffamily\fontsize{9bp}{10.8bp}\selectfont\bfseries}] at (335.000,48.000) {Partners shuffled};
\draw[color=tgc9cdd1,line width=0.5bp,line cap=round] (279.000,63.000) -- (279.000,175.000);
\draw[color=tgc9cdd1,line width=0.5bp,line cap=round] (279.000,175.000) -- (391.000,175.000);
\draw[color=tge1e4e7,line width=0.45bp,line cap=round] (335.000,63.000) -- (335.000,175.000);
\draw[color=tge1e4e7,line width=0.45bp,line cap=round] (279.000,119.000) -- (391.000,119.000);
\fill[tg808080] (325.864,116.271) circle[radius=0.6bp];
\fill[tg808080] (305.994,119.683) circle[radius=0.6bp];
\fill[tg808080] (334.241,122.627) circle[radius=0.6bp];
\fill[tg808080] (329.487,124.418) circle[radius=0.6bp];
\fill[tg808080] (315.312,113.731) circle[radius=0.6bp];
\fill[tg808080] (340.594,130.952) circle[radius=0.6bp];
\fill[tg808080] (327.453,114.540) circle[radius=0.6bp];
\fill[tg808080] (341.685,125.878) circle[radius=0.6bp];
\fill[tg808080] (335.025,126.486) circle[radius=0.6bp];
\fill[tg808080] (342.456,104.183) circle[radius=0.6bp];
\fill[tg808080] (319.965,114.032) circle[radius=0.6bp];
\fill[tg808080] (336.365,108.619) circle[radius=0.6bp];
\fill[tg808080] (323.526,118.912) circle[radius=0.6bp];
\fill[tg808080] (327.584,119.362) circle[radius=0.6bp];
\fill[tg808080] (323.196,114.398) circle[radius=0.6bp];
\fill[tg808080] (332.130,113.611) circle[radius=0.6bp];
\fill[tg808080] (334.941,118.187) circle[radius=0.6bp];
\fill[tg808080] (332.535,112.145) circle[radius=0.6bp];
\fill[tg808080] (342.735,120.917) circle[radius=0.6bp];
\fill[tg808080] (347.222,120.100) circle[radius=0.6bp];
\fill[tg808080] (335.202,125.851) circle[radius=0.6bp];
\fill[tg808080] (326.715,109.497) circle[radius=0.6bp];
\fill[tg808080] (348.933,107.262) circle[radius=0.6bp];
\fill[tg808080] (348.765,96.529) circle[radius=0.6bp];
\fill[tg808080] (328.051,116.526) circle[radius=0.6bp];
\fill[tg808080] (332.734,119.159) circle[radius=0.6bp];
\fill[tg808080] (326.417,126.005) circle[radius=0.6bp];
\fill[tg808080] (333.271,132.028) circle[radius=0.6bp];
\fill[tg808080] (337.220,113.847) circle[radius=0.6bp];
\fill[tg808080] (328.688,110.412) circle[radius=0.6bp];
\fill[tg808080] (325.345,114.538) circle[radius=0.6bp];
\fill[tg808080] (327.815,122.760) circle[radius=0.6bp];
\fill[tg808080] (340.422,131.229) circle[radius=0.6bp];
\fill[tg808080] (340.530,119.313) circle[radius=0.6bp];
\fill[tg808080] (354.031,106.611) circle[radius=0.6bp];
\fill[tg808080] (335.437,111.690) circle[radius=0.6bp];
\fill[tg808080] (341.898,120.039) circle[radius=0.6bp];
\fill[tg808080] (337.257,113.266) circle[radius=0.6bp];
\fill[tg808080] (347.526,114.232) circle[radius=0.6bp];
\fill[tg808080] (340.434,110.432) circle[radius=0.6bp];
\fill[tg808080] (344.384,116.602) circle[radius=0.6bp];
\fill[tg808080] (323.337,120.897) circle[radius=0.6bp];
\fill[tg808080] (333.470,125.708) circle[radius=0.6bp];
\fill[tg808080] (340.921,112.461) circle[radius=0.6bp];
\fill[tg808080] (334.482,128.304) circle[radius=0.6bp];
\fill[tg808080] (343.241,117.542) circle[radius=0.6bp];
\fill[tg808080] (357.160,119.859) circle[radius=0.6bp];
\fill[tg808080] (330.949,111.484) circle[radius=0.6bp];
\fill[tg808080] (333.054,107.817) circle[radius=0.6bp];
\fill[tg808080] (324.756,109.192) circle[radius=0.6bp];
\fill[tg808080] (333.380,120.233) circle[radius=0.6bp];
\fill[tg808080] (333.301,110.339) circle[radius=0.6bp];
\fill[tg808080] (330.546,115.255) circle[radius=0.6bp];
\fill[tg808080] (337.673,111.540) circle[radius=0.6bp];
\fill[tg808080] (356.932,118.531) circle[radius=0.6bp];
\fill[tg808080] (337.169,112.353) circle[radius=0.6bp];
\fill[tg808080] (332.005,107.964) circle[radius=0.6bp];
\fill[tg808080] (348.449,105.839) circle[radius=0.6bp];
\fill[tg808080] (347.113,116.025) circle[radius=0.6bp];
\fill[tg808080] (331.619,119.408) circle[radius=0.6bp];
\fill[tg808080] (338.638,109.621) circle[radius=0.6bp];
\fill[tg808080] (332.812,96.194) circle[radius=0.6bp];
\fill[tg808080] (336.440,133.094) circle[radius=0.6bp];
\fill[tg808080] (330.088,112.335) circle[radius=0.6bp];
\fill[tg808080] (323.714,121.123) circle[radius=0.6bp];
\fill[tg808080] (331.407,120.564) circle[radius=0.6bp];
\fill[tg808080] (336.845,124.007) circle[radius=0.6bp];
\fill[tg808080] (328.023,111.121) circle[radius=0.6bp];
\fill[tg808080] (339.112,112.901) circle[radius=0.6bp];
\fill[tg808080] (343.635,117.701) circle[radius=0.6bp];
\fill[tg808080] (341.216,120.389) circle[radius=0.6bp];
\fill[tg808080] (345.311,120.489) circle[radius=0.6bp];
\fill[tg808080] (340.266,116.037) circle[radius=0.6bp];
\fill[tg808080] (342.407,119.764) circle[radius=0.6bp];
\fill[tg808080] (327.912,105.190) circle[radius=0.6bp];
\fill[tg808080] (336.139,123.835) circle[radius=0.6bp];
\fill[tg808080] (334.627,163.641) circle[radius=0.6bp];
\fill[tg808080] (329.430,131.424) circle[radius=0.6bp];
\fill[tg808080] (340.382,122.693) circle[radius=0.6bp];
\fill[tg808080] (336.730,113.011) circle[radius=0.6bp];
\fill[tg808080] (327.151,117.980) circle[radius=0.6bp];
\fill[tg808080] (337.579,115.358) circle[radius=0.6bp];
\fill[tg808080] (326.847,110.988) circle[radius=0.6bp];
\fill[tg808080] (349.018,157.434) circle[radius=0.6bp];
\fill[tg808080] (336.195,117.529) circle[radius=0.6bp];
\fill[tg808080] (354.689,123.188) circle[radius=0.6bp];
\fill[tg808080] (322.622,96.870) circle[radius=0.6bp];
\fill[tg808080] (342.581,110.696) circle[radius=0.6bp];
\fill[tg808080] (335.622,126.688) circle[radius=0.6bp];
\fill[tg808080] (331.899,124.020) circle[radius=0.6bp];
\fill[tg808080] (354.029,108.928) circle[radius=0.6bp];
\fill[tg808080] (341.204,114.303) circle[radius=0.6bp];
\fill[tg808080] (340.627,119.393) circle[radius=0.6bp];
\fill[tg808080] (340.901,133.010) circle[radius=0.6bp];
\fill[tg808080] (335.441,124.783) circle[radius=0.6bp];
\fill[tg808080] (324.507,137.522) circle[radius=0.6bp];
\fill[tg808080] (336.651,120.987) circle[radius=0.6bp];
\fill[tg808080] (342.591,117.326) circle[radius=0.6bp];
\fill[tg808080] (333.964,111.540) circle[radius=0.6bp];
\fill[tg808080] (337.571,109.007) circle[radius=0.6bp];
\fill[tg808080] (332.099,121.204) circle[radius=0.6bp];
\fill[tg808080] (324.353,112.664) circle[radius=0.6bp];
\fill[tg808080] (344.593,115.435) circle[radius=0.6bp];
\fill[tg808080] (333.800,128.032) circle[radius=0.6bp];
\fill[tg808080] (329.958,106.953) circle[radius=0.6bp];
\fill[tg808080] (334.198,124.106) circle[radius=0.6bp];
\fill[tg808080] (328.405,122.246) circle[radius=0.6bp];
\fill[tg808080] (335.931,116.127) circle[radius=0.6bp];
\fill[tg808080] (345.765,121.813) circle[radius=0.6bp];
\fill[tg808080] (341.905,115.256) circle[radius=0.6bp];
\fill[tg808080] (343.445,114.335) circle[radius=0.6bp];
\fill[tg808080] (334.022,129.181) circle[radius=0.6bp];
\fill[tg808080] (340.456,128.669) circle[radius=0.6bp];
\fill[tg808080] (328.835,125.843) circle[radius=0.6bp];
\fill[tg808080] (334.183,127.893) circle[radius=0.6bp];
\fill[tg808080] (322.909,113.808) circle[radius=0.6bp];
\fill[tg808080] (334.238,114.078) circle[radius=0.6bp];
\fill[tg808080] (336.711,110.538) circle[radius=0.6bp];
\fill[tg808080] (353.346,113.668) circle[radius=0.6bp];
\fill[tg808080] (337.476,107.031) circle[radius=0.6bp];
\fill[tg808080] (337.578,119.421) circle[radius=0.6bp];
\fill[tg808080] (349.338,113.951) circle[radius=0.6bp];
\fill[tg808080] (332.949,134.441) circle[radius=0.6bp];
\fill[tg808080] (340.380,120.586) circle[radius=0.6bp];
\fill[tg808080] (324.730,119.076) circle[radius=0.6bp];
\fill[tg808080] (340.176,115.067) circle[radius=0.6bp];
\fill[tg808080] (347.928,126.262) circle[radius=0.6bp];
\fill[tg808080] (341.730,111.717) circle[radius=0.6bp];
\fill[tg808080] (325.808,112.103) circle[radius=0.6bp];
\fill[tg808080] (331.221,117.830) circle[radius=0.6bp];
\fill[tg808080] (336.787,120.905) circle[radius=0.6bp];
\fill[tg808080] (314.757,119.328) circle[radius=0.6bp];
\fill[tg808080] (335.607,120.691) circle[radius=0.6bp];
\fill[tg808080] (335.321,124.246) circle[radius=0.6bp];
\fill[tg808080] (349.755,106.608) circle[radius=0.6bp];
\fill[tg808080] (346.555,118.426) circle[radius=0.6bp];
\fill[tg808080] (336.492,121.589) circle[radius=0.6bp];
\fill[tg808080] (332.931,116.099) circle[radius=0.6bp];
\fill[tg808080] (327.282,114.551) circle[radius=0.6bp];
\fill[tg808080] (345.583,117.614) circle[radius=0.6bp];
\fill[tg808080] (337.443,114.891) circle[radius=0.6bp];
\fill[tg808080] (328.870,123.283) circle[radius=0.6bp];
\fill[tg808080] (351.159,107.698) circle[radius=0.6bp];
\fill[tg808080] (327.349,112.670) circle[radius=0.6bp];
\fill[tg808080] (349.019,118.410) circle[radius=0.6bp];
\fill[tg808080] (336.254,114.570) circle[radius=0.6bp];
\fill[tg808080] (338.752,110.793) circle[radius=0.6bp];
\fill[tg808080] (334.337,100.406) circle[radius=0.6bp];
\fill[tg808080] (341.976,118.930) circle[radius=0.6bp];
\fill[tg808080] (345.756,155.161) circle[radius=0.6bp];
\fill[tg808080] (344.429,118.193) circle[radius=0.6bp];
\fill[tg808080] (349.508,125.195) circle[radius=0.6bp];
\fill[tg808080] (332.898,108.601) circle[radius=0.6bp];
\fill[tg808080] (348.970,110.019) circle[radius=0.6bp];
\fill[tg808080] (326.354,103.775) circle[radius=0.6bp];
\fill[tg808080] (335.689,96.560) circle[radius=0.6bp];
\fill[tg808080] (333.240,114.725) circle[radius=0.6bp];
\fill[tg808080] (345.114,121.637) circle[radius=0.6bp];
\fill[tg808080] (325.789,113.012) circle[radius=0.6bp];
\fill[tg808080] (339.637,115.999) circle[radius=0.6bp];
\fill[tg808080] (325.678,107.773) circle[radius=0.6bp];
\fill[tg808080] (344.880,130.153) circle[radius=0.6bp];
\fill[tg808080] (341.099,104.303) circle[radius=0.6bp];
\fill[tg808080] (335.077,109.281) circle[radius=0.6bp];
\fill[tg808080] (334.360,123.375) circle[radius=0.6bp];
\fill[tg808080] (348.621,119.795) circle[radius=0.6bp];
\fill[tg808080] (335.660,125.569) circle[radius=0.6bp];
\fill[tg808080] (343.257,108.310) circle[radius=0.6bp];
\fill[tg808080] (336.698,115.134) circle[radius=0.6bp];
\fill[tg808080] (338.530,113.506) circle[radius=0.6bp];
\fill[tg808080] (331.317,110.945) circle[radius=0.6bp];
\fill[tg808080] (347.960,124.373) circle[radius=0.6bp];
\fill[tg808080] (332.088,112.021) circle[radius=0.6bp];
\fill[tg808080] (321.053,100.262) circle[radius=0.6bp];
\fill[tg808080] (335.391,112.251) circle[radius=0.6bp];
\fill[tg808080] (351.981,106.317) circle[radius=0.6bp];
\fill[tg808080] (342.532,134.114) circle[radius=0.6bp];
\fill[tg808080] (326.797,110.899) circle[radius=0.6bp];
\fill[tg808080] (346.627,108.517) circle[radius=0.6bp];
\fill[tg808080] (330.350,118.932) circle[radius=0.6bp];
\fill[tg808080] (333.627,116.825) circle[radius=0.6bp];
\fill[tg808080] (334.683,108.671) circle[radius=0.6bp];
\fill[tg808080] (334.789,117.766) circle[radius=0.6bp];
\fill[tg808080] (340.084,121.487) circle[radius=0.6bp];
\fill[tg808080] (333.330,113.245) circle[radius=0.6bp];
\fill[tg808080] (337.729,116.988) circle[radius=0.6bp];
\fill[tg808080] (344.418,118.431) circle[radius=0.6bp];
\fill[tg808080] (332.480,134.284) circle[radius=0.6bp];
\fill[tg808080] (336.969,123.525) circle[radius=0.6bp];
\fill[tg808080] (334.149,115.547) circle[radius=0.6bp];
\fill[tg808080] (321.541,124.104) circle[radius=0.6bp];
\fill[tg808080] (329.426,129.765) circle[radius=0.6bp];
\fill[tg808080] (350.985,120.801) circle[radius=0.6bp];
\fill[tg808080] (339.175,134.545) circle[radius=0.6bp];
\fill[tg808080] (342.369,115.812) circle[radius=0.6bp];
\fill[tg808080] (318.751,120.330) circle[radius=0.6bp];
\fill[tg808080] (337.309,127.981) circle[radius=0.6bp];
\fill[tg808080] (344.834,117.767) circle[radius=0.6bp];
\fill[tg808080] (338.888,120.591) circle[radius=0.6bp];
\fill[tg808080] (337.432,112.472) circle[radius=0.6bp];
\fill[tg808080] (335.474,118.957) circle[radius=0.6bp];
\fill[tg808080] (343.501,130.471) circle[radius=0.6bp];
\fill[tg808080] (345.297,107.609) circle[radius=0.6bp];
\fill[tg808080] (334.784,110.791) circle[radius=0.6bp];
\fill[tg808080] (340.323,124.046) circle[radius=0.6bp];
\fill[tg808080] (334.599,120.629) circle[radius=0.6bp];
\fill[tg808080] (343.224,126.378) circle[radius=0.6bp];
\fill[tg808080] (330.717,110.745) circle[radius=0.6bp];
\fill[tg808080] (332.435,115.880) circle[radius=0.6bp];
\fill[tg808080] (346.059,116.936) circle[radius=0.6bp];
\fill[tg808080] (331.530,122.803) circle[radius=0.6bp];
\fill[tg808080] (342.615,113.675) circle[radius=0.6bp];
\fill[tg808080] (328.314,124.406) circle[radius=0.6bp];
\fill[tg808080] (338.890,116.634) circle[radius=0.6bp];
\fill[tg808080] (334.359,122.818) circle[radius=0.6bp];
\fill[tg808080] (330.933,130.292) circle[radius=0.6bp];
\fill[tg808080] (328.495,114.324) circle[radius=0.6bp];
\fill[tg808080] (335.226,115.276) circle[radius=0.6bp];
\fill[tg808080] (332.115,114.561) circle[radius=0.6bp];
\fill[tg808080] (341.893,121.612) circle[radius=0.6bp];
\fill[tg808080] (348.798,129.189) circle[radius=0.6bp];
\fill[tg808080] (328.228,130.029) circle[radius=0.6bp];
\fill[tg808080] (335.574,113.555) circle[radius=0.6bp];
\fill[tg808080] (340.380,119.987) circle[radius=0.6bp];
\fill[tg808080] (344.306,114.723) circle[radius=0.6bp];
\fill[tg808080] (353.038,116.300) circle[radius=0.6bp];
\fill[tg808080] (340.750,106.771) circle[radius=0.6bp];
\fill[tg808080] (336.404,114.189) circle[radius=0.6bp];
\fill[tg808080] (327.057,115.696) circle[radius=0.6bp];
\fill[tg808080] (335.390,126.129) circle[radius=0.6bp];
\fill[tg808080] (340.759,108.836) circle[radius=0.6bp];
\fill[tg808080] (333.846,113.765) circle[radius=0.6bp];
\fill[tg808080] (341.323,102.837) circle[radius=0.6bp];
\fill[tg808080] (330.589,115.269) circle[radius=0.6bp];
\fill[tg808080] (323.313,115.016) circle[radius=0.6bp];
\fill[tg808080] (339.429,111.874) circle[radius=0.6bp];
\fill[tg808080] (327.797,119.695) circle[radius=0.6bp];
\fill[tg808080] (348.990,100.302) circle[radius=0.6bp];
\fill[tg808080] (338.873,127.840) circle[radius=0.6bp];
\fill[tg808080] (338.409,117.021) circle[radius=0.6bp];
\fill[tg808080] (342.227,116.597) circle[radius=0.6bp];
\fill[tg808080] (336.452,121.389) circle[radius=0.6bp];
\fill[tg808080] (319.190,126.072) circle[radius=0.6bp];
\fill[tg808080] (347.844,127.315) circle[radius=0.6bp];
\fill[tg808080] (352.162,110.490) circle[radius=0.6bp];
\fill[tg808080] (334.583,119.891) circle[radius=0.6bp];
\fill[tg808080] (344.807,122.685) circle[radius=0.6bp];
\fill[tg808080] (338.157,118.603) circle[radius=0.6bp];
\fill[tg808080] (322.722,113.318) circle[radius=0.6bp];
\fill[tg808080] (332.995,125.732) circle[radius=0.6bp];
\fill[tg808080] (346.505,128.355) circle[radius=0.6bp];
\fill[tg808080] (336.149,110.139) circle[radius=0.6bp];
\fill[tg808080] (332.605,118.718) circle[radius=0.6bp];
\fill[tg808080] (331.245,118.091) circle[radius=0.6bp];
\fill[tg808080] (328.116,121.755) circle[radius=0.6bp];
\fill[tg808080] (327.494,121.966) circle[radius=0.6bp];
\fill[tg808080] (350.019,110.009) circle[radius=0.6bp];
\fill[tg808080] (332.136,110.852) circle[radius=0.6bp];
\fill[tg808080] (333.544,119.901) circle[radius=0.6bp];
\fill[tg808080] (337.452,111.121) circle[radius=0.6bp];
\fill[tg808080] (321.056,116.210) circle[radius=0.6bp];
\fill[tg808080] (349.020,118.404) circle[radius=0.6bp];
\fill[tg808080] (339.538,117.879) circle[radius=0.6bp];
\fill[tg808080] (342.999,120.759) circle[radius=0.6bp];
\fill[tg808080] (337.230,121.029) circle[radius=0.6bp];
\fill[tg808080] (295.962,118.936) circle[radius=0.6bp];
\fill[tg808080] (344.316,124.497) circle[radius=0.6bp];
\fill[tg808080] (341.398,118.216) circle[radius=0.6bp];
\fill[tg808080] (301.851,115.948) circle[radius=0.6bp];
\fill[tg808080] (336.432,112.794) circle[radius=0.6bp];
\fill[tg808080] (332.097,117.585) circle[radius=0.6bp];
\fill[tg808080] (330.373,126.625) circle[radius=0.6bp];
\fill[tg808080] (319.575,110.422) circle[radius=0.6bp];
\fill[tg808080] (340.213,131.059) circle[radius=0.6bp];
\fill[tg808080] (346.404,119.750) circle[radius=0.6bp];
\fill[tg808080] (329.898,126.483) circle[radius=0.6bp];
\fill[tg808080] (343.675,114.869) circle[radius=0.6bp];
\fill[tg808080] (338.310,117.432) circle[radius=0.6bp];
\fill[tg808080] (336.908,103.639) circle[radius=0.6bp];
\fill[tg808080] (327.136,117.740) circle[radius=0.6bp];
\fill[tg808080] (281.659,122.982) circle[radius=0.6bp];
\fill[tg808080] (322.247,116.027) circle[radius=0.6bp];
\fill[tg808080] (331.967,117.506) circle[radius=0.6bp];
\fill[tg808080] (338.798,129.791) circle[radius=0.6bp];
\fill[tg808080] (336.513,116.284) circle[radius=0.6bp];
\fill[tg808080] (335.291,105.518) circle[radius=0.6bp];
\fill[tg808080] (337.269,121.329) circle[radius=0.6bp];
\fill[tg808080] (330.460,111.303) circle[radius=0.6bp];
\fill[tg808080] (338.586,117.194) circle[radius=0.6bp];
\fill[tg808080] (332.253,127.923) circle[radius=0.6bp];
\fill[tg808080] (332.795,117.402) circle[radius=0.6bp];
\fill[tg808080] (318.989,122.112) circle[radius=0.6bp];
\fill[tg808080] (329.760,117.968) circle[radius=0.6bp];
\fill[tg808080] (345.383,110.176) circle[radius=0.6bp];
\fill[tg808080] (341.842,115.595) circle[radius=0.6bp];
\fill[tg808080] (335.617,123.051) circle[radius=0.6bp];
\fill[tg808080] (347.337,110.051) circle[radius=0.6bp];
\fill[tg808080] (356.360,118.093) circle[radius=0.6bp];
\fill[tg808080] (351.439,113.845) circle[radius=0.6bp];
\fill[tg808080] (338.401,128.900) circle[radius=0.6bp];
\fill[tg808080] (328.982,109.752) circle[radius=0.6bp];
\fill[tg808080] (331.613,125.360) circle[radius=0.6bp];
\fill[tg808080] (321.713,127.728) circle[radius=0.6bp];
\fill[tg808080] (351.239,103.793) circle[radius=0.6bp];
\fill[tg808080] (310.725,102.172) circle[radius=0.6bp];
\fill[tg808080] (343.592,120.677) circle[radius=0.6bp];
\fill[tg808080] (323.036,110.683) circle[radius=0.6bp];
\fill[tg808080] (336.826,123.097) circle[radius=0.6bp];
\fill[tg808080] (336.744,122.378) circle[radius=0.6bp];
\fill[tg808080] (322.858,116.186) circle[radius=0.6bp];
\fill[tg808080] (331.473,116.257) circle[radius=0.6bp];
\fill[tg808080] (326.967,128.947) circle[radius=0.6bp];
\fill[tg808080] (343.844,104.908) circle[radius=0.6bp];
\fill[tg808080] (372.464,117.352) circle[radius=0.6bp];
\fill[tg808080] (337.823,125.991) circle[radius=0.6bp];
\fill[tg808080] (330.214,118.382) circle[radius=0.6bp];
\fill[tg808080] (330.485,106.561) circle[radius=0.6bp];
\fill[tg808080] (333.019,117.912) circle[radius=0.6bp];
\fill[tg808080] (334.151,128.770) circle[radius=0.6bp];
\fill[tg808080] (337.507,122.740) circle[radius=0.6bp];
\fill[tg808080] (343.840,112.201) circle[radius=0.6bp];
\fill[tg808080] (342.091,121.476) circle[radius=0.6bp];
\fill[tg808080] (346.950,121.365) circle[radius=0.6bp];
\fill[tg808080] (343.741,126.937) circle[radius=0.6bp];
\fill[tg808080] (341.698,123.899) circle[radius=0.6bp];
\fill[tg808080] (338.431,115.010) circle[radius=0.6bp];
\fill[tg808080] (337.678,113.467) circle[radius=0.6bp];
\fill[tg808080] (343.421,118.247) circle[radius=0.6bp];
\fill[tg808080] (340.233,119.088) circle[radius=0.6bp];
\fill[tg808080] (343.391,118.259) circle[radius=0.6bp];
\fill[tg808080] (317.766,124.989) circle[radius=0.6bp];
\fill[tg808080] (307.325,114.407) circle[radius=0.6bp];
\fill[tg808080] (319.543,113.684) circle[radius=0.6bp];
\fill[tg808080] (335.916,97.738) circle[radius=0.6bp];
\fill[tg808080] (338.934,118.365) circle[radius=0.6bp];
\fill[tg808080] (332.913,123.042) circle[radius=0.6bp];
\fill[tg808080] (332.557,126.422) circle[radius=0.6bp];
\fill[tg808080] (331.072,110.697) circle[radius=0.6bp];
\fill[tg808080] (344.811,103.676) circle[radius=0.6bp];
\fill[tg808080] (359.575,115.604) circle[radius=0.6bp];
\fill[tg808080] (345.125,123.934) circle[radius=0.6bp];
\fill[tg808080] (345.847,118.231) circle[radius=0.6bp];
\fill[tg808080] (329.952,114.583) circle[radius=0.6bp];
\fill[tg808080] (327.172,102.492) circle[radius=0.6bp];
\fill[tg808080] (324.434,118.544) circle[radius=0.6bp];
\fill[tg808080] (341.042,133.556) circle[radius=0.6bp];
\fill[tg808080] (338.352,123.075) circle[radius=0.6bp];
\fill[tg808080] (332.299,112.670) circle[radius=0.6bp];
\fill[tg808080] (338.535,117.372) circle[radius=0.6bp];
\fill[tg808080] (350.622,121.973) circle[radius=0.6bp];
\fill[tg808080] (341.229,110.398) circle[radius=0.6bp];
\fill[tg808080] (340.219,111.152) circle[radius=0.6bp];
\fill[tg808080] (339.701,118.736) circle[radius=0.6bp];
\fill[tg808080] (332.540,109.725) circle[radius=0.6bp];
\fill[tg808080] (344.963,118.987) circle[radius=0.6bp];
\fill[tg808080] (341.062,109.289) circle[radius=0.6bp];
\fill[tg808080] (341.783,126.597) circle[radius=0.6bp];
\fill[tg808080] (332.714,129.123) circle[radius=0.6bp];
\fill[tg808080] (338.228,101.491) circle[radius=0.6bp];
\fill[tg808080] (345.504,107.350) circle[radius=0.6bp];
\fill[tg808080] (326.971,113.086) circle[radius=0.6bp];
\fill[tg808080] (344.947,167.353) circle[radius=0.6bp];
\fill[tg808080] (329.939,116.813) circle[radius=0.6bp];
\fill[tg808080] (337.272,113.442) circle[radius=0.6bp];
\fill[tg808080] (327.187,111.871) circle[radius=0.6bp];
\fill[tg808080] (339.341,119.714) circle[radius=0.6bp];
\fill[tg808080] (334.024,126.431) circle[radius=0.6bp];
\fill[tg808080] (340.262,137.020) circle[radius=0.6bp];
\fill[tg808080] (332.872,129.991) circle[radius=0.6bp];
\fill[tg808080] (364.750,114.981) circle[radius=0.6bp];
\fill[tg808080] (333.396,112.268) circle[radius=0.6bp];
\fill[tg808080] (336.368,133.027) circle[radius=0.6bp];
\fill[tg808080] (326.883,121.543) circle[radius=0.6bp];
\fill[tg808080] (340.583,124.313) circle[radius=0.6bp];
\fill[tg808080] (337.926,126.508) circle[radius=0.6bp];
\fill[tg808080] (326.795,123.749) circle[radius=0.6bp];
\fill[tg808080] (346.200,114.388) circle[radius=0.6bp];
\fill[tg808080] (348.513,125.600) circle[radius=0.6bp];
\fill[tg808080] (329.063,113.164) circle[radius=0.6bp];
\fill[tg808080] (328.857,109.351) circle[radius=0.6bp];
\draw[color=tg0000ff,line width=0.7bp,line cap=round] (332.800,116.800) -- (337.200,121.200);
\draw[color=tg0000ff,line width=0.7bp,line cap=round] (332.800,121.200) -- (337.200,116.800);
\node[anchor=center,inner sep=0pt,outer sep=0pt,text=tg18314b,font={\sffamily\fontsize{6.5bp}{7.8bp}\selectfont}] at (285.222,183.000) {-4};
\node[anchor=east,inner sep=0pt,outer sep=0pt,text=tg18314b,font={\sffamily\fontsize{6.5bp}{7.8bp}\selectfont}] at (274.000,168.778) {-4};
\node[anchor=center,inner sep=0pt,outer sep=0pt,text=tg18314b,font={\sffamily\fontsize{6.5bp}{7.8bp}\selectfont}] at (335.000,183.000) {0};
\node[anchor=east,inner sep=0pt,outer sep=0pt,text=tg18314b,font={\sffamily\fontsize{6.5bp}{7.8bp}\selectfont}] at (274.000,119.000) {0};
\node[anchor=center,inner sep=0pt,outer sep=0pt,text=tg18314b,font={\sffamily\fontsize{6.5bp}{7.8bp}\selectfont}] at (384.778,183.000) {4};
\node[anchor=east,inner sep=0pt,outer sep=0pt,text=tg18314b,font={\sffamily\fontsize{6.5bp}{7.8bp}\selectfont}] at (274.000,69.222) {4};
\node[anchor=center,inner sep=0pt,outer sep=0pt,text=tg18314b,font={\sffamily\fontsize{7bp}{8.4bp}\selectfont}] at (335.000,195.000) {Standardized target 0};
\node[anchor=west,inner sep=0pt,outer sep=0pt,text=tg626b73,font={\sffamily\fontsize{6.5bp}{7.8bp}\selectfont}] at (279.000,57.000) {Target 1};
\node[anchor=center,inner sep=0pt,outer sep=0pt,text=tg18314b,font={\sffamily\fontsize{8bp}{9.6bp}\selectfont}] at (335.000,209.000) {Across horizons and targets};
\path[fill=tg67001f,draw=none,line width=0.5bp] (279.000,232.000) rectangle (288.383,241.383);
\path[fill=tgf0f4f6,draw=none,line width=0.5bp] (288.333,232.000) rectangle (297.717,241.383);
\path[fill=tgba2832,draw=none,line width=0.5bp] (297.667,232.000) rectangle (307.050,241.383);
\path[fill=tgf9ebe3,draw=none,line width=0.5bp] (307.000,232.000) rectangle (316.383,241.383);
\path[fill=tgde735c,draw=none,line width=0.5bp] (316.333,232.000) rectangle (325.717,241.383);
\path[fill=tgf9efe9,draw=none,line width=0.5bp] (325.667,232.000) rectangle (335.050,241.383);
\path[fill=tgea8e70,draw=none,line width=0.5bp] (335.000,232.000) rectangle (344.383,241.383);
\path[fill=tgf9eee7,draw=none,line width=0.5bp] (344.333,232.000) rectangle (353.717,241.383);
\path[fill=tgf4a683,draw=none,line width=0.5bp] (353.667,232.000) rectangle (363.050,241.383);
\path[fill=tgf9ebe3,draw=none,line width=0.5bp] (363.000,232.000) rectangle (372.383,241.383);
\path[fill=tgf5ac8b,draw=none,line width=0.5bp] (372.333,232.000) rectangle (381.717,241.383);
\path[fill=tgfaeae1,draw=none,line width=0.5bp] (381.667,232.000) rectangle (391.050,241.383);
\path[fill=tgf0f4f6,draw=none,line width=0.5bp] (279.000,241.333) rectangle (288.383,250.717);
\path[fill=tg67001f,draw=none,line width=0.5bp] (288.333,241.333) rectangle (297.717,250.717);
\path[fill=tgf2f5f6,draw=none,line width=0.5bp] (297.667,241.333) rectangle (307.050,250.717);
\path[fill=tgc2383a,draw=none,line width=0.5bp] (307.000,241.333) rectangle (316.383,250.717);
\path[fill=tgf0f4f6,draw=none,line width=0.5bp] (316.333,241.333) rectangle (325.717,250.717);
\path[fill=tgdf765e,draw=none,line width=0.5bp] (325.667,241.333) rectangle (335.050,250.717);
\path[fill=tgedf2f5,draw=none,line width=0.5bp] (335.000,241.333) rectangle (344.383,250.717);
\path[fill=tgea8e70,draw=none,line width=0.5bp] (344.333,241.333) rectangle (353.717,250.717);
\path[fill=tgf2f5f6,draw=none,line width=0.5bp] (353.667,241.333) rectangle (363.050,250.717);
\path[fill=tgf3a481,draw=none,line width=0.5bp] (363.000,241.333) rectangle (372.383,250.717);
\path[fill=tgf2f5f6,draw=none,line width=0.5bp] (372.333,241.333) rectangle (381.717,250.717);
\path[fill=tgf6af8e,draw=none,line width=0.5bp] (381.667,241.333) rectangle (391.050,250.717);
\path[fill=tgba2832,draw=none,line width=0.5bp] (279.000,250.667) rectangle (288.383,260.050);
\path[fill=tgf2f5f6,draw=none,line width=0.5bp] (288.333,250.667) rectangle (297.717,260.050);
\path[fill=tg67001f,draw=none,line width=0.5bp] (297.667,250.667) rectangle (307.050,260.050);
\path[fill=tgfaeae1,draw=none,line width=0.5bp] (307.000,250.667) rectangle (316.383,260.050);
\path[fill=tgc13639,draw=none,line width=0.5bp] (316.333,250.667) rectangle (325.717,260.050);
\path[fill=tgf9f0eb,draw=none,line width=0.5bp] (325.667,250.667) rectangle (335.050,260.050);
\path[fill=tgd55d4c,draw=none,line width=0.5bp] (335.000,250.667) rectangle (344.383,260.050);
\path[fill=tgf9f0eb,draw=none,line width=0.5bp] (344.333,250.667) rectangle (353.717,260.050);
\path[fill=tge17860,draw=none,line width=0.5bp] (353.667,250.667) rectangle (363.050,260.050);
\path[fill=tgf9f0eb,draw=none,line width=0.5bp] (363.000,250.667) rectangle (372.383,260.050);
\path[fill=tge58368,draw=none,line width=0.5bp] (372.333,250.667) rectangle (381.717,260.050);
\path[fill=tgf9efe9,draw=none,line width=0.5bp] (381.667,250.667) rectangle (391.050,260.050);
\path[fill=tgf9ebe3,draw=none,line width=0.5bp] (279.000,260.000) rectangle (288.383,269.383);
\path[fill=tgc2383a,draw=none,line width=0.5bp] (288.333,260.000) rectangle (297.717,269.383);
\path[fill=tgfaeae1,draw=none,line width=0.5bp] (297.667,260.000) rectangle (307.050,269.383);
\path[fill=tg67001f,draw=none,line width=0.5bp] (307.000,260.000) rectangle (316.383,269.383);
\path[fill=tgf9efe9,draw=none,line width=0.5bp] (316.333,260.000) rectangle (325.717,269.383);
\path[fill=tgba2832,draw=none,line width=0.5bp] (325.667,260.000) rectangle (335.050,269.383);
\path[fill=tgf8f1ed,draw=none,line width=0.5bp] (335.000,260.000) rectangle (344.383,269.383);
\path[fill=tgce4f45,draw=none,line width=0.5bp] (344.333,260.000) rectangle (353.717,269.383);
\path[fill=tgf8f1ed,draw=none,line width=0.5bp] (353.667,260.000) rectangle (363.050,269.383);
\path[fill=tgdb6b55,draw=none,line width=0.5bp] (363.000,260.000) rectangle (372.383,269.383);
\path[fill=tgf8f2ef,draw=none,line width=0.5bp] (372.333,260.000) rectangle (381.717,269.383);
\path[fill=tge17860,draw=none,line width=0.5bp] (381.667,260.000) rectangle (391.050,269.383);
\path[fill=tgde735c,draw=none,line width=0.5bp] (279.000,269.333) rectangle (288.383,278.717);
\path[fill=tgf0f4f6,draw=none,line width=0.5bp] (288.333,269.333) rectangle (297.717,278.717);
\path[fill=tgc13639,draw=none,line width=0.5bp] (297.667,269.333) rectangle (307.050,278.717);
\path[fill=tgf9efe9,draw=none,line width=0.5bp] (307.000,269.333) rectangle (316.383,278.717);
\path[fill=tg67001f,draw=none,line width=0.5bp] (316.333,269.333) rectangle (325.717,278.717);
\path[fill=tgf7f5f4,draw=none,line width=0.5bp] (325.667,269.333) rectangle (335.050,278.717);
\path[fill=tg991027,draw=none,line width=0.5bp] (335.000,269.333) rectangle (344.383,278.717);
\path[fill=tgf7f5f4,draw=none,line width=0.5bp] (344.333,269.333) rectangle (353.717,278.717);
\path[fill=tgbd2d35,draw=none,line width=0.5bp] (353.667,269.333) rectangle (363.050,278.717);
\path[fill=tgf8f4f2,draw=none,line width=0.5bp] (363.000,269.333) rectangle (372.383,278.717);
\path[fill=tgc53e3d,draw=none,line width=0.5bp] (372.333,269.333) rectangle (381.717,278.717);
\path[fill=tgf8f3f0,draw=none,line width=0.5bp] (381.667,269.333) rectangle (391.050,278.717);
\path[fill=tgf9efe9,draw=none,line width=0.5bp] (279.000,278.667) rectangle (288.383,288.050);
\path[fill=tgdf765e,draw=none,line width=0.5bp] (288.333,278.667) rectangle (297.717,288.050);
\path[fill=tgf9f0eb,draw=none,line width=0.5bp] (297.667,278.667) rectangle (307.050,288.050);
\path[fill=tgba2832,draw=none,line width=0.5bp] (307.000,278.667) rectangle (316.383,288.050);
\path[fill=tgf7f5f4,draw=none,line width=0.5bp] (316.333,278.667) rectangle (325.717,288.050);
\path[fill=tg67001f,draw=none,line width=0.5bp] (325.667,278.667) rectangle (335.050,288.050);
\path[fill=tgf6f7f7,draw=none,line width=0.5bp] (335.000,278.667) rectangle (344.383,288.050);
\path[fill=tg991027,draw=none,line width=0.5bp] (344.333,278.667) rectangle (353.717,288.050);
\path[fill=tgf5f6f7,draw=none,line width=0.5bp] (353.667,278.667) rectangle (363.050,288.050);
\path[fill=tgb82531,draw=none,line width=0.5bp] (363.000,278.667) rectangle (372.383,288.050);
\path[fill=tgf6f7f7,draw=none,line width=0.5bp] (372.333,278.667) rectangle (381.717,288.050);
\path[fill=tgc13639,draw=none,line width=0.5bp] (381.667,278.667) rectangle (391.050,288.050);
\path[fill=tgea8e70,draw=none,line width=0.5bp] (279.000,288.000) rectangle (288.383,297.383);
\path[fill=tgedf2f5,draw=none,line width=0.5bp] (288.333,288.000) rectangle (297.717,297.383);
\path[fill=tgd55d4c,draw=none,line width=0.5bp] (297.667,288.000) rectangle (307.050,297.383);
\path[fill=tgf8f1ed,draw=none,line width=0.5bp] (307.000,288.000) rectangle (316.383,297.383);
\path[fill=tg991027,draw=none,line width=0.5bp] (316.333,288.000) rectangle (325.717,297.383);
\path[fill=tgf6f7f7,draw=none,line width=0.5bp] (325.667,288.000) rectangle (335.050,297.383);
\path[fill=tg67001f,draw=none,line width=0.5bp] (335.000,288.000) rectangle (344.383,297.383);
\path[fill=tgf6f7f7,draw=none,line width=0.5bp] (344.333,288.000) rectangle (353.717,297.383);
\path[fill=tg960f27,draw=none,line width=0.5bp] (353.667,288.000) rectangle (363.050,297.383);
\path[fill=tgf7f6f6,draw=none,line width=0.5bp] (363.000,288.000) rectangle (372.383,297.383);
\path[fill=tgab162a,draw=none,line width=0.5bp] (372.333,288.000) rectangle (381.717,297.383);
\path[fill=tgf8f4f2,draw=none,line width=0.5bp] (381.667,288.000) rectangle (391.050,297.383);
\path[fill=tgf9eee7,draw=none,line width=0.5bp] (279.000,297.333) rectangle (288.383,306.717);
\path[fill=tgea8e70,draw=none,line width=0.5bp] (288.333,297.333) rectangle (297.717,306.717);
\path[fill=tgf9f0eb,draw=none,line width=0.5bp] (297.667,297.333) rectangle (307.050,306.717);
\path[fill=tgce4f45,draw=none,line width=0.5bp] (307.000,297.333) rectangle (316.383,306.717);
\path[fill=tgf7f5f4,draw=none,line width=0.5bp] (316.333,297.333) rectangle (325.717,306.717);
\path[fill=tg991027,draw=none,line width=0.5bp] (325.667,297.333) rectangle (335.050,306.717);
\path[fill=tgf6f7f7,draw=none,line width=0.5bp] (335.000,297.333) rectangle (344.383,306.717);
\path[fill=tg67001f,draw=none,line width=0.5bp] (344.333,297.333) rectangle (353.717,306.717);
\path[fill=tgf5f6f7,draw=none,line width=0.5bp] (353.667,297.333) rectangle (363.050,306.717);
\path[fill=tg930e26,draw=none,line width=0.5bp] (363.000,297.333) rectangle (372.383,306.717);
\path[fill=tgf5f6f7,draw=none,line width=0.5bp] (372.333,297.333) rectangle (381.717,306.717);
\path[fill=tga51429,draw=none,line width=0.5bp] (381.667,297.333) rectangle (391.050,306.717);
\path[fill=tgf4a683,draw=none,line width=0.5bp] (279.000,306.667) rectangle (288.383,316.050);
\path[fill=tgf2f5f6,draw=none,line width=0.5bp] (288.333,306.667) rectangle (297.717,316.050);
\path[fill=tge17860,draw=none,line width=0.5bp] (297.667,306.667) rectangle (307.050,316.050);
\path[fill=tgf8f1ed,draw=none,line width=0.5bp] (307.000,306.667) rectangle (316.383,316.050);
\path[fill=tgbd2d35,draw=none,line width=0.5bp] (316.333,306.667) rectangle (325.717,316.050);
\path[fill=tgf5f6f7,draw=none,line width=0.5bp] (325.667,306.667) rectangle (335.050,316.050);
\path[fill=tg960f27,draw=none,line width=0.5bp] (335.000,306.667) rectangle (344.383,316.050);
\path[fill=tgf5f6f7,draw=none,line width=0.5bp] (344.333,306.667) rectangle (353.717,316.050);
\path[fill=tg67001f,draw=none,line width=0.5bp] (353.667,306.667) rectangle (363.050,316.050);
\path[fill=tgf3f5f6,draw=none,line width=0.5bp] (363.000,306.667) rectangle (372.383,316.050);
\path[fill=tg7c0722,draw=none,line width=0.5bp] (372.333,306.667) rectangle (381.717,316.050);
\path[fill=tgf5f6f7,draw=none,line width=0.5bp] (381.667,306.667) rectangle (391.050,316.050);
\path[fill=tgf9ebe3,draw=none,line width=0.5bp] (279.000,316.000) rectangle (288.383,325.383);
\path[fill=tgf3a481,draw=none,line width=0.5bp] (288.333,316.000) rectangle (297.717,325.383);
\path[fill=tgf9f0eb,draw=none,line width=0.5bp] (297.667,316.000) rectangle (307.050,325.383);
\path[fill=tgdb6b55,draw=none,line width=0.5bp] (307.000,316.000) rectangle (316.383,325.383);
\path[fill=tgf8f4f2,draw=none,line width=0.5bp] (316.333,316.000) rectangle (325.717,325.383);
\path[fill=tgb82531,draw=none,line width=0.5bp] (325.667,316.000) rectangle (335.050,325.383);
\path[fill=tgf7f6f6,draw=none,line width=0.5bp] (335.000,316.000) rectangle (344.383,325.383);
\path[fill=tg930e26,draw=none,line width=0.5bp] (344.333,316.000) rectangle (353.717,325.383);
\path[fill=tgf3f5f6,draw=none,line width=0.5bp] (353.667,316.000) rectangle (363.050,325.383);
\path[fill=tg67001f,draw=none,line width=0.5bp] (363.000,316.000) rectangle (372.383,325.383);
\path[fill=tgf5f6f7,draw=none,line width=0.5bp] (372.333,316.000) rectangle (381.717,325.383);
\path[fill=tg790622,draw=none,line width=0.5bp] (381.667,316.000) rectangle (391.050,325.383);
\path[fill=tgf5ac8b,draw=none,line width=0.5bp] (279.000,325.333) rectangle (288.383,334.717);
\path[fill=tgf2f5f6,draw=none,line width=0.5bp] (288.333,325.333) rectangle (297.717,334.717);
\path[fill=tge58368,draw=none,line width=0.5bp] (297.667,325.333) rectangle (307.050,334.717);
\path[fill=tgf8f2ef,draw=none,line width=0.5bp] (307.000,325.333) rectangle (316.383,334.717);
\path[fill=tgc53e3d,draw=none,line width=0.5bp] (316.333,325.333) rectangle (325.717,334.717);
\path[fill=tgf6f7f7,draw=none,line width=0.5bp] (325.667,325.333) rectangle (335.050,334.717);
\path[fill=tgab162a,draw=none,line width=0.5bp] (335.000,325.333) rectangle (344.383,334.717);
\path[fill=tgf5f6f7,draw=none,line width=0.5bp] (344.333,325.333) rectangle (353.717,334.717);
\path[fill=tg7c0722,draw=none,line width=0.5bp] (353.667,325.333) rectangle (363.050,334.717);
\path[fill=tgf5f6f7,draw=none,line width=0.5bp] (363.000,325.333) rectangle (372.383,334.717);
\path[fill=tg67001f,draw=none,line width=0.5bp] (372.333,325.333) rectangle (381.717,334.717);
\path[fill=tgf6f7f7,draw=none,line width=0.5bp] (381.667,325.333) rectangle (391.050,334.717);
\path[fill=tgfaeae1,draw=none,line width=0.5bp] (279.000,334.667) rectangle (288.383,344.050);
\path[fill=tgf6af8e,draw=none,line width=0.5bp] (288.333,334.667) rectangle (297.717,344.050);
\path[fill=tgf9efe9,draw=none,line width=0.5bp] (297.667,334.667) rectangle (307.050,344.050);
\path[fill=tge17860,draw=none,line width=0.5bp] (307.000,334.667) rectangle (316.383,344.050);
\path[fill=tgf8f3f0,draw=none,line width=0.5bp] (316.333,334.667) rectangle (325.717,344.050);
\path[fill=tgc13639,draw=none,line width=0.5bp] (325.667,334.667) rectangle (335.050,344.050);
\path[fill=tgf8f4f2,draw=none,line width=0.5bp] (335.000,334.667) rectangle (344.383,344.050);
\path[fill=tga51429,draw=none,line width=0.5bp] (344.333,334.667) rectangle (353.717,344.050);
\path[fill=tgf5f6f7,draw=none,line width=0.5bp] (353.667,334.667) rectangle (363.050,344.050);
\path[fill=tg790622,draw=none,line width=0.5bp] (363.000,334.667) rectangle (372.383,344.050);
\path[fill=tgf6f7f7,draw=none,line width=0.5bp] (372.333,334.667) rectangle (381.717,344.050);
\path[fill=tg67001f,draw=none,line width=0.5bp] (381.667,334.667) rectangle (391.050,344.050);
\node[anchor=east,inner sep=0pt,outer sep=0pt,text=tg18314b,font={\sffamily\fontsize{6bp}{7.2bp}\selectfont}] at (275.000,236.667) {h1:y0};
\node[anchor=center,inner sep=0pt,outer sep=0pt,text=tg18314b,font={\sffamily\fontsize{6bp}{7.2bp}\selectfont}] at (283.667,352.000) {h1:y0};
\node[anchor=east,inner sep=0pt,outer sep=0pt,text=tg18314b,font={\sffamily\fontsize{6bp}{7.2bp}\selectfont}] at (275.000,274.000) {h3:y0};
\node[anchor=center,inner sep=0pt,outer sep=0pt,text=tg18314b,font={\sffamily\fontsize{6bp}{7.2bp}\selectfont}] at (321.000,352.000) {h3:y0};
\node[anchor=east,inner sep=0pt,outer sep=0pt,text=tg18314b,font={\sffamily\fontsize{6bp}{7.2bp}\selectfont}] at (275.000,330.000) {h6:y0};
\node[anchor=center,inner sep=0pt,outer sep=0pt,text=tg18314b,font={\sffamily\fontsize{6bp}{7.2bp}\selectfont}] at (377.000,352.000) {h6:y0};
\node[anchor=center,inner sep=0pt,outer sep=0pt,text=tg18314b,font={\sffamily\fontsize{8bp}{9.6bp}\selectfont\bfseries}] at (335.000,369.000) {Aggregate SD: 0.67};
\node[anchor=center,inner sep=0pt,outer sep=0pt,text=tg18314b,font={\sffamily\fontsize{7bp}{8.4bp}\selectfont}] at (224.000,223.000) {Correlation coordinates: (h1,y0), (h1,y1), (h2,y0), ...};
\path[fill=tg67001f,draw=none,line width=0.5bp] (421.000,232.000) rectangle (427.000,233.450);
\path[fill=tg700320,draw=none,line width=0.5bp] (421.000,233.400) rectangle (427.000,234.850);
\path[fill=tg790622,draw=none,line width=0.5bp] (421.000,234.800) rectangle (427.000,236.250);
\path[fill=tg810823,draw=none,line width=0.5bp] (421.000,236.200) rectangle (427.000,237.650);
\path[fill=tg8a0b25,draw=none,line width=0.5bp] (421.000,237.600) rectangle (427.000,239.050);
\path[fill=tg960f27,draw=none,line width=0.5bp] (421.000,239.000) rectangle (427.000,240.450);
\path[fill=tg9f1228,draw=none,line width=0.5bp] (421.000,240.400) rectangle (427.000,241.850);
\path[fill=tga81529,draw=none,line width=0.5bp] (421.000,241.800) rectangle (427.000,243.250);
\path[fill=tgb1182b,draw=none,line width=0.5bp] (421.000,243.200) rectangle (427.000,244.650);
\path[fill=tgb72230,draw=none,line width=0.5bp] (421.000,244.600) rectangle (427.000,246.050);
\path[fill=tgbb2a34,draw=none,line width=0.5bp] (421.000,246.000) rectangle (427.000,247.450);
\path[fill=tgbf3338,draw=none,line width=0.5bp] (421.000,247.400) rectangle (427.000,248.850);
\path[fill=tgc43b3c,draw=none,line width=0.5bp] (421.000,248.800) rectangle (427.000,250.250);
\path[fill=tgc94741,draw=none,line width=0.5bp] (421.000,250.200) rectangle (427.000,251.650);
\path[fill=tgce4f45,draw=none,line width=0.5bp] (421.000,251.600) rectangle (427.000,253.050);
\path[fill=tgd25849,draw=none,line width=0.5bp] (421.000,253.000) rectangle (427.000,254.450);
\path[fill=tgd6604d,draw=none,line width=0.5bp] (421.000,254.400) rectangle (427.000,255.850);
\path[fill=tgdb6b55,draw=none,line width=0.5bp] (421.000,255.800) rectangle (427.000,257.250);
\path[fill=tgde735c,draw=none,line width=0.5bp] (421.000,257.200) rectangle (427.000,258.650);
\path[fill=tge27b62,draw=none,line width=0.5bp] (421.000,258.600) rectangle (427.000,260.050);
\path[fill=tge58368,draw=none,line width=0.5bp] (421.000,260.000) rectangle (427.000,261.450);
\path[fill=tgea8e70,draw=none,line width=0.5bp] (421.000,261.400) rectangle (427.000,262.850);
\path[fill=tgee9677,draw=none,line width=0.5bp] (421.000,262.800) rectangle (427.000,264.250);
\path[fill=tgf19e7d,draw=none,line width=0.5bp] (421.000,264.200) rectangle (427.000,265.650);
\path[fill=tgf4a683,draw=none,line width=0.5bp] (421.000,265.600) rectangle (427.000,267.050);
\path[fill=tgf6af8e,draw=none,line width=0.5bp] (421.000,267.000) rectangle (427.000,268.450);
\path[fill=tgf7b596,draw=none,line width=0.5bp] (421.000,268.400) rectangle (427.000,269.850);
\path[fill=tgf8bb9e,draw=none,line width=0.5bp] (421.000,269.800) rectangle (427.000,271.250);
\path[fill=tgf9c2a7,draw=none,line width=0.5bp] (421.000,271.200) rectangle (427.000,272.650);
\path[fill=tgfac8af,draw=none,line width=0.5bp] (421.000,272.600) rectangle (427.000,274.050);
\path[fill=tgfbd0b9,draw=none,line width=0.5bp] (421.000,274.000) rectangle (427.000,275.450);
\path[fill=tgfcd7c2,draw=none,line width=0.5bp] (421.000,275.400) rectangle (427.000,276.850);
\path[fill=tgfddcc9,draw=none,line width=0.5bp] (421.000,276.800) rectangle (427.000,278.250);
\path[fill=tgfcdfcf,draw=none,line width=0.5bp] (421.000,278.200) rectangle (427.000,279.650);
\path[fill=tgfbe4d6,draw=none,line width=0.5bp] (421.000,279.600) rectangle (427.000,281.050);
\path[fill=tgfae7dc,draw=none,line width=0.5bp] (421.000,281.000) rectangle (427.000,282.450);
\path[fill=tgfaeae1,draw=none,line width=0.5bp] (421.000,282.400) rectangle (427.000,283.850);
\path[fill=tgf9eee7,draw=none,line width=0.5bp] (421.000,283.800) rectangle (427.000,285.250);
\path[fill=tgf8f2ef,draw=none,line width=0.5bp] (421.000,285.200) rectangle (427.000,286.650);
\path[fill=tgf7f5f4,draw=none,line width=0.5bp] (421.000,286.600) rectangle (427.000,288.050);
\path[fill=tgf5f6f7,draw=none,line width=0.5bp] (421.000,288.000) rectangle (427.000,289.450);
\path[fill=tgf0f4f6,draw=none,line width=0.5bp] (421.000,289.400) rectangle (427.000,290.850);
\path[fill=tgeaf1f5,draw=none,line width=0.5bp] (421.000,290.800) rectangle (427.000,292.250);
\path[fill=tge6eff4,draw=none,line width=0.5bp] (421.000,292.200) rectangle (427.000,293.650);
\path[fill=tge1edf3,draw=none,line width=0.5bp] (421.000,293.600) rectangle (427.000,295.050);
\path[fill=tgddebf2,draw=none,line width=0.5bp] (421.000,295.000) rectangle (427.000,296.450);
\path[fill=tgd7e8f1,draw=none,line width=0.5bp] (421.000,296.400) rectangle (427.000,297.850);
\path[fill=tgd2e6f0,draw=none,line width=0.5bp] (421.000,297.800) rectangle (427.000,299.250);
\path[fill=tgcce2ef,draw=none,line width=0.5bp] (421.000,299.200) rectangle (427.000,300.650);
\path[fill=tgc5dfec,draw=none,line width=0.5bp] (421.000,300.600) rectangle (427.000,302.050);
\path[fill=tgbbdaea,draw=none,line width=0.5bp] (421.000,302.000) rectangle (427.000,303.450);
\path[fill=tgb3d6e8,draw=none,line width=0.5bp] (421.000,303.400) rectangle (427.000,304.850);
\path[fill=tgacd2e5,draw=none,line width=0.5bp] (421.000,304.800) rectangle (427.000,306.250);
\path[fill=tga5cee3,draw=none,line width=0.5bp] (421.000,306.200) rectangle (427.000,307.650);
\path[fill=tg9dcbe1,draw=none,line width=0.5bp] (421.000,307.600) rectangle (427.000,309.050);
\path[fill=tg93c6de,draw=none,line width=0.5bp] (421.000,309.000) rectangle (427.000,310.450);
\path[fill=tg8ac0db,draw=none,line width=0.5bp] (421.000,310.400) rectangle (427.000,311.850);
\path[fill=tg81bad8,draw=none,line width=0.5bp] (421.000,311.800) rectangle (427.000,313.250);
\path[fill=tg78b4d5,draw=none,line width=0.5bp] (421.000,313.200) rectangle (427.000,314.650);
\path[fill=tg6bacd1,draw=none,line width=0.5bp] (421.000,314.600) rectangle (427.000,316.050);
\path[fill=tg62a7ce,draw=none,line width=0.5bp] (421.000,316.000) rectangle (427.000,317.450);
\path[fill=tg59a1ca,draw=none,line width=0.5bp] (421.000,317.400) rectangle (427.000,318.850);
\path[fill=tg4f9bc7,draw=none,line width=0.5bp] (421.000,318.800) rectangle (427.000,320.250);
\path[fill=tg4393c3,draw=none,line width=0.5bp] (421.000,320.200) rectangle (427.000,321.650);
\path[fill=tg3f8ec0,draw=none,line width=0.5bp] (421.000,321.600) rectangle (427.000,323.050);
\path[fill=tg3b88be,draw=none,line width=0.5bp] (421.000,323.000) rectangle (427.000,324.450);
\path[fill=tg3783bb,draw=none,line width=0.5bp] (421.000,324.400) rectangle (427.000,325.850);
\path[fill=tg327cb7,draw=none,line width=0.5bp] (421.000,325.800) rectangle (427.000,327.250);
\path[fill=tg2e77b5,draw=none,line width=0.5bp] (421.000,327.200) rectangle (427.000,328.650);
\path[fill=tg2a71b2,draw=none,line width=0.5bp] (421.000,328.600) rectangle (427.000,330.050);
\path[fill=tg266caf,draw=none,line width=0.5bp] (421.000,330.000) rectangle (427.000,331.450);
\path[fill=tg2065ab,draw=none,line width=0.5bp] (421.000,331.400) rectangle (427.000,332.850);
\path[fill=tg1d5fa2,draw=none,line width=0.5bp] (421.000,332.800) rectangle (427.000,334.250);
\path[fill=tg1a5899,draw=none,line width=0.5bp] (421.000,334.200) rectangle (427.000,335.650);
\path[fill=tg175290,draw=none,line width=0.5bp] (421.000,335.600) rectangle (427.000,337.050);
\path[fill=tg124984,draw=none,line width=0.5bp] (421.000,337.000) rectangle (427.000,338.450);
\path[fill=tg0f437b,draw=none,line width=0.5bp] (421.000,338.400) rectangle (427.000,339.850);
\path[fill=tg0c3d73,draw=none,line width=0.5bp] (421.000,339.800) rectangle (427.000,341.250);
\path[fill=tg08366a,draw=none,line width=0.5bp] (421.000,341.200) rectangle (427.000,342.650);
\path[fill=tg053061,draw=none,line width=0.5bp] (421.000,342.600) rectangle (427.000,344.050);
\node[anchor=west,inner sep=0pt,outer sep=0pt,text=tg18314b,font={\sffamily\fontsize{6.2bp}{7.44bp}\selectfont}] at (431.000,232.000) {+1};
\node[anchor=west,inner sep=0pt,outer sep=0pt,text=tg18314b,font={\sffamily\fontsize{6.2bp}{7.44bp}\selectfont}] at (431.000,288.000) {0};
\node[anchor=west,inner sep=0pt,outer sep=0pt,text=tg18314b,font={\sffamily\fontsize{6.2bp}{7.44bp}\selectfont}] at (431.000,344.000) {-1};
\node[anchor=center,inner sep=0pt,outer sep=0pt,text=tg18314b,font={\sffamily\fontsize{7bp}{8.4bp}\selectfont}] at (224.000,390.000) {Each series keeps its values and temporal paths; only cross-series pairing changes.};
\node[anchor=center,inner sep=0pt,outer sep=0pt,text=tg18314b,font={\sffamily\fontsize{7bp}{8.4bp}\selectfont}] at (224.000,405.000) {Aggregate = average standardized error at horizon 4. Marginal CRPS is unchanged.};
\end{tikzpicture}%
}
\endgroup

% END VECTOR FIGURE paper-joint-stresses
\caption{\textbf{Preserve each series; break their joint future.}
A two-target Merton log-price reference supplies 768 joint six-step futures
for one fixed case from the joint financial uncertainty tutorial. The right column re-pairs target 1's
complete trajectories across sample identities. The top panels show the
first 384 draws at horizon four, standardized by reference means and standard
deviations, within a fixed $\pm4.5$ display window. Blue crosses mark the
unchanged conditional means. Correlation maps and aggregate standard
deviations use all 768 draws; matrix coordinates alternate targets within
each horizon. The aggregate is the equal-weight average of
the two standardized forecast errors, not a portfolio-return estimate.
Every empirical per-target marginal and temporal path is preserved; marginal
CRPS is unchanged up to floating-point roundoff. These reference ensembles
illustrate a controlled dependence fault.}
\label{fig:joint-stresses}
\end{figure}

\paragraph{Limits of these extensions.}
Jump-regime filtering is approximate and cannot calibrate the exact
covariance rank check. CCC-GARCH cross-covariance integration has Monte
Carlo error, explicitly separate from its exact conditional simulator.
The joint extensions implement shared jump clocks and shared hidden chains,
not arbitrary idiosyncratic-jump, cross-lag, or dynamic-correlation
systems.


\end{document}
